\beginnumbering
\eledsection*{Adhyāya 32}
\pstart \ledsidenote{\passref{1.32.1}}\skt{अथातः स्वभावविप्रतिपत्तिं }\edtext{\skt{व्याख्यास्यामः}}{\Afootnote{H; \emph{om.}\ K; \skt{अध्यायं व्याख्यास्यामः ॥}\ A1938 N}}\skt{ ॥}\pend

\pstart \ledsidenote{\passref{1.32.3}}\edtext{\skt{स्वभावप्रसिद्धानां}}{\Afootnote{K N H; \skt{स्वभावसिद्धानां}\ A1938}}\skt{ }\edtext{\edtext{\skt{शरीरैकदेशानामन्यथात्वं}}{\lemma{\skt{शरीरैकदेशानाम्}}\Afootnote{A1938 K H; \skt{शरीरैकदेशनाम्}\ N}}}{\lemma{\skt{अन्यथात्वं}}\Afootnote{K H; \skt{अन्यभावित्वं}\ A1938; \skt{मन्यथात्वम्}\ N}}\skt{ मरणाय । तद्यथा शुक्लानां }\edtext{\skt{कृष्णत्वं}}{\Afootnote{A1938 K H; \skt{कृष्णानां}\ N}}\skt{ }\edtext{\skt{कृष्णानां}}{\Afootnote{A1938 K H; \skt{कृष्णाना}\ N}}\skt{ }\edtext{\skt{शुक्लत्वं}}{\Afootnote{K N H; \skt{शुक्लता, रक्तानामन्यवर्णत्वं,}\ A1938}}\skt{ । }\edtext{\skt{स्थिराणाम्मृदुत्वं}}{\lemma{\skt{मृदुत्वं}}\Afootnote{A1938 K\textsuperscript{pc} N H; \skt{मृ}\ K\textsuperscript{ac}}}\skt{ । }\edtext{\edtext{\edtext{\skt{चलानामचलत्वमचलानाञ्चलत्वं}}{\lemma{\skt{चलानाम्}}\Afootnote{K\textsuperscript{pc} N H; \skt{मृदूनां स्थिरता, चलाना}\ A1938; \skt{त्वं चलानाम्}\ K\textsuperscript{ac}}}}{\lemma{\skt{अचलत्वम्}}\Afootnote{K H; \skt{मचलत्वं,}\ A1938; \skt{चलत्वम्}\ N}}}{\lemma{\skt{चलत्वं}}\Afootnote{K N H; \skt{चलता,}\ A1938}}\skt{ । पृथूनां }\edtext{\skt{संक्षिप्तत्वं}}{\Afootnote{A1938 K\textsuperscript{pc} N H; \skt{संप्तत्वं}\ K\textsuperscript{ac}}}\skt{ संक्षिप्तानां }\edtext{\skt{पृथुत्वं}}{\Afootnote{K N H; \skt{पृथुता,}\ A1938}}\skt{ । दीर्घाणां }\edtext{\skt{ह्रस्वत्वं च}}{\Afootnote{H; \skt{ह्रस्वत्वं,}\ A1938; \skt{ह्रस्वं}\ K; \skt{ह्रस्वं वा}\ N}}\skt{ ह्रस्वानां दीर्घता । अपतनधर्म्मिणां }\edtext{\edtext{\skt{पतनमकस्माच्छैत्योष्णस्नेहरौक्ष्यप्रस्तम्भवैवर्ण्ण्यानि}}{\lemma{\skt{पतनम्}}\Afootnote{K N H; \skt{पतनधर्मित्वं , पतनधर्मिणामपतन}\ A1938}}}{\lemma{\skt{अकस्माच्छैत्योष्णस्नेहरौक्ष्यप्रस्तम्भवैवर्ण्ण्यानि}}\Afootnote{K N H; \skt{धर्मित्वमकस्मात्, शैत्यौष्ण्यस्नैग्ध्यरौक्ष्यप्रस्तम्भवैवर्ण्या}\ A1938}}\skt{ }\edtext{\edtext{\skt{अपसर्प्पणमङ्गानां}}{\lemma{\skt{अपसर्प्पणम्}}\Afootnote{K N H; \skt{वसादनं}\ A1938}}}{\lemma{\skt{अङ्गानां}}\Afootnote{K N H; \skt{चाङ्गानाम् ॥३॥}\ A1938}}\pend

\pstart \ledsidenote{\passref{1.32.4}}\edtext{\skt{स्वेभ्यः}}{\Afootnote{A1938 K\textsuperscript{pc} N H; \emph{om.}\ K\textsuperscript{ac}}}\skt{ }\edtext{\skt{स्थानेभ्यः}}{\Afootnote{A1938 K\textsuperscript{pc} N H; \skt{स्वेभ्यः}\ K\textsuperscript{ac}}}\skt{ }\edtext{\edtext{\skt{शरीरैकदेशानामवस्रंसोत्क्षिप्तभ्रान्तपतितविमुक्तनिर्ग्गमातिगमगुरुलघुत्वानि}}{\lemma{\skt{शरीरैकदेशानाम्}}\Afootnote{K N H; \skt{शरीरैकदेशाना}\ A1938}}}{\lemma{\skt{अवस्रंसोत्क्षिप्तभ्रान्तपतितविमुक्तनिर्ग्गमातिगमगुरुलघुत्वानि}}\Afootnote{H; \skt{मवस्त्रस्तोत्क्षिप्तभ्रान्तावक्षिप्तपतितविमुक्तनिर्गतान्तर्गतगुरुलघुत्वानि,}\ A1938; \skt{मवस्रंसोत्क्षिप्तभ्रान्तपतितविमुक्तनिर्गमातिगमगुरुलघुत्वानि}\ K; \skt{अवस्रंसोत्क्षिप्तभ्रान्तपतितविमुक्तानिर्गमातिगमगुरुलघुत्वानि}\ N}}\skt{ प्रवालवर्ण्णव्यङ्गप्रादुर्ब्भावो }\edtext{\skt{वाकस्मात्सिराणां च}}{\Afootnote{K N H; \skt{वाऽप्यकस्मात्, सिराणां च}\ A1938}}\skt{ }\edtext{\skt{दर्शन}}{\Afootnote{H; \skt{दर्शनं}\ A1938 K N}}\skt{ ललाटे नासावंशे वा }\edtext{\edtext{\edtext{\skt{पिटकोत्पत्तिरुदकोत्पत्तिर्न्नेत्ररोगम्विनास्रुप्रवृत्तिर्ललाटे}}{\lemma{\skt{पिटकोत्पत्तिर्}}\Afootnote{K\textsuperscript{pc} N H; \skt{पिड}\ A1938; \skt{पि}\ K\textsuperscript{ac}}}}{\lemma{\skt{उदकोत्पत्तिर्}}\Afootnote{K N H; \skt{कोत्प}\ A1938}}\skt{ वा प्रभातकाले}}{\lemma{\skt{त्तिर्… विनास्रुप्रवृत्तिर्}}\Afootnote{\emph{om.}\ A1938}}\skt{ }\edtext{\skt{स्वेदप्रवृत्तिः}}{\Afootnote{K N H; \skt{स्वेदः, नेत्ररोगाद्विना वाऽश्रुप्रवृत्तिः,}\ A1938}}\skt{ । }\edtext{\skt{गोमयचूर्ण्णप्रकाशस्य}}{\Afootnote{N H; \skt{गोमयचूर्णप्रकाशस्य वा}\ A1938; \skt{गोमयचूर्ण्णप्रकासस्य}\ K}}\skt{ रजसो }\edtext{\edtext{\skt{दर्शनमुत्तमाङ्गे}}{\lemma{\skt{दर्शनम्}}\Afootnote{K N H; \skt{दर्शन}\ A1938}}}{\lemma{\skt{उत्तमाङ्गे}}\Afootnote{K N H; \skt{मुत्तमाङ्गे}\ A1938}}\skt{ }\edtext{\skt{लीयनम्वा}}{\Afootnote{K N H; \skt{निलयनं वा}\ A1938}}\skt{ }\edtext{\skt{कपोतकङ्कगृद्ध्रप्रभृतीनां}}{\Afootnote{K N H; \skt{कपोतकङ्ककाकप्रभृतीनां,}\ A1938}}\skt{ }\edtext{\edtext{\skt{मूत्रपुरीषप्रवृद्धिरभुञ्जनानां}}{\lemma{\skt{मूत्रपुरीषप्रवृद्धिर्}}\Afootnote{K N H; \skt{मूत्रपुरीषवृद्धि}\ A1938}}}{\lemma{\skt{अभुञ्जनानां}}\Afootnote{K N H; \skt{रभुञ्जानानां,}\ A1938}}\skt{ }\edtext{\skt{स्तनमूलहृदयोरस्सु च}}{\Afootnote{K H; \skt{तत्प्रणाशो वा भुञ्जानानां, स्तनमूलहृदयोरःसु च}\ A1938; \skt{तनमूलहृदयोः सूच}\ N}}\skt{ शूलोत्पत्तयः । मध्ये }\edtext{\edtext{\skt{शूनत्वमन्तेषु}}{\lemma{\skt{शूनत्वम्}}\Afootnote{K N H; \skt{शूनत्व}\ A1938}}}{\lemma{\skt{अन्तेषु}}\Afootnote{K N H; \skt{मन्तेषु}\ A1938}}\skt{ }\edtext{\skt{परिम्लायित्वं}}{\Afootnote{A1938 K H; \skt{परिम्ला}\ N}}\skt{ । }\edtext{\skt{विपर्ययो}}{\Afootnote{K H; \skt{विपर्ययो वा,}\ A1938; \emph{om.}\ N}}\skt{ }\edtext{\skt{वा}}{\Afootnote{K H; \skt{तथाऽर्धाङ्गे श्वयथुः शोषोऽङ्गपक्षयोर्वा,}\ A1938; \skt{नयित्वम् । विपर्ययो वा}\ N}}\skt{ }\edtext{\skt{नष्टहीनविकलविहृतस्वरं वा}}{\Afootnote{K\textsuperscript{pc} N H; \skt{नष्टहीनविकलविकृतस्वरता वा,}\ A1938; \skt{नष्टहीनविकलविहृतस्वरं}\ K\textsuperscript{ac}}}\skt{ }\edtext{\skt{विवर्ण्णपुष्पप्रादुर्ब्भावो वा}}{\Afootnote{A1938 K\textsuperscript{pc} H; \skt{विवर्ण्णपुष्पप्रादुर्भावो}\ N}}\skt{ }\edtext{\edtext{\skt{दन्तनखशरीरेषु}}{\Afootnote{K N H; \skt{दन्तमुखनखशरीरेषु,}\ A1938}}\skt{ }\edtext{\skt{पुष्पदर्शनं}}{\Afootnote{H; \emph{om.}\ A1938 K\textsuperscript{ac} N; \skt{पुष्पदर्शन}\ K\textsuperscript{pc}}}\skt{ यस्य }\edtext{\skt{चाप्सु}}{\Afootnote{K N H; \skt{वाऽप्सु}\ A1938}}\skt{ }\edtext{\skt{कफशकृद्रेतांसि}}{\Afootnote{K N H; \skt{कफपुरीषरेतांसि}\ A1938}}}{\lemma{\skt{वा… वा}}\Afootnote{\emph{om.}\ K\textsuperscript{ac}}}\skt{ निमज्जन्ति }\edtext{\skt{यस्य च}}{\Afootnote{K\textsuperscript{pc} N H; \skt{यस्य वा}\ A1938; \skt{यस्य च न}\ K\textsuperscript{ac}}}\skt{ दृष्टिमण्डले भिन्नविकृतानि }\edtext{\edtext{\skt{रूपाण्यालोच्यन्ते}}{\lemma{\skt{रूपाण्य्}}\Afootnote{K N H; \skt{रूपा}\ A1938}}}{\lemma{\skt{आलोच्यन्ते}}\Afootnote{K N H; \skt{ण्यालोक्यन्ते,}\ A1938}}\skt{ स्नेहाभ्यक्तकेशाङ्ग इव यो भाति यश्च दुर्ब्बलो भक्तद्वेषातिसाराभ्याम्पीड्यते । कासमानश्च }\edtext{\skt{तृष्णातिभूतः}}{\Afootnote{K N H; \skt{तृष्णाभिभूतः,}\ A1938}}\skt{ }\edtext{\edtext{\skt{क्षीणश्छर्द्दिभक्तदेशयुक्तः}}{\lemma{\skt{क्षीणश्}}\Afootnote{K N H; \skt{क्षीण}\ A1938}}}{\lemma{\skt{छर्द्दिभक्तदेशयुक्तः}}\Afootnote{H; \skt{श्छर्दिभक्तद्वेषयुक्तः}\ A1938; \skt{छर्द्दिभक्तद्वेषयुक्तः}\ K N}}\skt{ }\edtext{\skt{सफेनरुधिरोद्वामी}}{\Afootnote{K N H; \skt{सफेनपूयरुधिरोद्वामी}\ A1938}}\skt{ हतस्वरः }\edtext{\skt{शूलाभिहतश्च}}{\Afootnote{K N H; \skt{शूलाभिपन्नश्च}\ A1938}}\skt{ }\edtext{\skt{मनुष्यशूनकरचरणो}}{\Afootnote{K H; \skt{मनुष्यः, शूनकरचरणवदनःक्षीणो}\ A1938; \skt{मनुष्यशूनकरचरणोन्}\ N}}\skt{ }\edtext{\skt{न्नद्वेषी}}{\Afootnote{A1938 K H; \skt{नद्वेषी}\ N}}\skt{ }\edtext{\skt{स्रस्तपिण्डकांसपाणिपादः}}{\Afootnote{K H; \skt{स्रस्तपिण्डिकांसपाणिपादो}\ A1938; \skt{शसुपिण्डिकान्सपाणिपादः ।}\ N}}\skt{ । }\edtext{\skt{यश्च}}{\Afootnote{K N H; \skt{ज्वरकासाभिभूतः, यस्तु}\ A1938}}\skt{ पूर्वाह्णे }\edtext{\edtext{\skt{भुक्तमपराह्णे}}{\lemma{\skt{भुक्तम्}}\Afootnote{K N H; \skt{भुक्त}\ A1938}}}{\lemma{\skt{अपराह्णे}}\Afootnote{K N H; \skt{मपराह्णे}\ A1938}}\skt{ }\edtext{\edtext{\skt{श्छर्द्दयत्यविदग्धं}}{\lemma{\skt{श्छर्द्दयत्य्}}\Afootnote{K N H; \skt{छर्दय}\ A1938}}}{\lemma{\skt{अविदग्धं}}\Afootnote{K N H; \skt{त्यविदग्ध}\ A1938}}\skt{ }\edtext{\skt{सार्यते वा}}{\Afootnote{K N H; \skt{मतिसार्यते वा}\ A1938}}\skt{ }\edtext{\skt{ज्वरकासाभिभूतः स}}{\Afootnote{K H; \skt{स}\ A1938; \skt{ज्वरकासाराभिभूतः स}\ N}}\skt{ }\edtext{\skt{श्वासान्म्रियते}}{\Afootnote{K H; \skt{श्वासान्म्रियते,}\ A1938; \skt{श्वासा मयते ।}\ N}}\skt{ । }\edtext{\skt{वस्तवद्विलपमानो}}{\Afootnote{K N H; \skt{बस्तवद्विलपन्यश्च}\ A1938}}\skt{ भूमौ पतति स्रस्तमुष्कः }\edtext{\skt{स्तब्धशेपो}}{\Afootnote{K N H; \skt{स्तब्धमेढ्रो}\ A1938}}\skt{ भग्नग्रीवः }\edtext{\skt{प्रणष्टमेहश्च}}{\Afootnote{K N H; \skt{प्रनष्टमेहनश्च}\ A1938}}\skt{ मनुष्यः }\edtext{\skt{प्राग्विशुष्यमाणहृदय}}{\Afootnote{K\textsuperscript{pc} N H; \skt{प्राग्विशुष्यमाणाहृदय}\ A1938; \skt{प्राग्विशुष्यमाण}\ K\textsuperscript{ac}}}\skt{ आर्द्द्रशरीरः । यश्च लोष्टं लोष्टेनाभिहन्ति । }\edtext{\skt{काष्ठम्वा}}{\Afootnote{K N H; \skt{काष्ठं}\ A1938}}\skt{ काष्ठेन तृणानि वा }\edtext{\skt{च्छिन्नन्ति}}{\Afootnote{H; \skt{छिनत्ति}\ A1938 K N}}\skt{ । }\edtext{\skt{अधरोष्ठम्वा}}{\Afootnote{K N H; \skt{अधरोष्ठं}\ A1938}}\skt{ दशति उत्तरोष्ठम्वा लेढि । }\edtext{\skt{आलुञ्चति}}{\Afootnote{K N H; \skt{आलुञ्चति वा}\ A1938}}\skt{ कर्ण्णौ केशाश्च । }\edtext{\skt{देवद्विजगुरुसुहृद्वैद्यान्वा}}{\Afootnote{H; \skt{देवद्विजगुरुसुहृद्वैद्यांश्च}\ A1938; \skt{देवद्विजगुरुसुहृद्वैद्यां वा}\ K; \skt{देवद्विजगुरुसुहृद्वैद्यां}\ N}}\skt{ }\edtext{\skt{विद्वेष्टि}}{\Afootnote{K N H; \skt{द्वेष्टि,}\ A1938}}\skt{ }\edtext{\skt{यस्य च}}{\Afootnote{K N H; \skt{यस्य}\ A1938}}\skt{ वक्रानुवक्रगा ग्रहा }\edtext{\skt{गर्हितस्थानगताः}}{\Afootnote{K N H; \skt{गर्हितस्थानगताः पीडयन्ति,}\ A1938}}\skt{ । }\edtext{\edtext{\edtext{\skt{जन्मऋक्षम्वास्योल्काशनिभ्यामभिहन्यते}}{\lemma{\skt{जन्मऋक्षम्}}\Afootnote{K N H; \skt{जन्मर्क्षं वा,}\ A1938}}}{\lemma{\skt{वास्योल्काशनिभ्याम्}}\Afootnote{K N H; \skt{यस्योल्काशनिभ्या}\ A1938}}}{\lemma{\skt{अभिहन्यते}}\Afootnote{K N H; \skt{मभिहन्यते}\ A1938}}\skt{ । }\edtext{\skt{रात्रौ वा}}{\Afootnote{N H; \skt{होरा वा,}\ A1938; \emph{om.}\ K}}\skt{ }\edtext{\skt{गृहद्वारशयनासनयानवाहनमणिरत्नोपकरणगर्हितलक्षणनिमित्तप्रादुर्ब्भावो}}{\Afootnote{N; \skt{गृहदारशयनासनयानवाहनमणिरत्नोपकरणगर्हितलक्षणनिमित्तप्रादुर्भावो}\ A1938; \emph{om.}\ K; \skt{गृहद्वारशयनाशनयानवाहनमणिरत्नोपकरणगर्हितलक्षणनिमित्तप्रादुर्ब्भावो}\ H}}\skt{ }\edtext{\skt{वेति}}{\Afootnote{A1938 N H; \emph{om.}\ K}}\skt{ ॥}\pend

\pstart\versepar \ledsidenote{\passref{1.32.5}}\edtext{\skt{भवन्ति}}{\Afootnote{A1938; \emph{om.}\ N H}}\skt{ }\edtext{\skt{चात्र}}{\Afootnote{A1938; \emph{om.}\ N H}}\skt{ ॥ }\pend
\pstart\versepar \skt{}\edtext{\skt{चिकित्स्यमानः}}{\Afootnote{A1938 H; \skt{भवन्ति चात्र ॥ चिकित्स्यमान}\ N}}\skt{ }\edtext{\skt{सम्यक्तु}}{\Afootnote{N H; \skt{सम्यक्च}\ A1938}}\skt{ }\edtext{\skt{विकारो }\edtext{\skt{योऽभिवर्द्धते}}{\lemma{\skt{ऽभिवर्द्धते}}\Afootnote{A1938; \skt{भिवर्धते}\ N H}}}{\lemma{\skt{विकारो यो}}\Afootnote{A1938 H; \skt{विकारो}\ N}}\skt{ । }\pend
\pstart\versepar \skt{प्रक्षीणबलमांसस्य लक्षणन्तद्गतायुषः ॥}\pend

\pstart\versepar \ledsidenote{\passref{1.32.6}}\skt{निवर्त्तते महाव्याधिः सहसा यस्य देहिनः । }\pend
\pstart\versepar \skt{न चाहारफलं यस्य दृश्यते स }\edtext{\skt{विनश्यति}}{\Afootnote{A1938 H; \skt{विनस्यति ॥}\ N}}\skt{ ॥}\pend

\pstart\versepar \ledsidenote{\passref{1.32.7}}\edtext{\skt{एतानि}}{\Afootnote{A1938 H; \skt{एता}\ N}}\skt{ }\edtext{\skt{रिष्टरूपाणि}}{\Afootnote{N H; \skt{अरिष्टरूपाणि}\ A1938}}\skt{ सम्यग्बुध्येत यो भिषक् । }\pend
\pstart\versepar \skt{साध्यासाध्यपरीक्षञ्च स }\edtext{\skt{राज्ञः}}{\Afootnote{A1938 H; \emph{om.}\ N}}\skt{ सम्मतो }\edtext{\skt{भवेदिति}}{\Afootnote{N H; \skt{भवेत् ॥७॥}\ A1938}}\skt{ ॥}\edtext{}{\lemma{\skt{भिषक् । साध्यासाध्यपरीक्षञ्च स}}\Afootnote{H; \skt{भिषक् । साध्यासाध्यपरीक्षायां स}\ A1938; \skt{भिषक् । साध्यासाध्यपरीक्षञ्च}\ N}}\pend

\pstart \ledsidenote{\passref{1.32.trailer}}\edtext{\skt{॥ ३१ ॥ {\symfont ❈}}}{\Afootnote{H; \skt{इति सुश्रुतसंहितायां सूत्रस्थाने स्वभावविप्रतिपत्तिर्नाम द्वात्रिंशोऽध्यायः ॥३२॥}\ A1938; \skt{३१ ॥ ० ॥}\ N}}\skt{ ॥}\pend

\endnumbering
\beginnumbering
\eledsection*{Adhyāya 33}
\pstart \ledsidenote{\passref{1.33.1}}\skt{अथातो }\edtext{\edtext{\skt{वारणीयमध्यायं}}{\lemma{\skt{वारणीयम्}}\Afootnote{N H; \skt{ऽवारणीयम्}\ A1938}}}{\lemma{\skt{अध्यायं}}\Afootnote{A1938 H; \skt{अध्याय}\ N}}\skt{ व्याख्यास्यामः ॥}\pend

\pstart\versepar \ledsidenote{\passref{1.33.3}}\skt{उपद्रवैस्तु ये जुष्टा-व्याधयो यान्त्यवार्यतां । }\pend
\pstart\versepar \skt{}\edtext{\edtext{\skt{रसायनैर्व्विना}}{\lemma{\skt{रसायनैर्}}\Afootnote{N H; \skt{रसायना}\ A1938}}}{\lemma{\skt{व्विना}}\Afootnote{N H; \skt{द्विना}\ A1938}}\skt{ }\edtext{\skt{वत्स तां}}{\Afootnote{N H; \skt{वत्स! ता}\ A1938}}\skt{ }\edtext{\skt{शृण्वेकमना}}{\Afootnote{H; \skt{न्शृण्वेकमना}\ A1938; \skt{च्छृण्वैकमना}\ N}}\skt{ }\edtext{\skt{मम}}{\Afootnote{A1938 H; \skt{ममः ॥}\ N}}\skt{ ॥}\pend

\pstart\versepar \ledsidenote{\passref{1.33.4}}\skt{वातव्याधिः }\edtext{\skt{प्रमेहश्च}}{\Afootnote{A1938 H; \skt{प्रमेहस्र}\ N}}\skt{ }\edtext{\skt{कुष्ठान्यथ}}{\Afootnote{N H; \skt{कुष्ठमर्शो}\ A1938}}\skt{ भगन्दरं । }\pend
\pstart\versepar \skt{}\edtext{\skt{अर्शोश्मरी}}{\Afootnote{H; \skt{अश्मरी}\ A1938; \skt{अशीष्मरी}\ N}}\skt{ मूढगर्भो }\edtext{\edtext{\skt{भवत्युदरमष्टमम्}}{\lemma{\skt{भवत्य्}}\Afootnote{N H; \skt{च तथै}\ A1938}}}{\lemma{\skt{उदरम्}}\Afootnote{N H; \skt{वोदरम्}\ A1938}}\skt{ ॥}\pend

\pstart \ledsidenote{\passref{1.33.5}}\skt{अष्टावेते }\edtext{\skt{महान्तः}}{\Afootnote{H; \skt{प्रकृत्यैव दुश्चि}\ A1938; \skt{म}\ N}}\skt{ }\edtext{\edtext{\skt{स्यूर्व्व्याधयो}}{\lemma{\skt{स्यूर्}}\Afootnote{H; \skt{कित्स्या म}\ A1938; \skt{हान्तस्युर्}\ N}}}{\lemma{\skt{व्व्याधयो}}\Afootnote{N H; \emph{om.}\ A1938}}\skt{ }\edtext{\skt{दुस्तराः}}{\Afootnote{H; \skt{हा}\ A1938; \skt{दुस्तरा}\ N}}\skt{ }\edtext{\skt{सदा}}{\Afootnote{H; \skt{गदाः}\ A1938; \skt{मदा ।}\ N}}\skt{ । प्राणमांसक्षयः }\edtext{\edtext{\skt{शोफस्तृष्णा}}{\lemma{\skt{शोफस्}}\Afootnote{N H; \skt{शोष}\ A1938}}}{\lemma{\skt{तृष्णा}}\Afootnote{N H; \skt{स्तृष्णा}\ A1938}}\skt{ }\edtext{\edtext{\skt{च्छर्द्दिर्ज्वरस्तथा}}{\lemma{\skt{च्छर्द्दिर्}}\Afootnote{H; \skt{च्छर्दि}\ A1938; \skt{च्छर्दि}\ N}}}{\lemma{\skt{ज्वरस्}}\Afootnote{N H; \skt{र्ज्वरस्}\ A1938}}\skt{ ॥}\pend

\pstart\versepar \ledsidenote{\passref{1.33.6}}\skt{अतीसारश्च मूर्च्छा च हिक्काश्वासस्तथैव च । }\pend
\pstart\versepar \skt{}\edtext{\skt{एतैरुपद्रवैर्ज्जुष्टान्सर्व्वानेतान्विवर्जयेत्}}{\lemma{\skt{ज्जुष्टान्सर्व्वानेतान्विवर्जयेत्}}\Afootnote{H; \skt{जुष्टान्सर्वानेव विवर्जयेत् ॥६॥}\ A1938; \skt{जुष्टा सर्वानेतान्विवर्जयेत् ।}\ N}}\skt{ ॥}\pend

\pstart\versepar \ledsidenote{\passref{1.33.7}}\skt{शूनं }\edtext{\skt{सुप्तत्वचम्भग्नङ्कम्पाध्माननिपीडितं}}{\lemma{\skt{सुप्तत्वचम्}}\Afootnote{A1938 H; \skt{सुप्तं त्वचं}\ N}}\skt{ । }\pend
\pstart\versepar \skt{}\edtext{\edtext{\edtext{\skt{रुजार्त्तमन्तञ्च}}{\lemma{\skt{रुजार्त्तम्}}\Afootnote{H; \skt{नरं रुजार्त}\ A1938}}}{\lemma{\skt{अन्तञ्च}}\Afootnote{H; \skt{मन्तश्च}\ A1938}}\skt{ }\edtext{\edtext{\skt{नरम्वातव्याधिर्व्विनाशयेत्}}{\lemma{\skt{नरम्}}\Afootnote{H; \emph{om.}\ A1938}}}{\lemma{\skt{व्विनाशयेत्}}\Afootnote{A1938 H; \skt{विवर्ययेदुदरिणं विरिक्तो यो भिपूर्यते ।}\ N}}\skt{ ॥}}{\lemma{\skt{रुजार्त्तम्… व्वि}}\Afootnote{\emph{om.}\ N}}\pend

\pstart\versepar \ledsidenote{\passref{1.33.8}}\edtext{\edtext{\edtext{\skt{यथोक्तोपद्रवाविष्टमतिप्रस्रुतमेव}}{\lemma{\skt{यथोक्तोपद्रवाविष्टम्}}\Afootnote{A1938 H; \skt{यथोक्तपद्रवाविष्ट}\ N}}}{\lemma{\skt{अतिप्रस्रुतम्}}\Afootnote{A1938 H; \skt{मतिप्रस्रुतम्}\ N}}\skt{ च}}{\lemma{\skt{एव च}}\Afootnote{N H; \skt{एव वा ।}\ A1938}}\skt{ । }\pend
\pstart\versepar \skt{}\edtext{\skt{पिटकापीडितं}}{\Afootnote{N H; \skt{पिडकापीडितं}\ A1938}}\skt{ }\edtext{\skt{गाढं}}{\Afootnote{A1938 H; \skt{गाढ}\ N}}\skt{ प्रमेहो हन्ति मानवं ॥}\pend

\pstart\versepar \ledsidenote{\passref{1.33.9}}\edtext{\skt{प्रभिन्नं}}{\Afootnote{N H; \skt{प्रभिन्न}\ A1938}}\skt{ प्रस्रुताङ्गञ्च रक्तनेत्रं }\edtext{\skt{हतस्वरं}}{\Afootnote{A1938 H; \skt{हतं स्वरम् ।}\ N}}\skt{ । }\pend
\pstart\versepar \skt{}\edtext{\skt{पञ्चकर्म्मगुणातीतं}}{\Afootnote{A1938 H; \skt{कञ्चकर्मगुणातीतं}\ N}}\skt{ कुष्ठं }\edtext{\skt{हन्ति हि}}{\Afootnote{N H; \skt{हन्तीह}\ A1938}}\skt{ कुष्ठिनं ॥}\pend

\pstart\versepar \ledsidenote{\passref{1.33.11}}\skt{वातमूत्रपुरीषाणि क्रिमयः शुक्रमेव च । }\pend
\pstart\versepar \skt{भगन्दरात्प्रस्रवन्ति }\edtext{\skt{स}}{\Afootnote{H; \skt{यस्य तं परिवर्जयेत् ॥११॥}\ A1938; \skt{तृष्णारोचकशूला त्रमतिप्रस्रुतशोणितम् । शोफातीसारसंयुक्तन्दु}\ N}}\skt{ }\edtext{\skt{नश्यति}}{\Afootnote{H; \emph{om.}\ A1938; \skt{नामाक्ष}\ N}}\skt{ }\edtext{\skt{भगन्दरी}}{\Afootnote{H; \emph{om.}\ A1938; \skt{पयेन्नरम् ॥}\ N}}\skt{ ॥}\edtext{}{\lemma{\skt{… प्रस्रवन्ति}}\Afootnote{\emph{om.}\ N}}\pend

\pstart\versepar \ledsidenote{\passref{1.33.10}}\edtext{\skt{तृष्णरोचकशूलार्त्तमतिप्रस्रुतशोणितं}}{\lemma{\skt{तृष्णरोचकशूलार्त्तम्}}\Afootnote{H; \skt{तृष्णारोचकशूलार्तम्}\ A1938}}\skt{ । }\pend
\pstart\versepar \skt{}\edtext{\skt{शोफातीसारसंयुक्तं}}{\Afootnote{H; \skt{शोफातीसारसंयुक्तमर्शोव्याधिर्विनाशयेत् ॥१०॥}\ A1938}}\skt{ }\edtext{\skt{दुर्नामा}}{\Afootnote{H; \emph{om.}\ A1938; \skt{दरात्प्र}\ N}}\skt{ }\edtext{\edtext{\skt{क्षपयेन्नरं}}{\lemma{\skt{क्षपयेन्}}\Afootnote{H; \emph{om.}\ A1938; \skt{स्रवन्ति}\ N}}}{\lemma{\skt{नरं}}\Afootnote{H; \emph{om.}\ A1938; \skt{सनेस्यति भगन्दरी ॥}\ N}}\skt{ ॥}\edtext{}{\lemma{\skt{… शोफातीसारसंयु}}\Afootnote{\emph{om.}\ N}}\pend

\pstart\versepar \ledsidenote{\passref{1.33.12}}\edtext{\skt{प्रशूननाभिवृषणं}}{\Afootnote{A1938 H; \skt{प्रसूननाभिवृषण}\ N}}\skt{ }\edtext{\skt{वद्धमूत्ररुजातुरं}}{\Afootnote{H; \skt{रुद्धमूत्रं रुगन्वितम् ।}\ A1938; \skt{बद्धमूत्ररुजातुरम् ।}\ N}}\skt{ । }\pend
\pstart\versepar \skt{अश्मरी }\edtext{\edtext{\skt{क्षपयत्याशु}}{\lemma{\skt{क्षपयत्य्}}\Afootnote{H; \skt{क्षपय}\ A1938; \skt{क्षापय}\ N}}}{\lemma{\skt{आशु}}\Afootnote{H; \skt{त्याशु}\ A1938; \skt{त्याशुः}\ N}}\skt{ }\edtext{\skt{सिकताशर्क्करान्वितं}}{\Afootnote{N H; \skt{सिकता शर्करान्विता ॥१२॥}\ A1938}}\skt{ ॥}\pend

\pstart\versepar \ledsidenote{\passref{1.33.14}}\edtext{\skt{पार्श्वभङ्गान्नविद्वेषः}}{\Afootnote{H; \skt{पार्श्वभङ्गान्नविद्वेष}\ A1938; \skt{पार्श्वभङ्गानविद्वेषः}\ N}}\skt{ शोफातीसारपीडितम् । }\pend
\pstart\versepar \skt{}\edtext{\edtext{\skt{विवर्जयेदुदरिणम्विरिक्तो}}{\lemma{\skt{विवर्जयेदुदरिणम्}}\Afootnote{N H; \skt{विरिक्तं पूर्यमाणं च वर्जयेदुदरार्दि}\ A1938}}\skt{ }\edtext{\skt{योऽभिपूर्यते}}{\lemma{\skt{ऽभिपूर्यते}}\Afootnote{H; \emph{om.}\ A1938; \skt{भिपूज्यते ।}\ N}}}{\lemma{\skt{विरिक्तो यो}}\Afootnote{H; \skt{तम् ॥१४॥}\ A1938; \skt{विरियो}\ N}}\skt{ ॥}\pend

\pstart\versepar \ledsidenote{\passref{1.33.13}}\edtext{\skt{योनीसम्वरणं}}{\Afootnote{N H; \skt{गर्भकोषपरासङ्गो मक्कल्लो योनिसंवृ}\ A1938}}\skt{ सङ्गः }\edtext{\skt{कुक्षो}}{\Afootnote{H; \skt{कुक्षौ}\ N}}\skt{ }\edtext{\skt{मक्कल्लसंज्ञितः}}{\Afootnote{H; \skt{मक्कलसज्ञितः ।}\ N}}\skt{ । }\pend
\pstart\versepar \skt{}\edtext{\skt{हन्युः}}{\Afootnote{N H; \skt{हन्या}\ A1938}}\skt{ }\edtext{\skt{स्त्रियं}}{\Afootnote{H; \skt{त्स्त्रियं}\ A1938; \skt{त्रियं}\ N}}\skt{ }\edtext{\skt{मूढग-र्व्भे}}{\Afootnote{H; \skt{मूढगर्भे}\ A1938; \skt{मूढगर्भो}\ N}}\skt{ }\edtext{\skt{यथोक्ताश्चाप्युपद्रवाः}}{\lemma{\skt{यथोक्ताश्}}\Afootnote{A1938 H; \skt{यथो}\ N}}\skt{ ॥}\edtext{}{\lemma{\skt{रणं… मक्कल्लसंज्ञितः}}\Afootnote{\emph{om.}\ A1938}}\pend

\pstart\versepar \ledsidenote{\passref{1.33.15}}\edtext{\edtext{\skt{विसंज्ञस्ताम्यते}}{\lemma{\skt{विसंज्ञस्ताम्}}\Afootnote{H; \skt{यस्ताम्यति विसञ्ज्ञश्}\ A1938; \skt{विसज्ञस्ता}\ N}}\skt{ यस्तु}}{\lemma{\skt{यते यस्तु}}\Afootnote{H; \skt{च}\ A1938; \skt{म्यते यस्तु}\ N}}\skt{ शेते निपति-तो }\edtext{\skt{यथा}}{\Afootnote{N H; \skt{ऽपि वा ।}\ A1938}}\skt{ । }\pend
\pstart\versepar \skt{शीतार्द्दितोऽन्तरुष्णश्च ज्वरेण म्रियते }\edtext{\skt{नरः}}{\Afootnote{A1938 H; \skt{ननः ।}\ N}}\skt{ ॥}\pend

\pstart\versepar \ledsidenote{\passref{1.33.16}}\edtext{\skt{यो हृष्टरोमा}}{\Afootnote{A1938 H; \skt{यो हृष्टरक्तो}\ N\textsuperscript{ac}; \skt{यो हृष्ट}\ N\textsuperscript{pc}}}\skt{ रक्ताक्षो हृदि संघातशूलवान् । }\pend
\pstart\versepar \skt{वक्त्रेण }\edtext{\skt{चोच्छ्वसिति तं}}{\Afootnote{H; \skt{चोच्छ्वस्यात्तं}\ A1938; \skt{चैवोच्छ्विसितितं}\ N}}\skt{ ज्वरो हन्ति मानवं ॥}\edtext{}{\lemma{\skt{संघातशूलवान् । वक्त्रेण}}\Afootnote{N H; \skt{सङ्घातशूलवान् । नित्यं वक्त्रेण}\ A1938}}\pend

\pstart\versepar \ledsidenote{\passref{1.33.17}}\skt{हिक्काश्वास }\edtext{\skt{समायुक्तं}}{\Afootnote{N H; \skt{पिपासार्तं}\ A1938}}\skt{ }\edtext{\skt{मूढम्विभ्रान्तलोचनं}}{\lemma{\skt{विभ्रान्तलोचनं}}\Afootnote{A1938 H; \skt{चित्रान्तलोचनं ।}\ N}}\skt{ । }\pend
\pstart\versepar \skt{}\edtext{\skt{सन्ततोच्छ्वासिनं}}{\Afootnote{A1938 H; \skt{सन्ततोच्छ्वासितं}\ N}}\skt{ क्षीणं नरं }\edtext{\skt{क्षपयते}}{\Afootnote{N H; \skt{क्षपयति}\ A1938}}\skt{ ज्वरः ॥}\pend

\pstart\versepar \ledsidenote{\passref{1.33.19}}\edtext{\skt{श्वासशूलपिपासार्तं-क्षीणञ्ज्वरनिपीडितं}}{\lemma{\skt{ज्वरनिपीडितं}}\Afootnote{A1938 H; \skt{ज्वरनिपीडनं ।}\ N}}\skt{ । }\pend
\pstart\versepar \skt{विशेषेण नरं }\edtext{\edtext{\skt{वृद्धमतीसारो}}{\lemma{\skt{वृद्धम्}}\Afootnote{N H; \skt{वृद्ध}\ A1938}}}{\lemma{\skt{अतीसारो}}\Afootnote{N H; \skt{मतीसारो}\ A1938}}\skt{ }\edtext{\skt{विनाशयेत्}}{\Afootnote{A1938 H; \skt{विनाशयेत ।}\ N}}\skt{ ॥}\pend

\pstart\versepar \ledsidenote{\passref{1.33.20}}\edtext{\edtext{\skt{शुक्लाक्षमन्नद्वेष्टारमूर्द्धश्वासनिपीडितम्}}{\lemma{\skt{शुक्लाक्षमन्नद्वेष्टारम्}}\Afootnote{H; \skt{शुक्लाक्षमन्नद्वेष्टार}\ A1938; \skt{शुक्लाक्षमन्नद्विष्टारम्}\ N}}}{\lemma{\skt{ऊर्द्धश्वासनिपीडितम्}}\Afootnote{H; \skt{मूर्ध्वश्वासनिपीडितम् ।}\ A1938; \skt{ऊर्ध्वश्वासनिपीडितं ।}\ N}}\skt{ । }\pend
\pstart\versepar \skt{}\edtext{\skt{कृच्छ्रेण}}{\Afootnote{A1938 H; \skt{कृच्छेण}\ N}}\skt{ }\edtext{\skt{बहु}}{\Afootnote{A1938 N; \skt{वहु}\ H}}\skt{ मेहन्तं यक्ष्मा }\edtext{\skt{हन्तीह}}{\Afootnote{A1938 H; \skt{हन्तीव}\ N}}\skt{ मानवं ॥}\pend

\pstart\versepar \ledsidenote{\passref{1.33.21}}\edtext{\skt{श्वासशूलपिपासान्नविद्वेषो}}{\Afootnote{H; \skt{श्वासशूलपिपासान्नविद्वेष}\ A1938; \skt{श्वासशूलपिपासान्नविद्वेषो}\ N}}\skt{ }\edtext{\skt{ग्रन्थिमूढता}}{\Afootnote{A1938 H; \skt{ग्रन्थिरूढता ॥}\ N}}\skt{ । }\pend
\pstart\versepar \skt{}\edtext{\skt{जायते}}{\Afootnote{N H; \skt{भवन्ति}\ A1938}}\skt{ }\edtext{\skt{दुर्व्वलत्वं च}}{\Afootnote{H; \skt{दुर्बलत्वं च}\ A1938 N}}\skt{ }\edtext{\skt{गुल्मिनो}}{\Afootnote{A1938 H; \skt{गुल्मिणो}\ N}}\skt{ }\edtext{\skt{मरणाय}}{\Afootnote{N H; \skt{मृत्युमेष्य}\ A1938}}\skt{ }\edtext{\skt{वै॥}}{\lemma{\skt{वै}}\Afootnote{N H; \skt{तः ॥२१॥}\ A1938}}\pend

\pstart\versepar \ledsidenote{\passref{1.33.22}}\edtext{\skt{आध्मानं}}{\Afootnote{H; \skt{आध्मातं}\ A1938; \skt{आध्मान}\ N}}\skt{ }\edtext{\skt{वद्धनिष्यन्दं}}{\Afootnote{H; \skt{बद्धनिष्यन्दं}\ A1938 N}}\skt{ }\edtext{\skt{च्छर्द्दिहिक्कातृषान्वितम्}}{\Afootnote{H; \skt{छर्दिहिक्कातृडन्वितम् ।}\ A1938; \skt{च्छर्दिहिक्काकृशान्वितं ॥}\ N}}\skt{ । }\pend
\pstart\versepar \skt{}\edtext{\skt{रुजाश्वाससमायुक्तं}}{\Afootnote{N H; \skt{रुजाश्वाससमाविष्टं}\ A1938}}\skt{ }\edtext{\edtext{\skt{विद्रधिर्न्नाशयेत्नरं}}{\lemma{\skt{विद्रधिर्}}\Afootnote{H; \skt{विद्रधि}\ A1938 N}}}{\lemma{\skt{न्नाशयेत्नरं}}\Afootnote{H; \skt{र्नाशयेन्नरम्}\ A1938 N}}\skt{ ॥}\pend

\pstart\versepar \ledsidenote{\passref{1.33.23}}\edtext{\skt{पाण्डुदन्तनखो}}{\Afootnote{A1938 H; \skt{पाण्डुदत्तनखो}\ N}}\skt{ }\edtext{\skt{यस्तु}}{\Afootnote{H; \skt{यश्च}\ A1938; \skt{यस्तु}\ N}}\skt{ पाण्डुनेत्रश्च }\edtext{\skt{मानवः}}{\Afootnote{A1938 H; \skt{मानवं ।}\ N}}\skt{ । }\pend
\pstart\versepar \skt{}\edtext{\skt{पाण्डुसंघातदर्शी च}}{\Afootnote{A1938 H; \skt{पाण्डुसद्यातदशी च}\ N}}\skt{ पाण्डुरोगी विनश्यति ॥}\pend

\pstart\versepar \ledsidenote{\passref{1.33.24}}\edtext{\skt{लोहितञ्छर्द्दयेद्यस्तु}}{\lemma{\skt{लोहितञ्}}\Afootnote{A1938 H; \skt{लोहितं च्छर्दितं}\ N}}\skt{ }\edtext{\skt{बहुशो}}{\Afootnote{A1938 N; \skt{वहुशो}\ H}}\skt{ }\edtext{\skt{लोहितेक्षणः}}{\Afootnote{A1938 H; \skt{लोहितेक्षणेः ।}\ N}}\skt{ । }\pend
\pstart\versepar \skt{}\edtext{\edtext{\skt{लोहितोऽङ्गारदर्शी}}{\lemma{\skt{लोहितो}}\Afootnote{H; \emph{om.}\ A1938; \skt{लोहिताङ्}\ N}}\skt{ च}}{\lemma{\skt{ऽङ्गारदर्शी च}}\Afootnote{H; \skt{रक्तानां च}\ A1938; \skt{गरदशीच}\ N}}\skt{ }\edtext{\skt{म्रियते}}{\Afootnote{N H; \skt{दिशां द्रष्टा}\ A1938}}\skt{ }\edtext{\skt{रक्तपैत्तिकः}}{\Afootnote{H; \skt{रक्तपित्ती विनश्यति ॥२४॥}\ A1938; \skt{रक्तपित्तकः ।}\ N}}\skt{ ॥}\pend

\pstart\versepar \ledsidenote{\passref{1.33.25}}\edtext{\skt{अवाचो पि}}{\Afootnote{N H; \emph{om.}\ A1938}}\skt{ }\edtext{\edtext{\skt{सवाचोऽपि}}{\lemma{\skt{सवाचो}}\Afootnote{N H; \skt{अवाङ्}\ A1938}}}{\lemma{\skt{ऽपि}}\Afootnote{N H; \skt{मुखस्तून्मुखो वा}\ A1938}}\skt{ }\edtext{\skt{क्षीणमान्सवलो}}{\Afootnote{H; \skt{क्षीणमांसबलो}\ A1938 N}}\skt{ नरः । }\pend
\pstart\versepar \skt{}\edtext{\skt{जागरूको}}{\Afootnote{N H; \skt{जागरिष्णुर्}\ A1938}}\skt{ }\edtext{\skt{क्ष्यासन्देहमुन्मादेन}}{\lemma{\skt{क्ष्यासन्देहम्}}\Afootnote{H; \skt{असन्देहम्}\ A1938; \skt{ह्यसन्देहम्}\ N}}\skt{ विनश्यति ॥}\pend

\pstart\versepar \ledsidenote{\passref{1.33.26}}\edtext{\edtext{\skt{अपस्मरन्तम्बहुशः}}{\lemma{\skt{अपस्मरन्तम्}}\Afootnote{H; \skt{बहुशोऽपस्मरन्तं}\ A1938; \skt{अपरस्मरन्त}\ N}}}{\lemma{\skt{बहुशः}}\Afootnote{\emph{om.}\ A1938; \skt{बहुशः प्र}\ N; \skt{वहुशः}\ H}}\skt{ }\edtext{\edtext{\skt{प्रक्षीणञ्चलितं}}{\lemma{\skt{प्रक्षीणञ्}}\Afootnote{H; \skt{तु प्रक्षीणं}\ A1938; \skt{भ}\ N}}}{\lemma{\skt{चलितं}}\Afootnote{H; \skt{चलित}\ A1938; \skt{लित}\ N}}\skt{ भ्रुवं । }\pend
\pstart\versepar \skt{नेत्राभ्याञ्च विकुर्व्वाणमपस्मारो }\edtext{\skt{विनाशयेत्}}{\Afootnote{A1938 H; \skt{विनाशनात् ॥}\ N}}\skt{ ॥}\pend

\pstart \ledsidenote{\passref{1.33.trailer}}\edtext{\skt{३३ ॥}}{\Afootnote{\skt{इति सुश्रुतसंहितायां सूत्रस्थानेऽवारणीयो नाम त्रयस्त्रिंशत्तमोऽध्यायः ॥३३॥}\ A1938; \skt{ल ॥}\ N; \skt{३३ ॥}\ H}}\pend

\endnumbering
\beginnumbering
\eledsection*{Adhyāya 34}
\pstart \ledsidenote{\passref{1.34.1}}\skt{अथातो }\edtext{\skt{युक्तसेनीयं}}{\Afootnote{H; \skt{युक्तसेनीयमध्यायं}\ A1938}}\skt{ व्याख्यास्यामः ॥}\pend

\pstart\versepar \ledsidenote{\passref{1.34.1.add}}\edtext{\edtext{\skt{धन्वन्तरिम्महाप्राज्ञं}}{\lemma{\skt{धन्वन्तरिम्}}\Afootnote{H\textsuperscript{pc}; \skt{धन्वन्तरि}\ N; \emph{om.}\ H\textsuperscript{ac}}}}{\lemma{\skt{महाप्राज्ञं}}\Afootnote{H\textsuperscript{pc}; \skt{महाप्रज्ञं}\ N; \emph{om.}\ H\textsuperscript{ac}}}\skt{ }\edtext{\skt{सर्व्वशास्त्रविशारदं}}{\Afootnote{N H\textsuperscript{pc}; \emph{om.}\ H\textsuperscript{ac}}}\skt{ । }\pend
\pstart\versepar \skt{}\edtext{\edtext{\skt{चरणावुपसंगृह्य}}{\lemma{\skt{चरणाव्}}\Afootnote{N H\textsuperscript{pc}; \emph{om.}\ H\textsuperscript{ac}}}}{\lemma{\skt{उपसंगृह्य}}\Afootnote{N H\textsuperscript{pc}; \emph{om.}\ H\textsuperscript{ac}}}\skt{ }\edtext{\skt{सुश्रुतः}}{\Afootnote{N H\textsuperscript{pc}; \emph{om.}\ H\textsuperscript{ac}}}\skt{ }\edtext{\skt{परिपृच्छति}}{\Afootnote{N H\textsuperscript{pc}; \emph{om.}\ H\textsuperscript{ac}}}\skt{ ॥}\pend

\pstart\versepar \ledsidenote{\passref{1.34.3}}\edtext{\skt{युक्तसेनस्य}}{\Afootnote{N H; \skt{नृपतेर्युक्तसेनस्य}\ A1938}}\skt{ }\edtext{\skt{नृपतेः}}{\Afootnote{N H; \emph{om.}\ A1938}}\skt{ }\edtext{\skt{परानभिजिघांसतः}}{\Afootnote{H; \skt{परानभिजिगीषतः ।}\ A1938; \skt{परानपि जिघांसतः ।}\ N}}\skt{ । }\pend
\pstart\versepar \skt{भिषजा रक्षणं कार्यं यथा }\edtext{\skt{तद्ब्रूहि मे}}{\Afootnote{H; \skt{तदुपदे}\ A1938; \skt{ब्रूहि मे}\ N}}\skt{ }\edtext{\skt{मुने}}{\Afootnote{N H; \skt{क्ष्यते ॥३॥}\ A1938}}\skt{ ॥}\pend

\pstart \ledsidenote{\passref{1.34.3.add}}\edtext{\skt{तस्य}}{\Afootnote{N H\textsuperscript{pc}; \emph{om.}\ H\textsuperscript{ac}}}\skt{ }\edtext{\skt{तद्वचनं}}{\Afootnote{N H\textsuperscript{pc}; \emph{om.}\ H\textsuperscript{ac}}}\skt{ }\edtext{\skt{श्रुत्वा}}{\Afootnote{N H\textsuperscript{pc}; \emph{om.}\ H\textsuperscript{ac}}}\skt{ }\edtext{\edtext{\skt{प्राब्रवीद्भिषजाम्वरः}}{\lemma{\skt{प्राब्रवीद्भिषजाम्}}\Afootnote{N H\textsuperscript{pc}; \emph{om.}\ H\textsuperscript{ac}}}}{\lemma{\skt{वरः}}\Afootnote{N H\textsuperscript{pc}; \emph{om.}\ H\textsuperscript{ac}}}\skt{ ॥}\pend

\pstart\versepar \ledsidenote{\passref{1.34.4}}\edtext{\skt{विजिगीषुः}}{\Afootnote{A1938 H; \skt{विजिगीषु}\ N}}\skt{ }\edtext{\edtext{\skt{सहामात्यैर्यात्रायुक्तः}}{\lemma{\skt{सहामात्यैर्}}\Afootnote{A1938 H; \skt{महामात्यैर्}\ N}}}{\lemma{\skt{यात्रायुक्तः}}\Afootnote{A1938 H; \skt{यात्रायुक्त}\ N}}\skt{ प्रयत्नतः । }\pend
\pstart\versepar \skt{रक्षितव्यो विशेषेण }\edtext{\skt{विषादेव}}{\Afootnote{H; \skt{विषादेव}\ A1938; \skt{विषादेन}\ N}}\skt{ नराधिपः ।}\pend

\pstart\versepar \ledsidenote{\passref{1.34.5}}\edtext{\skt{पन्थानमुदकञ्छायां}}{\lemma{\skt{छायां}}\Afootnote{A1938 H; \skt{स्थाया}\ N}}\skt{ भक्तं }\edtext{\skt{यवसमिन्धनं}}{\Afootnote{N H; \skt{यवसमिन्धनम् ।}\ A1938}}\skt{ ॥ }\pend
\pstart\versepar \skt{}\edtext{\skt{दुषयन्त्यरयस्तानि}}{\lemma{\skt{दुषयन्त्य्}}\Afootnote{H; \skt{दूषयन्}\ A1938; \skt{दूषयन्त्य्}\ N}}\skt{ }\edtext{\skt{जानीयाच्छोधनीय}}{\Afootnote{H; \skt{जानीयाच्छोरयीति}\ N}}\skt{ च । }\pend
\pstart\versepar \skt{तस्य लिङ्गं चिकित्साञ्च कल्पस्थाने प्रवक्ष्यते ॥}\edtext{}{\lemma{\skt{त्य्… छोधनीय}}\Afootnote{\emph{om.}\ A1938}}\edtext{}{\lemma{\skt{च । तस्य लिङ्गं चिकित्साञ्च कल्पस्थाने प्रवक्ष्यते}}\Afootnote{\skt{रयस्तच्च जानीयाच्छोधयेत्तथा ॥५॥}\ A1938; \skt{च ।}\ N H}}\pend

\pstart\versepar \ledsidenote{\passref{1.34.6}}\edtext{\skt{एकोत्तरं}}{\Afootnote{\skt{तस्य लिङ्गं चिकित्सा च कल्पस्थाने प्रवक्ष्यते एकोत्तरं}\ A1938; \skt{तस्य लिङ्गं चिकित्साञ्च कल्पस्थाने प्रवक्ष्यते । एकोत्तर}\ N; \skt{तस्य लिङ्गं चिकित्साञ्च कल्पस्थाने प्रवक्ष्यते ॥ एकोत्तरं}\ H}}\skt{ }\edtext{\skt{मृत्युशतमथर्व्वाणः}}{\lemma{\skt{अथर्व्वाणः}}\Afootnote{A1938 H; \skt{अथर्वाण}\ N}}\skt{ प्रचक्षते । }\pend
\pstart\versepar \skt{}\edtext{\skt{तत्रैकः कालसंयुक्तः शेषास्त्वागन्तवः स्मृताः}}{\Afootnote{\emph{om.}\ A1938 N H}}\skt{ ।}\pend

\pstart\versepar \ledsidenote{\passref{1.34.7}}\edtext{\skt{दोषागन्तुनिमित्तेभ्यो}}{\Afootnote{\skt{तत्रैकः कालसंयुक्तः शेषा आगन्तवः स्मृताः दोषागन्तुजमृत्युभ्यो}\ A1938; \skt{तत्रैक कालसयुक्तः शेषास्त्वागन्तवः स्मृताः । दोषागन्त्रुनिमित्तेभ्यो}\ N; \skt{तत्रैकः कालसंयुक्तः शेषास्त्वागन्तवः स्मृताः । दोषागन्तुनिमित्तेभ्यो}\ H}}\skt{ }\edtext{\skt{रसमन्त्रविशारदौ}}{\Afootnote{A1938 H; \skt{रसमविशारदौ ।}\ N}}\skt{ । }\pend
\pstart\versepar \skt{}\edtext{\skt{रक्षतां नृपतिन्नित्यं यत्नाद्वैद्यपुरोहितौ}}{\Afootnote{\emph{om.}\ A1938 N H}}\skt{ ॥}\pend

\pstart\versepar \ledsidenote{\passref{1.34.8}}\edtext{\skt{ब्रह्मावेदाङ्गमष्टाङ्गमायुर्व्वेदम्प्रचक्षते}}{\lemma{\skt{ब्रह्मावेदाङ्गमष्टाङ्गम्}}\Afootnote{\skt{रक्षेतां नृपतिं नित्यं यत्तौ वैद्यपुरोहितौ ब्रह्मा वेदाङ्गमष्टाङ्गम्}\ A1938; \skt{रक्षेतां नृपतिनिपति न्नाद्वैद्यपुरोहितौ । ब्रह्मा वेदाङ्गमष्टाङ्गम्}\ N; \skt{रक्षतां नृपतिन्नित्यं यत्नाद्वैद्यपुरोहितौ ॥ ब्रह्मावेदाङ्गमष्टाङ्गम्}\ H}}\skt{ । }\pend
\pstart\versepar \skt{तस्मात्पुरोहितमते वर्त्तेत भिषगात्मवान् ॥}\edtext{}{\lemma{\skt{प्रचक्षते… वर्त्तेत}}\Afootnote{\emph{om.}\ A1938}}\edtext{}{\lemma{\skt{प्रचक्षते । तस्मात्पुरोहितमते वर्त्तेत भिषगात्मवान्}}\Afootnote{\skt{प्रचक्षते ।}\ N H}}\pend

\pstart\versepar \ledsidenote{\passref{1.34.9}}\edtext{\skt{सङ्करः}}{\Afootnote{\skt{पुरोहितमते तस्माद्वर्तेत भिषगात्मवान्सङ्करः}\ A1938; \skt{तस्मात्पुरोहितमते वर्त्तेत भिषगात्मवान् । शङ्करः}\ N; \skt{तस्मात्पुरोहितमते वर्त्तेत भिषगात्मवान् ॥ सङ्करः}\ H}}\skt{ सर्व्ववर्ण्णानां }\edtext{\skt{विनाशो}}{\Afootnote{N H; \skt{प्रणाशो}\ A1938}}\skt{ }\edtext{\skt{धर्म्मकर्म्मिणां}}{\Afootnote{N H; \skt{धर्मकर्मणाम् ॥९॥}\ A1938}}\skt{ । }\pend
\pstart\versepar \skt{}\edtext{\skt{प्रजनामपि कृच्छ्राणि भवन्ति नृपनाशतः}}{\Afootnote{\emph{om.}\ A1938 N H}}\skt{ ॥}\pend

\pstart\versepar \ledsidenote{\passref{1.34.10}}\edtext{\edtext{\skt{पुरुषाणान्नृपानाञ्च}}{\lemma{\skt{पुरुषाणान्}}\Afootnote{\skt{प्रजानामपि चोच्छित्तिर्नृपव्यसनहेतुतः पुरुषाणां}\ A1938; \skt{प्रजानामपि कृच्छ्राणि भवन्ति नृपनाशतः ॥ पुरुषाणान्}\ N; \skt{प्रजनामपि कृच्छ्राणि भवन्ति नृपनाशतः ॥ पुरुषाणान्}\ H}}}{\lemma{\skt{नृपानाञ्च}}\Afootnote{H; \skt{नृपाणां च}\ A1938 N}}\skt{ केवलन्तुल्यमूर्त्तिता । }\pend
\pstart\versepar \skt{}\edtext{\skt{आज्ञात्यागः क्षमा धैर्यम्विक्रमश्चाप्यमानुषः}}{\Afootnote{\emph{om.}\ A1938 N H}}\skt{ ॥}\pend

\pstart\versepar \ledsidenote{\passref{1.34.11}}\edtext{\edtext{\skt{तस्माद्देवमिवाभीक्ष्णम्वाङ्मनः}}{\lemma{\skt{तस्माद्देवम्}}\Afootnote{\skt{आज्ञा त्यागः क्षमा धैर्यं विक्रमश्चाप्यमानुषः । तस्माद्देवम्}\ A1938; \skt{आज्ञात्यागः क्षमा धेर्यम्विक्रमश्चाप्यमानुषः । तस्मादेवम्}\ N; \skt{आज्ञात्यागः क्षमा धैर्यम्विक्रमश्चाप्यमानुषः ॥ तस्माद्देवम्}\ H}}}{\lemma{\skt{इवाभीक्ष्णम्}}\Afootnote{A1938 H; \skt{इवाक्ष क्ष्णम्}\ N}}\skt{ }\edtext{\skt{कर्मभिः}}{\Afootnote{A1938 H; \skt{कर्म}\ N}}\skt{ शुभैः । }\pend
\pstart\versepar \skt{}\edtext{\skt{चिन्तयेत्नृपतिम्वैद्यः श्रेयांसीच्छन्विचक्षणः}}{\Afootnote{\emph{om.}\ A1938 H}}\skt{ ॥}\edtext{}{\lemma{\skt{भिः… विच}}\Afootnote{\emph{om.}\ N}}\pend

\pstart\versepar \ledsidenote{\passref{1.34.12}}\edtext{\skt{स्कन्धावारे च}}{\Afootnote{\skt{चिन्तयेन्नृपतिं वैद्यः श्रेयांसीच्छन्विचक्षणः स्कन्धावारे च}\ A1938; \skt{चिन्तयन्नृपतिवैद्यः श्रेयान्सिच्छविचक्षणः । स्कन्धावारे च}\ N; \skt{चिन्तयेत्नृपतिम्वैद्यः श्रेयांसीच्छन्विचक्षणः ॥ स्कन्धावारे च}\ H}}\skt{ महति }\edtext{\skt{राजवेश्म}}{\Afootnote{H; \skt{राज}\ A1938; \skt{राजचेश्मः}\ N}}\skt{ समीपतः । }\pend
\pstart\versepar \skt{भवेत्सन्निहितो वैद्यः सर्व्वोपकरणान्वितः ॥}\edtext{}{\lemma{\skt{वेश्म… स}}\Afootnote{\emph{om.}\ A1938}}\edtext{}{\lemma{\skt{समीपतः । भवेत्सन्निहितो वैद्यः सर्व्वोपकरणान्वितः}}\Afootnote{\skt{समिपतः ।}\ N; \skt{समीपतः ।}\ H}}\pend

\pstart\versepar \ledsidenote{\passref{1.34.13}}\edtext{\edtext{\edtext{\edtext{\skt{तत्रस्थमेनन्ध्वजवद्यशःख्यातिभिरुच्छ्रितं}}{\lemma{\skt{तत्रस्थम्}}\Afootnote{\skt{भवेत्सन्निहितो वैद्यः सर्वोपकरणान्वितः । तत्रस्थ}\ A1938; \skt{भवैत्सन्निहितो वैद्यः सर्वोपकरणान्वितः । तत्रस्थ}\ N; \skt{भवेत्सन्निहितो वैद्यः सर्व्वोपकरणान्वितः ॥ तत्रस्थम्}\ H}}}{\lemma{\skt{एनन्}}\Afootnote{H; \skt{मेनं}\ A1938; \skt{मैन}\ N}}}{\lemma{\skt{ध्वजवद्यशःख्यातिभिर्}}\Afootnote{N H; \skt{ध्वजवद्यशःख्यातिस}\ A1938}}}{\lemma{\skt{उच्छ्रितं}}\Afootnote{\skt{मुच्छ्रितम् ॥१३॥}\ A1938; \skt{उच्छ्रितः ।}\ N; \skt{उच्छ्रितंः ॥}\ H}}\skt{ ॥ }\pend
\pstart\versepar \skt{}\edtext{\skt{उपसर्प्पन्त्यमोहेन विषशल्यामयार्द्दिताः}}{\Afootnote{\emph{om.}\ A1938 N H}}\skt{ ॥}\pend

\pstart\versepar \ledsidenote{\passref{1.34.14}}\edtext{\edtext{\skt{स्वतन्त्रकुशलोऽन्येषु}}{\lemma{\skt{स्वतन्त्रकुशलो}}\Afootnote{\skt{उपसर्पन्त्यमोहेन विषशल्यामयार्दिताः स्वतन्त्रकुशलो}\ A1938 H; \skt{उपसर्पन्ति मोहेन विषशल्यामर्यार्दिता स्वतन्त्रभूशलो}\ N}}}{\lemma{\skt{ऽन्येषु}}\Afootnote{A1938 H; \skt{न्येषु}\ N}}\skt{ }\edtext{\skt{शास्त्रार्थेष्वबहिः}}{\lemma{\skt{अबहिः}}\Afootnote{H; \skt{अबहि}\ A1938; \skt{बहिः}\ N}}\skt{ कृतः । }\pend
\pstart\versepar \skt{वैद्यो ध्वजमिवाभाति नृपतद्विद्यपूजितः ॥}\edtext{}{\lemma{\skt{कृतः । वैद्यो ध्वजमिवाभाति नृपतद्विद्यपूजितः}}\Afootnote{\skt{ष्कृतः ॥१४॥}\ A1938; \skt{कृतः ।}\ N H}}\pend

\pstart\versepar \ledsidenote{\passref{1.34.15}}\edtext{\skt{वैद्यो}}{\Afootnote{\skt{वैद्यो ध्वज इवाभाति नृपतद्विद्यपूजितः वैद्यो}\ A1938; \skt{वैद्योमिवाभाति नृपतद्विद्यपूजितं । वैद्यो}\ N; \skt{वैद्यो ध्वजमिवाभाति नृपतद्विद्यपूजितः ॥ वैद्यो}\ H}}\skt{ व्याध्युपसृष्टश्च }\edtext{\edtext{\skt{भैषजम्परिचारकः}}{\lemma{\skt{भैषजम्}}\Afootnote{H; \skt{भेषजं}\ A1938; \skt{भैषजुम्}\ N}}}{\lemma{\skt{परिचारकः}}\Afootnote{\skt{परिचारकः}\ A1938 N H}}\skt{ । }\pend
\pstart\versepar \skt{}\edtext{\skt{एते पादाश्चिकित्सायाः कर्म्मसाधनहेतवः}}{\Afootnote{\emph{om.}\ A1938 N H}}\skt{ ॥}\pend

\pstart\versepar \ledsidenote{\passref{1.34.16}}\edtext{\skt{गुणवद्भिस्त्रिभिः}}{\lemma{\skt{गुणवद्भिस्}}\Afootnote{\skt{एते पादाश्चिकित्सायाः कर्मसाधनहेतवः । गुणवद्भिस्}\ A1938; \skt{एते पादाचिकित्सायाः कर्मसाधवः । गुणवद्भिस्}\ N; \skt{एते पादाश्चिकित्सायाः कर्म्मसाधनहेतवः ॥ गुणवद्भिस्}\ H}}\skt{ }\edtext{\skt{पादैश्चतुर्थ्यो}}{\lemma{\skt{चतुर्थ्यो}}\Afootnote{H; \skt{चतुर्थो}\ A1938 N}}\skt{ गुणवान्भिषक् । }\pend
\pstart\versepar \skt{व्याधिमल्पेन कालेन महान्तमपि साधयेत् ॥}\edtext{}{\lemma{\skt{गुणवान्भिषक् । व्याधिमल्पेन कालेन महान्तमपि साधयेत्}}\Afootnote{\skt{गुणवान्भिषक्}\ A1938 N H}}\pend

\pstart\versepar \ledsidenote{\passref{1.34.17}}\edtext{\edtext{\skt{वैद्यहीनास्त्रयः}}{\lemma{\skt{वैद्यहीनास्}}\Afootnote{\skt{व्याधिमल्पेन कालेन महान्तमपि साधयेत्वैद्यहीना}\ A1938; \skt{व्याधिमल्पेन काल्पेमहान्तमहान्तमपि साधयेत् । वैद्यहीनास्}\ N; \skt{व्याधिमल्पेन कालेन महान्तमपि साधयेत् ॥ वैद्यहीनास्}\ H}}}{\lemma{\skt{त्रयः}}\Afootnote{N H; \skt{स्त्रयः}\ A1938}}\skt{ पादा }\edtext{\skt{गुणवन्तोऽप्यपार्थकाः}}{\lemma{\skt{ऽप्यपार्थकाः}}\Afootnote{\skt{ऽप्यपार्थकाः ॥१७॥}\ A1938; \skt{ऽप्यपार्थकाः ।}\ N H}}\skt{ । }\pend
\pstart\versepar \skt{}\edtext{\skt{उद्गातृहोतृब्रह्माणो यथाध्वर्यम्विनाध्वरे}}{\Afootnote{\emph{om.}\ A1938 N H}}\skt{ ॥}\pend

\pstart\versepar \ledsidenote{\passref{1.34.18}}\edtext{\skt{वैद्यस्तु}}{\Afootnote{\skt{उद्गातृहोतृब्रह्माणो यथाऽध्वर्युं विनाऽध्वरे वैद्यस्तु}\ A1938; \skt{उद्गातृहोतृब्रह्माणो यथाध्वर्यं विनाध्वरे वैद्यस्तु}\ N H}}\skt{ }\edtext{\skt{गुणवानेको}}{\Afootnote{N H; \skt{गुणवानेक}\ A1938}}\skt{ }\edtext{\skt{यापयेदातुरं}}{\Afootnote{N H; \skt{स्तारयेदातुरां}\ A1938}}\skt{ सदा । }\pend
\pstart\versepar \skt{}\edtext{\skt{प्लवम्परिचरन्धीरः कर्ण्णधारमिवाम्भसि}}{\Afootnote{\emph{om.}\ A1938 N H}}\skt{ ॥}\pend

\pstart\versepar \ledsidenote{\passref{1.34.19}}\edtext{\skt{तत्वोपगतशास्त्रार्थो}}{\Afootnote{\skt{प्लवं प्रतितरैर्हीनं कर्णधार इवाम्भसि तत्त्वाधिगतशास्त्रार्थो}\ A1938; \skt{प्लवं परिचरं धीरं कर्णधारमिवाम्भसि । तत्वापगशास्त्रार्थो}\ N; \skt{प्लवम्परिचरन्धीरः कर्ण्णधारमिवाम्भसि ॥ तत्वोपगतशास्त्रार्थो}\ H}}\skt{ दृष्टकर्म्मा स्वयङ्कृती । }\pend
\pstart\versepar \skt{}\edtext{\skt{लघुहस्तः शुचिः सूरः सज्जोपस्करभेषजः}}{\Afootnote{\emph{om.}\ A1938 N H}}\skt{ ॥}\pend

\pstart\versepar \ledsidenote{\passref{1.34.20}}\edtext{\edtext{\skt{प्रत्युत्पन्नमतिर्द्धीमान्व्यवसायी}}{\lemma{\skt{प्रत्युत्पन्नमतिर्}}\Afootnote{\skt{लघुहस्तः शुचिः शूरः सज्जोपस्करभेषजः । प्रत्युत्पन्नमति}\ A1938; \skt{लघुहस्तः शुचिः शूरः सज्जोपस्करभेषजः । प्रत्युत्पन्नमति}\ N; \skt{लघुहस्तः शुचिः सूरः सज्जोपस्करभेषजः ॥ प्रत्युत्पन्नमतिर्}\ H}}}{\lemma{\skt{द्धीमान्व्यवसायी}}\Afootnote{H; \skt{र्धीमान्व्यवसायी}\ A1938; \skt{धीमां व्यवसायी}\ N}}\skt{ विशारदः । }\pend
\pstart\versepar \skt{}\edtext{\skt{सत्यधर्म्मपरो यश्च सभिषक्पाद उच्यते}}{\Afootnote{\emph{om.}\ A1938 N H}}\skt{ ॥}\pend

\pstart\versepar \ledsidenote{\passref{1.34.21}}\edtext{\skt{आयुष्मान्सत्ववान्साध्यो}}{\Afootnote{\skt{सत्यधर्मपरो यश्च स भिषक्पाद उच्यते आयुष्मान्सत्त्ववान्साध्यो}\ A1938 H; \skt{सत्यधर्मपरो यश्च स भिव्याद उच्यते । आयुष्मान्सत्यवां साध्यो}\ N}}\skt{ }\edtext{\skt{द्रव्यवान्मित्रवानपि}}{\Afootnote{H; \skt{द्रव्यवानात्मवानपि ॥२१॥}\ A1938; \skt{द्रव्यवांमितिवान्नपि ।}\ N}}\skt{ । }\pend
\pstart\versepar \skt{}\edtext{\skt{वैद्यवाक्यकृदास्तिष्को व्याधितः पाद उच्यते}}{\Afootnote{\emph{om.}\ A1938 N H}}\skt{ ॥}\pend

\pstart\versepar \ledsidenote{\passref{1.34.22}}\edtext{\skt{प्रशस्तदेशसंभूतं}}{\Afootnote{\skt{आस्तिको वैद्यवाक्यस्थो व्याधितः पाद उच्यते प्रशस्तदेशसम्भूतं}\ A1938; \skt{वैद्यवाकृक्षदातिक्यो व्याधितः पाद उच्यते । प्रशस्तदेशसम्भूतं}\ N; \skt{वैद्यवाक्यकृदास्तिष्को व्याधितः पाद उच्यते ॥ प्रशस्तदेशसंभूतं}\ H}}\skt{ प्रशस्ते काल उद्धृतं । }\pend
\pstart\versepar \skt{युक्तमात्रम्मनस्कान्तङ्गन्धवर्ण्णरसान्वितः ॥}\edtext{}{\lemma{\skt{काल… गन्}}\Afootnote{\emph{om.}\ A1938}}\edtext{}{\lemma{\skt{काल… म}}\Afootnote{\emph{om.}\ N}}\edtext{}{\lemma{\skt{उद्धृतं । युक्तमात्रम्मनस्कान्तङ्गन्धवर्ण्णरसान्वितः}}\Afootnote{\skt{उद्धृतं ।}\ H}}\pend

\pstart\versepar \ledsidenote{\passref{1.34.23}}\edtext{\edtext{\skt{दोषघ्नमग्लानिनिकरमविकार्य}}{\lemma{\skt{दोषघ्नमग्लानिनिकरम्}}\Afootnote{\skt{युक्तमात्रं मनस्कान्तं गन्धवर्णरसान्वितम् । दोषघ्नमग्लानिकर}\ A1938; \skt{युक्तमात्रमनस्कान्तं गन्धवर्ण्णरसान्वितम् । दोषघ्नमग्लानिकरम्}\ N; \skt{युक्तमात्रम्मनस्कान्तङ्गन्धवर्ण्णरसान्वितः ॥ दोषघ्नमग्लानिनिकरम्}\ H}}}{\lemma{\skt{अविकार्य}}\Afootnote{H; \skt{मविकारि}\ A1938; \skt{अविकारि}\ N}}\skt{ }\edtext{\skt{विपर्ययं}}{\Afootnote{H; \skt{विपर्यये ।}\ A1938 N}}\skt{ । }\pend
\pstart\versepar \skt{समीक्ष्य }\edtext{\skt{दत्तङ्काले}}{\lemma{\skt{दत्तङ्}}\Afootnote{A1938 H; \skt{दन्त}\ N}}\skt{ च भेषजम्पाद उच्यते ॥}\pend

\pstart\versepar \ledsidenote{\passref{1.34.24}}\edtext{\edtext{\skt{स्निग्धोऽजुगुप्सुर्ब्बलवान्युक्तो}}{\lemma{\skt{ऽजुगुप्सुर्}}\Afootnote{H; \skt{ऽजुगुप्सु}\ A1938; \skt{जुगुप्सु}\ N}}}{\lemma{\skt{ब्बलवान्युक्तो}}\Afootnote{H; \skt{र्बलवान्युक्तो}\ A1938; \skt{बलवां युक्तो}\ N}}\skt{ व्याधितरक्षणे । }\pend
\pstart\versepar \skt{}\edtext{\skt{वैद्यवाक्यकृदश्रान्तः}}{\Afootnote{A1938 H; \skt{वैद्यवाक्यदशान्तः}\ N}}\skt{ पादः परिचरः स्मृतः ॥}\pend

\pstart \ledsidenote{\passref{1.34.trailer}}\edtext{\skt{लय्क}}{\Afootnote{H; \skt{इति सुश्रुतसंहितायां सूत्रस्थाने युक्तसेनीयो नाम चतुस्त्रिंशत्तमोऽध्यायः ॥३४॥}\ A1938; \skt{लप्क ॥}\ N}}\skt{ ॥}\pend

\endnumbering
\beginnumbering
\eledsection*{Adhyāya 35}
\pstart \ledsidenote{\passref{1.35.1}}\skt{अथात }\edtext{\skt{आतुरोपक्रमणीयम्व्याख्यास्यामः}}{\lemma{\skt{आतुरोपक्रमणीयम्}}\Afootnote{\skt{आतुरोपक्रमणीयमध्यायं}\ A1938; \skt{आतुरोप्रक्रमणीयमध्यायं}\ N; \skt{आतुरोपक्रमनीयम्}\ H}}\skt{ ॥}\pend

\pstart \ledsidenote{\passref{1.35.3}}\edtext{\skt{आतुरमुपक्रममाणेन}}{\lemma{\skt{उपक्रममाणेन}}\Afootnote{A1938 H; \skt{उपक्रममानेन}\ N}}\skt{ }\edtext{\edtext{\skt{वैद्येनायुरादावेव}}{\lemma{\skt{वैद्येनायुरादाव्}}\Afootnote{N H; \skt{भिषजाऽऽयुरादा}\ A1938}}}{\lemma{\skt{एव}}\Afootnote{N H; \skt{वेव}\ A1938}}\skt{ }\edtext{\skt{तावत्परीसत्यायुषि}}{\Afootnote{H; \skt{परीक्षितव्यं ; सत्यायुषि}\ A1938; \skt{तावत्परीक्ष्यं सत्यायुषि}\ N}}\skt{ }\edtext{\skt{व्याध्यृत्वग्निवयोदेहसत्ववलसात्म्यप्रकृतिभेषजदेशान्परीक्षेत}}{\Afootnote{H; \skt{व्याध्यृत्वग्निवयोदेहबलसत्त्वसात्म्यप्रकृतिभेषजदेशान्परीक्षेत ॥३॥}\ A1938; \skt{व्याध्यृन्ववयोदेहसत्वबलसात्म्यप्रकृतिभेषजं पाद उच्यते । क्षेत ।}\ N}}\skt{ ।}\pend

\pstart \ledsidenote{\passref{1.35.4}}\skt{तत्र }\edtext{\edtext{\edtext{\skt{महापाणिपादपृष्ठस्तनाग्रसुखदर्शनस्कन्धललाटदीर्घाङ्गुलिपर्व्वोच्छ्वासप्रेक्षणवाहुम्विस्तीर्ण्णभ्रूस्तनान्तरोरस्कं}}{\lemma{\skt{महापाणिपादपृष्ठस्तनाग्रसुखदर्शनस्कन्धललाटदीर्घाङ्गुलिपर्व्वोच्छ्वासप्रेक्षणवाहुम्}}\Afootnote{H; \skt{महापाणिपादपार्श्वपृष्ठस्तनाग्रदशनवदनस्कन्धललाटं दीर्घाङ्गुलिपर्वोच्छ्वासप्रेक्षणबाहुं,}\ A1938; \skt{महापाणिपादपृष्ठस्तनाग्रदशनस्कन्धललाट दीर्घाङ्गुलिपर्वौच्छ्वासप्रेक्षणबाहुं}\ N}}}{\lemma{\skt{विस्तीर्ण्णभ्रूस्}}\Afootnote{H; \skt{विस्तीर्णभ्रू}\ A1938; \skt{विस्तीर्ण्णवाभ्रू}\ N}}}{\lemma{\skt{तनान्तरोरस्कं}}\Afootnote{H; \skt{स्तनान्तरोरस्कं}\ A1938 N}}\skt{ }\edtext{\edtext{\edtext{\edtext{\skt{हृस्वजङ्घामेढ्रग्रीवङ्गम्भीरसत्वस्वरनासिमनुच्चैर्वद्धस्तन्मुपचितमहारोमं}}{\lemma{\skt{हृस्वजङ्घामेढ्रग्रीवङ्}}\Afootnote{H; \skt{ह्रस्वजङ्घामेढ्रग्रीवं}\ A1938 N}}}{\lemma{\skt{गम्भीरसत्वस्वरनासिमनुच्चैर्}}\Afootnote{\skt{गम्भीरसत्त्वस्वरनाभिमनुच्चैर्}\ A1938 N; \skt{गम्भीरसत्वस्व रनासिमनुच्चैर्}\ H}}}{\lemma{\skt{वद्धस्तन्म्}}\Afootnote{\skt{बद्धस्तनम्}\ A1938 N; \skt{वद्धस्तन्म्}\ H}}}{\lemma{\skt{उपचितमहारोमं}}\Afootnote{H; \skt{उपचितमहा}\ A1938; \skt{उपचितमहारोग}\ N}}\skt{ }\edtext{\skt{रोमसकर्ण्णे}}{\Afootnote{H; \skt{रोमशकर्णं,}\ A1938; \skt{रोमसकर्ण्ण}\ N}}\skt{ }\edtext{\skt{पार्श्व}}{\Afootnote{N H; \skt{पश्चा}\ A1938}}\skt{ }\edtext{\skt{मस्तकं}}{\Afootnote{N H; \skt{न्मस्तिष्कं,}\ A1938}}\skt{ }\edtext{\skt{स्नातानुलिप्तं}}{\Afootnote{A1938 H; \skt{स्नातानुप्तिप्न}\ N}}\skt{ }\edtext{\edtext{\skt{मूर्द्धमानूपूर्व्यो}}{\lemma{\skt{मूर्द्धम्}}\Afootnote{H; \skt{मूर्धानुपूर्व्या विशुष्य}\ A1938; \skt{मूध्नानुपूर्व्या}\ N}}}{\lemma{\skt{आनूपूर्व्यो}}\Afootnote{H; \skt{माणशरीरं}\ A1938; \emph{om.}\ N}}\skt{ पश्चाच्च }\edtext{\skt{विशुष्यमाणहृदयम्पुरुषञ्जानीयाद्दीर्घायुः}}{\lemma{\skt{जानीयाद्दीर्घायुः}}\Afootnote{H; \skt{जानीयाद्दीर्घायुः}\ A1938; \skt{जानीयाद्दीर्घायु}\ N}}\skt{ खल्वयमिति । }\edtext{\skt{तमेकान्तेन}}{\lemma{\skt{एकान्तेन}}\Afootnote{H; \skt{एकान्तेनो}\ A1938; \skt{एकान्तनो}\ N}}\skt{ }\edtext{\skt{उपक्रमेत् । तत्र}}{\Afootnote{H; \skt{पक्रमेत् । एभिर्}\ A1938; \skt{पक्रमेत् । तत्रे}\ N}}\skt{ }\edtext{\edtext{\edtext{\edtext{\edtext{\skt{तैल्लैक्षणैर्विपरीतैरल्पयुम्मिश्रैर्म्मध्यमायुरिति}}{\lemma{\skt{तैल्लैक्षणैर्}}\Afootnote{H; \skt{लक्षणैर्}\ A1938; \skt{तैलक्षणैर्}\ N}}}{\lemma{\skt{अल्पयुम्}}\Afootnote{H; \skt{अल्पायुः}\ A1938 N}}}{\lemma{\skt{म्मध्यम्}}\Afootnote{N H; \skt{मध्य}\ A1938}}}{\lemma{\skt{आयुर्}}\Afootnote{N H; \skt{मायुर्}\ A1938}}}{\lemma{\skt{इति}}\Afootnote{A1938 H; \emph{om.}\ N}}\skt{ ॥}\pend

\pstart\versepar \ledsidenote{\passref{1.35.5}}\edtext{\skt{भवन्ति}}{\Afootnote{A1938 N; \emph{om.}\ H}}\skt{ }\edtext{\skt{चात्र}}{\Afootnote{A1938 N; \emph{om.}\ H}}\skt{ ॥ }\pend
\pstart\versepar \skt{}\edtext{\skt{गूढसन्धिसिरास्नायुः}}{\Afootnote{A1938 H; \skt{मूढसन्धिशिरास्नायुः}\ N}}\skt{ }\edtext{\skt{संहताङ्गः}}{\Afootnote{A1938 H; \skt{संहताङ्ग}\ N}}\skt{ स्थिरेन्द्रियः । }\pend
\pstart\versepar \skt{उत्तरोत्तरसुक्षेत्रो यः स दीर्घायुरुच्यते ॥}\pend

\pstart\versepar \ledsidenote{\passref{1.35.6}}\edtext{\skt{गर्भात्प्रभृत्यरोगो यः}}{\Afootnote{A1938 H; \skt{गर्भात्प्रभृत्यरोगोयः}\ N}}\skt{ शनैः }\edtext{\skt{समभिवर्द्धते}}{\Afootnote{N H; \skt{समुपचीयते ।}\ A1938}}\skt{ । }\pend
\pstart\versepar \skt{शरीरज्ञानविज्ञानैः स दीर्घायुः }\edtext{\skt{समासतः}}{\Afootnote{A1938 H; \skt{समागतः ।}\ N}}\skt{ ॥}\pend

\pstart\versepar \ledsidenote{\passref{1.35.7}}\edtext{\skt{मध्यमस्यायुषो}}{\Afootnote{A1938 H; \skt{मध्यमयुषो}\ N}}\skt{ ज्ञान मत }\edtext{\edtext{\skt{ऊर्ध्वन्निबोध}}{\lemma{\skt{ऊर्ध्वन्}}\Afootnote{A1938 H; \skt{मूर्ध्व}\ N}}\skt{ मे}}{\lemma{\skt{निबोध मे}}\Afootnote{A1938 N; \skt{निवोध मे ।}\ H}}\skt{ । }\pend
\pstart\versepar \skt{}\edtext{\edtext{\skt{अधस्तादक्षयोर्यस्य}}{\lemma{\skt{अधस्तादक्षयोर्}}\Afootnote{A1938; \skt{अधस्ताद्वक्षयोर्}\ N; \skt{अधस्ता द्वक्ष द्रस्त यो}\ H}}}{\lemma{\skt{यस्य}}\Afootnote{A1938 N; \skt{त्यस्य}\ H}}\skt{ }\edtext{\skt{रेखाः}}{\Afootnote{N H; \skt{लेखाः}\ A1938}}\skt{ }\edtext{\skt{सुव्यक्तमायताः}}{\Afootnote{\skt{स्युर्व्यक्तमायताः ॥७॥}\ A1938; \skt{सुव्यक्तमायताः}\ N H}}\skt{ ॥}\pend

\pstart\versepar \ledsidenote{\passref{1.35.8}}\skt{द्वे वा }\edtext{\edtext{\skt{तिस्रोऽधिका}}{\lemma{\skt{तिस्रो}}\Afootnote{A1938 H; \skt{तिस्रो}\ N}}}{\lemma{\skt{ऽधिका}}\Afootnote{A1938 H; \skt{धिका}\ N}}\skt{ वापि पादौ कर्णौ च मांसलौ । }\pend
\pstart\versepar \skt{}\edtext{\edtext{\skt{नासाग्रमूर्ध्वञ्च}}{\lemma{\skt{नासाग्रम्}}\Afootnote{A1938 H; \skt{नासाग्र}\ N}}}{\lemma{\skt{ऊर्ध्वञ्च}}\Afootnote{A1938 H; \skt{मूर्द्वञ्च}\ N}}\skt{ }\edtext{\skt{भवेत्रेखाश्च}}{\Afootnote{H; \skt{भवेदूर्ध्वं लेखाश्च}\ A1938; \skt{भवेद्रेखा व्यक्ताश्च}\ N}}\skt{ पृष्ठतः ॥}\pend

\pstart\versepar \ledsidenote{\passref{1.35.9}}\skt{यस्य स्युस्तस्य }\edtext{\skt{विज्ञेयमायुर्भवति}}{\lemma{\skt{विज्ञेयम्}}\Afootnote{N H; \skt{परमम्}\ A1938}}\skt{ सप्ततिः । }\pend
\pstart\versepar \skt{}\edtext{\edtext{\skt{जघन्यश्चायुषो}}{\lemma{\skt{जघन्यश्}}\Afootnote{H; \skt{जघन्य्}\ A1938; \skt{जघन्य}\ N}}}{\lemma{\skt{चायुषो}}\Afootnote{H; \skt{अस्यायुषो}\ A1938; \skt{स्वायुषो}\ N}}\skt{ }\edtext{\skt{ज्ञान}}{\Afootnote{N H; \skt{ज्ञानम्}\ A1938}}\skt{ }\edtext{\skt{मत}}{\Afootnote{N H; \skt{अत}\ A1938}}\skt{ }\edtext{\edtext{\skt{ऊर्ध्वन्निबोध}}{\lemma{\skt{ऊर्ध्वन्}}\Afootnote{A1938 H; \skt{मूर्द्व}\ N}}}{\lemma{\skt{निबोध}}\Afootnote{A1938 N; \skt{निवोध}\ H}}\skt{ मे ॥}\pend

\pstart\versepar \ledsidenote{\passref{1.35.10}}\edtext{\skt{हृस्वानि}}{\Afootnote{H; \skt{ह्रस्वानि}\ A1938 N}}\skt{ यस्य }\edtext{\skt{पर्व्वाणि}}{\Afootnote{A1938 H; \skt{सर्वाणि}\ N}}\skt{ सुमहच्चापि मेहनं । }\pend
\pstart\versepar \skt{}\edtext{\skt{अवलीढमुरोयश्च न च}}{\Afootnote{\skt{तथोरस्यवलीढानि न च}\ A1938; \skt{अवलीढमुरोयश्च न च}\ N H}}\skt{ }\edtext{\edtext{\skt{स्यात्पृष्ठमायतं}}{\lemma{\skt{स्यात्पृष्ठम्}}\Afootnote{A1938 H; \skt{स्यात्पृष्ठ}\ N}}}{\lemma{\skt{आयतं}}\Afootnote{A1938 H; \skt{मायतम् ।}\ N}}\skt{ ॥}\pend

\pstart\versepar \ledsidenote{\passref{1.35.11}}\edtext{\skt{ऊर्ध्वञ्च}}{\Afootnote{A1938 H; \skt{उर्द्वञ्च}\ N}}\skt{ }\edtext{\skt{श्रवणौ}}{\Afootnote{A1938 H; \skt{श्रवणो}\ N}}\skt{ स्थानान्नासा }\edtext{\skt{चोच्चा}}{\Afootnote{A1938 H; \skt{च्चोचा}\ N}}\skt{ शरीरिणः । }\pend
\pstart\versepar \skt{हसतो जल्पतो वापि दन्तमान्सम्प्रदृश्यते । }\pend
\pstart\versepar \skt{प्रेक्षते यश्च }\edtext{\skt{विभ्रान्तः स}}{\Afootnote{N H; \skt{विभ्रान्तं स}\ A1938}}\skt{ }\edtext{\skt{जीवेत्पञ्चविङ्शतिः}}{\Afootnote{H; \skt{जीवेत्पञ्चविंशतिम् ॥११॥}\ A1938; \skt{जीवेत्पञ्चविङ्शति ॥}\ N}}\skt{ ॥}\pend

\pstart \ledsidenote{\passref{1.35.12ab}}\skt{अथ }\edtext{\edtext{\skt{पुनरायुषो}}{\lemma{\skt{पुनर्}}\Afootnote{N H; \skt{पुन}\ A1938}}}{\lemma{\skt{आयुषो}}\Afootnote{N H; \skt{रायुषो}\ A1938}}\skt{ }\edtext{\edtext{\skt{विज्ञानार्थमङ्गप्रत्यङ्गप्रमाणसार}}{\lemma{\skt{विज्ञानार्थम्}}\Afootnote{A1938 H; \skt{विज्ञानार्थाम्}\ N}}}{\lemma{\skt{अङ्गप्रत्यङ्गप्रमाणसार}}\Afootnote{H; \skt{अङ्गप्रत्यङ्गप्रमाणसारा}\ A1938 N}}\skt{ }\edtext{\skt{अनुपदेक्ष्यामः}}{\Afootnote{H; \skt{नुपदेक्ष्यामः}\ A1938 N}}\skt{ ॥ }\edtext{\skt{तत्र}}{\Afootnote{H; \skt{तत्राङ्}\ A1938 N}}\skt{ }\edtext{\skt{अङ्गात्यन्तराधिशक्थिवाहुशिरांसि}}{\Afootnote{H; \skt{गान्यन्तराधिसक्थिबाहुशिरांसि,}\ A1938; \skt{गास्वन्तराधिसक्थिबाहुशिरांसि}\ N}}\skt{ }\edtext{\skt{तदवयवाः}}{\Afootnote{A1938 H; \skt{तदवयवाः ।}\ N}}\skt{ }\edtext{\skt{प्रत्यङ्गानि}}{\Afootnote{H; \skt{प्रत्यङ्गानीति ।}\ A1938; \skt{प्रत्याङ्गानि}\ N}}\skt{ }\edtext{\skt{स्वैरङ्गुलैः}}{\lemma{\skt{स्वैर्}}\Afootnote{N H; \skt{तत्र, स्वैर्}\ A1938}}\skt{ पादाङ्गुष्ठप्रदेशिन्यौ }\edtext{\skt{द्वङ्गुलायते}}{\Afootnote{H; \skt{द्व्यङ्गुलायते}\ A1938 N}}\skt{ । }\edtext{\skt{प्रदेशिन्यास्तु}}{\Afootnote{N H\textsuperscript{pc}; \skt{प्रदेशिन्यास्तु}\ A1938; \skt{प्रदेन्यास्तु}\ H\textsuperscript{ac}}}\skt{ }\edtext{\skt{मध्यमानामिकाकनिष्ठिका}}{\Afootnote{A1938 H; \skt{मध्यामानामिका कनिष्ठका}\ N}}\skt{ }\edtext{\skt{पूर्व्वतः}}{\Afootnote{H; \skt{यथोत्तरं}\ A1938; \skt{पुरतः}\ N}}\skt{ । }\edtext{\skt{पञ्चञ्भागहीनाचतुरङ्गुलायते}}{\Afootnote{H; \skt{पञ्चमभागहीनाः; चतुरङ्गुलायते}\ A1938; \skt{पञ्चभागहीनाचतुरङ्गुलायते}\ N}}\skt{ । }\edtext{\skt{पञ्चङ्गुलविस्मृते}}{\Afootnote{H; \skt{पञ्चाङ्गुलविस्तृते}\ A1938; \skt{षडङ्गुलेविस्तृस्तृते ।}\ N}}\skt{ ।}\pend

\pstart \ledsidenote{\passref{1.35.16cd}}\edtext{\skt{स्निग्धताम्रनखनयनजिहोष्ठपाणिपादतलं}}{\Afootnote{H; \skt{स्निग्धताम्रनखनयनतालुजिह्वौष्ठपाणिपादतलं}\ A1938; \skt{स्निग्धताम्रनखनयजिह्वौष्ठपाणिपादतलं}\ N}}\skt{ । }\edtext{\skt{रक्त्नेन}}{\Afootnote{H; \skt{रक्तेन}\ A1938 N}}\skt{ सुप्रसन्न मृदुत्वग्रोमाणं }\edtext{\skt{त्वषां}}{\Afootnote{H; \skt{त्वक्सारं विद्यादिति । एषां पूर्वं}\ A1938; \skt{त्ववत्येषां}\ N}}\skt{ }\edtext{\skt{पूर्व्व}}{\Afootnote{N H; \skt{पूर्वं}\ A1938}}\skt{ }\edtext{\skt{प्राधान्यादायुः}}{\Afootnote{H; \skt{प्रधानमायुः}\ A1938; \skt{प्राधान्यादायुः}\ N}}\skt{ }\edtext{\skt{सौभाग्यापयोगायेति}}{\Afootnote{H; \skt{सौभाग्ययोरिति ॥१६॥}\ A1938; \skt{सौभाग्योयोगायेति ।}\ N}}\skt{ ॥}\pend

\pstart\versepar \ledsidenote{\passref{1.35.17}}\edtext{\edtext{\skt{सामान्यतोऽङ्गप्रत्यङ्गप्रमाणादथ}}{\lemma{\skt{सामान्यतो}}\Afootnote{N H; \skt{विशेषतो}\ A1938}}}{\lemma{\skt{ऽङ्गप्रत्यङ्गप्रमाणादथ}}\Afootnote{A1938 H; \skt{ङ्गप्रत्यङ्गाप्रमाणादम}\ N}}\skt{ सारतः । }\pend
\pstart\versepar \skt{परीक्ष्यायुः }\edtext{\skt{सुनिपुणो}}{\Afootnote{A1938 H; \skt{सुनिपुणौ}\ N}}\skt{ }\edtext{\skt{भिषक्सिध्यति}}{\Afootnote{A1938 H; \skt{भिषग्भिध्यति}\ N}}\skt{ कर्म्मसु ॥}\pend

\pstart \ledsidenote{\passref{1.35.18}}\skt{व्याधिविशेषास्तु }\edtext{\skt{प्रागभिहिताः}}{\Afootnote{A1938 H; \skt{प्रागविहिताः ।}\ N}}\skt{ }\edtext{\skt{सर्व्वं}}{\Afootnote{H; \skt{सर्व}\ A1938 N}}\skt{ }\edtext{\skt{एवैते}}{\Afootnote{A1938 H; \skt{एव एते}\ N}}\skt{ }\edtext{\skt{त्रिविधा}}{\Afootnote{N H; \skt{त्रिविधाः-}\ A1938}}\skt{ }\edtext{\skt{भवन्ति}}{\Afootnote{N H; \emph{om.}\ A1938}}\skt{ । साध्या याप्याः }\edtext{\edtext{\skt{प्रत्याख्येयाश्चेति}}{\lemma{\skt{प्रत्याख्येयाश्}}\Afootnote{A1938 H; \skt{प्रत्याख्येया}\ N}}}{\lemma{\skt{चेति}}\Afootnote{H; \emph{om.}\ A1938; \skt{श्चेति ।}\ N}}\skt{ । }\edtext{\skt{तत्र}}{\Afootnote{N H; \skt{च ।}\ A1938}}\skt{ }\edtext{\skt{साध्य}}{\Afootnote{N H; \emph{om.}\ A1938}}\skt{ }\edtext{\skt{याप्य}}{\Afootnote{H; \skt{त}\ A1938; \skt{याप्य}\ N}}\skt{ }\edtext{\skt{प्रत्याख्ये}}{\Afootnote{N H; \skt{त्रै}\ A1938}}\skt{ }\edtext{\skt{यान्व्याधीन्भूयत्रि}}{\Afootnote{H; \skt{तान्भूयस्त्रि}\ A1938; \skt{या व्याधीभूयस्त्रि}\ N}}\skt{ }\edtext{\skt{विधा}}{\Afootnote{N H; \skt{धा}\ A1938}}\skt{ }\edtext{\skt{परीक्षेत}}{\Afootnote{A1938 H; \skt{परीक्षते ।}\ N}}\skt{ । }\edtext{\edtext{\skt{किमयमौपसर्गिकः}}{\lemma{\skt{किमयम्}}\Afootnote{H; \skt{किमय}\ A1938; \skt{किमयम्}\ N}}}{\lemma{\skt{औपसर्गिकः}}\Afootnote{H; \skt{मौपसर्गिकः,}\ A1938; \skt{औषधसगिकः ।}\ N}}\skt{ }\edtext{\edtext{\skt{प्राक्केवलोऽन्यलक्षण}}{\lemma{\skt{प्राक्केवलो}}\Afootnote{A1938 H; \skt{प्राक्कवलो}\ N}}}{\lemma{\skt{ऽन्यलक्षण}}\Afootnote{A1938 H; \skt{न्यलक्षण}\ N}}\skt{ इति । }\edtext{\skt{तरौपसर्गिको}}{\Afootnote{H; \skt{तत्र, औपसर्गिको}\ A1938; \skt{तत्रौपसर्गिको}\ N}}\skt{ नाम यः }\edtext{\skt{पूर्व्वोत्पन्नं}}{\Afootnote{A1938 H; \skt{पूर्वोत्पन्न}\ N}}\skt{ }\edtext{\edtext{\skt{व्याधिञ्जघन्यकालजातो}}{\lemma{\skt{व्याधिञ्}}\Afootnote{A1938 H; \skt{व्याधि}\ N}}}{\lemma{\skt{जघन्यकालजातो}}\Afootnote{A1938 H; \skt{ज पुन्यकालजातो}\ N}}\skt{ }\edtext{\edtext{\skt{व्याधिरुपसृजति}}{\lemma{\skt{व्याधिर्}}\Afootnote{N H; \skt{व्याधि}\ A1938}}\skt{ । स}}{\lemma{\skt{उपसृजति । स}}\Afootnote{N H; \skt{रुपसृजति, स}\ A1938}}\skt{ }\edtext{\skt{तत्मूले स}}{\Afootnote{H; \skt{तन्मूल}\ A1938 N}}\skt{ }\edtext{\skt{एवोपद्रवसंज्ञः}}{\Afootnote{A1938 H; \skt{एवोपद्रवसज्ञः ।}\ N}}\skt{ । }\edtext{\skt{प्राक्केवलो}}{\Afootnote{A1938 H; \skt{प्राक्कवलो}\ N}}\skt{ }\edtext{\skt{नाम । यः}}{\Afootnote{A1938 H; \skt{नामः यः}\ N}}\skt{ }\edtext{\skt{प्रागेवोत्पन्नो}}{\Afootnote{A1938 H; \skt{प्रागेवोपन्नो}\ N}}\skt{ }\edtext{\edtext{\skt{व्याधिररपूर्वरूपोपद्रवश्च}}{\lemma{\skt{व्याधिरर्}}\Afootnote{H; \skt{व्याधिर}\ A1938; \skt{व्याधिर्}\ N}}}{\lemma{\skt{अपूर्वरूपोपद्रवश्च}}\Afootnote{N H; \skt{पूर्वरूपोऽनुपद्रवश्च;}\ A1938}}\skt{ । अन्य लक्षणो नाम यो }\edtext{\skt{भविष्यत्ख्यापकः स}}{\Afootnote{H; \skt{भविष्यद्व्याधिख्यापकः, स}\ A1938; \skt{भविष्यतख्यापकः । स}\ N}}\skt{ पूर्व्वरूपसंज्ञः । तत्र }\edtext{\edtext{\skt{सोपद्रवमन्योन्याविरोधेनोपचरेत्}}{\lemma{\skt{सोपद्रवम्}}\Afootnote{H; \skt{सोपद्रव}\ A1938; \skt{सोपद्रवसोप}\ N}}\skt{ । वलवन्तरमुपद्रवम्वा}}{\lemma{\skt{अन्योन्याविरोधेनोपचरेत् । वलवन्तरम्}}\Afootnote{H; \skt{मन्योन्याविरोधेनोपक्रमेत, बलवन्तम्}\ A1938; \skt{द्रमन्योन्याविरोधेनोपचरेत् । बलवत्तरम्}\ N}}\skt{ }\edtext{\skt{प्राक्केवलं}}{\Afootnote{A1938 H; \skt{प्राक्केवल}\ N}}\skt{ यथास्वम्प्रतिकुर्वीत । }\edtext{\skt{अन्ये}}{\Afootnote{H; \skt{अन्य}\ A1938 N}}\skt{ }\edtext{\skt{लक्षणे}}{\Afootnote{A1938 H; \skt{लक्षण}\ N}}\skt{ }\edtext{\skt{त्वादिव्याधौ}}{\Afootnote{H; \skt{त्वादिव्याधौ}\ A1938; \skt{त्वादिव्याधो}\ N}}\skt{ }\edtext{\skt{यथा}}{\Afootnote{N H; \emph{om.}\ A1938}}\skt{ }\edtext{\skt{यथावत्प्रयतेत}}{\Afootnote{N H; \skt{प्रयतेत ॥१८॥}\ A1938}}\skt{ ।}\pend

\pstart\versepar \ledsidenote{\passref{1.35.19}}\edtext{\skt{भवति}}{\Afootnote{A1938 H; \skt{भवन्ति}\ N}}\skt{ }\edtext{\skt{चात्र}}{\Afootnote{A1938 N; \emph{om.}\ H}}\skt{ । }\pend
\pstart\versepar \skt{नास्ति रोगो विना }\edtext{\skt{दोषं}}{\Afootnote{N H; \skt{दोषैर्}\ A1938}}\skt{ }\edtext{\skt{यस्मात्तस्माच्चिकित्सकः}}{\Afootnote{H; \skt{यस्मात्तस्माद्विचक्षणः ।}\ A1938; \skt{यस्मात्तस्माचिकित्सकः ।}\ N}}\skt{ । }\pend
\pstart\versepar \skt{अनुक्तमपि दोषाणां }\edtext{\edtext{\skt{लिङ्गैर्दोषमुपाचरेत्}}{\lemma{\skt{दोषम्}}\Afootnote{N; \skt{व्याधिम्}\ A1938; \skt{दोष}\ H}}}{\lemma{\skt{उपाचरेत्}}\Afootnote{A1938 N; \skt{मुपाचरेत् ॥}\ H}}\skt{ ॥}\pend

\pstart \ledsidenote{\passref{1.35.20}}\edtext{\skt{प्रागभिहिता}}{\Afootnote{A1938 H; \skt{प्रागविहिता}\ N}}\skt{ ऋतवः ।}\pend

\pstart\versepar \ledsidenote{\passref{1.35.21}}\edtext{\skt{शीत}}{\Afootnote{H; \skt{शीते}\ A1938 N}}\skt{ शीतप्रतीकारमुष्णे चोष्णनिवारणं । }\pend
\pstart\versepar \skt{}\edtext{\edtext{\skt{कालप्राप्ताङ्क्रियां}}{\lemma{\skt{कालप्राप्ताङ्}}\Afootnote{H; \skt{कृ}\ A1938; \skt{कालप्राप्तं}\ N}}}{\lemma{\skt{क्रियां}}\Afootnote{N H; \skt{त्वा}\ A1938}}\skt{ }\edtext{\skt{कुर्यात्क्रियाकालन्न}}{\Afootnote{H; \skt{कुर्यात्क्रियां प्राप्तां क्रियाकालं न}\ A1938; \skt{कुर्यात्क्रिया कालन}\ N}}\skt{ हापयेत् ॥}\pend

\pstart\versepar \ledsidenote{\passref{1.35.22}}\edtext{\skt{अपाप्ते वा}}{\Afootnote{H; \skt{अप्राप्ते वा}\ A1938 N}}\skt{ }\edtext{\skt{क्रियाकाले}}{\Afootnote{A1938 H; \skt{क्रियाप्राप्ते}\ N}}\skt{ प्राप्ते वा न कृता }\edtext{\skt{क्रिया}}{\Afootnote{N H; \skt{क्रिया । क्रिया}\ A1938}}\skt{ । }\pend
\pstart\versepar \skt{}\edtext{\skt{हीनातिरिक्ता च}}{\Afootnote{H; \skt{हीनाऽतिरिक्ता}\ A1938; \skt{हीनातिक्ता च}\ N}}\skt{ }\edtext{\skt{कृता}}{\Afootnote{N H; \skt{वा}\ A1938}}\skt{ }\edtext{\edtext{\skt{साद्येष्वपि}}{\lemma{\skt{साद्येष्व्}}\Afootnote{H; \skt{साध्येष्व्}\ A1938; \skt{साध्ये}\ N}}\skt{ न}}{\lemma{\skt{अपि न}}\Afootnote{A1938 H; \skt{ष्वपि न}\ N}}\skt{ सिध्यति ॥}\pend

\pstart \ledsidenote{\passref{1.35.24}}\edtext{\edtext{\edtext{\skt{प्रागभिहितोऽग्निरन्नाभिपाचकः}}{\lemma{\skt{प्रागभिहितो}}\Afootnote{A1938 H; \skt{प्राग्विविहितो}\ N}}}{\lemma{\skt{ऽग्निर्}}\Afootnote{A1938 H; \skt{ग्नि}\ N}}\skt{ स}}{\lemma{\skt{अन्नाभिपाचकः स}}\Afootnote{H; \skt{अन्नस्य पाचकः । स}\ A1938; \skt{रन्नाग्निपाचकः । स}\ N}}\skt{ }\edtext{\skt{पञ्चविधो}}{\Afootnote{N H; \skt{चतुर्विधो}\ A1938}}\skt{ भवति ॥ }\edtext{\skt{समो}}{\Afootnote{N H; \emph{om.}\ A1938}}\skt{ }\edtext{\skt{दोषाभिपन्नो}}{\Afootnote{N H; \skt{दोषानभिपन्न एकः,}\ A1938}}\skt{ }\edtext{\skt{विक्रियामापन्न}}{\lemma{\skt{आपन्न}}\Afootnote{N H; \skt{आपन्नस्}\ A1938}}\skt{ }\edtext{\skt{इति}}{\Afootnote{N H; \skt{त्रिविधो भवति-}\ A1938}}\skt{ ॥ विषमो वातेन । }\edtext{\skt{तीक्ष्ण}}{\Afootnote{H; \skt{तीक्ष्णः}\ A1938 N}}\skt{ पित्तेन । मन्दः श्लेष्मणा । }\edtext{\skt{समः}}{\Afootnote{N H; \skt{चतुर्थः समः}\ A1938}}\skt{ }\edtext{\skt{सर्व्वैः}}{\Afootnote{H; \skt{सर्व}\ A1938 N}}\skt{ }\edtext{\edtext{\skt{समरिति}}{\lemma{\skt{समर्}}\Afootnote{H; \skt{साम्या}\ A1938; \skt{समैर्}\ N}}}{\lemma{\skt{इति}}\Afootnote{N H; \skt{दिति ।}\ A1938}}\skt{ ॥ तत्र यो यथा }\edtext{\edtext{\skt{कालमन्नमुपयुक्तं}}{\lemma{\skt{अन्नम्}}\Afootnote{N H; \emph{om.}\ A1938}}}{\lemma{\skt{उपयुक्तं}}\Afootnote{N H; \skt{उपयुक्त}\ A1938}}\skt{ }\edtext{\skt{सम्य}}{\Afootnote{N H; \skt{मन्नं सम्य}\ A1938}}\skt{ }\edtext{\skt{पचति स}}{\Afootnote{H; \skt{क्पचति स}\ A1938 N}}\skt{ समः । }\edtext{\skt{यस्तु}}{\Afootnote{N H; \skt{समैर्दोषैः; यः}\ A1938}}\skt{ कदाचित्सम्यक्पचति }\edtext{\edtext{\skt{कदाचिदसम्यक्पचत्याध्मानशूलातीसारप्रवाहणानि}}{\lemma{\skt{कदाचिदसम्यक्पचत्य्}}\Afootnote{H; \skt{कदाचि}\ A1938; \skt{कदाचित्स}\ N}}}{\lemma{\skt{आध्मानशूलातीसारप्रवाहणानि}}\Afootnote{H; \skt{दाध्मानशूलोदावर्तातिसारजठरगौरवान्त्रकूजनप्रवाहणानि}\ A1938; \skt{म्यम्कचत्याध्मानसूलातीसारप्रव्रहणानि}\ N}}\skt{ कृत्वा }\edtext{\skt{पचति स}}{\Afootnote{N H; \skt{स}\ A1938}}\skt{ }\edtext{\skt{विषमः । यस्तु}}{\Afootnote{N H; \skt{विषमः; यः}\ A1938}}\skt{ }\edtext{\edtext{\edtext{\skt{प्रभूतमप्यन्नमुपयुक्तमाशु}}{\lemma{\skt{अप्य्}}\Afootnote{A1938 N; \emph{om.}\ H}}}{\lemma{\skt{अन्नम्}}\Afootnote{N; \emph{om.}\ A1938; \skt{प्यन्नम्}\ H}}}{\lemma{\skt{आशु}}\Afootnote{N H; \skt{अन्नमाशु}\ A1938}}\skt{ पचति स तीक्ष्णः । }\edtext{\edtext{\edtext{\skt{सर्वाभिवर्धमानोऽत्यग्निर्भवति}}{\lemma{\skt{सर्वाभिवर्धमानो}}\Afootnote{H; \skt{स एवाभिवर्धमानो}\ A1938 N}}}{\lemma{\skt{ऽत्यग्निर्}}\Afootnote{H; \skt{ऽत्यग्नि}\ A1938 N}}}{\lemma{\skt{भवति}}\Afootnote{N H; \skt{रित्या}\ A1938}}\skt{ ॥ }\edtext{\edtext{\skt{मुहुम्र्म्मुहुः}}{\lemma{\skt{मुहुम्र्}}\Afootnote{H; \skt{भाष्यते, स मुहुर्}\ A1938; \skt{मुहुर्मुः ॥}\ N}}}{\lemma{\skt{म्मुहुः}}\Afootnote{A1938 H; \emph{om.}\ N}}\skt{ }\edtext{\edtext{\edtext{\edtext{\skt{प्रभूततरमन्नमुपयुक्तमाशुतरं}}{\lemma{\skt{प्रभूततरम्}}\Afootnote{N H; \skt{प्रभूतम्}\ A1938}}}{\lemma{\skt{अन्नम्}}\Afootnote{N H; \skt{अप्य्}\ A1938}}}{\lemma{\skt{उपयुक्तम्}}\Afootnote{A1938 H; \skt{उपयुक्ताम्}\ N}}}{\lemma{\skt{आशुतरं}}\Afootnote{N H; \skt{अन्नमाशुतरं}\ A1938}}\skt{ पचति पाकान्ते च }\edtext{\skt{गलताल्वोष्ठप्रशोषदाहसन्तापान्जनयति}}{\Afootnote{H; \skt{गलताल्वोष्ठशोषदाहसन्तापा ञ्जनय}\ A1938; \skt{गलताल्वोष्ठप्रशोषदाहसन्तापां जनयत्य्}\ N}}\skt{ }\edtext{\skt{अस्यैव}}{\Afootnote{H; \skt{अस्त्येव}\ N}}\skt{ भस्मक इति }\edtext{\skt{व्यपदिशन्ति ॥ }\edtext{\edtext{\edtext{\edtext{\edtext{\skt{यस्त्वल्पमप्यन्नमुपयुक्तमुदरशिरोगौरवकासश्वासच्छर्दिप्रसेकगात्रसदनानि}}{\lemma{\skt{त्वल्पम्}}\Afootnote{H; \skt{स्त्वल्प}\ A1938; \skt{स्त्वन्नम्}\ N}}}{\lemma{\skt{अप्य्}}\Afootnote{H; \skt{मप्यु}\ A1938; \skt{अ}\ N}}}{\lemma{\skt{अन्नम्}}\Afootnote{H; \emph{om.}\ A1938; \skt{प्यन्नम्}\ N}}}{\lemma{\skt{उपयुक्तम्}}\Afootnote{H; \skt{पयुक्त}\ A1938; \skt{उपयुक्ताम्}\ N}}}{\lemma{\skt{उदरशिरोगौरवकासश्वासच्छर्दिप्रसेकगात्रसदनानि}}\Afootnote{N H; \skt{मुदरशिरोगौरवकासश्वासप्रसेकच्छर्दिगात्रसदनानि}\ A1938}}}{\lemma{\skt{व्यपदिशन्ति ॥ यस्}}\Afootnote{H; \skt{व्यपदिशन्ति । य}\ N}}\skt{ कृत्वा }\edtext{\skt{महा}}{\Afootnote{H; \skt{महता}\ A1938 N}}\skt{ कालेन पचति स }\edtext{\skt{मन्दः ॥ भ}}{\Afootnote{H; \skt{मन्दः ॥२४॥}\ A1938; \skt{मर्र्दः ॥ भवति चात्र ॥}\ N}}\skt{ ॥}\edtext{}{\lemma{\skt{ति… व्यपदिशन्ति}}\Afootnote{\emph{om.}\ A1938}}\pend

\pstart\versepar \ledsidenote{\passref{1.35.25}}\skt{विषमो }\edtext{\skt{वातजान्रोगाञ्स्तीक्ष्णः}}{\Afootnote{H; \skt{वातजान्रोगांस्तीक्ष्णः}\ A1938; \skt{वातजां रोगान्स्तीक्ष्ण}\ N}}\skt{ पित्तनिमित्तजान् । }\pend
\pstart\versepar \skt{}\edtext{\edtext{\skt{करोत्यग्निस्तथा}}{\lemma{\skt{अग्निस्}}\Afootnote{N H; \skt{त्यग्निस्}\ A1938}}}{\lemma{\skt{तथा}}\Afootnote{A1938 H; \skt{तथो}\ N}}\skt{ }\edtext{\skt{मन्दो}}{\Afootnote{A1938 H; \skt{मन्दो}\ N}}\skt{ }\edtext{\skt{विकारान्कफसम्भवान्}}{\Afootnote{A1938 H; \skt{विकाराङ्कफसम्भवाम् ॥}\ N}}\skt{ ॥}\edtext{}{\lemma{\skt{पित्तनिमित्तजान् । करोत्य्}}\Afootnote{H; \skt{पित्तनिमित्तजान् । करो}\ A1938; \skt{पित्तनिमित्तजां । करोत्य्}\ N}}\pend

\pstart \ledsidenote{\passref{1.35.26}}\skt{तत्र समे }\edtext{\skt{रक्षणं}}{\Afootnote{N H; \skt{परिरक्षणं}\ A1938}}\skt{ कुर्व्वीत ॥ }\edtext{\skt{विषमं}}{\Afootnote{N H; \skt{विषमे}\ A1938}}\skt{ }\edtext{\skt{लवणाम्लैः}}{\Afootnote{H; \skt{स्निग्धाम्ललव}\ A1938; \skt{ललवणाम्लोष्णैः}\ N}}\skt{ }\edtext{\skt{स्नेहयुक्तैः}}{\Afootnote{N H; \skt{णैः}\ A1938}}\skt{ }\edtext{\skt{क्रियाविशेषैरुपचरेत् ॥ }\edtext{\edtext{\skt{तीक्ष्णम्मधुरस्निग्धयुक्तैर्व्विरेचनैश्च}}{\lemma{\skt{मधुरस्निग्धयुक्तैर्}}\Afootnote{H; \skt{मधुरस्निग्धशीतैर्}\ A1938; \skt{मधुर युक्तै}\ N}}}{\lemma{\skt{व्विरेचनैश्च}}\Afootnote{H; \skt{विरेकैश्च;}\ A1938; \skt{र्विरेचनैविरेचनैश्च ।}\ N}}}{\lemma{\skt{उपचरेत् ॥ तीक्ष्णम्}}\Afootnote{H; \skt{प्रतिकुर्वीत; तीक्ष्णे}\ A1938; \skt{उपरेत् । तीक्ष्ण}\ N}}\skt{ । }\edtext{\edtext{\skt{एवमेवात्यग्निम्विशेषेण}}{\lemma{\skt{एवात्य्}}\Afootnote{H; \skt{एवा}\ A1938 N}}}{\lemma{\skt{अग्निम्}}\Afootnote{H; \skt{त्यग्नौ,}\ A1938; \skt{त्यग्निं}\ N}}\skt{ }\edtext{\edtext{\skt{माहियैर्द्दधिक्षीरसर्प्पिभिरिति}}{\lemma{\skt{माहियैर्द्दधिक्षीरसर्प्पिभिर्}}\Afootnote{H; \skt{माहिषैश्च क्षीरदधिसर्पिर्भिः;}\ A1938; \skt{माहिषैर्दधिक्षीरसर्पिर्भिर्}\ N}}}{\lemma{\skt{इति}}\Afootnote{N H; \emph{om.}\ A1938}}\skt{ ॥ }\edtext{\skt{मन्दे}}{\Afootnote{A1938 H; \skt{मन्दङ्}\ N}}\skt{ }\edtext{\skt{कटुतिक्तकषायैर्व्वमनैश्च ॥ भ}}{\lemma{\skt{व्वमनैश्च ॥ भ}}\Afootnote{H; \skt{वमनैश्च ॥२६॥}\ A1938; \skt{वमनैश्च ॥ भवतश्चात्र ॥}\ N}}\skt{ ॥}\pend

\pstart\versepar \ledsidenote{\passref{1.35.27}}\edtext{\skt{औदार्यो}}{\Afootnote{H; \skt{जाठरो}\ A1938; \skt{ओदार्यो}\ N}}\skt{ भगवानग्निः }\edtext{\edtext{\skt{पावकोऽन्नस्य}}{\lemma{\skt{पावको}}\Afootnote{H; \skt{ईश्वरो}\ A1938; \emph{om.}\ N}}}{\lemma{\skt{ऽन्नस्य}}\Afootnote{H; \skt{ऽन्नस्य पाच}\ A1938; \skt{पाचकोन्नस्य}\ N}}\skt{ }\edtext{\skt{चेश्वरः}}{\Afootnote{N H; \skt{कः ।}\ A1938}}\skt{ । }\pend
\pstart\versepar \skt{सौक्ष्म्याद्रसानाददानो }\edtext{\skt{विवेक्तुन्नेह}}{\lemma{\skt{नेह}}\Afootnote{N H; \skt{नैव}\ A1938}}\skt{ शक्यते ॥}\pend

\pstart\versepar \ledsidenote{\passref{1.35.28}}\edtext{\skt{प्राणापानसमानैस्तु}}{\Afootnote{A1938 H; \skt{प्राणोपाणसमानैस्तु}\ N}}\skt{ सर्व्वतः }\edtext{\edtext{\skt{पवनैस्त्रिभिः}}{\lemma{\skt{पवनैस्}}\Afootnote{N H; \skt{पवनै}\ A1938}}}{\lemma{\skt{त्रिभिः}}\Afootnote{N H; \skt{स्त्रिभिः ।}\ A1938}}\skt{ । }\pend
\pstart\versepar \skt{}\edtext{\skt{ध्यामते}}{\Afootnote{H; \skt{ध्मायते}\ A1938; \skt{ध्माप्यते}\ N}}\skt{ पाल्यते }\edtext{\skt{चापि स्वां स्वां}}{\Afootnote{A1938 H; \skt{चापि स्चांस्वाङ्}\ N}}\skt{ गतिमवस्थितैः ॥}\pend

\pstart \ledsidenote{\passref{1.35.29}}\skt{वयस्तु }\edtext{\skt{त्रिविधम्वाल्यं}}{\lemma{\skt{वाल्यं}}\Afootnote{H; \skt{बाल्यं,}\ A1938; \skt{ललं}\ N}}\skt{ मध्यम्वृद्धमिति । }\edtext{\skt{तत्रोनषोडशवर्षाद्वालाभवन्ति । ते}}{\Afootnote{H; \skt{तत्रोनषोडशवर्षीया बालाः । ते}\ A1938; \skt{तत्रोनसोडषवशाद्बाला भवन्ति । ते}\ N}}\skt{ }\edtext{\skt{त्रिविधाः}}{\Afootnote{A1938 H; \skt{त्रिविधा}\ N}}\skt{ क्षीरपाः }\edtext{\skt{क्षीरान्नादा}}{\Afootnote{H; \skt{क्षीरान्नादा, अन्नादा}\ A1938; \skt{क्षीरान्}\ N}}\skt{ }\edtext{\skt{अन्नादा}}{\Afootnote{H; \skt{इति । तेषु संवत्सरपराः क्षीरपाः, द्विसंवत्सरपराः क्षीरान्नादाः, परतोऽन्नादा}\ A1938; \skt{नदा अन्नादा}\ N}}\skt{ इति । षोडशसप्तत्योरन्तरे }\edtext{\edtext{\skt{मध्यस्वयस्तस्यद्विकल्पो}}{\lemma{\skt{मध्यस्वयस्}}\Afootnote{H; \skt{मध्यं वयः ।}\ A1938; \skt{मध्यम्व}\ N}}}{\lemma{\skt{तस्यद्विकल्पो}}\Afootnote{H; \skt{तस्य विकल्पो-}\ A1938; \skt{यस्कन्विकल्पा}\ N}}\skt{ }\edtext{\edtext{\skt{वृद्धिर्यौवनं}}{\lemma{\skt{वृद्धिर्}}\Afootnote{A1938 H; \skt{वृद्धि}\ N}}}{\lemma{\skt{यौवनं}}\Afootnote{A1938 H; \skt{र्यौवनं}\ N}}\skt{ }\edtext{\skt{सम्पूर्ण्णतेति}}{\Afootnote{N H; \skt{सम्पूर्णता, परिहाणिरिति । त}\ A1938}}\skt{ । }\edtext{\edtext{\edtext{\skt{तत्राविंशतेर्वृद्धिरात्रिंशतो}}{\lemma{\skt{तत्राविंशतेर्}}\Afootnote{H; \skt{त्र, आविंशते}\ A1938; \skt{विंशते}\ N}}}{\lemma{\skt{वृद्धिर्}}\Afootnote{H; \skt{र्वृद्धिः,}\ A1938; \skt{रा}\ N}}}{\lemma{\skt{आत्रिंशतो}}\Afootnote{A1938 H; \skt{त्रिंशतो}\ N}}\skt{ यौवनमाचत्वारिंशतः }\edtext{\skt{सर्व्वधात्विन्द्रियसम्पूर्ण्णेतिति}}{\Afootnote{H; \skt{सर्वधात्विन्द्रियबलवीर्यसम्पूर्णता,}\ A1938; \skt{सर्वधात्विन्द्रियसंपूर्ण्णतेति ।}\ N}}\skt{ ॥ अत }\edtext{\skt{ऊर्ध्वं}}{\Afootnote{H; \skt{ऊर्ध्व}\ A1938; \skt{ऊर्द्वं}\ N}}\skt{ }\edtext{\skt{हा }\edtext{\skt{निर्भवति}}{\lemma{\skt{भवति}}\Afootnote{N H; \skt{णि}\ A1938}}}{\lemma{\skt{हा निर्}}\Afootnote{N H; \skt{मीषत्परिहा}\ A1938}}\skt{ }\edtext{\skt{यावदा}}{\Afootnote{N H; \skt{र्यावत्स}\ A1938}}\skt{ }\edtext{\skt{सप्ततेः}}{\Afootnote{N H; \skt{प्ततिरिति ।}\ A1938}}\skt{ । }\edtext{\skt{सप्तत्यास्रूर्धं}}{\Afootnote{H; \skt{सप्ततेरूर्ध्वं}\ A1938; \skt{सप्तत्यास्तूर्द्वं}\ N}}\skt{ }\edtext{\edtext{\skt{क्षीयमाणधात्विन्द्रियवलवीर्योत्साहोऽन्यहनि}}{\lemma{\skt{क्षीयमाणधात्विन्द्रियवलवीर्योत्साहो}}\Afootnote{H; \skt{क्षीयमाणधात्विन्द्रियबलवीर्योत्साह}\ A1938; \skt{क्षीयमाणधात्वन्द्रियबलवीर्योत्साहो}\ N}}}{\lemma{\skt{ऽन्यहनि}}\Afootnote{H; \skt{महन्यहनि}\ A1938; \skt{हन्यहनि}\ N}}\skt{ }\edtext{\edtext{\edtext{\skt{वलीपलितखालित्यकासश्वासप्रभृतिभिरुपद्रवैरभिभूयमानं}}{\lemma{\skt{वलीपलितखालित्य्}}\Afootnote{A1938 H; \skt{वलीपलितखालि}\ N}}}{\lemma{\skt{अकासश्वासप्रभृतिभिर्}}\Afootnote{H; \skt{अजुष्टं कासश्वासप्रभृतिभिर्}\ A1938; \skt{त्यकासश्वासप्रभृतिभिर्}\ N}}}{\lemma{\skt{अभिभूयमानं}}\Afootnote{A1938 H; \skt{अभिभूयमानः}\ N}}\skt{ }\edtext{\skt{सर्व्वक्रियास्वसमर्था}}{\Afootnote{H; \skt{सर्वक्रियास्वसमर्थं}\ A1938; \skt{सर्वक्रियास्वसमर्थो}\ N}}\skt{ }\edtext{\edtext{\skt{जीर्ण्णागारमिवावसीदति}}{\lemma{\skt{जीर्ण्णागारम्}}\Afootnote{H; \skt{जीर्णागार}\ A1938; \skt{जीरागार}\ N}}\skt{ । }\edtext{\edtext{\skt{तम्वृद्धमाचक्षते}}{\lemma{\skt{वृद्धम्}}\Afootnote{N H; \skt{वृद्ध}\ A1938}}}{\lemma{\skt{आचक्षते}}\Afootnote{N H; \skt{माचक्षते ॥२९॥}\ A1938}}}{\lemma{\skt{इवावसीदति । तम्}}\Afootnote{H; \skt{मिवाभिवृष्टमवसीदन्तं}\ A1938; \skt{मिचावसीदति । तं}\ N}}\skt{ ।}\pend

\pstart \ledsidenote{\passref{1.35.30}}\edtext{\skt{उत्तरोत्तरासु च}}{\Afootnote{N H; \skt{तत्रोत्तरोत्तरासु}\ A1938}}\skt{ }\edtext{\edtext{\skt{वयोऽवस्त्जासु}}{\lemma{\skt{वयो}}\Afootnote{H; \skt{वयो}\ A1938 N}}}{\lemma{\skt{ऽवस्त्जासु}}\Afootnote{H; \skt{वस्थासू}\ A1938 N}}\skt{ }\edtext{\skt{उत्तरोत्तरे}}{\Afootnote{H; \skt{त्तरोत्तरा}\ A1938; \skt{त्तरोत्तर}\ N}}\skt{ }\edtext{\skt{भेषजमात्रप्रयोगः}}{\Afootnote{N H; \skt{भेषजमात्राविशेषा भवन्ति,}\ A1938}}\skt{ । }\edtext{\skt{ऋते}}{\Afootnote{N H; \skt{ऋते च}\ A1938}}\skt{ }\edtext{\edtext{\skt{परिहानिम्कुर्वन्ति}}{\lemma{\skt{परिहानिम्}}\Afootnote{H; \skt{परिहाणेः; तत्राद्यापेक्षया प्रति}\ A1938; \skt{परिहाणिर्भवति ॥}\ N}}}{\lemma{\skt{कुर्वन्ति}}\Afootnote{H; \skt{कुर्वीत ॥३०॥}\ A1938; \emph{om.}\ N}}\skt{ ॥}\pend

\pstart\versepar \ledsidenote{\passref{1.35.31}}\edtext{\skt{भवन्ति}}{\Afootnote{A1938; \skt{भवतः}\ N; \emph{om.}\ H}}\skt{ }\edtext{\skt{चात्र}}{\Afootnote{A1938 N; \emph{om.}\ H}}\skt{ । }\pend
\pstart\versepar \skt{}\edtext{\skt{वाले}}{\Afootnote{H; \skt{बाले}\ A1938 N}}\skt{ विवर्द्धते }\edtext{\skt{श्लेमा}}{\Afootnote{H; \skt{श्लेष्मा}\ A1938 N}}\skt{ मध्यमे }\edtext{\edtext{\skt{पित्तमेव}}{\lemma{\skt{पित्तम्}}\Afootnote{A1938 H; \skt{पित्त}\ N}}\skt{ तु}}{\lemma{\skt{एव तु}}\Afootnote{A1938 H; \skt{मेव तु ।}\ N}}\skt{ । }\pend
\pstart\versepar \skt{}\edtext{\skt{भूयिष्ठम्वर्धते}}{\lemma{\skt{वर्धते}}\Afootnote{A1938 H; \emph{om.}\ N}}\skt{ वायुर्व्वृद्धे }\edtext{\skt{तद्वीक्ष्य}}{\Afootnote{H; \skt{तद्वीक्ष्य}\ A1938 N}}\skt{ योजयेत् ॥}\pend

\pstart\versepar \ledsidenote{\passref{1.35.32}}\edtext{\skt{अग्निक्षारविरेकैस्तु}}{\Afootnote{A1938 H; \skt{अग्निक्षारविरेकैस्तु}\ N}}\skt{ }\edtext{\skt{वालवृद्धौ}}{\Afootnote{H; \skt{बालवृद्धौ}\ A1938 N}}\skt{ विवर्ज्जयेत् । }\pend
\pstart\versepar \skt{तत्साध्येषु विकारेषु }\edtext{\skt{मृद्वीं}}{\Afootnote{A1938 H; \skt{मृद्वी}\ N}}\skt{ }\edtext{\skt{कुर्यात्क्रियां}}{\Afootnote{A1938 H; \skt{कुर्यात्क्रिया}\ N}}\skt{ शनैः ॥}\pend

\pstart \ledsidenote{\passref{1.35.33}}\skt{देहः स्थूलः कृशो मध्य इति }\edtext{\skt{प्रागभितं}}{\Afootnote{H; \skt{प्रागुपदिष्टः ॥३३॥}\ A1938; \skt{प्रागविहितं ॥}\ N}}\skt{ ॥ }\edtext{\skt{च}}{\Afootnote{H; \emph{om.}\ A1938 N}}\skt{ ॥}\pend

\pstart\versepar \ledsidenote{\passref{1.35.34}}\edtext{\skt{कर्षयेद्वृहयेद्च्चापि}}{\Afootnote{H; \skt{कर्शयेद्बृंहयेच्चापि}\ A1938; \skt{भवन्ति चात्र ॥ कर्शयेद्बृंहयेच्चापि}\ N}}\skt{ सदा स्थूलकृशौ नरौ । }\pend
\pstart\versepar \skt{}\edtext{\skt{मध्यस्य}}{\Afootnote{N H; \emph{om.}\ A1938}}\skt{ }\edtext{\skt{रक्षणञ्चापि}}{\lemma{\skt{चापि}}\Afootnote{N H; \skt{चैव म}\ A1938}}\skt{ }\edtext{\skt{कुर्व्वीत}}{\Afootnote{N H; \skt{ध्यस्य कुर्वीत}\ A1938}}\skt{ सततम्भिषक् ॥}\pend

\pstart \ledsidenote{\passref{1.35.35}}\edtext{\edtext{\skt{वलन्त्वनभिहतगुणं}}{\lemma{\skt{वलन्}}\Afootnote{H; \skt{बलम्}\ A1938 N}}}{\lemma{\skt{त्वनभिहतगुणं}}\Afootnote{H; \skt{अभिहितगुणं;}\ A1938; \skt{त्वभिहितगुणं}\ N}}\skt{ }\edtext{\skt{दौर्व्वल्यन्तु}}{\Afootnote{H; \skt{दौर्बल्यं तु}\ A1938 N}}\skt{ }\edtext{\edtext{\edtext{\skt{स्वभावदोषजराभिरापेक्षिकम्भवति}}{\lemma{\skt{स्वभावदोषजराभिर्}}\Afootnote{H; \skt{स्वभावदोषजरादिभिर्}\ A1938; \skt{स्वभावओषजराभिर्}\ N}}}{\lemma{\skt{आपेक्षिकम्}}\Afootnote{H; \skt{अवेक्षि}\ A1938; \skt{आवेक्षिम्}\ N}}}{\lemma{\skt{भवति}}\Afootnote{N H; \skt{तव्यम् ।}\ A1938}}\skt{ ॥ }\edtext{\skt{यस्माद्वलवतः}}{\Afootnote{H; \skt{यस्माद्बलवतः}\ A1938; \skt{यस्माद्बलवः}\ N}}\skt{ }\edtext{\edtext{\skt{सर्व्वक्रियाप्रतिपत्तिस्तस्माद्वलप्रधानमधिकरणानां}}{\lemma{\skt{सर्व्वक्रियाप्रतिपत्तिस्}}\Afootnote{H; \skt{सर्वक्रियाप्रवृत्तिस्}\ A1938; \skt{सर्वक्रियाप्रतिपत्ति}\ N}}}{\lemma{\skt{तस्माद्वलप्रधानम्}}\Afootnote{H; \skt{तस्माद्बलमेव प्रधानम्}\ A1938; \skt{तस्माद्बलम्प्रधानम्}\ N}}\skt{ ।}\pend

\pstart\versepar \ledsidenote{\passref{1.35.36}}\skt{केचित्कृशाः प्राणवन्तः }\edtext{\edtext{\skt{स्थूलाश्चाल्पवला}}{\lemma{\skt{स्थूलाश्}}\Afootnote{A1938 H; \skt{स्थूला}\ N}}}{\lemma{\skt{चाल्पवला}}\Afootnote{H; \skt{चाल्पबला}\ A1938 N}}\skt{ नराः । }\pend
\pstart\versepar \skt{}\edtext{\edtext{\edtext{\skt{तस्मात्स्थिरत्वव्यायामेर्व्वलम्वैद्यः}}{\lemma{\skt{तस्मात्स्थिरत्वव्यायामेर्}}\Afootnote{H; \skt{तस्मात्स्थिरत्वं व्यायामैर्}\ A1938; \skt{तस्मात्स्थिरत्वव्यायामैर्}\ N}}}{\lemma{\skt{व्वलम्}}\Afootnote{H; \skt{बलं}\ A1938 N}}}{\lemma{\skt{वैद्यः}}\Afootnote{A1938 H; \skt{वैद्य}\ N}}\skt{ }\edtext{\skt{प्रतर्क्कयेत्}}{\Afootnote{A1938 H; \skt{प्रतर्पयेत् ॥}\ N}}\skt{ ॥}\pend

\pstart \ledsidenote{\passref{1.35.37}}\skt{सत्वन्तु }\edtext{\edtext{\edtext{\edtext{\skt{व्यसनभयाभ्युदयक्रियादिस्थानेष्वविकम्वकरम्भवति}}{\lemma{\skt{व्यसनभयाभ्युदयक्रियादिस्थानेष्व्}}\Afootnote{N H; \skt{व्यसनाभ्युदयक्रियादिस्थानेष्व्}\ A1938}}}{\lemma{\skt{अविकम्}}\Afootnote{N H; \skt{अविक्ल}\ A1938}}}{\lemma{\skt{वकरम्}}\Afootnote{A1938 H; \skt{पकरम्}\ N}}}{\lemma{\skt{भवति}}\Afootnote{N H; \emph{om.}\ A1938}}\skt{ ॥ }\edtext{\skt{च}}{\Afootnote{\emph{om.}\ A1938 N H}}\skt{ ॥}\pend

\pstart\versepar \ledsidenote{\passref{1.35.38}}\edtext{\skt{भवति}}{\Afootnote{N; \emph{om.}\ A1938 H}}\skt{ }\edtext{\skt{चात्र}}{\Afootnote{N; \emph{om.}\ A1938 H}}\skt{ ॥ }\pend
\pstart\versepar \skt{सत्ववान्सहते सर्व्वं }\edtext{\skt{संस्तभ्यमानः}}{\Afootnote{H; \skt{संस्तभ्यात्मान}\ A1938; \skt{स्तम्भ्यमानः}\ N}}\skt{ }\edtext{\edtext{\skt{परैर्न्नरः}}{\lemma{\skt{परैर्न्}}\Afootnote{H; \skt{मात्म}\ A1938; \skt{परैर्}\ N}}}{\lemma{\skt{नरः}}\Afootnote{N H; \skt{ना ।}\ A1938}}\skt{ । }\pend
\pstart\versepar \skt{}\edtext{\skt{राजसः}}{\Afootnote{A1938; \skt{राजास}\ N; \skt{राजसं}\ H}}\skt{ }\edtext{\edtext{\skt{स्तभ्यमानोऽपि}}{\lemma{\skt{स्तभ्यमानो}}\Afootnote{A1938 H; \skt{स्तम्भ्यमानो}\ N}}}{\lemma{\skt{ऽपि}}\Afootnote{N H; \skt{ऽन्यैः}\ A1938}}\skt{ }\edtext{\skt{सहते}}{\Afootnote{A1938 H; \skt{हसते सहते}\ N}}\skt{ नैव तामसः ॥}\pend

\pstart \ledsidenote{\passref{1.35.41}}\edtext{\edtext{\edtext{\skt{प्रकृतिम्भेषजञ्चान्यत्रोपदेक्ष्यामः}}{\lemma{\skt{प्रकृतिम्}}\Afootnote{A1938 H; \skt{प्रकृति}\ N}}}{\lemma{\skt{भेषजञ्}}\Afootnote{A1938 H; \skt{भेषंजञ्}\ N}}}{\lemma{\skt{चान्यत्रोपदेक्ष्यामः}}\Afootnote{N H; \skt{चोपरिष्टाद्वक्ष्यामः ॥४१॥}\ A1938}}\skt{ ।}\pend

\pstart \ledsidenote{\passref{1.35.39}}\skt{सात्म्यानि तु }\edtext{\skt{देशजानि}}{\Afootnote{H; \skt{देशकालजात्यृतु}\ A1938; \skt{देसजानि}\ N}}\skt{ }\edtext{\skt{रोगर्तुव्यायामोदकरसदिवास्वप्नप्रभृतीनि}}{\Afootnote{\skt{रोगव्यायामोदकदिवास्वप्नरसप्रभृतीनि}\ A1938; \skt{रोगार्त्तुव्यायामोदकरसदिठास्वप्नप्रभृतीनि}\ N; \skt{रोगर्तुव्यायामोदकरसदिवास्प्नप्रभृतीनि}\ H}}\skt{ }\edtext{\skt{प्रकृतिविरुद्धान्यपि}}{\lemma{\skt{प्रकृतिविरुद्धान्य्}}\Afootnote{A1938 H; \skt{विरुद्धान्य्}\ N}}\skt{ }\edtext{\skt{यान्यनावाधकराणि}}{\lemma{\skt{अनावाधकराणि}}\Afootnote{H; \skt{अबाधकराणि}\ A1938; \skt{अनाबाधकराणि}\ N}}\skt{ भवन्ति ॥}\pend

\pstart\versepar \ledsidenote{\passref{1.35.40}}\edtext{\skt{भवति}}{\Afootnote{N; \emph{om.}\ A1938 H}}\skt{ चात्र ॥ }\pend
\pstart\versepar \skt{यो रसः }\edtext{\skt{कल्पितो}}{\Afootnote{N H; \skt{कल्पते}\ A1938}}\skt{ यस्य }\edtext{\skt{सुखतां}}{\Afootnote{N H; \skt{सुखायै}\ A1938}}\skt{ }\edtext{\skt{याति}}{\Afootnote{H; \skt{व नि}\ A1938; \skt{यांति}\ N}}\skt{ }\edtext{\skt{सेवितः}}{\Afootnote{N H; \skt{षेवितः ।}\ A1938}}\skt{ । }\pend
\pstart\versepar \skt{}\edtext{\edtext{\skt{व्यायामजातमन्यद्वा}}{\lemma{\skt{व्यायामजातम्}}\Afootnote{A1938 H; \skt{व्यायामजात}\ N}}}{\lemma{\skt{अन्यद्वा}}\Afootnote{A1938 H; \skt{सन्यद्वा}\ N}}\skt{ तत्सात्म्यमिति निर्द्दिशेत् ॥}\edtext{}{\lemma{\skt{चात्र ॥ यो}}\Afootnote{N; \skt{यो}\ A1938 H}}\pend

\pstart \ledsidenote{\passref{1.35.42}}\edtext{\edtext{\skt{देशस्त्वानूपो}}{\lemma{\skt{देशस्}}\Afootnote{A1938 H; \skt{देश}\ N}}}{\lemma{\skt{त्वानूपो}}\Afootnote{A1938 H; \skt{स्त्वार्नून्क्रपो}\ N}}\skt{ जाङ्गलः साधारण }\edtext{\skt{इति}}{\Afootnote{N H; \skt{इ}\ A1938}}\skt{ }\edtext{\skt{त्रिविधा}}{\Afootnote{H; \emph{om.}\ A1938; \skt{त्रिविधो}\ N}}\skt{ }\edtext{\skt{भवति}}{\Afootnote{N H; \skt{ति ।}\ A1938}}\skt{ । तत्र }\edtext{\skt{निम्नोन्नतवहूदकनदीवर्षगहनमृदूशीतलो}}{\Afootnote{H; \skt{बहूदकनिम्नोन्नतनदीवर्षगहनो मृदुशीतानिलो}\ A1938; \skt{निम्नोन्नतबहूदकनदीवर्षगहण मृदुशीतलो}\ N}}\skt{ }\edtext{\skt{वहूमहापर्वतवृक्षो}}{\Afootnote{H; \skt{बहुमहापर्वतवृक्षो}\ A1938 N}}\skt{ मृदुसुकुमारोपचितशरीरमनुष्यप्रायो }\edtext{\edtext{\skt{वातकफरोगभूयिष्ठश्चानूपः}}{\lemma{\skt{वातकफरोगभूयिष्ठश्}}\Afootnote{H; \skt{कफवातरोगभूयिष्ठश्}\ A1938; \skt{वातकफलरोगभूयिष्ठो}\ N\textsuperscript{ac}; \skt{वातकफरोगभूयिष्ठो}\ N\textsuperscript{pc}}}}{\lemma{\skt{चानूपः}}\Afootnote{H; \skt{चानूपः;}\ A1938; \skt{स्वात्क्रपः ॥}\ N}}\skt{ ॥ }\edtext{\edtext{\skt{वक्ताकासमवितरलाल्पकण्टकवृक्षवर्षप्रस्रवणोदपानप्रायोष्णदारूणवातप्रविरलाल्पशैलस्थिरकृशशरीरमनुष्यप्रायो}}{\lemma{\skt{वक्ताकासम्}}\Afootnote{H; \skt{आकाशसमः}\ A1938; \skt{बङ्काकाशसम्}\ N}}}{\lemma{\skt{अवितरलाल्पकण्टकवृक्षवर्षप्रस्रवणोदपानप्रायोष्णदारूणवातप्रविरलाल्पशैलस्थिरकृशशरीरमनुष्यप्रायो}}\Afootnote{H; \skt{प्रविरलाल्पकण्टकिवृक्षप्रायोऽल्पवर्षप्रस्रवणोदपानोदकप्राय उष्णदारुणवातः प्रविरलाल्पशैलः स्थिरकृशशरीरमनुष्यप्रायो}\ A1938; \skt{अविरलाल्पकण्ठकवृक्षवर्षप्रस्रवण्येदपानप्रायोष्णदारुणवातप्रविरलाल्पशैलस्थिरकृशसरी॥॥॥॥॥।रमनुष्यप्रायो}\ N}}\skt{ }\edtext{\skt{वातपित्तरो गभूयिष्ठश्च}}{\Afootnote{\skt{वातपित्तरोगभूयिष्ठश्च}\ A1938 H; \skt{वातपित्तरोगभूयिष्ठष्च}\ N}}\skt{ जाङ्गलः । }\edtext{\edtext{\skt{उभयलक्षणस्साधारणस्य}}{\lemma{\skt{उभयलक्षणस्}}\Afootnote{N H; \skt{उभयदेशलक्षणः}\ A1938}}\skt{ प्रधान}}{\lemma{\skt{साधारणस्य प्रधान}}\Afootnote{H; \skt{साधारण}\ A1938; \skt{साधारण सप्रधान}\ N}}\skt{ }\edtext{\skt{इति}}{\Afootnote{A1938 N; \skt{इति ॥ भ ॥}\ H}}\skt{ ॥}\pend

\pstart\versepar \ledsidenote{\passref{1.35.43}}\edtext{\skt{भवति}}{\Afootnote{\skt{भवन्ति}\ A1938 N; \emph{om.}\ H}}\skt{ }\edtext{\skt{चात्र}}{\Afootnote{A1938 N; \emph{om.}\ H}}\skt{ ॥ }\pend
\pstart\versepar \skt{समाः साधारणे }\edtext{\skt{यस्माद्वर्षशीतोष्णमारूताः}}{\Afootnote{H; \skt{यस्माच्छीतवर्षोष्ममारुताः ।}\ A1938; \skt{यस्माद्वर्षाशीतोष्णमारुताः ।}\ N}}\skt{ । }\pend
\pstart\versepar \skt{दोषाणां समता }\edtext{\skt{वापि}}{\Afootnote{H; \skt{जन्तो}\ A1938; \skt{चापि}\ N}}\skt{ }\edtext{\skt{तस्मात्साधारणो}}{\Afootnote{N H; \skt{स्तस्मात्साधारणो}\ A1938}}\skt{ }\edtext{\skt{वरः}}{\Afootnote{N H; \skt{मतः ॥४३॥}\ A1938}}\skt{ ॥}\pend

\pstart\versepar \ledsidenote{\passref{1.35.46}}\skt{देशप्रकृति }\edtext{\skt{सात्म्यर्त्व}}{\Afootnote{N H; \skt{सात्म्यर्तु}\ A1938}}\skt{ }\edtext{\skt{विपरीताचिरोत्थितः}}{\Afootnote{N H; \skt{विपरीतोऽचिरोत्थितः ।}\ A1938}}\skt{ । }\pend
\pstart\versepar \skt{सम्पत्तौ }\edtext{\edtext{\edtext{\skt{भेषजादीनाम्वलसत्वायुषान्तथा}}{\lemma{\skt{भेषजादीनाम्}}\Afootnote{N H; \skt{भिषगादीनां}\ A1938}}}{\lemma{\skt{वलसत्वायुषान्}}\Afootnote{H; \skt{बलसत्त्वायुषां}\ A1938; \skt{बलसत्वायुषा}\ N}}}{\lemma{\skt{तथा}}\Afootnote{A1938 H; \skt{त्तथां ॥}\ N}}\skt{ ॥}\pend

\pstart\versepar \ledsidenote{\passref{1.35.47}}\skt{केवलः समदेहाग्नेः सुखसाध्यतमो }\edtext{\skt{मतः}}{\Afootnote{N H; \skt{गदः ।}\ A1938}}\skt{ । }\pend
\pstart\versepar \skt{}\edtext{\edtext{\skt{अतोऽन्यथाप्यसाध्यः}}{\lemma{\skt{ऽन्यथाप्य्}}\Afootnote{\skt{ऽन्यथा}\ A1938; \skt{न्यथाम्य्}\ N; \skt{न्यथाप्य्}\ H}}}{\lemma{\skt{असाध्यः}}\Afootnote{N H; \skt{त्वसाध्यः}\ A1938}}\skt{ स्यात्कृच्छ्रो }\edtext{\skt{व्यामिश्रलक्षणः}}{\Afootnote{A1938 H; \skt{व्यामिसलक्षणः ।}\ N}}\skt{ ॥}\pend

\pstart\versepar \ledsidenote{\passref{1.35.48}}\skt{क्रियायास्तु }\edtext{\skt{गुणालाभे}}{\Afootnote{A1938 H; \skt{गुण्यलाभे}\ N}}\skt{ क्रियामन्यां प्रयोजयेत् । }\pend
\pstart\versepar \skt{पूर्व्वस्यां }\edtext{\skt{शान्तवेगायान्न}}{\Afootnote{A1938 H; \skt{शन्तवेगायान्न}\ N}}\skt{ }\edtext{\skt{क्रियासंअरो}}{\Afootnote{H; \skt{क्रियासङ्करो}\ A1938 N}}\skt{ हितः ॥}\edtext{}{\lemma{\skt{प्रयोजयेत् । पूर्व्वस्यां}}\Afootnote{A1938 H; \skt{प्रयोजयेत् ॥ पूर्वस्या}\ N}}\pend

\pstart\versepar \ledsidenote{\passref{1.35.44}}\edtext{\skt{न तथा}}{\Afootnote{A1938 N; \skt{न तह्हा}\ H}}\skt{ }\edtext{\skt{बलवन्तः स्युः}}{\Afootnote{N; \skt{बलवन्तः स्यु}\ A1938; \skt{वलवन्तस्य वन्तः स्युः}\ H}}\skt{ }\edtext{\skt{जलजा वा}}{\Afootnote{N H; \skt{र्जलजा वा}\ A1938}}\skt{ }\edtext{\skt{स्थलाहृताः}}{\Afootnote{A1938 N; \skt{स्थलाहताः ।}\ H}}\skt{ । }\pend
\pstart\versepar \skt{}\edtext{\skt{स्वदेशे}}{\Afootnote{A1938; \skt{त्येदेशे}\ N; \skt{स्वेदेणे}\ H}}\skt{ }\edtext{\skt{निचिता}}{\Afootnote{A1938 H; \skt{निविता}\ N}}\skt{ }\edtext{\edtext{\skt{दोषास्त्वन्यस्मिन्कोपमागताः}}{\lemma{\skt{दोषास्}}\Afootnote{N H; \skt{दोषा}\ A1938}}}{\lemma{\skt{त्वन्यस्मिन्कोपम्}}\Afootnote{\skt{अन्यस्मिन्कोपम्}\ A1938; \skt{त्वन्यस्मिं कोपम्}\ N; \skt{त्वन्यस्मिन्कायम्}\ H}}\skt{ ॥}\pend

\pstart\versepar \ledsidenote{\passref{1.35.45}}\skt{उचिते वर्त्तमानस्य नास्ति }\edtext{\edtext{\skt{देशकृतम्भयं}}{\lemma{\skt{देशकृतम्}}\Afootnote{A1938 H; \skt{दोषकृतम्}\ N}}}{\lemma{\skt{भयं}}\Afootnote{N H; \skt{भयम् । ૉ}\ A1938}}\skt{ । }\pend
\pstart\versepar \skt{आहारस्वप्नचेष्टादौ }\edtext{\skt{तद्देशः}}{\Afootnote{N; \skt{तद्देश}\ A1938; \skt{तद्देशः}\ H}}\skt{ }\edtext{\skt{सगुण}}{\Afootnote{H; \skt{स्य गुणे}\ A1938; \skt{यगुणे}\ N}}\skt{ }\edtext{\skt{सतीति}}{\Afootnote{N H; \skt{सति ॥४५॥}\ A1938}}\skt{ ॥}\pend

\pstart \ledsidenote{\passref{1.35.trailer}}\edtext{\edtext{\skt{इत्यध्याय}}{\lemma{\skt{इत्य्}}\Afootnote{\skt{इति सुश्रुतसंहितायां सूत्रस्थाने आतुरोपक्रमणीयो नाम पञ्चत्रिंशो}\ A1938; \skt{॥ ३७ ॥ ओ ॥}\ N; \skt{ल हृ ३० ५}\ H}}}{\lemma{\skt{अध्याय}}\Afootnote{\skt{ऽध्यायः ॥३५॥}\ A1938; \emph{om.}\ N H}}\skt{ ३५}\pend

\endnumbering
\beginnumbering
\eledsection*{Adhyāya 37}
\pstart \ledsidenote{\passref{1.37.1}}\skt{अथातो }\edtext{\skt{मिश्रकन्नामाध्यायं}}{\Afootnote{H; \skt{मिश्रकमध्यायं}\ A1938; \skt{विस्रकन्नामाध्यायं}\ N}}\skt{ व्याख्यास्यामः ॥}\pend

\pstart\versepar \ledsidenote{\passref{1.37.3}}\edtext{\skt{मातुलुङ्गाग्निमन्थौ च}}{\Afootnote{N H; \skt{मातुलुङ्ग्यग्निमन्थौ च}\ A1938}}\skt{ भद्रदारु महौषधं । }\pend
\pstart\versepar \skt{}\edtext{\skt{अहिंसा}}{\Afootnote{H; \skt{अहिंस्रा}\ A1938; \skt{अहिय्रां}\ N}}\skt{ चैव }\edtext{\skt{रास्ना च}}{\Afootnote{A1938 H; \skt{रात्स्नाच}\ N}}\skt{ प्रलेपो }\edtext{\skt{वातशोफहा}}{\Afootnote{N H; \skt{वातशोफजित् ॥३॥}\ A1938}}\skt{ ॥}\pend

\pstart\versepar \ledsidenote{\passref{1.37.4}}\skt{दूर्व्वा च नलमूलञ्च }\edtext{\edtext{\skt{मधुकञ्चन्दनन्तथा}}{\lemma{\skt{चन्दनन्}}\Afootnote{A1938 H; \skt{चन्दन}\ N}}}{\lemma{\skt{तथा}}\Afootnote{A1938 H; \skt{त्तथा ।}\ N}}\skt{ । }\pend
\pstart\versepar \skt{शीतलाश्च गणाः सर्व्वे }\edtext{\skt{प्रलेपः}}{\Afootnote{A1938 H; \skt{प्रलेप}\ N}}\skt{ }\edtext{\skt{पित्तशोफहा}}{\Afootnote{N H; \skt{पित्तशोफहृत् ॥४॥}\ A1938}}\skt{ ॥}\pend

\pstart\versepar \ledsidenote{\passref{1.37.5}}\edtext{\skt{आगन्तुजे}}{\Afootnote{A1938 H; \skt{आगंतु}\ N}}\skt{ }\edtext{\skt{रक्तजे च}}{\Afootnote{A1938 H; \skt{जोवे}\ N}}\skt{ }\edtext{\skt{लेप}}{\Afootnote{N H; \skt{ह्येष}\ A1938}}\skt{ }\edtext{\edtext{\skt{एषोऽभिपूजितः}}{\lemma{\skt{एषो}}\Afootnote{H; \skt{ए}\ A1938; \skt{एषो}\ N}}}{\lemma{\skt{ऽभिपूजितः}}\Afootnote{H; \skt{व विधिः स्मृतः ।}\ A1938; \skt{भिपूजितः ।}\ N}}\skt{ । }\pend
\pstart\versepar \skt{}\edtext{\skt{विधिर्व्विषघ्नो}}{\Afootnote{H; \skt{विधिर्विषघ्नो}\ A1938 N}}\skt{ }\edtext{\skt{विषजे}}{\Afootnote{A1938 H; \skt{विषजो}\ N}}\skt{ }\edtext{\edtext{\skt{पित्तघ्नोऽभिहितस्तथा}}{\lemma{\skt{पित्तघ्नो}}\Afootnote{A1938 H; \skt{पित्तघ्नो}\ N}}}{\lemma{\skt{ऽभिहितस्}}\Afootnote{H; \skt{ऽपि हितस्}\ A1938; \skt{भिहितस्}\ N}}\skt{ ॥}\pend

\pstart\versepar \ledsidenote{\passref{1.37.6}}\edtext{\skt{अजगन्धाश्वगन्धा च}}{\Afootnote{A1938 H; \skt{अजगन्धा च}\ N}}\skt{ काला }\edtext{\skt{शरल}}{\Afootnote{H; \skt{सरल}\ A1938 N}}\skt{ }\edtext{\skt{एव च}}{\Afootnote{N H; \skt{या सह ।}\ A1938}}\skt{ । }\pend
\pstart\versepar \skt{}\edtext{\skt{एके}}{\Afootnote{N H; \skt{एकै}\ A1938}}\skt{ }\edtext{\skt{षिकाजशृङ्गी च}}{\Afootnote{A1938 H; \skt{षिकाङ्गसृङ्गीच}\ N}}\skt{ प्रलेपः }\edtext{\skt{श्लेष्मशोफहा}}{\Afootnote{N H; \skt{श्लेष्मशोफहृत् ॥६॥}\ A1938}}\skt{ ॥}\pend

\pstart\versepar \ledsidenote{\passref{1.37.7}}\edtext{\skt{एते च}}{\Afootnote{N H; \skt{एते व}\ A1938}}\skt{ }\edtext{\skt{वर्ग्गा}}{\Afootnote{N H; \skt{र्गास्त्रयो}\ A1938}}\skt{ }\edtext{\skt{लोध्रञ्च}}{\Afootnote{N H; \skt{लोध्रं}\ A1938}}\skt{ }\edtext{\skt{पथ्या}}{\Afootnote{A1938 H; \skt{पर्व्या}\ N}}\skt{ }\edtext{\skt{पिण्डीकृतास्तुयः}}{\Afootnote{H; \skt{पिण्डीतकानि च ।}\ A1938; \skt{पिण्डीकृताः त्रयः ।}\ N}}\skt{ । }\pend
\pstart\versepar \skt{अनन्ता चेति }\edtext{\edtext{\skt{लेपोऽयं}}{\lemma{\skt{लेपो}}\Afootnote{A1938 H; \skt{ले}\ N}}}{\lemma{\skt{ऽयं}}\Afootnote{A1938 H; \emph{om.}\ N}}\skt{ }\edtext{\skt{शोफेसर्व्वकृतेहितः}}{\Afootnote{\skt{सान्निपातिकशोफहृत् ॥७॥}\ A1938; \skt{पोयं शोफे सर्वकृते हितः ॥}\ N; \skt{… शोफेसर्व्वकृतेहितः ॥}\ H}}\skt{ ॥}\pend

\pstart\versepar \ledsidenote{\passref{1.37.8}}\edtext{\edtext{\skt{स्निग्धोऽम्ललवणो}}{\lemma{\skt{स्निग्धो}}\Afootnote{H; \skt{स्निग्धा}\ A1938 N}}}{\lemma{\skt{ऽम्ललवणो}}\Afootnote{H; \skt{म्ललवणो}\ A1938 N}}\skt{ }\edtext{\skt{वाते}}{\Afootnote{A1938 H; \skt{वातो}\ N}}\skt{ }\edtext{\skt{कोष्णः}}{\Afootnote{A1938 H; \skt{कोष्णाः}\ N}}\skt{ शीतः पयोयुतः । }\pend
\pstart\versepar \skt{}\edtext{\skt{पित्तेऽथोष्णेह्कफे}}{\lemma{\skt{ऽथोष्णेह्कफे}}\Afootnote{H; \skt{चोष्णः कफे}\ A1938; \skt{त्थोष्णः कफे}\ N}}\skt{ }\edtext{\skt{क्षारमूत्राढ्यस्तत्प्रशान्तये}}{\Afootnote{A1938 H; \skt{क्षारमूत्राद्यस्तत्प्रशान्तयेत् ॥}\ N}}\skt{ ॥}\pend

\pstart\versepar \ledsidenote{\passref{1.37.9}}\edtext{\skt{शणशिग्रुफलातसी}}{\Afootnote{H; \skt{शणमूलकशिग्रूणां फला}\ A1938; \skt{सणसिग्रुफलातसी}\ N}}\skt{ }\edtext{\skt{तिलकल्काश्च}}{\Afootnote{H; \skt{नि तिल}\ A1938; \skt{तिलकल्काच}\ N}}\skt{ }\edtext{\skt{सर्षषाः}}{\Afootnote{H; \skt{सर्षपाः ।}\ A1938 N}}\skt{ । }\pend
\pstart\versepar \skt{}\edtext{\skt{शक्तवः}}{\Afootnote{H; \skt{सक्तवः}\ A1938 N}}\skt{ }\edtext{\skt{किण्वमुष्णानि}}{\Afootnote{H; \skt{किण्वमतसी}\ A1938; \skt{किण्वामुष्णानि}\ N}}\skt{ }\edtext{\skt{द्रव्याण्यपि च}}{\Afootnote{H; \skt{द्रव्याण्युष्णानि}\ A1938; \skt{द्रव्यान्यिपिच}\ N}}\skt{ पाचनम् ॥}\pend

\pstart\versepar \ledsidenote{\passref{1.37.10}}\edtext{\skt{चिरिबिल्वाग्निको}}{\Afootnote{H; \skt{चिरबिल्वोऽग्निको}\ A1938; \skt{चिरिबिल्वाग्निको}\ N}}\skt{ दन्ती चित्रको हयमारकः । }\pend
\pstart\versepar \skt{कपोतकङ्कगृद्ध्राणाम्पुरीषाणि }\edtext{\skt{विदारणे}}{\Afootnote{N H; \skt{च दारणम् ।}\ A1938}}\skt{ ॥ }\pend
\pstart\versepar \skt{क्षारद्रव्याणि वा यानि क्षारो वा दारणं परं ।}\pend

\pstart\versepar \ledsidenote{\passref{1.37.11}}\edtext{\skt{द्रव्याणाम्पिच्छिलानां}}{\lemma{\skt{द्रव्याणाम्}}\Afootnote{A1938 H; \skt{द्रव्याण्यम्}\ N}}\skt{ तु }\edtext{\skt{त्वह्यूलानि}}{\Afootnote{H; \skt{त्वङ्मूलानि}\ A1938 N}}\skt{ प्रपीडनं ॥ }\pend
\pstart\versepar \skt{}\edtext{\skt{यवगोधूममाषाणां}}{\Afootnote{A1938 H; \skt{यवगोधूममाषाण्यञ्}\ N}}\skt{ चूर्ण्णानि च समासतः ।}\pend

\pstart\versepar \ledsidenote{\passref{1.37.12}}\edtext{\skt{शङ्खिन्यं}}{\Afootnote{N H; \skt{शाङ्खिन्यङ्}\ A1938}}\skt{ }\edtext{\skt{कोठसुमनाःकरवीरंसुवर्च्चला}}{\Afootnote{H; \skt{कोठसुमनःकरवीरसुवर्चलाः ।}\ A1938; \skt{कोठसुमनाः करवीरं सुवर्च्वला ॥}\ N}}\skt{ ॥ }\pend
\pstart\versepar \skt{शोधनानि कषायाणि वर्ग्गश्चारग्वधादिकः ।}\pend

\pstart\versepar \ledsidenote{\passref{1.37.13}}\skt{अजगन्धाजशृङ्गी च गवाक्षी लाङ्गलाह्वया ॥ }\pend
\pstart\versepar \skt{}\edtext{\skt{पूतीकश्चित्रकः}}{\lemma{\skt{चित्रकः}}\Afootnote{A1938 H; \skt{चित्रक}\ N}}\skt{ }\edtext{\skt{पाठा}}{\Afootnote{A1938 H; \skt{पोठा}\ N}}\skt{ }\edtext{\skt{विडङ्गैलाहरेणवः}}{\Afootnote{A1938 H; \skt{विडिगैलाहरेणवः ।}\ N}}\skt{ ।}\pend

\pstart\versepar \ledsidenote{\passref{1.37.14}}\edtext{\skt{कटुत्रिकं}}{\Afootnote{N H; \emph{om.}\ A1938}}\skt{ यवक्षारो }\edtext{\skt{लवणानि}}{\Afootnote{A1938 H; \skt{लवण्यनि}\ N}}\skt{ }\edtext{\skt{मनःशिला}}{\Afootnote{A1938 H; \skt{मनःसिलाः ॥}\ N}}\skt{ ॥ }\pend
\pstart\versepar \skt{}\edtext{\edtext{\skt{काशीसन्तृवृता}}{\lemma{\skt{काशीसन्}}\Afootnote{H; \skt{कासीसं}\ A1938 N}}}{\lemma{\skt{तृवृता}}\Afootnote{N H; \skt{त्रिवृता}\ A1938}}\skt{ दन्ती हरितालं }\edtext{\skt{सुराष्ट्रिका}}{\Afootnote{H; \skt{सुराष्ट्रजा ॥१४॥}\ A1938; \skt{सुराष्क्रिका ।}\ N}}\skt{ ।}\pend

\pstart\versepar \ledsidenote{\passref{1.37.15}}\skt{संशोधनीनां वर्त्तीनां द्रव्याण्येतानि निर्द्दिशेत् ॥ }\pend
\pstart\versepar \skt{}\edtext{\skt{द्रव्येष्वेतेषु}}{\lemma{\skt{एतेषु}}\Afootnote{H; \skt{एव उ}\ A1938; \skt{ष्वेतेषु}\ N}}\skt{ }\edtext{\skt{कुर्व्वीत}}{\Afootnote{N H; \skt{षधैः कु}\ A1938}}\skt{ }\edtext{\skt{कल्कानपि च}}{\Afootnote{H; \skt{र्यात्कल्कानपि च}\ A1938; \skt{कल्प्तानपिव}\ N}}\skt{ }\edtext{\skt{शोधनान्}}{\Afootnote{A1938 H; \skt{शोधनीं ॥}\ N}}\skt{ ।}\edtext{}{\lemma{\skt{निर्द्दिशेत् ॥ द्रव्येष्व्}}\Afootnote{H; \skt{निर्दिशेतेतैर्}\ A1938; \skt{निर्द्दिशेत् ॥ द्रव्ये}\ N}}\pend

\pstart\versepar \ledsidenote{\passref{1.37.16}}\edtext{\skt{अर्कोत्तमासुधाक्षीरं}}{\Afootnote{N H; \skt{अर्कोत्तमां स्नुहीक्षीरं}\ A1938}}\skt{ पिष्ट्वा }\edtext{\edtext{\skt{क्षारोत्तमामपि}}{\lemma{\skt{क्षारोत्तमाम्}}\Afootnote{N H; \skt{क्षारोत्तमा}\ A1938}}}{\lemma{\skt{अपि}}\Afootnote{N H; \skt{नि च ।}\ A1938}}\skt{ ॥ }\pend
\pstart\versepar \skt{जातीमूलं हरिद्रे द्वे कासीसं }\edtext{\skt{कटुरोहिणी}}{\Afootnote{N H; \skt{कटुरोहिणीम् ॥१६॥}\ A1938}}\skt{ ।}\pend

\pstart\versepar \ledsidenote{\passref{1.37.17}}\edtext{\skt{पूर्व्वोद्दिष्टेषु}}{\Afootnote{H; \skt{पूर्वोद्दिष्टानि}\ A1938; \skt{पूर्वोत्दिष्टेषु}\ N}}\skt{ }\edtext{\skt{चाङ्गेषु}}{\Afootnote{N H; \skt{चान्यानि}\ A1938}}\skt{ कुर्यात्संशोधनं घृतं ॥ }\pend
\pstart\versepar \skt{मयूरको राजवृक्षो निम्बः }\edtext{\skt{कोशातकी}}{\Afootnote{A1938 H; \skt{कोसातकी}\ N}}\skt{ }\edtext{\skt{तिला}}{\Afootnote{H; \skt{तिलाः}\ A1938 N}}\skt{ ।}\pend

\pstart\versepar \ledsidenote{\passref{1.37.18}}\skt{बृहती }\edtext{\skt{कण्टकारी च}}{\Afootnote{A1938 H; \skt{कष्काकारी च}\ N}}\skt{ }\edtext{\skt{हरितालम्मनःशिला}}{\lemma{\skt{मनःशिला}}\Afootnote{A1938 H; \skt{मनःसिला ॥}\ N}}\skt{ ॥ }\pend
\pstart\versepar \skt{शोधनानि च }\edtext{\skt{यानि स्युस्तैले}}{\lemma{\skt{यानि स्युस्}}\Afootnote{H; \skt{योज्यानि}\ A1938; \skt{यानिः स्युस्}\ N}}\skt{ }\edtext{\skt{योज्यानि}}{\Afootnote{H; \skt{द्रव्याणि}\ A1938; \skt{पोज्यानि}\ N}}\skt{ शोधने ।}\pend

\pstart\versepar \ledsidenote{\passref{1.37.19}}\skt{कासीसे सैन्धवे किण्वे }\edtext{\skt{वचायां}}{\Afootnote{A1938 H; \skt{वचाया}\ N}}\skt{ }\edtext{\skt{रजनीद्वये}}{\Afootnote{A1938 H; \skt{रर्ज्जनद्वये ॥}\ N}}\skt{ ॥ }\pend
\pstart\versepar \skt{शोधनाङ्गेषु चान्येषु चूर्ण्णं कुर्व्वीत शोधनं ।}\pend

\pstart\versepar \ledsidenote{\passref{1.37.20}}\edtext{\skt{सालसारादिसारेषु}}{\Afootnote{A1938 H; \skt{सालसारदिसारेषु}\ N}}\skt{ पटोलत्रिफलासु च ॥ }\pend
\pstart\versepar \skt{रसक्रिया विधातव्या शोधनी शोधनेषु च ।}\pend

\pstart\versepar \ledsidenote{\passref{1.37.21}}\skt{श्रीवेष्टके सर्ज्जरसे सरले देवदारुणि ॥ }\pend
\pstart\versepar \skt{सारेष्वपि च कुर्व्वीत }\edtext{\skt{मतिमान्व्रणधूपनं}}{\Afootnote{A1938 H; \skt{मतिमां व्रणधूपनंं ।}\ N}}\skt{ ।}\pend

\pstart\versepar \ledsidenote{\passref{1.37.22}}\edtext{\skt{कषायाणामनुष्णानान्द्र्व्याणान्त्वक्षु}}{\lemma{\skt{द्र्व्याणान्}}\Afootnote{H; \skt{वृक्षाणां}\ A1938; \skt{द्रव्याणान्}\ N}}\skt{ }\edtext{\skt{साधितं}}{\Afootnote{H; \skt{साधितः ।}\ A1938; \skt{साधाधितं ॥}\ N}}\skt{ ॥ }\pend
\pstart\versepar \skt{}\edtext{\skt{शृतशीत}}{\Afootnote{H; \skt{शृतः शीत}\ A1938; \skt{सृतसीत}\ N}}\skt{ }\edtext{\skt{कषायम्वा}}{\Afootnote{N H; \skt{कषायो वा}\ A1938}}\skt{ }\edtext{\skt{रोपनार्थे}}{\Afootnote{H; \skt{रोपणार्थं}\ A1938; \skt{रोपणार्थे}\ N}}\skt{ प्रशस्यते ।}\pend

\pstart\versepar \ledsidenote{\passref{1.37.23}}\skt{सोमामृताश्वगन्धासु काकोल्यादौ }\edtext{\skt{गणे}}{\Afootnote{A1938 H; \skt{गणेस्}\ N}}\skt{ तथा ॥ }\pend
\pstart\versepar \skt{}\edtext{\skt{क्षीरीप्ररोहेष्वपि}}{\lemma{\skt{क्षीरीप्ररोहेष्व्}}\Afootnote{N H; \skt{क्षीरिप्ररोहेष्व्}\ A1938}}\skt{ च }\edtext{\skt{वर्त्तयो}}{\Afootnote{A1938 H; \skt{वर्त्तयोर्}\ N}}\skt{ }\edtext{\skt{रोपणा}}{\Afootnote{H; \skt{रोपणाः}\ A1938; \skt{ओपणा}\ N}}\skt{ }\edtext{\skt{हिताः}}{\Afootnote{N H; \skt{स्मृताः ॥२३॥}\ A1938}}\skt{ ।}\pend

\pstart\versepar \ledsidenote{\passref{1.37.24}}\skt{समङ्गा }\edtext{\skt{सोमसरला}}{\Afootnote{A1938 H; \skt{सोमसरलाः}\ N}}\skt{ }\edtext{\skt{सोमवल्काः}}{\Afootnote{N H; \skt{सोमवल्कः}\ A1938}}\skt{ }\edtext{\skt{सचन्दनाः}}{\Afootnote{H; \skt{सचन्दनः ।}\ A1938; \skt{सचन्दना ॥}\ N}}\skt{ । }\pend
\pstart\versepar \skt{}\edtext{\skt{काकोल्यादिश्च}}{\Afootnote{A1938 H; \skt{काकोल्यादिश्च}\ N}}\skt{ }\edtext{\skt{कल्काः स्युः}}{\Afootnote{H; \skt{कल्कः स्या}\ A1938; \skt{कल्कास्युः}\ N}}\skt{ }\edtext{\skt{प्रशस्ता}}{\Afootnote{N H; \skt{त्प्रशस्तो}\ A1938}}\skt{ }\edtext{\skt{व्रणरोपणे}}{\Afootnote{A1938 H; \skt{प्रणरोपणे ।}\ N}}\skt{ ।}\pend

\pstart\versepar \ledsidenote{\passref{1.37.25}}\edtext{\skt{पृथक्पर्ण्ण्यात्मगुप्ता च}}{\Afootnote{A1938 H; \skt{पृथक्पण्यात्मगुप्ता च}\ N}}\skt{ हरिद्रे मालती सिता ॥ }\pend
\pstart\versepar \skt{काकोल्यादिश्च }\edtext{\skt{योज्याः स्युः}}{\Afootnote{H; \skt{योज्यः स्या}\ A1938; \skt{योज्यः स्युः}\ N}}\skt{ }\edtext{\skt{प्रशस्ता}}{\Afootnote{N H; \skt{द्भिषजा}\ A1938}}\skt{ रोपणे घृते ।}\pend

\pstart\versepar \ledsidenote{\passref{1.37.26}}\edtext{\skt{कालानुसार्यगुरुणी}}{\Afootnote{\skt{कालानुसार्यागुरुणी}\ A1938; \skt{कालातुसार्यगुरुणी}\ N; \skt{कालानुसार्यगुरुणी}\ H}}\skt{ }\edtext{\skt{हरिद्रा}}{\Afootnote{N H; \skt{हरिद्रे}\ A1938}}\skt{ }\edtext{\skt{भद्रदारु च}}{\Afootnote{N H; \skt{देवदारु च ।}\ A1938}}\skt{ ॥ }\pend
\pstart\versepar \skt{प्रियङ्गवश्च }\edtext{\skt{लोध्रञ्च}}{\Afootnote{N H; \skt{रोध्रं च}\ A1938}}\skt{ तैले योज्यानि रोपणे ।}\pend

\pstart\versepar \ledsidenote{\passref{1.37.27}}\edtext{\skt{कङ्गुक}}{\Afootnote{H; \skt{कङ्गुका}\ A1938 N}}\skt{ त्रिफला }\edtext{\skt{लोध्रं}}{\Afootnote{N H; \skt{रोध्रं}\ A1938}}\skt{ }\edtext{\skt{काशीसं}}{\Afootnote{H; \skt{कासीसं}\ A1938 N}}\skt{ }\edtext{\skt{श्रवणाह्वया}}{\Afootnote{A1938 H; \skt{स्रवणाह्वया ॥}\ N}}\skt{ ॥ }\pend
\pstart\versepar \skt{}\edtext{\skt{धवोश्वकर्ण्णयोस्त्वक्च}}{\lemma{\skt{धवोश्वकर्ण्णयोस्}}\Afootnote{H; \skt{धवाश्वकर्णयोस्}\ A1938 N}}\skt{ }\edtext{\skt{रोपणे}}{\Afootnote{N H; \skt{रोपणं}\ A1938}}\skt{ }\edtext{\skt{चूर्ण्णमिष्यते}}{\Afootnote{A1938 H; \skt{चूर्ण्णमिष्यते ।}\ N}}\skt{ ।}\pend

\pstart\versepar \ledsidenote{\passref{1.37.29}}\skt{त्वक्षु न्यग्रोधवर्ग्गस्य }\edtext{\skt{त्रिफलास्तथैव}}{\lemma{\skt{त्रिफलास्}}\Afootnote{H; \skt{त्रिफलायास्}\ A1938 N}}\skt{ च ॥ }\pend
\pstart\versepar \skt{रसक्रियां रोपणार्थे विदधीत यथाक्रमं ।}\pend

\pstart\versepar \ledsidenote{\passref{1.37.30}}\edtext{\edtext{\skt{अपामार्ग्गोऽस्वगन्धा}}{\lemma{\skt{अपामार्ग्गो}}\Afootnote{A1938 N\textsuperscript{pc} H; \skt{आपामार्गो}\ N\textsuperscript{ac}}}\skt{ च}}{\lemma{\skt{ऽस्वगन्धा च}}\Afootnote{H; \skt{ऽश्वगन्धा च}\ A1938 N}}\skt{ तालपत्त्री सुवर्च्चला ॥ }\pend
\pstart\versepar \skt{उत्सादने }\edtext{\skt{प्रशस्यन्ते}}{\Afootnote{A1938 H; \skt{प्रशस्यत्ते}\ N}}\skt{ }\edtext{\skt{काकोल्यादिश्च यो}}{\Afootnote{A1938 N H\textsuperscript{pc}; \skt{काकोल्यादिश्च}\ H\textsuperscript{ac}}}\skt{ गणः ।}\pend

\pstart\versepar \ledsidenote{\passref{1.37.31}}\skt{कासीसं }\edtext{\skt{सैन्धवन्कुत्थं}}{\Afootnote{H; \skt{सैन्धवं किण्वं}\ A1938; \skt{सेन्धवन्तुत्थं}\ N}}\skt{ }\edtext{\edtext{\skt{कुरुविन्दम्मनःशिला}}{\lemma{\skt{कुरुविन्दम्}}\Afootnote{N H; \skt{कुरुविन्दो}\ A1938}}}{\lemma{\skt{मनःशिला}}\Afootnote{A1938 H; \skt{मनःसिलाः ॥}\ N}}\skt{ ॥ }\pend
\pstart\versepar \skt{कुक्कुटाण्डकपालानि }\edtext{\skt{सुमनोमुकुलानि}}{\Afootnote{A1938 H; \skt{सुमनामुकुलानि}\ N}}\skt{ च ।}\pend

\pstart\versepar \ledsidenote{\passref{1.37.32}}\skt{फले }\edtext{\skt{शैरीषकारञ्जैर्द्धातुचूर्णानि}}{\lemma{\skt{शैरीषकारञ्जैर्}}\Afootnote{H; \skt{शैरीषकारञ्जे}\ A1938; \skt{सैरीषकारंजे}\ N}}\skt{ }\edtext{\skt{यानि च}}{\Afootnote{N H; \skt{यानि}\ A1938}}\skt{ ॥ }\pend
\pstart\versepar \skt{}\edtext{\skt{सुवर्च्चिकाथ}}{\Afootnote{H; \skt{सुवर्च्चिकाष}\ N}}\skt{ }\edtext{\skt{काशीसं}}{\Afootnote{H; \skt{कासीसं}\ N}}\skt{ सैन्धवं क्षारमेव च । }\pend
\pstart\versepar \skt{}\edtext{\skt{प्रणेषूत्सन्नमांसेषु}}{\Afootnote{N H; \skt{व्रणेषूत्सन्नमांसेषु}\ A1938}}\skt{ }\edtext{\edtext{\skt{प्रशस्तान्यवसादने}}{\lemma{\skt{प्रशस्तान्य्}}\Afootnote{N H; \skt{प्रशस्ता}\ A1938}}}{\lemma{\skt{अवसादने}}\Afootnote{N H; \skt{न्यवसादने ॥३२॥}\ A1938}}\skt{ ॥ }\pend
\pstart\versepar \skt{}\edtext{\skt{आलेपन}}{\Afootnote{H; \emph{om.}\ A1938; \skt{आलेपना}\ N}}\skt{ }\edtext{\skt{निचूर्ण्णानि}}{\Afootnote{N H; \emph{om.}\ A1938}}\skt{ }\edtext{\skt{प्रदद्यात्सकलानि}}{\Afootnote{H; \emph{om.}\ A1938; \skt{प्रदृद्यात्सकलानि}\ N}}\skt{ }\edtext{\skt{वै}}{\Afootnote{N H; \emph{om.}\ A1938}}\skt{ ।}\edtext{}{\lemma{\skt{च… च}}\Afootnote{\emph{om.}\ A1938}}\pend

\pstart\versepar \ledsidenote{\passref{1.37.33}}\skt{समस्तं वर्ग्गमर्द्धम्वा }\edtext{\edtext{\skt{यथालाभमथापि}}{\lemma{\skt{यथालाभम्}}\Afootnote{N H; \skt{यथालाभ}\ A1938}}\skt{ वा}}{\lemma{\skt{अथापि वा}}\Afootnote{N H; \skt{मथापि वा ।}\ A1938}}\skt{ ॥ }\pend
\pstart\versepar \skt{प्रयुञ्जीत भिषक्प्राज्ञो }\edtext{\skt{यथोद्दिष्टेषु}}{\Afootnote{A1938 H; \skt{यथोदिष्टेषु}\ N}}\skt{ }\edtext{\skt{कर्म्मसु}}{\Afootnote{A1938 H; \skt{कर्मषु ॥}\ N}}\skt{ ॥}\pend

\pstart \ledsidenote{\passref{1.37.trailer}}\edtext{\skt{३७}}{\Afootnote{\skt{इति सुश्रुतसंहितायां सूत्रस्थाने मिश्रकाध्यायो नाम सप्तत्रिंशत्तमोऽध्यायः ॥३७॥}\ A1938; \skt{ल भ ३० ९}\ N; \skt{ल भ्र ३०७}\ H}}\pend

\endnumbering
\beginnumbering
\eledsection*{Adhyāya 36}
\pstart \ledsidenote{\passref{1.36.1}}\skt{अथातो }\edtext{\skt{भूमिप्रविभागविज्ञानीयम्व्याख्यास्यामः}}{\lemma{\skt{भूमिप्रविभागविज्ञानीयम्}}\Afootnote{N H; \skt{भूमिप्रविभागीयमध्यायं}\ A1938}}\skt{ ॥}\pend

\pstart \ledsidenote{\passref{1.36.3}}\edtext{\edtext{\skt{स्वभ्रशर्क्कराश्मविषमवल्मीकश्मशानाघतदेवतायतनसिकतोषैररनुपहतामनूषरामभंगुरामदूरोदकां}}{\lemma{\skt{स्वभ्रशर्क्कराश्मविषमवल्मीकश्मशानाघतदेवतायतनसिकतोषैरर्}}\Afootnote{\skt{श्वभ्रशर्कराश्मविषमवल्मीकश्मशानाघातनदेवतायतनसिकताभिर्}\ A1938; \skt{श्वभ्रशर्क्कराश्मविषमवल्मीश्मशानघनदेवायतनसिकतोषरैर्}\ N; \skt{स्वभ्रशर्क्कराश्मविषमवल्मीकश्मशानाघतदेवतायतनसिकतोषैरर्}\ H}}}{\lemma{\skt{अनुपहतामनूषरामभंगुरामदूरोदकां}}\Afootnote{A1938 H; \skt{अनुपहतामत्क्रषरामभङ्गुरामदूरोदका}\ N}}\skt{ स्निग्धां }\edtext{\skt{प्ररोहवतीं}}{\Afootnote{A1938 H; \skt{प्ररोहवतीर्}\ N}}\skt{ }\edtext{\skt{मृद्वीं}}{\Afootnote{A1938 H; \skt{म्मृद्वी}\ N}}\skt{ }\edtext{\skt{स्थिरां}}{\Afootnote{A1938 H; \skt{स्थिरा}\ N}}\skt{ समां }\edtext{\skt{गौरीं}}{\Afootnote{N H; \skt{कृष्णां गौरीं}\ A1938}}\skt{ लोहितां वा भूमिमौषध ग्रहणाय परीक्षेत ॥ }\edtext{\edtext{\skt{तस्याङ्जातमपि}}{\lemma{\skt{जातम्}}\Afootnote{N H; \skt{जात}\ A1938}}\skt{ हि}}{\lemma{\skt{अपि हि}}\Afootnote{N H; \skt{मपि}\ A1938}}\skt{ }\edtext{\edtext{\skt{क्रिमविषशस्त्रातपपवनदहनतोयसम्बाधमार्ग्गैरनुपहतमेकसारं}}{\lemma{\skt{क्रिमविषशस्त्रातपपवनदहनतोयसम्बाधमार्ग्गैर्}}\Afootnote{H; \skt{कृमिविषशस्त्रातपपवनदहनतोयसम्बाधमार्गैर्}\ A1938; \skt{क्रिमिविषशस्त्रातपपवनदहनतोयसम्बाधमार्गैर्}\ N}}}{\lemma{\skt{अनुपहतमेकसारं}}\Afootnote{H; \skt{अनुपहतमेकरसं}\ A1938; \skt{अनुपहतमेकसारम्}\ N}}\skt{ }\edtext{\skt{पुष्टमवगाढमूलं}}{\lemma{\skt{अवगाढमूलं}}\Afootnote{N H; \skt{पृथ्ववगाढमूल}\ A1938}}\skt{ }\edtext{\edtext{\skt{चौषधमादद्यादित्येषां}}{\lemma{\skt{चौषधमादद्यादित्य्}}\Afootnote{H; \skt{मुदीच्यां चौषधमाददीते}\ A1938; \skt{चौषधमादद्यादित्य्}\ N}}}{\lemma{\skt{एषां}}\Afootnote{N H; \skt{त्येष}\ A1938}}\skt{ }\edtext{\skt{भूमिपरीक्षाविभागः}}{\Afootnote{N H; \skt{भूमिपरीक्षाविशेषः}\ A1938}}\skt{ सामान्यः ॥}\pend

\pstart \ledsidenote{\passref{1.36.4}}\skt{विशेषतस्तु }\edtext{\skt{तत्र}}{\Afootnote{A1938 H; \skt{तत्रा}\ N}}\skt{ }\edtext{\skt{अश्मवती}}{\Afootnote{A1938 H; \skt{श्मवती}\ N}}\skt{ स्थिरा गुर्व्वी }\edtext{\skt{श्धामा}}{\Afootnote{H; \skt{श्यामा}\ A1938; \skt{स्यामा}\ N}}\skt{ कृष्णा वा }\edtext{\skt{स्थूलतृणशस्यप्रायाः}}{\Afootnote{H; \skt{स्थूलवृक्षशस्यप्राया}\ A1938; \skt{स्थूलतृणसस्य प्रायाः ।}\ N}}\skt{ स्वगुणभूयिष्ठा । स्निग्धा शीतलासन्नोदका }\edtext{\skt{स्निग्धतृणशस्याकोमलवृक्षप्राया}}{\Afootnote{H; \skt{स्निग्धशस्यतृणकोमलवृक्षप्राया}\ A1938; \skt{स्निग्धभृणसस्याकोमलवृक्षप्राया}\ N}}\skt{ }\edtext{\skt{शुक्लाम्बुगुणभूयिष्ठा}}{\Afootnote{A1938 H; \skt{शुल्काचाम्बुगुणभूयिष्ठा ॥}\ N}}\skt{ ॥ }\edtext{\skt{नानावर्ण्णा}}{\Afootnote{A1938 H; \skt{नानावर्ण्ण}\ N}}\skt{ }\edtext{\skt{लघ्वाश्मवती}}{\Afootnote{H; \skt{लघ्वश्मवती}\ A1938; \skt{लप्म्व्यश्मवती}\ N}}\skt{ }\edtext{\skt{विरलाल्यपाण्डुरप्ररोहाग्निगुणभूयिष्ठा}}{\Afootnote{H; \skt{प्रविरलाल्पपाण्डुवृक्षप्ररोहाऽग्निगुणभूयिष्ठा;}\ A1938; \skt{विरलाल्पपाण्डुरप्ररोहाग्निगुणभूयिष्ठा ।}\ N}}\skt{ ॥ }\edtext{\skt{रूक्षा}}{\Afootnote{A1938 H; \skt{रूक्ष}\ N}}\skt{ भस्मरासभवर्ण्णा }\edtext{\skt{क्षतनुवृक्षकोटराल्परसप्रायानिलगुणभूयिष्ठा}}{\Afootnote{H; \skt{तनुवृक्षाल्परसकोटरवृक्षप्रायाऽनिलगुणभूयिष्ठा;}\ A1938; \skt{रूक्षतनुवृक्षकोटराल्परसप्रायानिलगुणभूयिष्ठा ॥}\ N}}\skt{ ॥ मृद्वी }\edtext{\skt{समाव्यभ्रवत्यव्यक्तरसजला}}{\Afootnote{N H; \skt{समा श्वभ्रवत्यव्यक्तरसजला}\ A1938}}\skt{ }\edtext{\skt{महापर्व्वतवृक्षप्रार्या}}{\Afootnote{H; \skt{सर्वतोऽसारवृक्षा महापर्वतवृक्षप्राया}\ A1938; \skt{महापर्वतवृक्षः । प्राया}\ N}}\skt{ }\edtext{\skt{स्यामा}}{\Afootnote{N H; \skt{श्यामा}\ A1938}}\skt{ }\edtext{\skt{चाकाशगुणभूयिष्ठा}}{\Afootnote{A1938 H; \skt{चाकाशगुणभूयिष्ठं ॥}\ N}}\skt{ ॥}\pend

\pstart \ledsidenote{\passref{1.36.5}}\edtext{\skt{तत्र}}{\Afootnote{N H; \skt{अत्र}\ A1938}}\skt{ }\edtext{\skt{केचिदाचार्याः}}{\Afootnote{H; \skt{केचिदाहुराचार्याः-}\ A1938; \skt{केचिदावायाः}\ N}}\skt{ }\edtext{\skt{प्रावृट्वर्षाशरद्धेमन्तवसन्तग्रीष्मेषु}}{\Afootnote{H; \skt{प्रावृड्वर्षाशरद्धेमन्तवसन्तग्रीष्मेषु}\ A1938; \skt{प्रावृद्वर्षाशरद्धेमन्तेवसन्ते ग्रीष्मेषु}\ N}}\skt{ }\edtext{\edtext{\skt{यथासंख्यम्मूलपत्त्रत्वचक्षीरसारफलान्याददीत}}{\lemma{\skt{मूलपत्त्रत्वचक्षीरसारफलान्य्}}\Afootnote{H; \skt{मूलपत्रत्वक्क्षीरसारफलान्य्}\ A1938; \skt{मूलपत्रत्वचक्षीरसारफला}\ N}}}{\lemma{\skt{आददीत}}\Afootnote{H; \skt{आददीतेति;}\ A1938; \skt{न्याददीत ॥}\ N}}\skt{ ॥ तत्तु न }\edtext{\skt{सम्यक्कस्मात्सौम्याग्नेयत्वाज्जगतः}}{\Afootnote{H; \skt{सम्यक्, सौम्याग्नेयत्वाज्जगतः ।}\ A1938; \skt{सम्यक् । कस्मात्सौम्याग्नेयत्वाज्जागतः ।}\ N}}\skt{ । सौम्यान्यौषधानि }\edtext{\skt{सौम्येष्वृतुष्वाददीत}}{\lemma{\skt{आददीत}}\Afootnote{A1938 H; \skt{आददीत् । आ}\ N}}\skt{ । }\edtext{\edtext{\edtext{\skt{आग्नेयान्याग्नेयेष्वेवमव्यापन्नानि}}{\lemma{\skt{आग्नेयान्य्}}\Afootnote{A1938 H; \skt{ग्नेयाः}\ N}}}{\lemma{\skt{आग्नेयेष्व्}}\Afootnote{A1938 H; \skt{न्याग्नेयेष्व्}\ N}}}{\lemma{\skt{अव्यापन्नानि}}\Afootnote{N H; \skt{अव्यापन्नगुणानि}\ A1938}}\skt{ भवन्ति ॥ }\edtext{\skt{सौम्यान्यौषधानि}}{\lemma{\skt{औषधानि}}\Afootnote{H; \skt{औषध्मनि}\ A1938; \skt{ओषधानि}\ N}}\skt{ सौम्येष्वृतुषु गृहीतानि सोमगुणभूयिष्ठायां भूमौ }\edtext{\skt{मधुरतरास्निग्धतराणिशीततराणि}}{\Afootnote{N H\textsuperscript{pc}; \skt{जातान्यतिमधुरस्निग्धशीतानि}\ A1938; \skt{मधुरतआरास्निग्धतराणिशीततराणि}\ H\textsuperscript{ac}}}\skt{ }\edtext{\skt{भवन्ति}}{\Afootnote{\skt{जायन्ते । एतेन शेषं व्याख्यातम् ॥५॥}\ A1938; \skt{शीततराणि च भवन्ति ॥}\ N; \skt{भवन्ति । …}\ H}}\skt{ ।}\pend

\pstart \ledsidenote{\passref{1.36.5.add1}}\edtext{\skt{आग्नेयान्यौषधानि}}{\lemma{\skt{औषधानि}}\Afootnote{H; \skt{ओषधानि}\ N}}\skt{ आग्नेयेष्वृतुषु }\edtext{\skt{गृहीतान्यग्निगुणभूयिष्ठायां}}{\Afootnote{H; \skt{गृहीतान्यग्निगुणभूयिष्ठायाम्}\ N}}\skt{ भूमौ उष्णतराणि कटुतराणि रूक्षतराणि भवन्ति ।}\pend

\pstart \ledsidenote{\passref{1.36.6}}\edtext{\skt{तत्रा}}{\Afootnote{N H; \skt{तत्र}\ A1938}}\skt{ }\edtext{\skt{पृथिव्ῑगुणभूयिष्ठायां}}{\Afootnote{H; \skt{पृथिव्यम्बुगुणभूयिष्ठायां}\ A1938; \skt{प्पृथिवीगुणभूयिष्ठायाम्}\ N}}\skt{ भूमौ जातानि }\edtext{\skt{विरेचनद्रव्याण्याददीत}}{\Afootnote{H; \skt{विरेचनद्रव्याण्याददीत ,}\ A1938; \skt{विरेचनदव्याण्याददीत ॥}\ N}}\skt{ । }\edtext{\skt{अग्न्यनिलगुणभूयिष्ठायाञ्च}}{\Afootnote{N H; \skt{अग्न्याकाशमारुतगुणभूयिष्ठायां}\ A1938}}\skt{ वमनद्रव्याणि । }\edtext{\edtext{\skt{उभयगुणभूयिष्ठायामुभयतोभागहराणि}}{\lemma{\skt{उभयगुणभूयिष्ठायाम्}}\Afootnote{A1938 H; \skt{उभयसुणभूयिष्ठायां}\ N}}}{\lemma{\skt{उभयतोभागहराणि}}\Afootnote{N H; \skt{उभयतो भागानि,}\ A1938}}\skt{ । आकाशगुणभूयिष्ठायां }\edtext{\edtext{\skt{संशमनान्येवम्बलवत्तराणि}}{\lemma{\skt{संशमनान्य्}}\Afootnote{A1938 H; \skt{संसमना}\ N}}}{\lemma{\skt{एवम्}}\Afootnote{A1938 H; \skt{न्येवम्}\ N}}\skt{ भवन्ति ।}\pend

\pstart \ledsidenote{\passref{1.36.7}}\edtext{\skt{सर्व्वाणि}}{\Afootnote{N H; \skt{सर्वाण्येव}\ A1938}}\skt{ }\edtext{\skt{चाभिनवान्यत्र}}{\Afootnote{N H; \skt{चाभिनवानि, अन्यत्र}\ A1938}}\skt{ }\edtext{\skt{मधुघृतपिप्पलीविडङ्गेभ्य}}{\Afootnote{N H; \skt{मधुघृतगुडपिप्पलीविडङ्गेभ्यः ॥७॥}\ A1938}}\skt{ }\edtext{\skt{इति}}{\Afootnote{N H; \emph{om.}\ A1938}}\skt{ ॥}\pend

\pstart\versepar \ledsidenote{\passref{1.36.8}}\edtext{\skt{विडङ्गपिप्पलीक्षौद्रं}}{\Afootnote{N H; \skt{विडङ्गं पिप्पली क्षौद्रं}\ A1938}}\skt{ }\edtext{\edtext{\skt{सर्पिश्चाप्यनवं}}{\lemma{\skt{सर्पिश्}}\Afootnote{A1938 H; \skt{सपिश्}\ N}}}{\lemma{\skt{चाप्यनवं}}\Afootnote{H\textsuperscript{pc}; \skt{चाप्यनवं}\ A1938 N; \skt{चानवं}\ H\textsuperscript{ac}}}\skt{ हितं । }\pend
\pstart\versepar \skt{}\edtext{\skt{शेषमन्यत्त्वभिनवं}}{\lemma{\skt{अन्यत्त्वभिनवं}}\Afootnote{H; \skt{अन्यत्वभिनवं}\ A1938 N}}\skt{ }\edtext{\skt{गृह्णीयाद्दोषवर्ज्जितं}}{\Afootnote{H; \skt{गृह्णीयाद्दोषवर्जितम् ॥८॥}\ A1938; \skt{गृह्णीया दोषवर्जितं ॥}\ N}}\skt{ ॥}\pend

\pstart\versepar \ledsidenote{\passref{1.36.16}}\skt{जङ्गमानां }\edtext{\edtext{\skt{वयःस्थानाञ्चर्मरोमनखादिकं}}{\lemma{\skt{वयःस्थानाञ्}}\Afootnote{A1938 H; \skt{वयस्थानाञ्}\ N}}}{\lemma{\skt{चर्मरोमनखादिकं}}\Afootnote{H; \skt{रक्तरोमनखादिकम् ।}\ A1938; \skt{चर्मरोमनखादिक्रं ।}\ N}}\skt{ । }\pend
\pstart\versepar \skt{क्षीरमूत्रपुरीषाणि जीर्ण्णाहारेषु संहरेत् ॥}\pend

\pstart\versepar \ledsidenote{\passref{1.36.17}}\edtext{\skt{प्लोतमृद्भाण्डफलकं}}{\Afootnote{N H; \skt{प्लोतमृद्भाण्डफलक}\ A1938}}\skt{ सङ्कुविन्यस्तभेषजान् । }\pend
\pstart\versepar \skt{प्रशस्तायान्दिशि शुचौ }\edtext{\edtext{\skt{भेषजागारमिष्यते}}{\lemma{\skt{भेषजागारम्}}\Afootnote{N H; \skt{भेषजागार}\ A1938}}}{\lemma{\skt{इष्यते}}\Afootnote{N H; \skt{मिष्यते ॥१७॥}\ A1938}}\skt{ ॥}\edtext{}{\lemma{\skt{सङ्कुविन्यस्तभेषजान् । प्रशस्तायान्}}\Afootnote{H; \skt{शङ्कुविन्यस्तभेषजम् । प्रशस्तायां}\ A1938; \skt{संकुर्विन्यस्तभेषजान् । प्रशस्तायान्}\ N}}\pend

\pstart\versepar \ledsidenote{\passref{1.36.10}}\edtext{\skt{गोपालास्तापसा}}{\lemma{\skt{गोपालास्}}\Afootnote{N H; \skt{भवन्ति चात्र - गोपालास्}\ A1938}}\skt{ व्याधा ये }\edtext{\skt{चान्ये}}{\Afootnote{A1938 H; \skt{वान्ये}\ N}}\skt{ वनचारिणः । }\pend
\pstart\versepar \skt{मूलाहाराश्च ये }\edtext{\skt{तेषां}}{\Afootnote{N H; \skt{तेभ्यो}\ A1938}}\skt{ }\edtext{\edtext{\skt{भेषजव्यक्तिरिष्यते}}{\lemma{\skt{भेषजव्यक्तिर्}}\Afootnote{A1938 H; \skt{भेषज्यव्यक्तर्}\ N}}}{\lemma{\skt{इष्यते}}\Afootnote{A1938 H; \skt{इष्यत}\ N}}\skt{ }\edtext{\skt{इति}}{\Afootnote{N H; \emph{om.}\ A1938}}\skt{ ॥}\pend

\pstart \ledsidenote{\passref{1.36.trailer}}\edtext{\skt{३७}}{\Afootnote{\skt{इति सुश्रुतसंहितायां सूत्रस्थाने भूमिप्रविभागीयो नाम षट्त्रिंशोऽध्यायः ॥३६॥}\ A1938; \skt{उन्देfइनेद्}\ N; \skt{ल ऊ ३० ७}\ H}}\pend

\endnumbering
\beginnumbering
\eledsection*{Adhyāya 38}
\pstart \ledsidenote{\passref{1.38.1}}\skt{अथातो }\edtext{\skt{द्रव्यसङ्ग्रहणीयम्व्याख्यास्यामः}}{\lemma{\skt{द्रव्यसङ्ग्रहणीयम्}}\Afootnote{H; \skt{द्रवयसंग्रहणीयमध्यायं}\ A1938; \skt{द्रवयसङ्ग्राहणीयमध्यायं}\ N}}\skt{ ॥}\pend

\pstart \ledsidenote{\passref{1.38.3}}\edtext{\skt{पञ्चत्रिंशद्द्रव्यगणा}}{\Afootnote{H; \skt{समासेन सप्तत्रिंशद्द्रव्यगणा}\ A1938; \skt{पञ्चत्रिङ्शद्रव्यगणा}\ N}}\skt{ भवन्ति ॥}\pend

\pstart \ledsidenote{\passref{1.38.4}}\skt{तद्यथा । विदारिगन्धा विदारीविश्वदेवा }\edtext{\skt{सहदेवाश्वदंष्ट्रा}}{\Afootnote{A1938 H; \skt{सहदेवाश्च द्रंष्टा}\ N}}\skt{ }\edtext{\skt{पृथक्पर्ण्णी}}{\Afootnote{A1938 H; \skt{पृषक्यण्णी}\ N}}\skt{ }\edtext{\skt{शतावरीए}}{\Afootnote{H; \skt{शतावरी सारिवा कृष्ण}\ A1938; \skt{सतावरी}\ N}}\skt{ }\edtext{\skt{शारिवा}}{\Afootnote{H; \skt{सारिवा}\ A1938 N}}\skt{ }\edtext{\skt{जीवकर्षभो}}{\Afootnote{H; \skt{जीवकर्षभकौ म}\ A1938; \skt{जीवकऋषभकौ क्षुद्रस}\ N}}\skt{ }\edtext{\skt{क्षुद्रसहा}}{\Afootnote{H; \skt{हासहा क्षुद्रसहा}\ A1938; \skt{हा महासहा}\ N}}\skt{ बृहत्यौ }\edtext{\skt{पुनर्न्नैवेरण्डहंसपादीवृश्चिकाल्यृषभी}}{\Afootnote{H; \skt{पुनर्नवैरण्डो हंसपादी वृश्चिकाल्यृषभी}\ A1938; \skt{पुनर्ण्णवैरण्डहंसपादी वृश्चिकाल्मृषभी}\ N}}\skt{ }\edtext{\skt{शृगालविन्ना}}{\Afootnote{N H; \emph{om.}\ A1938}}\skt{ चेति ॥}\pend

\pstart\versepar \ledsidenote{\passref{1.38.5}}\edtext{\skt{विदारिगन्धादिरयङ्गणः}}{\lemma{\skt{विदारिगन्धादिरयङ्}}\Afootnote{N H; \skt{विदारीगन्धादिरयं}\ A1938}}\skt{ पित्तानिलापहः । }\pend
\pstart\versepar \skt{}\edtext{\skt{शोफगुल्माङ्गमर्द्दोर्द्ध्वश्वासकासविनाशनः}}{\Afootnote{H; \skt{शोषगुल्माङ्गमर्दोर्ध्वश्वासकासविनाशनः ॥५॥}\ A1938; \skt{शोफगुल्माङ्गमर्द्दोध्वं श्वासकासविनाशनः ॥}\ N}}\skt{ ॥}\pend

\pstart \ledsidenote{\passref{1.38.6}}\edtext{\skt{आरग्वधमदनफलगोपघोण्टाकण्टकी कण्टकारीकुटजफलपाठापाटलीमूर्व्वासप्तपर्ण्णनिम्बकुरणअटकगुडूचीचित्रकशार्ङ्गष्टाद्विकरञ्जपटोलकिराततिक्तकाः}}{\Afootnote{H\textsuperscript{pc}; \skt{आरग्वधमदनगोपघोण्टाकण्टकीकुटजपाठापाटलामूर्वेन्द्रयवसप्तपर्णनिम्बकुरुण्टकदासीकुरुण्टक- गुडूचीचित्रकशार्ङ्ग(र्ङ्गे)ष्टाकरञ्जद्वयपटोलकिराततिक्तकानि}\ A1938; \skt{आरग्वधमदनफलगोपघोण्टकण्टकारीकुटाजफलपाठापाटलीमूर्वासप्तपर्ण्णनिम्बकुरण्टकगुडूचीचित्रकशार्ङ्गष्वाद्विकरण्डपटोलकिराततिक्तकाः}\ N; \skt{आरग्वधमदनफलगोपघोण्टा कण्टकारीकुटजफलपाठापाटलीमूर्व्वासप्तपर्ण्णनिम्बकुरणअटकगुडूचीचित्रकशार्ङ्गष्टाद्विकरञ्जपटोलकिराततिक्तकाः}\ H\textsuperscript{ac}}}\skt{ सुषवी चेति ॥}\pend

\pstart\versepar \ledsidenote{\passref{1.38.7}}\edtext{\skt{आरग्वधादिरित्येष}}{\lemma{\skt{आरग्वधादिरित्य्}}\Afootnote{H; \skt{आरग्वधादिरित्य्}\ A1938 N}}\skt{ गणः श्लेष्मविषापहः । }\pend
\pstart\versepar \skt{}\edtext{\skt{मेहकुष्टज्वरवमीकण्डूघ्नो}}{\Afootnote{H; \skt{मेहकुष्ठज्वरवमीकण्डूघ्नो}\ A1938; \skt{मेहकुष्वज्वरवमीकण्डूघ्नो}\ N}}\skt{ व्रणशोधनः ॥}\pend

\pstart \ledsidenote{\passref{1.38.8}}\edtext{\skt{शालसाराजकर्ण्णखदिरक्रमुभूर्ज्जमेषशृङ्गीति}}{\Afootnote{H; \skt{सालसाराजकर्णखदिरकदरकालस्कन्धक्रमुकभूर्जमेषशृङ्गति}\ A1938; \skt{सालसाराजकर्ण्णखदिरक्रमुकभूर्जामेषशृङ्गीति}\ N}}\skt{ }\edtext{\skt{निसचन्दनकुचन्दनशिंशपा}}{\Afootnote{H; \skt{निशचन्दनकुचन्दनशिंशपा-}\ A1938; \skt{षचन्दनकुचन्दनशिंशपा}\ N}}\skt{ }\edtext{\skt{शिरीषासनधवार्ज्जुननक्तमालपूतीकाश्च}}{\Afootnote{H; \skt{शिरीषासनधवार्जुनतालशाकनक्तमालपूतीकाश्व}\ A1938; \skt{सिरीषासनधर्वार्जुननक्तमालपूतीकाश्व}\ N}}\skt{ }\edtext{\skt{कर्ण्णगुरूणि}}{\Afootnote{H; \skt{कर्णागुरूणि}\ A1938; \skt{कर्ण्णआगुरूणि}\ N}}\skt{ }\edtext{\skt{कालीयकञ्चेति}}{\lemma{\skt{चेति}}\Afootnote{A1938 H; \skt{चेत् ॥}\ N}}\skt{ ॥}\pend

\pstart\versepar \ledsidenote{\passref{1.38.9}}\edtext{\skt{शालसारादिरित्येष}}{\lemma{\skt{शालसारादिर्}}\Afootnote{H; \skt{सालसारादिर्}\ A1938 N}}\skt{ गणः }\edtext{\skt{कुष्ठविनाशनः}}{\Afootnote{A1938 H; \skt{कुष्ठपिनाशनः ।}\ N}}\skt{ । }\pend
\pstart\versepar \skt{मेहपाण्ड्वामयहरः }\edtext{\skt{कफमेदोविशोषणः}}{\Afootnote{A1938 N; \skt{कफमेदो विशोषणः ॥}\ H}}\skt{ ॥}\pend

\pstart \ledsidenote{\passref{1.38.10}}\edtext{\skt{वरुणार्त्तगलमधुशिग्रुतर्क्कारीमेषशृङ्गीपूतीकनक्तमालपाटलाग्निमन्थशैरीयकद्वयं}}{\Afootnote{H; \skt{वरुणार्तगलशिग्रुमधुशिग्रुतर्कारीमेषशृङ्गीपूतीकनक्तमालमोरटाग्निमन्थसैरेयकद्वय}\ A1938; \skt{वरुणार्त्तगलमधुशिग्रुतर्क्कारीमेषशृङ्गीपूतीकनक्वक्तमालपाटलग्निमन्थशैरीषकद्वयम्}\ N}}\skt{ }\edtext{\edtext{\skt{बिम्बीचित्रशतावरीबिल्वाजशृङ्गीदर्ब्भाबृहतीद्वयञ्चेति}}{\lemma{\skt{बिम्बीचित्रशतावरीबिल्वाजशृङ्गीदर्ब्भाबृहतीद्वयञ्}}\Afootnote{H; \skt{बिम्बीवसुकवसिरचित्रक- शतावरीबिल्वाजशृङ्गीदर्भा बृहतीद्वयं}\ A1938; \skt{चिम्बीचित्रसतावरीबिल्वाजशृङ्गीदर्भाबृहतीद्वयञ्}\ N}}}{\lemma{\skt{चेति}}\Afootnote{A1938 H; \skt{चेत् ॥}\ N}}\skt{ ॥}\pend

\pstart\versepar \ledsidenote{\passref{1.38.11}}\edtext{\skt{वरुणादिगणोप्येष}}{\Afootnote{H; \skt{वरुणादिर्गणो ह्येष}\ A1938 N}}\skt{ }\edtext{\skt{कफमेदो}}{\Afootnote{H; \skt{कफमेदो}\ A1938 N}}\skt{ }\edtext{\skt{विशोषणः}}{\Afootnote{N H; \skt{निवारणः ।}\ A1938}}\skt{ । }\pend
\pstart\versepar \skt{विनिहन्ति }\edtext{\skt{शिरः}}{\Afootnote{A1938 H; \skt{सिरः}\ N}}\skt{ }\edtext{\skt{शूलं}}{\Afootnote{H; \skt{शूल}\ A1938; \skt{सूलं}\ N}}\skt{ }\edtext{\skt{गुल्माभ्यन्तरविद्रधीः}}{\Afootnote{H; \skt{गुल्माभ्यन्तरविद्रधीन् ॥११॥}\ A1938; \skt{गुल्माभ्यन्तरविद्रधी ॥}\ N}}\skt{ ॥}\pend

\pstart \ledsidenote{\passref{1.38.12}}\edtext{\skt{वीरतरसहचरद्वयसैरीयकदर्ब्भवृक्षादनीगुन्द्रानलकुशकाशाश्मभेदकाग्नि}}{\Afootnote{H; \skt{वीरतरुसहचरद्वयदर्भवृक्षादनीगुन्द्रानलकुशकाशाश्मभेदकाग्नि}\ A1938; \skt{वीरतरसहचरद्वयसैरीयकदर्भवृक्षादनीगुन्द्रानलकुशकासास्मभेदकाग्नि}\ N}}\skt{ }\edtext{\skt{मन्थमोरटवसुकविसरभल्लूककुरूटिकेन्दीवरकपोतवङ्काश्वदंष्ट्रा}}{\Afootnote{H\textsuperscript{pc}; \skt{मन्थमोरटावसुकवसिरभल्लूककुरण्टिकेन्दीवरकपोतवङ्काः श्वदंष्ट्रा}\ A1938; \skt{मन्थमोरटवसुकवसिककुरूटिकैन्दीवरकपोतवङ्काश्च दंष्ट्रा}\ N; \skt{मन्थमोरटवसुकविसरकुरूटिकेन्दीवरकपोतवङ्काश्वदंष्ट्रा}\ H\textsuperscript{ac}}}\skt{ चेति ॥}\pend

\pstart\versepar \ledsidenote{\passref{1.38.13}}\edtext{\skt{वीरतरादिरित्येष}}{\lemma{\skt{वीरतरादिरित्य्}}\Afootnote{\skt{वीरतर्वादिरित्य्}\ A1938; \skt{वीरतरादिरित्य्}\ N H}}\skt{ गणो वातविकारहृत् । }\pend
\pstart\versepar \skt{शर्क्कराश्मरिहामूत्रकृच्छ्राघातरुजापहः ॥}\edtext{}{\lemma{\skt{वातविकारहृत् । शर्क्कराश्मरिहामूत्रकृच्छ्राघातरुजापहः}}\Afootnote{H; \skt{वातविकारनुत् । अश्मरीशर्करामूत्रकृच्छ्राघातरुजापहः ॥१३॥}\ A1938; \skt{वातविकारहृत् ॥ शर्कराश्मरिहामूत्रकृच्छ्रघातरुजापहः ॥}\ N}}\pend

\pstart \ledsidenote{\passref{1.38.14}}\edtext{\skt{लोध्रशावरलोध्रपलाशकुटन्नटाशोककट्फलैलवालुशल्लकीजिङ्जिणीक}}{\Afootnote{H; \skt{लोध्रसावरलोध्रपलाशकुटन्नटाशोकफञ्जीकट्फलैलवालुकशल्लकीजिङ्गिनीक}\ A1938; \skt{लोध्रसवरलोध्रपलाशकुटन्नटाशोककफ्हलैलवालुकसल्लकीझिञ्झिणीक}\ N}}\skt{ }\edtext{\skt{दम्बाकदली}}{\Afootnote{N H; \skt{दम्बसालाः कदली}\ A1938}}\skt{ चेति ॥}\pend

\pstart\versepar \ledsidenote{\passref{1.38.15}}\skt{एष }\edtext{\skt{लोध्रादिको}}{\Afootnote{N H; \skt{रोध्रादिरि}\ A1938}}\skt{ }\edtext{\skt{नाम्ना}}{\Afootnote{H; \skt{त्युक्तो}\ A1938; \skt{नास्ना}\ N}}\skt{ मेदः }\edtext{\skt{कफहरोगणः}}{\Afootnote{H; \skt{कफहरो गणः ।}\ A1938 N}}\skt{ । }\pend
\pstart\versepar \skt{}\edtext{\skt{योनिदोषहरस्तम्भी}}{\Afootnote{N H; \skt{योनिदोषहरः स्तम्भी}\ A1938}}\skt{ }\edtext{\skt{वण्यो}}{\Afootnote{H; \skt{वर्ण्यो}\ A1938; \skt{वर्ल्यो}\ N}}\skt{ विषविनाशनः ॥}\pend

\pstart \ledsidenote{\passref{1.38.16}}\edtext{\skt{अर्क्कालर्क्ककरञ्जद्वयनागदन्तीमयूरकभार्ग्गीरास्नेन्द्रपुष्पीक्षुद्रस्वेताभल्लूकमहास्वेतावृश्चिकाल्या}}{\Afootnote{H; \skt{अर्कालर्ककरञ्जद्वयनागदन्तीमयूरकभार्गीरास्नेन्द्रपुष्पीक्षुद्रश्वेतामहाश्वेतावृश्चिकाल्य}\ A1938; \skt{अर्कालर्ककरंजद्वयनागदन्तीमयूरकभार्गीरास्तेन्द्रपुष्पीक्षुद्रश्वेताभल्लूकमहाश्वेतावृश्चिकाल्या}\ N}}\skt{ }\edtext{\skt{लवणातापसवृक्षश्चेति}}{\lemma{\skt{लवणातापसवृक्षश्}}\Afootnote{N H; \skt{लवणास्तापसवृक्षश्}\ A1938}}\skt{ ॥}\pend

\pstart\versepar \ledsidenote{\passref{1.38.17}}\edtext{\skt{अर्क्कादिस्तु}}{\Afootnote{H; \skt{अर्कादिको}\ A1938; \skt{अर्क्कीदिस्तु}\ N}}\skt{ गणो ह्येष }\edtext{\skt{कफमेदोविषापहः}}{\Afootnote{A1938 H; \skt{केफमेदोविषापहः ॥}\ N}}\skt{ ॥ }\pend
\pstart\versepar \skt{}\edtext{\skt{क्रिमिकुष्ठप्रशमनो}}{\Afootnote{H; \skt{कृमिकुष्ठप्रशमनो}\ A1938; \skt{कृमिकुष्टप्रशमनो}\ N}}\skt{ }\edtext{\skt{विशेषाद्व्रणशोधनः}}{\Afootnote{A1938 H; \skt{विशेषोद्ध्रेणशोधनः ॥}\ N}}\skt{ ॥}\pend

\pstart \ledsidenote{\passref{1.38.18}}\edtext{\skt{सुरसास्वेतसुरसाफणिझ्जकार्ज्जकभूस्तृणसुगन्धककालमालकुठेरकक्षवकखरपुसविडङ्ग ङ्गपाठलेतिकट्फलसुरसीनिर्ग्गुण्डीफुलुफलोन्दुरुकर्ण्णी}}{\Afootnote{H\textsuperscript{pc}; \skt{सुरसाश्वेतसुरसाफणिज्झकार्जकभूस्तृणसुगन्धकसुमुखकालमालकुठेरककासमर्दक्षवकखरपुष्पाविडङ्गकट्फल- सुरसीनिर्गुण्डीकुलाहलोन्दुरुकर्णि}\ A1938; \skt{षरमः श्वेतसुरसः फणिज्झकार्जकभूस्तृणसुगन्धककालमालकुठोरकक्षवकखरपुष्पविडङ्गकप्फलसुरसीनिगुण्ठीफुलुह्रलोन्दुरुपर्ण्णी}\ N; \skt{सुरसास्वेतसुरसाफणिझ्जकार्ज्जकभूस्तृणसुगन्धककालमालकुठेरकक्षवकखरपुसविडङ्ग कट्फलसुरसीनिर्ग्गुण्डीफुलुफलोन्दुरुकर्ण्णी}\ H\textsuperscript{ac}}}\skt{ }\edtext{\skt{फञ्जीप्रजीबलाकाकमाच्यो}}{\Afootnote{H; \skt{काफञ्जीप्राचीबलकाकमाच्यो}\ A1938; \skt{फंजीप्रजीवलाकाकमाव्यो}\ N}}\skt{ }\edtext{\skt{विषमुष्टिश्चेति}}{\lemma{\skt{विषमुष्टिश्}}\Afootnote{N H; \skt{विषमुष्टिकश्}\ A1938}}\skt{ ॥}\pend

\pstart\versepar \ledsidenote{\passref{1.38.19}}\skt{सुरसादिर्गणो ह्येष }\edtext{\skt{कफहृत्क्रिमिसूदनः}}{\Afootnote{H; \skt{कफहृत्कृमिसूदनः ।}\ A1938 N}}\skt{ । }\pend
\pstart\versepar \skt{}\edtext{\skt{प्रतिस्यायारुचिः}}{\Afootnote{N H; \skt{प्रतिश्यायारुचि}\ A1938}}\skt{ }\edtext{\skt{कासश्वासघ्नो}}{\Afootnote{N H; \skt{श्वासकासघ्नो}\ A1938}}\skt{ व्रणशोधनः ॥}\pend

\pstart \ledsidenote{\passref{1.38.20}}\edtext{\skt{मुष्ककपलाशधवचित्रकमदनवृक्षशिंशपावज्रवृक्षस्त्रिफला}}{\Afootnote{H; \skt{मुष्ककपलाशधवचित्रकमदनवृक्षकशिंशपावज्रवृक्षास्त्रिफला}\ A1938; \skt{मुष्ककपलासएवचित्रकमदनवृक्षसिंसपाव्रज्रवृक्षस्त्रिफला}\ N}}\skt{ चेति ॥}\pend

\pstart\versepar \ledsidenote{\passref{1.38.21}}\edtext{\edtext{\skt{मुष्ककादिर्ग्गणो}}{\lemma{\skt{मुष्ककादिर्}}\Afootnote{H; \skt{मुष्ककादि}\ A1938 N}}}{\lemma{\skt{ग्गणो}}\Afootnote{H; \skt{र्गणो}\ A1938; \skt{र्गणो}\ N}}\skt{ ह्येष मेदोघ्नः }\edtext{\skt{शुक्रदोषहा}}{\Afootnote{N H; \skt{शुक्रदोषहृ}\ A1938}}\skt{ । }\pend
\pstart\versepar \skt{}\edtext{\skt{मेहार्शः}}{\Afootnote{H; \skt{त् । मेहार्शः}\ A1938; \skt{मेहार्सः}\ N}}\skt{ }\edtext{\skt{पाण्डुरोगघ्नः}}{\Afootnote{N H; \skt{पाण्डुरोगाश्म}\ A1938}}\skt{ }\edtext{\skt{शर्क्राश्मरिनाशनः}}{\Afootnote{H; \skt{शर्करानाशनः परः ॥२१॥}\ A1938; \skt{शर्करार्स्मरिनाशनः ॥}\ N}}\skt{ ॥}\pend

\pstart \ledsidenote{\passref{1.38.22}}\edtext{\skt{पिप्पलीपिप्पलीमूलंचव्यचित्रकशृङ्गवेरमरिचहरेणुकैलाजमोदेन्द्रयवपाठाजीरकसर्षप}}{\Afootnote{H; \skt{पिप्पलीपिप्पलीमूलचव्यचित्रकशृङ्गवेरमरिचहस्तिपिप्पलीहरेणुकैलाजमोदेन्द्रयवपाठाजीरकसर्षप-}\ A1938; \skt{पिप्पली पिप्पलीमूलञ्च व्यचित्रकशृङ्गवेरमरिचहरेण्वकैलाजमोदेन्द्रयवपाठाजीरकसर्षप}\ N}}\skt{ }\edtext{\skt{महानिम्बहिंगुभार्ग्गामधुरसातिविषाविडङ्गकटुरोहिणी}}{\Afootnote{H; \skt{महानिम्बफलहिङ्गुभार्गीमधुरसातिविषावचाविडङ्गानि कटुरोहिणी}\ A1938; \skt{माहानिम्बहिङ्गुभार्गीमधुरसातिविषाविडिङ्गकटुरोहिणी}\ N}}\skt{ चेति ॥}\pend

\pstart\versepar \ledsidenote{\passref{1.38.23}}\skt{पिप्पल्यादिः कफहरः }\edtext{\skt{प्रतिस्याय}}{\Afootnote{N H; \skt{प्रतिश्याया}\ A1938}}\skt{ }\edtext{\skt{मरीचकं}}{\Afootnote{H; \skt{निलारु}\ A1938; \skt{मरोचकं ॥}\ N}}\skt{ । }\pend
\pstart\versepar \skt{}\edtext{\edtext{\skt{अनिलञ्चापि}}{\lemma{\skt{अनिलञ्}}\Afootnote{H; \skt{चीः । निह}\ A1938; \skt{अ}\ N}}}{\lemma{\skt{चापि}}\Afootnote{H; \skt{न्याद्दीपनो गुल्मशू}\ A1938; \skt{निलञ्चापि}\ N}}\skt{ }\edtext{\skt{गुल्मञ्च}}{\Afootnote{N H; \skt{लघ्नश्चा}\ A1938}}\skt{ }\edtext{\skt{दीपनस्त्वामपाचनः}}{\Afootnote{N H; \skt{मपाचनः ॥२३॥}\ A1938}}\skt{ ॥}\pend

\pstart \ledsidenote{\passref{1.38.24}}\edtext{\skt{एलातगरकुष्ठमांसीध्यामकत्वपत्त्रनागपुष्पप्रियङ्गुहरेणुकाव्याघ्रनखशुक्तिचण्डास्थौणेयक}}{\Afootnote{H; \skt{एलातगरकुष्ठमांसीध्यामकत्वक्पत्रनागपुष्पप्रियङ्गुहरेणुकाव्याघ्रनखशुक्तिचण्डास्थौणेयक}\ A1938; \skt{एलातगरकुष्ठमांसीध्यामकत्वक्पत्रनागपुष्पप्रियङ्गुहरेणुकाव्याघ्रनखशुक्तिचण्डास्थौणैयक}\ N}}\skt{ }\edtext{\skt{श्रीवेष्टकचोरकवालकगुग्गुलुसर्ज्जरसतुरुष्ककुन्दुरुकागुरुस्पृक्काभद्रदारुकुङ्कुमानि}}{\Afootnote{H; \skt{श्रीवेष्टक- चोचचोरकवालुकगुग्गुलुसर्जरसतुरुष्ककुन्दुरुकागुरुस्पृक्कोशीरभद्रदारुकुङ्कुमानि}\ A1938; \skt{श्रीवेष्टकचोवचोरकवालकगुग्गुलुसर्ज्जरसतुरष्ककुन्दुरुकागुरुस्पृर्क्काभद्रदारुकुङ्कुमानि}\ N}}\skt{ }\edtext{\skt{पुन्नागकेसरञ्चेति}}{\Afootnote{N H; \skt{पुन्नागकेशरं चेति ॥२४॥}\ A1938}}\skt{ ॥}\pend

\pstart\versepar \ledsidenote{\passref{1.38.25}}\edtext{\skt{एलादिको}}{\Afootnote{A1938 H; \skt{एलदिकौ}\ N}}\skt{ }\edtext{\skt{वातकफो}}{\Afootnote{H; \skt{वातकफौ}\ A1938 N}}\skt{ }\edtext{\skt{निहन्याद्विषमेव}}{\lemma{\skt{निहन्याद्विषम्}}\Afootnote{N H; \skt{निहन्याद्विषम्}\ A1938}}\skt{ च । }\pend
\pstart\versepar \skt{वर्ण्णप्रसादनः }\edtext{\skt{कण्डूपिटका कोठनाशनः}}{\Afootnote{N H\textsuperscript{pc}; \skt{कण्डूपिडकाकोठनाशनः ॥२५॥}\ A1938; \skt{कण्डूपिट कोठनाशनः ॥}\ H\textsuperscript{ac}}}\skt{ ॥}\pend

\pstart \ledsidenote{\passref{1.38.26}}\edtext{\skt{वचा}}{\Afootnote{A1938 H; \skt{एचा}\ N}}\skt{ }\edtext{\edtext{\skt{मुस्तातिविषाभद्रदारुनागकेसरञ्चेति}}{\lemma{\skt{मुस्तातिविषाभद्रदारुनागकेसरञ्}}\Afootnote{N H; \skt{मुस्तातिविषाभयाभद्रदारूणि नागकेशरं}\ A1938}}}{\lemma{\skt{चेति}}\Afootnote{A1938 H; \skt{चेत् ॥}\ N}}\skt{ ।}\pend

\pstart \ledsidenote{\passref{1.38.27}}\edtext{\skt{हरिद्रादारुहरिद्राकलशीकुटजबीजानि}}{\Afootnote{A1938 H; \skt{हरिद्रादारुहरिद्रातलसीकुटजबीजानि}\ N}}\skt{ मधुकं चेति ॥}\pend

\pstart\versepar \ledsidenote{\passref{1.38.28}}\skt{एतौ वचा हरिद्रादी गणौ स्तन्यविशोधनौ । }\pend
\pstart\versepar \skt{}\edtext{\skt{आमातीसारशमनौ}}{\Afootnote{H; \skt{आमातिसारशमनौ}\ A1938 N}}\skt{ }\edtext{\skt{कफमेदो}}{\Afootnote{H; \emph{om.}\ A1938; \skt{कफनमेदो}\ N}}\skt{ }\edtext{\skt{विशोषणौ}}{\Afootnote{N H; \skt{विशेषाद्दोषपाचनौ ॥२८॥}\ A1938}}\skt{ ।}\pend

\pstart \ledsidenote{\passref{1.38.29}}\edtext{\skt{श्यामामहाश्यामातृवृद्दन्तीबिल्वकम्पिल्यकरम्यककेवुकपुत्रश्रेणीगवाक्षीराजवृक्षकरञ्जद्वय}}{\Afootnote{H; \skt{श्यामामहाश्यामात्रिवृद्दन्तीशङ्खिनीतिल्वककम्पिल्लकरम्यकक्रमुकपुत्रश्रेणीगवाक्षी- राजवृक्षकरञ्जद्वय}\ A1938; \skt{श्यामामहाश्यामातृवृद्दन्तीतिल्वकम्पिल्यकरम्यककेवुकपत्रश्रेणी गवाक्षीराजवृक्षकरञ्जद्वय}\ N}}\skt{ }\edtext{\skt{गुडूचीसप्तवर्ण्णच्छगलान्त्रीपीलुस्नुहाक्षीसुवर्ण्णक्षीरी}}{\Afootnote{H; \skt{गुडूचीसप्तलाच्छगलान्त्रीसुधाः सुवर्णक्षीरी}\ A1938; \skt{गुडूचीसप्तपर्ण्णच्छगलान्त्रीपीलुस्तुहासुवर्ण्णक्षीरी}\ N}}\skt{ चेति ॥}\pend

\pstart\versepar \ledsidenote{\passref{1.38.30}}\edtext{\skt{एष}}{\Afootnote{N H; \skt{उक्तः}\ A1938}}\skt{ }\edtext{\skt{श्यामादिको}}{\Afootnote{N H; \skt{श्यामादिरित्येष}\ A1938}}\skt{ }\edtext{\skt{नाम्ना}}{\Afootnote{N H; \emph{om.}\ A1938}}\skt{ गणो गुल्मविषापहः । }\pend
\pstart\versepar \skt{}\edtext{\skt{आनाहोदरहाभेदी}}{\Afootnote{N H; \skt{आनाहोदरविड्भेदी}\ A1938}}\skt{ तथोदावर्त्तनाशनः ॥}\pend

\pstart \ledsidenote{\passref{1.38.31}}\edtext{\skt{बृहतीकण्टकारिकाकुटजफलपाठामधुकञ्चेति}}{\Afootnote{A1938 H; \skt{बृहतीकण्ठकारिकाकुटजफलपाठा मधुकञ्चेत् ॥}\ N}}\skt{ ॥}\pend

\pstart\versepar \ledsidenote{\passref{1.38.32}}\skt{पाचनीयो }\edtext{\edtext{\skt{बृहत्यादिर्ग्गणः}}{\lemma{\skt{बृहत्यादिर्}}\Afootnote{H; \skt{बृहत्यादि}\ A1938 N}}}{\lemma{\skt{ग्गणः}}\Afootnote{H; \skt{र्गणः}\ A1938 N}}\skt{ पित्तानिलापहः । }\pend
\pstart\versepar \skt{}\edtext{\skt{कफारोचककृद्रोगमूत्रकृच्छ्रविनाशनः}}{\Afootnote{H; \skt{कफारोचकहृद्रोगमूत्रकृच्छ्ररुजापहः ॥३२॥}\ A1938; \skt{कफारोचकहृद्रोगमूत्रकृच्छ्रविनाशनः ॥}\ N}}\skt{ ॥}\pend

\pstart \ledsidenote{\passref{1.38.33}}\edtext{\skt{पटोलाचन्दनमूर्व्वागुडूचीपाठा}}{\Afootnote{N H; \skt{पटोलचन्दनकुचन्दनमूर्वागुडूचीपाठाः}\ A1938}}\skt{ कटुरोहिणी चेति ॥}\pend

\pstart\versepar \ledsidenote{\passref{1.38.34}}\edtext{\edtext{\skt{पटोलादिर्ग्गणः}}{\lemma{\skt{पटोलादिर्}}\Afootnote{H; \skt{पटोलादि}\ A1938 N}}}{\lemma{\skt{ग्गणः}}\Afootnote{H; \skt{र्गणः}\ A1938 N}}\skt{ }\edtext{\skt{पित्तकफारोचकनाशनः}}{\Afootnote{A1938 H; \skt{पित्तः कफरोचकनाशनः ।}\ N}}\skt{ । }\pend
\pstart\versepar \skt{ज्वरोपशमनो }\edtext{\edtext{\skt{व्रण्यश्छर्द्दिकण्डूविषापहः}}{\lemma{\skt{व्रण्यश्}}\Afootnote{H; \skt{व्रण्य}\ A1938 N}}}{\lemma{\skt{छर्द्दिकण्डूविषापहः}}\Afootnote{H; \skt{श्छर्दिकण्डूविषापहः ॥३४॥}\ A1938; \skt{च्छद्दीकण्डूविषोपहः ॥}\ N}}\skt{ ॥}\pend

\pstart \ledsidenote{\passref{1.38.35}}\edtext{\skt{काकोलीक्षीरकाकोलीजीवकऋषभकमुद्गपर्ण्णीमेदामहामेदाच्छिन्नरुहाकर्क्कटाशृङ्गीतुगाक्षीपद्मकप्रपौण्डरीक}}{\Afootnote{H; \skt{काकोलीक्षीरकाकोलीजीवकर्षभकमुद्गपर्णीमाषपर्णीमेदामहामेदाच्छिन्नरुहा कर्कटशृङ्गीतुगाक्षीरीपद्मकप्रपौण्डरीक}\ A1938; \skt{काकोलीक्षीरकाकोलीजीवकऋषभकमुद्गपण्णीमाषपल्लीमेदोमहामेदाच्छिन्नरुहाकर्क्कटाशृङ्गीतुगाक्षीरीपद्मकप्रणण्डरीक}\ N}}\skt{ }\edtext{\skt{ऋद्धिवृद्धिजीवन्त्यौ}}{\Afootnote{H; \skt{र्धिवृद्धिमृद्वीकाजीवन्त्यो}\ A1938; \skt{ऋद्धिवृद्धिजीवन्ती}\ N}}\skt{ }\edtext{\skt{मधुकञ्चेति}}{\lemma{\skt{चेति}}\Afootnote{A1938 H; \skt{चेत् ॥}\ N}}\skt{ ॥}\pend

\pstart\versepar \ledsidenote{\passref{1.38.36}}\edtext{\skt{काकोल्यादिरयम्पित्तशोणितानिलनाशनः}}{\lemma{\skt{काकोल्यादिरयम्}}\Afootnote{A1938 H; \skt{काकोल्यादिरपं}\ N}}\skt{ । }\pend
\pstart\versepar \skt{जीवनो }\edtext{\skt{बृंहणो}}{\Afootnote{A1938 H; \skt{बृंहनो}\ N}}\skt{ }\edtext{\skt{वृष्यः}}{\Afootnote{A1938 H; \skt{वृष्य}\ N}}\skt{ स्तन्यश्लेष्मकरस्तथा ॥}\pend

\pstart \ledsidenote{\passref{1.38.37}}\edtext{\edtext{\skt{ऊषसैन्धवशिलाजतुकाशीसद्वयहिंगुतुत्थकञ्चेति}}{\lemma{\skt{ऊषसैन्धवशिलाजतुकाशीसद्वयहिंगुतुत्थकञ्}}\Afootnote{H; \skt{ऊषकसैन्धवशिलाजतुकासीसद्वयहिङ्गूनि तुत्थकं}\ A1938; \skt{ऊषसैन्धवसिलाजतुकासीसद्वयहिङ्गूतुत्थकञ्}\ N}}}{\lemma{\skt{चेति}}\Afootnote{A1938 H; \skt{चेत् ॥}\ N}}\skt{ ॥}\pend

\pstart\versepar \ledsidenote{\passref{1.38.38}}\skt{ऊषकादिः कफं हन्ति गणो }\edtext{\skt{मेदा}}{\Afootnote{H; \skt{मेदो}\ A1938 N}}\skt{ विशोषणः । }\pend
\pstart\versepar \skt{}\edtext{\skt{शर्क्कराश्मरिहामूत्रकृच्छ्रगुल्मप्रमर्द्दनः}}{\Afootnote{N H; \skt{अश्मरीशर्करामूत्रकृच्छ्रगुल्मप्रणाशनः ॥३८॥}\ A1938}}\skt{ ॥}\pend

\pstart \ledsidenote{\passref{1.38.39}}\edtext{\edtext{\skt{शारिवामधुकचन्दनकुचन्दनपद्मककाश्मर्यमधूकपुष्पाणअयुशीरञ्चेति}}{\lemma{\skt{शारिवामधुकचन्दनकुचन्दनपद्मककाश्मर्यमधूकपुष्पाणअयुशीरञ्}}\Afootnote{H; \skt{सारिवामधुकचन्दनकुचन्दनपद्मककाश्मरीफलमधूकपुष्पाण्युशीरं}\ A1938; \skt{सारिवामधुकचन्दनकुचन्दनपद्मककास्मपमधूकपुष्पाण्युशीरञ्}\ N}}}{\lemma{\skt{चेति}}\Afootnote{A1938 H; \skt{चेत् ॥}\ N}}\skt{ ॥}\pend

\pstart\versepar \ledsidenote{\passref{1.38.40}}\edtext{\skt{शारिवादिः}}{\Afootnote{H; \skt{सारिवादिः}\ A1938 N}}\skt{ }\edtext{\skt{पिपासार्त्ति}}{\Afootnote{H; \skt{पिपासाघ्नो}\ A1938; \skt{पिपासन्ति}\ N}}\skt{ रक्तपित्तहरो गणः । }\pend
\pstart\versepar \skt{पित्तज्वरप्रशमनो }\edtext{\skt{विशेषाद्दाहनाशनः}}{\Afootnote{H; \skt{विशेषाद्दाहनाशनः}\ A1938 N}}\skt{ ॥}\pend

\pstart \ledsidenote{\passref{1.38.41}}\edtext{\skt{अञ्जनरसाञ्जननागपुष्पप्रियङ्गुनलिनकेसरं}}{\Afootnote{N H; \skt{अञ्जनरसाञ्जननागपुष्पप्रियङ्गुनीलोत्पलनलदनलिनकेशराणि मधुकं}\ A1938}}\skt{ }\edtext{\skt{चेति}}{\Afootnote{A1938 H; \skt{चेत् ॥}\ N}}\skt{ ॥}\pend

\pstart\versepar \ledsidenote{\passref{1.38.42}}\edtext{\edtext{\skt{अञ्जनादिर्न्नेत्ररोग}}{\lemma{\skt{अञ्जनादिर्}}\Afootnote{H; \skt{अञ्जनादि}\ A1938 N}}}{\lemma{\skt{न्नेत्ररोग}}\Afootnote{H; \skt{र्गणो ह्येष}\ A1938; \skt{र्नेत्ररोग}\ N}}\skt{ रक्तपित्तनिबर्हणः । }\pend
\pstart\versepar \skt{}\edtext{\skt{विषोपशमनो}}{\Afootnote{A1938 H; \skt{विशोपशमनो}\ N}}\skt{ दाहं }\edtext{\skt{हन्यादभ्यन्तरं}}{\Afootnote{H; \skt{निहन्त्याभ्यन्तरं}\ A1938; \skt{हन्यादाभ्यन्तरन्}\ N}}\skt{ }\edtext{\skt{नृणां}}{\Afootnote{N H; \skt{भृशम् ॥४२॥}\ A1938}}\skt{ ॥}\pend

\pstart \ledsidenote{\passref{1.38.43}}\edtext{\skt{फरूषकद्राक्षाकट्फलराजादनदाडिमकतकफलानि}}{\Afootnote{H; \skt{परूषकद्राक्षाकट्फलदाडिमराजादनकतकफलशाकफलानि}\ A1938; \skt{फरूषकद्राक्षाकप्फलराजादनदाडिमकतकफलानितृफलानि}\ N}}\skt{ }\edtext{\skt{त्रिफला}}{\Afootnote{A1938 H; \emph{om.}\ N}}\skt{ चेति ॥}\pend

\pstart\versepar \ledsidenote{\passref{1.38.44}}\edtext{\skt{फरूषकादिको}}{\Afootnote{N H; \skt{परूषकादिरि}\ A1938}}\skt{ }\edtext{\skt{नाम्ना}}{\Afootnote{N H; \skt{त्ये}\ A1938}}\skt{ }\edtext{\skt{गण}}{\Afootnote{N H; \skt{ष ग}\ A1938}}\skt{ }\edtext{\edtext{\skt{एषोऽनिलापहः}}{\lemma{\skt{एषो}}\Afootnote{H; \skt{णो}\ A1938; \skt{ए}\ N}}}{\lemma{\skt{ऽनिलापहः}}\Afootnote{H; \skt{ऽनिलविनाशनः ।}\ A1938; \skt{षोनिलापहः ।}\ N}}\skt{ । }\pend
\pstart\versepar \skt{मूत्रदोषहरो }\edtext{\skt{हृद्यः}}{\Afootnote{A1938 H; \skt{हृद्य}\ N}}\skt{ }\edtext{\edtext{\skt{पिपासघ्नोऽरुचो}}{\lemma{\skt{पिपासघ्नो}}\Afootnote{N H; \skt{पिपासाघ्नो}\ A1938}}}{\lemma{\skt{ऽरुचो}}\Afootnote{H; \skt{रुचि}\ A1938; \skt{रुचो}\ N}}\skt{ }\edtext{\skt{हितः}}{\Afootnote{N H; \skt{प्रदः ॥४४॥}\ A1938}}\skt{ ॥}\pend

\pstart \ledsidenote{\passref{1.38.45}}\edtext{\skt{प्रियङ्गुसमङ्गाधातकीपुन्नागनागपुष्यचन्दनमोचरसाञ्जनकुम्भीकपद्मकेसरयोजनवल्यो}}{\Afootnote{H; \skt{प्रियङ्गुसमङ्गाधातकीपुन्नागनागपुष्पचन्दनकुचन्दनमोचरसरसाञ्जनकुम्भीकस्रोतोजपद्मकेसरयोजनवल्ल्यो}\ A1938; \skt{प्रियङ्गुसमङ्गाधातकीपुन्नागनागपुष्पचन्दनमोचरसाञ्जनकुम्भीकापद्मकेरसयोजनवल्यो}\ N}}\skt{ दीर्घमूला चेति ॥}\pend

\pstart \ledsidenote{\passref{1.38.46}}\edtext{\skt{अम्बष्ठाधातकीकुसुमसमङ्गाकट्वङ्गमधुकबिल्वपेसिकाशवरलोध्रपलाशनन्दीवृक्षाः}}{\Afootnote{H; \skt{अम्बष्ठाधातकीकुसुमसमङ्गाकट्वङ्गमधुकबिल्वपेशिकासावररोध्रपलाशनन्दीवृक्षाः}\ A1938; \skt{अम्बष्ठाधातकीकुसुमसमङ्गाकट्वाङ्गमधुकबिल्वपेसिकासवरलोध्रपलाशनन्दीवृक्षाः}\ N}}\skt{ }\edtext{\edtext{\skt{पद्मकेशरञ्चेति}}{\lemma{\skt{पद्मकेशरञ्}}\Afootnote{H; \skt{पद्मकेशराणि}\ A1938; \skt{पद्मकेसरञ्}\ N}}}{\lemma{\skt{चेति}}\Afootnote{A1938 H; \skt{चेत् ॥}\ N}}\skt{ ॥}\pend

\pstart\versepar \ledsidenote{\passref{1.38.47}}\skt{गणौ }\edtext{\skt{प्रिययङ्ग्वम्बष्ठादी}}{\Afootnote{H; \skt{प्रियङ्ग्वम्बष्ठादी}\ A1938; \skt{प्रियङ्ग्वाम्बष्ठादी}\ N}}\skt{ }\edtext{\skt{पक्वातीसारनाशनौ}}{\Afootnote{A1938 H; \skt{पत्कातीसारनाशनौ ।}\ N}}\skt{ । }\pend
\pstart\versepar \skt{सन्धानीयौ हितौ }\edtext{\skt{पित्ते}}{\Afootnote{A1938 H; \skt{पित्तौ}\ N}}\skt{ }\edtext{\skt{व्रणानाञ्चापि}}{\lemma{\skt{व्रणानाञ्}}\Afootnote{A1938 H; \skt{प्रेणानाञ्}\ N}}\skt{ रोपणौ ॥}\pend

\pstart \ledsidenote{\passref{1.38.48}}\edtext{\skt{न्यग्रोधोदुम्बरोश्वत्थप्लक्षमधूकककुभाम्रम्बूद्वयपियालरोहिणीवञ्जलक}}{\Afootnote{H; \skt{न्यग्रोधोदुम्बराश्वत्थप्लक्षमधुककपीतनककुभाम्रकोशाम्रचोरकपत्रजम्बूद्वयप्रियाल- मधूकरोहिणीवञ्जुलक}\ A1938; \skt{न्यग्रोधोदुम्बराश्वत्थप्लक्षमधूकककुभाम्व्रजंबूद्वयपियालरोहिणीवञ्जुलक}\ N}}\skt{ }\edtext{\skt{कदम्बबदरीतिन्दुकलोध्रपलाशनन्दीवृक्षाश्चेति}}{\Afootnote{H; \skt{दम्बबदरीतिन्दुकीसल्लकीरोध्रसावररोध्रभल्लातकपलाशा नन्दीवृक्षश्चेति ॥४८॥}\ A1938; \skt{दम्बयदरीतिन्दुकलोध्रपलाशनन्दीवृक्षाचेति ॥}\ N}}\skt{ ॥}\pend

\pstart\versepar \ledsidenote{\passref{1.38.49}}\edtext{\edtext{\skt{न्यग्रोधादिर्ग्गणो}}{\lemma{\skt{न्यग्रोधादिर्}}\Afootnote{H; \skt{न्यग्रोधादि}\ A1938 N}}}{\lemma{\skt{ग्गणो}}\Afootnote{H; \skt{र्गणो}\ A1938; \skt{गणो}\ N}}\skt{ }\edtext{\skt{व्रण्यः}}{\Afootnote{A1938 H; \skt{व्रणः}\ N}}\skt{ संग्राही भग्नसाधकः । }\pend
\pstart\versepar \skt{रक्तपित्तहरो }\edtext{\skt{दाहमेदघ्नो}}{\Afootnote{N H; \skt{दाहमेदोघ्नो}\ A1938}}\skt{ }\edtext{\skt{योनिदोषहा}}{\Afootnote{N H; \skt{योनिदोषहृत् ॥४९॥}\ A1938}}\skt{ ॥}\pend

\pstart \ledsidenote{\passref{1.38.50}}\skt{गुडूचीनिम्बकुस्तुम्बुरुचन्दनानि पद्मकञ्चेति ॥}\pend

\pstart\versepar \ledsidenote{\passref{1.38.51}}\skt{एष }\edtext{\skt{सर्व्वज्वरान्हन्ति}}{\Afootnote{A1938 H; \skt{सर्वज्वरां हन्ति}\ N}}\skt{ गुडूच्यादिस्तु दीपनः । }\pend
\pstart\versepar \skt{}\edtext{\skt{हल्लासोरोचकवमी}}{\Afootnote{H; \skt{हृल्लासारोचकवमी}\ A1938 N}}\skt{ पिपासा दाहनाशनः ॥}\pend

\pstart \ledsidenote{\passref{1.38.52}}\edtext{\skt{उत्पलरक्तोत्पलकुमुदकुवलयसौगन्धिकपुण्डरीकानि}}{\Afootnote{N H; \skt{उत्पलरक्तोत्पलकुमुदसौगन्धिककुवलयपुण्डरीकाणि}\ A1938}}\skt{ मधुकञ्चेति ॥}\pend

\pstart\versepar \ledsidenote{\passref{1.38.53}}\skt{उत्पलादिरयन्दाह }\edtext{\skt{रक्तपित्तविनाशनः}}{\Afootnote{N H; \skt{पित्तरक्तविनाशनः ।}\ A1938}}\skt{ ॥ }\pend
\pstart\versepar \skt{पिपासा }\edtext{\edtext{\skt{विषहृद्रोगश्छर्दिमूर्च्छाहरो}}{\lemma{\skt{विषहृद्रोगश्}}\Afootnote{H; \skt{विषहृद्रोग}\ A1938 N}}}{\lemma{\skt{छर्दिमूर्च्छाहरो}}\Afootnote{H; \skt{च्छर्दिमूर्च्छाहरो}\ A1938; \skt{स्थर्दिमूर्च्छाहरो}\ N}}\skt{ गणः ॥}\pend

\pstart \ledsidenote{\passref{1.38.54}}\edtext{\skt{मुस्ताहरिद्रादारुहरिद्राहरीतक्यामलकबिभीतकहैमवतीवचापाठाकटुकरोहिणी}}{\Afootnote{H\textsuperscript{pc}; \skt{मुस्ताहरिद्रादारुहरिद्राहरीतक्यामलकबिभीतककुष्ठहैमवतीवचापाठाकटुरोहिणी}\ A1938; \skt{मुस्ताहरिद्रादारुहरिद्राहरीतक्यामलकविभीतकहेमवतीवचापाठाकटुरोहिणी}\ N; \skt{मुस्ताहरिद्राहरीत क्यामलकबिभीतकहैमवतीवचापाठाकटुकरोहिणी}\ H\textsuperscript{ac}}}\skt{ }\edtext{\skt{शार्ङ्गाष्टातिविषाद्रमिडी}}{\Afootnote{H; \skt{शार्ङ्गष्टातिविषाद्राविडी}\ A1938; \skt{शार्ङ्गष्वातिविषाद्रमिडी}\ N}}\skt{ }\edtext{\skt{चेति}}{\Afootnote{N H; \skt{भल्लातकानि चित्रकश्चेति ॥५४॥}\ A1938}}\skt{ ॥}\pend

\pstart\versepar \ledsidenote{\passref{1.38.55}}\skt{एष मुस्तादिको नाम्ना }\edtext{\skt{गणः}}{\Afootnote{A1938 H; \skt{गः}\ N}}\skt{ }\edtext{\skt{श्लेष्मानिलापहः}}{\Afootnote{N H; \skt{श्लेष्मनिषूदनः ।}\ A1938}}\skt{ । }\pend
\pstart\versepar \skt{}\edtext{\skt{योनिदोषहरश्चैव}}{\lemma{\skt{चैव}}\Afootnote{N H; \skt{स्तन्य}\ A1938}}\skt{ शोधनः पाचनस्तथा ॥}\pend

\pstart \ledsidenote{\passref{1.38.56}}\edtext{\skt{हरीतक्यामलकबिभीतकानि}}{\Afootnote{H; \skt{हरीतक्यामलकबिभीतकानीति त्रिफला ॥५६॥}\ A1938; \skt{हरीतक्यामलकविभीतकानि ।}\ N}}\skt{ ॥}\pend

\pstart\versepar \ledsidenote{\passref{1.38.57}}\skt{त्रिफला कफपित्तघ्नीमेहकुष्ठविनाशनी । }\pend
\pstart\versepar \skt{}\edtext{\skt{चक्षुष्यादीपनीपथ्या}}{\Afootnote{H; \skt{चक्षुष्या दीपनी चैव}\ A1938; \skt{चक्षुष्यादीपनीपष्या}\ N}}\skt{ विषमज्वरनाशनी ॥}\pend

\pstart \ledsidenote{\passref{1.38.62}}\skt{त्रपुसीसताम्ररजतसुवर्ण्णकृष्णलोहानि }\edtext{\edtext{\skt{लोहमलाश्चेति}}{\lemma{\skt{लोहमलाश्}}\Afootnote{H; \skt{लोहमल}\ A1938; \skt{लोहमला}\ N}}}{\lemma{\skt{चेति}}\Afootnote{H; \skt{श्चेति}\ A1938 N}}\skt{ ॥}\pend

\pstart\versepar \ledsidenote{\passref{1.38.63}}\edtext{\skt{गणस्त्रप्वादिरित्येष}}{\lemma{\skt{गणस्त्रप्वादिरित्य्}}\Afootnote{A1938 H; \skt{गणस्त्रव्वादिरित्य्}\ N}}\skt{ }\edtext{\skt{गरक्रिमिहरः}}{\Afootnote{A1938 H; \skt{गरः क्रिमिहरः}\ N}}\skt{ परः । }\pend
\pstart\versepar \skt{}\edtext{\skt{पिपासागुल्महृद्रोगः}}{\Afootnote{N H; \skt{पिपासाविषहृद्रोग}\ A1938}}\skt{ }\edtext{\skt{पाण्डुमेहहरस्तथा}}{\lemma{\skt{पाण्डुमेहहरस्}}\Afootnote{A1938 H; \skt{पण्डुमेहहरस्}\ N}}\skt{ ॥}\pend

\pstart \ledsidenote{\passref{1.38.64}}\edtext{\skt{लाक्षारेवतककुटजाश्वमारककट्फलहरिद्रानिम्बसप्तच्छदमालत्यस्त्रायमाणा}}{\Afootnote{\skt{लाक्षारेवतकुटजाश्वमारकट्फलहरिद्राद्वयनिम्बसप्तच्छदमालत्यस्त्रायमाणा}\ A1938; \skt{लाक्षारेवतककुटजाश्वमारककप्फलहरिद्रानिम्बसप्तच्छदमालत्यस्त्रायमाणा}\ N; \skt{लाक्षारेवतककुटजाश्चमारककट्फलहरिद्रानिम्बसप्तच्छदमालत्यस्त्रायमाणा}\ H}}\skt{ चेति ॥}\pend

\pstart\versepar \ledsidenote{\passref{1.38.65}}\skt{कषायतिक्तमधुरः }\edtext{\skt{कफपित्तविनाशनः}}{\Afootnote{N H; \skt{कफपित्तार्तिनाशनः ।}\ A1938}}\skt{ । }\pend
\pstart\versepar \skt{कुष्ठक्रिमिहरश्चैव दुष्टव्रणविशोधनः ॥}\pend

\pstart\versepar \ledsidenote{\passref{1.38.80}}\edtext{\edtext{\skt{एभिर्ल्लेपास्तथा}}{\lemma{\skt{ल्लेपास्}}\Afootnote{H; \skt{लेपा}\ A1938 N}}}{\lemma{\skt{तथा}}\Afootnote{N H; \skt{न्कषायांश्च}\ A1938}}\skt{ तैलान्सर्प्पींष्यपिचपानकान् । }\pend
\pstart\versepar \skt{रसक्रियां }\edtext{\skt{कषायाश्च}}{\Afootnote{H; \skt{यथा}\ A1938; \skt{कषायाश्च}\ N}}\skt{ }\edtext{\skt{भिषक्कुर्व्वीत}}{\Afootnote{N H; \skt{न्यायं कुर्वीत}\ A1938}}\skt{ }\edtext{\skt{कर्म्मसु}}{\Afootnote{N H; \skt{मतिमान्भिषक् ॥८०॥}\ A1938}}\skt{ ॥}\edtext{}{\lemma{\skt{तैलान्सर्प्पींष्यपिचपानकान् । रसक्रियां}}\Afootnote{H; \skt{तैलं सर्पींषि पानकान् । प्रविभज्य}\ A1938; \skt{तैलां सर्पीष्यपि च पानकां । रसक्रियां}\ N}}\pend

\pstart\versepar \ledsidenote{\passref{1.38.66}}\edtext{\skt{पञ्चपञ्चकान्वक्ष्यामः}}{\Afootnote{H; \skt{पञ्च पञ्चमूलान्यत ऊर्ध्वं वक्ष्यामः ।}\ A1938; \skt{पञ्चपञ्चकाम्वक्ष्यामः ॥}\ N}}\skt{ ॥ }\pend
\pstart\versepar \skt{}\edtext{\skt{तत्रैरण्डो द्वे}}{\Afootnote{H; \skt{तत्र त्रिकण्टक}\ A1938; \skt{तत्रेरण्डौ द्वौ}\ N}}\skt{ }\edtext{\skt{बृहत्यौ}}{\Afootnote{N H; \skt{बृहतीद्वय}\ A1938}}\skt{ }\edtext{\skt{पृथक्यर्ण्णाविदारीगन्धा}}{\Afootnote{H; \skt{पृथक्पर्ण्यो विदारिगन्धा}\ A1938; \skt{पृथक्पर्ण्णीविदारिगन्धा}\ N}}\skt{ चेति }\edtext{\skt{कनीयान्}}{\Afootnote{N H; \skt{कनीयः ॥६६॥}\ A1938}}\skt{ ॥}\pend

\pstart \ledsidenote{\passref{1.38.68}}\edtext{\skt{बिल्वाग्निमन्थौटुश्टुकपाटलीकाश्मर्याणि}}{\Afootnote{H; \skt{बिल्वाग्निमन्थटिण्टुकपाटलाः काश्मरी चेति}\ A1938; \skt{बिल्वाग्निमन्थौ टुण्टूकपाटलीकाश्मर्याणि}\ N}}\skt{ }\edtext{\skt{महान्}}{\Afootnote{N H; \skt{महत् ॥६८॥}\ A1938}}\skt{ ॥}\pend

\pstart \ledsidenote{\passref{1.38.72}}\edtext{\skt{विदारीशारिवारजनीगुडूचीआजशृङ्गी च}}{\Afootnote{H; \skt{विदारीसारिवारजनीगुडूच्योऽजशृङ्गी चेति}\ A1938; \skt{विदारीसारिवारजनीगुडूची अंजशृङ्गी च}\ N}}\skt{ }\edtext{\skt{बल्लीसंज्ञाः}}{\Afootnote{H; \skt{वल्लीसञ्ज्ञः ॥७२॥}\ A1938; \skt{वल्लीसज्ञाः ॥}\ N}}\skt{ ।}\pend

\pstart \ledsidenote{\passref{1.38.73}}\edtext{\skt{करमर्द्दीत्रिकण्टकसैरीयकमातुलुङ्गीगृध्रनख्यः}}{\Afootnote{H; \skt{करमर्दीत्रिकण्टकसैरीयकशतावरीगृध्रनख्य}\ A1938; \skt{करमर्द्दीतृकण्टकसैरीयकमातुलुंगीगृध्रनख्यः}\ N}}\skt{ }\edtext{\skt{कणअटकसंज्ञाः}}{\Afootnote{H; \skt{इति कण्टकसञ्ज्ञः ॥७३॥}\ A1938; \skt{कण्टकसज्ञः ॥}\ N}}\skt{ ॥}\pend

\pstart \ledsidenote{\passref{1.38.75}}\edtext{\skt{कुशकाशकाण्टेक्षुदर्ब्भास्तृणसंज्ञाः}}{\Afootnote{H; \skt{कुशकाशनलदर्भकाण्डेक्षुका इति तृणसञ्ज्ञकः ॥७५॥}\ A1938; \skt{कुशकासकाण्डेक्षुदुर्भास्तृणसंज्ञाः ॥}\ N}}\skt{ ॥}\pend

\pstart \ledsidenote{\passref{1.38.77}}\edtext{\skt{तेषाम्वातहरावाद्या}}{\lemma{\skt{तेषाम्}}\Afootnote{N H; \skt{एषां}\ A1938}}\skt{ }\edtext{\skt{वन्त्यः}}{\Afootnote{A1938 H; \skt{वड्यः}\ N}}\skt{ पित्तविनाशनः । पञ्चकौ }\edtext{\skt{श्लेष्मनावितरौ}}{\Afootnote{H; \skt{श्लेष्मशमनावितरौ}\ A1938 N}}\skt{ परिकीर्त्तितौ ॥}\pend

\pstart\versepar \ledsidenote{\passref{1.38.79}}\skt{समासेन }\edtext{\edtext{\skt{गणाप्येताः}}{\lemma{\skt{गणाप्य्}}\Afootnote{\skt{गणा ह्य्}\ A1938; \skt{गुणा ह्य्}\ N; \skt{गणा}\ H}}}{\lemma{\skt{एताः}}\Afootnote{N; \skt{एते}\ A1938; \skt{प्येताः}\ H}}\skt{ प्रोक्तास्तेषान्तु विस्तरं । }\pend
\pstart\versepar \skt{चिकित्सितेषु वक्ष्यामि ज्ञात्वा }\edtext{\skt{रोगबलाबलम्}}{\Afootnote{N H; \skt{दोषबलाबलम् ॥७९॥}\ A1938}}\pend

\pstart \ledsidenote{\passref{1.38.trailer}}\edtext{\skt{इति}}{\Afootnote{\skt{इति सुश्रुतसंहितायां सूत्रस्थाने द्रव्यसङ्ग्रहणीयो नामाष्टत्रिंशोऽध्यायः ॥३८॥}\ A1938; \skt{इति ॥ ल ड ॥ ॥ ३० ८ ॥}\ N; \skt{इति ल ड ॥ ॥ ३० ८ ॥}\ H}}\skt{ ॥ ३८ ॥}\pend

\endnumbering
\beginnumbering
\eledsection*{Adhyāya 39}
\pstart \ledsidenote{\passref{1.39.1}}\skt{अथातः }\edtext{\edtext{\skt{संशोधनसंशमनीयम्व्याख्यास्यामः}}{\lemma{\skt{संशोधनसंशमनीयम्}}\Afootnote{H; \skt{संशोधनसंशमनीयमध्यायं}\ A1938; \skt{संशोधनसंसमनीयमध्यायं व्याख्यास्यामः ॥}\ N}}}{\lemma{\skt{व्याख्यास्यामः}}\Afootnote{A1938 H; \emph{om.}\ N}}\skt{ ।}\pend

\pstart \ledsidenote{\passref{1.39.3}}\edtext{\skt{मदनकुटजजीमूतकेक्ष्वाकुधामार्ग्गवकृतवेधनलोध्रसर्षपविडङ्गपिप्पलीकरञ्जप्रपूनाडकोविदारारिष्टाश्वगन्धाविदुलबन्धुमधुकजीवन्ती}}{\Afootnote{H; \skt{मदनकुटजजीमूतकेक्ष्वाकुधामार्गवकृतवेधनसर्षपविडङ्गपिप्पलीकरञ्जप्रपुन्नाड कोविदारकर्बुदारारिष्टाश्वगन्धाविदुलबन्धुजीवकश्वेताशणपुष्पी}\ A1938; \skt{मदनकुटजजीमूतकेक्ष्वाकुधामार्गवकृतवेधनलोध्रसर्षपविडङ्गपिप्पलीकरंजप्रपुरानाटकोविदारारिष्टाश्वगन्धाविदुलबन्धुमधुकजीवन्ती}\ N}}\skt{ }\edtext{\skt{बिम्बीफलामृगेर्व्वारु}}{\Afootnote{H; \skt{बिम्बीवचामृगेर्वार}\ A1938; \skt{बिम्बीवलामृगेर्वारु}\ N}}\skt{ }\edtext{\skt{चित्रस्वेता}}{\Afootnote{H; \skt{वश्चित्रा}\ A1938; \skt{चित्राश्चेता}\ N}}\skt{ }\edtext{\skt{शलपुष्पीवचाचेत्यूर्द्ध्वभागहराणि}}{\Afootnote{H; \skt{चेत्यूर्ध्वभागहराणि ।}\ A1938; \skt{सणपुष्पिवचाचेत्युर्द्धभागहराणि ॥}\ N}}\skt{ ॥ तत्र कोविदारपूर्व्वाणाम्फलानि । कोविदारादीनाम्मूलानि ॥}\pend

\pstart \ledsidenote{\passref{1.39.4}}\edtext{\skt{तृवृच्छ्यामादन्तीद्रवन्तीसफलाविषाणिगवाक्षीच्छगलान्त्रीपीलुस्नेहासुवर्ण्णक्षीरीचित्रककिणिहीकुशकासतिल्वकम्पिल्यकरम्पकपाटला}}{\Afootnote{H; \skt{विवृताश्यामादन्तीद्रवन्तीसप्तलाशङ्खिनीविषाणिकागवाक्षीच्छगलान्त्रीस्नुक्सु वर्णक्षीरीचित्रककिणिहीकुशकाशतिल्वककम्पिल्लकरम्यकपाटला}\ A1938; \skt{तृवृच्छ्यामादन्तीद्रवन्तीसप्तलाविषाणि गवाक्षीच्छगलान्त्रीपीलुस्नुहासुवर्ण्णक्षीरीचित्रककिणिहीकुशकासतिल्वकम्पिल्यकरम्यकपाटला}\ N}}\skt{ }\edtext{\skt{पूगहरीतक्यामलकविभीतकनीलिनीपञ्चङ्गुलपूतीकारग्वधमहावृक्षसप्तच्छदार्क्करेणु}}{\Afootnote{H; \skt{पूगहरीतक्या मलकबिभीतकनीलिनीचतुरङ्गुलैरण्डपूतीकमहावृक्षसप्तच्छदार्का}\ A1938; \skt{पूगहरीतक्यामलकबिभीतकनीलिनीपञ्चाङ्गुलपूतीकारग्वधमहावृक्षसप्तच्छदार्केरण्ड}\ N}}\skt{ }\edtext{\skt{ज्योतिष्मती}}{\Afootnote{A1938 H; \skt{ज्योतिश्मती}\ N}}\skt{ चेत्यधोभागहराणि ॥तत्र }\edtext{\skt{तिल्वकपूर्व्वाणां}}{\Afootnote{A1938 H; \skt{बिल्वकषूर्र्वानाम्}\ N}}\skt{ मूलानि । }\edtext{\edtext{\skt{तिल्वकादीनाम्पाटल्यन्तानाञ्च}}{\lemma{\skt{पाटल्यन्तानाञ्}}\Afootnote{N H; \skt{पाटलान्तानां}\ A1938}}}{\lemma{\skt{च}}\Afootnote{N H; \skt{त्वचः, कम्पिल्लकफल}\ A1938}}\skt{ }\edtext{\skt{त्वचः}}{\Afootnote{N H; \skt{रजः,}\ A1938}}\skt{ । }\edtext{\edtext{\skt{पूगादीनामेरण्डान्तानाम्फलानि}}{\lemma{\skt{पूगादीनाम्}}\Afootnote{N H; \skt{पूगादीना}\ A1938}}}{\lemma{\skt{एरण्डान्तानाम्}}\Afootnote{N H; \skt{मेरण्डान्तानां}\ A1938}}\skt{ । पूतीकारग्वधयोः पत्राणि । शेषाणां क्षीराणीति ॥}\pend

\pstart \ledsidenote{\passref{1.39.5}}\skt{कोशातकी सप्तला }\edtext{\skt{देवदाली}}{\Afootnote{H; \skt{शङ्खिनी देवदाली}\ A1938; \skt{देवताली}\ N}}\skt{ }\edtext{\skt{करवल्लिका}}{\Afootnote{H; \skt{कारवेल्लिका}\ A1938 N}}\skt{ }\edtext{\skt{चेत्युभयतो}}{\lemma{\skt{उभयतो}}\Afootnote{N H; \skt{उभयतो}\ A1938}}\skt{ भागहराणि । }\edtext{\skt{एषां}}{\Afootnote{A1938 H; \skt{एषा}\ N}}\skt{ }\edtext{\skt{स्वरसाः}}{\Afootnote{N H; \skt{स्वरसा इति ॥५॥}\ A1938}}\skt{ ॥}\pend

\pstart \ledsidenote{\passref{1.39.6}}\edtext{\skt{पिप्पलीविडङ्गापामार्ग्गशिग्रुशिरीषसिद्धार्थकमरिचकरवीरबिम्बीगिरिकर्ण्णिकाकिणिहीवचाज्योतिष्मतीकरञ्जार्क्कलशुनातिविषाशृ ङ्गवेरतालीसतमालसुरसार्ज्जकेङ्गुदीमेषशृङ्गीमातुलुङ्गीमुरुङ्गीपीलुजातीसालतालमधूकलाक्षाहिंगुलवणमद्यगोशकृद्रसोमूत्राणि}}{\Afootnote{H; \skt{पिप्पलीविडङ्गापामार्गशिग्रुसिद्धार्थकशिरीषमरिचकरवीरबिम्बीगिरिकर्णिकाकिणिहीवचाज्योतिष्मती- करञ्जार्कालर्कलशुनातिविषाशृङ्गवेरतालीशतमालसुरसार्जकेङ्गुदीमेषशृङ्गीमातुलुङ्गीमुरङ्गीपीलु- जातीशालतालमधूकलाक्षाहिङ्गुलवणमद्यगोशकृद्रसमूत्राणीति}\ A1938; \skt{पिप्पलीविडङ्गपामार्गशिग्रूशिरीषसिद्धार्थकमरिचकरवीरबिम्बिगिरिकर्ण्णिकाकिणिहीवचाज्योतिष्मतीकरञ्जार्कलसुनातिविषाशृङ्गवेरतालीशतमालसुरसार्जकेंगुदीमेषशृङ्गीमातुलुं।गीमुरुंगीपीलुजातीसालतालमधूकलाक्षाहिंगुलवणमद्यगोसकृद्रसोमूत्राणि}\ N}}\skt{ शिरोविरेचनानि । तत्र }\edtext{\skt{करवीरपूर्व्वाणाम्फलानि}}{\lemma{\skt{फलानि}}\Afootnote{H; \skt{फलानि, करवीरादीनामर्कान्तानां मूलानि,}\ A1938; \skt{रसानि ।}\ N}}\skt{ । }\edtext{\edtext{\skt{तालीसपूर्व्वाणाङ्कन्दाः}}{\lemma{\skt{तालीसपूर्व्वाणाङ्}}\Afootnote{N H; \skt{तालीशपूर्वाणां}\ A1938}}}{\lemma{\skt{कन्दाः}}\Afootnote{A1938 H; \skt{स्कन्धाः ॥}\ N}}\skt{ । }\edtext{\skt{तालीसकादीनामर्ज्जकान्तानाम्पत्राणि}}{\lemma{\skt{तालीसकादीनामर्ज्जकान्तानाम्}}\Afootnote{H; \skt{तालीशादीनामर्जकान्तानां}\ A1938; \skt{तालीसकादीनामर्जकान्ता ना}\ N}}\skt{ । }\edtext{\edtext{\skt{इङ्गुदीमेषशृङ्ग्योस्त्वक्}}{\lemma{\skt{इङ्गुदीमेषशृङ्ग्योस्}}\Afootnote{H; \skt{इङ्गुदीमेषशृङ्ग्यो}\ A1938; \skt{इंगुदीमेषस्र्ंग्योस्}\ N}}\skt{ । }\edtext{\skt{मातुलुङ्गीमुरुङ्गीपीलुजातीनाम्पुष्पाणि}}{\lemma{\skt{पुष्पाणि}}\Afootnote{A1938 H; \skt{पूस्पाणि ।}\ N}}}{\lemma{\skt{त्वक् । मातुलुङ्गीमुरुङ्गीपीलुजातीनाम्}}\Afootnote{H; \skt{स्त्वचः, मातुलुङ्गीमुरङ्गीपीलुजातीनां}\ A1938; \skt{त्वक् । मातुलुंगीमुरुगीपीलुजातीनां}\ N}}\skt{ । शालतालमधूकानां साराः । }\edtext{\skt{लाक्षा}}{\Afootnote{N H; \emph{om.}\ A1938}}\skt{ }\edtext{\skt{हिंगु च}}{\Afootnote{N H; \skt{हिङ्गु लाक्षे}\ A1938}}\skt{ निर्यासौ । लवणानि पार्थिवविशेषाः । }\edtext{\edtext{\skt{मद्यमासुतसंयोगाः}}{\lemma{\skt{मद्यम्}}\Afootnote{H; \skt{मद्या}\ A1938; \skt{मद्य}\ N}}}{\lemma{\skt{आसुतसंयोगाः}}\Afootnote{H; \skt{न्यासुतसंयोगाः,}\ A1938; \skt{मासुतसंयोगा}\ N}}\skt{ । }\edtext{\skt{शकृद्रसमूत्रे}}{\Afootnote{A1938 H; \skt{सकृद्रसमूत्रे}\ N}}\skt{ }\edtext{\skt{मलाविति}}{\Afootnote{\skt{मलाविति}\ A1938 N H}}\skt{ ॥}\pend

\pstart \ledsidenote{\passref{1.39.7}}\edtext{\skt{संशमनानि}}{\Afootnote{H; \skt{संशमना}\ A1938; \skt{संसमना}\ N}}\skt{ }\edtext{\skt{अत}}{\Afootnote{H; \skt{न्यत}\ A1938 N}}\skt{ }\edtext{\skt{उर्ध्व्वम्वक्ष्यामः}}{\lemma{\skt{उर्ध्व्वम्}}\Afootnote{H; \skt{ऊर्ध्वं}\ A1938; \skt{ऊर्द्धम्}\ N}}\skt{ ॥ तत्र }\edtext{\skt{भद्रदारुकुष्ठहरिद्रावरुणमेषशृङ्गीबलातिबलार्त्तगलकच्छुराशल्लकीकुबेराक्षीवीरतरसहचराग्निमन्थवत्सादनीश्व}}{\Afootnote{H; \skt{भद्रदारुकुष्ठहरिद्रावरुणमेषशृङ्गीबलातिबलार्तगलकच्छुराशल्लकीकुबेराक्षीवीरतरुसहचराग्निमन्थ- वत्सादन्ये}\ A1938; \skt{भद्रदारुकुष्ठहरिद्रावरुणमेषश्र्ङ्गीबलातिबलातगलकन्थूरासल्लकीकुवेराक्षीवीरतरसहचराग्निमन्थवच्छादनीस्व}\ N}}\skt{ }\edtext{\skt{दंष्ट्राश्मभेदकार्क्कालर्क्कशतावरीपुनर्न्नवासुकवसिरकाञ्चनभार्ग्गीवृश्चिकालीवदरयवकुलत्तप्रभृतीनि}}{\Afootnote{H\textsuperscript{pc}; \skt{रण्डाश्मभेदकालर्कार्कशतावरीपुनर्नवावसुकवशिरकाञ्चनकभार्गीकार्पासीवृश्चिकालीपत्तूरबदरयवकोलकुलत्थप्रभृतीनि}\ A1938; \skt{द्रष्टास्मभेदकार्कालर्क्कशतानवरीपुनर्न्णवावसुकवसिरकाञ्चनभागीवृश्चीकालीबदरयवकुलत्थप्रभृतीनि}\ N; \skt{दंष्ट्राश्मभेदर्क्कालर्क्कशतावरीपुनर्न्नवासुकवसिरकाञ्चनभार्ग्गीवृश्चिकालीवदरयवकुलत्तप्रभृतीनि}\ H\textsuperscript{ac}}}\skt{ }\edtext{\skt{द्वे}}{\Afootnote{N H; \skt{विदारिगन्धादिश्च द्वे}\ A1938}}\skt{ }\edtext{\skt{चाद्ये}}{\Afootnote{A1938 H; \skt{च्याद्ये}\ N}}\skt{ }\edtext{\skt{पञ्चमूल्यौ}}{\Afootnote{A1938 H; \skt{पञ्चमूल्यो}\ N}}\skt{ समासेन }\edtext{\skt{वातसं}}{\Afootnote{A1938 H; \skt{वातस}\ N}}\skt{ }\edtext{\skt{शमनो}}{\Afootnote{A1938 H; \skt{समनो}\ N}}\skt{ वर्ग्गः ॥}\pend

\pstart \ledsidenote{\passref{1.39.8}}\edtext{\skt{चन्दनह्रीवेरोशीरमञ्जिष्ठापयस्यागुन्द्रशेवालकह्लारकोत्पलमूर्व्वाप्रभृतीनि}}{\Afootnote{H; \skt{चन्दनकुचन्दनह्रीबेरोशीरमञ्जिष्ठापयस्याविदारीशतावरीगुन्द्राशैवलकह्लारकुमुदोत्पलकन्द(द)लीदूर्वामूर्वाप्रभृतीनि}\ A1938; \skt{चन्दनह्रीवेरोशीरमञ्जिष्ठापयस्यागुन्द्रासेवालकह्लारकोत्पलमूर्व्वाप्रभृतीनि}\ N}}\skt{ }\edtext{\skt{न्यग्रोधादिरुत्पलादिरिति}}{\Afootnote{N H; \skt{काकोल्यादिः सारिवादिरञ्जनादिरुत्पलादिर्न्यग्रोधादिस्तृणपञ्चमूलमिति}\ A1938}}\skt{ समासेन }\edtext{\skt{पित्तसंशमनो}}{\Afootnote{A1938 H; \skt{पित्तसंसमनो}\ N}}\skt{ वर्गः ॥}\pend

\pstart \ledsidenote{\passref{1.39.9}}\edtext{\skt{कालेयकागुरुतैलपर्ण्णिकाकुष्ठहरिद्राशीतशिवाशतपुष्पासरलारास्नाप्रकीर्योदकीचेङ्गुदीसुमनाः}}{\Afootnote{H; \skt{कालेयकागुरुतिलपर्णीकुष्ठहरिद्राशीतशिवशतपुष्पासरलारास्नाप्रकीर्योदकीर्येङ्गुदीसुमना-}\ A1938; \skt{कालेयकागुरुतैलपर्ण्णिकाकुष्ठहरिद्रासीतिसिवाशतपुष्पासरलारात्स्नाप्रकीर्योदकीचेङ्गुदीसुमना}\ N}}\skt{ }\edtext{\skt{काकादनीलाङ्गलकीहस्तिपर्ण्णमुञ्जाललामज्ज}}{\Afootnote{H; \skt{काकादनीलाङ्गलकीहस्तिकर्णमुञ्जातकलामज्ज}\ A1938; \skt{काकादनीलाङ्गलकीहस्तिपर्ण्णमुञ्जालामज्जा}\ N}}\skt{ }\edtext{\skt{कप्रभृतीनिवल्लीकण्टकपञ्चमूल्यौ द्वे}}{\Afootnote{H; \skt{कप्रभृतीनि वल्लीकण्टकपञ्चमूल्यौ}\ A1938; \skt{कप्रभृतीनि वल्लीकण्ठकपञ्चमूल्योद्वे}\ N}}\skt{ }\edtext{\edtext{\edtext{\skt{पिप्पल्यादिर्म्मुष्ककादिर्व्वचादिः}}{\lemma{\skt{पिप्पल्यादिर्}}\Afootnote{N H; \skt{पिप्पल्यादिर्बृहत्यादि}\ A1938}}}{\lemma{\skt{म्मुष्ककादिर्}}\Afootnote{N H; \skt{र्मुष्ककादि}\ A1938}}}{\lemma{\skt{व्वचादिः}}\Afootnote{H; \skt{र्वचादिः}\ A1938; \skt{वचादि}\ N}}\skt{ }\edtext{\edtext{\edtext{\skt{सुरसादिरारग्वधादिरुत्पलादिरिति}}{\lemma{\skt{सुरसादिरारग्वधादिर्}}\Afootnote{H; \skt{सुरसादिरारग्वधादि}\ A1938; \skt{सुरसादिरारग्वधा}\ N}}}{\lemma{\skt{उत्पलादिर्}}\Afootnote{H; \emph{om.}\ A1938; \skt{दिर्}\ N}}}{\lemma{\skt{इति}}\Afootnote{N H; \skt{रिति}\ A1938}}\skt{ समासेन }\edtext{\skt{श्लेष्मसंशमनो}}{\Afootnote{A1938 H; \skt{श्लेष्मसंसमनो}\ N}}\skt{ वर्ग्गः ॥}\pend

\pstart \ledsidenote{\passref{1.39.10}}\skt{तत्र }\edtext{\edtext{\skt{सर्व्वाण्येवौषधानि}}{\lemma{\skt{सर्व्वाण्य्}}\Afootnote{H; \skt{सर्वा}\ A1938 N}}}{\lemma{\skt{एवौषधानि}}\Afootnote{H; \skt{ण्येवौषधानि}\ A1938; \skt{ण्येवौषध्यनि}\ N}}\skt{ }\edtext{\skt{व्याधग्निपुरुषबलान्येवेक्ष्य}}{\Afootnote{\skt{व्याध्यग्निपुरुषबलान्यभिसमीक्ष्य}\ A1938; \skt{व्याध्यग्निपुरुषबलान्यवेक्ष्य}\ N; \skt{व्याधग्निपुरुषबलान्येवेक्ष्य}\ H}}\skt{ }\edtext{\edtext{\edtext{\edtext{\edtext{\skt{विदध्याद्व्याध्यबलादधिकमौषधमुपयुक्तन्तमुपसाम्यव्याधिमन्यमावहति}}{\lemma{\skt{विदध्याद्व्याध्यबलादधिकम्}}\Afootnote{H; \skt{विदध्यात् । तत्र, व्याधिबलादधिक}\ A1938; \skt{विदध्याद्व्याध्यबलादधिक}\ N}}}{\lemma{\skt{औषधम्}}\Afootnote{H; \skt{मौषधम्}\ A1938 N}}}{\lemma{\skt{उपयुक्तन्तम्}}\Afootnote{A1938 H; \skt{उपयुक्तं त}\ N}}}{\lemma{\skt{उपसाम्यव्याधिम्}}\Afootnote{H; \skt{उपशम्य व्याधिं}\ A1938; \skt{मुसाम्य व्याधिम्}\ N}}}{\lemma{\skt{अन्यमावहति}}\Afootnote{N; \skt{व्याधिमन्यमावहति;}\ A1938; \skt{अन्यमावहति ॥}\ H}}\skt{ ॥ }\edtext{\edtext{\skt{अग्न्यबलादधिकमजीर्ण्णम्विष्टभ्य}}{\lemma{\skt{अग्न्यबलादधिकम्}}\Afootnote{H; \skt{अग्निबलादधिक}\ A1938; \skt{अग्न्यबलादधिकम्}\ N}}}{\lemma{\skt{अजीर्ण्णम्}}\Afootnote{N H; \skt{मजीर्णं}\ A1938}}\skt{ वा पच्यते । }\edtext{\skt{पुरुषबलादधिकं}}{\Afootnote{H; \skt{पुरुषबलादधिकं}\ A1938; \skt{पुरुषाबलादधिकं}\ N}}\skt{ }\edtext{\skt{ग्लानिमूर्च्छना}}{\Afootnote{H; \skt{ग्लानिमूर्च्छामदाना}\ A1938; \skt{ग्लानिमूर्च्छाना}\ N}}\skt{ वहति । }\edtext{\skt{शमनमेवं}}{\lemma{\skt{शमनम्}}\Afootnote{N H; \skt{संशमनम्;}\ A1938}}\skt{ }\edtext{\edtext{\skt{संशोधनमतिपातयति}}{\lemma{\skt{संशोधनम्}}\Afootnote{N H; \skt{संशोधन}\ A1938}}}{\lemma{\skt{अतिपातयति}}\Afootnote{N H; \skt{मतिपातयति ।}\ A1938}}\skt{ । }\edtext{\edtext{\skt{हीनमेभ्यो}}{\lemma{\skt{हीनम्}}\Afootnote{H; \skt{हीन}\ A1938 N}}}{\lemma{\skt{एभ्यो}}\Afootnote{H; \skt{मेभ्यो}\ A1938 N}}\skt{ }\edtext{\edtext{\skt{दत्तमकिञ्चित्करम्भवति}}{\lemma{\skt{दत्तम्}}\Afootnote{H; \skt{दत्त}\ A1938; \skt{दन्तम्}\ N}}}{\lemma{\skt{अकिञ्चित्करम्}}\Afootnote{N H; \skt{मकिञ्चित्करं}\ A1938}}\skt{ । }\edtext{\edtext{\skt{तस्मात्सममेव}}{\lemma{\skt{तस्मात्समम्}}\Afootnote{A1938 H; \skt{तस्मात्सम}\ N}}}{\lemma{\skt{एव}}\Afootnote{A1938 H; \skt{येव}\ N}}\skt{ }\edtext{\skt{विध्यात्}}{\Afootnote{N H; \skt{विदध्यात् ॥१०॥}\ A1938}}\skt{ ॥}\pend

\pstart\versepar \ledsidenote{\passref{1.39.11}}\edtext{\skt{भवति}}{\Afootnote{\skt{भवन्ति}\ A1938; \emph{om.}\ N H}}\skt{ }\edtext{\skt{चात्र}}{\Afootnote{A1938; \emph{om.}\ N H}}\skt{ ॥ }\pend
\pstart\versepar \skt{रोगे }\edtext{\skt{शोधनसाध्ये तु }\edtext{\skt{यम्विद्याद्दोषमुल्वनं}}{\lemma{\skt{विद्याद्दोषम्}}\Afootnote{H; \skt{भवेद्दोष}\ A1938; \skt{विंद्याद्दोष}\ N}}}{\lemma{\skt{शोधनसाध्ये तु यम्}}\Afootnote{N H; \skt{शोधनसाध्ये तु यो}\ A1938}}\skt{ }\pend
\pstart\versepar \skt{तं }\edtext{\skt{समीक्ष्य}}{\Afootnote{N H; \skt{स्मै दद्या}\ A1938}}\skt{ }\edtext{\skt{भिषक्कुर्याद्दोषपच्यावनं}}{\Afootnote{N H; \skt{द्भिषक्प्राज्ञो दोषप्रच्यावानं}\ A1938}}\skt{ }\edtext{\skt{मृदुः}}{\Afootnote{N H; \skt{मृदु ॥११॥}\ A1938}}\skt{ ॥}\edtext{}{\lemma{\skt{उल्वनं तं}}\Afootnote{H; \skt{दुर्बलः । त}\ A1938; \skt{मुञ्चनं ॥ तं}\ N}}\pend

\pstart\versepar \ledsidenote{\passref{1.39.11.add1}}\edtext{\edtext{\skt{जलोद्धृतौषधपलस्तोयद्विकुडवायुतः}}{\lemma{\skt{जलोद्धृतौषधपलस्}}\Afootnote{H\textsuperscript{pc}; \skt{भवन्ति चात्र ॥ जलौद्ध्रितौषधपलस्}\ N; \emph{om.}\ H\textsuperscript{ac}}}}{\lemma{\skt{तोयद्विकुडवायुतः}}\Afootnote{H\textsuperscript{pc}; \skt{तोयद्विल्कुडवायुतः}\ N; \emph{om.}\ H\textsuperscript{ac}}}\skt{ । }\pend
\pstart\versepar \skt{}\edtext{\skt{पादावशेषितः}}{\Afootnote{H\textsuperscript{pc}; \skt{पादावशेषित}\ N; \emph{om.}\ H\textsuperscript{ac}}}\skt{ }\edtext{\skt{क्वाथः}}{\Afootnote{H\textsuperscript{pc}; \skt{त्काथः}\ N; \emph{om.}\ H\textsuperscript{ac}}}\skt{ }\edtext{\skt{पूतः}}{\Afootnote{H\textsuperscript{pc}; \skt{पूत}\ N; \emph{om.}\ H\textsuperscript{ac}}}\skt{ }\edtext{\skt{पानाय}}{\Afootnote{N H\textsuperscript{pc}; \emph{om.}\ H\textsuperscript{ac}}}\skt{ }\edtext{\skt{शस्यते}}{\Afootnote{H\textsuperscript{pc}; \skt{सस्यते ॥}\ N; \emph{om.}\ H\textsuperscript{ac}}}\skt{ ॥}\pend

\pstart\versepar \ledsidenote{\passref{1.39.11.add2}}\edtext{\skt{तद्वत्क्षीरार्द्धकुडवः}}{\Afootnote{N H\textsuperscript{pc}; \emph{om.}\ H\textsuperscript{ac}}}\skt{ }\edtext{\skt{सौषधस्त्रिगुणोदकः}}{\Afootnote{N H\textsuperscript{pc}; \emph{om.}\ H\textsuperscript{ac}}}\skt{ ॥ }\pend
\pstart\versepar \skt{}\edtext{\edtext{\skt{कल्काक्षमात्रिकश्चूर्ण्णो}}{\lemma{\skt{कल्काक्षमात्रिकश्}}\Afootnote{N; \emph{om.}\ H\textsuperscript{ac}; \skt{कल्काक्षमात्रिकञ्}\ H\textsuperscript{pc}}}}{\lemma{\skt{चूर्ण्णो}}\Afootnote{H\textsuperscript{pc}; \skt{चूर्ण्णा}\ N; \emph{om.}\ H\textsuperscript{ac}}}\skt{ }\edtext{\skt{बिडालपदकान्वितः}}{\Afootnote{H\textsuperscript{pc}; \skt{विडालपदकान्वितः ।}\ N; \emph{om.}\ H\textsuperscript{ac}}}\skt{ ॥}\pend

\pstart\versepar \ledsidenote{\passref{1.39.11.add3}}\edtext{\skt{भवेत्पाणितलं}}{\Afootnote{N H\textsuperscript{pc}; \emph{om.}\ H\textsuperscript{ac}}}\skt{ }\edtext{\skt{लेहः}}{\Afootnote{N H\textsuperscript{pc}; \emph{om.}\ H\textsuperscript{ac}}}\skt{ }\edtext{\skt{स्वरसस्तु}}{\Afootnote{N H\textsuperscript{pc}; \emph{om.}\ H\textsuperscript{ac}}}\skt{ }\edtext{\skt{पलद्वयं}}{\Afootnote{N H\textsuperscript{pc}; \emph{om.}\ H\textsuperscript{ac}}}\skt{ । }\pend
\pstart\versepar \skt{}\edtext{\edtext{\skt{व्याध्यादीनाम्बले}}{\lemma{\skt{व्याध्यादीनाम्}}\Afootnote{N H\textsuperscript{pc}; \emph{om.}\ H\textsuperscript{ac}}}}{\lemma{\skt{बले}}\Afootnote{N H\textsuperscript{pc}; \emph{om.}\ H\textsuperscript{ac}}}\skt{ }\edtext{\skt{मध्ये}}{\Afootnote{N H\textsuperscript{pc}; \emph{om.}\ H\textsuperscript{ac}}}\skt{ }\edtext{\skt{मात्रैषा}}{\Afootnote{N H\textsuperscript{pc}; \emph{om.}\ H\textsuperscript{ac}}}\skt{ }\edtext{\skt{शमनौषधै}}{\Afootnote{H\textsuperscript{pc}; \skt{समनौषधे ।}\ N; \emph{om.}\ H\textsuperscript{ac}}}\skt{ ॥}\pend

\pstart\versepar \ledsidenote{\passref{1.39.11.add4}}\edtext{\skt{अतो}}{\Afootnote{N H\textsuperscript{pc}; \emph{om.}\ H\textsuperscript{ac}}}\skt{ }\edtext{\skt{ज्ञात्वेष्टतो}}{\Afootnote{H\textsuperscript{pc}; \skt{ज्ञाच्चेष्टतो}\ N; \emph{om.}\ H\textsuperscript{ac}}}\skt{ }\edtext{\skt{योज्यो}}{\Afootnote{H\textsuperscript{pc}; \skt{योज्या}\ N; \emph{om.}\ H\textsuperscript{ac}}}\skt{ }\edtext{\skt{हीमही}}{\Afootnote{H\textsuperscript{pc}; \skt{हीने ही}\ N; \emph{om.}\ H\textsuperscript{ac}}}\skt{ }\edtext{\skt{नाधिकेधिकाः}}{\Afootnote{N H\textsuperscript{pc}; \emph{om.}\ H\textsuperscript{ac}}}\skt{ । }\pend
\pstart\versepar \skt{}\edtext{\skt{ज्ञात्वा}}{\Afootnote{N H\textsuperscript{pc}; \emph{om.}\ H\textsuperscript{ac}}}\skt{ }\edtext{\edtext{\skt{कोष्ठौषधबलम्मात्राकल्प्या}}{\lemma{\skt{कोष्ठौषधबलम्}}\Afootnote{N H\textsuperscript{pc}; \emph{om.}\ H\textsuperscript{ac}}}}{\lemma{\skt{मात्राकल्प्या}}\Afootnote{H\textsuperscript{pc}; \skt{मात्राकल्प्ये}\ N; \emph{om.}\ H\textsuperscript{ac}}}\skt{ }\edtext{\skt{सुशोभने}}{\Afootnote{H\textsuperscript{pc}; \skt{सुशोधने ।}\ N; \emph{om.}\ H\textsuperscript{ac}}}\skt{ ॥}\pend

\pstart\versepar \ledsidenote{\passref{1.39.12}}\skt{चले दोषे मृदौ कोष्ठे }\edtext{\edtext{\skt{नेक्षन्तत्र}}{\lemma{\skt{नेक्षन्}}\Afootnote{H; \skt{नेक्षे}\ A1938; \skt{नेक्ष}\ N}}}{\lemma{\skt{तत्र}}\Afootnote{H; \skt{तात्र}\ A1938; \skt{त्तत्र}\ N}}\skt{ बलन्नृणां । }\pend
\pstart\versepar \skt{}\edtext{\skt{अव्यापद्दुर्ब्बलस्यापि}}{\Afootnote{N; \skt{अव्याधिदुर्बलस्यापि}\ A1938; \skt{अव्यापदुर्ब्बलस्यापि}\ H}}\skt{ शोधनं हि तदा }\edtext{\skt{भवेत्}}{\Afootnote{A1938 H; \skt{भवेदिति ॥}\ N}}\skt{ ॥}\pend

\pstart \ledsidenote{\passref{1.39.trailer}}\edtext{\skt{इति}}{\Afootnote{H; \skt{इति सुश्रुतसंहितायां सूत्रस्थाने संशोधनसंशमनीयो नामैकोनचत्वारिंशोऽध्यायः ॥३९॥}\ A1938; \skt{॥ ल डे३० ९ ॥}\ N}}\skt{ ॥ ३९ ॥}\pend

\endnumbering
\beginnumbering
\eledsection*{Adhyāya 40}
\pstart \ledsidenote{\passref{1.40.1}}\skt{अथातो }\edtext{\edtext{\skt{द्रव्यरसवीर्यविपाकविज्ञानीयम्व्याख्यास्यामः}}{\lemma{\skt{द्रव्यरसवीर्यविपाकविज्ञानीयम्}}\Afootnote{H; \skt{द्रव्यरसगुणवीर्यविपाकविज्ञानीयमध्यायं}\ A1938; \skt{द्रव्यरसवीर्यविपाकविज्ञानीयं व्याव्याख्यास्यामः ॥}\ N}}}{\lemma{\skt{व्याख्यास्यामः}}\Afootnote{A1938 H; \emph{om.}\ N}}\skt{ ॥}\pend

\pstart \ledsidenote{\passref{1.40.3}}\skt{केचिदाचार्या }\edtext{\skt{व्रुवते}}{\Afootnote{H; \skt{ब्रुवते-}\ A1938; \skt{व्रूवते}\ N}}\skt{ । }\edtext{\skt{द्रव्यम्प्रधानम्कस्मात्व्यवस्थितत्वात् । इह}}{\lemma{\skt{कस्मात्व्यवस्थितत्वात् । इह}}\Afootnote{A1938 H; \skt{कस्मात्व्यस्थितत्वात् । इह}\ N}}\skt{ खलु द्रव्यम्व्यवस्थितन्न रसादयः । }\edtext{\skt{यथा}}{\Afootnote{A1938 H; \skt{यत}\ N}}\skt{ }\edtext{\skt{ये}}{\Afootnote{H; \skt{आमे फले}\ A1938; \skt{किञ्चिदिशे}\ N}}\skt{ }\edtext{\skt{फले}}{\Afootnote{N H; \skt{ये}\ A1938}}\skt{ }\edtext{\skt{रसादयस्ते}}{\Afootnote{H; \skt{रसादयस्ते}\ A1938 N}}\skt{ पक्वे न }\edtext{\skt{भवन्ति}}{\Afootnote{N H; \skt{सन्ति; नि}\ A1938}}\skt{ ॥ }\edtext{\skt{नित्यत्वाच्च}}{\Afootnote{N H; \skt{त्यत्वाच्च,}\ A1938}}\skt{ नित्यं हि }\edtext{\edtext{\skt{द्रव्यमनित्यं}}{\lemma{\skt{द्रव्यम्}}\Afootnote{N; \skt{द्रव्य}\ A1938 H}}\skt{ हि}}{\lemma{\skt{अनित्यं हि}}\Afootnote{N; \skt{मनित्या}\ A1938; \skt{मनित्यं हि}\ H}}\skt{ गुणाः । यथा कल्कादि }\edtext{\edtext{\skt{प्रविभागास्तदेव}}{\lemma{\skt{प्रविभागास्}}\Afootnote{H; \skt{प्रविभागः,}\ A1938; \skt{प्रविभागा}\ N}}}{\lemma{\skt{तदेव}}\Afootnote{H; \skt{स एव}\ A1938; \skt{तेव}\ N}}\skt{ }\edtext{\skt{सम्पन्नरसगन्धं}}{\Afootnote{H; \skt{सम्पन्नरसगन्धो}\ A1938; \skt{सम्पन्नरसगन्ध}\ N}}\skt{ }\edtext{\skt{व्यापन्नरसगन्धवान्भवतीति}}{\Afootnote{H; \skt{व्यापन्नरसगन्धो वा भवति;}\ A1938; \skt{व्यापन्न रसगन्ध सम्भवतीति ।}\ N}}\skt{ । }\edtext{\skt{स्वजातिव्यवस्थानाच्च}}{\Afootnote{H; \skt{स्वजात्यवस्थानाच्च,}\ A1938; \skt{स्वजाजातिव्यवस्थानाच्च}\ N}}\skt{ }\edtext{\skt{यथा}}{\Afootnote{N H; \skt{यथा हि}\ A1938}}\skt{ पार्थिवं }\edtext{\edtext{\skt{द्रव्यमन्यभावन्न}}{\lemma{\skt{द्रव्यम्}}\Afootnote{N H; \skt{द्रव्य}\ A1938}}}{\lemma{\skt{अन्यभावन्न}}\Afootnote{N H; \skt{मन्यभावं न}\ A1938}}\skt{ गच्छति । एवं }\edtext{\skt{शोषाणि}}{\Afootnote{H; \skt{शेषाणि}\ A1938 N}}\skt{ । }\edtext{\skt{पञ्चेन्द्रियग्रहणाच्च}}{\Afootnote{H; \skt{पञ्चेन्द्रियग्रहणाश्च,}\ A1938; \skt{पञ्चैन्द्रियग्रहणाच्च ।}\ N}}\skt{ । }\edtext{\edtext{\skt{पञ्चैन्द्रियैर्गृह्यते}}{\lemma{\skt{पञ्चैन्द्रियैर्}}\Afootnote{H; \skt{पञ्चभिरिन्द्रियैर्}\ A1938; \skt{पञ्चन्द्रियै}\ N}}}{\lemma{\skt{गृह्यते}}\Afootnote{A1938 H; \skt{र्गृह्यते}\ N}}\skt{ }\edtext{\skt{द्रष्टव्यन्न}}{\Afootnote{H; \skt{द्रव्यं न}\ A1938 N}}\skt{ रसादयः । }\edtext{\skt{आश्रयत्वाच्च}}{\Afootnote{N H; \skt{आश्रयत्वाच्च,}\ A1938}}\skt{ }\edtext{\edtext{\skt{द्रव्यमाश्रितारसा}}{\lemma{\skt{द्रव्यम्}}\Afootnote{N H; \skt{द्रव्य}\ A1938}}}{\lemma{\skt{आश्रितारसा}}\Afootnote{N H; \skt{माश्रिता रसा}\ A1938}}\skt{ }\edtext{\skt{इति}}{\Afootnote{N H; \skt{दयः;}\ A1938}}\skt{ । आरम्भ }\edtext{\skt{सामर्थ्याच्च}}{\Afootnote{N H; \skt{सामर्थ्याच्च,}\ A1938}}\skt{ द्रव्याश्रित आरम्भः ॥ }\edtext{\skt{यथा}}{\Afootnote{A1938 H; \skt{यह्ता}\ N}}\skt{ }\edtext{\edtext{\edtext{\skt{विदारिगन्धादिमाहृत्यावद्य}}{\lemma{\skt{विदारिगन्धादिम्}}\Afootnote{N H; \skt{‘विदारिगन्धादि}\ A1938}}}{\lemma{\skt{आहृत्याव्}}\Afootnote{\skt{माहृत्य}\ A1938; \skt{आहृत्या}\ N H}}}{\lemma{\skt{अद्य}}\Afootnote{\skt{सङ्क्षुद्य}\ A1938; \skt{वक्तुद्य}\ N; \skt{वद्य}\ H}}\skt{ }\edtext{\skt{विपचेदित्येवमादिषु न}}{\lemma{\skt{आदिषु न}}\Afootnote{A1938 H; \skt{आदिर्}\ N}}\skt{ }\edtext{\edtext{\skt{रसादिषारम्भः}}{\lemma{\skt{रसादिष्}}\Afootnote{H; \skt{रसादि}\ A1938 N}}}{\lemma{\skt{आरम्भः}}\Afootnote{H; \skt{ष्वारम्भः;}\ A1938; \skt{ष्वारम्भः ।}\ N}}\skt{ । }\edtext{\skt{शास्त्रप्रामाण्याच्च}}{\Afootnote{N H; \skt{शास्त्रप्रामाण्याच्च,}\ A1938}}\skt{ }\edtext{\skt{शास्त्रं ह}}{\Afootnote{H; \skt{शास्त्रे हि}\ A1938; \skt{शास्त्र हा ह ।}\ N}}\skt{ । }\edtext{\skt{द्विविधन्द्रव्यं}}{\lemma{\skt{द्विविधन्}}\Afootnote{N H; \emph{om.}\ A1938}}\skt{ }\edtext{\edtext{\skt{स्थावरञ्जञ्गम्ति}}{\lemma{\skt{स्थावरञ्}}\Afootnote{N H; \skt{प्रधानमुपदेशे योगानां, यथा- ‘मातुलुङ्गाग्निमन्थौ च’ इत्यादौ न र}\ A1938}}}{\lemma{\skt{जञ्गम्ति}}\Afootnote{H; \skt{सादय उपदिश्यन्ते;}\ A1938; \skt{जङ्गम चेति ।}\ N}}\skt{ । }\edtext{\skt{क्रमापेक्षिकत्वाच्च}}{\Afootnote{N H; \skt{क्रमापेक्षितत्वाच्च}\ A1938}}\skt{ रसादीनां रसादयो हि }\edtext{\edtext{\skt{द्रव्यक्रममपेक्षन्ते}}{\lemma{\skt{द्रव्यक्रमम्}}\Afootnote{N H; \skt{द्रव्यक्रम}\ A1938}}}{\lemma{\skt{अपेक्षन्ते}}\Afootnote{N H; \skt{मपेक्षन्ते,}\ A1938}}\skt{ यथा तरुणे तरुणाः । }\edtext{\skt{संपूर्ण}}{\Afootnote{N H; \skt{सम्पूर्णे}\ A1938}}\skt{ सम्पूर्णा इति । }\edtext{\skt{एकदेशसाध्यत्वाच्च}}{\Afootnote{N H; \skt{एकदेशसाध्यत्वाच्च,}\ A1938}}\skt{ }\edtext{\edtext{\skt{द्रव्याणामेकदेशेनापि}}{\lemma{\skt{द्रव्याणाम्}}\Afootnote{H; \skt{द्रव्याणा}\ A1938; \skt{द्रव्यणां}\ N}}}{\lemma{\skt{एकदेशेनापि}}\Afootnote{N H; \skt{मेकदेशेनापि}\ A1938}}\skt{ व्याधयः साध्यन्ते । यथा महावृक्षक्षीरेणेति । }\edtext{\skt{तस्माद्द्रव्यम्प्रधानन्रसादयः}}{\lemma{\skt{प्रधानन्रसादयः}}\Afootnote{H; \skt{प्रधानं न रसादयः}\ A1938 N}}\skt{ }\edtext{\skt{कस्मात्निरवयवत्वात्द्रव्यस्य}}{\Afootnote{N H; \skt{कस्मात्? निरवयवत्वात् । द्र}\ A1938}}\skt{ । }\edtext{\skt{द्रव्यलक्षणन्तु}}{\Afootnote{N H; \skt{व्यलक्षणं तु}\ A1938}}\skt{ }\edtext{\skt{क्रियागुणसमवायः}}{\Afootnote{H; \skt{‘क्रियागुणवत्स}\ A1938; \skt{क्रियागुणसमवायः क्रियागुणसमवायः}\ N}}\skt{ क्रियाद्रव्ये विद्यते । गुणाः }\edtext{\skt{समवायकारणञ्चेति}}{\lemma{\skt{चेति}}\Afootnote{N H; \skt{इति ॥३॥}\ A1938}}\skt{ ॥}\pend

\pstart \ledsidenote{\passref{1.40.4}}\edtext{\edtext{\edtext{\skt{नेत्याहुरन्ये}}{\lemma{\skt{नेत्य्}}\Afootnote{N H; \skt{ने}\ A1938}}}{\lemma{\skt{आहुर्}}\Afootnote{N H; \skt{त्याहु}\ A1938}}}{\lemma{\skt{अन्ये}}\Afootnote{N H; \skt{रन्ये ,}\ A1938}}\skt{ }\edtext{\skt{रसास्तु}}{\Afootnote{A1938 H; \skt{रसस्तु}\ N}}\skt{ }\edtext{\skt{प्रधानन्कस्मातागमातागमो हि}}{\Afootnote{H; \skt{प्रधानं; कस्मात्? आगमात्, आगमो हि}\ A1938; \skt{प्रदानङ्कस्मादागमाच्च आगमो हि}\ N}}\skt{ }\edtext{\edtext{\edtext{\skt{शास्त्रमित्युच्यते}}{\lemma{\skt{शास्त्रम्}}\Afootnote{H; \skt{शास्त्र}\ A1938; \skt{शात्रम्}\ N}}}{\lemma{\skt{इत्य्}}\Afootnote{H; \emph{om.}\ A1938; \skt{इय्}\ N}}}{\lemma{\skt{उच्यते}}\Afootnote{N H; \skt{मुच्यते;}\ A1938}}\skt{ । }\edtext{\skt{शास्त्रे}}{\Afootnote{N H; \skt{शास्त्रे हि}\ A1938}}\skt{ रसा }\edtext{\skt{अधिकृता}}{\Afootnote{A1938 H; \skt{आधिकृ}\ N}}\skt{ }\edtext{\skt{यथा}}{\Afootnote{A1938 H; \skt{त ।}\ N}}\skt{ रसायत्त }\edtext{\skt{आहार}}{\Afootnote{A1938 H; \skt{आहाराः}\ N}}\skt{ इति । }\edtext{\skt{उपदेशाच्च}}{\Afootnote{N H; \skt{तस्मिंस्तु प्राणाः; उपदेशाच्च,}\ A1938}}\skt{ । }\edtext{\skt{उपदिश्यते हि}}{\Afootnote{H; \skt{उपदिश्यन्ते हि}\ A1938; \skt{उपदेशा}\ N}}\skt{ । }\edtext{\skt{यथा}}{\Afootnote{H; \skt{रसाः, यथा-}\ A1938; \skt{हि । यथा}\ N}}\skt{ }\edtext{\skt{मधुराम्ललवणा}}{\Afootnote{A1938 H; \skt{मधुराम्वललवणा}\ N}}\skt{ वातं शमयन्ति । }\edtext{\skt{अतो}}{\Afootnote{N H; \emph{om.}\ A1938}}\skt{ }\edtext{\skt{रसाः}}{\Afootnote{N H; \emph{om.}\ A1938}}\skt{ }\edtext{\skt{प्रधानं}}{\Afootnote{N H; \emph{om.}\ A1938}}\skt{ । }\edtext{\skt{अनुमाच्च}}{\Afootnote{H; \skt{अनुमानाच्च,}\ A1938; \skt{अनुमानाच्च ।}\ N}}\skt{ । }\edtext{\skt{रसतो ह्यनुमीयते}}{\lemma{\skt{रसतो ह्य्}}\Afootnote{N H; \skt{रसेन ह्य्}\ A1938}}\skt{ द्रव्यं । यथा }\edtext{\edtext{\skt{मधुरमिति}}{\lemma{\skt{मधुरम्}}\Afootnote{N H; \skt{मधुर}\ A1938}}}{\lemma{\skt{इति}}\Afootnote{N H; \skt{मिति;}\ A1938}}\skt{ ॥ }\edtext{\skt{ऋषिवचनाच्च}}{\Afootnote{N H; \skt{ऋषिवचनाच्च,}\ A1938}}\skt{ । }\edtext{\skt{ऋषिवचनो}}{\Afootnote{N H; \skt{ऋषिवचनं}\ A1938}}\skt{ वेदः । यथा }\edtext{\edtext{\edtext{\skt{किञ्चिदीज्यार्थम्मधुरमाहरेदिति}}{\lemma{\skt{किञ्चिदीज्यार्थम्}}\Afootnote{\skt{किञ्चिदिज्यार्थं}\ A1938; \skt{किञ्चिदीज्यार्थ}\ N; \skt{किञ्चिदीज्यार्थम्}\ H}}}{\lemma{\skt{मधुरम्}}\Afootnote{N; \skt{मधुर}\ A1938 H}}}{\lemma{\skt{आहरेदिति}}\Afootnote{\skt{माहरेदिति,}\ A1938; \skt{आहारेदिति ।}\ N; \skt{माहरेदिति ।}\ H}}\skt{ । तस्मात्रसाः }\edtext{\skt{प्रधानं}}{\Afootnote{A1938 H; \skt{प्रधान}\ N}}\skt{ । }\edtext{\skt{रसेषु तु}}{\Afootnote{N H; \skt{रसेषु}\ A1938}}\skt{ गुणसंज्ञा }\edtext{\edtext{\skt{रसलक्षणमत्रोपदेक्ष्यामः}}{\lemma{\skt{रसलक्षणम्}}\Afootnote{N H; \skt{रसलक्षण}\ A1938}}}{\lemma{\skt{अत्रोपदेक्ष्यामः}}\Afootnote{\skt{मन्यत्रोपदेक्ष्यामः ॥४॥}\ A1938; \skt{अन्यत्रोपदेक्ष्यामः ॥}\ N; \skt{अत्रो पदेक्ष्यामः ।}\ H}}\skt{ ।}\pend

\pstart \ledsidenote{\passref{1.40.5}}\skt{नेत्याहुरन्ये वीर्यम्प्रधानमिति । }\edtext{\skt{कस्मात्प्राधान्यात् । इहत्वौषधकर्म्माणि}}{\Afootnote{H; \skt{कस्मात्? तद्वशेनौषधकर्मनि}\ A1938; \skt{कस्मात्प्रधान्यात् । इहन्नौषध कर्माणि}\ N}}\skt{ }\edtext{\skt{वीर्य}}{\Afootnote{N H; \skt{ष्पत्तेः ।}\ A1938}}\skt{ }\edtext{\skt{प्राधान्यं न}}{\Afootnote{H; \skt{इहौषध}\ A1938; \skt{प्रदान्ये न}\ N}}\skt{ }\edtext{\skt{भवन्ति}}{\Afootnote{N H; \skt{कर्मा}\ A1938}}\skt{ । }\edtext{\skt{तद्यथा}}{\Afootnote{N H; \skt{ण्यू}\ A1938}}\skt{ । }\edtext{\skt{उर्द्धभागोधोभागोभयभागसंशमनं}}{\Afootnote{H; \skt{र्ध्वाधोभागोभयभागसंशोधनसंशमन}\ A1938; \skt{उर्ध्वभागो भागोभयभागसंशमनं}\ N}}\skt{ }\edtext{\skt{साङ्ग्राहिकदीपनलेखनवृंहणरसायनवाजीकरणश्च}}{\Afootnote{H; \skt{सङ्ग्राहिकाग्निदीपनपीडनलेखनबृंहण- रसायनवाजीकरणश्व}\ A1938; \skt{सग्राहिकदीपनलेखनवृंहणरसायनवाजीकरणश्व}\ N}}\skt{ }\edtext{\skt{यथुहरविलयनदारणमर्द्दनप्राणघ्नविषप्रशमनानि}}{\Afootnote{H; \skt{यथुकरविलयनदहनदारणमादनप्राणघ्नविषप्रशमनादीनि}\ A1938; \skt{यथुहरविलयनदारुणमर्द्दनप्राणघ्नविषप्रशमनादीनि}\ N}}\skt{ वीर्य }\edtext{\skt{प्राधान्येन}}{\Afootnote{H; \skt{प्राधान्या}\ A1938; \skt{प्रधान्येन}\ N}}\skt{ }\edtext{\skt{भवन्ति । }\edtext{\skt{तन्तुवीर्यं}}{\lemma{\skt{तुवीर्यं}}\Afootnote{N H; \skt{च्च वीर्यं}\ A1938}}}{\lemma{\skt{भवन्ति । तन्}}\Afootnote{N H; \skt{द्भवन्ति । त}\ A1938}}\skt{ }\edtext{\skt{द्विविधं}}{\Afootnote{N H; \skt{द्विविध}\ A1938}}\skt{ । }\edtext{\skt{उष्णं}}{\Afootnote{N H; \skt{मुष्णं}\ A1938}}\skt{ शीतं }\edtext{\skt{चाग्निषोमीयत्वाज्जगतः}}{\Afootnote{N H; \skt{च, अग्नीषोमीयत्वाज्जगतः ।}\ A1938}}\skt{ । }\edtext{\edtext{\skt{केचिदष्टविधमाहुः}}{\lemma{\skt{केचिदष्टविधम्}}\Afootnote{H; \skt{केचिदष्टविध}\ A1938; \skt{कचिदष्टविधम्}\ N}}}{\lemma{\skt{आहुः}}\Afootnote{N H; \skt{माहुः-}\ A1938}}\skt{ । }\edtext{\skt{उष्णं}}{\Afootnote{N H; \skt{शीत}\ A1938}}\skt{ }\edtext{\skt{शीतं}}{\Afootnote{N H; \skt{मुष्णं}\ A1938}}\skt{ }\edtext{\skt{स्निग्धं}}{\Afootnote{A1938 H; \skt{स्निग्ध}\ N}}\skt{ रूक्षम्वि }\edtext{\edtext{\edtext{\skt{षदम्पिच्चिलम्मृदुस्तीक्ष्णञ्चेति}}{\lemma{\skt{षदम्}}\Afootnote{N H; \skt{शदं}\ A1938}}}{\lemma{\skt{पिच्चिलम्}}\Afootnote{H; \skt{पिच्छिलं}\ A1938 N}}}{\lemma{\skt{मृदुस्}}\Afootnote{N H; \skt{मृदु}\ A1938}}\skt{ । एतानि }\edtext{\skt{खलु}}{\Afootnote{N H; \emph{om.}\ A1938}}\skt{ वीर्याणि }\edtext{\edtext{\skt{स्वभावगुणोत्कर्षाद्रसमभिभूयात्मकर्म्म}}{\lemma{\skt{स्वभावगुणोत्कर्षाद्रसम्}}\Afootnote{N H; \skt{स्वबलगुणोत्कर्षाद्रस}\ A1938}}}{\lemma{\skt{अभिभूयात्मकर्म्म}}\Afootnote{N H; \skt{मभिभूयात्मकर्म}\ A1938}}\skt{ कुर्व्वन्ति । यथा }\edtext{\skt{तावद्वृहत्यञ्च}}{\Afootnote{N H; \skt{तावन्महत्पञ्च}\ A1938}}\skt{ }\edtext{\skt{मूलङ्कषाय}}{\lemma{\skt{कषाय}}\Afootnote{N H; \skt{कषायं}\ A1938}}\skt{ }\edtext{\skt{तिक्तं}}{\Afootnote{N H; \skt{तिक्तानुरसं}\ A1938}}\skt{ वातं }\edtext{\edtext{\skt{शमयत्युष्णवीर्यत्वात्}}{\lemma{\skt{शमयत्य्}}\Afootnote{A1938 H; \skt{समयत्य्}\ N}}\skt{ । कटुका}}{\lemma{\skt{उष्णवीर्यत्वात् । कटुका}}\Afootnote{H; \skt{उष्णवीर्यत्वात् ; तथा कुलत्थः कषायः, कटुकः पलाण्डुः, स्नेहभावाच्च; मधुरश्चेक्षुरसो वातं वर्धयति, शीतवीर्यत्वात् ; कटुका}\ A1938; \skt{उष्णवीर्यात्वात् । कटुका}\ N}}\skt{ }\edtext{\skt{पिप्पलीपित्तं}}{\Afootnote{A1938 H; \skt{पिप्पलीपित्त}\ N}}\skt{ शमयति । }\edtext{\skt{मृदुशीतवीर्यत्वात् । तिक्ता}}{\Afootnote{N H; \skt{मृदुशीतवीर्यत्वात् ; अम्लमामलकं लवणं सैन्धवं च; तिक्ता}\ A1938}}\skt{ काकमाची पित्तं }\edtext{\skt{वर्द्धयत्युष्णवीर्यत्वात् । कटुकं}}{\lemma{\skt{उष्णवीर्यत्वात् । कटुकं}}\Afootnote{N H; \skt{उष्णवीर्यत्वात् , मधुरा मत्स्याश्च; कटुकं}\ A1938}}\skt{ }\edtext{\skt{मूलं}}{\Afootnote{N H; \skt{मूलकं}\ A1938}}\skt{ श्लेष्माणं वर्द्धयति }\edtext{\skt{स्निग्धवीर्यत्वात् । अम्लं}}{\Afootnote{A1938 H; \skt{स्निग्धवीर्यत्वात् । अम्वलं}\ N}}\skt{ कपित्थं }\edtext{\skt{श्लेष्माणं}}{\Afootnote{H; \skt{श्लेष्माणां}\ A1938; \skt{लेष्माणं}\ N}}\skt{ }\edtext{\skt{ति}}{\Afootnote{H; \skt{शमयति,}\ A1938; \skt{न्जपयति}\ N}}\skt{ }\edtext{\skt{रुक्षवीर्यत्वात् ॥ भ}}{\Afootnote{H; \skt{रूक्षवीर्यत्वात् , मधुरं क्षौद्रं च; तदेतन्निदर्शनमात्रमुक्तम् ॥५॥}\ A1938; \skt{रुक्षवीर्यत्वात् ॥}\ N}}\skt{ ॥}\pend

\pstart\versepar \ledsidenote{\passref{1.40.6}}\edtext{\skt{ये}}{\Afootnote{H; \skt{भवन्ति चात्र ये}\ A1938; \skt{भ ॥ ये}\ N}}\skt{ रसा वातशमना भवन्ति यदि तेषु वै । }\pend
\pstart\versepar \skt{}\edtext{\skt{रुक्षलाघवशैत्यानि}}{\lemma{\skt{रुक्षलाघवशैत्य्}}\Afootnote{H; \skt{रौक्ष्यलाघवशैत्य्}\ A1938 N}}\skt{ न ते हन्युः समीरणम् ॥}\pend

\pstart\versepar \ledsidenote{\passref{1.40.7}}\edtext{\skt{ये रसाः}}{\Afootnote{A1938 H; \skt{ये रसो}\ N}}\skt{ पित्तशमना भवन्ति यदि तेषुवै । }\pend
\pstart\versepar \skt{}\edtext{\skt{तैक्ष्ण्यौष्ण्यलघुता}}{\Afootnote{N H; \skt{तैक्ष्ण्यौष्ण्यलघुताश्}\ A1938}}\skt{ }\edtext{\skt{चैव न ते}}{\Afootnote{A1938 H; \skt{चैव नं ते}\ N}}\skt{ }\edtext{\skt{तल्कर्म्मकारिणः}}{\Afootnote{H; \skt{तत्कर्मकारिणः ॥७॥}\ A1938; \skt{तत्कार्मिकारिणः ।}\ N}}\skt{ ॥}\pend

\pstart\versepar \ledsidenote{\passref{1.40.8}}\edtext{\skt{ये रसाः}}{\Afootnote{A1938 H; \skt{ये रसा}\ N}}\skt{ }\edtext{\skt{श्लेष्मशमना}}{\Afootnote{A1938 H; \skt{श्लेष्मो शमना}\ N}}\skt{ भवन्ति }\edtext{\skt{यदि}}{\Afootnote{A1938 H; \skt{येदि}\ N}}\skt{ तेषु वै । }\pend
\pstart\versepar \skt{}\edtext{\edtext{\skt{स्नेहगौरवशैत्यानु}}{\lemma{\skt{स्नेहगौरवशैत्य्}}\Afootnote{N H; \skt{स्नेहगौरवशै}\ A1938}}\skt{ न ते}}{\lemma{\skt{आनु न ते}}\Afootnote{H; \skt{त्यानि न ते}\ A1938; \skt{आनि न ते}\ N}}\skt{ तत्कर्म्मकारिणः ॥}\pend

\pstart \ledsidenote{\passref{1.40.9}}\edtext{\skt{तस्माद्वीर्यं}}{\Afootnote{A1938 H; \skt{तस्मा वीर्यां}\ N}}\skt{ }\edtext{\skt{प्रधानमिति}}{\lemma{\skt{प्रधानम्}}\Afootnote{A1938 H; \skt{प्रधान्यम्}\ N}}\skt{ ॥}\pend

\pstart \ledsidenote{\passref{1.40.10}}\skt{नेत्याहुरन्ये विपाकः प्रधानमिति । कस्मात्सम्यङ्मिथ्या }\edtext{\skt{विपक्वत्वात् । इह हि}}{\Afootnote{N H; \skt{विपक्वत्वात् ; इह स}\ A1938}}\skt{ }\edtext{\skt{द्रव्याण्यपह}}{\Afootnote{H; \skt{र्वद्रवाण्यभ्यवहृ}\ A1938; \skt{द्रव्याण्यवहृ}\ N}}\skt{ तानि सम्यङ्मिथ्या }\edtext{\skt{चाविपक्वानि}}{\Afootnote{N H; \skt{विपक्वानि}\ A1938}}\skt{ गुणन्दोषम्वा जनयंति । }\edtext{\edtext{\skt{तत्राहुरन्ये}}{\lemma{\skt{तत्राहुर्}}\Afootnote{N H; \skt{तत्राहु}\ A1938}}}{\lemma{\skt{अन्ये}}\Afootnote{N H; \skt{रन्ये-}\ A1938}}\skt{ प्रति र सम्पाक इति । }\edtext{\skt{विपाकः}}{\Afootnote{H; \emph{om.}\ A1938; \skt{विपाकं}\ N}}\skt{ }\edtext{\edtext{\skt{केचित्त्रिविधमिच्छन्ति}}{\lemma{\skt{केचित्त्रिविधम्}}\Afootnote{H; \skt{केचित्त्रिविध}\ A1938; \skt{के हिचित्त्रिविधम्}\ N}}}{\lemma{\skt{इच्छन्ति}}\Afootnote{N H; \skt{मिच्छन्ति-}\ A1938}}\skt{ । }\edtext{\edtext{\skt{मधुरमम्लम्कटुकञ्चेति}}{\lemma{\skt{मधुरम्}}\Afootnote{N H; \skt{मधुर}\ A1938}}}{\lemma{\skt{अम्लम्}}\Afootnote{H; \skt{मम्लं}\ A1938; \skt{अम्व्लं}\ N}}\skt{ । तत्तु न }\edtext{\skt{सम्यक् । भूतगुणादागमाच्चाम्लो}}{\Afootnote{\skt{सम्यक्, भूतगुणादामाच्चान्योऽम्लो}\ A1938; \skt{सम्यक् । भूतगुणादामाच्चान्योम्व्लो}\ N; \skt{सम्यक् । भूतगुणादागमाच्चाम्लो}\ H}}\skt{ विपाको नास्ति । पित्तं हि }\edtext{\edtext{\edtext{\skt{विदग्धमम्लतामुपैत्याग्नेयत्वात्}}{\lemma{\skt{विदग्धम्}}\Afootnote{N H; \skt{विदग्ध}\ A1938}}}{\lemma{\skt{अम्लताम्}}\Afootnote{H; \skt{मम्लताम्}\ A1938; \skt{अम्व्लताम्}\ N}}\skt{ । यदेवं}}{\lemma{\skt{आग्नेयत्वात् । यदेवं}}\Afootnote{N H; \skt{आग्नेयत्वात् ; यद्येवं}\ A1938}}\skt{ }\edtext{\skt{लवणोऽप्यन्यः}}{\lemma{\skt{ऽप्यन्यः}}\Afootnote{A1938 H; \skt{प्यन्यः}\ N}}\skt{ }\edtext{\skt{पाके}}{\Afootnote{H; \skt{पाको}\ A1938 N}}\skt{ भविष्यति । श्लेष्मा हि विदग्धो }\edtext{\skt{लवणतामुपैति}}{\lemma{\skt{उपैति}}\Afootnote{N H; \skt{उपै}\ A1938}}\skt{ । }\edtext{\skt{तमादसिद्धान्त}}{\Afootnote{H; \skt{तीति । म}\ A1938; \skt{तस्मादसिद्धान्त}\ N}}\skt{ }\edtext{\skt{एषः}}{\Afootnote{N H; \skt{धुरो}\ A1938}}\skt{ । }\edtext{\edtext{\skt{प्रतिरसमम्लाम्लस्यैवं}}{\lemma{\skt{प्रतिरसम्}}\Afootnote{N H; \skt{मधुरस्या}\ A1938}}}{\lemma{\skt{अम्लाम्लस्यैवं}}\Afootnote{H; \skt{म्लोऽम्लस्यैवं}\ A1938; \skt{अम्व्लाम्व्लस्यैवं}\ N}}\skt{ }\edtext{\skt{सर्व्वषामिति}}{\lemma{\skt{सर्व्वषाम्}}\Afootnote{N H; \skt{सर्वेषाम्}\ A1938}}\skt{ }\edtext{\edtext{\skt{दृष्टान्तञ्चोदाहरन्ति}}{\lemma{\skt{दृष्टान्तञ्}}\Afootnote{H; \skt{केचिदाहुः; दृष्टान्तं}\ A1938; \skt{दृष्तान्त}\ N}}}{\lemma{\skt{चोदाहरन्ति}}\Afootnote{N H; \skt{चोपदिशन्ति-}\ A1938}}\skt{ । }\edtext{\skt{यथा}}{\Afootnote{N H; \skt{यथा तावत्क्षीरमुखागतं पच्यमानं मधुरमेव स्यात्तथा}\ A1938}}\skt{ शालियवमुद्गादयः प्रकीर्णाः }\edtext{\edtext{\skt{स्वभावमुत्तरकाले}}{\lemma{\skt{स्वभावम्}}\Afootnote{N H; \skt{स्वभाव}\ A1938}}}{\lemma{\skt{उत्तरकाले}}\Afootnote{N H; \skt{मुत्तरकाले}\ A1938}}\skt{ }\edtext{\skt{पि न}}{\Afootnote{N H; \skt{ऽपि न}\ A1938}}\skt{ }\edtext{\edtext{\skt{परित्यजन्ति}}{\lemma{\skt{परित्य्}}\Afootnote{A1938 H; \emph{om.}\ N}}}{\lemma{\skt{अजन्ति}}\Afootnote{A1938 H; \skt{त्यजन्ति}\ N}}\skt{ }\edtext{\skt{तद्वतिति}}{\Afootnote{N H; \skt{तद्वदिति ।}\ A1938}}\skt{ । }\edtext{\edtext{\skt{केचित्पुनरवलवन्तो}}{\lemma{\skt{केचित्पुनर्}}\Afootnote{N H; \skt{केचिद्वदन्ति-}\ A1938}}}{\lemma{\skt{अवलवन्तो}}\Afootnote{N H; \skt{अबलवन्तो}\ A1938}}\skt{ }\edtext{\skt{वलवतां}}{\Afootnote{N H; \skt{बलवतां}\ A1938}}\skt{ }\edtext{\edtext{\skt{वशमायान्तितेवमवस्थितः}}{\lemma{\skt{वशमायान्तितेवम्}}\Afootnote{H; \skt{वशमायान्तीति । एवम}\ A1938; \skt{वशमायान्ति इतेवम्}\ N}}}{\lemma{\skt{अवस्थितः}}\Afootnote{N H; \skt{नवस्थितिः,}\ A1938}}\skt{ }\edtext{\skt{पाक}}{\Afootnote{N H; \skt{तस्मादसिद्धान्त}\ A1938}}\skt{ }\edtext{\skt{इति}}{\Afootnote{N H; \skt{एषः ।}\ A1938}}\skt{ । आगमे }\edtext{\skt{स्वाह}}{\Afootnote{H; \skt{हि}\ A1938; \skt{स्वाहा ।}\ N}}\skt{ । द्विविध एव पाको मधुरः कटुकश्च । तयोर्मधुराख्यो गुरुः कटुकाख्यो }\edtext{\edtext{\skt{लघुरिति}}{\lemma{\skt{लघुर्}}\Afootnote{N H; \skt{लघु}\ A1938}}}{\lemma{\skt{इति}}\Afootnote{N H; \skt{रिति ।}\ A1938}}\skt{ । तत्र }\edtext{\skt{पृथिव्यप्तेजोवाय्वाकाशानां}}{\Afootnote{A1938 H; \skt{पृथिव्याप्ततेजोवाय्य्वाकाशानां}\ N}}\skt{ }\edtext{\skt{द्वैविध्यम्भवति}}{\lemma{\skt{द्वैविध्यम्}}\Afootnote{A1938 H; \skt{द्वैविद्यम्}\ N}}\skt{ । }\edtext{\skt{गुणसाधर्म्यान्गुरुता}}{\Afootnote{\skt{गुणसाधर्म्याद्गुरुता}\ A1938; \skt{गुणसाधर्म्याद्गुरूता ।}\ N; \skt{गुणसाधर्म्यान्गुरुता ।}\ H}}\skt{ । लघुता च }\edtext{\skt{पृथिव्या यो}}{\Afootnote{N H; \skt{पृथिव्यापश्च}\ A1938}}\skt{ गुर्व्यः । शेषाणि लघूनि तस्माद्द्विविध एव }\edtext{\skt{पाको}}{\Afootnote{N H; \skt{पा}\ A1938}}\skt{ }\edtext{\skt{भवति}}{\Afootnote{N H; \skt{क इति ॥१०॥}\ A1938}}\skt{ ॥ }\edtext{\skt{भ}}{\Afootnote{H; \emph{om.}\ A1938 N}}\skt{ ॥}\pend

\pstart\versepar \ledsidenote{\passref{1.40.11}}\edtext{\skt{द्रव्येषु}}{\Afootnote{H; \skt{भवन्ति चात्र द्रव्येषु}\ A1938; \skt{भ ॥ द्रव्येषु}\ N}}\skt{ पच्यमानेषु }\edtext{\skt{येष्वम्बुपृथिवीगुणाः}}{\lemma{\skt{अम्बुपृथिवीगुणाः}}\Afootnote{A1938; \skt{अम्वुपृथिवीगुणाः}\ N H}}\skt{ । }\pend
\pstart\versepar \skt{}\edtext{\edtext{\edtext{\skt{निवर्त्तन्तेऽधिकास्तत्र}}{\lemma{\skt{निवर्त्तन्ते}}\Afootnote{H; \skt{निर्वर्तन्ते}\ A1938 N}}}{\lemma{\skt{ऽधिकास्}}\Afootnote{A1938 H; \skt{ऽअधिका}\ N}}}{\lemma{\skt{तत्र}}\Afootnote{A1938 H; \skt{स्त त्र}\ N}}\skt{ पाको मधुर उच्यते ॥}\pend

\pstart\versepar \ledsidenote{\passref{1.40.12}}\edtext{\skt{तेजोनिलाकाशगुणाः}}{\Afootnote{N H; \skt{तेजोऽनिलाकाशगुणाः}\ A1938}}\skt{ पच्यमानेषु येषु तु । }\pend
\pstart\versepar \skt{}\edtext{\edtext{\skt{निवर्त्तन्तेऽधिकास्तत्र}}{\lemma{\skt{निवर्त्तन्ते}}\Afootnote{H; \skt{निर्वर्तन्ते}\ A1938 N}}}{\lemma{\skt{ऽधिकास्}}\Afootnote{A1938 H; \skt{धिकास्}\ N}}\skt{ पाकः कटुक उच्यते ॥}\pend

\pstart\versepar \ledsidenote{\passref{1.40.13}}\edtext{\skt{पृथक्त्ववादिनामेषेव}}{\Afootnote{H; \skt{पृथक्त्वदर्शिनामेष}\ A1938; \skt{पृथक्त्व त्वादिनामेवं}\ N}}\skt{ वादिनाम्वादसंग्रहः । }\pend
\pstart\versepar \skt{}\edtext{\edtext{\skt{चतुर्ण्णामपि}}{\lemma{\skt{चतुर्ण्णाम्}}\Afootnote{N H; \skt{चतुर्णा}\ A1938}}}{\lemma{\skt{अपि}}\Afootnote{N H; \skt{मपि}\ A1938}}\skt{ }\edtext{\skt{सामग्र्यमिच्छन्त्यत्र}}{\lemma{\skt{इच्छन्त्यत्र}}\Afootnote{H; \skt{इच्छन्त्यत्र}\ A1938 N}}\skt{ विपश्चितः ॥}\pend

\pstart\versepar \ledsidenote{\passref{1.40.14}}\skt{तद्द्रव्यमात्मना किञ्चित्किञ्चिद्वीर्येण सेवितं । }\pend
\pstart\versepar \skt{}\edtext{\skt{किञ्चिद्रसविपाकाभ्यान्दोषं}}{\lemma{\skt{किञ्चिद्रसविपाकाभ्यान्}}\Afootnote{A1938 H; \skt{किञ्चिद्रसविपाकांभ्यां}\ N}}\skt{ हन्ति करोति वा ॥}\pend

\pstart\versepar \ledsidenote{\passref{1.40.15}}\skt{पाको नास्ति विना }\edtext{\skt{वीर्याद्वीर्यन्नास्ति}}{\lemma{\skt{वीर्याद्वीर्यन्}}\Afootnote{N H; \skt{वीर्याद्वीर्यं}\ A1938}}\skt{ विना रसात् । }\pend
\pstart\versepar \skt{रसो नास्ति विना }\edtext{\skt{द्रव्याद्द्रव्यं}}{\Afootnote{H; \skt{द्रव्याद्द्रव्यं}\ A1938; \skt{द्रव्याद्रव्य}\ N}}\skt{ श्रेष्ठतमं }\edtext{\skt{स्मृतं}}{\Afootnote{A1938 H; \skt{स्मृत ।}\ N}}\skt{ ॥}\pend

\pstart\versepar \ledsidenote{\passref{1.40.16}}\edtext{\skt{जन्मन्तु}}{\Afootnote{H; \skt{जन्म तु}\ A1938 N}}\skt{ द्रव्यरसयोरन्योन्यापेक्षिकं स्मृतं । }\pend
\pstart\versepar \skt{अन्योन्यापेक्षिकं जन्म यथा }\edtext{\skt{स्याद्देहदेहिनोः}}{\Afootnote{A1938 H; \skt{स्याद्देहदेहिनाः ॥}\ N}}\skt{ ॥}\pend

\pstart\versepar \ledsidenote{\passref{1.40.17}}\skt{वीर्यसंज्ञा गुणा }\edtext{\skt{येऽष्टौ }\edtext{\skt{तेऽपि}}{\lemma{\skt{ऽपि}}\Afootnote{A1938 H; \skt{पि}\ N}}}{\lemma{\skt{ऽष्टौ ते}}\Afootnote{A1938 H; \skt{ष्टौ ते}\ N}}\skt{ द्रव्याश्रयाः स्मृताः । }\pend
\pstart\versepar \skt{रसेषु न }\edtext{\skt{भवन्त्येते}}{\Afootnote{H; \skt{भवन्त्येते}\ A1938 N}}\skt{ }\edtext{\skt{निर्गुणाः}}{\Afootnote{H; \skt{निर्गुणास्तु गुणाः}\ A1938 N}}\skt{ स्मृताः ॥}\pend

\pstart\versepar \ledsidenote{\passref{1.40.18}}\skt{द्रव्ये }\edtext{\skt{द्रव्ये च}}{\Afootnote{N H; \skt{द्रव्याणि}\ A1938}}\skt{ यस्माद्धि विपच्यन्ते न }\edtext{\skt{षड्रसाः}}{\Afootnote{A1938 H; \skt{ड्रसाः ।}\ N}}\skt{ । }\pend
\pstart\versepar \skt{श्रेष्ठं द्रव्यमतो ज्ञेयं }\edtext{\skt{शेषा}}{\Afootnote{A1938 H; \skt{शेशा}\ N}}\skt{ }\edtext{\skt{भावास्तदाश्रयाः}}{\lemma{\skt{तदाश्रयाः}}\Afootnote{N H; \skt{तदाश्रयाः ॥१८॥}\ A1938}}\skt{ ॥}\pend

\pstart\versepar \ledsidenote{\passref{1.40.19}}\edtext{\skt{अमीमांस्यान्यचिन्त्यानि}}{\Afootnote{N H; \skt{अमीमांस्यान्यचिन्त्यानि}\ A1938}}\skt{ प्रसिद्धानि स्वभावतः । }\pend
\pstart\versepar \skt{आगमेनोपयोज्यानि भेषजानि विचक्षणैः ॥}\pend

\pstart\versepar \ledsidenote{\passref{1.40.20}}\skt{प्रत्यक्षलक्षणफलास्प्रसिद्धाश्च }\edtext{\skt{स्वभावतः}}{\Afootnote{N H; \skt{स्वभावतः ।}\ A1938}}\skt{ । }\pend
\pstart\versepar \skt{}\edtext{\edtext{\skt{ओषधीन्हेतुभिर्व्विद्वान्न}}{\lemma{\skt{ओषधीन्हेतुभिर्}}\Afootnote{H; \skt{नौषधीर्हेतुभिर्}\ A1938; \skt{ओषधीहेतुभिर्}\ N}}}{\lemma{\skt{व्विद्वान्न}}\Afootnote{N H; \skt{विद्वा}\ A1938}}\skt{ }\edtext{\skt{परीक्षेत्कथञ्चन}}{\Afootnote{N H; \skt{न्परीक्षेत कदाचन ॥२०॥}\ A1938}}\skt{ ॥}\pend

\pstart\versepar \ledsidenote{\passref{1.40.21}}\edtext{\skt{सहस्त्रेणापि}}{\Afootnote{H; \skt{सहस्रेणापि}\ A1938 N}}\skt{ }\edtext{\skt{हेतूनान्नाम्वष्ठादिर्विरेचयेत्}}{\lemma{\skt{नाम्वष्ठादिर्}}\Afootnote{H; \skt{नाम्बष्ठादिर्}\ A1938; \skt{नाम्वष्ठादि}\ N}}\skt{ । }\pend
\pstart\versepar \skt{तस्मात्तिष्ठेद्विमतिमानागमेश तु }\edtext{\skt{हेतुषु}}{\Afootnote{A1938 H; \skt{हेतुषु ॥ {\symfont ❈} ॥}\ N}}\skt{ ॥}\edtext{}{\lemma{\skt{विरेचयेत् । तस्मात्तिष्ठेद्विमतिमानागमेश तु}}\Afootnote{H; \skt{विरेचयेत् । तस्मात्तिष्ठेत्तु मतिमानागमे न तु}\ A1938; \skt{विरेचयेत् । तस्मात्तिष्ठेद्विमतिमनागमे न तु}\ N}}\pend

\pstart\versepar \ledsidenote{\passref{1.40.21.add1}}\edtext{\skt{या}}{\Afootnote{H; \skt{च्छा या}\ N}}\skt{ स्वभावविज्ञानं वा }\edtext{\skt{रशं}}{\Afootnote{H; \skt{रणं}\ N}}\skt{ }\edtext{\skt{युक्तसेनिकं}}{\Afootnote{H; \skt{क्ते सेनिकं ॥}\ N}}\skt{ । }\pend
\pstart\versepar \skt{आतुरोपक्रमम्मिश्रं भूमिज्ञानन्तथैव च ॥}\pend

\pstart\versepar \ledsidenote{\passref{1.40.21.add2}}\skt{द्रव्यसंग्रहणञ्चैव तथा }\edtext{\skt{शमनशोधनं}}{\Afootnote{H; \skt{शमनशोधन ।}\ N}}\skt{ । }\pend
\pstart\versepar \skt{रसवीर्यविपाकेन प्रोक्तमन्यद्दशैव तु ॥}\pend

\pstart \ledsidenote{\passref{1.40.trailer}}\edtext{\skt{चतुर्थोदश}}{\Afootnote{H; \skt{इति सुश्रुतसंहितायां सूत्रस्थाने द्रव्यगुणरसवीर्यविपाकविज्ञानीयो नाम चत्वारिंशत्तमोऽध्यायः ॥४०॥}\ A1938; \skt{चतुर्थो दश ॥}\ N}}\skt{ ॥ ॥ ॥}\pend

\endnumbering
\beginnumbering
\eledsection*{Adhyāya 41}
\pstart \ledsidenote{\passref{1.41.1}}\skt{अथातो }\edtext{\skt{द्रव्यविशेषविज्ञानीयम्व्याख्यास्यामः}}{\lemma{\skt{व्याख्यास्यामः}}\Afootnote{H; \skt{अध्यायं व्याख्यास्यामः ॥}\ A1938; \emph{om.}\ N}}\skt{ ॥}\pend

\pstart \ledsidenote{\passref{1.41.3}}\edtext{\skt{तत्र}}{\Afootnote{A1938 H; \emph{om.}\ N}}\skt{ पृथिव्यप्तेजोवाय्वाकाशानां }\edtext{\edtext{\skt{समुदायाद्द्रव्याभिनिर्वृत्तिर्ब्भवति}}{\lemma{\skt{समुदायाद्द्रव्याभिनिर्वृत्तिर्}}\Afootnote{H; \skt{समुदायाद्द्रव्याभिनिर्वृत्तिः,}\ A1938; \skt{समुदयाद्द्रव्याभिनिर्वृत्तिर्}\ N}}}{\lemma{\skt{ब्भवति}}\Afootnote{N H; \skt{उत्कर्षस्त्वभिव्यञ्जको भवति-}\ A1938}}\skt{ । इदं }\edtext{\edtext{\edtext{\skt{पार्थिवमिदमप्यमिदन्तेजसमिदम्वायव्यमिदमाकाश्यमिति}}{\lemma{\skt{अप्यम्}}\Afootnote{H; \skt{आप्यम्}\ A1938 N}}}{\lemma{\skt{तेजसम्}}\Afootnote{N H; \skt{तैजसम्}\ A1938}}}{\lemma{\skt{आकाश्यम्}}\Afootnote{N H; \skt{आकाशीयम्}\ A1938}}\skt{ ॥}\pend

\pstart \ledsidenote{\passref{1.41.4.1}}\skt{तत्र }\edtext{\edtext{\skt{स्थूलसान्द्रमन्दस्थिरगुरुकठिनगंधगुणबहुलमीषत्कषायप्रायसो}}{\lemma{\skt{स्थूलसान्द्रमन्दस्थिरगुरुकठिनगंधगुणबहुलम्}}\Afootnote{N H; \skt{स्थूलसान्द्रमन्दस्थिरगुरुकठिनं गन्धबहुल}\ A1938}}}{\lemma{\skt{ईषत्कषायप्रायसो}}\Afootnote{N H; \skt{मीषत्कषायं प्रायशो}\ A1938}}\skt{ मधुरमिति }\edtext{\skt{पार्थिवन्तत्स्थैर्यगौरवसंघातोपचयकरं}}{\lemma{\skt{तत्स्थैर्यगौरवसंघातोपचयकरं}}\Afootnote{N H; \skt{तत्स्थैर्यबलगौरवसङ्घातोपचयकरं}\ A1938}}\skt{ विशेषतश्चाधोगतिस्वभावमिति ॥}\pend

\pstart \ledsidenote{\passref{1.41.4.2}}\edtext{\edtext{\skt{शीतस्तिमितमन्दगुरुरससान्द्रमृदुपिच्छिलरसगुणबहुलमीषदम्लं}}{\lemma{\skt{शीतस्तिमितमन्दगुरुरससान्द्रमृदुपिच्छिलरसगुणबहुलम्}}\Afootnote{H\textsuperscript{pc}; \skt{शीतस्तिमितस्निग्धमन्दगुरुसरसान्द्रमृदुपिच्छिलं रसबहुलमीषत्कषायाम्ललवणं मधुर}\ A1938; \skt{शीतस्तिमितमन्दगुरुरससांशीतपिच्छिलावम्बुगुणभूयिष्ठौ । पृथिवी सोमगुणभूयिष्ठः । स्नेह तो।याकाशगुणभूयिष्ठं मृसान्द्रमृदुपिच्छिलरसगुणबहुलम्}\ N; \skt{शीतस्तिमितमन्दगुरुरससान्द्रमृदुपिच्छिलरसगुणरसगुणबहुलम्}\ H\textsuperscript{ac}}}}{\lemma{\skt{ईषदम्लं}}\Afootnote{H; \skt{रस}\ A1938; \skt{ईषदम्ब्म्लम्}\ N}}\skt{ }\edtext{\skt{प्रायसो}}{\Afootnote{N H; \skt{प्राय}\ A1938}}\skt{ }\edtext{\edtext{\edtext{\skt{मधुरमाप्यन्तत्स्नेहनह्लादनविष्यन्दनकरमिति}}{\lemma{\skt{मधुरम्}}\Afootnote{N H; \skt{मा}\ A1938}}}{\lemma{\skt{आप्यन्}}\Afootnote{N H; \skt{प्यं;}\ A1938}}}{\lemma{\skt{तत्स्नेहनह्लादनविष्यन्दनकरम्}}\Afootnote{N H; \skt{तत्स्नेहनह्लादनक्लेदनबन्धनविष्यन्दनकरम्}\ A1938}}\skt{ ।}\pend

\pstart \ledsidenote{\passref{1.41.4.3}}\edtext{\skt{तीक्ष्णोष्णरूक्षसूक्ष्मलघुविषदं}}{\Afootnote{N H; \skt{उष्णतीक्ष्णसूक्ष्मरूक्षखरलघुविशदं}\ A1938}}\skt{ }\edtext{\edtext{\edtext{\edtext{\skt{रूपगुणबहुलमीषदम्लप्रायम्विशेषतश्चोर्द्ध्वगतिस्वभावमिति}}{\lemma{\skt{रूपगुणबहुलम्}}\Afootnote{N H; \skt{रूपबहुल}\ A1938}}}{\lemma{\skt{ईषदम्लप्रायम्}}\Afootnote{H; \skt{मीषदम्ललवणं कटुकरसप्रायं}\ A1938; \skt{ईषदम्ब्लप्रायम्}\ N}}}{\lemma{\skt{विशेषतश्}}\Afootnote{N H; \skt{विशेषत}\ A1938}}}{\lemma{\skt{चोर्द्ध्वगतिस्वभावम्}}\Afootnote{N H; \skt{श्चोर्ध्वगतिस्वभावम्}\ A1938}}\skt{ }\edtext{\skt{तैजसं}}{\Afootnote{A1938 H; \skt{तेजसं ।}\ N}}\skt{ ॥ }\edtext{\skt{तद्दहनपचनतापनप्रकाशनप्रभावर्ण्णकरमिति}}{\lemma{\skt{तद्दहनपचनतापनप्रकाशनप्रभावर्ण्णकरम्}}\Afootnote{H; \skt{तद्दहनपचनदारणतापनप्रकाशनप्रभावर्णकरम्}\ A1938; \skt{तद्दहनपचनतापनप्रकासनप्रवर्ण्णकरम्}\ N}}\skt{ ।}\pend

\pstart \ledsidenote{\passref{1.41.4.4}}\edtext{\skt{सूक्ष्मरूक्षखरशिशिरलघुविषदं}}{\Afootnote{H; \skt{सूक्ष्मरूक्षखरशिशिरलघुविशदं}\ A1938; \skt{सूक्ष्मरूक्षखरशिशिरलघुविसदं}\ N}}\skt{ }\edtext{\edtext{\skt{स्पर्शगुणबहुलमीषत्तिक्तम्विशेषतः}}{\lemma{\skt{स्पर्शगुणबहुलम्}}\Afootnote{H; \skt{स्पर्शबहुल}\ A1938; \skt{स्पर्शगुणबहुलश्}\ N}}}{\lemma{\skt{ईषत्तिक्तम्}}\Afootnote{H; \skt{मीषत्तिक्तं}\ A1938; \skt{चेषत्तिक्तम्}\ N}}\skt{ कषायमिति }\edtext{\skt{वायव्यं}}{\Afootnote{N H; \skt{वायवीयं;}\ A1938}}\skt{ }\edtext{\skt{तद्वैषद्यं}}{\Afootnote{H; \skt{तद्वैशद्य}\ A1938; \skt{तद्वैसद्यं}\ N}}\skt{ }\edtext{\edtext{\skt{लाघवग्लपनविरूक्षणकरमिति}}{\lemma{\skt{लाघवग्लपनविरूक्षणकरम्}}\Afootnote{H; \skt{लाघवग्लपनविरूक्षणविचारणकरम्}\ A1938; \skt{लाघवग्लापनविरूक्षणमिति ॥}\ N}}}{\lemma{\skt{इति}}\Afootnote{A1938 H; \emph{om.}\ N}}\skt{ ॥}\pend

\pstart \ledsidenote{\passref{1.41.4.5}}\edtext{\skt{श्लक्ष्णसूक्ष्मव्यवायिविषदविविक्तरसं}}{\Afootnote{H; \skt{श्लक्ष्णसूक्ष्ममृदुव्यवायिविशदविविक्तमव्यक्तरसं}\ A1938; \skt{श्लक्ष्णसूक्ष्मव्यवायिविशदविविक्तरसं}\ N}}\skt{ }\edtext{\edtext{\edtext{\edtext{\skt{शब्दगुणबहुलमाकाश्यन्तस्मात्तत्मार्द्दवसौषिर्यलाघवकरमिति}}{\lemma{\skt{शब्दगुणबहुलम्}}\Afootnote{N H; \skt{शब्दबहुल}\ A1938}}}{\lemma{\skt{आकाश्यन्}}\Afootnote{N H; \skt{माकाशी}\ A1938}}}{\lemma{\skt{तस्मात्तत्मार्द्दवसौषिर्यलाघवकरम्}}\Afootnote{H; \skt{यं; तन्मार्दवशौषिर्यलाघवकर}\ A1938; \skt{तस्मात्तन्मार्दवसौषिर्यलाघवकरम्}\ N}}}{\lemma{\skt{इति}}\Afootnote{N H; \skt{मिति(५) ॥४॥}\ A1938}}\skt{ ॥}\pend

\pstart \ledsidenote{\passref{1.41.5}}\skt{अनेन निदर्शनेन }\edtext{\skt{नानौषधभूतञ्जगति}}{\lemma{\skt{नानौषधभूतञ्}}\Afootnote{N H; \skt{नानौषधीभूतं}\ A1938}}\skt{ किञ्चिद्द्रव्यमस्तीति }\edtext{\skt{कृत्वा तन्तं}}{\Afootnote{A1938 H; \skt{कृत्वा तं त}\ N}}\skt{ }\edtext{\edtext{\skt{युक्तिविशेषमर्थञ्चाभिप्रतीत्य}}{\lemma{\skt{युक्तिविशेषमर्थञ्}}\Afootnote{H; \skt{युक्तिविशेषमर्थं}\ A1938; \skt{युक्तिविसेषमर्थं}\ N}}}{\lemma{\skt{चाभिप्रतीत्य}}\Afootnote{N H; \skt{चाभिसमीक्ष्य}\ A1938}}\skt{ स्ववीर्यगुणयुक्तानि }\edtext{\skt{द्रव्याणि}}{\Afootnote{A1938 H; \skt{द्रव्यानि}\ N}}\skt{ }\edtext{\skt{कार्म्मुकानि}}{\Afootnote{N H; \skt{कार्मुकाणि}\ A1938}}\skt{ भवन्ति ॥ तानि यदा कुर्व्वन्ति स कालः । यत्कुर्वन्ति तत्कर्म्म । येन कुर्व्वन्ति }\edtext{\skt{तद्वीर्यं}}{\Afootnote{N H; \skt{तद्वीर्यं,}\ A1938}}\skt{ । }\edtext{\skt{यत्र}}{\Afootnote{A1938 H; \skt{यत}\ N}}\skt{ कुर्व्वन्ति }\edtext{\skt{तदधिकरणं}}{\Afootnote{A1938 N\textsuperscript{pc} H; \skt{तदधिकरण ः ।}\ N\textsuperscript{ac}}}\skt{ । यथा कुर्व्वन्ति स उपायः । }\edtext{\skt{यदभिनिष्पादयन्ति}}{\Afootnote{N H; \skt{यन्निष्पादयन्ति}\ A1938}}\skt{ }\edtext{\edtext{\skt{तत्फलमिति}}{\lemma{\skt{तत्फलम्}}\Afootnote{N H; \skt{तत्फल}\ A1938}}}{\lemma{\skt{इति}}\Afootnote{N H; \skt{मिति ॥५॥}\ A1938}}\skt{ ।}\pend

\pstart \ledsidenote{\passref{1.41.6}}\edtext{\skt{तदत्र}}{\Afootnote{N H; \skt{तत्र,}\ A1938}}\skt{ }\edtext{\skt{विरेचनद्रव्याण्यम्बुपृथिवीगुणभूयिष्ठानि}}{\lemma{\skt{अम्बुपृथिवीगुणभूयिष्ठानि}}\Afootnote{N H; \skt{पृथिव्यम्बुगुणभूयिष्ठानि,}\ A1938}}\skt{ । }\edtext{\skt{पृथिव्यापो हि}}{\Afootnote{N H; \skt{पृथिव्यापो}\ A1938}}\skt{ गुर्व्व्यः । }\edtext{\skt{गुरुत्वादधो}}{\Afootnote{N H; \skt{तानि गुरुत्वादधो}\ A1938}}\skt{ गच्छन्ति । }\edtext{\edtext{\skt{तस्माद्विरेचनद्रव्याण्यम्बुपृथिवीगुणभूयिष्ठानि}}{\lemma{\skt{तस्माद्विरेचनद्रव्याण्य्}}\Afootnote{N H; \skt{तस्माद्विरेचनम}\ A1938}}}{\lemma{\skt{अम्बुपृथिवीगुणभूयिष्ठानि}}\Afootnote{N H; \skt{धोगुणभूयिष्ठमनुमाना}\ A1938}}\skt{ । }\edtext{\edtext{\skt{वमनद्रव्याण्यग्निवायुगुणभूयिष्ठानि}}{\lemma{\skt{वमनद्रव्याण्य्}}\Afootnote{N H; \skt{त् ; वमनद्रव्या}\ A1938}}}{\lemma{\skt{अग्निवायुगुणभूयिष्ठानि}}\Afootnote{N H; \skt{ण्यग्निवायुगुणभूयिष्ठानि,}\ A1938}}\skt{ । }\edtext{\skt{अग्निवायू हि}}{\Afootnote{N H; \skt{अग्निवायु हि}\ A1938}}\skt{ लघू }\edtext{\edtext{\skt{लघुत्वाच्चोर्द्ध्वमुत्तिष्ठतः}}{\lemma{\skt{लघुत्वाच्चोर्द्ध्वम्}}\Afootnote{H; \skt{लघुत्वाच्च तान्यूर्ध्व}\ A1938; \skt{लघुत्वाच्चोर्द्वम्}\ N}}}{\lemma{\skt{उत्तिष्ठतः}}\Afootnote{N H; \skt{मुत्तिष्ठन्ति,}\ A1938}}\skt{ । }\edtext{\edtext{\edtext{\edtext{\skt{तस्माद्वमनप्रामाण्यादग्निवायुभूयिष्ठान्युभयगुणभूयिष्ठन्द्रव्यमुभयतो}}{\lemma{\skt{तस्माद्वमनप्रामाण्यादग्निवायुभूयिष्ठान्य्}}\Afootnote{H; \skt{तस्माद्वमनमप्यूर्ध्वगुणभूयिष्ठम्;}\ A1938; \skt{तस्माद्वचनप्रामाण्यादग्निवायुभूयिष्ठान्य्}\ N}}}{\lemma{\skt{उभयगुणभूयिष्ठन्}}\Afootnote{N H; \skt{उभयगुणभूयिष्ठ}\ A1938}}}{\lemma{\skt{द्रव्यम्}}\Afootnote{N H; \emph{om.}\ A1938}}}{\lemma{\skt{उभयतो}}\Afootnote{N H; \skt{मुभयतो}\ A1938}}\skt{ }\edtext{\skt{भागहरं}}{\Afootnote{N H; \skt{भागम्;}\ A1938}}\skt{ । आकाशगुणभूयिष्ठं संशमनं । }\edtext{\edtext{\skt{सांग्राहिकमनिलगुणभूयिष्ठं}}{\lemma{\skt{सांग्राहिकम्}}\Afootnote{H; \skt{साङ्ग्राहिक}\ A1938; \skt{संग्राहिकम्}\ N}}}{\lemma{\skt{अनिलगुणभूयिष्ठं}}\Afootnote{N H; \skt{मनिलगुणभूयिष्ठम्,}\ A1938}}\skt{ । }\edtext{\edtext{\skt{दीपनमग्निगुणभूयिष्ठं}}{\lemma{\skt{दीपनम्}}\Afootnote{N H; \skt{अनिलस्य शोषणात्मकत्वात् ; दीपन}\ A1938}}}{\lemma{\skt{अग्निगुणभूयिष्ठं}}\Afootnote{N H; \skt{मग्निगुणभूयिष्ठं,}\ A1938}}\skt{ । }\edtext{\edtext{\skt{लेखनमनिलानलगुणभूयिष्ठं}}{\lemma{\skt{लेखनम्}}\Afootnote{N H; \skt{तत्समानत्वात् ; लेखन}\ A1938}}}{\lemma{\skt{अनिलानलगुणभूयिष्ठं}}\Afootnote{N H; \skt{मनिलानलगुणभूयिष्ठं;}\ A1938}}\skt{ । }\edtext{\skt{बृंहनम्पृथिव्यम्बुगुणभूयिष्ठं}}{\lemma{\skt{बृंहनम्}}\Afootnote{H; \skt{बृंहणं}\ A1938; \skt{बृंहनि}\ N}}\skt{ । }\edtext{\edtext{\edtext{\skt{एवमौषधकर्म्माण्यनुमानात्साधयेत्}}{\lemma{\skt{एवम्}}\Afootnote{N H; \skt{एव}\ A1938}}}{\lemma{\skt{औषधकर्म्माण्य्}}\Afootnote{N H; \skt{मौषधकर्मा}\ A1938}}}{\lemma{\skt{अनुमानात्साधयेत्}}\Afootnote{N H; \skt{ण्यनुमानात्साधयेत् ॥६॥}\ A1938}}\skt{ ॥}\pend

\pstart\versepar \ledsidenote{\passref{1.41.7}}\edtext{\skt{भवन्ति}}{\Afootnote{A1938 H; \emph{om.}\ N}}\skt{ }\edtext{\skt{श्लोकाः}}{\Afootnote{N H; \emph{om.}\ A1938}}\skt{ ॥ }\pend
\pstart\versepar \skt{}\edtext{\edtext{\skt{मह्यम्ब्वग्न्यात्मकैर्द्द्रव्यैस्त्रिभिः}}{\lemma{\skt{मह्यम्ब्वग्न्यात्मकैर्}}\Afootnote{H; \skt{चात्र भूतेजोवारिजैर्}\ A1938; \skt{मह्यम्ब्वग्न्यात्मकै}\ N}}}{\lemma{\skt{त्रिभिः}}\Afootnote{N H; \skt{शमं}\ A1938}}\skt{ }\edtext{\skt{साम्यति}}{\Afootnote{N H; \skt{याति स}\ A1938}}\skt{ }\edtext{\skt{मारुतः}}{\Afootnote{N H; \skt{मीरणः ।}\ A1938}}\skt{ । }\pend
\pstart\versepar \skt{}\edtext{\skt{खभूम्यम्बुवायुजैः}}{\Afootnote{N H; \skt{भूम्यम्बुवायुजैः}\ A1938}}\skt{ }\edtext{\skt{पित्तञ्चतुर्ब्भिः}}{\lemma{\skt{चतुर्ब्भिः}}\Afootnote{N H; \emph{om.}\ A1938}}\skt{ }\edtext{\skt{संप्रशाम्यति}}{\Afootnote{N H; \skt{क्षिप्रमाप्नोति निर्वृतिम् ॥७॥}\ A1938}}\skt{ ॥}\pend

\pstart\versepar \ledsidenote{\passref{1.41.8}}\edtext{\skt{कफः}}{\Afootnote{N H; \emph{om.}\ A1938}}\skt{ }\edtext{\edtext{\skt{खतेजोऽनिलजैस्त्रिभिः}}{\lemma{\skt{खतेजोऽनिलजैस्}}\Afootnote{N H; \skt{खतेजोऽनिलजैः}\ A1938}}}{\lemma{\skt{त्रिभिः}}\Afootnote{N H; \skt{श्लेष्मा}\ A1938}}\skt{ }\edtext{\skt{शाम्यति}}{\Afootnote{N H; \skt{शममेति शरीरि}\ A1938}}\skt{ }\edtext{\skt{देहिनां}}{\Afootnote{N H; \skt{णाम् । विय}\ A1938}}\skt{ । }\pend
\pstart\versepar \skt{}\edtext{\edtext{\skt{खवायुक्ताभ्यान्द्रव्याभ्यां}}{\lemma{\skt{खवायुक्ताभ्यान्}}\Afootnote{H; \skt{त्पवनजाता}\ A1938; \skt{खवायुजाभ्यां}\ N}}}{\lemma{\skt{द्रव्याभ्यां}}\Afootnote{N H; \skt{भ्यां}\ A1938}}\skt{ वृद्धिमभ्येति मारुतः ॥}\pend

\pstart\versepar \ledsidenote{\passref{1.41.9}}\skt{आग्नेयमेव यद्द्रव्यन्तेन पित्तमुदीर्यते । }\pend
\pstart\versepar \skt{}\edtext{\edtext{\edtext{\edtext{\skt{मह्यम्बुजाभ्यान्द्रव्याभ्याङ्कफश्चाभिविवर्द्धते}}{\lemma{\skt{मह्यम्बुजाभ्यान्}}\Afootnote{N H; \skt{वसुधाजल}\ A1938}}}{\lemma{\skt{द्रव्याभ्याङ्}}\Afootnote{H; \skt{जाताभ्यां}\ A1938; \skt{द्रव्याभ्या}\ N}}}{\lemma{\skt{कफश्}}\Afootnote{N H; \skt{बला}\ A1938}}}{\lemma{\skt{चाभिविवर्द्धते}}\Afootnote{N H; \skt{सः परिवर्धते ॥९॥}\ A1938}}\skt{ ॥}\pend

\pstart\versepar \ledsidenote{\passref{1.41.10}}\edtext{\skt{एवमेव}}{\lemma{\skt{एव}}\Afootnote{H; \skt{एत}\ A1938; \skt{एवं}\ N}}\skt{ }\edtext{\skt{गुणाधिक्यं}}{\Afootnote{H; \skt{द्गुणाधिक्यं}\ A1938; \skt{गुणाधिक्य}\ N}}\skt{ द्रव्ये द्रव्ये }\edtext{\skt{व्यवस्थितं}}{\Afootnote{N H; \skt{विनिश्चितम् ।}\ A1938}}\skt{ । }\pend
\pstart\versepar \skt{द्विशो वा बहुशो वापि ज्ञात्वा }\edtext{\skt{दोषेऽवचाचरेत्}}{\lemma{\skt{ऽवचाचरेत्}}\Afootnote{H; \skt{षु चाचरेत् ॥१०॥}\ A1938; \skt{वचाचरेत् ॥}\ N}}\skt{ ॥}\pend

\pstart \ledsidenote{\passref{1.41.11}}\skt{तत्र य इमे }\edtext{\skt{गुणाष्टौवीर्यसंज्ञकाः}}{\Afootnote{N H; \skt{गुणा वीर्यसञ्ज्ञकाः}\ A1938}}\skt{ । }\edtext{\edtext{\edtext{\skt{शीतोष्णस्निग्धरूक्षमृदुतीक्ष्णपिच्छिलविषदास्तेषान्तीक्ष्णोष्णावाग्नेयौ}}{\lemma{\skt{शीतोष्णस्निग्धरूक्षमृदुतीक्ष्णपिच्छिलविषदास्}}\Afootnote{N H; \skt{शीतोष्णस्निग्धरूक्षमृदुतीक्ष्णपिच्छिलविशदास्}\ A1938}}}{\lemma{\skt{तीक्ष्णोष्णाव्}}\Afootnote{N H; \skt{तीक्ष्णोष्णा}\ A1938}}}{\lemma{\skt{आग्नेयौ}}\Afootnote{N H; \skt{वाग्नेयौ,}\ A1938}}\skt{ । }\edtext{\edtext{\skt{शीतपिच्छिलावम्बुगुणभूयिष्ठौ}}{\lemma{\skt{शीतपिच्छिलाव्}}\Afootnote{N H; \skt{शीतपिच्छिला}\ A1938}}}{\lemma{\skt{अम्बुगुणभूयिष्ठौ}}\Afootnote{N H; \skt{वम्बुगुणभूयिष्ठौ ,}\ A1938}}\skt{ । }\edtext{\skt{पृथिवीसोमगुणभूयिष्ठः}}{\Afootnote{H; \skt{पृथिव्यम्बुगुणभूयिष्ठः}\ A1938; \skt{पृथिवीसोमभूयिष्ठौः ।}\ N}}\skt{ । स्नेहस्तोयाकाशगुणभूयिष्ठं । मृदुत्वं । वायुगुणभूयिष्ठं रौक्ष्यं । }\edtext{\skt{क्षितिसमीरणभूयिष्ठं}}{\Afootnote{N H; \skt{क्षितिसमीरणगुणभूयिष्ठं}\ A1938}}\skt{ }\edtext{\skt{वैषद्यं}}{\Afootnote{H; \skt{वैशद्यं}\ A1938 N}}\skt{ । }\edtext{\edtext{\skt{विपाकावुक्तगुणौ}}{\lemma{\skt{विपाकाव्}}\Afootnote{N H; \skt{विपाका}\ A1938}}}{\lemma{\skt{उक्तगुणौ}}\Afootnote{N H; \skt{वुक्तगुणौ ।}\ A1938}}\skt{ । कर्म्माण्युष्णस्य }\edtext{\skt{दहनपाचनमूर्च्छनस्वेदनवमनविरेचनानि ॥ शीतस्य प्रह्लादनविष्कम्भनस्थिरीकरणप्रसादनक्लेदनजीवनानि । स्निग्धस्य }\edtext{\skt{स्नेहनबृंहणसन्तर्प्पणव्याजीकरणवयःस्थापनानि}}{\Afootnote{H; \skt{स्नेहनबृंहणसंतर्पणव्याजीकरणवयस्थापनानि ।}\ N}}\skt{ ॥ }\edtext{\skt{रूक्षस्यानिलवृद्धिसंग्रहणपीडनविरूक्षशोषरोहणानि}}{\Afootnote{H; \skt{रूक्षस्यानिलवृद्धिसंग्रहणपीडनविरूक्षशोपरोहणानि ॥}\ N}}\skt{ ॥ विषदस्य }\edtext{\skt{क्लेदावूषणविरूक्षशोषरोहणानि}}{\Afootnote{H; \skt{क्लेदावूषणविरूक्षशोपरोहणानि ॥}\ N}}\skt{ ॥ }\edtext{\skt{पिच्छिलस्योपलेपनपूरणबृंहणसंश्लेषणव्याजीकरणानि}}{\Afootnote{H; \skt{पिच्छिलस्योपलेपनपूरणबृंहणसंक्लेपनव्याजीकरणानि ।}\ N}}\skt{ ॥ मृदो }\edtext{\skt{रक्तमान्सप्रसादनमुख्यसंस्पर्शनानि}}{\Afootnote{H; \skt{रक्तमांसप्रसादनमुख्यसंस्पर्शणानि ॥}\ N}}\skt{ ॥ तीक्ष्णस्य }\edtext{\skt{सङ्ग्रहाचूषणावदारणश्रावणानि}}{\Afootnote{H; \skt{सङ्ग्रहाचूषणावदारणानि ॥}\ N}}\skt{ । }\edtext{\skt{तत्रोष्णस्निग्धौ}}{\Afootnote{N H; \skt{तत्र, उष्णस्निग्धौ}\ A1938}}\skt{ वातघ्नौ । शीतमृदुपिच्छिलाः पित्तघ्नाः । }\edtext{\skt{तीक्ष्णरूक्षविषदाः}}{\Afootnote{N H; \skt{तीक्ष्णरूक्षविशदाः}\ A1938}}}{\lemma{\skt{कर्म्माण्य्… त}}\Afootnote{\emph{om.}\ A1938}}\skt{ श्लेष्मघ्नाः । गुरुपाको }\edtext{\skt{वातपित्तघ्नः}}{\Afootnote{A1938 N H\textsuperscript{pc}; \skt{वातपित्तघ्न्ः ।}\ H\textsuperscript{ac}}}\skt{ । लघुपाकः }\edtext{\skt{श्लेष्मपित्तघ्न}}{\Afootnote{N H; \skt{श्लेष्म}\ A1938}}\skt{ }\edtext{\skt{इति}}{\Afootnote{N H; \skt{घ्नः ।}\ A1938}}\skt{ । }\edtext{\skt{तेषु}}{\Afootnote{N H; \skt{तेषां}\ A1938}}\skt{ मृदुशीतोष्णाः स्पर्शग्राह्याः । }\edtext{\skt{पिच्छिलविषदौ}}{\Afootnote{N\textsuperscript{pc} H; \skt{पिच्छिलविशदौ}\ A1938; \skt{पिच्छिलविदोषदौ}\ N\textsuperscript{ac}}}\skt{ }\edtext{\skt{चक्षुःस्पर्शाभ्यां}}{\Afootnote{A1938 H; \skt{चक्षुस्पर्शाभ्यां ।}\ N}}\skt{ । स्निग्धरूक्षौ चक्षुषा । तीक्ष्णः }\edtext{\skt{सुखदुःखोत्पादनात् । गुरुपाकः}}{\Afootnote{N H; \skt{मुखे दुःखोत्पादनात् । गुरुपाकः}\ A1938}}\skt{ }\edtext{\skt{सृष्टविण्मूत्रतया}}{\Afootnote{A1938 H; \skt{सृष्टविट्मूत्रतया}\ N}}\skt{ }\edtext{\skt{कफां}}{\Afootnote{H; \skt{कफो}\ A1938 N}}\skt{ }\edtext{\skt{क्लेशनञ्च}}{\Afootnote{H; \skt{त्क्लेशेन च,}\ A1938; \skt{क्लेदनञ्च ।}\ N}}\skt{ । }\edtext{\edtext{\skt{लघुर्ब्बद्धविण्मूत्रतया}}{\lemma{\skt{लघुर्}}\Afootnote{H; \skt{लघु}\ A1938 N}}}{\lemma{\skt{ब्बद्धविण्मूत्रतया}}\Afootnote{H; \skt{र्बद्धविण्मूत्रतया}\ A1938; \skt{बद्धविट्मूत्रतया}\ N}}\skt{ }\edtext{\skt{मारुतकोपनञ्च}}{\Afootnote{N H; \skt{मारुतकोपेन च ।}\ A1938}}\skt{ । तत्र }\edtext{\skt{तुल्यगुणेषु}}{\Afootnote{A1938 H; \skt{तु गुणैषु}\ N}}\skt{ भूतेषु }\edtext{\edtext{\skt{रसवैशेष्यमुपलक्षयेत्}}{\lemma{\skt{रसवैशेष्यम्}}\Afootnote{A1938 H; \skt{रसवैसेष्यम्}\ N}}\skt{ ॥ यथा}}{\lemma{\skt{उपलक्षयेत् ॥ यथा}}\Afootnote{N H; \skt{उपलक्षयेत् । तद्यथा-}\ A1938}}\skt{ मधुरो गुरुश्च पार्थिवः । }\edtext{\edtext{\skt{मधुरस्निग्धशीताश्चाप्या}}{\lemma{\skt{मधुरस्निग्धशीताश्}}\Afootnote{N H; \skt{मधुरः स्निग्धश्}\ A1938}}}{\lemma{\skt{चाप्या}}\Afootnote{N H; \skt{चाप्य}\ A1938}}\skt{ इति ॥}\pend

\pstart\versepar \ledsidenote{\passref{1.41.12}}\edtext{\skt{॥ भवति}}{\Afootnote{A1938; \emph{om.}\ N H}}\skt{ }\edtext{\skt{चात्र}}{\Afootnote{A1938; \emph{om.}\ N H}}\skt{ ॥ }\pend
\pstart\versepar \skt{}\edtext{\skt{गुणा ये}}{\Afootnote{N H; \skt{गुणा य}\ A1938}}\skt{ उक्ता द्रव्येषु शरीरेष्वपि }\edtext{\skt{तान्विदुः}}{\Afootnote{H; \skt{ते तथा ।}\ A1938; \skt{ताम्विदुः ।}\ N}}\skt{ । }\pend
\pstart\versepar \skt{}\edtext{\skt{स्थानवृद्धिक्षयात्तस्माद्देहिनां}}{\Afootnote{N H; \skt{स्थानवृद्धिक्षयास्तस्माद्देहिनां}\ A1938}}\skt{ }\edtext{\skt{द्रव्यहेतुका}}{\Afootnote{N H; \skt{द्रव्यहेतुकाः ॥१२॥}\ A1938}}\skt{ }\edtext{\skt{इति}}{\Afootnote{N H; \emph{om.}\ A1938}}\skt{ ॥}\pend

\pstart \ledsidenote{\passref{1.41.trailer}}\edtext{\skt{॥ ४१ ॥}}{\Afootnote{\skt{इति सुश्रुतसंहितायां सूत्रस्थाने द्रव्यविशेषविज्ञानीयो नामैकचत्वारिंशत्तमोऽध्यायः ॥४१॥}\ A1938; \skt{४१ ॥ {\symfont ❈} ॥}\ N; \skt{॥ ४१ ॥}\ H}}\pend

\endnumbering
\beginnumbering
\eledsection*{Adhyāya 42}
\pstart \ledsidenote{\passref{1.42.1}}\skt{अथातो }\edtext{\skt{रसविशेषवि-ज्ञानीयं}}{\Afootnote{H; \skt{रसविशेषविज्ञानीयमध्यायं}\ A1938; \skt{रसविशेषविज्ञानीयं । व्या ॥}\ N}}\skt{ }\edtext{\skt{व्याख्यास्यामः}}{\Afootnote{A1938 H; \emph{om.}\ N}}\skt{ ॥}\pend

\pstart \ledsidenote{\passref{1.42.3}}\skt{आकाशपवनदहनतोयभूमिषु यथासङ्ख्यमेकोत्तरपरिवृद्धाः ॥ श ब्दस्पर्शरूपरसगन्धाः ॥ तस्मादाप्यो रसः }\edtext{\skt{परस्परानुप्रवेशाच्च}}{\Afootnote{N H; \skt{परस्परसंसर्गात्परस्परानुग्रहात्परस्परानुप्रवेशाच्च}\ A1938}}\skt{ }\edtext{\skt{सर्व्वेषां}}{\Afootnote{N H; \skt{सर्वेषु सर्वेषां}\ A1938}}\skt{ }\edtext{\edtext{\skt{सान्निध्यमस्ति}}{\lemma{\skt{सान्निध्यम्}}\Afootnote{H; \skt{सान्निध्य}\ A1938; \skt{सानिध्यम्}\ N}}}{\lemma{\skt{अस्ति}}\Afootnote{H; \skt{मस्ति,}\ A1938; \skt{अ}\ N}}\skt{ }\edtext{\skt{उत्कर्षात्तु}}{\Afootnote{H; \skt{उत्कर्षापकर्षात्तु}\ A1938; \skt{त्युत्तर्षा}\ N}}\skt{ }\edtext{\skt{ग्रहणं । स}}{\Afootnote{A1938 H; \skt{त्सं}\ N}}\skt{ खल्वाप्यो रसः }\edtext{\skt{शेषभूतसंसर्ग्गाद्विदग्धः}}{\Afootnote{A1938 H; \skt{शेषभूतसंसर्गाद्विदग्धः}\ N}}\skt{ }\edtext{\skt{षड्विधो}}{\Afootnote{N H; \skt{षोढा वि}\ A1938}}\skt{ }\edtext{\skt{भवति}}{\Afootnote{N H; \skt{भज्यते;}\ A1938}}\skt{ । }\edtext{\skt{तद्यथा}}{\Afootnote{A1938 H; \skt{तद्था}\ N}}\skt{ । }\edtext{\skt{मधुरोऽम्लो}}{\lemma{\skt{ऽम्लो}}\Afootnote{A1938 H; \skt{म्व्लो}\ N}}\skt{ लवणः }\edtext{\skt{कटुकस्तिक्तः}}{\Afootnote{A1938 H; \skt{कटुकस्तिक्तः}\ N}}\skt{ कषाय }\edtext{\skt{इति}}{\Afootnote{N H; \skt{इति । ते च}\ A1938}}\skt{ । }\edtext{\skt{तत्रोदकवाहुल्यात्मधुरः}}{\Afootnote{N H; \skt{भूयः परस्परसंसर्गात्त्रिषष्टिधा भिद्यन्ते । तत्र, भूम्यम्बुगुणबाहुल्यान्मधुरः,}\ A1938}}\skt{ । }\edtext{\skt{तोयाग्निवाहुल्यादम्लः}}{\Afootnote{H; \skt{भूम्यग्निगुणबाहुल्यादम्लः,}\ A1938; \skt{तोयाग्निवाहुल्यादम्व्लः ।}\ N}}\skt{ }\edtext{\skt{भूम्यग्निवाहुल्याल्लवणः}}{\Afootnote{N H; \skt{तोयाग्निगुणबाहुल्याल्लवणः ,}\ A1938}}\skt{ }\edtext{\skt{वाय्वाकाशवाहुल्यात्तिक्तः}}{\Afootnote{N H; \skt{वाय्वग्निगुणबाहुल्यात्कटुकः, वाय्वाकाशगुणबाहुल्यात्तिक्तः,}\ A1938}}\skt{ । }\edtext{\skt{वाय्वग्निवाहुल्यात्कटुकः}}{\Afootnote{H; \skt{वाय्वग्निवाह्ŭल्यात्कटुकः ।}\ N}}\skt{ }\edtext{\edtext{\skt{पृथिव्यनिलवाहुल्यात्कषाय}}{\Afootnote{H; \skt{पृथिव्यनिलगुणबाहुल्यात्कषाय}\ A1938; \skt{पृथिव्यानिलावाहुल्यात्कषाय}\ N}}\skt{ इति ॥}}{\lemma{\skt{वाय्वग्निवाहुल्यात्कटुकः पृ}}\Afootnote{\emph{om.}\ A1938}}\pend

\pstart \ledsidenote{\passref{1.42.4}}\skt{तत्र }\edtext{\skt{मधुराम्ललवणा}}{\Afootnote{A1938 H; \skt{मधुराम्व्ललवणा}\ N}}\skt{ वातघ्नाः । मधुरतिक्तकषायाः पित्तघ्नाः कटुतिक्तकषायाः श्लेष्मघ्नाः ॥}\pend

\pstart \ledsidenote{\passref{1.42.5}}\skt{तत्र }\edtext{\skt{वायुरात्मैवात्मा}}{\Afootnote{H; \skt{वायोरात्मैवात्मा,}\ A1938; \skt{वायुरात्मैवात्मा ।}\ N}}\skt{ । }\edtext{\skt{पित्तमाग्नेयं}}{\Afootnote{A1938 H; \skt{पित्तमाग्नेयं ।}\ N}}\skt{ । }\edtext{\skt{श्लेष्म}}{\Afootnote{H; \skt{श्लेष्मा}\ A1938 N}}\skt{ सौम्य इति ॥}\pend

\pstart \ledsidenote{\passref{1.42.6}}\skt{त एते रसाः स्वयोनिवर्धना अन्ययोनिप्रशमनाश्च ।}\pend

\pstart \ledsidenote{\passref{1.42.7}}\edtext{\skt{केचिदाहुरग्नीषोमीयत्वाज्जगतः}}{\lemma{\skt{अग्नीषोमीयत्वाज्जगतः}}\Afootnote{H; \skt{अग्नीषोमीयत्वाज्जगतो}\ A1938 N}}\skt{ । }\edtext{\skt{द्विविधा}}{\Afootnote{H; \skt{रसा द्विवि}\ A1938; \skt{द्विधा}\ N}}\skt{ }\edtext{\skt{रसाः}}{\Afootnote{N H; \skt{धाः-}\ A1938}}\skt{ सौम्या }\edtext{\skt{आग्नेयाश्च}}{\Afootnote{A1938 H; \skt{ग्नेयाश्च ।}\ N}}\skt{ । }\edtext{\skt{तत्र}}{\Afootnote{N H; \emph{om.}\ A1938}}\skt{ मधुरतिक्तकषायाः }\edtext{\skt{सौम्याः}}{\Afootnote{H; \skt{सौम्याः}\ A1938 N}}\skt{ }\edtext{\skt{द्विविधा}}{\Afootnote{H; \emph{om.}\ A1938 N}}\skt{ }\edtext{\skt{कट्वम्ललवणा}}{\Afootnote{A1938 H; \skt{कट्वाम्व्ललवणास्}\ N}}\skt{ }\edtext{\skt{आग्नेयाः}}{\Afootnote{H; \skt{आग्नेयाः ।}\ A1938; \skt{त्वाग्नेयाः ॥}\ N}}\skt{ । }\edtext{\skt{मधुराम्ललवणाः}}{\Afootnote{H; \skt{तत्र मधुराम्ललवणाः}\ A1938; \skt{मधुराम्व्ललवणाः}\ N}}\skt{ स्निग्धा गुरवश्च }\edtext{\skt{द्विविधा}}{\Afootnote{H; \emph{om.}\ A1938 N}}\skt{ कटुतिक्तकषाया रूक्षा लघवश्च । सौम्याः शीताः । }\edtext{\edtext{\skt{आग्नेयास्तूष्णाः}}{\lemma{\skt{आग्नेयास्}}\Afootnote{N H; \skt{आग्नेया}\ A1938}}}{\lemma{\skt{तूष्णाः}}\Afootnote{N H; \skt{उष्णाः ॥७॥}\ A1938}}\skt{ ।}\pend

\pstart \ledsidenote{\passref{1.42.8.1}}\skt{तत्र }\edtext{\skt{शैत्यरूक्षलाघववैषद्यगुणलक्षणो}}{\Afootnote{H; \skt{शैत्यरौक्ष्यलाघववैशद्यवैष्टम्भ्यगुणलक्षणो}\ A1938; \skt{शैत्यरुक्षलाघववैशद्यगुणलक्षणो}\ N}}\skt{ वायुः । तस्य समानयोनिः कषायो रसः । }\edtext{\skt{सौम्यः}}{\Afootnote{H; \skt{सो}\ A1938; \skt{सौस्य}\ N}}\skt{ }\edtext{\edtext{\skt{शैत्याच्छैत्यमभिवर्धयति}}{\lemma{\skt{शैत्याच्छैत्यम्}}\Afootnote{H; \skt{ऽस्य शैत्याच्छैत्यं}\ A1938; \skt{शैत्याच्छैत्यम्}\ N}}}{\lemma{\skt{अभिवर्धयति}}\Afootnote{N H; \skt{वर्धयति,}\ A1938}}\skt{ । }\edtext{\skt{रौक्ष्याद्रौक्ष्यं}}{\Afootnote{A1938 H; \skt{रौक्ष्यारौक्ष्यं}\ N}}\skt{ लाघवाल्लाघवं }\edtext{\edtext{\skt{वैषद्याद्वैषद्यमिति}}{\lemma{\skt{वैषद्याद्वैषद्यम्}}\Afootnote{H; \skt{वैशद्याद्वैशद्यं,}\ A1938; \skt{वैशुद्याद्वैशद्यम्}\ N}}}{\lemma{\skt{इति}}\Afootnote{N H; \skt{वैष्टम्भ्याद्वैष्टम्भ्यमिति; (१) ।८।}\ A1938}}\skt{ ॥}\pend

\pstart \ledsidenote{\passref{1.42.8.2}}\edtext{\skt{औष्ण्यतैक्ष्ण्यरौक्ष्यलाघववैषद्यगुणलक्षणम्पित्तन्तस्य}}{\Afootnote{H; \skt{औष्ण्यतैक्ष्ण्यरौक्ष्यलाघववैशद्यगुणलक्षणं पित्तं, तस्य}\ A1938; \skt{औष्णतैक्ष्ण्यरौक्ष्यलाघववैशद्यगुणलक्षणं पित्तं तस्य}\ N}}\skt{ समानयोनिः }\edtext{\skt{कटुको}}{\Afootnote{A1938 H; \emph{om.}\ N}}\skt{ }\edtext{\skt{रसः}}{\Afootnote{A1938 H; \emph{om.}\ N}}\skt{ । }\edtext{\skt{सौम्य}}{\Afootnote{H; \skt{सोऽस्य}\ A1938; \skt{सौस्य}\ N}}\skt{ }\edtext{\skt{औष्ण्यादौष्ण्यमभिवर्धयति}}{\lemma{\skt{अभिवर्धयति}}\Afootnote{N H; \skt{वर्धयति,}\ A1938}}\skt{ । तैक्ष्ण्यात्तैक्ष्ण्यं रौक्ष्याद्रौक्ष्यं लाघवाल्लाघवं }\edtext{\skt{वैशद्याद्वैशद्यमिति}}{\lemma{\skt{वैशद्याद्वैशद्यम्}}\Afootnote{A1938 H; \skt{वैशद्याम्}\ N}}\skt{ ॥}\pend

\pstart \ledsidenote{\passref{1.42.8.3}}\edtext{\skt{माधुर्यस्नेहगौरवशैत्यपैच्छिल्यगुणलक्षणः}}{\Afootnote{A1938 H; \skt{माधुर्यस्नेहगौरवशैत्यपैच्छिल्यगुणलक्षण । ः}\ N}}\skt{ श्लेष्मा तस्य }\edtext{\skt{समानयो}}{\Afootnote{H; \skt{समानयो}\ A1938 N}}\skt{ }\edtext{\skt{निर्मधुरो}}{\Afootnote{H; \skt{निर्मधुरो}\ A1938 N}}\skt{ रसः । }\edtext{\skt{सौम्य}}{\Afootnote{H; \skt{सोऽस्य}\ A1938; \skt{सोस्य}\ N}}\skt{ }\edtext{\edtext{\skt{माधुर्यात्माधुर्यमभिवर्धयति}}{\lemma{\skt{माधुर्यात्माधुर्यम्}}\Afootnote{H; \skt{माधुर्यान्माधुर्यं}\ A1938; \skt{माधुर्याद्माधुर्यम्}\ N}}}{\lemma{\skt{अभिवर्धयति}}\Afootnote{N H; \skt{वर्धयति,}\ A1938}}\skt{ । }\edtext{\skt{स्नेहात्स्नेहं}}{\Afootnote{A1938 H; \skt{स्नेहात्सेनहं}\ N}}\skt{ गौरवाद्गौरवं }\edtext{\skt{शैत्याच्छैत्यं}}{\Afootnote{H; \skt{शैत्याच्छैत्यं,}\ A1938; \skt{शैत्याच्छैत्यं}\ N}}\skt{ पैच्छिल्यात्पैच्छिल्यमिति ॥}\pend

\pstart \ledsidenote{\passref{1.42.8.4}}\skt{तस्य }\edtext{\edtext{\skt{पुनर्व्विपरीतः}}{\lemma{\skt{पुनर्}}\Afootnote{N H; \skt{पुन}\ A1938}}}{\lemma{\skt{व्विपरीतः}}\Afootnote{N H; \skt{रन्ययोनिः}\ A1938}}\skt{ कटुको }\edtext{\skt{रसः}}{\Afootnote{N H; \skt{रसः; स}\ A1938}}\skt{ । श्लेष्मणः }\edtext{\skt{प्रत्यनीकत्वात् । }\edtext{\skt{कटुक्यात्माधुर्यमभिभवति}}{\lemma{\skt{अभिभवति}}\Afootnote{N H; \skt{मभिभवति,}\ A1938}}}{\lemma{\skt{प्रत्यनीकत्वात् । कटुक्यात्माधुर्यम्}}\Afootnote{H; \skt{प्रत्यनीकत्वात्कटुकत्वान्माधुर्य}\ A1938; \skt{प्रत्यनीकत्वात् । कटुत्वात्माधुर्यम्}\ N}}\skt{ । }\edtext{\skt{रौक्ष्यात्स्नेहं}}{\Afootnote{A1938 H; \skt{रौक्ष्यस्नेहं}\ N}}\skt{ }\edtext{\skt{लाघवाद्गौरवं}}{\Afootnote{A1938 H; \skt{लाघवाद्गौरव}\ N}}\skt{ }\edtext{\edtext{\edtext{\skt{औष्ण्यात्पैच्छित्यन्तदेतन्निदर्शनमात्रमुक्तं}}{\lemma{\skt{औष्ण्यात्पैच्छित्यन्}}\Afootnote{H; \skt{औष्ण्याच्छैत्यं, वैशद्यात्पैच्छिल्यमिति ।}\ A1938; \skt{औष्णात् । वैशद्यात्पैच्छिल्यं}\ N}}}{\lemma{\skt{तदेतन्}}\Afootnote{A1938 H; \skt{तदेतन्}\ N}}}{\lemma{\skt{उक्तं}}\Afootnote{N H; \skt{उक्तं भवति (४) ॥८॥}\ A1938}}\skt{ ।}\pend

\pstart \ledsidenote{\passref{1.42.9}}\skt{रसलक्षणमत }\edtext{\skt{ऊर्ध्वं}}{\Afootnote{A1938 H; \skt{ऊधे}\ N}}\skt{ वक्ष्यामः ॥ तत्र यः }\edtext{\edtext{\skt{परितोषमुत्पादयति}}{\lemma{\skt{परितोषम्}}\Afootnote{N H; \skt{परितोष}\ A1938}}}{\lemma{\skt{उत्पादयति}}\Afootnote{N H; \skt{मुत्पादयति प्रह्लादयति}\ A1938}}\skt{ }\edtext{\skt{तर्प्पयति}}{\Afootnote{N H; \skt{तर्पयति जीवयति}\ A1938}}\skt{ }\edtext{\skt{मुखोपलेपञ्जनयति}}{\lemma{\skt{जनयति}}\Afootnote{N H; \skt{जन}\ A1938}}\skt{ }\edtext{\skt{प्रह्लादयति}}{\Afootnote{N H; \skt{यति}\ A1938}}\skt{ श्लेष्माणञ्चाभिवर्धयति स मधुरः ॥ यो दन्तहर्षमुत्पादयति । }\edtext{\skt{मुखाश्रावञ्जन-यति}}{\lemma{\skt{मुखाश्रावञ्}}\Afootnote{H; \skt{मुखास्रावं}\ A1938 N}}\skt{ }\edtext{\skt{श्रद्धाञ्चोत्पादयति सोऽम्लः}}{\lemma{\skt{चोत्पादयति सो}}\Afootnote{A1938 H; \skt{चोत्पादयति सोम्}\ N}}\skt{ ॥ यो भक्तरुचिमुत्पादयति । कफप्रसेकञ्जनयति }\edtext{\skt{मार्द्दवञ्चोपादयति स}}{\lemma{\skt{चोपादयति स}}\Afootnote{H; \skt{चापादयति स}\ A1938; \skt{चोत्पादयति स}\ N}}\skt{ लवणः ॥ यो }\edtext{\edtext{\skt{जिह्वामुद्वेजयति}}{\lemma{\skt{जिह्वाम्}}\Afootnote{N H; \skt{जिह्वाग्रं}\ A1938}}}{\lemma{\skt{उद्वेजयति}}\Afootnote{N H; \skt{बाधते उद्वेगं जनयति}\ A1938}}\skt{ शिरो }\edtext{\skt{गृह्णाति}}{\Afootnote{N H; \skt{गृह्णीते}\ A1938}}\skt{ }\edtext{\skt{नासिकाञ्च}}{\Afootnote{H; \skt{नासिकां}\ A1938; \skt{नासिकाञ्चा}\ N}}\skt{ }\edtext{\skt{श्रावयति स}}{\Afootnote{H; \skt{स्रावयति स}\ A1938 N}}\skt{ कटुकः ॥ यो गले }\edtext{\edtext{\skt{शोषमुत्पाद-यति}}{\lemma{\skt{शोषम्}}\Afootnote{N H; \skt{चोष}\ A1938}}}{\lemma{\skt{उत्पाद-यति}}\Afootnote{N H; \skt{मुत्पादयति}\ A1938}}\skt{ }\edtext{\edtext{\skt{मुखवैषद्यमुत्पादयति}}{\lemma{\skt{मुखवैषद्यम्}}\Afootnote{H; \skt{मुखवैशद्यं}\ A1938; \skt{मुखवैशम्यम्}\ N}}}{\lemma{\skt{उत्पादयति}}\Afootnote{N H; \skt{जनयति}\ A1938}}\skt{ }\edtext{\skt{भक्तरुचिञ्चोपाद्यति स}}{\lemma{\skt{चोपाद्यति स}}\Afootnote{H; \skt{चापादयति हर्षं च स}\ A1938; \skt{चापादयति स}\ N}}\skt{ }\edtext{\skt{तिक्तः ॥ य}}{\Afootnote{N H; \skt{तिक्तः; यो}\ A1938}}\skt{ }\edtext{\skt{आस्यम्परिशोषयति-}}{\lemma{\skt{आस्यम्}}\Afootnote{N H; \skt{वक्त्रं}\ A1938}}\skt{ }\edtext{\skt{जिह्वां}}{\Afootnote{A1938 H; \skt{जिह्वा}\ N}}\skt{ स्तम्भयति }\edtext{\skt{कण्ठम्बध्नाति}}{\lemma{\skt{बध्नाति}}\Afootnote{A1938 H; \skt{वध्नाति}\ N}}\skt{ }\edtext{\skt{हृदयम्पीडयति स}}{\lemma{\skt{पीडयति स}}\Afootnote{N H; \skt{कर्षति पीडयति च स}\ A1938}}\skt{ }\edtext{\skt{कषायः}}{\Afootnote{N H; \skt{कषाय इति ॥९॥}\ A1938}}\skt{ ॥}\edtext{}{\lemma{\skt{ऽम्लः ॥ यो}}\Afootnote{A1938 H; \skt{व्ल ॥ यो}\ N}}\pend

\pstart \ledsidenote{\passref{1.42.10.1}}\skt{रसगुणानत ऊर्ध्वम्वक्ष्यामः ॥ तत्र मधुरो रसो }\edtext{\edtext{\skt{रक्तमांसदोऽस्थिमज्जोजः}}{\lemma{\skt{रक्तमांसदो}}\Afootnote{H; \skt{रसरक्तमांसमेदो}\ A1938; \skt{रक्तमांसमेदो}\ N}}}{\lemma{\skt{ऽस्थिमज्जोजः}}\Afootnote{H; \skt{ऽस्थिमज्जौजः}\ A1938; \skt{स्थिमजोजः}\ N}}\skt{ }\edtext{\edtext{\skt{शुक्रवर्धनश्चाक्षुष्यः}}{\lemma{\skt{शुक्रवर्धनश्}}\Afootnote{N H; \skt{शुक्रस्तन्यवर्धनश्}\ A1938}}}{\lemma{\skt{चाक्षुष्यः}}\Afootnote{H; \skt{चक्षुष्यः}\ A1938 N\textsuperscript{pc}; \skt{चक्षुशष्यः}\ N\textsuperscript{ac}}}\skt{ }\edtext{\skt{केश्यौवलकृत्सन्धानः}}{\Afootnote{H; \skt{केश्यो वर्ण्यो बलकृत्सन्धानः}\ A1938; \skt{केश्यो वलकृत्सन्धान}\ N}}\skt{ }\edtext{\skt{शोणितः शोणितप्रसादो}}{\Afootnote{H; \skt{शोणितरसप्रसादनो}\ A1938; \skt{शोणितप्रसादो}\ N}}\skt{ }\edtext{\skt{वालवृद्धक्षतक्षीणहितः}}{\Afootnote{H; \skt{बालवृद्धक्षतक्षीणहितः}\ A1938; \skt{वालवृद्धक्षक्षीणहितः}\ N}}\skt{ }\edtext{\skt{षट्पदपिपीलिकानामिष्टतमः}}{\Afootnote{H; \skt{षट्पदपिपीलिकानामिष्टतम}\ A1938; \skt{षट्पदपिपीलिकानामिष्टतमः}\ N}}\skt{ ॥ }\edtext{\skt{तृष्णामूर्च्छादाहः प्रशमनः}}{\Afootnote{H; \skt{स्तृष्णामूर्च्छादाहप्रशमनः}\ A1938; \skt{तृष्णामूर्च्छादाहप्रशमनः}\ N}}\skt{ । }\edtext{\skt{षडिन्द्रियप्रसादनश्चेति । स}}{\lemma{\skt{चेति । स}}\Afootnote{N H; \skt{कृमिकफकरश्चेति; स}\ A1938}}\skt{ }\edtext{\edtext{\skt{एवंगुणोऽप्येक}}{\lemma{\skt{एवंगुणो}}\Afootnote{A1938 H; \skt{एवंगु}\ N}}}{\lemma{\skt{ऽप्येक}}\Afootnote{H; \skt{ऽप्येक}\ A1938; \skt{णोप्येक}\ N}}\skt{ }\edtext{\skt{एवात्यर्थमुपयुज्यमानः}}{\lemma{\skt{उपयुज्यमानः}}\Afootnote{N H; \skt{आसेव्यमानः}\ A1938}}\skt{ }\edtext{\skt{कासः}}{\Afootnote{N H; \skt{कास}\ A1938}}\skt{ }\edtext{\skt{श्वासो लसकवमथुर्व्वदनमाधुर्यस्वरोपघातकृमिगलगण्डानापादयति}}{\Afootnote{H; \skt{श्वासालसकवमथुवदनमाधुर्यस्वरोपघातकृमिगलगण्डानापादयति}\ A1938; \skt{श्वासालसकवमथुर्वदन माधुर्यस्वरोपघातक्रिमिगलगण्डानापादयति ।}\ N\textsuperscript{ac}; \skt{श्वासालसकवमथुर्वदतन माधुर्यस्वरोपघातक्रिमिगलगण्डानापादयति ।}\ N\textsuperscript{pc}}}\skt{ । }\edtext{\edtext{\skt{तथार्व्वुदश्लीपदबस्तिगुदोपलेपाभिष्यन्दप्रभृतीन्नयनविकारान्नुपजयति}}{\lemma{\skt{तथार्व्वुदश्लीपदबस्तिगुदोपलेपाभिष्यन्दप्रभृतीन्}}\Afootnote{\skt{तथाऽर्बुदश्लीपदबस्तिगुदोपलेपाभिष्यन्दप्रभृती ञ्}\ A1938; \skt{तथार्बुदश्लीपदवस्तिगुडोपलेपाभिष्यन्दप्रभृतीन्}\ N; \skt{तथार्व्वुदश्लीपदवस्तिगुदोपलेपाभिष्यन्दप्रभृतीन्}\ H}}}{\lemma{\skt{नयनविकारान्नुपजयति}}\Afootnote{H; \skt{जनयति(१) ।१०।}\ A1938; \skt{नयत विकारानुपजनयति ॥}\ N}}\skt{ ॥}\pend

\pstart \ledsidenote{\passref{1.42.10.2}}\edtext{\skt{अम्लस्तु}}{\Afootnote{H; \skt{अम्लो}\ A1938; \skt{अम्व्लस्तु}\ N}}\skt{ }\edtext{\skt{जरणः}}{\Afootnote{A1938 H; \skt{जरण}\ N}}\skt{ }\edtext{\skt{पाचनपवननिग्रहणो}}{\Afootnote{H; \skt{पाचनो दीपनः पवननिग्रहणो}\ A1938; \skt{पाचणपवननिग्रह}\ N}}\skt{ }\edtext{\skt{नुलोमनो}}{\Afootnote{H; \skt{ऽनुलोमनः}\ A1938; \skt{णोनुलोमनो}\ N}}\skt{ }\edtext{\skt{विदाही}}{\Afootnote{N H; \skt{कोष्ठविदाही}\ A1938}}\skt{ }\edtext{\skt{वहिःशीतः}}{\Afootnote{N H; \skt{बहिःशीतः}\ A1938}}\skt{ }\edtext{\skt{क्लेदनः}}{\Afootnote{A1938 H; \skt{क्लेदनः प्रीणनः}\ N}}\skt{ प्रायशो }\edtext{\edtext{\skt{हृद्यश्चेति}}{\lemma{\skt{हृद्यश्}}\Afootnote{N H; \skt{ह्यद्य}\ A1938}}\skt{ । स}}{\lemma{\skt{चेति । स}}\Afootnote{N H; \skt{श्चेति; स}\ A1938}}\skt{ }\edtext{\edtext{\skt{खल्वेवंगुणोऽप्येक}}{\lemma{\skt{खल्वेवंगुणो}}\Afootnote{H; \skt{एवङ्गुणो}\ A1938; \skt{खल्वेवं गु}\ N}}}{\lemma{\skt{ऽप्येक}}\Afootnote{H; \skt{ऽप्येक}\ A1938; \skt{णोप्येक}\ N}}\skt{ }\edtext{\skt{एवात्यर्थमुपयुज्यमानो}}{\lemma{\skt{उपयुज्यमानो}}\Afootnote{N H; \skt{उपसेव्यमानो}\ A1938}}\skt{ }\edtext{\edtext{\skt{दन्तहर्षनयनसम्मीलनरोमसम्वेजनकफविलायनशरीरप्रशिथिलतामापादयति}}{\lemma{\skt{दन्तहर्षनयनसम्मीलनरोमसम्वेजनकफविलायनशरीरप्रशिथिलताम्}}\Afootnote{H; \skt{दन्तहर्षनयनसम्मीलनरोमसंवेजनकफविलयनशरीरशैथिल्या}\ A1938; \skt{दन्तहर्षनयनसंमीलनरोमसंवेजनकफविलोपनशरीरप्रशिथिलताम्}\ N}}}{\lemma{\skt{आपादयति}}\Afootnote{N H; \skt{न्यापादयति,}\ A1938}}\skt{ । तथा }\edtext{\skt{क्षतविहतदग्धदष्टभग्नशू-नच्युतावमथितच्छियविसर्प्पितच्छिन्नविद्धोत्पिष्टादीनि}}{\Afootnote{H; \skt{क्षताभिहतदग्धदष्टभग्नरुग्णशूनप्रच्युतावमूत्रितविसर्पितच्छिन्नभिन्नविद्धोत्पिष्टादीनि}\ A1938; \skt{क्षतविहतदग्धदष्टभग्नशूनच्युताठमत्रित परिसर्पितच्छिन्नविद्धोत्पिष्टादीनि}\ N\textsuperscript{ac}; \skt{क्षतविहतदग्धदष्टभग्नशूच्युताठमत परिसर्पितच्छिन्नविद्धोत्पिष्टादीनि}\ N\textsuperscript{pc}}}\skt{ }\edtext{\skt{पाच-यत्याग्नेयस्वभावत्व्वात्}}{\Afootnote{H; \skt{पाचयत्याग्नेयस्वभावात्परिदहति कण्ठमुरो हृदयं चेति(२) ।१०।}\ A1938; \skt{पाचयत्याग्नेयस्वभावत्वात् ॥}\ N}}\skt{ ॥}\pend

\pstart \ledsidenote{\passref{1.42.10.3}}\edtext{\skt{लवणस्तु}}{\Afootnote{N H; \skt{लवणः}\ A1938}}\skt{ संशोधनः पाचनो }\edtext{\skt{विश्लेषणस्तर्प्पणः}}{\lemma{\skt{तर्प्पणः}}\Afootnote{H; \emph{om.}\ A1938; \skt{तर्पण}\ N}}\skt{ क्लेदनः }\edtext{\skt{शैथिल्यकृत्सर्व्वरसप्रत्यनीकभूतो}}{\Afootnote{H; \skt{शैथिल्यकृदुष्णः सर्वरसप्रत्यनीको}\ A1938; \skt{शैथिल्यकृत्सर्वरसप्रत्यनीभूतो}\ N}}\skt{ मार्ग्गविशोधनः }\edtext{\edtext{\skt{सर्व्वशरीरावयवमार्द्दवकरश्चेति}}{\lemma{\skt{सर्व्वशरीरावयवमार्द्दवकरश्}}\Afootnote{N H; \skt{सर्वशरीरावयवमार्दवकर}\ A1938}}\skt{ । स}}{\lemma{\skt{चेति । स}}\Afootnote{N H; \skt{श्चेति; स}\ A1938}}\skt{ }\edtext{\skt{खलु}}{\Afootnote{H; \emph{om.}\ A1938; \skt{खल्वे}\ N}}\skt{ }\edtext{\edtext{\edtext{\skt{एवंगुणोऽप्येक}}{\lemma{\skt{एवंगुणो}}\Afootnote{A1938 H; \skt{वंगु}\ N}}}{\lemma{\skt{ऽप्य्}}\Afootnote{H; \skt{ऽ}\ A1938; \skt{णो}\ N}}}{\lemma{\skt{एक}}\Afootnote{H; \skt{प्येक}\ A1938 N}}\skt{ }\edtext{\edtext{\skt{एवात्यर्थमुपसेव्यमानो}}{\lemma{\skt{एवात्यर्थम्}}\Afootnote{N H; \skt{एवात्यर्थ}\ A1938}}}{\lemma{\skt{उपसेव्यमानो}}\Afootnote{N H; \skt{मासेव्यमानो}\ A1938}}\skt{ }\edtext{\skt{गात्रकण्डूको थ}}{\Afootnote{N H; \skt{गात्रकण्डूकोठ}\ A1938}}\skt{ }\edtext{\skt{शोफवैवर्ण्णकरः}}{\Afootnote{H; \skt{शोफवैवर्ण्य}\ A1938; \skt{शोफवैवर्ण्ण्यकरः}\ N}}\skt{ }\edtext{\skt{स्वरोपघातेन्द्रियोपतापानुपजनयति}}{\Afootnote{N H; \skt{पुंस्त्वोपघातेन्द्रियोपतापमु}\ A1938}}\skt{ । }\edtext{\skt{तथाक्षिमुखपाकरक्तपित्तवातशोणिताम्लीकाप्रभृतीन्विकारानुपजनयति}}{\Afootnote{H; \skt{खाक्षिपाकरक्तपित्तवातशोणिताम्लीकाप्रभृतीनापादयति(३) ।१०।}\ A1938; \skt{तथा क्षिमुखपाकरक्तपित्तवात्तपित्तशोणिताम्व्लीकाप्रभृतीन्विकारानुपजनयति ॥}\ N}}\skt{ ॥}\pend

\pstart \ledsidenote{\passref{1.42.10.4}}\edtext{\skt{कटुकस्तु}}{\Afootnote{N H; \skt{कटुको}\ A1938}}\skt{ दीपनः पाचनो रोचनः }\edtext{\skt{शोधनः}}{\Afootnote{A1938 H; \skt{शोधन}\ N}}\skt{ }\edtext{\skt{स्थौल्या ल्यालसकविषकुष्ठकण्डूप्रशमनः}}{\Afootnote{H; \skt{स्थौल्यालस्यकफकृमिविषकुष्ठकण्डूप्रशमनः}\ A1938; \skt{स्थौल्यालसकविषकुष्ठकण्डूप्रशमनः ।}\ N}}\skt{ }\edtext{\skt{सन्धिबन्धच्छेदनो}}{\Afootnote{H; \skt{सन्धिबन्धविच्छेदनो}\ A1938; \skt{सन्धिवन्धच्छेद}\ N}}\skt{ }\edtext{\edtext{\skt{वसादनस्स्तन्यशुक्रकफमेदसामुपहन्ता}}{\lemma{\skt{वसादनस्}}\Afootnote{H; \skt{ऽवसादनः}\ A1938; \skt{नोवसादन}\ N}}}{\lemma{\skt{स्तन्यशुक्रकफमेदसाम्}}\Afootnote{N H; \skt{स्तन्यशुक्रमेदसाम्}\ A1938}}\skt{ चेति । स }\edtext{\edtext{\edtext{\skt{खल्वेवंगुणोऽप्येक}}{\lemma{\skt{खल्वेवंगुणो}}\Afootnote{H; \skt{एवङ्गुणो}\ A1938; \skt{खल्वेवंगु}\ N}}}{\lemma{\skt{ऽप्य्}}\Afootnote{H; \skt{ऽ}\ A1938; \skt{णो}\ N}}}{\lemma{\skt{एक}}\Afootnote{H; \skt{प्येक}\ A1938 N}}\skt{ }\edtext{\skt{एवात्यर्थमुपयुज्यमानो}}{\lemma{\skt{उपयुज्यमानो}}\Afootnote{N H; \skt{उपसेव्यमानो}\ A1938}}\skt{ }\edtext{\skt{भ्रममदगलताल्वोष्ठशोषदाहसन्तापानापादयति}}{\Afootnote{H; \skt{भ्रममदगलताल्वोष्ठशोषदाहसन्ता}\ A1938; \skt{भ्रममदगलताल्वोष्ठप्रशोषदाहसन्तापानापादयति ॥}\ N}}\skt{ । }\edtext{\skt{तथा}}{\Afootnote{N H; \skt{प}\ A1938}}\skt{ }\edtext{\skt{वलविघातकम्पतोदभेदकृत्करचरणपार्श्वपृष्ठप्रभृतिषु च}}{\Afootnote{N H; \skt{बलविघातकम्पतोदभेदकृत्करचरणपार्श्वपृष्ठप्रभृतिषु च}\ A1938}}\skt{ }\edtext{\skt{वातशूलानुपजनयति}}{\Afootnote{N H; \skt{वातशूलानापादयति(४) ।१०।}\ A1938}}\skt{ ॥}\pend

\pstart \ledsidenote{\passref{1.42.10.5}}\edtext{\edtext{\skt{तिक्तस्तु}}{\lemma{\skt{तिक्तस्}}\Afootnote{H; \skt{तिक्त}\ A1938; \emph{om.}\ N}}}{\lemma{\skt{तु}}\Afootnote{H; \skt{श्छेदनो}\ A1938; \skt{तिक्तस्तु}\ N}}\skt{ रोचनो दीपनः शोधनः }\edtext{\skt{कण्डूकोष्ठतृष्णामूर्च्छाप्रशनः}}{\Afootnote{H; \skt{कण्डूकोठतृष्णामूर्च्छाज्वरप्रशमनः}\ A1938; \skt{कण्डूकोथतृष्णामूर्च्छाज्वरप्रशमनः ॥}\ N}}\skt{ स्तन्यशोधनो }\edtext{\edtext{\skt{विण्मूत्रक्लेदमेदोवसापूयोपशोषणश्चेति}}{\lemma{\skt{विण्मूत्रक्लेदमेदोवसापूयोपशोषणश्}}\Afootnote{N H; \skt{विण्मूत्रक्लेदमेदोवसापूयोपशोषण}\ A1938}}\skt{ । स}}{\lemma{\skt{चेति । स}}\Afootnote{N H; \skt{श्चेति; स}\ A1938}}\skt{ }\edtext{\edtext{\edtext{\skt{खल्वेवङ्गुणोऽप्येक}}{\lemma{\skt{खल्वेवङ्गुणो}}\Afootnote{H; \skt{एवङ्गुणो}\ A1938; \skt{खल्वेवंगु}\ N}}}{\lemma{\skt{ऽप्य्}}\Afootnote{H; \skt{ऽ}\ A1938; \skt{णो}\ N}}}{\lemma{\skt{एक}}\Afootnote{H; \skt{प्येक}\ A1938 N}}\skt{ }\edtext{\skt{एवात्यर्थमुपयुज्यमानो}}{\lemma{\skt{उपयुज्यमानो}}\Afootnote{N H; \skt{उपसेव्यमानो}\ A1938}}\skt{ }\edtext{\skt{गात्रमन्यास्तम्भाक्षेपकोर्द्दितशिरःशूलानुप-जनयति}}{\Afootnote{H; \skt{गात्रमन्यास्तम्भाक्षेपकार्दितशिरःशू}\ A1938; \skt{गात्रमन्यास्तम्भाक्षेपकार्दितशिरःशूलानुपजनयति ।}\ N}}\skt{ । }\edtext{\skt{तथा}}{\Afootnote{N H; \skt{ल}\ A1938}}\skt{ }\edtext{\edtext{\skt{भ्रमतोदभेदच्छेदास्यवैरस्यान्यापादयति}}{\lemma{\skt{भ्रमतोदभेदच्छेदास्यवैरस्यान्य्}}\Afootnote{H; \skt{भ्रमतोदभेदच्छेदास्यवैरस्यान्य्}\ A1938; \skt{भ्रमतोदभेदच्छेदास्यवैरस्या}\ N}}}{\lemma{\skt{आपादयति}}\Afootnote{A1938 H; \skt{न्यापादयति ॥}\ N}}\skt{ ॥}\pend

\pstart \ledsidenote{\passref{1.42.10.6}}\edtext{\skt{कषायस्तु}}{\Afootnote{N H; \skt{कषायः}\ A1938}}\skt{ }\edtext{\skt{सङ्ग्राहिको}}{\Afootnote{N H; \skt{सङ्ग्राहको}\ A1938}}\skt{ रोपणः स्तम्भनः }\edtext{\skt{शोधनोल्लेखनः}}{\Afootnote{H; \skt{शोधनो लेखनः}\ A1938; \skt{शोधणो लेखनः ।}\ N}}\skt{ }\edtext{\skt{शोषणः}}{\Afootnote{A1938 H; \emph{om.}\ N}}\skt{ }\edtext{\skt{पीडनः}}{\Afootnote{A1938 H; \skt{पीडन}\ N}}\skt{ क्लेदोपशोषणश्चेति । स }\edtext{\edtext{\skt{खल्वेवङ्गुणोऽप्येक}}{\lemma{\skt{खल्वेवङ्गुणो}}\Afootnote{H; \skt{एवङ्गुणो}\ A1938; \skt{खल्वेवंगुणो}\ N\textsuperscript{ac}; \skt{खल्वेवंगु}\ N\textsuperscript{pc}}}}{\lemma{\skt{ऽप्येक}}\Afootnote{H; \skt{ऽप्येक}\ A1938; \skt{प्येक}\ N\textsuperscript{ac}; \skt{णोप्येक}\ N\textsuperscript{pc}}}\skt{ }\edtext{\skt{एवात्यर्थमुपयुज्यमानो}}{\lemma{\skt{उपयुज्यमानो}}\Afootnote{N H; \skt{उपसेव्यमानो}\ A1938}}\skt{ }\edtext{\edtext{\skt{हृदयपीडास्यशोषोदराध्मानवाक्यग्रहमन्यास्तम्भप्रभृतीन्विकारानुपजनयति}}{\Afootnote{H; \skt{हृत्पीडास्यशोषोदराध्मानवाक्यग्रहमन्यास्तम्भ-}\ A1938; \skt{हृदयपीडास्यशोषोदराध्मानाह वाक्यग्रहणमन्यास्तम्भप्रभृतीन्विकारानुपजनयति}\ N\textsuperscript{ac} N\textsuperscript{pc}}}\skt{ । तथा }\edtext{\skt{गात्रस्फुरणचिमिचिमायनाकुञ्चनाक्षेपणप्रभृतीन्विकारानुपजनयति}}{\Afootnote{H; \skt{गात्रस्फुरणचुमुचुमायनाकुञ्चनाक्षेपणप्रभृती ञ्जनयति ॥१०॥}\ A1938; \skt{गात्रस्फुरणविमिविमायनाकुञ्चनाक्षेपणप्रभृतीन्विकारानुपजनयति ॥}\ N}}}{\lemma{\skt{प्रभृतीन्… गा}}\Afootnote{\emph{om.}\ A1938}}\skt{ ॥}\pend

\pstart \ledsidenote{\passref{1.42.11}}\edtext{\skt{सर्व्वेषामेव}}{\lemma{\skt{सर्व्वेषाम्}}\Afootnote{N H; \skt{अतः सर्वेषाम्}\ A1938}}\skt{ }\edtext{\edtext{\skt{द्रव्याण्युपदेक्ष्यामः}}{\lemma{\skt{द्रव्याण्य्}}\Afootnote{N H; \skt{द्रव्या}\ A1938}}}{\lemma{\skt{उपदेक्ष्यामः}}\Afootnote{N H; \skt{ण्युपदेक्ष्यामः ।}\ A1938}}\skt{ । तद्यथा }\edtext{\skt{काकोल्यादि}}{\Afootnote{N H; \skt{काकोल्यादिः}\ A1938}}\skt{ }\edtext{\skt{क्षीरघृतवसामज्जशालियवगोधूममाषशृङ्गाटककशेरुक कालाङ्कालुकपियालपुष्करवीजकाश्मर्यमधूकद्राक्षाखर्जूरराजादननालिकेरेक्षुविकाराः}}{\Afootnote{H; \skt{क्षीरघृतवसामज्जशालिषष्टिकयवगोधूममाषशृङ्गाटककसेरुकत्रपुसैर्वारुकर्कारुकालाबूकालिन्दकतक- गिलोड्यप्रियालपुष्करबीजकाश्मर्यमधूकद्राक्षाखर्जूरराजादनतालनालिकेरेक्षुविकार}\ A1938; \skt{क्षीरघृतवसामज्जाशालियवगोधूममाष शृङ्गाटककसेरुककलोड्यलुवुपियालपुष्करवीजकाष्मर्यमधूकध्रद्राक्षाखर्ज्जूर राजादननालिकेरेक्षुविकाराः ।}\ N\textsuperscript{ac}; \skt{क्षीरघृतवसामज्जाशालियवगोधूममाष शृङ्गाटककसेरुककलोड्यलुवुपियालपुष्करवीजकाष्मर्यमधूकद्राक्षाखर्ज्जूर राजादननालिकेरेक्षुविकाराः ।}\ N\textsuperscript{pc}}}\skt{ । }\edtext{\skt{बलातिबलात्मगुप्ताविदारीपयस्यागोक्षुरकक्षीरमोरटमधूलिकाकूष्माण्डादिः}}{\Afootnote{H; \skt{बलातिबलात्मगुप्ता- विदारीपयस्यागोक्षुरकक्षीरमोरटमधूलिकाकूष्माण्डप्रभृतीनि}\ A1938; \skt{वलातिवलात्मगुप्ताविदारीपयस्यागोक्षुरकक्षीरमोरटामधूकलिका कुष्माजण्डादिः ।}\ N\textsuperscript{ac}; \skt{वलातिवलात्मगुप्ताविदारीपयस्यागोक्षुरकक्षीरमोरटामधूलिका कुष्माण्डादिः ।}\ N\textsuperscript{pc}}}\skt{ समासेन मधुरो वर्ग्गः ॥ }\edtext{\skt{दाडिमामलकमातुलुङ्गाम्राम्रातककपित्थकरमर्द्दवदरप्राचीनामलककोशाम्रभव्यपारावतवेत्रफलतिन्तिडीकलजचाम्लवेतसदन्तशठदधितक्रसुरासौवीरतुषोदकधान्याम्लप्रभृतीनि}}{\Afootnote{H; \skt{दाडिमामलकमातुलुङ्गाम्रातककपित्थकरमर्दबदरकोलप्राचीनामलकतिन्तिडीककोशाम्रकभव्यपारावत- वेत्रफललकुचाम्लवेतसदन्तशठदधितक्रसुराशुक्तसौवीरकतुषोदकधान्याम्लप्रभृतीनि}\ A1938; \skt{दाडिमामलकमातलुङ्गाम्राम्रातककपित्थकरमर्दवदरप्राचीर्ण्णमकलकोशाम्रभव्यआपारावतवेत्रफलतिन्तिडीक कुचाम्व्लवेतसदन्तशठदधितक्रसुरासौवीरतुषोदकधान्याम्व्लप्रभृतीनि}\ N\textsuperscript{ac}; \skt{दाडिमामलकमातलुङ्गाम्राम्रातककपित्थकरमर्दवदरप्राचीर्ण्णमकलकोशाम्रभव्यपारावतवेत्रफलतिन्तिडीक कुचाम्व्लवेतसदन्तशठदधितक्रसुरासौवीरतुषोदकधान्याम्व्लप्रभृतीनि}\ N\textsuperscript{pc}}}\skt{ }\edtext{\skt{समासेनाम्लो}}{\Afootnote{A1938 H; \skt{समासेनाम्व्लो}\ N}}\skt{ वर्ग्गः ॥ }\edtext{\skt{सैन्धवसौवर्च्चलविडपाक्यरोमकसामुद्रपाचिमक्षारोषरसुवर्च्चिलप्रभृतीनि}}{\Afootnote{N H; \skt{सैन्धवसौवर्चलविडपाक्यरोमकसामुद्रकपक्त्रिमयवक्षारोषरप्रसूतसुवर्चिकाप्रभृतीनि}\ A1938}}\skt{ समासेन लवणो वर्ग्गः ॥ पिप्प-ल्यादिः }\edtext{\edtext{\edtext{\skt{सुरसादिर्म्मधुशिग्रुमूलकलशुनसुमुखशीतशिवकुष्ठदेवदारुहरेणुकाऽवल्गुजफलचण्डागुग्गुलुमुस्तलाङ्गलकीशुकनासीपीलुप्रभृती-नि}}{\lemma{\skt{सुरसादिर्}}\Afootnote{A1938 H; \skt{सुरसा}\ N}}}{\lemma{\skt{म्मधुशिग्रुमूलकलशुनसुमुखशीतशिवकुष्ठदेवदारुहरेणुका}}\Afootnote{H; \skt{शिग्रुमधुशिग्रुमूलकलशुनसुमुखशीतशिवकुष्ठदेवदारुहरेणुका}\ A1938; \skt{दिर्मधुशिग्रुमूलकलशुनसुमुखशीतशिवकुष्ठदेवदारुकरेणुका}\ N}}}{\lemma{\skt{ऽवल्गुजफलचण्डागुग्गुलुमुस्तलाङ्गलकीशुकनासीपीलुप्रभृती-नि}}\Afootnote{H; \skt{वाल्गुजफल- चण्डागुग्गुलुमुस्तलाङ्गलकीशुकनासापीलुप्रभृतीनि}\ A1938; \skt{वल्गुजफलचण्डागुग्गुलु मुस्तालाङ्गलकीशुकनाषीपीलुप्रभृतीनि}\ N}}\skt{ }\edtext{\skt{शालसारादिश्च}}{\Afootnote{H; \skt{सालसारादिश्च}\ A1938 N}}\skt{ }\edtext{\skt{प्रायशः}}{\Afootnote{A1938 H; \skt{प्रायषः}\ N}}\skt{ कटुको वर्ग्गः ॥ }\edtext{\skt{आरग्वधादिर्ग्गुडूच्यादिर्म्मण्डूकपर्ण्णीवेत्रकरीरहरिद्राद्वयेन्द्रयववरुणस्वादुकण्टकसप्तपर्ण्णवृहतीद्वयशंखिनीतृवृत्कृतवेधनकर्क्कोटककारवेल्लवार्त्ताकुकरवीरसुमनाःशंखपुष्प्यपामार्ग्गत्रायमाणाशोकरोहिवैजयन्तीसुवर्च्चलावृश्चिकालीज्योतिष्मतीप्रभृतीनि}}{\Afootnote{H; \skt{आरग्वधादिर्गुडूच्यादिर्मण्डूकपर्णीवेत्रकरीरहरिद्राद्वयेन्द्रयववरुणस्वादुकण्टकसप्तपर्णबृहतीद्वयशङ्खिनीद्रवन्ती- त्रिवृत्कृतवेधनकर्कोटककारवेल्लवार्ताककरीरकरवीरसुमनःशङ्खपुष्प्यपामार्गत्रायमाणाशोकरोहिणीवैजयन्ती- सुवर्चलापुनर्नवावृश्चिकालीज्योतिष्मतीप्रभृतीनि}\ A1938; \skt{आरग्वधादिर्गुडुच्यादिर्मण्डूकपर्ण्णीवेत्रकरीरहरिद्राद्वयेन्द्रयववरुणस्वादुकणअटकसप्तपर्ण्णवृहतीद्वयशंङ्खिनी तृवृत्कृतवेधनकर्कोटककारवेल्लकवार्त्ताकरवीसुमना शङ्खपुष्प्यपामार्गत्रायमाणाशोकरोहि वेजयन्तीसुवर्च्चलावृश्चिकालीज्योतिष्मतीप्रभृतीनि}\ N}}\skt{ }\edtext{\skt{तिक्तो}}{\Afootnote{N H; \skt{समासेन तिक्तो}\ A1938}}\skt{ वर्ग्गः ॥ }\edtext{\edtext{\skt{न्यग्रोधादिरम्बष्ठादिर्लोध्रादिः}}{\lemma{\skt{न्यग्रोधादिरम्बष्ठादिर्}}\Afootnote{H; \skt{न्यग्रोधादिरम्बष्ठा}\ A1938; \skt{न्यग्रोधादिरम्वष्ठादि}\ N}}}{\lemma{\skt{लोध्रादिः}}\Afootnote{H; \skt{दिः}\ A1938; \skt{र्लोध्रादिः}\ N}}\skt{ }\edtext{\skt{प्रियङ्वादिस्त्रिफला}}{\Afootnote{H; \skt{प्रियङ्ग्वादी रोध्रादिस्त्रिफला}\ A1938; \skt{प्रियङ्ग्वादिः त्रिफला}\ N}}\skt{ }\edtext{\skt{शल्लकीजम्व्वस्थितिन्दुकादीनि}}{\Afootnote{H; \skt{शल्लकीजम्ब्वाम्रबकुलतिन्दुकफलानि}\ A1938; \skt{शल्लकीजम्वस्थितिन्दुकादीनि}\ N}}\skt{ }\edtext{\skt{कतकशाकफलानि}}{\Afootnote{N H; \skt{कतकशाकफलपाषाणभेदकवनस्पतिफलानि}\ A1938}}\skt{ सालसारादिश्च }\edtext{\skt{प्रायसः}}{\Afootnote{H; \skt{प्रायशः}\ A1938 N}}\skt{ }\edtext{\skt{कुरवककोविदारजीवन्तीसुनिषण्णकप्रभृतीनि}}{\Afootnote{N H; \skt{कुरुवककोविदारकजीवन्तीचिल्लीपालङ्क्यासुनिषण्णकप्रभृतीनि}\ A1938}}\skt{ वरकादयो }\edtext{\skt{मुद्गादयश्च}}{\Afootnote{A1938 H; \skt{मुद्गादयौश्च}\ N}}\skt{ समासेन कषा-यो वर्ग्गः ॥}\pend

\pstart \ledsidenote{\passref{1.42.12}}\skt{तत्रैतेषां }\edtext{\skt{रसानां}}{\Afootnote{A1938 H; \skt{रषानां}\ N}}\skt{ }\edtext{\edtext{\skt{संयोगास्त्रिषष्टिर्व्भवन्ति}}{\lemma{\skt{संयोगास्}}\Afootnote{H; \skt{संयोगा}\ A1938 N}}}{\lemma{\skt{त्रिषष्टिर्व्भवन्ति}}\Afootnote{H; \skt{स्त्रिषष्टिर्भवन्ति ।}\ A1938; \skt{त्त्रिषष्टिर्भवति ।}\ N}}\skt{ ॥ तद्यथा पञ्चदश द्विकाः । }\edtext{\edtext{\skt{विंशतिस्त्रिकाः}}{\lemma{\skt{विंशतिस्}}\Afootnote{H; \skt{विंशति}\ A1938 N}}}{\lemma{\skt{त्रिकाः}}\Afootnote{N H; \skt{स्त्रिकाः,}\ A1938}}\skt{ । पञ्चदश चतुष्काः । षट्पञ्चकाः । }\edtext{\skt{एकषट्काः}}{\Afootnote{H; \skt{ए}\ A1938; \skt{एकः षट्काः ।}\ N}}\skt{ । }\edtext{\skt{एकैकशश्च}}{\Afootnote{H; \skt{कशः}\ A1938; \skt{कैकसश्च}\ N}}\skt{ }\edtext{\skt{शड्रसा}}{\Afootnote{H; \skt{षड्रसाः}\ A1938 N}}\skt{ }\edtext{\skt{इति}}{\Afootnote{N H; \skt{एकः षट्क इति ।}\ A1938}}\skt{ । तेषामन्यत्र }\edtext{\skt{वक्ष्यामः}}{\Afootnote{N H; \skt{प्रयोजनानि वक्ष्यामः ॥१२॥}\ A1938}}\skt{ ॥}\pend

\pstart\versepar \ledsidenote{\passref{1.42.13}}\edtext{\skt{भवति}}{\Afootnote{A1938 H; \skt{भ वति चात्र}\ N}}\skt{ }\edtext{\skt{चात्र}}{\Afootnote{A1938 H; \emph{om.}\ N}}\skt{ ॥ }\pend
\pstart\versepar \skt{}\edtext{\skt{युक्ताः}}{\Afootnote{N H; \skt{जग्धाः}\ A1938}}\skt{ }\edtext{\skt{षडधिगच्छन्ति}}{\Afootnote{A1938 H; \skt{षडधि गच्छति}\ N}}\skt{ }\edtext{\skt{वलिनो}}{\Afootnote{N H; \skt{बलिनो}\ A1938}}\skt{ वश्यतां रसाः । }\pend
\pstart\versepar \skt{यथा }\edtext{\skt{दोषाः}}{\Afootnote{H; \emph{om.}\ A1938; \skt{दोषा}\ N}}\skt{ प्रकुपिता }\edtext{\skt{वशं}}{\Afootnote{N H; \skt{दोषा वशं}\ A1938}}\skt{ यान्ति }\edtext{\skt{वलीयस}}{\Afootnote{N H; \skt{बलीयसः ॥१३॥}\ A1938}}\skt{ }\edtext{\skt{इति}}{\Afootnote{N H; \emph{om.}\ A1938}}\skt{ ॥}\pend

\pstart \ledsidenote{\passref{1.42.trailer}}\edtext{\skt{॥ ४२ ॥}}{\Afootnote{\skt{इति सुश्रुतसंहितायां सूत्रस्थाने रसविशेषविज्ञानीयो नाम द्विचत्वारिंशत्तमोऽध्यायः ॥४२॥}\ A1938; \skt{फ २ ४२}\ N; \skt{४२ ॥}\ H}}\pend

\endnumbering
\beginnumbering
\eledsection*{Adhyāya 43}
\pstart \ledsidenote{\passref{1.43.1}}\skt{अथातो }\edtext{\skt{वमनद्रव्यविकल्पविज्ञानीयम्व्याख्यास्यामः}}{\lemma{\skt{व्याख्यास्यामः}}\Afootnote{H; \skt{अध्यायं व्याख्यास्यामः ॥}\ A1938; \emph{om.}\ N}}\skt{ ॥}\pend

\pstart \ledsidenote{\passref{1.43.3.part1}}\edtext{\edtext{\edtext{\skt{वमनद्रव्यानाम्फलाख्यानाम्मध्ये}}{\lemma{\skt{वमनद्रव्यानाम्}}\Afootnote{H; \skt{वमनद्रव्याणां}\ A1938 N}}}{\lemma{\skt{फलाख्यानाम्}}\Afootnote{N H; \skt{फलादीनां}\ A1938}}}{\lemma{\skt{मध्ये}}\Afootnote{N H; \emph{om.}\ A1938}}\skt{ मदनफलानि श्रेष्ठतमानि भवन्ति ॥ }\edtext{\skt{तत्र}}{\Afootnote{N H; \skt{अथ}\ A1938}}\skt{ }\edtext{\edtext{\skt{मदनपुष्पाणामातपशुष्काणा}}{\lemma{\skt{मदनपुष्पाणाम्}}\Afootnote{N H; \skt{मदनपुष्पाणा}\ A1938}}}{\lemma{\skt{आतपशुष्काणा}}\Afootnote{H; \skt{मातपपरिशुष्काणां}\ A1938; \skt{आतपशुष्काणां}\ N}}\skt{ }\edtext{\skt{चूर्णप्रकुञ्चं}}{\Afootnote{H; \skt{चूर्णप्रकुञ्जं}\ A1938; \skt{चूर्ण्णं प्रकुंचं}\ N}}\skt{ }\edtext{\edtext{\skt{प्रत्यक्पुष्पीनिम्वकषायाणामन्यतमेनालो}}{\lemma{\skt{प्रत्यक्पुष्पीनिम्वकषायाणाम्}}\Afootnote{H; \skt{प्रत्यक्पुष्पासदापुष्पीनिम्बकषायाणा}\ A1938; \skt{प्रत्यक्पुष्पीमदापुष्पीनिम्बकषायाणांम्}\ N}}}{\lemma{\skt{अन्यतमेनालो}}\Afootnote{H; \skt{मन्यतमेनालोड्य मधुसैन्धवयुक्तां पुष्पचूर्णमात्रां}\ A1938; \skt{अन्यतमेनालोड्य}\ N}}\skt{ पाययित्वा }\edtext{\skt{वामयेत् । }\edtext{\edtext{\edtext{\skt{मदनशत्वाटूचूर्ण्णान्येवम्मधुलवणयुक्तान्यभिप्रतप्तानि}}{\lemma{\skt{एवम्}}\Afootnote{N H; \skt{वं वा}\ A1938}}}{\lemma{\skt{मधुलवणयुक्तान्य्}}\Afootnote{N H; \skt{बकुलरम्यकोपयुक्तानि मधुलवणयुक्ता}\ A1938}}}{\lemma{\skt{अभिप्रतप्तानि}}\Afootnote{N H; \skt{न्यभिप्रतप्तानि;}\ A1938}}}{\lemma{\skt{वामयेत् । मदनशत्वाटूचूर्ण्णान्य्}}\Afootnote{H; \skt{वामयेत् । मदनशलाटुचूर्णान्ये}\ A1938; \skt{वामयेत् । मदनसलाटूचूर्ण्णान्य्}\ N}}\skt{ । }\edtext{\skt{मदनशलाटुचूर्ण्णसिद्धाम्वा}}{\Afootnote{A1938 H; \skt{मदनसलाटुचूर्ण्णसिद्धाम्वा}\ N}}\skt{ }\edtext{\skt{तिलयवागून्}}{\Afootnote{N H; \skt{तिलतण्डुलयवागूम् ।}\ A1938}}\skt{ । }\edtext{\skt{निवृत्ताम्वा}}{\Afootnote{H; \skt{निर्वृत्तानां वा}\ A1938; \skt{निर्वृत्ताम्वा}\ N}}\skt{ }\edtext{\skt{नातिहरिति}}{\Afootnote{H; \skt{नातिहरित}\ A1938 N}}\skt{ पाण्डूनां }\edtext{\skt{कुशमूढाववद्धमृद्गोमयप्रलिप्तानां}}{\Afootnote{H; \skt{कुशमूढावबद्धमृद्गोमयप्रलिप्तानां}\ A1938; \skt{कुसुमूटाववद्धमृद्गोमयप्रलिप्तानां}\ N}}\skt{ }\edtext{\edtext{\skt{यववूसमाषशाल्यादिधान्यराशावष्ठरात्रोषितक्लिन्नभिन्नानाम्फलपिप्पलीराहृत्यातपे}}{\lemma{\skt{यववूसमाषशाल्यादिधान्यराशावष्ठरात्रोषितक्लिन्नभिन्नानाम्}}\Afootnote{H; \skt{यवबु(तु)षमुद्गमाषशाल्यादिधान्यराशावष्टरात्रोषितक्लिन्नभिन्नानां फलानां}\ A1938; \skt{यववुसमाषसाल्यादिधान्यराशावष्टरात्रोषितक्लिन्नभिन्नानां}\ N}}}{\lemma{\skt{फलपिप्पलीराहृत्यातपे}}\Afootnote{H; \skt{फलपिप्पलीरुद्धृत्यातपे}\ A1938; \skt{फलं पिप्पलीभिराहृत्यातपे}\ N}}\skt{ }\edtext{\skt{परिशोषयेत्तासां}}{\Afootnote{N H; \skt{शोषयेत्, तासां}\ A1938}}\skt{ }\edtext{\skt{दधिपलविमृदितपरिशुष्काणां}}{\Afootnote{N H; \skt{दधिमधुपललविमृदितपरिशुष्काणां}\ A1938}}\skt{ }\edtext{\edtext{\edtext{\skt{सुभाजनस्थानामन्तर्न्नखमुष्ठिमुष्णेन}}{\lemma{\skt{सुभाजनस्थानाम्}}\Afootnote{N H; \skt{सुभाजनस्थाना}\ A1938}}}{\lemma{\skt{अन्तर्न्नखमुष्ठिम्}}\Afootnote{N H; \skt{मन्तर्नखमुष्टि}\ A1938}}}{\lemma{\skt{उष्णेन}}\Afootnote{N H; \skt{मुष्णे}\ A1938}}\skt{ }\edtext{\skt{यष्ठीमधुकषायेण}}{\Afootnote{N H; \skt{यष्टीमधुककषाये}\ A1938}}\skt{ }\edtext{\skt{कोविदाराध्यन्यतमकाषायेण वा}}{\Afootnote{H; \skt{कोविदारादीनामन्यतमे वा कषाये प्र}\ A1938; \skt{कोविदाराध्वन्यतमकषायेण वा}\ N}}\skt{ }\edtext{\skt{विमृद्य}}{\Afootnote{N H; \skt{मृद्य}\ A1938}}\skt{ }\edtext{\edtext{\edtext{\edtext{\edtext{\skt{त्रिरात्रपर्युषितम्मधुसैन्धवयुक्तमाशीर्भिरभिमन्त्रितमुदङ्मुखः}}{\lemma{\skt{त्रिरात्रपर्युषितम्}}\Afootnote{N H; \skt{रात्रिपर्युषितं}\ A1938}}}{\lemma{\skt{मधुसैन्धवयुक्तम्}}\Afootnote{N H; \skt{मधुसैन्धवयुक्त}\ A1938}}}{\lemma{\skt{आशीर्भिर्}}\Afootnote{H; \skt{माशीर्भिर्}\ A1938; \skt{आसीर्भिर्}\ N}}}{\lemma{\skt{अभिमन्त्रितम्}}\Afootnote{N H; \skt{अभिमन्त्रित}\ A1938}}}{\lemma{\skt{उदङ्मुखः}}\Afootnote{N H; \skt{मुदङ्मुखः}\ A1938}}\skt{ }\edtext{\edtext{\skt{प्राङ्मुखमातुरम्पाययेतनेन}}{\lemma{\skt{प्राङ्मुखम्}}\Afootnote{N H; \skt{प्राङ्मुख}\ A1938}}}{\lemma{\skt{आतुरम्}}\Afootnote{N H; \skt{मातुरं}\ A1938}}\skt{ }\edtext{\skt{मन्त्रेण}}{\Afootnote{H; \skt{मन्त्रेणा}\ A1938 N}}\skt{ }\edtext{\skt{अभिमन्त्र्य}}{\Afootnote{H; \skt{भिमन्त्र्य}\ A1938 N}}\skt{ ॥}\pend

\pstart\versepar \ledsidenote{\passref{1.43.3.1}}\edtext{\skt{ब्रह्मदक्षाश्विरूद्रेन्द्रभूचन्द्रार्क्कानिलानलाः}}{\Afootnote{H; \skt{“ब्रह्मदक्षाश्विरुद्रेन्द्रभूचन्द्रार्कानलानिलाः।}\ A1938; \skt{ब्रह्मादक्षासिरुद्रेन्द्रभूचन्द्रार्कानिलानलाः ।}\ N}}\skt{ । }\pend
\pstart\versepar \skt{ऋषयः }\edtext{\skt{सौषधिग्रामा}}{\Afootnote{N H; \skt{सौषधीग्रामा}\ A1938}}\skt{ भूतसंघाश्च पान्तु ते ॥}\pend

\pstart\versepar \ledsidenote{\passref{1.43.3.2}}\edtext{\skt{रसायनेव}}{\Afootnote{H; \skt{रसायनमि}\ A1938; \skt{रसायनेन}\ N}}\skt{ }\edtext{\skt{सिद्धानान्देवानाममृतं}}{\lemma{\skt{सिद्धानान्}}\Afootnote{N H; \skt{वर्षीणां}\ A1938}}\skt{ यथा । }\pend
\pstart\versepar \skt{}\edtext{\skt{सुधेवोत्तमनागानास्भैषज्यमिदमस्तु}}{\lemma{\skt{सुधेवोत्तमनागानास्भैषज्यम्}}\Afootnote{H; \skt{सुधेवोत्तमनागानां भैषज्यम्}\ A1938 N\textsuperscript{ac}; \skt{सु दवोत्तमनागानां भैषज्यम्}\ N\textsuperscript{pc}}}\skt{ ते ॥}\pend

\pstart \ledsidenote{\passref{1.43.3.part2}}\skt{विशेषेण }\edtext{\skt{श्लेष्मज्वरप्रतिश्यायान्तर्व्विद्रधिष्वप्रवर्त्तमाने}}{\lemma{\skt{अप्रवर्त्तमाने}}\Afootnote{N H; \skt{अप्रवर्तमाने वा}\ A1938}}\skt{ दोषे }\edtext{\edtext{\skt{पिप्पलीवचासर्षपकल्कोन्मिश्रैरुष्णाम्वुभिः}}{\lemma{\skt{पिप्पलीवचासर्षपकल्कोन्मिश्रैर्}}\Afootnote{N H; \skt{पिप्पलीवचागौरसर्षपकल्कोन्मिश्रैः}\ A1938}}}{\lemma{\skt{उष्णाम्वुभिः}}\Afootnote{H; \skt{सलवणैरुष्णाम्बुभिः}\ A1938; \skt{उष्णाम्बुभिः}\ N}}\skt{ पुनः पुनः प्रवर्त्तयेदासम्यग्वान्तलक्षणादिति । मदनफलमज्जचूर्ण्णम्वा }\edtext{\skt{तत्क्वाथपरिभावितम्वसेन}}{\lemma{\skt{वसेन}}\Afootnote{H; \skt{मदन}\ A1938; \skt{वमन}\ N}}\skt{ }\edtext{\skt{द्रव्यकषायेण}}{\Afootnote{N H; \skt{फलकषायेण;}\ A1938}}\skt{ । }\edtext{\skt{मदनफलमज्जा}}{\Afootnote{H; \skt{मदनफलमज्जसिद्धस्य वा पयसः सन्तानिकां क्षौद्रयुक्तां , मदनफलमज्ज}\ A1938; \skt{मदनफलमज्ज}\ N}}\skt{ सिद्धम्वा पयः । मदनफलमज्जसिद्धेन वा }\edtext{\skt{पयसा वा}}{\Afootnote{N H; \skt{पयसा}\ A1938}}\skt{ }\edtext{\skt{यवागूमधोभागासृहृद्रोगयोः}}{\lemma{\skt{अधोभागासृहृद्रोगयोः}}\Afootnote{H; \skt{अधोभागासृक्पित्तहृद्दाहयोः;}\ A1938; \skt{अधोभागासृक्पित्तहृद्रोगयोः ।}\ N}}\skt{ । मदनफलमज्जसिद्धस्य वा पयसो दधिभावमुपगतस्य । }\edtext{\skt{दध्युत्तरकं}}{\Afootnote{N H; \skt{दध्युत्तरं}\ A1938}}\skt{ दधि वा }\edtext{\skt{कफप्रसेकशर्द्दिस्तमेषु}}{\Afootnote{H; \skt{कफप्रसेकच्छर्दिमूर्च्छातमकेषु;}\ A1938; \skt{कफप्रसेकश्छर्दिस्तमकेषु ।}\ N}}\skt{ । }\edtext{\skt{मदनफलमज्जाचूर्ण्ण}}{\Afootnote{H; \skt{मदनफलमज्जरसं}\ A1938; \skt{मदनफलमज्जचूर्ण्णं}\ N}}\skt{ }\edtext{\skt{रसम्वा}}{\Afootnote{N H; \skt{स्नेहं वा}\ A1938}}\skt{ }\edtext{\skt{भल्लातकसेन्हवदादाय}}{\Afootnote{H; \skt{भल्लातकस्नेहवदादाय}\ A1938; \skt{भल्लातकस्नेहवदादाय}\ N}}\skt{ फाणितीभूतं }\edtext{\skt{लेहयेत् । आतपशुष्कम्वा}}{\Afootnote{H; \skt{लेहयेत्, आतपपरिशुष्कं वा}\ A1938; \skt{लेहयेत् । आतप}\ N}}\skt{ }\edtext{\skt{जीवन्ती}}{\Afootnote{H; \skt{तमेव जीवन्ती}\ A1938; \skt{सुष्कम्वा जीवन्ती}\ N}}\skt{ कषायेण }\edtext{\skt{पित्तकफस्थानगते}}{\Afootnote{H; \skt{पित्ते कफस्थानगते;}\ A1938; \skt{पित्ते स्थानगते ।}\ N}}\skt{ । मदनफल }\edtext{\skt{मज्जाक्वाथम्वा}}{\Afootnote{H; \skt{मज्जक्वाथं वा}\ A1938 N}}\skt{ पिप्पल्यादि }\edtext{\skt{प्रतिवापं}}{\Afootnote{N H; \skt{प्रतीवापं,}\ A1938}}\skt{ }\edtext{\skt{ऊषणनिम्वकषाय }\edtext{\skt{योरन्यतमेन}}{\lemma{\skt{अन्यतमेन}}\Afootnote{N H; \skt{अन्यतरेण}\ A1938}}}{\lemma{\skt{ऊषणनिम्वकषाय योर्}}\Afootnote{H; \skt{तच्चूर्णं वा निम्बरूपिकाकषाययोर्}\ A1938; \skt{ऊषणनिम्बकषाययोर्}\ N}}\skt{ }\edtext{\skt{सन्तर्प्पणं}}{\Afootnote{N H; \skt{सन्तर्पण}\ A1938}}\skt{ । }\edtext{\edtext{\skt{कफसर्व्वव्याधिहरमिति}}{\lemma{\skt{कफसर्व्वव्याधिहरम्}}\Afootnote{N H; \skt{कफजव्याधिहरं,}\ A1938}}}{\lemma{\skt{इति}}\Afootnote{N H; \skt{मदनफलमज्जचूर्णं वा मधूककाश्मर्यद्राक्षाकषायेण ।}\ A1938}}\skt{ }\edtext{\skt{मदनफलमुक्तं}}{\lemma{\skt{मदनफलम्}}\Afootnote{N H; \skt{मदनफलविधानम्}\ A1938}}\skt{ ।}\pend

\pstart \ledsidenote{\passref{1.43.4}}\edtext{\skt{जीमूतककुसुमचूर्ण्णम्वा}}{\Afootnote{H; \skt{जीमूतककुसुमचूर्णं}\ A1938; \skt{जीमूतककुस्भुमचूर्ण्णं वा}\ N}}\skt{ }\edtext{\skt{पूर्व्ववदेव}}{\Afootnote{A1938 H; \skt{पूर्वावदेव}\ N}}\skt{ क्षीरेण }\edtext{\skt{निवृत्तेषु}}{\Afootnote{N H; \skt{निर्वृत्तेषु}\ A1938}}\skt{ }\edtext{\skt{यवागूं}}{\Afootnote{N H; \skt{क्षीरयवागूं,}\ A1938}}\skt{ }\edtext{\skt{रोमसेषु}}{\Afootnote{N H; \skt{रोमशेषु}\ A1938}}\skt{ । }\edtext{\skt{सन्तानिकाम्वा}}{\Afootnote{N H; \skt{सन्तानिकां,}\ A1938}}\skt{ }\edtext{\skt{निर्ल्लोमशेषु}}{\Afootnote{H; \skt{अरोमशेषु}\ A1938; \skt{निर्लोमसेषु}\ N}}\skt{ । }\edtext{\skt{दध्युत्तरकं}}{\Afootnote{H; \skt{दध्युत्तरं,}\ A1938; \skt{दध्युत्तरक}\ N}}\skt{ }\edtext{\skt{हरितकपाण्डुषु}}{\Afootnote{N H; \skt{हरितपाण्डुषु}\ A1938}}\skt{ । दधि }\edtext{\skt{तत्कषायसंसृटां}}{\Afootnote{H; \skt{तत्कषायसंसृष्टां वा सुरां}\ A1938; \skt{तत्कषायं संसृष्टां}\ N}}\skt{ }\edtext{\skt{सुराम्वा}}{\Afootnote{H; \skt{कफारोचककासश्वासपाण्डुरोग}\ A1938; \emph{om.}\ N}}\skt{ }\edtext{\skt{पर्यागतेषु}}{\Afootnote{H; \skt{यक्ष्मसु; पर्यागतेषु}\ A1938; \skt{सुराम्वा पर्यागतेषु ।}\ N}}\skt{ । }\edtext{\skt{मदनफलमज्ज}}{\Afootnote{N H; \skt{मदनफलमज्जवदुपयोगः ॥४॥}\ A1938}}\pend

\pstart \ledsidenote{\passref{1.43.5}}\edtext{\skt{विधानवत्कुटजफलमज्जविधानं}}{\Afootnote{H; \skt{तद्वदेव कुटजफलविधानम् ॥५॥}\ A1938; \skt{विधानवत्कूटजफलमज्जविधानं ॥}\ N}}\skt{ ।}\pend

\pstart \ledsidenote{\passref{1.43.7}}\skt{इक्ष्वाकुकुसुमचूर्ण्णम्वा }\edtext{\skt{पूर्व्ववदेव}}{\Afootnote{N H; \skt{पूर्ववत् , एवं}\ A1938}}\skt{ }\edtext{\skt{क्षीरेण}}{\Afootnote{N H; \skt{क्षीरेण, कासश्वासच्छर्दिकफरोगेषूपयोगः ॥७॥}\ A1938}}\skt{ ।}\pend

\pstart \ledsidenote{\passref{1.43.8}}\skt{धामार्गवस्यापि }\edtext{\skt{मदनफलमज्जा}}{\Afootnote{H; \skt{मदनफलमज्ज}\ A1938 N}}\skt{ }\edtext{\skt{विधानवदुपयोगः}}{\Afootnote{N H; \skt{वदुपयोगः; विशेषतस्तु गरगुल्मोदरकासश्वासश्लेष्मामयेषु वायौ च कफस्थानगते ॥८॥}\ A1938}}\skt{ ॥}\pend

\pstart \ledsidenote{\passref{1.43.9}}\edtext{\edtext{\edtext{\edtext{\edtext{\skt{कृतवेधनमदनफलमज्जपिप्पलीमाम्वमनद्रव्यपरिभावितानाम्बहुशश्चूर्ण्णमुत्पलादिषु}}{\lemma{\skt{कृतवेधनमदनफलमज्जपिप्पलीमाम्}}\Afootnote{H; \skt{कृतवेधनफलपिप्पलीनां}\ A1938; \skt{कृतवेधनमदनफलमज्जपिप्पलीनां}\ N}}}{\lemma{\skt{वमनद्रव्यपरिभावितानाम्}}\Afootnote{N H; \skt{वमनद्रव्यकषायपरिपीतानां}\ A1938}}}{\lemma{\skt{बहुशश्}}\Afootnote{\skt{बहुश}\ A1938; \skt{बहुसश्}\ N; \skt{वहुशश्}\ H}}}{\lemma{\skt{चूर्ण्णम्}}\Afootnote{N H; \skt{श्चूर्ण}\ A1938}}}{\lemma{\skt{उत्पलादिषु}}\Afootnote{H; \skt{मुत्पलादिषु}\ A1938; \skt{उत्पलादि}\ N}}\skt{ }\edtext{\edtext{\skt{दत्तमाघ्रा}}{\lemma{\skt{दत्तम्}}\Afootnote{N H; \skt{दत्त}\ A1938}}\skt{ य}}{\lemma{\skt{आघ्रा य}}\Afootnote{N H; \skt{माघ्रातं}\ A1938}}\skt{ }\edtext{\skt{वसति}}{\Afootnote{N H; \skt{वामयति,}\ A1938}}\skt{ }\edtext{\skt{तत्त्वनववद्धदोषेषु}}{\Afootnote{N H; \skt{तत्त्वनवबद्धदोषेषु}\ A1938}}\skt{ }\edtext{\edtext{\skt{यवागूमाकण्ठम्पीतवत्सु}}{\lemma{\skt{यवागूमाकण्ठम्}}\Afootnote{H; \skt{यवागूमाकण्ठा}\ A1938; \skt{यवागूमाकण्ठं}\ N}}\skt{ च}}{\lemma{\skt{पीतवत्सु च}}\Afootnote{N H; \skt{त्पीतवत्सु च}\ A1938}}\skt{ }\edtext{\skt{दद्यादिति}}{\Afootnote{N H; \skt{विदध्यात् । वमनविरेचन}\ A1938}}\skt{ । }\edtext{\edtext{\skt{शिरोविरेचनान्येवं}}{\lemma{\skt{शिरोविरेचनान्य्}}\Afootnote{H; \skt{शिरोविरोचनद्रव्या}\ A1938; \skt{सिरोविरेचनान्य्}\ N}}}{\lemma{\skt{एवं}}\Afootnote{N H; \skt{ण्येवं वा}\ A1938}}\skt{ प्रधानतमानि भवन्ति ॥}\pend

\pstart\versepar \ledsidenote{\passref{1.43.11}}\edtext{\skt{भवति}}{\Afootnote{H; \emph{om.}\ A1938 N}}\skt{ }\edtext{\skt{चात्र}}{\Afootnote{H; \emph{om.}\ A1938 N}}\skt{ }\pend
\pstart\versepar \skt{}\edtext{\skt{कल्कैः}}{\Afootnote{N H; \emph{om.}\ A1938}}\skt{ कषायैः }\edtext{\skt{स्वरसैश्च}}{\Afootnote{H; \skt{स्वरसैः कल्कैश्}\ A1938; \skt{स्वरसैश्}\ N}}\skt{ चूर्ण्णैरपि च वुद्धिमान् । }\pend
\pstart\versepar \skt{पेयलेह्यादिभोज्येषु }\edtext{\skt{वमनान्यूपकल्पयेत्}}{\Afootnote{H; \skt{वमनान्युपकल्पयेत् ॥११॥}\ A1938; \skt{वमनान्युपकल्पयेदिति ॥}\ N}}\skt{ ॥}\edtext{}{\lemma{\skt{वुद्धिमान् । पेयलेह्यादिभोज्येषु}}\Afootnote{H; \skt{बुद्धिमान् । पेयलेह्याद्यभोज्येषु}\ A1938; \skt{बुद्धिमान् । पेयलेह्यादिभोज्येषु}\ N}}\pend

\pstart \ledsidenote{\passref{1.43.trailer}}\edtext{\skt{॥ ४३ ॥}}{\Afootnote{\skt{इति सुश्रुतसंहितायां सूत्रस्थाने वमनद्रव्यविकल्पविज्ञानीयो नाम त्रिचत्वारिंशत्तमोऽध्यायः ॥४३॥}\ A1938; \skt{४३ ॥ {\symfont ❈} ॥}\ N; \skt{त्रि ॥ ॥}\ H}}\pend

\endnumbering
\beginnumbering
\eledsection*{Adhyāya 44}
\pstart \ledsidenote{\passref{1.44.1}}\skt{अथातो }\edtext{\skt{विरेचनद्रव्यविकल्पविज्ञानीयम्व्याख्यास्यामः}}{\lemma{\skt{व्याख्यास्यामः}}\Afootnote{H; \skt{अध्यायं व्याख्यास्यामः ॥}\ A1938; \emph{om.}\ N}}\skt{ ॥}\pend

\pstart\versepar \ledsidenote{\passref{1.44.3}}\edtext{\skt{अरुणाभन्तृवृन्मूलं}}{\lemma{\skt{तृवृन्मूलं}}\Afootnote{H; \skt{त्रिवृन्मूलं}\ A1938; \skt{तृवृन्मूलं}\ N}}\skt{ श्रेष्ठं मूलविरेचने । }\pend
\pstart\versepar \skt{}\edtext{\edtext{\edtext{\skt{प्रधानस्तिल्वकस्त्वक्षु}}{\lemma{\skt{प्रधानस्}}\Afootnote{H; \skt{प्रधानं}\ A1938; \skt{प्रधानन}\ N}}}{\lemma{\skt{तिल्वकस्}}\Afootnote{\skt{तिल्वक}\ A1938 H; \skt{स्तिल्वकस्}\ N}}}{\lemma{\skt{त्वक्षु}}\Afootnote{N; \skt{स्त्वक्षु}\ A1938 H}}\skt{ फलेष्वपि हरीतकी ॥}\pend

\pstart\versepar \ledsidenote{\passref{1.44.4}}\edtext{\skt{स्नेहापयः}}{\Afootnote{H; \skt{तैलेष्वेरण्डजं तैलं स्वरसे कारवेल्लिका । सुधापयः}\ A1938; \skt{स्नुहापयः}\ N}}\skt{ पयस्सूक्तमिति प्राधान्यसङ्ग्रहः । }\pend
\pstart\versepar \skt{तेषाम्विधानं वक्ष्यामि }\edtext{\skt{यथावदनुपूर्व्वशः}}{\Afootnote{H; \skt{यथावदनुपूर्वशः ॥४॥}\ A1938; \skt{यथावदनुपूर्व्वसः ॥}\ N}}\skt{ ॥}\pend

\pstart\versepar \ledsidenote{\passref{1.44.5}}\edtext{\skt{वैरेचनद्रव्यकषायपीतम्मूलम्}}{\lemma{\skt{वैरेचनद्रव्यकषायपीतम्}}\Afootnote{N H; \skt{वैरेचनद्रव्यरसानुपीतं}\ A1938}}\skt{ }\pend
\pstart\versepar \skt{}\edtext{\skt{महत्त्रैवृतमत्र}}{\Afootnote{H; \skt{महत्त्रैवृतमस्त}\ A1938; \skt{महत्तैवृतमत्र}\ N}}\skt{ }\edtext{\skt{शुद्धं}}{\Afootnote{N H; \skt{दोषम् ।}\ A1938}}\skt{ । }\pend
\pstart\versepar \skt{}\edtext{\skt{चूर्ण्णीकृतं}}{\Afootnote{A1938 H; \skt{चूर्ण्णीकृत}\ N}}\skt{ }\edtext{\skt{सैन्धवनागराढ्यम्}}{\Afootnote{N H; \skt{सैन्धवनागराढ्य}\ A1938}}\skt{ }\pend
\pstart\versepar \skt{}\edtext{\edtext{\skt{पिबेत्तदम्लैरनिलाधिजुष्टः}}{\lemma{\skt{पिबेत्तदम्लैर्}}\Afootnote{\skt{मम्लैः पिबेन्मारुत}\ A1938; \skt{पिबेतस्च्लै}\ N; \skt{पिबेत्तदम्लैर्}\ H}}}{\lemma{\skt{अनिलाधिजुष्टः}}\Afootnote{H; \skt{रोगजुष्टः ॥५॥}\ A1938; \skt{रनिलादिजुष्तः ।}\ N}}\skt{ ॥}\pend

\pstart\versepar \ledsidenote{\passref{1.44.6}}\edtext{\skt{स्वादुक्वाथैरपि चेक्षोर्व्विकारैः}}{\lemma{\skt{स्वादुक्वाथैरपि चेक्षोर्}}\Afootnote{\skt{इक्षोर्}\ A1938; \skt{स्वादुत्काथैरपि चेक्षो}\ N; \skt{स्वादुक्वाथैरपिचेक्षोर्}\ H}}\skt{ }\edtext{\skt{पैत्तेरोगे}}{\Afootnote{H; \skt{मधुरै रसैस्तत्पैत्ते गदे}\ A1938; \skt{पैत्तरोगै}\ N}}\skt{ }\edtext{\skt{क्षीरयुक्तान्निहन्यात्}}{\lemma{\skt{क्षीरयुक्तान्}}\Afootnote{H; \skt{क्षी}\ A1938; \skt{क्षीरयुक्ता}\ N}}\skt{ । }\pend
\pstart\versepar \skt{}\edtext{\edtext{\edtext{\skt{द्राक्षारसत्रिफलाक्वाथमूत्रैर्युक्तम्पिबेत्कफजव्योषगाढं}}{\lemma{\skt{सत्रिफलाक्वाथमूत्रैर्यु}}\Afootnote{\emph{om.}\ A1938}}}{\lemma{\skt{युक्तम्}}\Afootnote{H; \skt{युतं}\ A1938; \skt{युक्तां}\ N}}}{\lemma{\skt{पिबेत्कफजव्योषगाढं}}\Afootnote{H\textsuperscript{pc}; \skt{पिबेच्च गुडूच्यरिष्टत्रिफलारसेन सव्योषमूत्रं कफजे पिबेत्तत् ॥६॥}\ A1938; \skt{पिबेत्कफजेध्योषगढां ।}\ N; \skt{पिबेत्कफव्योषगाढं ॥}\ H\textsuperscript{ac}}}\skt{ ॥}\edtext{}{\lemma{\skt{रयुक्तान्… द्राक्षार}}\Afootnote{\emph{om.}\ A1938}}\edtext{}{\lemma{\skt{निहन्यात् । द्राक्षारसत्रिफलाक्वाथमूत्रैर्}}\Afootnote{H; \skt{निहन्यात् । द्राक्षरसत्रिफलात्काथमूत्रैर्}\ N}}\pend

\pstart\versepar \ledsidenote{\passref{1.44.7}}\edtext{\skt{त्रिवर्ण्णकत्रिकटुकाभ्यान्युक्तं}}{\Afootnote{H; \skt{त्रिवर्णकत्र्यूषणयुक्तम्}\ A1938; \skt{त्रिवर्ण्णकत्रिदुकाख्यायुक्तं}\ N}}\skt{ }\edtext{\edtext{\skt{लिह्याच्चूर्ण्णन्तद्गुडेनाभियोज्यं}}{\lemma{\skt{लिह्याच्चूर्ण्णन्}}\Afootnote{N H; \skt{ए}\ A1938}}}{\lemma{\skt{तद्गुडेनाभियोज्यं}}\Afootnote{H; \skt{तद्गुडेन लिह्याद}\ A1938; \skt{तद्गुडेनाभियोज्यं ।}\ N}}\skt{ । }\pend
\pstart\versepar \skt{}\edtext{\skt{क्वाथप्रस्थे}}{\Afootnote{H; \skt{नवेन चूर्णम् । प्रस्थे च}\ A1938; \skt{थ प्रस्थे}\ N}}\skt{ }\edtext{\edtext{\skt{कुडवन्तस्य}}{\lemma{\skt{कुडवन्}}\Afootnote{N H; \skt{तन्मूलर}\ A1938}}}{\lemma{\skt{तस्य}}\Afootnote{N H; \skt{सस्य}\ A1938}}\skt{ }\edtext{\skt{दद्यात्युक्त्या}}{\Afootnote{H; \skt{दत्वा तन्मू}\ A1938; \skt{दद्याद्यु}\ N}}\skt{ }\edtext{\skt{दद्यान्नागरं}}{\lemma{\skt{दद्यान्}}\Afootnote{H; \skt{ल}\ A1938; \skt{दन्}\ N}}\skt{ }\edtext{\edtext{\skt{सैन्धवञ्च}}{\lemma{\skt{च}}\Afootnote{N H; \skt{कल्कं कुडवप्रमाणम् ॥७॥}\ A1938}}\skt{ ॥}}{\Afootnote{\emph{om.}\ A1938}}\pend

\pstart\versepar \ledsidenote{\passref{1.44.8}}\skt{पचेत्सर्व्वं }\edtext{\skt{यावदेतद्घनं}}{\Afootnote{H; \skt{यावदेव तघनं}\ N}}\skt{ }\edtext{\skt{स्याल्लेहीभूतन्तत्प्रयोज्यं}}{\lemma{\skt{तत्प्रयोज्यं}}\Afootnote{H; \skt{तत्प्रयोज्य}\ N}}\skt{ }\edtext{\skt{ततस्तु}}{\Afootnote{H; \skt{त्ततस्तु}\ N}}\skt{ । }\pend
\pstart\versepar \skt{}\edtext{\skt{तृवृत्कल्कोनागरार्द्धेन}}{\Afootnote{H; \skt{कर्षोन्मिते सैन्धवनागरे च विपाच्य कल्कीकृतमेतदद्यात् । तत्कल्कभागः समहौष}\ A1938; \skt{कर्षोन्मिते सैन्धनागरे च विपाच्य कल्कीकृतमेतदद्यात् । तृवृत्कल्कानागरार्द्धेन}\ N}}\skt{ }\edtext{\skt{युक्तः}}{\Afootnote{N H; \skt{धार्धः}\ A1938}}\skt{ ससैन्धवो }\edtext{\skt{मूत्रयुक्तस्तु}}{\Afootnote{N H; \skt{मूत्रयुतश्च}\ A1938}}\skt{ पेयः ॥}\edtext{}{\lemma{\skt{… तु}}\Afootnote{\emph{om.}\ A1938}}\pend

\pstart\versepar \ledsidenote{\passref{1.44.9}}\edtext{\edtext{\skt{समेतृवृन्नागरे}}{\lemma{\skt{समेतृवृन्}}\Afootnote{H; \skt{समास्त्रिवृ}\ A1938; \skt{समेतृवृ}\ N}}}{\lemma{\skt{नागरे}}\Afootnote{N H; \skt{न्नागर}\ A1938}}\skt{ }\edtext{\skt{चाभयाञ्च}}{\Afootnote{N H; \skt{काभयाः स्यु}\ A1938}}\skt{ }\edtext{\edtext{\skt{दद्यादर्द्धम्पूगफलं}}{\lemma{\skt{दद्यादर्द्धम्}}\Afootnote{H; \skt{र्भागार्धकं}\ A1938; \skt{दद्यादर्ध}\ N}}}{\lemma{\skt{पूगफलं}}\Afootnote{A1938 H; \skt{पूराफलं}\ N}}\skt{ }\edtext{\skt{सदारुः}}{\Afootnote{N H; \skt{सुपक्वम् ।}\ A1938}}\skt{ । }\pend
\pstart\versepar \skt{}\edtext{\skt{विडंगसारं}}{\Afootnote{H; \skt{विडङ्गसारो}\ A1938; \skt{विडंगसार}\ N}}\skt{ }\edtext{\skt{मरिच}}{\Afootnote{N H; \skt{मरिचं स}\ A1938}}\skt{ }\edtext{\skt{वैषयोगः}}{\Afootnote{N H; \skt{दारु योगः}\ A1938}}\skt{ }\edtext{\skt{ससैन्धवो}}{\Afootnote{N H; \skt{ससिन्धूद्भव}\ A1938}}\skt{ मूत्रयुक्तः }\edtext{\skt{प्रधानः}}{\Afootnote{N H; \emph{om.}\ A1938}}\skt{ ॥}\pend

\pstart\versepar \ledsidenote{\passref{1.44.12}}\edtext{\skt{वैरेचनद्रव्यचूर्ण्णस्य}}{\Afootnote{N H; \skt{वैरेकीयद्रव्यचूर्णस्य}\ A1938}}\skt{ }\edtext{\edtext{\skt{भागस्ततः}}{\lemma{\skt{भागस्}}\Afootnote{H; \skt{भा}\ A1938; \skt{द्वाभिः ।}\ N}}}{\lemma{\skt{ततः}}\Afootnote{N H; \emph{om.}\ A1938}}\skt{ }\edtext{\skt{क्वाथः}}{\Afootnote{H; \emph{om.}\ A1938; \skt{त्काथः}\ N}}\skt{ }\edtext{\skt{सम्मितञ्चास्य}}{\Afootnote{H; \skt{गं सिद्धं}\ A1938; \skt{सस्मितञ्चास्य}\ N}}\skt{ }\edtext{\skt{तुल्यं}}{\Afootnote{H; \skt{सार्धं क्वाथभागैश्चतुर्भिः ।}\ A1938; \skt{तुल्पं । समर्द्दितं स}\ N}}\skt{ । }\pend
\pstart\versepar \skt{}\edtext{\skt{सुमर्द्दितं}}{\Afootnote{H; \skt{आमृद्नीया}\ A1938; \emph{om.}\ N}}\skt{ }\edtext{\skt{सर्प्पिषा}}{\Afootnote{H; \skt{त्सर्पिषा}\ A1938; \skt{र्प्पिषा}\ N}}\skt{ }\edtext{\skt{तच्छ्रितेन}}{\Afootnote{N H; \skt{तच्छृतेन}\ A1938}}\skt{ }\edtext{\skt{तत्क्वाथोश्मस्वेदिता}}{\Afootnote{H; \skt{तत्क्वाथोष्मस्वेदितं}\ A1938; \skt{तत्क्वाथोस्य स्वेदिता}\ N}}\skt{ }\edtext{\skt{वर्त्तितञ्च}}{\lemma{\skt{वर्त्तितञ्}}\Afootnote{N H; \skt{सामितं}\ A1938}}\skt{ ॥}\pend

\pstart\versepar \ledsidenote{\passref{1.44.13}}\edtext{\skt{पाकप्राप्ते}}{\Afootnote{A1938 H; \skt{प्राकप्राप्तै}\ N}}\skt{ }\edtext{\skt{फाणितं}}{\Afootnote{N H; \skt{फाणिते}\ A1938}}\skt{ }\edtext{\skt{चूर्ण्णिताक्तं}}{\Afootnote{N H; \skt{चूर्णितं त}\ A1938}}\skt{ }\edtext{\skt{क्षिप्त्वा}}{\Afootnote{N H; \skt{त्क्षिप्तं}\ A1938}}\skt{ पक्वं }\edtext{\skt{चावतार्या}}{\Afootnote{N H; \skt{चावतार्य}\ A1938}}\skt{ प्रमादात् । }\pend
\pstart\versepar \skt{शीतीकृत्वा }\edtext{\skt{मोदकं}}{\Afootnote{H; \skt{मोदका}\ A1938; \skt{मोदकां}\ N}}\skt{ }\edtext{\edtext{\edtext{\skt{सौरभाज्यङ्कुर्यादेवम्भक्ष्यकल्पा}}{\lemma{\skt{सौरभाज्यङ्}}\Afootnote{H; \skt{हृद्यगन्धाः}\ A1938; \skt{सौरभाढ्यं}\ N}}}{\lemma{\skt{कुर्यादेवम्}}\Afootnote{N H; \skt{कार्यास्त्वेते}\ A1938}}}{\lemma{\skt{भक्ष्यकल्पा}}\Afootnote{H\textsuperscript{pc}; \skt{भक्ष्यकल्पाः}\ A1938; \skt{भक्षकल्पं}\ N; \skt{भश्ठक्ष्यकल्पा}\ H\textsuperscript{ac}}}\skt{ }\edtext{\skt{समासान्}}{\Afootnote{H; \skt{समासात् ॥१३॥}\ A1938; \skt{समांसं ।}\ N}}\skt{ ।}\edtext{}{\lemma{\skt{प्रमादात् । शीतीकृत्वा}}\Afootnote{H; \skt{प्रयत्नात् । शीतीभूता}\ A1938; \skt{प्रमादात् । सीतीकृउत्वा}\ N}}\pend

\pstart\versepar \ledsidenote{\passref{1.44.14}}\edtext{\edtext{\skt{मुद्गानाम्वातैः}}{\lemma{\skt{मुद्गानाम्}}\Afootnote{H; \skt{रसेन तेषां परिभाव्य मुद्गा}\ A1938; \skt{मुद्गाना}\ N}}\skt{ समं }\edtext{\skt{साधितानां}}{\Afootnote{N\textsuperscript{pc} H; \skt{साधितांनां}\ N\textsuperscript{ac}}}\skt{ }\edtext{\skt{यूषो}}{\Afootnote{N H; \skt{न्यू}\ A1938}}\skt{ }\edtext{\skt{हृद्यः}}{\Afootnote{N H; \skt{षः}\ A1938}}}{\lemma{\skt{नाम्… यू}}\Afootnote{\emph{om.}\ A1938}}\skt{ }\edtext{\skt{सधृतः}}{\Afootnote{H; \skt{ससिन्धू}\ A1938; \skt{सघृतः}\ N}}\skt{ }\edtext{\skt{सैन्धवाढ्यः}}{\Afootnote{H; \skt{द्भवसर्पिरिष्टः ।}\ A1938; \skt{सैन्ध}\ N}}\skt{ ॥ }\pend
\pstart\versepar \skt{}\edtext{\skt{विरेचयेद्वैदलेष्वेषु}}{\Afootnote{H; \skt{वैरेचनेऽन्यैरपि वैदलैः स्यादेवं विदध्याद्वमनौषधैश्च ॥१४॥}\ A1938; \skt{वाढ्यः । विरेच द्वैदलेष}\ N}}\skt{ }\edtext{\skt{कल्पः}}{\Afootnote{N H; \emph{om.}\ A1938}}\skt{ }\edtext{\edtext{\skt{कार्यास्तज्ञैर्व्वामनीयेषु}}{\lemma{\skt{कार्यास्तज्ञैर्}}\Afootnote{H; \emph{om.}\ A1938; \skt{कायास्तज्ञैर्}\ N}}}{\lemma{\skt{व्वामनीयेषु}}\Afootnote{N H; \emph{om.}\ A1938}}\skt{ }\edtext{\skt{चैव}}{\Afootnote{N H; \emph{om.}\ A1938}}\skt{ ॥}\pend

\pstart\versepar \ledsidenote{\passref{1.44.15}}\edtext{\edtext{\skt{इक्षुर्द्विधा}}{\lemma{\skt{इक्षुर्}}\Afootnote{\emph{om.}\ A1938; \skt{इक्ष्य}\ N; \skt{इक्षु}\ H}}}{\lemma{\skt{द्विधा}}\Afootnote{N; \emph{om.}\ A1938; \skt{र्द्विधा}\ H}}\skt{ }\edtext{\skt{पाटयित्वावलिप्य}}{\Afootnote{N H; \skt{भित्त्वा द्विधेक्षुं परिलिप्य}\ A1938}}\skt{ }\edtext{\skt{तृवृत्कल्कैः}}{\Afootnote{N H; \skt{कल्कैस्त्रिभण्डिजातैः}\ A1938}}\skt{ }\edtext{\skt{प्रविसन्धाय}}{\Afootnote{H; \skt{प्रतिबध्य र}\ A1938; \skt{प्रविसत्वाय}\ N}}\skt{ }\edtext{\skt{चापि}}{\Afootnote{N H; \skt{ज्ज्वा । प}\ A1938}}\skt{ । }\pend
\pstart\versepar \skt{}\edtext{\skt{पक्त्वा}}{\Afootnote{H; \skt{क्वं च}\ A1938; \skt{पक्वा}\ N}}\skt{ }\edtext{\skt{सम्यक्पुटपाकं}}{\Afootnote{H; \skt{सम्यक्पुटपा}\ A1938; \skt{सम्य-क्पूटपाकं}\ N}}\skt{ }\edtext{\skt{क्रमेण}}{\Afootnote{N H; \skt{कयुक्त्या}\ A1938}}\skt{ }\edtext{\skt{खादेच्छीतं}}{\Afootnote{N H; \skt{खादेत्तु तं}\ A1938}}\skt{ }\edtext{\skt{पित्तरोगाभिभूतः}}{\Afootnote{N H; \skt{पित्तगदी सुशीतम् ॥१५॥}\ A1938}}\skt{ ॥}\pend

\pstart\versepar \ledsidenote{\passref{1.44.28}}\skt{वैरेचनिकनिःक्वाथ भागाः }\edtext{\edtext{\skt{शीतास्त्रयो}}{\lemma{\skt{शीतास्}}\Afootnote{H; \skt{शीता}\ A1938; \skt{सीतास्}\ N}}}{\lemma{\skt{त्रयो}}\Afootnote{N H; \skt{स्त्रयो}\ A1938}}\skt{ मिताः । }\pend
\pstart\versepar \skt{द्वौ }\edtext{\skt{फाणितस्य तां}}{\Afootnote{N H; \skt{फाणितस्य त}\ A1938}}\skt{ }\edtext{\skt{सर्व्वा}}{\Afootnote{H; \skt{च्चापि}\ A1938; \skt{मर्चाम्}\ N}}\skt{ }\edtext{\skt{न्यूनरग्नावधिश्रयेत्}}{\Afootnote{H; \skt{पुनरग्नावधिश्रयेत् ॥२८॥}\ A1938; \skt{पुरग्नावधिश्रयेत् ।}\ N}}\skt{ ॥}\edtext{}{\lemma{\skt{मिताः । द्वौ}}\Afootnote{N H; \skt{मताः । द्वौ}\ A1938}}\pend

\pstart\versepar \ledsidenote{\passref{1.44.29}}\skt{तत्साधुसिद्धम्विज्ञाय }\edtext{\skt{शीती}}{\Afootnote{N H; \skt{शीतं}\ A1938}}\skt{ कृत्वा निधापयेत् । }\pend
\pstart\versepar \skt{कलशे कृतसंस्कारे }\edtext{\skt{विभज्यन्तु}}{\Afootnote{H; \skt{विभज्य र्तू}\ A1938; \skt{विभज्य तू}\ N}}\skt{ }\edtext{\skt{हिमाहिमौ}}{\Afootnote{A1938 H; \skt{हिसाहिसौ ।}\ N}}\skt{ ॥}\edtext{}{\lemma{\skt{निधापयेत् । कलशे}}\Afootnote{H; \skt{निधापयेत्कलसे}\ A1938 N}}\pend

\pstart\versepar \ledsidenote{\passref{1.44.30}}\edtext{\skt{ऊर्द्ध्वञ्जातरसं}}{\lemma{\skt{ऊर्द्ध्वञ्}}\Afootnote{N H; \skt{मासादूर्ध्वं}\ A1938}}\skt{ }\edtext{\skt{मासादासवं}}{\Afootnote{H; \emph{om.}\ A1938; \skt{मासादासवं}\ N}}\skt{ }\edtext{\skt{मधुगन्धिकं}}{\Afootnote{N H; \skt{मधुगन्धं व}\ A1938}}\skt{ । }\pend
\pstart\versepar \skt{}\edtext{\skt{मात्रया}}{\Afootnote{H; \skt{रासवम् । पिबेदसावेव विधिः क्षारमूत्रासवेष्व्}\ A1938; \skt{मार्त्रयो}\ N}}\skt{ }\edtext{\skt{प्रपिबेत्प्रातः}}{\Afootnote{H; \skt{अपि ॥३०॥}\ A1938; \skt{प्रपिचेत्प्रातस्}\ N}}\skt{ }\edtext{\skt{ततः}}{\Afootnote{N H; \emph{om.}\ A1938}}\skt{ }\edtext{\skt{सम्यग्विरिच्यते}}{\Afootnote{N H; \emph{om.}\ A1938}}\skt{ ॥}\pend

\pstart\versepar \ledsidenote{\passref{1.44.31}}\edtext{\skt{वैरेचनिकमूलानां}}{\Afootnote{A1938 H; \skt{वैरेचनिकमूलाणां}\ N}}\skt{ क्वाथे मासान्सुसाधितान् । }\pend
\pstart\versepar \skt{सुधौतं तत्कषायेण शालीनाञ्चापि }\edtext{\skt{तण्डुलं}}{\Afootnote{N H; \skt{तण्डुलान् ॥३१॥}\ A1938}}\skt{ ॥}\edtext{}{\lemma{\skt{मासान्सुसाधितान् । सुधौतं}}\Afootnote{H; \skt{माषान्सुभावितान् । सुधौतांस्}\ A1938; \skt{माषां सुसाधितां । सुधौतं}\ N}}\pend

\pstart\versepar \ledsidenote{\passref{1.44.32}}\skt{अवक्षुद्यैकतः }\edtext{\skt{पिण्डां}}{\Afootnote{N H; \skt{पिण्डा}\ A1938}}\skt{ }\edtext{\skt{कृत्वा}}{\Afootnote{N H; \skt{न्कृत्वा}\ A1938}}\skt{ }\edtext{\skt{शुष्कां}}{\Afootnote{H; \skt{शुष्का}\ A1938; \skt{शुस्कां}\ N}}\skt{ }\edtext{\skt{सुचूर्ण्णितां}}{\Afootnote{N H; \skt{न्सुचूर्णिता}\ A1938}}\skt{ । }\pend
\pstart\versepar \skt{}\edtext{\edtext{\skt{शालितण्डुलचूर्ण्णन्तुसम्यक्स्विन्नं}}{\lemma{\skt{शालितण्डुलचूर्ण्णन्}}\Afootnote{H; \skt{न् । शालितण्डुलचूर्णं}\ A1938; \skt{सालितण्डुलचूर्ण्णन्}\ N}}}{\lemma{\skt{तुसम्यक्स्विन्नं}}\Afootnote{N H; \skt{च तत्कषायो}\ A1938}}\skt{ }\edtext{\skt{सुशीतलं}}{\Afootnote{H; \skt{ष्मसाधितम् ॥३२॥}\ A1938; \skt{सुसीतलं ।}\ N}}\skt{ ॥}\pend

\pstart\versepar \ledsidenote{\passref{1.44.33}}\skt{तस्य पिष्टस्य भागास्त्रीन्किण्वभागविमिश्रितान् । }\pend
\pstart\versepar \skt{मण्डोदकार्थे }\edtext{\skt{क्वाथन्तु}}{\Afootnote{H; \skt{क्वाथं च}\ A1938 N}}\skt{ दद्यात्तत्सर्वमेकतः ॥}\edtext{}{\lemma{\skt{भागास्त्रीन्किण्वभागविमिश्रितान् । मण्डोदकार्थे}}\Afootnote{A1938 H; \skt{भागास्त्री किण्चभागविमिश्र्तां । मण्डोदकार्थे}\ N}}\pend

\pstart\versepar \ledsidenote{\passref{1.44.34}}\edtext{\skt{निदध्यात्कलसे }\edtext{\skt{तान्तुसुराञ्च}}{\lemma{\skt{तुसुराञ्च}}\Afootnote{H; \skt{तु सुरां जा}\ A1938; \skt{तु सुरां ज}\ N}}}{\lemma{\skt{निदध्यात्कलसे तान्}}\Afootnote{A1938 H; \skt{निंदध्यात्कलसैतान्}\ N}}\skt{ }\edtext{\skt{तरसाम्पिबेत्}}{\lemma{\skt{तरसाम्}}\Afootnote{A1938 H; \skt{तरसा}\ N}}\skt{ । }\pend
\pstart\versepar \skt{एष एव सुराकल्पो वमनेष्वपि कीर्त्तितः ॥}\edtext{}{\lemma{\skt{पिबेत् । एष}}\Afootnote{A1938 H; \skt{ंस्पिबेत् । एष}\ N}}\pend

\pstart\versepar \ledsidenote{\passref{1.44.35}}\skt{मूलानि }\edtext{\skt{तृवृतादीनाम्प्रथमस्य}}{\lemma{\skt{तृवृतादीनाम्}}\Afootnote{N H; \skt{त्रिवृदादीनां}\ A1938}}\skt{ गणस्य च । }\pend
\pstart\versepar \skt{महतः पञ्चमूलस्य }\edtext{\skt{भार्ग्गी}}{\Afootnote{N H; \skt{मूर्वा}\ A1938}}\skt{ }\edtext{\skt{शार्ङ्गष्वयोरपि}}{\lemma{\skt{शार्ङ्गष्वयोर्}}\Afootnote{H; \skt{शार्ङ्गष्टयोर्}\ A1938; \skt{सार्ङ्गष्टयोर्}\ N}}\skt{ ॥}\pend

\pstart\versepar \ledsidenote{\passref{1.44.36}}\edtext{\edtext{\skt{त्रिफलाम्वचामतिविषां}}{\lemma{\skt{त्रिफलाम्}}\Afootnote{H; \skt{तृफलां}\ N}}\skt{ }\edtext{\edtext{\skt{सुधाहेमवतीन्तथा}}{\lemma{\skt{सुधाहेमवतीन्}}\Afootnote{H; \skt{सुधां हैमवतीं}\ A1938; \skt{सुधां हेमवतीन्}\ N}}}{\lemma{\skt{तथा}}\Afootnote{H; \skt{चैव त्रिफलातिविषे}\ A1938; \emph{om.}\ N}}}{\lemma{\skt{… सु}}\Afootnote{\emph{om.}\ A1938}}\skt{ । }\pend
\pstart\versepar \skt{}\edtext{\skt{संहृत्यैतानि}}{\Afootnote{H; \skt{वचाम् । संहृत्यैतानि}\ A1938; \skt{तेषान् । संहृत्यैतानि}\ N}}\skt{ सर्व्वाणि }\edtext{\edtext{\skt{कुर्याद्भागावुभौ}}{\Afootnote{H; \skt{कुर्याद्भागावुभौ}\ N}}\skt{ पृथक् ॥}}{\Afootnote{\emph{om.}\ A1938}}\pend

\pstart\versepar \ledsidenote{\passref{1.44.37}}\edtext{\edtext{\skt{कुर्यान्निःक्वाथमेकस्मिनेकस्मिन्चूर्ण्णमेव}}{\lemma{\skt{एकस्मिनेकस्मिन्चूर्ण्णम्}}\Afootnote{H; \skt{एकस्मिन्नेकस्मिंश्चूर्णम्}\ A1938; \skt{एकस्मिन्नेकस्मिञ्चूर्ण्णम्}\ N}}}{\lemma{\skt{एव}}\Afootnote{N H; \skt{एव तु ।}\ A1938}}\skt{ }\edtext{\skt{च}}{\Afootnote{N H; \skt{क्षुण्णांस्तस्मिंस्}\ A1938}}\skt{ । }\pend
\pstart\versepar \skt{}\edtext{\skt{तेन}}{\Afootnote{N H; \skt{तु निः}\ A1938}}\skt{ }\edtext{\skt{क्वाथेन}}{\Afootnote{N H; \skt{क्वाथे}\ A1938}}\skt{ }\edtext{\skt{बहुशो}}{\Afootnote{H; \emph{om.}\ A1938; \skt{बहुसो}\ N}}\skt{ }\edtext{\skt{विशुद्धां}}{\Afootnote{N H; \emph{om.}\ A1938}}\skt{ }\edtext{\skt{भावयेद्यवान्}}{\Afootnote{N H; \skt{भावयेद्बहुशो यवान् ॥३७॥}\ A1938}}\skt{ ॥}\pend

\pstart\versepar \ledsidenote{\passref{1.44.38}}\edtext{\edtext{\edtext{\skt{शुष्काणाम्मृदुभृष्टानान्तेषाम्भागास्त्रयो}}{\lemma{\skt{मृदुभृष्टानान्}}\Afootnote{A1938 H; \skt{मधुसृष्टानां स्}\ N}}}{\lemma{\skt{भागास्}}\Afootnote{N H; \skt{भागा}\ A1938}}}{\lemma{\skt{त्रयो}}\Afootnote{N H; \skt{स्त्रयो}\ A1938}}\skt{ }\edtext{\skt{मिताः}}{\Afootnote{N H; \skt{मताः ।}\ A1938}}\skt{ । }\pend
\pstart\versepar \skt{}\edtext{\edtext{\skt{चतुर्थम्भागमावाप्य}}{\lemma{\skt{भागम्}}\Afootnote{H; \skt{भाग}\ A1938 N}}}{\lemma{\skt{आवाप्य}}\Afootnote{H; \skt{मावाप्य}\ A1938; \skt{वा वाप्य}\ N}}\skt{ }\edtext{\edtext{\skt{चूर्ण्णमप्यत्र}}{\lemma{\skt{चूर्ण्णम्}}\Afootnote{N H; \skt{चूर्णा}\ A1938}}}{\lemma{\skt{अप्यत्र}}\Afootnote{H; \skt{नामत्र}\ A1938; \skt{अप्यत्र}\ N}}\skt{ कीर्त्तितं ॥}\pend

\pstart\versepar \ledsidenote{\passref{1.44.39}}\edtext{\skt{कलसे}}{\Afootnote{N H; \skt{प्रक्षिप्य कलसे}\ A1938}}\skt{ }\edtext{\edtext{\edtext{\skt{प्रक्षियेद्विद्वान्ततस्तन्तदनन्तरं}}{\lemma{\skt{प्रक्षियेद्विद्वान्}}\Afootnote{H; \skt{सम्य}\ A1938; \skt{प्रक्षिपेद्विद्व}\ N}}}{\lemma{\skt{ततस्तन्}}\Afootnote{H; \skt{क्ततस्तं}\ A1938; \skt{ततस्त}\ N}}}{\lemma{\skt{तदनन्तरं}}\Afootnote{A1938 H; \skt{त्तदनन्तरं ।}\ N}}\skt{ । }\pend
\pstart\versepar \skt{तेषामेव कषायेण }\edtext{\skt{शीतेनाभिप्रपूरयेत्}}{\Afootnote{N H; \skt{शीतलेन सुयोजितम् ॥३९॥}\ A1938}}\skt{ ॥}\pend

\pstart\versepar \ledsidenote{\passref{1.44.40}}\skt{पूर्व्ववत्सन्निदध्यात्तु ज्ञेयं }\edtext{\skt{सौवीरकं}}{\Afootnote{A1938 H; \skt{सौवीरक}\ N}}\skt{ }\edtext{\skt{हितं}}{\Afootnote{N H; \skt{हि त}\ A1938}}\skt{ । }\pend
\pstart\versepar \skt{}\edtext{\skt{पूर्व्वोक्त}}{\Afootnote{N H; \skt{त् । पूर्वोक्तं}\ A1938}}\skt{ वर्ग्गमाहृत्य द्विधा कृत्वैकमेतयोः ॥}\pend

\pstart\versepar \ledsidenote{\passref{1.44.41}}\edtext{\skt{भागं}}{\Afootnote{A1938 H; \skt{भाग}\ N}}\skt{ संक्षुद्य }\edtext{\skt{संसृज्य}}{\Afootnote{A1938 H; \skt{संसृत्य}\ N}}\skt{ यवैः }\edtext{\skt{स्थाल्यां}}{\Afootnote{N H; \skt{स्थाल्या}\ A1938}}\skt{ सहाक्षिपेत् । }\pend
\pstart\versepar \skt{अजशृग्याः कषायेण }\edtext{\skt{तेनाभ्यासिच्य}}{\Afootnote{N H; \skt{तमभ्यासिच्य}\ A1938}}\skt{ साधयेत् ॥}\edtext{}{\lemma{\skt{सहाक्षिपेत् । अजशृग्याः}}\Afootnote{H; \skt{मधिश्रयेत् । अजशृङ्ग्याः}\ A1938; \skt{सहा क्षिपेत् ॥ अजशृङ्ग्या}\ N}}\pend

\pstart\versepar \ledsidenote{\passref{1.44.42}}\edtext{\skt{सुसिद्धां}}{\Afootnote{N H; \skt{सुसिद्धांश्}\ A1938}}\skt{ }\edtext{\edtext{\skt{चावचार्यैतामोषधिभ्यो}}{\lemma{\skt{चावचार्यैताम्}}\Afootnote{N H; \skt{चावतार्यैता}\ A1938}}}{\lemma{\skt{ओषधिभ्यो}}\Afootnote{H; \skt{नौषधिभ्यो}\ A1938; \skt{औषधिभ्यो}\ N}}\skt{ विमोक्षयेत् । }\pend
\pstart\versepar \skt{विमृद्य }\edtext{\edtext{\skt{सतुषानेतान्ततस्तां}}{\lemma{\skt{सतुषानेतान्}}\Afootnote{N H; \skt{सतुषान्सम्य}\ A1938}}}{\lemma{\skt{ततस्तां}}\Afootnote{H; \skt{क्ततस्ता}\ A1938; \skt{ततस्तां}\ N}}\skt{ }\edtext{\skt{पूर्व्ववद्यवां}}{\Afootnote{H; \skt{न्पूर्ववन्मितान् ॥४२॥}\ A1938; \skt{पूर्व्ववद्यवा ।}\ N}}\skt{ ॥}\edtext{}{\lemma{\skt{विमोक्षयेत् । विमृद्य}}\Afootnote{N H; \skt{विवेचयेत् । विमृद्य}\ A1938}}\pend

\pstart\versepar \ledsidenote{\passref{1.44.43}}\skt{पूर्व्वोक्तौषधभागस्य चूर्ण्णं }\edtext{\skt{दत्वा}}{\Afootnote{H; \skt{दत्त्वा तु}\ A1938 N}}\skt{ । }\pend
\pstart\versepar \skt{}\edtext{\skt{तेनैव}}{\Afootnote{\skt{पूर्ववत्तेनैव}\ A1938 N; \skt{… । तेनैव}\ H}}\skt{ सहयूषेण कलसे पूर्व्व }\edtext{\skt{सुक्षिपेत्}}{\Afootnote{H; \skt{वत्क्षिपेत्}\ A1938 N}}\skt{ ॥}\pend

\pstart\versepar \ledsidenote{\passref{1.44.44}}\skt{ज्ञात्वा जातरसञ्चापि तत्तुषोदकमाप्नुयात् । }\pend
\pstart\versepar \skt{}\edtext{\edtext{\skt{तुषोदसौवीरकयोर्व्विधिरेष}}{\lemma{\skt{व्विधिर्}}\Afootnote{N H; \skt{विधि}\ A1938}}}{\lemma{\skt{एष}}\Afootnote{N H; \skt{रेष}\ A1938}}\skt{ प्रकीर्त्तितः ॥}\edtext{}{\lemma{\skt{आप्नुयात् । तुषोदसौवीरकयोर्}}\Afootnote{H; \skt{आदिशेत् । तुषाम्बुसौवीरकयोर्}\ A1938; \skt{आप्नुयात् ॥ तुषोदकसौवीरयो}\ N}}\pend

\pstart\versepar \ledsidenote{\passref{1.44.45}}\edtext{\skt{षड्रात्त्रात्सप्तरात्राद्वा ते च}}{\Afootnote{H; \skt{षड्रात्रात्सप्तरात्राद्वा ते च}\ A1938; \skt{षट्रात्रात्सप्तरात्राद्वा ते च}\ N}}\skt{ पेये }\edtext{\skt{सुरे}}{\Afootnote{N H; \skt{प्रकी}\ A1938}}\skt{ }\edtext{\skt{स्मृते}}{\Afootnote{N H; \skt{र्तिते ।}\ A1938}}\skt{ । }\pend
\pstart\versepar \skt{वैरेचनेषु }\edtext{\skt{द्रव्येषु}}{\Afootnote{N H; \skt{सर्व्येषु}\ A1938}}\skt{ }\edtext{\skt{तृवृत्कल्पविधिः}}{\Afootnote{N H; \skt{त्रिवृन्मूलविधिः}\ A1938}}\skt{ स्मृतः ॥}\pend

\pstart\versepar \ledsidenote{\passref{1.44.82}}\edtext{\skt{महावृक्षपयः}}{\Afootnote{A1938 N H\textsuperscript{pc}; \emph{om.}\ H\textsuperscript{ac}}}\skt{ }\edtext{\edtext{\skt{पीतैर्यवागूस्तण्डुलैः}}{\lemma{\skt{पीतैर्}}\Afootnote{A1938 N H\textsuperscript{pc}; \emph{om.}\ H\textsuperscript{ac}}}}{\lemma{\skt{यवागूस्तण्डुलैः}}\Afootnote{\skt{यवागूस्तण्डुलैः}\ A1938 H\textsuperscript{pc}; \skt{यवागूस्तण्डुलै}\ N; \emph{om.}\ H\textsuperscript{ac}}}\skt{ }\edtext{\skt{कृताः}}{\Afootnote{N H\textsuperscript{pc}; \skt{कृता । पीता}\ A1938; \emph{om.}\ H\textsuperscript{ac}}}\skt{ । }\pend
\pstart\versepar \skt{}\edtext{\skt{विरेचयेदाशुपीतागुडेनोत्कारिका}}{\Afootnote{N H\textsuperscript{pc}; \skt{विरेचयत्याशु गुडेनोत्कारिका}\ A1938; \emph{om.}\ H\textsuperscript{ac}}}\skt{ }\edtext{\skt{कृता}}{\Afootnote{A1938 H\textsuperscript{pc}; \skt{कृताः ।}\ N; \emph{om.}\ H\textsuperscript{ac}}}\skt{ ॥}\pend

\pstart\versepar \ledsidenote{\passref{1.44.83}}\skt{लेहो वा साधितः }\edtext{\skt{सम्यक्स्नुहीक्षीरसुराघृतैः}}{\Afootnote{H; \skt{सम्यक्स्नुहीक्षीरपयोघृतैः ।}\ A1938; \skt{सम्यक्स्तुहीक्षीरसुराघृतैः ।}\ N}}\skt{ । }\pend
\pstart\versepar \skt{}\edtext{\skt{विभाविता}}{\Afootnote{H; \skt{भावितास्तु}\ A1938; \skt{विभावितास्तु}\ N}}\skt{ }\edtext{\skt{स्नुहीक्षीरेपिप्पल्यो}}{\Afootnote{A1938 H; \skt{हीक्षीरे पिप्पल्यो}\ N}}\skt{ }\edtext{\skt{लवणानि}}{\Afootnote{N H; \skt{लवणान्वि}\ A1938}}\skt{ }\edtext{\skt{च}}{\Afootnote{N H; \skt{ताः ॥८३॥}\ A1938}}\skt{ ॥}\pend

\pstart \ledsidenote{\passref{1.44.84ab}}\edtext{\skt{चूर्ण्णङ्कम्पिल्यकस्यापि}}{\Afootnote{N H; \skt{चूर्णं काम्पिल्लकं वाऽपि}\ A1938}}\skt{ }\edtext{\skt{तत्पीतङ्गुडिकाकृतं}}{\lemma{\skt{गुडिकाकृतं}}\Afootnote{N H; \skt{गुटिकीकृतम्}\ A1938}}\skt{ ।}\pend

\pstart\versepar \ledsidenote{\passref{1.44.64}}\edtext{\skt{हरीतकी}}{\Afootnote{A1938 H; \skt{हरीतकीस्}\ N}}\skt{ विडङ्गानि }\edtext{\skt{सैन्धवन्नागरन्तथा}}{\lemma{\skt{तथा}}\Afootnote{N H; \skt{त्रिवृ}\ A1938}}\skt{ ॥ }\pend
\pstart\versepar \skt{}\edtext{\skt{मरीचानि च}}{\Afootnote{H; \skt{त् । मरिचानि च}\ A1938; \skt{मरिचानि च}\ N}}\skt{ }\edtext{\skt{तत्सर्व्वङ्गोमूत्रेण}}{\lemma{\skt{गोमूत्रेण}}\Afootnote{N H; \skt{गोमुत्रेण}\ A1938}}\skt{ विरेचनं ।}\pend

\pstart\versepar \ledsidenote{\passref{1.44.65}}\edtext{\skt{हरीतकीम्भद्रदारु}}{\lemma{\skt{हरीतकीम्}}\Afootnote{N\textsuperscript{pc} H; \skt{हरीतकी}\ A1938; \skt{हईरीतकीं}\ N\textsuperscript{ac}}}\skt{ कुष्ठं पूगफलन्तथा ॥ }\pend
\pstart\versepar \skt{सैन्धवं }\edtext{\skt{शृङ्गवेरञ्च}}{\Afootnote{A1938 H; \skt{सृंगवेरञ्च}\ N}}\skt{ गोमूत्रेण विरेचनं ।}\pend

\pstart\versepar \ledsidenote{\passref{1.44.66}}\edtext{\skt{नीलीफलनाञ्चूर्ण्णन्तुनागराभययोस्तथा}}{\lemma{\skt{नीलीफलनाञ्}}\Afootnote{H; \skt{नीलिनीफल}\ A1938; \skt{नीलीफलाना}\ N}}\skt{ ॥ }\pend
\pstart\versepar \skt{}\edtext{\skt{लीढ्वा}}{\Afootnote{H; \skt{लिह्या}\ A1938 N}}\skt{ }\edtext{\skt{गुडैण}}{\Afootnote{H; \skt{द्गुडेन}\ A1938 N}}\skt{ }\edtext{\edtext{\edtext{\skt{सलिलम्पश्चादुष्णम्पिवेन्नरः}}{\lemma{\skt{सलिलम्पश्}}\Afootnote{H; \skt{सलिलं प}\ A1938 N}}}{\lemma{\skt{चादुष्णम्}}\Afootnote{H; \skt{श्चादुष्णं}\ A1938; \skt{श्चादुष्णं}\ N}}}{\lemma{\skt{पिवेन्}}\Afootnote{H; \skt{पिबेन्}\ A1938 N}}\skt{ ।}\pend

\pstart\versepar \ledsidenote{\passref{1.44.67}}\skt{पिप्पल्यादि कषायेण }\edtext{\skt{पिबेत्पिष्टां}}{\Afootnote{A1938 H; \skt{पिबेत्पिष्टं}\ N}}\skt{ हरीतकीं ॥ }\pend
\pstart\versepar \skt{सैन्धवोपहितां }\edtext{\skt{सम्यगेष}}{\Afootnote{N H; \skt{सद्य एष}\ A1938}}\skt{ योगो }\edtext{\skt{विरेचनः}}{\Afootnote{N H; \skt{विरेचयेत् ॥६७॥}\ A1938}}\skt{ ।}\pend

\pstart \ledsidenote{\passref{1.44.71cd}}\skt{त्रिफला सर्व्वरोगघ्नी }\edtext{\skt{घृतभोगेन}}{\Afootnote{H; \skt{त्रिभागघृत}\ A1938; \skt{घृतरोगेन}\ N}}\skt{ मूर्च्छिता ॥}\pend

\pstart \ledsidenote{\passref{1.44.72ab}}\edtext{\skt{वयः}}{\Afootnote{H; \skt{वय}\ A1938; \skt{यवस्}\ N}}\skt{ }\edtext{\skt{संस्थापनञ्चापि}}{\lemma{\skt{संस्थापनञ्}}\Afootnote{H; \skt{सः स्थापनं}\ A1938; \skt{सं श्थापणञ्}\ N}}\skt{ }\edtext{\skt{कुर्यात्सततसेविता}}{\Afootnote{H; \skt{कुर्यात्सन्ततसेविता ।७२।}\ A1938; \skt{कुर्यात्सततः सेविता ॥}\ N}}\skt{ ।}\pend

\pstart\versepar \ledsidenote{\passref{1.44.68}}\skt{हरीतकी भक्ष्यमाणा नागरेण गुडेन वा ॥ }\pend
\pstart\versepar \skt{}\edtext{\skt{सैन्धवोपहिता}}{\Afootnote{A1938 H; \skt{सेन्धवोपहिता}\ N}}\skt{ वापि सातत्येनाग्निदीपनी ।}\pend

\pstart \ledsidenote{\passref{1.44.54cd}}\skt{व्योषं }\edtext{\skt{त्रिजातकम्मुस्ता}}{\lemma{\skt{मुस्ता}}\Afootnote{A1938 H; \emph{om.}\ N}}\skt{ }\edtext{\skt{विडङ्गामलके}}{\Afootnote{A1938 H; \skt{विडङ्गामलकन्}\ N}}\skt{ तथा ॥}\pend

\pstart\versepar \ledsidenote{\passref{1.44.55}}\skt{नवैतानि }\edtext{\skt{समानानि}}{\Afootnote{H; \skt{समांशानि}\ A1938 N}}\skt{ }\edtext{\skt{तृवृदष्टगाणानि वा}}{\Afootnote{H; \skt{त्रिवृदष्टगुणानि वै ।}\ A1938; \skt{तृवृदष्टगुणानि वा ।}\ N}}\skt{ । }\pend
\pstart\versepar \skt{}\edtext{\skt{सुश्लक्ष्णचूर्ण्णितानीह गाडितानि}}{\Afootnote{H; \skt{श्लक्ष्णचूर्णीकृतानीह दन्तीभागद्वयं तथा ॥५५॥}\ A1938; \skt{सुश्लक्ष्ण चूर्ण्णडानीह गाढिडानि}\ N}}\skt{ }\edtext{\skt{विमिश्रयेत्}}{\Afootnote{N H; \emph{om.}\ A1938}}\skt{ ॥}\pend

\pstart \ledsidenote{\passref{1.44.56cd}}\skt{षड्भिश्च }\edtext{\edtext{\skt{शर्क्कराभागैरीषत्सैन्धवमाक्षिकैः}}{\lemma{\skt{शर्क्कराभागैर्}}\Afootnote{A1938 H; \skt{शर्करागैर्}\ N}}}{\lemma{\skt{ईषत्सैन्धवमाक्षिकैः}}\Afootnote{A1938 H; \skt{ईषत्सेन्धवमाक्षिकेः ।}\ N}}\skt{ ।}\pend

\pstart \ledsidenote{\passref{1.44.57ab}}\edtext{\edtext{\skt{पिण्डीकृतम्भक्षयित्वाततः}}{\lemma{\skt{पिण्डीकृतम्}}\Afootnote{H; \skt{पिण्डितं}\ A1938; \skt{पिण्डीकृत}\ N}}}{\lemma{\skt{भक्षयित्वाततः}}\Afootnote{N H; \skt{भक्षयित्वा तु ततः}\ A1938}}\skt{ शीताम्बु पाययेत् ॥}\pend

\pstart \ledsidenote{\passref{1.44.58cd}}\edtext{\edtext{\skt{अविपत्तिरयं}}{\lemma{\skt{अविपत्तिर्}}\Afootnote{N H; \skt{त्रिवृदष्टकसञ्ज्ञो}\ A1938}}}{\lemma{\skt{अयं}}\Afootnote{H; \skt{ऽ}\ A1938; \skt{अप}\ N}}\skt{ }\edtext{\skt{योगः}}{\Afootnote{N H; \skt{यं}\ A1938}}\skt{ }\edtext{\skt{प्रशस्तः}}{\Afootnote{A1938 H; \skt{प्रसस्त}\ N}}\skt{ पित्तरोगिणां ।}\pend

\pstart\versepar \ledsidenote{\passref{1.44.59}}\edtext{\edtext{\skt{क्षारानुपानम्भोक्तव्यो}}{\lemma{\skt{क्षारानुपानम्}}\Afootnote{H\textsuperscript{pc}; \skt{भक्ष्यः क्षीरानुपानो}\ A1938; \skt{क्षारामुपान}\ N}}}{\lemma{\skt{भोक्तव्यो}}\Afootnote{H\textsuperscript{pc}; \skt{वा}\ A1938; \skt{भोक्त्व्यो}\ N}}\skt{ }\edtext{\skt{वातश्लेमार्तुरैर्न्नरैः}}{\lemma{\skt{वातश्लेमार्तुरैर्}}\Afootnote{\skt{पित्तश्लेष्मातुरैर्}\ A1938; \skt{वातश्लेष्मातुरैर्}\ N; \skt{वातश्लेमर्तुरैर्}\ H\textsuperscript{pc}}}\skt{ ॥ }\pend
\pstart\versepar \skt{}\edtext{\edtext{\skt{भक्ष्यरूपसधर्म्मत्वादाढ्येष्वेव}}{\lemma{\skt{भक्ष्यरूपसधर्म्मत्वादाढ्येष्व्}}\Afootnote{H; \skt{भक्ष्यरूपसधर्मत्वादाढ्ये}\ A1938; \skt{भक्षरूपसधर्मत्वादाढ्येष्व्}\ N}}}{\lemma{\skt{एव}}\Afootnote{N H; \skt{ष्वेव}\ A1938}}\skt{ विधीयते ।}\edtext{}{\lemma{\skt{… भ}}\Afootnote{\emph{om.}\ H\textsuperscript{ac}}}\pend

\pstart \ledsidenote{\passref{1.44.86ef}}\edtext{\skt{अवेक्ष्य}}{\Afootnote{A1938 N H\textsuperscript{pc}; \emph{om.}\ H\textsuperscript{ac}}}\skt{ }\edtext{\skt{सम्यग्रोगादीन्यथावदुपयोजयेत्}}{\Afootnote{A1938 N H\textsuperscript{pc}; \emph{om.}\ H\textsuperscript{ac}}}\skt{ ॥}\pend

\pstart \ledsidenote{\passref{1.44.84cd}}\edtext{\skt{सप्तलां}}{\Afootnote{N H; \skt{सप्तला}\ A1938}}\skt{ }\edtext{\edtext{\edtext{\edtext{\skt{शाङ्खिनीन्दन्तीन्तृवृदारग्वधम्वचां}}{\lemma{\skt{शाङ्खिनीन्}}\Afootnote{H; \skt{शङ्खिनी}\ A1938; \skt{साङ्खिनी}\ N}}}{\lemma{\skt{दन्तीन्}}\Afootnote{H; \skt{दन्ती}\ A1938 N}}}{\lemma{\skt{तृवृदारग्वधम्}}\Afootnote{H; \skt{त्रिवृदारग्वधं}\ A1938; \skt{तृवृदारग्वं}\ N}}}{\lemma{\skt{वचां}}\Afootnote{N H; \skt{गवाम् ॥८४॥}\ A1938}}\skt{ ।}\pend

\pstart\versepar \ledsidenote{\passref{1.44.85}}\edtext{\skt{मूत्रेण}}{\Afootnote{N H; \skt{मूत्रेणा}\ A1938}}\skt{ }\edtext{\skt{भाव्यं}}{\Afootnote{N H; \skt{प्लाव्य}\ A1938}}\skt{ सप्ताहं }\edtext{\skt{स्नुहाक्षीरे}}{\Afootnote{N H; \skt{स्नुहीक्षीरे}\ A1938}}\skt{ }\edtext{\skt{तथैव च}}{\Afootnote{N H; \skt{ततः परम् ।}\ A1938}}\skt{ ॥ }\pend
\pstart\versepar \skt{}\edtext{\edtext{\skt{कीर्ण्णन्तेनाथचूर्ण्णेन}}{\lemma{\skt{कीर्ण्णन्}}\Afootnote{A1938 H; \skt{कीर्ण्ण}\ N}}}{\lemma{\skt{तेनाथचूर्ण्णेन}}\Afootnote{H; \skt{तेनैव चूर्णेन}\ A1938; \skt{तेनावे चूर्ण्णेन}\ N}}\skt{ माल्यं वसनमेव }\edtext{\skt{वा}}{\Afootnote{H; \skt{च}\ A1938 N}}\skt{ ।}\pend

\pstart \ledsidenote{\passref{1.44.86ab}}\edtext{\skt{आघ्रायावृत्य वा}}{\Afootnote{A1938 H; \skt{आघ्रां यावृत्य वा}\ N}}\skt{ }\edtext{\skt{सम्यक्मृदुकोष्ठो}}{\Afootnote{N H; \skt{सम्यङ्मृदुकोष्ठो}\ A1938}}\skt{ }\edtext{\skt{विच्यते}}{\Afootnote{N H; \skt{विरिच्यते ।}\ A1938}}\skt{ ॥}\pend

\pstart\versepar \ledsidenote{\passref{1.44.90}}\skt{घृतेषु तैलेषु पयःसु }\edtext{\skt{चापि}}{\Afootnote{A1938 N\textsuperscript{pc} H; \skt{चापिः ।}\ N\textsuperscript{ac}}}\skt{ मद्येषु मूत्रेषु तथा रसेषु । }\pend
\pstart\versepar \skt{}\edtext{\skt{अन्नेषु}}{\Afootnote{N H; \skt{भक्ष्यान्न}\ A1938}}\skt{ }\edtext{\skt{भक्ष्येषु}}{\Afootnote{N H; \skt{लेह्येषु}\ A1938}}\skt{ }\edtext{\skt{तथैव}}{\Afootnote{N H; \skt{च तेषु}\ A1938}}\skt{ }\edtext{\skt{लेहे}}{\Afootnote{N H; \skt{तेषु}\ A1938}}\skt{ }\edtext{\skt{विरेचनं}}{\Afootnote{N H; \skt{विरेच}\ A1938}}\skt{ }\edtext{\skt{साधु}}{\Afootnote{N H; \emph{om.}\ A1938}}\skt{ }\edtext{\skt{नियोजनीय}}{\Afootnote{H; \skt{नान्यग्रमतिर्विदध्यात् ॥९०॥}\ A1938; \skt{नियोजनीय इति ॥}\ N}}\pend

\pstart \ledsidenote{\passref{1.44.trailer}}\edtext{\skt{॥ इति}}{\Afootnote{H; \skt{इति सुश्रुतसंहितायां सूत्रस्थाने विरेचनद्रव्यविकल्पविज्ञानीयो नाम चतुश्चत्वारिंशोऽध्यायः ॥४४॥}\ A1938; \emph{om.}\ N}}\skt{ ॥ ॥}\pend

\endnumbering
\beginnumbering
\eledsection*{Adhyāya 45}
\pstart \ledsidenote{\passref{1.45.1}}\skt{अथातो }\edtext{\skt{द्रवद्रव्यविशेषविज्ञानीयं}}{\Afootnote{H; \skt{द्रवद्रव्यविधिमध्यायं}\ A1938; \skt{द्रव्यविशेषविज्ञानीयं व्याख्यास्यामः ॥}\ N}}\skt{ }\edtext{\skt{व्याख्यास्यामः}}{\Afootnote{A1938 H; \emph{om.}\ N}}\skt{ ॥}\pend

\pstart \ledsidenote{\passref{1.45.3}}\edtext{\edtext{\skt{पानीयंमान्तरीक्ष्यमनिर्द्देश्यरसममृतं}}{\lemma{\skt{आन्तरीक्ष्यम्}}\Afootnote{N H; \skt{आन्तरीक्षम्}\ A1938}}}{\lemma{\skt{अनिर्द्देश्यरसम्}}\Afootnote{A1938 H; \skt{अनिर्देश्यं रसम्}\ N}}\skt{ }\edtext{\edtext{\skt{जीवनन्तर्प्पणन्धारणमाश्वासजननं}}{\lemma{\skt{जीवनन्}}\Afootnote{A1938 H; \skt{जीवन}\ N}}}{\lemma{\skt{तर्प्पणन्}}\Afootnote{A1938 H; \skt{न्वर्पणं}\ N}}\skt{ }\edtext{\skt{श्रमक्लमपिपासामदमूर्च्छातन्द्राप्रशमनमेकान्ततः}}{\lemma{\skt{श्रमक्लमपिपासामदमूर्च्छातन्द्राप्रशमनम्}}\Afootnote{N H; \skt{श्रमक्लमपिपासामदमूर्च्छातन्द्रानिद्रादाहप्रशमनम्}\ A1938}}\skt{ पथ्यतमञ्च ॥}\pend

\pstart \ledsidenote{\passref{1.45.4}}\edtext{\skt{तदेवावनौ}}{\Afootnote{N H; \skt{तदेवावनि}\ A1938}}\skt{ }\edtext{\skt{पतितमन्यतममुपलभ्यते}}{\lemma{\skt{उपलभ्यते}}\Afootnote{H; \skt{रसमुपलभते}\ A1938; \skt{उपलभते ।}\ N}}\skt{ । }\edtext{\skt{नदीसरस्तडागकूपवापीप्रश्रवणोद्भिजचौण्ट्यादिषु}}{\Afootnote{H; \skt{स्थानविशेषान्नदीनदसरस्तडागवापीकूपचुण्टीप्रस्रवणोद्भिदविकिरकेदारपल्वलादिषु}\ A1938; \skt{नदीसरस्तडागकूपवापीप्रश्रवणाद्भिजचौण्ट्यादिषु}\ N}}\skt{ }\edtext{\skt{स्थानेष्वन्यतमं}}{\lemma{\skt{अन्यतमं}}\Afootnote{N H; \skt{अवस्थितमिति ॥४॥}\ A1938}}\skt{ }\edtext{\edtext{\skt{रसमुपभते}}{\lemma{\skt{रसम्}}\Afootnote{N H; \emph{om.}\ A1938}}}{\lemma{\skt{उपभते}}\Afootnote{N H; \emph{om.}\ A1938}}\skt{ ।}\pend

\pstart \ledsidenote{\passref{1.45.5}}\skt{तत्र }\edtext{\skt{लोहितकपिलपाण्डुपीतनीलशुक्लेष्ववनिप्रदेशेषु}}{\lemma{\skt{लोहितकपिलपाण्डुपीतनीलशुक्लेष्व्}}\Afootnote{N H; \skt{लोहितकपिलपाण्डुनीलपीतशुक्लेष्व्}\ A1938}}\skt{ }\edtext{\skt{मधुराम्ललवणकटुतिक्तकषाया}}{\Afootnote{H; \skt{मधुराम्ललवणकटुतिक्तकषायाणि}\ A1938; \skt{मधुराम्ब्ललवणकटुतिक्तकषायाः ।}\ N}}\skt{ }\edtext{\skt{यथासंख्यमुदकरसा}}{\lemma{\skt{उदकरसा}}\Afootnote{N H; \skt{उदकानि सम्}\ A1938}}\skt{ }\edtext{\skt{भवन्तीत्येके}}{\lemma{\skt{भवन्तीत्य्}}\Afootnote{N H; \skt{भवन्ति’ इत्य्}\ A1938}}\skt{ भाषन्ते ॥}\pend

\pstart \ledsidenote{\passref{1.45.6}}\skt{तत्तु न }\edtext{\edtext{\skt{सम्यक्पृथिव्यादीनामन्योऽन्यानुप्रवेशकृतः}}{\lemma{\skt{अन्यो}}\Afootnote{H; \skt{अन्यो}\ A1938 N}}\skt{ स}}{\lemma{\skt{ऽन्यानुप्रवेशकृतः स}}\Afootnote{H; \skt{न्यानुप्रवेशकृतः स}\ A1938 N}}\skt{ }\edtext{\skt{खलूदकरसो}}{\Afootnote{N H; \skt{लिलरसो}\ A1938}}\skt{ }\edtext{\skt{भवत्युत्कर्षापकर्षकृतेन}}{\lemma{\skt{उत्कर्षापकर्षकृतेन}}\Afootnote{N H; \skt{उत्कर्षापकर्षे}\ A1938}}\skt{ }\edtext{\skt{संयोगविशेषेण}}{\Afootnote{N H; \skt{ण ।}\ A1938}}\skt{ ॥ तत्र }\edtext{\edtext{\skt{स्वलक्षणोदकगुणभूयिष्ठायाम्मधुरं}}{\lemma{\skt{स्वलक्षणोदकगुणभूयिष्ठायाम्}}\Afootnote{N H; \skt{स्वलक्ष}\ A1938}}\skt{ । }\edtext{\skt{तोयाग्निगुणभूयिष्ठायामम्लं।}}{\lemma{\skt{अम्लं}}\Afootnote{H; \skt{भूमावम्लं}\ A1938; \skt{अम्ल्बं।}\ N}}}{\lemma{\skt{णोदकगुणभूयिष्ठायाम्… तोयाग्निगुण}}\Afootnote{\emph{om.}\ A1938}}\skt{ }\edtext{\skt{पृथिव्यग्निगुणभूयिष्ठायां लवणं । }\edtext{\edtext{\skt{वाय्वग्निगुणभूयिष्ठायाङ्कटुकं।}}{\lemma{\skt{वाय्वग्निगुणभूयिष्ठायाङ्}}\Afootnote{N H; \skt{च अम्बुगुणभूयिष्ठायां}\ A1938}}}{\lemma{\skt{कटुकं}}\Afootnote{N H; \skt{मधुरं;}\ A1938}}}{\lemma{\skt{पृथिव्यग्निगुणभूयिष्ठायां ल}}\Afootnote{\emph{om.}\ A1938}}\skt{ }\edtext{\skt{वाय्वाकाशगुणभूयिष्ठायां}}{\Afootnote{N H; \skt{तेजोगुणभूयिष्ठायां}\ A1938}}\skt{ }\edtext{\skt{तिक्तं}}{\Afootnote{N H; \skt{कटुकं तिक्तं}\ A1938}}\skt{ । }\edtext{\edtext{\skt{पृथिव्यनिलगुणभूयिष्ठायाङ्कषायं}}{\lemma{\skt{पृथिव्यनिलगुणभूयिष्ठायाङ्}}\Afootnote{H; \skt{च वायुगुणभूयिष्ठायां}\ A1938; \skt{पृथिव्यंनिलगुणभूयिष्ठायां}\ N}}}{\lemma{\skt{कषायं}}\Afootnote{A1938 H; \skt{कषायः ॥}\ N}}\skt{ । आकाशगुणभूयिष्ठायामव्यक्तरसमव्यक्तं }\edtext{\edtext{\edtext{\skt{ह्याकाशमित्यतस्तत्प्रधानमव्यक्तरसत्वात्तत्पेयमान्तरीक्षालाभे}}{\lemma{\skt{अतस्}}\Afootnote{A1938 H\textsuperscript{pc}; \skt{आतस्}\ N; \skt{अ}\ H\textsuperscript{ac}}}}{\lemma{\skt{तत्प्रधानम्}}\Afootnote{A1938 H; \skt{तत्प्राधानम्}\ N}}}{\lemma{\skt{अव्यक्तरसत्वात्तत्पेयम्}}\Afootnote{A1938 H; \skt{अव्यक्तरसात्वात्तत्पेयम्}\ N}}\skt{ ।}\pend

\pstart \ledsidenote{\passref{1.45.7}}\edtext{\skt{तत्रान्त-रिक्षञ्चतुर्व्विधं}}{\lemma{\skt{तत्रान्त-रिक्षञ्}}\Afootnote{H; \skt{तत्रान्तरीक्षं}\ A1938 N}}\skt{ ॥ }\edtext{\skt{तद्याथा}}{\Afootnote{N H; \skt{तद्यथा}\ A1938}}\skt{ । धारं कारं तौषारं हैममिति । तेषान्धारं प्रधानं }\edtext{\edtext{\skt{लघुत्वात्तत्पुनर्द्विविधं}}{\lemma{\skt{लघुत्वात्तत्पुनर्}}\Afootnote{N H; \skt{लघुत्वात् ; तत्पुन}\ A1938}}}{\lemma{\skt{द्विविधं}}\Afootnote{N H; \skt{र्द्विविधं-}\ A1938}}\skt{ । गाङ्गं }\edtext{\skt{सामुद्रञ्च}}{\lemma{\skt{च}}\Afootnote{N H; \skt{चेति । तत्र गाङ्गमाश्वयुजे मासि प्रायशो व}\ A1938}}\skt{ }\edtext{\skt{तयोः}}{\Afootnote{N H; \skt{र्षति । तयोर्द्वयोरपि}\ A1938}}\skt{ परीक्षणं । }\edtext{\skt{तत्र}}{\Afootnote{N H; \skt{कुर्वीत-}\ A1938}}\skt{ }\edtext{\edtext{\edtext{\skt{शाल्योदनपिण्डमकुथितमविदग्धं}}{\lemma{\skt{शाल्योदनपिण्डम्}}\Afootnote{N H; \skt{शाल्योदनपिण्ड}\ A1938}}}{\lemma{\skt{अकुथितम्}}\Afootnote{H; \skt{मकुथित}\ A1938; \skt{अप्रकुथितम्}\ N}}}{\lemma{\skt{अविदग्धं}}\Afootnote{N H; \skt{मविदग्धं}\ A1938}}\skt{ । }\edtext{\skt{रजतभाजनोपहितं}}{\Afootnote{A1938 H; \skt{रजतभानोपहितं}\ N}}\skt{ वर्षति देवे }\edtext{\skt{कुर्व्वीत स}}{\Afootnote{N H; \skt{बहिष्कुर्वीत, स}\ A1938}}\skt{ }\edtext{\edtext{\skt{चेत्मुहूर्त्तस्थितस्तादृश}}{\lemma{\skt{चेत्मुहूर्त्तस्थितस्}}\Afootnote{H; \skt{यदि मुहूर्तं स्थित}\ A1938; \skt{चेन्मुहूर्त्तस्थितस्}\ N}}}{\lemma{\skt{तादृश}}\Afootnote{N H; \skt{स्तादृश}\ A1938}}\skt{ एव भवति । }\edtext{\edtext{\skt{तद्गाङ्गमिति}}{\lemma{\skt{तद्गाङ्गम्}}\Afootnote{H; \skt{तदा गाङ्गं}\ A1938; \skt{तद्गाङ्}\ N}}}{\lemma{\skt{इति}}\Afootnote{H; \skt{पतती}\ A1938; \skt{गमिति}\ N}}\skt{ }\edtext{\skt{अवगन्तव्यं}}{\Afootnote{H; \skt{त्यवगन्तव्यं;}\ A1938; \skt{अवगन्तव्यमि}\ N}}\skt{ । }\edtext{\skt{वर्ण्णान्यतासित्थप्रक्लेदो वा}}{\Afootnote{H; \skt{वर्णान्यत्वे सिक्थप्रक्लेदे}\ A1938; \skt{वर्ण्णान्यतासिक्थप्रक्लेदो वा}\ N}}\skt{ }\edtext{\skt{तत्सामुद्रमिति}}{\lemma{\skt{तत्सामुद्रम्}}\Afootnote{H; \skt{च सामुद्रम्}\ A1938; \skt{तं सामुद्रम्}\ N}}\skt{ विद्यात् । तन्नोपादेयं । }\edtext{\edtext{\edtext{\skt{सामुद्रमप्याश्वयुजि}}{\lemma{\skt{सामुद्रम्}}\Afootnote{N H; \skt{सामुद्र}\ A1938}}}{\lemma{\skt{अप्य्}}\Afootnote{N H; \skt{म}\ A1938}}}{\lemma{\skt{आश्वयुजि}}\Afootnote{N H; \skt{प्याश्वयुजे}\ A1938}}\skt{ मासि गृहीतं }\edtext{\skt{गाङ्गवद्भवति}}{\Afootnote{N H; \skt{गाङ्गवद्भवति ।}\ A1938}}\skt{ । }\edtext{\edtext{\skt{तदुपादेयमिति}}{\lemma{\skt{तदुपादेयम्}}\Afootnote{N H; \skt{गाङ्गं पुनः प्रधानं, तदुपाददीताश्वयुजे मासि ।}\ A1938}}}{\lemma{\skt{इति}}\Afootnote{H; \skt{शुचि}\ A1938; \skt{इ}\ N}}\skt{ । }\edtext{\skt{शुक्लपटैकदेशप्रच्युतमथवा}}{\lemma{\skt{शुक्लपटैकदेशप्रच्युतम्}}\Afootnote{H; \skt{शुक्लविततपटैकदेशच्युतम्}\ A1938; \skt{ति । शुक्लपटैकदेशप्रच्युतम्}\ N}}\skt{ }\edtext{\edtext{\skt{हर्म्म्यतलपरिभ्रष्टमन्यैर्व्वा}}{\lemma{\skt{हर्म्म्यतलपरिभ्रष्टम्}}\Afootnote{N H; \skt{हर्म्यतलपरिभ्रष्ट}\ A1938}}}{\lemma{\skt{अन्यैर्व्वा}}\Afootnote{N H; \skt{मन्यैर्वा}\ A1938}}\skt{ }\edtext{\skt{प्रयोगैर्ग्गृहीतं}}{\lemma{\skt{प्रयोगैर्}}\Afootnote{N H; \skt{शुचिभिर्भाजनैर्}\ A1938}}\skt{ }\edtext{\skt{शैलभाजने}}{\Afootnote{N H; \skt{सौवर्णे राजते}\ A1938}}\skt{ }\edtext{\skt{शैलवन्मृन्मये वा}}{\Afootnote{N H; \skt{मृन्मये वा}\ A1938}}\skt{ }\edtext{\edtext{\skt{पात्रेष्वनुगुप्तन्निदध्यात्}}{\lemma{\skt{पात्रेष्व्}}\Afootnote{N H; \skt{पात्रे}\ A1938}}}{\lemma{\skt{अनुगुप्तन्}}\Afootnote{N H; \emph{om.}\ A1938}}\skt{ । तत्सर्व्वकालमुपयुञ्जीत । तस्यालाभे भौमं । }\edtext{\skt{तत्पुनः}}{\Afootnote{N H; \skt{तच्चाकाशगुणबहुलम् । तत्पुनः}\ A1938}}\skt{ सप्तविधं । तद्यथा कौपन्नादेयं }\edtext{\skt{सारसन्ताडाग }\edtext{\edtext{\skt{म्प्रास्रवणमौद्भिजं}}{\lemma{\skt{प्रास्रवणम्}}\Afootnote{A1938 H; \skt{प्राश्रावणम्}\ N}}}{\lemma{\skt{औद्भिजं}}\Afootnote{N H; \skt{औद्भिदं,}\ A1938}}}{\lemma{\skt{ताडाग म्}}\Afootnote{A1938 H; \skt{ताडाकं}\ N}}\skt{ चौण्ट्यमिति ॥}\pend

\pstart \ledsidenote{\passref{1.45.8}}\skt{तत्र }\edtext{\skt{वर्षावर्षास्वान्तरिक्षमौद्भिजं वा}}{\Afootnote{H; \skt{वर्षास्वान्तरिक्षमौद्भिदं वा}\ A1938; \skt{वर्षावर्षास्वान्तरिक्षमौद्भिजं वा}\ N}}\skt{ सेवेत महागुणत्वा-त् । शरदि सर्व्वम्प्रसन्नत्वात् । हेमन्ते सारसन्ताडागम्वा । वसन्ते कौपं }\edtext{\skt{चौण्ट्यं}}{\Afootnote{N H; \emph{om.}\ A1938}}\skt{ }\edtext{\skt{प्रास्रावणम्वा}}{\Afootnote{H; \skt{प्रास्रवणं वा;}\ A1938; \skt{प्राश्रावणम्वा ।}\ N}}\skt{ । }\edtext{\skt{ग्रीष्मेऽप्येवं}}{\lemma{\skt{ऽप्येवं}}\Afootnote{A1938 H; \skt{प्येवं ।}\ N}}\skt{ । }\edtext{\edtext{\skt{प्रावृष्येन्द्रमनभिवृष्टं}}{\lemma{\skt{प्रावृष्येन्द्रम्}}\Afootnote{N H; \skt{प्रावृषि चौण्}\ A1938}}}{\lemma{\skt{अनभिवृष्टं}}\Afootnote{N H; \skt{ट्यमनभिवृष्टं}\ A1938}}\skt{ सर्व्वं }\edtext{\skt{एव}}{\Afootnote{N H; \skt{चेति ॥८॥}\ A1938}}\skt{ ।}\pend

\pstart \ledsidenote{\passref{1.45.9}}\skt{कीटमूत्रपुरीषाण्डशव-कोथप्रदूषितं । }\edtext{\skt{तृणपर्ण्णोदकयुतं}}{\Afootnote{N H; \skt{तृणपर्णोत्करयुतं}\ A1938}}\skt{ कलुषम्विषसंयुतं ॥}\pend

\pstart\versepar \ledsidenote{\passref{1.45.10}}\skt{योऽवगाहेत वर्षासु }\edtext{\skt{पिबेद्वापि}}{\Afootnote{A1938 N; \skt{पिवेद्वापि}\ H}}\skt{ नवञ्जलं । }\pend
\pstart\versepar \skt{स }\edtext{\skt{बाह्याभ्यन्तरान्रोगान्लभते}}{\Afootnote{\skt{बाह्याभ्यन्तरान्रोगान्प्राप्नुया}\ A1938; \skt{बाह्याभ्यन्तरान्रोगां लभते}\ N; \skt{वाह्याभ्यन्तरान्रोगान्लभते}\ H}}\skt{ }\edtext{\skt{क्षिप्रमेव}}{\lemma{\skt{क्षिप्रम्}}\Afootnote{N H; \skt{त्क्षिप्रम्}\ A1938}}\skt{ तु ॥}\pend

\pstart \ledsidenote{\passref{1.45.11}}\skt{तत्र }\edtext{\skt{यत्पङ्कशेवालतृणहठपद्मप्रभृतिभिरवच्छन्नं}}{\lemma{\skt{यत्पङ्कशेवालतृणहठपद्मप्रभृतिभिर्}}\Afootnote{N H; \skt{यत्पङ्कशैवलहठतृणपद्मपत्रप्रभृतिभिर्}\ A1938}}\skt{ }\edtext{\edtext{\skt{रविशशिकिरणानिलैर्व्वातिजुष्टां}}{\lemma{\skt{रविशशिकिरणानिलैर्}}\Afootnote{A1938 H; \skt{रविससिकिरणानिलैर्}\ N}}}{\lemma{\skt{व्वातिजुष्टां}}\Afootnote{H; \skt{नाभिजुष्टं}\ A1938; \skt{वाभिजुष्टं}\ N}}\skt{ }\edtext{\edtext{\skt{गन्धरसोपसृष्टन्तद्व्यापन्नमिति}}{\lemma{\skt{गन्धरसोपसृष्टन्}}\Afootnote{N H; \skt{गन्धवर्णरसोपसृष्टं}\ A1938}}}{\lemma{\skt{तद्व्यापन्नम्}}\Afootnote{N H; \skt{तद्व्यापन्नम्}\ A1938}}\skt{ }\edtext{\skt{विद्यात् । स्पर्शरसरूपगन्धवीर्यविपाकाः}}{\Afootnote{H; \skt{विद्यात् । तस्य स्पर्शरूपरसगन्धवीर्यविपाक}\ A1938; \skt{विद्यात् ॥ स्पर्शरसरूपवीर्यविपाकाः}\ N}}\skt{ । दोषाः }\edtext{\skt{षट्खरता}}{\Afootnote{N H; \skt{षट्सम्भवन्ति । तत्र, खरता}\ A1938}}\skt{ }\edtext{\edtext{\skt{पैच्छिल्यमौष्ण्यं}}{\lemma{\skt{पैच्छिल्यम्}}\Afootnote{A1938 H; \skt{पैच्छिल्याम्}\ N}}}{\lemma{\skt{औष्ण्यं}}\Afootnote{A1938 H; \skt{औष्णं}\ N}}\skt{ दन्तग्राहिता च स्पर्शदोषः । }\edtext{\skt{पङ्कसिकताशैवालवाहुल्याद्विकृतवर्ण्णता}}{\Afootnote{H; \skt{पङ्कसिकताशैवालबहुवर्णता}\ A1938; \skt{पङ्कसिकतासैवाबाहुल्याद्विकृतवर्ण्ण्यता}\ N\textsuperscript{ac}; \skt{पङ्कसिकतासैवालबाहुल्याद्विकृतवर्ण्ण्यता}\ N\textsuperscript{pc}}}\skt{ रूपदोषः । }\edtext{\skt{व्यक्तरसता}}{\Afootnote{A1938 H; \skt{वक्तरसता}\ N}}\skt{ रसदोषः । अनिष्टगन्धता गन्धदोषः । }\edtext{\skt{यदुपयुक्तन्तृष्णागौरवशूलकफप्रसेकानापादयति स}}{\lemma{\skt{तृष्णागौरवशूलकफप्रसेकानापादयति स}}\Afootnote{A1938 H; \skt{तृष्णागौरवशूलकफप्रसेकानापादयति सर्}\ N}}\skt{ }\edtext{\skt{वीर्यदोषः}}{\Afootnote{H; \skt{विर्यदोषः}\ A1938; \skt{वीयदोषः ।}\ N}}\skt{ । यदुपयुक्तं }\edtext{\skt{चिराद्विपच्यते}}{\Afootnote{A1938 H; \skt{चिराद्वा विपच्यते}\ N}}\skt{ }\edtext{\skt{विष्टम्भयति स}}{\Afootnote{N H; \skt{विष्टम्भयति वा स}\ A1938}}\skt{ }\edtext{\skt{विपाकदोषः}}{\Afootnote{N H; \skt{विपाकदोष}\ A1938}}\skt{ ॥ }\edtext{\skt{त}}{\Afootnote{N H; \skt{इति । त}\ A1938}}\skt{ एते }\edtext{\skt{आन्तरिक्षे}}{\Afootnote{A1938 H; \skt{अक्षरिक्षे}\ N}}\skt{ }\edtext{\skt{दोषा न}}{\Afootnote{N H; \emph{om.}\ A1938}}\skt{ }\edtext{\skt{भवन्ति}}{\Afootnote{H; \skt{न सन्ति ॥११॥}\ A1938; \skt{सन्ति ॥}\ N}}\skt{ ।}\pend

\pstart \ledsidenote{\passref{1.45.12}}\edtext{\skt{व्यापन्ने}}{\Afootnote{N H; \skt{व्यापन्नस्य}\ A1938}}\skt{ चाग्निक्वथनं }\edtext{\edtext{\edtext{\skt{सूर्यातपनमयम्पिण्डतप्तनिर्व्वापणम्वा}}{\lemma{\skt{सूर्यातपनम्}}\Afootnote{N H; \skt{सूर्यातपप्रतापनं}\ A1938}}}{\lemma{\skt{अयम्}}\Afootnote{N H; \skt{तप्तायः}\ A1938}}}{\lemma{\skt{पिण्डतप्तनिर्व्वापणम्वा}}\Afootnote{N H; \skt{पिण्डसिकतालोष्ट्राणां वा निर्वापणं}\ A1938}}\skt{ }\edtext{\edtext{\skt{प्रसाधकम्भवति}}{\lemma{\skt{प्रसाधकम्}}\Afootnote{N H; \skt{प्रसादनं}\ A1938}}}{\lemma{\skt{भवति}}\Afootnote{N H; \skt{च कर्तव्यं, नागचम्पकोत्पलपाटलापुष्पप्रभृतिभिश्चाधिवासनमिति ॥१२॥}\ A1938}}\pend

\pstart \ledsidenote{\passref{1.45.17}}\edtext{\skt{सप्तकलुषस्य}}{\Afootnote{N H; \skt{तत्र सप्त कलुषस्य}\ A1938}}\skt{ }\edtext{\skt{प्रसाधकानि}}{\Afootnote{H; \skt{प्रसादनानि}\ A1938 N}}\skt{ }\edtext{\skt{भ्वन्ति}}{\Afootnote{N H; \skt{भवन्ति ।}\ A1938}}\skt{ । तद्यथा ॥ }\edtext{\skt{कतकगोमेदकविसग्रन्थिपर्ण्णीमूलशेवालवस्त्राणिमुक्तामणिश्चेति}}{\lemma{\skt{कतकगोमेदकविसग्रन्थिपर्ण्णीमूलशेवालवस्त्राणिमुक्तामणिश्}}\Afootnote{H; \skt{कतकगोमेदकबिसग्रन्थिशैवालमूलवस्त्राणि मुक्तामणिश्}\ A1938; \skt{कतकगोमेदकविसग्रन्थिपर्ण्णीमूलसेवालवस्त्राणि मुक्ता मणिश्}\ N}}\skt{ ॥}\pend

\pstart \ledsidenote{\passref{1.45.19}}\edtext{\skt{सप्तशीतीकरणानि}}{\Afootnote{A1938 H; \skt{सप्तसीतीकरणानि}\ N}}\skt{ भवन्ति । तद्यथा ॥ }\edtext{\skt{प्रवातस्थापनम्दकप्रक्षेपनं}}{\lemma{\skt{दकप्रक्षेपनं}}\Afootnote{H; \skt{उदकप्रक्षेपणं}\ A1938 N}}\skt{ }\edtext{\skt{यष्टिभ्रामणम्व्यजनं}}{\lemma{\skt{यष्टिभ्रामणम्}}\Afootnote{H; \skt{यष्टिकाभ्रामणं,}\ A1938; \skt{यष्टिभ्रामण}\ N}}\skt{ । वस्त्रोद्धरणं । }\edtext{\skt{वालुकाप्रक्षेपणं}}{\Afootnote{A1938 H; \skt{बालुकाक्षेपणं}\ N}}\skt{ }\edtext{\edtext{\skt{सिक्यावलम्वनञ्चेति}}{\lemma{\skt{सिक्यावलम्}}\Afootnote{H; \skt{शिक्यावलम्}\ A1938 N}}}{\lemma{\skt{वनञ्}}\Afootnote{H; \skt{बनं}\ A1938; \skt{बणञ्}\ N}}\skt{ ॥}\pend

\pstart \ledsidenote{\passref{1.45.18}}\skt{पञ्च निक्षेपणानि भवन्ति । तद्यथा । फलकं त्र्यष्टकं }\edtext{\skt{मुञ्जावलयं}}{\Afootnote{H; \skt{मुञ्जवलय,}\ A1938; \skt{मुञ्जवलयं}\ N}}\skt{ }\edtext{\edtext{\skt{दकमञ्चिकाशिक्यकञ्चेति}}{\lemma{\skt{दकमञ्}}\Afootnote{N H; \skt{उदकमञ्}\ A1938}}}{\lemma{\skt{चिकाशिक्यकञ्}}\Afootnote{N H; \skt{चिका, शिक्यं}\ A1938}}\skt{ ।}\pend

\pstart \ledsidenote{\passref{1.45.12a}}\edtext{\skt{नागपुष्पचम्पकोत्पलपाटलाप्रभृतिभिश्चाधिवासनमिति}}{\lemma{\skt{नागपुष्पचम्पकोत्पलपाटलाप्रभृतिभिश्}}\Afootnote{H; \skt{नागपुष्पचकोत्पलपाटलाप्रभृतिभिश्}\ N}}\skt{ ॥}\pend

\pstart\versepar \ledsidenote{\passref{1.45.20}}\edtext{\skt{सुगन्धविस्पष्टरसं}}{\Afootnote{N H; \skt{निर्गन्धमव्यक्तरसं}\ A1938}}\skt{ }\edtext{\edtext{\skt{सुशीतन्तर्षनाशनं}}{\lemma{\skt{सुशीतन्}}\Afootnote{N H; \skt{तृष्णाघ्नं शुचि शी}\ A1938}}}{\lemma{\skt{तर्षनाशनं}}\Afootnote{H; \skt{तलम् ।}\ A1938; \skt{तर्षनासनं ।}\ N}}\skt{ । }\pend
\pstart\versepar \skt{अच्छं }\edtext{\skt{लघु च}}{\Afootnote{A1938 H; \skt{लघुञ्च}\ N}}\skt{ हृद्यञ्च तोयङ्गुणवदुच्यते ॥}\pend

\pstart \ledsidenote{\passref{1.45.21}}\skt{तत्र नद्यः पश्चिमाभिमुखाः पथ्या }\edtext{\skt{लघूदकत्वात् । पूर्व्वाभिमुखा न}}{\Afootnote{N H; \skt{लघूदकत्वात् ; पूर्वाभिमुखास्तु न}\ A1938}}\skt{ प्रशस्यन्ते गुरूदकत्वात् । दक्षिणाभिमुखा }\edtext{\skt{नातिदोषला}}{\Afootnote{N H; \skt{नातिदोषलाः,}\ A1938}}\skt{ }\edtext{\skt{भवन्ति}}{\Afootnote{N H; \emph{om.}\ A1938}}\skt{ । }\edtext{\skt{साधारणत्वात् । तत्र}}{\Afootnote{A1938 H; \skt{साधारध्वात्वात् । तत्र}\ N}}\skt{ सह्यप्रभवाः कुष्ठञ्जनयन्ति । }\edtext{\skt{विन्ध्यप्रभवा}}{\Afootnote{H; \skt{विन्ध्यप्रभवाः}\ A1938; \skt{विन्ध्यप्रभावा}\ N}}\skt{ कुष्ठं }\edtext{\skt{हृद्रोगञ्च}}{\Afootnote{N H; \skt{पाण्डुरोगं च,}\ A1938}}\skt{ । }\edtext{\skt{मलयप्रभवाः}}{\Afootnote{A1938 H; \skt{मलयप्रभवा}\ N}}\skt{ }\edtext{\skt{क्रिमीन् । महेन्द्रप्रभवाः}}{\Afootnote{H; \skt{कृमीन्, महेन्द्रप्रभवाः}\ A1938; \skt{क्रिमीन् । महेन्द्रप्रभवा}\ N}}\skt{ श्लीपदोदराणि । हिमवत्प्रभवाः । }\edtext{\skt{गलगण्डं}}{\Afootnote{N\textsuperscript{pc} H; \skt{हृद्रोगश्वयथुशिरोरोगश्लीपदगलगण्}\ A1938; \skt{गलग}\ N\textsuperscript{ac}}}\skt{ }\edtext{\skt{क्वचिद्गण्डमालां}}{\Afootnote{N H; \skt{डा}\ A1938}}\skt{ । }\edtext{\edtext{\edtext{\edtext{\skt{प्राच्यावन्त्यापरावन्त्याश्चार्शांस्युपजनयन्ति}}{\lemma{\skt{प्राच्यावन्त्य्}}\Afootnote{N H; \skt{न्, प्राच्यावन्त्या}\ A1938}}}{\lemma{\skt{आपरावन्त्याश्}}\Afootnote{N H; \skt{अपरावन्त्याश्}\ A1938}}}{\lemma{\skt{चार्शांस्य्}}\Afootnote{H; \skt{चार्शां}\ A1938; \skt{वार्स्वांस्य्}\ N}}}{\lemma{\skt{उपजनयन्ति}}\Afootnote{N H; \skt{स्युपजनयन्ति,}\ A1938}}\skt{ । }\edtext{\skt{तत्र}}{\Afootnote{N H; \emph{om.}\ A1938}}\skt{ }\edtext{\skt{पारियात्रप्रभवा}}{\Afootnote{N H; \skt{पारियात्रप्रभवाः}\ A1938}}\skt{ }\edtext{\skt{वलारोग्यकरा}}{\Afootnote{N H; \skt{पथ्या बलारोग्यकर्य}\ A1938}}\skt{ इति ॥}\pend

\pstart\versepar \ledsidenote{\passref{1.45.22}}\skt{नद्यः }\edtext{\skt{शीघ्रवहा}}{\Afootnote{A1938 H; \skt{सीघ्रवहा}\ N}}\skt{ लघ्व्यः प्रोक्ता याश्चामलोदकाः । }\pend
\pstart\versepar \skt{}\edtext{\skt{गुर्व्व्यः}}{\Afootnote{A1938 H; \skt{गुर्व्य}\ N}}\skt{ }\edtext{\skt{शैवालकलुषजलौद्यामन्दगाश्च}}{\Afootnote{H; \skt{शैवालसञ्छन्नाः कलुषा मन्दगाश्च}\ A1938; \skt{सेवालकलुषजलौघा मन्दगाश्च}\ N}}\skt{ याः ॥}\pend

\pstart \ledsidenote{\passref{1.45.24}}\skt{तत्र सर्व्वेषां }\edtext{\edtext{\skt{भौमानाम्प्रत्यूषसि}}{\lemma{\skt{भौमानाम्}}\Afootnote{N H; \skt{एव भौमानां}\ A1938}}}{\lemma{\skt{प्रत्यूषसि}}\Afootnote{N H; \skt{ग्रहणं प्रत्यूषसि,}\ A1938}}\skt{ }\edtext{\skt{ग्रहणं}}{\Afootnote{N H; \emph{om.}\ A1938}}\skt{ । तत्र हि }\edtext{\edtext{\skt{शैत्यलाघवमधिकम्भवति}}{\lemma{\skt{शैत्यलाघवम्}}\Afootnote{\skt{अमलत्वं शैत्यं चा}\ A1938; \skt{सैत्यलाघवम्}\ N; \skt{शैत्यलाघवम्}\ H}}}{\lemma{\skt{अधिकम्}}\Afootnote{N H; \skt{धिकं}\ A1938}}\skt{ । }\edtext{\edtext{\skt{शैत्यञ्चापरो}}{\lemma{\skt{शैत्यञ्}}\Afootnote{\skt{स एव}\ A1938; \skt{सैत्यं}\ N; \skt{शैत्यञ्}\ H}}}{\lemma{\skt{चापरो}}\Afootnote{H; \skt{चापां परो}\ A1938; \skt{चापापरो}\ N}}\skt{ गुण इति ॥}\pend

\pstart\versepar \ledsidenote{\passref{1.45.25}}\skt{दिवार्क्ककिरणैर्ज्जुष्टं }\edtext{\edtext{\edtext{\skt{जुष्तमिन्दुकरैर्न्निशि}}{\lemma{\skt{जुष्तम्}}\Afootnote{N H; \emph{om.}\ A1938}}}{\lemma{\skt{इन्दुकरैर्}}\Afootnote{N H; \emph{om.}\ A1938}}}{\lemma{\skt{न्निशि}}\Afootnote{N H; \skt{निशा}\ A1938}}\skt{ । }\pend
\pstart\versepar \skt{}\edtext{\skt{अरूक्षमनभिष्यन्दि}}{\lemma{\skt{अरूक्षम्}}\Afootnote{N H; \skt{यामिन्दुरश्मिभिः । अरूक्षम्}\ A1938}}\skt{ तत्तुल्यं }\edtext{\skt{गगनाम्बुना}}{\Afootnote{A1938 H; \skt{गगणाम्बुरा ॥}\ N}}\skt{ ॥}\pend

\pstart\versepar \ledsidenote{\passref{1.45.26}}\edtext{\skt{गगनाम्बु}}{\Afootnote{A1938 H; \skt{गगणाम्बुण}\ N\textsuperscript{ac}; \skt{गगणाम्बु}\ N\textsuperscript{pc}}}\skt{ त्रिदोषघ्नं गृहीतं यत्सुभाजने । }\pend
\pstart\versepar \skt{बल्यं रसायनं }\edtext{\skt{शीतं}}{\Afootnote{N H; \skt{मेध्यं}\ A1938}}\skt{ }\edtext{\skt{मात्रापेक्षन्ततः}}{\lemma{\skt{मात्रापेक्षन्}}\Afootnote{H; \skt{पात्रापेक्षि}\ A1938; \skt{मात्रापेक्ष्यन्}\ N}}\skt{ परं ॥}\pend

\pstart \ledsidenote{\passref{1.45.27ab}}\edtext{\skt{रक्तोघ्नं}}{\Afootnote{N H; \skt{रक्षोघ्नं}\ A1938}}\skt{ }\edtext{\skt{शीतलं}}{\Afootnote{A1938 H; \skt{सीतलं}\ N}}\skt{ ह्लादि }\edtext{\skt{ज्वरदाहस्तृषापहं}}{\Afootnote{H; \skt{ज्वरदाहविषापहम् ।}\ A1938; \skt{ज्वरदाहस्तृषापहं ।}\ N}}\skt{ ।}\pend

\pstart\versepar \ledsidenote{\passref{1.45.28}}\edtext{\skt{मूर्च्छापित्तोष्मदाहेषु}}{\Afootnote{N H; \skt{मूर्च्छापित्तोष्णदाहेषु}\ A1938}}\skt{ }\edtext{\skt{विषे}}{\Afootnote{A1938 H; \emph{om.}\ N}}\skt{ रक्ते }\edtext{\skt{मादत्यये}}{\Afootnote{N H; \skt{मदात्यये ।}\ A1938}}\skt{ ॥ }\pend
\pstart\versepar \skt{भ्रमक्लमपरीतेषु तमके वमथौ तथा ।}\pend

\pstart\versepar \ledsidenote{\passref{1.45.29}}\edtext{\skt{ऊर्द्धगे}}{\Afootnote{N H; \skt{ऊर्ध्वगे}\ A1938}}\skt{ रक्तपित्ते च }\edtext{\skt{शीतमम्भः}}{\Afootnote{A1938 H; \skt{सीतमम्भः}\ N}}\skt{ प्रशस्यते ॥ }\pend
\pstart\versepar \skt{}\edtext{\skt{पार्श्वशूलप्रतिश्याये}}{\Afootnote{\skt{पार्श्वशूले प्रतिश्याये}\ A1938; \skt{प्रार्श्वसूले प्रतिस्याये}\ N; \skt{प्रार्श्वशूल प्रतिश्याये}\ H}}\skt{ वातरोगे गलग्रहे ।}\pend

\pstart\versepar \ledsidenote{\passref{1.45.30}}\edtext{\skt{आध्माते}}{\Afootnote{N H; \skt{आध्माने}\ A1938}}\skt{ }\edtext{\skt{स्तिमिते}}{\Afootnote{A1938 H; \skt{तिमिते}\ N}}\skt{ }\edtext{\skt{काष्ठे}}{\Afootnote{N H; \skt{कोष्ठे}\ A1938}}\skt{ }\edtext{\skt{सद्यःशुद्धे}}{\Afootnote{A1938 H; \skt{सद्यशुद्धे}\ N}}\skt{ नवज्वरे ॥ }\pend
\pstart\versepar \skt{हिक्कायां }\edtext{\skt{स्नेहपीते च}}{\Afootnote{A1938 H; \skt{स्नेहपीते}\ N}}\skt{ शीताम्बु परिवर्ज्जयेत् ।}\pend

\pstart \ledsidenote{\passref{1.45.45cd}}\skt{अरोचके }\edtext{\skt{प्रतिश्याये}}{\Afootnote{A1938 H; \skt{प्रतिस्याये}\ N}}\skt{ }\edtext{\skt{प्रमेहश्च}}{\Afootnote{H; \skt{प्रसेके श्व}\ A1938; \skt{प्रमेहश्व}\ N}}\skt{ यथौ }\edtext{\skt{तथा}}{\Afootnote{N H; \skt{क्षये ॥४५॥}\ A1938}}\skt{ ॥}\pend

\pstart\versepar \ledsidenote{\passref{1.45.46}}\edtext{\skt{मन्देऽग्नावुदरे}}{\lemma{\skt{ऽग्नावुदरे}}\Afootnote{A1938 H; \skt{ग्नावुदरे}\ N}}\skt{ }\edtext{\skt{कोष्ठे}}{\Afootnote{N H; \skt{कुष्ठे}\ A1938}}\skt{ ज्वरे }\edtext{\skt{नेत्रामयेषु च}}{\Afootnote{N H; \skt{नेत्रामये तथा ।}\ A1938}}\skt{ । }\pend
\pstart\versepar \skt{व्रणे च मधुमेहे च पानीयम्मन्दमाचरेत् ॥}\pend

\pstart \ledsidenote{\passref{1.45.27cd}}\edtext{\skt{चन्द्रकान्तभवम्वारि}}{\lemma{\skt{चन्द्रकान्तभवम्}}\Afootnote{H; \skt{चन्द्रकान्तोद्भवं}\ A1938; \skt{चन्द्राकान्ताभवम्}\ N}}\skt{ पित्तघ्नम्विमलं स्मृतम् ।}\pend

\pstart \ledsidenote{\passref{1.45.33cd}}\edtext{\skt{सक्षारम्प्रायसः}}{\Afootnote{N H; \skt{सक्षारं पित्तलं}\ A1938}}\skt{ }\edtext{\skt{कौपन्नादेयं}}{\lemma{\skt{नादेयं}}\Afootnote{N H; \skt{श्लेष्मघ्नं}\ A1938}}\skt{ }\edtext{\edtext{\skt{कीर्त्तितन्तथा॥}}{\lemma{\skt{कीर्त्तितन्}}\Afootnote{N H; \skt{दीपनं}\ A1938}}}{\lemma{\skt{तथा}}\Afootnote{N H; \skt{लघु ॥३३॥}\ A1938}}\pend

\pstart \ledsidenote{\passref{1.45.32cd}}\edtext{\skt{ताडाकम्वातलं}}{\lemma{\skt{ताडाकम्}}\Afootnote{N H; \skt{ताडागं}\ A1938}}\skt{ }\edtext{\skt{रूक्षं}}{\Afootnote{N H; \emph{om.}\ A1938}}\skt{ }\edtext{\edtext{\skt{सारसञ्चैव}}{\lemma{\skt{सारसञ्}}\Afootnote{N H; \emph{om.}\ A1938}}}{\lemma{\skt{चैव}}\Afootnote{N H; \emph{om.}\ A1938}}\skt{ }\edtext{\skt{तादृशं}}{\Afootnote{N H; \skt{स्वादु कषायं कटुपाकि च ॥३२॥}\ A1938}}\skt{ ।}\pend

\pstart \ledsidenote{\passref{1.45.32.1}}\edtext{\edtext{\skt{औत्समस्मसमाश्लेषादुष्णं}}{\lemma{\skt{अस्मसम्}}\Afootnote{H; \skt{अश्मसम्}\ N}}}{\lemma{\skt{आश्लेषादुष्णं}}\Afootnote{H; \skt{आश्लेषादुष्णं}\ N}}\skt{ पित्तेन शस्यते ॥}\pend

\pstart \ledsidenote{\passref{1.45.35.1}}\skt{सर्व्वदा सर्व्वदोषेषु }\edtext{\skt{पथ्यम्प्रास्रावणम्पयः}}{\lemma{\skt{प्रास्रावणम्}}\Afootnote{H; \skt{प्राश्रावणम्}\ N}}\skt{ ।}\pend

\pstart \ledsidenote{\passref{1.45.35ab}}\edtext{\edtext{\skt{अविदाह्युद्भिजं}}{\lemma{\skt{अविदाह्य्}}\Afootnote{N H; \skt{मधुरं पित्तशमनमविदाह्यौ}\ A1938}}}{\lemma{\skt{उद्भिजं}}\Afootnote{N H; \skt{द्भि}\ A1938}}\skt{ }\edtext{\edtext{\edtext{\skt{तोयम्पित्तघ्नम्मधुरं}}{\lemma{\skt{तोयम्}}\Afootnote{N H; \skt{दं}\ A1938}}}{\lemma{\skt{पित्तघ्नम्}}\Afootnote{N H; \emph{om.}\ A1938}}}{\lemma{\skt{मधुरं}}\Afootnote{H; \emph{om.}\ A1938; \skt{मधुर}\ N}}\skt{ स्मृतं ॥}\pend

\pstart \ledsidenote{\passref{1.45.39cd}}\edtext{\skt{कफमेदोनिलहरं}}{\Afootnote{N H; \skt{कफमेदोऽनिलामघ्नं}\ A1938}}\skt{ दीपनं }\edtext{\skt{बस्तिशोधनं}}{\Afootnote{A1938; \skt{वस्तिसोधनं ।}\ N; \skt{वस्तिशोधनं ।}\ H}}\skt{ ।}\pend

\pstart \ledsidenote{\passref{1.45.40ab}}\edtext{\skt{श्वासकासज्वरहरम्पथ्यमुष्णोदकं}}{\lemma{\skt{श्वासकासज्वरहरम्}}\Afootnote{A1938 H; \skt{श्वासकासज्वरंहर}\ N}}\skt{ सदा ॥}\pend

\pstart\versepar \ledsidenote{\passref{1.45.40.1}}\skt{दीपनी पाचनी लघ्वी पथ्या }\edtext{\skt{बस्तिविशोधनी}}{\Afootnote{\skt{बस्तिविसोधनी ।}\ N; \skt{वस्तिविशोधनी ।}\ H}}\skt{ । }\pend
\pstart\versepar \skt{}\edtext{\skt{वातानुलोमनी}}{\Afootnote{H; \skt{वातानुलोमिआनी}\ N\textsuperscript{ac}; \skt{वातानुलोमिनी}\ N\textsuperscript{pc}}}\skt{ }\edtext{\skt{पेयाक्षत्पिपासहरा}}{\Afootnote{H; \skt{पेया क्षुत्पिपासहरा}\ N}}\skt{ }\edtext{\skt{स्मृताः}}{\Afootnote{N\textsuperscript{pc} H; \skt{स्मृतिआः ।}\ N\textsuperscript{ac}}}\skt{ ॥}\pend

\pstart \ledsidenote{\passref{1.45.34cd}}\edtext{\skt{कफघ्नो}}{\Afootnote{N H; \skt{कफघ्नं}\ A1938}}\skt{ }\edtext{\skt{दीपनो}}{\Afootnote{N H; \skt{दीपनं}\ A1938}}\skt{ }\edtext{\skt{हृद्यः}}{\Afootnote{N H; \skt{हृद्यं}\ A1938}}\skt{ }\edtext{\edtext{\edtext{\skt{शुद्धानाम्व्रणिनामपि}}{\lemma{\skt{शुद्धानाम्}}\Afootnote{N H; \emph{om.}\ A1938}}}{\lemma{\skt{व्रणिनाम्}}\Afootnote{H; \emph{om.}\ A1938; \skt{व्राणिना}\ N\textsuperscript{ac}; \skt{व्राणिनाम्}\ N\textsuperscript{pc}}}}{\lemma{\skt{अपि}}\Afootnote{N\textsuperscript{pc} H; \skt{लघु प्रस्रवणोद्भवम् ॥३४॥}\ A1938; \skt{पिमपि ।}\ N\textsuperscript{ac}}}\skt{ ।}\pend

\pstart \ledsidenote{\passref{1.45.34x}}\skt{ज्ञेयः पथ्यतमश्चापि मुद्गयूषः कृताकृतः ॥}\pend

\pstart\versepar \ledsidenote{\passref{1.45.b}}\skt{प्रीणनः }\edtext{\skt{प्राणजननः}}{\Afootnote{N\textsuperscript{pc} H; \skt{प्राणंतुजननः}\ N\textsuperscript{ac}}}\skt{ श्वासकासक्षयापः । }\pend
\pstart\versepar \skt{वातहन्ताश्रमहरो हृद्यो }\edtext{\skt{मांसरसः}}{\Afootnote{H; \skt{मंसरसः}\ N}}\skt{ स्मृतः ॥}\pend

\pstart\versepar \ledsidenote{\passref{1.45.c}}\edtext{\skt{खलाः}}{\Afootnote{H; \skt{खला}\ N}}\skt{ खलयवाग्वश्च }\edtext{\skt{रागषाडवषट्टकाः}}{\Afootnote{H; \skt{रागसाडवषट्टकाः ।}\ N}}\skt{ । }\pend
\pstart\versepar \skt{एवमादीनि चान्यानि क्रियन्ते वैद्यवाक्यतः ॥}\pend

\pstart\versepar \ledsidenote{\passref{1.45.d}}\edtext{\skt{यदाकारणमासाद्य}}{\lemma{\skt{यदाकारणम्}}\Afootnote{H; \skt{यदा कारणम्}\ N}}\skt{ भोक्तॄणां च्छन्दतोऽपि वा । }\pend
\pstart\versepar \skt{}\edtext{\skt{अनेकद्रव्ययोनित्वाच्छास्त्र-तस्ताम्विनिर्द्दिशेत्}}{\lemma{\skt{अनेकद्रव्ययोनित्वाच्छास्त्र-तस्ताम्}}\Afootnote{H; \skt{अनेकद्रव्ययोनित्वात्छास्त्रतस्तां}\ N}}\skt{ ॥}\pend

\pstart \ledsidenote{\passref{1.45.43cd}}\skt{स्निग्धं स्वादु }\edtext{\skt{रसं}}{\Afootnote{N H; \skt{हिमं}\ A1938}}\skt{ हृद्यं }\edtext{\edtext{\skt{वृंहणम्वलवत्तरं}}{\lemma{\skt{वृंहणम्}}\Afootnote{\skt{दीपनं}\ A1938; \skt{बृंहणम्}\ N; \skt{… वृंहणम्}\ H}}}{\lemma{\skt{वलवत्तरं}}\Afootnote{H; \skt{बस्तिशोधनम् ॥४३॥}\ A1938; \skt{बलवत्तरं ।}\ N}}\skt{ ।}\pend

\pstart \ledsidenote{\passref{1.45.44ab}}\edtext{\skt{वृष्यम्पित्तपिपासाघ्नं}}{\lemma{\skt{पित्तपिपासाघ्नं}}\Afootnote{A1938 H; \skt{पित्तपिपासघ्नं}\ N}}\skt{ }\edtext{\skt{नालिकेरोदकङ्गुरुः}}{\lemma{\skt{गुरुः}}\Afootnote{N H; \skt{गुरु ।}\ A1938}}\skt{ ॥}\pend

\pstart\versepar \ledsidenote{\passref{1.45.47}}\skt{गव्यमाजन्तथा }\edtext{\skt{चोष्ट्रमाविकं}}{\Afootnote{H; \skt{चौष्ट्रमाविकं}\ A1938; \skt{चोष्ट्रमाविकं}\ N}}\skt{ माहिषञ्च यत् । }\pend
\pstart\versepar \skt{अश्वायाश्चैव नार्याश्च }\edtext{\edtext{\skt{नाग्योश्चैव}}{\lemma{\skt{नाग्योश्}}\Afootnote{\skt{करेणूनां}\ A1938; \skt{नाग्याश्}\ N; \skt{… नाग्योश्}\ H}}\skt{ तु}}{\lemma{\skt{चैव तु}}\Afootnote{N H; \skt{च यत्प}\ A1938}}\skt{ }\edtext{\skt{यत्स्मृतं}}{\Afootnote{N H; \skt{यः ॥४७॥}\ A1938}}\skt{ ॥}\pend

\pstart \ledsidenote{\passref{1.45.48ab}}\edtext{\skt{तत्वनेकौषधिरसः}}{\Afootnote{N H; \skt{तत्त्वनेकौषधिरस}\ A1938}}\skt{ }\edtext{\edtext{\skt{प्रसादक्षीरताङ्गतः}}{\lemma{\skt{प्रसादक्षीरताङ्}}\Afootnote{N H; \skt{प्रसादं प्राणदं}\ A1938}}}{\lemma{\skt{गतः}}\Afootnote{N H; \skt{गुरु ।}\ A1938}}\skt{ ।}\pend

\pstart \ledsidenote{\passref{1.45.48ef}}\edtext{\skt{सर्व्वप्राणभृतान्तस्मात्सात्म्यं}}{\lemma{\skt{तस्मात्सात्म्यं}}\Afootnote{A1938 H; \skt{तस्मात्सात्म्य}\ N}}\skt{ }\edtext{\edtext{\skt{क्षीरमिहोच्यते}}{\lemma{\skt{क्षीरम्}}\Afootnote{A1938; \skt{क्षीर}\ N H}}}{\lemma{\skt{इहोच्यते}}\Afootnote{A1938; \skt{मिहोच्यते}\ N H}}\skt{ ॥}\pend

\pstart \ledsidenote{\passref{1.45.49}}\edtext{\skt{गव्यन्तु}}{\Afootnote{\skt{तत्र सर्वमेव क्षीरं प्राणिनामप्रति}\ A1938; \skt{गव्यन्तु शित}\ N; \skt{… गव्यन्तु}\ H}}\skt{ }\edtext{\edtext{\skt{शितस्निग्धमधुरमविदाहि}}{\lemma{\skt{शितस्निग्धमधुरम्}}\Afootnote{H\textsuperscript{pc}; \skt{षिद्धं जातिसात्म्या}\ A1938; \skt{स्निग्धमधुरम्}\ N; \skt{शितस्निमधुरम्}\ H\textsuperscript{ac}}}}{\lemma{\skt{अविदाहि}}\Afootnote{N H; \emph{om.}\ A1938}}\skt{ }\edtext{\skt{वातपित्तश्लेष्मशोणितमानसेषु}}{\Afootnote{N H; \skt{त्, वातपित्तशोणितमानसेष्व्}\ A1938}}\skt{ }\edtext{\skt{विकारेष्वविरुद्धं}}{\lemma{\skt{विकारेष्व्}}\Afootnote{N H; \skt{अपि विकारेष्व्}\ A1938}}\skt{ । }\edtext{\skt{जीर्ण्णज्वरकासशोफक्षयरक्तपित्तगुल्मोदरमूर्च्छाभ्रममदपिपासापाण्डुरोगार्श}}{\Afootnote{H\textsuperscript{ac}; \skt{जीर्णज्वरकासश्वासशोषक्षयगुल्मोन्मादोदरमूर्च्छाभ्रममददाहपिपासाहृद्बस्तिदोषपाण्डुरोगग्रहणीदोषार्शः}\ A1938; \skt{जीर्ण्णज्वरकासशोफक्षयरक्तपित्तगुल्मौ दरमूर्च्छाभ्रममदपिपासापाण्डुरोगार्श}\ N\textsuperscript{ac}; \skt{जीर्ण्णज्वरकासशोफक्षयरक्तपित्तगुल्मोदरमूर्च्छाभ्रममदपिपासापाण्डुरोगार्श}\ N\textsuperscript{pc}; \skt{जीर्ण्णज्वरकासशोफक्षयरक्तपित्तगुल्मोदरमूर्च्छाभ्रममदपिपासापाण्डुरोगा}\ H\textsuperscript{pc}}}\skt{ }\edtext{\edtext{\skt{उदावर्त्तातीसारयोनिरोगगर्भ्भास्रावश्रमक्लमहरम्वल्यं}}{\lemma{\skt{उदावर्त्तातीसारयोनिरोगगर्भ्भास्रावश्रमक्लमहरम्}}\Afootnote{H\textsuperscript{ac}; \skt{शूलोदावर्तातिसारप्रवाहिकायोनिरोगगर्भास्रावरक्तपित्तश्रमक्लमहरं,}\ A1938; \skt{उदावर्त्तातीसारयोणिरोगगर्भास्रावकश्रमक्लमहरम्}\ N\textsuperscript{ac}; \skt{उदावर्त्तातीसारयोणिरोगगर्भास्रावश्रमक्लमहरम्}\ N\textsuperscript{pc}; \skt{र्श उदावर्त्तातीसारयोनिरोगगर्भ्भास्रावश्रमक्लमहरम्}\ H\textsuperscript{pc}}}}{\lemma{\skt{वल्यं}}\Afootnote{\skt{पाप्मापहं बल्यं वृष्यं}\ A1938; \skt{बल्यं वृष्यं}\ N; \skt{व}\ H\textsuperscript{ac}; \emph{om.}\ H\textsuperscript{pc}}}\skt{ }\edtext{\skt{वृष्यं रसायनं}}{\Afootnote{\skt{वाजीकरणं रसायनं}\ A1938; \skt{रसायनं}\ N; \skt{ल्यं वृष्यं रसायनं}\ H\textsuperscript{ac}; \skt{वल्यं वृष्यं रसायनं}\ H\textsuperscript{pc}}}\skt{ मेध्यं }\edtext{\skt{व्याजीककरणं}}{\Afootnote{H; \emph{om.}\ A1938; \skt{व्याजीककरण}\ N}}\skt{ }\edtext{\edtext{\edtext{\edtext{\edtext{\skt{सन्धानमास्थापनमायुष्यम्वमनम्विरेचनञ्चेति}}{\lemma{\skt{सन्धानम्}}\Afootnote{H; \skt{व}\ A1938; \skt{न्सन्धानम्}\ N}}}{\lemma{\skt{आस्थापनम्}}\Afootnote{N H; \skt{यःस्थापनम्}\ A1938}}}{\lemma{\skt{वमनम्}}\Afootnote{H; \skt{जीवनं}\ A1938; \emph{om.}\ N}}}{\lemma{\skt{विरेचनञ्}}\Afootnote{H; \skt{बृंहणं सन्}\ A1938; \emph{om.}\ N}}}{\lemma{\skt{चेति}}\Afootnote{H; \skt{धानं}\ A1938; \emph{om.}\ N}}\skt{ ॥ }\edtext{\skt{स्थापनं}}{\Afootnote{\skt{वमनविरेचनास्थापनं}\ A1938; \skt{वमनम्विरेच}\ N; \skt{…थापनं}\ H}}\skt{ }\edtext{\edtext{\skt{तुल्यगुणत्वाच्चौजसो वर्धनं बालवृद्धक्षतक्षीणानां क्षुद्व्यवायव्यायामकर्शितानां च }\edtext{\skt{पथ्यतमम्}}{\Afootnote{A1938 H; \skt{ति ॥}\ N}}\skt{ ॥}}{\lemma{\skt{स्थापनं… अकर्शितानां}}\Afootnote{\emph{om.}\ N\textsuperscript{ac}}}}{\lemma{\skt{स्थापनं… अकर्शितानां}}\Afootnote{\emph{om.}\ N\textsuperscript{pc}}}\pend

\pstart\versepar \ledsidenote{\passref{1.45.52}}\skt{दीपनं लघु सङ्ग्राहि }\edtext{\edtext{\skt{विशेषञ्चात्र}}{\lemma{\skt{विशेषञ्}}\Afootnote{\skt{श्वासका}\ A1938; \skt{विसेषञ्}\ N; \skt{… विशेषञ्}\ H}}\skt{ मे}}{\lemma{\skt{चात्र मे}}\Afootnote{N H; \skt{सास्रपि}\ A1938}}\skt{ }\edtext{\skt{शृणु}}{\Afootnote{H; \skt{त्तनु}\ A1938; \skt{सृणु ।}\ N}}\skt{ ॥ }\pend
\pstart\versepar \skt{}\edtext{\edtext{\skt{अजानामल्पकायत्वात्कटुतिक्तनिषेवणात्}}{\lemma{\skt{अजानाम्}}\Afootnote{N H; \skt{त् । अजानाम्}\ A1938}}}{\lemma{\skt{अल्पकायत्वात्कटुतिक्तनिषेवणात्}}\Afootnote{A1938 H; \skt{अल्पकायत्वात्कटुतिक्तनिसेवनात् ।}\ N}}\skt{ ।}\pend

\pstart\versepar \ledsidenote{\passref{1.45.53}}\edtext{\skt{नात्यम्बुपानाद्व्यायामात्सर्व्वदोषहरम्पयः}}{\lemma{\skt{नात्यम्बुपानाद्व्यायामात्सर्व्वदोषहरम्}}\Afootnote{N; \skt{नात्यम्बुपानाद्व्यायामात्सर्वव्याधिहरं}\ A1938; \skt{नात्यम्बुपानाद्व्यायामात्सर्व्वदोषहरम्}\ H}}\skt{ ॥ }\pend
\pstart\versepar \skt{रूक्षोष्णं }\edtext{\skt{लवणङ्किञ्चिदौष्ट्रं}}{\lemma{\skt{किञ्चिदौष्ट्रं}}\Afootnote{A1938 H; \skt{किञ्चिदौष्टं}\ N}}\skt{ }\edtext{\skt{स्वादुरसं}}{\Afootnote{A1938 N\textsuperscript{pc} H; \skt{स्वादुरघसं}\ N\textsuperscript{ac}}}\skt{ }\edtext{\skt{लघु}}{\Afootnote{A1938 H; \skt{लघुं ।}\ N}}\skt{ ।}\pend

\pstart\versepar \ledsidenote{\passref{1.45.54}}\skt{शोफगुल्मोदरार्शोघ्नं }\edtext{\skt{क्रिमिकुष्ठकफापहं}}{\Afootnote{N H; \skt{कृमिकुष्ठविषापहम्}\ A1938}}\skt{ ॥ }\pend
\pstart\versepar \skt{आविकं मधुरं स्निग्धं गुरु }\edtext{\skt{पित्तकफपहं}}{\Afootnote{H; \skt{पित्तकफावहम् ॥५४॥}\ A1938; \skt{पित्तकफापहं ।}\ N}}\skt{ ।}\pend

\pstart\versepar \ledsidenote{\passref{1.45.55}}\skt{पथ्यं केवलवातेषु कासे }\edtext{\skt{चानिलसम्भवे}}{\Afootnote{A1938 H; \skt{वाऽनिलसम्भवे ॥}\ N}}\skt{ ॥ }\pend
\pstart\versepar \skt{महाभिष्यन्दि मधुरं }\edtext{\skt{माहिषम्वह्निशादनं}}{\lemma{\skt{वह्निशादनं}}\Afootnote{H; \skt{वह्निनाशनम् ॥५५॥}\ A1938; \skt{वह्निसादनं ।}\ N}}\skt{ ।}\pend

\pstart\versepar \ledsidenote{\passref{1.45.56}}\skt{निद्राकरं }\edtext{\skt{शीतकरं}}{\Afootnote{N H; \skt{शीततरं}\ A1938}}\skt{ गव्यात्स्निग्धतरं }\edtext{\skt{गुरु}}{\Afootnote{A1938 H; \skt{गुरुः ॥}\ N}}\skt{ ॥ }\pend
\pstart\versepar \skt{}\edtext{\skt{उष्णमैकसफं}}{\Afootnote{H; \skt{उष्णमैकशफं}\ A1938; \skt{उष्णमेकशफं}\ N}}\skt{ बल्यं शाखावातहरं पयः ।}\pend

\pstart\versepar \ledsidenote{\passref{1.45.57}}\edtext{\skt{मधुराम्लग्सं}}{\Afootnote{H; \skt{मधुराम्लरसं}\ A1938; \skt{मधुराम्ब्लरसं}\ N}}\skt{ रूक्षं लवणानुरसं }\edtext{\skt{लघु}}{\Afootnote{A1938 H; \skt{लघुः ॥}\ N}}\skt{ ॥ }\pend
\pstart\versepar \skt{}\edtext{\skt{नार्यास्तु}}{\Afootnote{N H; \skt{नार्यास्तु}\ A1938}}\skt{ }\edtext{\skt{मध्रुअं}}{\Afootnote{N H; \skt{मधुरं}\ A1938}}\skt{ स्तन्यङ्कषायानुरसं }\edtext{\skt{हिमं}}{\Afootnote{A1938 H; \skt{हिम ।}\ N}}\skt{ ।}\pend

\pstart\versepar \ledsidenote{\passref{1.45.58}}\edtext{\edtext{\skt{नस्याश्च्योतनयोः}}{\lemma{\skt{नस्याश्}}\Afootnote{N H; \skt{नस्या}\ A1938}}}{\lemma{\skt{च्योतनयोः}}\Afootnote{N H; \skt{श्च्योतनयोः}\ A1938}}\skt{ पथ्यञ्जीवनं लघु दीपनं ॥ }\pend
\pstart\versepar \skt{हस्तिन्या मधुरं वृष्यं कषायानुरसं }\edtext{\skt{गुरु}}{\Afootnote{A1938 H; \skt{गुरुः}\ N}}\skt{ ।}\pend

\pstart \ledsidenote{\passref{1.45.59ab}}\skt{स्निग्धं }\edtext{\skt{स्थैर्यकरं}}{\Afootnote{N H; \skt{स्थैर्यकरं}\ A1938}}\skt{ शीतं चक्षुष्यम्बलवर्द्धनं ॥}\pend

\pstart \ledsidenote{\passref{1.45.61cd}}\skt{पयोऽभिष्यन्दि }\edtext{\skt{गुर्वामं}}{\Afootnote{A1938 H; \skt{गुर्वाम}\ N}}\skt{ }\edtext{\skt{प्रायसः}}{\Afootnote{N H; \skt{प्रायशः}\ A1938}}\skt{ परिकीर्तितं ॥}\pend

\pstart\versepar \ledsidenote{\passref{1.45.62}}\skt{तदेवोक्तं लघुतरं }\edtext{\skt{मन्दाभिष्यन्दि च}}{\Afootnote{N H; \skt{अनभिष्यन्दि वै}\ A1938}}\skt{ }\edtext{\skt{श्रितं}}{\Afootnote{N H; \skt{शृतम्}\ A1938}}\skt{ ॥ }\pend
\pstart\versepar \skt{वर्जयित्वा }\edtext{\skt{स्त्रियाः}}{\Afootnote{A1938 H; \skt{स्त्रिया}\ N}}\skt{ }\edtext{\skt{स्तन्यम्ममेव}}{\lemma{\skt{मम्}}\Afootnote{N H; \skt{आमम्}\ A1938}}\skt{ हि }\edtext{\skt{तधितं}}{\Afootnote{N H; \skt{तद्धितम् ॥६२॥}\ A1938}}\skt{ ॥}\pend

\pstart \ledsidenote{\passref{1.45.63ab}}\edtext{\skt{धारोष्णन्गुणवत्क्षीरं}}{\Afootnote{H; \skt{धारोष्णं गुणवत्क्षीरं}\ A1938 N}}\skt{ }\edtext{\skt{विपरीतमतोऽन्यथा}}{\lemma{\skt{विपरीतमतो}}\Afootnote{H; \skt{विपरीतमतो}\ A1938 N}}\skt{ ॥}\pend

\pstart\versepar \ledsidenote{\passref{1.45.64}}\edtext{\edtext{\edtext{\skt{अनिष्टमम्लगन्धञ्च}}{\lemma{\skt{अनिष्टम्}}\Afootnote{N H; \skt{अनिष्ट}\ A1938}}}{\lemma{\skt{अम्लगन्धञ्}}\Afootnote{H; \skt{गन्धम्}\ A1938; \skt{अम्ल्बगन्धञ्}\ N}}}{\lemma{\skt{च}}\Afootnote{N H; \skt{अम्लं च}\ A1938}}\skt{ }\edtext{\edtext{\skt{विवर्ण्णम्विरसन्तु}}{\lemma{\skt{विवर्ण्णम्}}\Afootnote{A1938 H; \skt{विवर्ण्ण}\ N}}}{\lemma{\skt{विरसन्तु}}\Afootnote{N H; \skt{विरसं च}\ A1938}}\skt{ यत् । }\pend
\pstart\versepar \skt{वर्ज्यं सलवणं क्षीरं यच्च }\edtext{\skt{विग्रथितम्भवेत्}}{\lemma{\skt{भवेत्}}\Afootnote{N H; \skt{भवेत् ॥६४॥ इति क्षीरवर्गः ।}\ A1938}}\skt{ ॥}\pend

\pstart \ledsidenote{\passref{1.45.65}}\edtext{\skt{दधि तु}}{\Afootnote{\skt{अथ दधिवर्गः । दधि तु}\ A1938; \emph{om.}\ N H}}\skt{ }\edtext{\edtext{\skt{सृष्तमूत्रपुरीषं}}{\Afootnote{\skt{मधुरमम्लमत्यम्लं चेति; तत्कषायानुरसं स्निग्धमुष्णं पीनसविषमज्वरातिसारारोचकमूत्र}\ A1938; \skt{दधि तु सृष्टमूत्रपुरीषं}\ N; \skt{दधि तु सृष्तमूत्रपुरीष}\ H}}\skt{ }\edtext{\edtext{\skt{गुरूम्श्लेष्मभिष्यन्दि}}{\lemma{\skt{गुरूम्}}\Afootnote{\skt{गुर्वम्}\ N; \skt{ञ्गुरूम्}\ H}}}{\lemma{\skt{श्लेष्मभिष्यन्दि}}\Afootnote{H; \skt{ब्लमभिष्यन्दि}\ N}}\skt{ श्लेष्मपित्तशोफवर्द्धनं}}{\lemma{\skt{पुरीषं… श्लेष्मपित्तशोफव}}\Afootnote{\emph{om.}\ A1938}}\skt{ }\edtext{\skt{कार्श्यापहं}}{\Afootnote{A1938 H; \skt{कार्स्यापहं}\ N}}\skt{ }\edtext{\skt{रोचनम्मङ्गल्यञ्च}}{\lemma{\skt{रोचनम्}}\Afootnote{H; \skt{वृष्यं प्राणकरं}\ A1938; \skt{रोचकं}\ N}}\skt{ ॥}\pend

\pstart \ledsidenote{\passref{1.45.65.add1}}\edtext{\skt{तदुत}}{\Afootnote{H; \skt{तदुधृत}\ N}}\skt{ सारं ग्राह्यमनभिष्यन्दि च । सरः }\edtext{\skt{कफमेदशुक्रकृत्त्रिदोषम्मन्दजातं}}{\lemma{\skt{कफमेदशुक्रकृत्त्रिदोषम्}}\Afootnote{H; \skt{कफमेदशुक्रकृत् । तिदोषम्}\ N}}\skt{ ॥}\pend

\pstart \ledsidenote{\passref{1.45.84}}\edtext{\edtext{\edtext{\skt{तक्रङ्कषायानुरसमुष्णवीर्यमतीसारगरघ्नं}}{\lemma{\skt{तक्रङ्}}\Afootnote{N H; \skt{अथ तक्रवर्गः । तक्रं तु मधुरमम्लं}\ A1938}}}{\lemma{\skt{उष्णवीर्यम्}}\Afootnote{A1938 H; \skt{अम्ब्लमुष्णवीर्यम्}\ N}}\skt{ }\edtext{\edtext{\skt{लघुशुक्रवलासक्षयकरमर्शोघ्नञ्च}}{\lemma{\skt{लघुशुक्रवलासक्षयकरम्}}\Afootnote{H; \skt{लघुशुक्रवलासयक्षयकरम्}\ N}}}{\lemma{\skt{अर्शोघ्नञ्च}}\Afootnote{H; \skt{अर्सोघ्नञ्च ॥}\ N}}\skt{ ॥ दधिमण्डो विष्टम्भारुचिहरो }\edtext{\edtext{\skt{लघुतरश्चाग्निदीपनः}}{\lemma{\skt{लघुतरश्}}\Afootnote{N H; \skt{लघुरूक्षम्}\ A1938}}}{\lemma{\skt{चाग्निदीपनः}}\Afootnote{N H; \skt{अग्निदीपनं गरशोफातिसारग्रहणीपाण्डुरोगार्शःप्लीहगुल्मारोचकविषमज्वरतृष्णाच्छर्दिप्रसेकशूलभेदःश्लेष्मानिलहरं मधुरविपाकं हृद्यं मूत्रकृच्छ्रस्नेहव्यापत्प्रशमनमवृष्यं च ॥८४॥}\ A1938}}}{\lemma{\skt{अतीसारगरघ्नं… ल}}\Afootnote{\emph{om.}\ A1938}}\skt{ ।}\pend

\pstart\versepar \ledsidenote{\passref{1.45.89}}\edtext{\edtext{\skt{वातेऽम्लसैन्धवोपेतं}}{\lemma{\skt{वाते}}\Afootnote{A1938 H; \skt{वाते म्}\ N}}}{\lemma{\skt{ऽम्लसैन्धवोपेतं}}\Afootnote{H; \skt{ऽम्लं सैन्धवोपेतं,}\ A1938; \skt{ब्लं सेन्धवोपेतं}\ N}}\skt{ स्वादु पित्ते सशर्क्करं । }\pend
\pstart\versepar \skt{पिबेत्तक्रङ्कफे चापि }\edtext{\skt{सव्योषक्षारसंयुतं}}{\Afootnote{N H; \skt{व्योषक्षारसमन्वितम् ॥८९॥}\ A1938}}\skt{ ॥}\pend

\pstart\versepar \ledsidenote{\passref{1.45.90}}\skt{ग्राहिणी वातला रूक्षा }\edtext{\skt{विज्ञेया}}{\Afootnote{N H; \skt{दुर्जरा}\ A1938}}\skt{ }\edtext{\skt{तक्रकूचिका}}{\Afootnote{H; \skt{तक्रकूर्चिका}\ A1938; \skt{तु क्रकूचिका ।}\ N}}\skt{ । }\pend
\pstart\versepar \skt{}\edtext{\edtext{\skt{तद्वत्किरातम्मथितं}}{\lemma{\skt{तद्वत्किरातम्}}\Afootnote{H; \skt{त}\ A1938; \skt{तद्वत्किराटम्}\ N}}\skt{ }\edtext{\skt{बृंहणं}}{\Afootnote{H; \skt{बृंहनं}\ N}}\skt{ }\edtext{\skt{क्षीरमारटं}}{\lemma{\skt{आरटं}}\Afootnote{H; \skt{ओरटं ॥}\ N}}\skt{ ॥}}{\lemma{\skt{द्वत्… आरटं}}\Afootnote{\emph{om.}\ A1938}}\pend

\pstart \ledsidenote{\passref{1.45.92}}\edtext{\skt{नवनीतन्तु}}{\lemma{\skt{तु}}\Afootnote{N H; \skt{पुनः सद्यस्कं लघु}\ A1938}}\skt{ }\edtext{\edtext{\skt{सुकुमारमधुरमम्लानुरसंग्राहि}}{\lemma{\skt{सुकुमारमधुरम्}}\Afootnote{N H; \skt{सुकुमारं मधुरं}\ A1938}}\skt{ व्रणशोफार्द्दितापहं}}{\lemma{\skt{अम्लानुरसंग्राहि व्रणशोफार्द्दितापहं}}\Afootnote{H; \skt{कषायमीषदम्लं शीतलं मेध्यं दीपनं हृद्यं सङ्ग्राहि पित्तानिलहरं वृष्यमविदाहि क्षयकासव्रणशोषार्शोऽर्दितापहं, चिरोत्थितं गुरु कफमेदोविवर्धनं बलकरं बृंहणं शोषघ्नं विशेषेण बालानां प्रशस्यते ॥९२॥}\ A1938; \skt{अम्ब्लानुरसंग्राहि व्रणसोफार्दितापहं ॥}\ N}}\skt{ ।}\pend

\pstart \ledsidenote{\passref{1.45.96}}\edtext{\skt{घृतन्तु}}{\Afootnote{N H; \skt{अथ घृतम् । घृतं तु}\ A1938}}\skt{ }\edtext{\edtext{\edtext{\skt{शीतवीर्यम्मधुरमभिष्यन्दि}}{\lemma{\skt{शीतवीर्यम्}}\Afootnote{N H; \skt{मधुरं सौम्यं मृदु शीतवीर्य}\ A1938}}}{\lemma{\skt{मधुरम्}}\Afootnote{N H; \skt{म}\ A1938}}}{\lemma{\skt{अभिष्यन्दि}}\Afootnote{N H; \skt{नभिष्यन्दि}\ A1938}}\skt{ }\edtext{\skt{त्रिदोषापकर्षणमग्निदीपनं}}{\lemma{\skt{त्रिदोषापकर्षणम्}}\Afootnote{N H; \skt{स्नेहनमुदावर्तोन्मादापस्मारशूलज्वरानाहवातपित्तप्रशमनम्}\ A1938}}\skt{ }\edtext{\edtext{\skt{सौकुमार्योजस्तेजोबलकरमायुष्यं}}{\lemma{\skt{सौकुमार्योजस्}}\Afootnote{H; \skt{स्मृतिमतिमेधाकान्तिस्वरलावण्यसौकुमार्यौज}\ A1938; \skt{सौकुमार्योज}\ N}}}{\lemma{\skt{तेजोबलकरम्}}\Afootnote{H; \skt{स्तेजोबलकरम्}\ A1938 N}}\skt{ }\edtext{\edtext{\skt{वृष्यम्मेध्यम्वयस्थापनञ्च}}{\lemma{\skt{वयस्थापनञ्}}\Afootnote{N H; \skt{वयःस्थापनं}\ A1938}}}{\lemma{\skt{च}}\Afootnote{H; \skt{गुरु च}\ A1938; \emph{om.}\ N}}\skt{ }\edtext{\edtext{\skt{क्षुष्यम्पापोपशमनं}}{\lemma{\skt{क्षुष्यम्}}\Afootnote{A1938 H; \emph{om.}\ N}}}{\lemma{\skt{पापोपशमनं}}\Afootnote{H; \skt{श्लेष्माभिवर्धनं पाप्मालक्ष्मीप्रशमनं विषहरं}\ A1938; \skt{चक्षुष्यम्पापोपसमनं}\ N}}\skt{ }\edtext{\skt{रक्षोघ्नञ्चेति}}{\lemma{\skt{चेति}}\Afootnote{N H; \skt{च ॥९६॥}\ A1938}}\skt{ ।}\pend

\pstart\versepar \ledsidenote{\passref{1.45.108}}\edtext{\skt{सर्प्पिः}}{\Afootnote{N H; \emph{om.}\ A1938}}\skt{ }\edtext{\skt{पुराणन्तिमिरप्रतिश्याश्वासकासनुत्}}{\lemma{\skt{पुराणन्}}\Afootnote{N H; \emph{om.}\ A1938}}\skt{ । }\pend
\pstart\versepar \skt{मूर्च्छाकुष्ठविषोन्मादग्रहापस्मारनाशनं ॥}\edtext{}{\lemma{\skt{तिमिरप्रतिश्याश्वासकासनुत् । मूर्च्छाकुष्ठविषोन्मादग्रहापस्मारनाशनं}}\Afootnote{H\textsuperscript{ac}; \skt{भवति चात्र पुराणं तिमिरश्वासपीनसज्वरकासनुत् । मूर्च्छाकुष्ठविषोन्मादग्रहापस्मारनाशनम् ॥१०८॥}\ A1938; \skt{तिमिरप्रतिश्यायश्वासकासनुत् । मूर्च्छाकुष्ठविषोन्मादग्रहापस्मारनाशनं ।}\ N; \skt{तिमिरप्रप्रतिश्यायस्यापि प्रतिश्यावास्तितिश्याश्वासकासनुत् । मूर्च्छाकुष्ठविषोन्मादग्रहापस्मारनाशनं ॥}\ H\textsuperscript{pc}}}\pend

\pstart\versepar \ledsidenote{\passref{1.45.95}}\skt{विकल्प एष }\edtext{\skt{दध्यादि}}{\Afootnote{N H; \skt{दध्यादिः}\ A1938}}\skt{ श्रेष्ठो }\edtext{\skt{गव्योऽनुवर्ण्णितः}}{\lemma{\skt{ऽनुवर्ण्णितः}}\Afootnote{N H; \skt{ऽभिवर्णितः ।}\ A1938}}\skt{ । }\pend
\pstart\versepar \skt{}\edtext{\skt{विकल्पानवशिष्टानां}}{\Afootnote{H; \skt{विकल्पानवशिष्टांस्तु}\ A1938; \skt{विकल्पानवशिष्टाणां}\ N}}\skt{ }\edtext{\skt{क्षीरवीर्यात्समादिशेत्}}{\Afootnote{H; \skt{क्षीरवीर्यात्समादिशेत् ॥९५॥ इति तक्रवर्गः ।}\ A1938; \skt{क्षीरवीर्या समादिशेत् ॥}\ N}}\skt{ ॥}\pend

\pstart \ledsidenote{\passref{1.45.112}}\edtext{\skt{तिलतैलं}}{\Afootnote{H; \skt{अथ तैलानि । तैलं त्वाग्नेयमुष्णं तीक्ष्णं मधुरं मधुरविपाकं बृंहणं प्रीणनं व्यवायि सूक्ष्मं विशदं गुरु सरं विकासि वृष्यं त्वक्प्रसादनं मेधामार्दवमांसस्थैर्यवर्णबलकरं}\ A1938; \skt{तिलतैलम्मधुर}\ N}}\skt{ }\edtext{\skt{मधुरं}}{\Afootnote{H; \skt{चक्षुष्यं बद्धमूत्रं लेखनं}\ A1938; \emph{om.}\ N}}\skt{ }\edtext{\edtext{\edtext{\skt{तिक्तानुरसन्तीक्ष्णमनिल}}{\lemma{\skt{तिक्तानुरसन्}}\Afootnote{H; \skt{तिक्तकषायानुरसं}\ A1938; \skt{तिक्तानुसमन्}\ N}}}{\lemma{\skt{तीक्ष्णम्}}\Afootnote{N H; \skt{पाच}\ A1938}}}{\lemma{\skt{अनिल}}\Afootnote{N H; \skt{नमनिल}\ A1938}}\skt{ }\edtext{\edtext{\skt{वलासक्षौअकरमशीतम्पित्तजननं}}{\lemma{\skt{वलासक्षौअकरम्}}\Afootnote{H; \skt{बलासक्षयकरं}\ A1938; \skt{वलासक्षयकरम्}\ N}}}{\lemma{\skt{अशीतम्}}\Afootnote{N H; \skt{क्रिमिघ्नमशित}\ A1938}}\skt{ । योनिशिरः }\edtext{\skt{शूलप्रशमनन्तथा}}{\lemma{\skt{शूलप्रशमनन्}}\Afootnote{N H; \skt{कर्णशूलप्रशमनं गर्भाशयशोधनं च,}\ A1938}}\skt{ }\edtext{\skt{छिन्नभिन्नक्षताग्निदग्धपिच्चितभग्नच्यूतावमथितहतव्याडविदष्टपतनप्रभृतिषु}}{\Afootnote{H; \skt{छिन्नभिन्नविद्धोत्पिष्टच्युतमथितक्षतपिच्चितभग्नस्फुटितक्षाराग्निदग्धविश्लिष्टदारिता भिहतदुर्भग्नमृगव्यालविदष्टप्रभृतिषु च}\ A1938; \skt{छिन्नभिक्षताग्निदग्धपिच्चितरुग्नयूतावमथितहतव्याडविदष्टपतनप्रभृतिषु}\ N}}\skt{ }\edtext{\skt{परिषेकाभ्यङ्गनयोः}}{\Afootnote{H; \skt{परिषेकाभ्यङ्गावगाहादिषु तिलतैलं}\ A1938; \skt{परिसेकाभ्यञ्जनयोः}\ N}}\skt{ }\edtext{\edtext{\skt{प्रशस्तमिति}}{\lemma{\skt{प्रशस्तम्}}\Afootnote{N H; \skt{प्रशस्य}\ A1938}}}{\lemma{\skt{इति}}\Afootnote{N H; \skt{ते ॥११२॥}\ A1938}}\skt{ ॥}\pend

\pstart\versepar \ledsidenote{\passref{1.45.113}}\edtext{\skt{तद्बस्तिषु च}}{\Afootnote{N; \skt{तद्बस्तिषु च}\ A1938; \skt{तद्वस्तिषु च}\ H}}\skt{ }\edtext{\skt{पाने च}}{\Afootnote{H; \skt{पानेषु}\ A1938; \skt{पवने च}\ N}}\skt{ नस्ये }\edtext{\skt{कर्ण्णादिपूरणे}}{\Afootnote{N H; \skt{कर्णाक्षिपूरणे ।}\ A1938}}\skt{ । }\pend
\pstart\versepar \skt{अन्नपानविधौ चापि प्रयोज्यम्वातशान्तये ॥}\pend

\pstart \ledsidenote{\passref{1.45.115}}\edtext{\skt{निम्वातसीकुसुम्भसर्षपपीलुकरङ्गुदीशिग्रुसुवर्ण्णलाफलतैलानि}}{\Afootnote{H; \skt{निम्बातसीकुसुम्भमूलकजीमूतकवृक्षककृतवेधनार्ककम्पिल्लकहस्तिकर्णपृथ्वीकापीलु- करञ्जेङ्गुदीशिग्रुसर्षपसुवर्चलाविडङ्गज्योतिष्मतीफलतैलानि}\ A1938; \skt{निम्बातसीकुसुम्भषर्सपपीलुकरञ्जेङ्गुदीसिग्रुसुवर्च्चलाफलतैलानि}\ N}}\skt{ }\edtext{\skt{तीक्ष्ण}}{\Afootnote{N H; \skt{तीक्ष्णा}\ A1938}}\skt{ }\edtext{\edtext{\skt{कटुकान्युष्णवीन्याणि}}{\lemma{\skt{कटुकान्य्}}\Afootnote{N; \skt{नि}\ A1938; \skt{कटुका}\ H}}}{\lemma{\skt{उष्णवीन्याणि}}\Afootnote{N; \skt{लघून्युष्णवीर्याणि}\ A1938; \skt{न्युष्णवीन्याणि}\ H}}\skt{ }\edtext{\skt{क्रिमिकफमेहोपहरणानि}}{\Afootnote{H; \skt{कटूनि कटुविपाकानि सराण्यनिलकफकृमिकुष्ठप्रमेहशिरोरोगापहराणि चेति ॥११५॥}\ A1938; \skt{क्रिमिकफमेहोपहराणानि ॥}\ N}}\skt{ ॥}\pend

\pstart \ledsidenote{\passref{1.45.120}}\edtext{\skt{अतिमुक्तकपिलाक्षत्रपुसेर्व्वारुकुष्माण्डकतैलानि}}{\Afootnote{H; \skt{किराततिक्तकातिमुक्तकबिभीतकनालिकेरकोलाक्षोडजीवन्तीप्रियालकर्बुदार- सूर्यवल्लीत्रपुसैर्वारुककर्कारुकूष्माण्डप्रभृतीनां तैलानि मधुराणि मधुरविपाकानि वातपित्तप्रशमनानि शीतवीर्याण्यभिष्यन्दीनि सृष्टमूत्राण्यग्निसादनानि चेति ॥१२०॥}\ A1938; \skt{अतिमुक्तकपियालाक्षत्रपुसेर्वारुकर्कारुकुष्माण्डकतैलानि ।}\ N}}\skt{ ।}\pend

\pstart \ledsidenote{\passref{1.45.121}}\edtext{\skt{मधुरकषायाणि}}{\Afootnote{N H; \skt{मधूककाश्मर्यपलाशतैलानि मधुरकषायाणि}\ A1938}}\skt{ कफपित्तप्रशमनानि ॥}\pend

\pstart \ledsidenote{\passref{1.45.122}}\edtext{\skt{तुरुवककर्क्कोटकतैलमधुरकषाये}}{\Afootnote{H; \skt{तुवरकभल्लातकतैले उष्णे मधुरकषाये}\ A1938; \skt{तुरुवककर्कोटकतैलमधुरकषाय}\ N}}\skt{ तिक्तानुरसे }\edtext{\skt{क्रिमिकफकुष्ठमेदोदरहरे}}{\Afootnote{\skt{वातकफकुष्ठमेदोमेहकृमिप्रशमने उभयतोभागदोषहरे}\ A1938; \skt{क्रिमिकफकुष्ठमेदोहरे च ॥ एरण्डतैलम्मधुरमुष्णं वातकफापहं मेदोहरञ्}\ N; \skt{क्रिमिकफकुष्ठमेदोदरहरे च ॥}\ H}}\skt{ }\edtext{\skt{च}}{\Afootnote{A1938 N; \emph{om.}\ H}}\skt{ ॥}\pend

\pstart \ledsidenote{\passref{1.45.123}}\edtext{\skt{सरलदेवदारुशिंशपागुरुसारस्नेहाः}}{\Afootnote{H; \skt{सरलदेवदारुशिंशपागुरुगण्डीरसारस्नेहा}\ A1938; \skt{सरलदेवदारुसिंसपागुरुसारस्नेहाः ॥}\ N}}\skt{ । }\edtext{\skt{तिक्तककटुकषायादुष्टव्रणशोधनाः}}{\Afootnote{H; \skt{स्तिक्तकटुकषाया दुष्टव्रणशोधनाः}\ A1938; \skt{तिक्तककटुकशायाः । दुष्टव्रणशोधनाः}\ N}}\skt{ }\edtext{\skt{क्रिमिकफकुष्ठहराश्च}}{\lemma{\skt{क्रिमिकफकुष्ठहराश्}}\Afootnote{N\textsuperscript{pc} H; \skt{कृमिकफकुष्ठानिलहराश्}\ A1938; \skt{क्रिमिकिअफकुष्ठहराश्}\ N\textsuperscript{ac}}}\skt{ ॥}\pend

\pstart \ledsidenote{\passref{1.45.131}}\edtext{\skt{ग्राम्यानूपौदकानां}}{\Afootnote{H; \skt{ग्राम्यानूपौदकानां च वसामेदोमज्जानो गुरूष्णमधुरा वातघ्नाः, जाङ्गलैकशफक्रव्यादादीनां}\ A1938; \skt{ग्राम्यान्यानूपौदकानाम्वसामज्जानो गुरूष्णमधुरा वातघ्नाः । जाङ्गलैकशफक्रव्यादानां}\ N\textsuperscript{ac}; \skt{ग्राम्यानूपौदकानाम्वसामज्जानो गुरूष्णमधुरा वातघ्नाः । जाङ्गलैकशफक्रव्यादानां}\ N\textsuperscript{pc}}}\skt{ }\edtext{\skt{सर्व्वे }\edtext{\skt{दोषघ्नाश्चाग्निदीपनाः}}{\lemma{\skt{चाग्निदीपनाः}}\Afootnote{N H; \emph{om.}\ A1938}}}{\lemma{\skt{सर्व्वे दोषघ्नाश्}}\Afootnote{H; \skt{लघुशीतकषाया रक्तपित्तघ्नाः, प्रतुदविष्किराणां श्लेष्मघ्नाः ।}\ A1938; \skt{लघुशीतकषाया रपित्तघ्नाः । विष्किराः प्रमेहानां श्लेष्मघ्नाः । प्रतुदानां सर्वदोषघ्नाश्}\ N\textsuperscript{ac}; \skt{लघुशीतकषाया रक्तपित्तघ्नाः । विष्किराः प्रमेहानां श्लेष्मघ्नाः । प्रतुदानां सर्वदोषघ्नाश्}\ N\textsuperscript{pc}}}\skt{ । तत्र }\edtext{\skt{तैलघृतवसामज्जानो}}{\Afootnote{N H; \skt{घृततैलवसामेदोमज्जानो यथोत्तरं}\ A1938}}\skt{ }\edtext{\skt{गुरुपाका}}{\Afootnote{N H; \skt{गुरुविपाका}\ A1938}}\skt{ }\edtext{\skt{वातहराश्चेति}}{\lemma{\skt{चेति}}\Afootnote{N H; \skt{च ॥१३१॥ इति तैलवर्गः ।}\ A1938}}\skt{ ।}\pend

\pstart\versepar \ledsidenote{\passref{1.45.129}}\skt{यावन्तः }\edtext{\skt{स्थावराः}}{\Afootnote{\skt{स्थावराः स्ने}\ A1938; \skt{स्थावारा}\ N; \skt{स्थावरस्था वा}\ H}}\skt{ }\edtext{\skt{देहाः}}{\Afootnote{N; \skt{हाः}\ A1938; \skt{हाराः}\ H}}\skt{ }\edtext{\skt{समासेन}}{\Afootnote{N H; \skt{समासात्प}\ A1938}}\skt{ }\edtext{\skt{प्रकीर्त्तिताः}}{\Afootnote{N H; \skt{रिकीर्तिताः ।}\ A1938}}\skt{ । }\pend
\pstart\versepar \skt{सर्व्वे तैलगुणा ज्ञेयाः सर्व्वे }\edtext{\skt{चानिलनाशनाः}}{\Afootnote{A1938 H; \skt{चानिलनासनाः ॥}\ N}}\skt{ ॥}\pend

\pstart\versepar \ledsidenote{\passref{1.45.132}}\edtext{\skt{क्षौद्रन्तु}}{\Afootnote{N H; \skt{अथ मधुवर्गः । मधु तु}\ A1938}}\skt{ मधुरङ्कषायानुरसं रूक्षं }\edtext{\skt{लघुसुकुमारं}}{\Afootnote{N H; \skt{शीतमग्निदीपनं वर्ण्यं स्वर्यं लघु सुकुमारं}\ A1938}}\skt{ । }\edtext{\skt{सन्धानीयं}}{\Afootnote{N H; \skt{लेखनं हृद्यं वाजीकरणं सन्धानं}\ A1938}}\skt{ }\edtext{\skt{रोपणं}}{\Afootnote{H; \skt{शोधनं रोपणं(सङ्}\ A1938; \skt{सोधनं रोपनणं}\ N\textsuperscript{ac}; \skt{सोधनं रोपणं}\ N\textsuperscript{pc}}}\skt{ }\edtext{\edtext{\edtext{\edtext{\skt{लेखनञ्चक्षुष्यम्वर्ण्यैश्वर्यम्विषघ्नङ्कृमिच्छर्द्यतीसारमेहपित्तकफप्रशमनं}}{\lemma{\skt{लेखनञ्}}\Afootnote{N H; \skt{ग्राहि )}\ A1938}}}{\lemma{\skt{वर्ण्यैश्वर्यम्}}\Afootnote{H; \skt{प्र}\ A1938; \skt{वर्ण्यं श्वर्यम्}\ N}}}{\lemma{\skt{विषघ्नङ्}}\Afootnote{H; \skt{सादनं}\ A1938; \skt{विषघ्न}\ N}}}{\lemma{\skt{कृमिच्छर्द्यतीसारमेहपित्तकफप्रशमनं}}\Afootnote{H; \skt{सूक्ष्ममार्गानुसारि पित्तश्लेष्ममेदोमेहहिक्काश्वासकासातिसारच्छर्दितृष्णाकृमिविषप्रशमनं}\ A1938; \skt{क्रिमिच्छर्द्यतीसारमेहपित्तकफरक्तप्रशमनं}\ N}}\skt{ }\edtext{\skt{साग्राहिकं}}{\Afootnote{H; \skt{ह्लादि त्रिदोष}\ A1938; \skt{सङ्ग्राहिकम्}\ N}}\skt{ }\edtext{\skt{प्रह्लादनं}}{\Afootnote{H; \skt{प्रशमनं}\ A1938; \skt{प्रह्लादन}\ N}}\skt{ }\pend
\pstart\versepar \skt{}\edtext{\skt{तत्तु}}{\Afootnote{N H; \skt{च; त}\ A1938}}\skt{ }\edtext{\skt{त्रिविधं}}{\Afootnote{N H; \skt{त्तु}\ A1938}}\skt{ माक्षिकं }\edtext{\edtext{\skt{पौत्तिकं}}{\Afootnote{H; \skt{पोतिकं}\ N}}\skt{ भ्रामरं }\edtext{\skt{पूर्व्व}}{\Afootnote{H; \skt{पूर्वं}\ N}}\skt{ }\edtext{\skt{पूर्व्व}}{\Afootnote{H; \skt{पूर्वं}\ N}}\skt{ }\edtext{\skt{लघुतरं}}{\Afootnote{N H; \skt{लघुत्वात्कफघ्नं, पैच्छिल्यान्माधुर्यात्कषायभावाच्च वातपित्तघ्नम् ॥१३२॥}\ A1938}}\skt{ । }\edtext{\edtext{\edtext{\edtext{\skt{तत्पुराणम्प्रधानमनम्लञ्च}}{\lemma{\skt{तत्पुराणम्}}\Afootnote{H; \emph{om.}\ A1938; \skt{तत्पुराणां}\ N}}}{\lemma{\skt{प्रधानम्}}\Afootnote{N H; \emph{om.}\ A1938}}}{\lemma{\skt{अनम्लञ्}}\Afootnote{H; \emph{om.}\ A1938; \skt{अनम्ब्लञ्}\ N}}}{\lemma{\skt{च}}\Afootnote{N H; \emph{om.}\ A1938}}}{\lemma{\skt{विधं… ल}}\Afootnote{\emph{om.}\ A1938}}\skt{ ।}\pend

\pstart \ledsidenote{\passref{1.45.143}}\edtext{\skt{नानाद्रव्येभ्यो}}{\Afootnote{N H; \skt{तत्तु नानाद्रव्य}\ A1938}}\skt{ }\edtext{\skt{विरुद्धरसवीर्यविषपुष्परसपानात्मक्षिकासम्भृतत्वाच्च}}{\Afootnote{H; \skt{रसगुणवीर्यविपाकविरुद्धानां पुष्परसानां मक्षिकासम्भवत्वाच्चा}\ A1938; \skt{विरुद्धरसवीर्यविषपुष्परसपानान्मक्षिकासम्भृतत्वाच्च ।}\ N}}\skt{ }\edtext{\skt{उष्णोपचारयोगवाहि}}{\Afootnote{N H; \skt{नुष्णोपचारम् ॥१४३॥}\ A1938}}\skt{ }\edtext{\skt{च}}{\Afootnote{N H; \emph{om.}\ A1938}}\skt{ ॥}\pend

\pstart\versepar \ledsidenote{\passref{1.45.144}}\edtext{\skt{भवति}}{\Afootnote{N H; \emph{om.}\ A1938}}\skt{ }\edtext{\skt{चात्र}}{\Afootnote{H; \emph{om.}\ A1938 N}}\skt{ ॥ }\pend
\pstart\versepar \skt{उष्णैर्व्विरुध्यते सर्वम्विषान्वयतया }\edtext{\skt{मधुः}}{\Afootnote{N H; \skt{मधु ।}\ A1938}}\skt{ । }\pend
\pstart\versepar \skt{}\edtext{\edtext{\edtext{\skt{उष्णार्त्तमुष्णमुष्णे}}{\lemma{\skt{उष्णार्त्तम्}}\Afootnote{N H; \skt{उष्णार्त}\ A1938}}}{\lemma{\skt{उष्णम्}}\Afootnote{N H; \skt{मुष्णैर्}\ A1938}}\skt{ वा }\edtext{\skt{तन्निहन्ति}}{\lemma{\skt{निहन्ति}}\Afootnote{A1938 H; \skt{ल्लिहन्ति}\ N}}}{\lemma{\skt{उष्णे वा तन्}}\Afootnote{A1938 H; \skt{उष्णे वा त}\ N}}\skt{ }\edtext{\skt{विषं}}{\Afootnote{N H; \emph{om.}\ A1938}}\skt{ }\edtext{\skt{यथा}}{\Afootnote{N H; \skt{यथा विषम् ॥१४४॥}\ A1938}}\skt{ ॥}\pend

\pstart \ledsidenote{\passref{1.45.148}}\edtext{\skt{इक्षवो}}{\Afootnote{N H; \skt{अथेक्षुवर्गः । इक्षवो}\ A1938}}\skt{ मधुरा मधुरविपाका }\edtext{\skt{शीता}}{\Afootnote{N H; \skt{गुरवः शीताः}\ A1938}}\skt{ स्निग्धा }\edtext{\skt{वृष्या}}{\Afootnote{N H; \skt{बल्या वृष्या}\ A1938}}\skt{ मूत्रलाः । }\edtext{\edtext{\edtext{\skt{रक्तपित्तप्रशमनाङ्कफकराश्चेति}}{\lemma{\skt{रक्तपित्तप्रशमनाङ्}}\Afootnote{H; \skt{रक्तपित्तप्रशमनाः}\ A1938 N}}}{\lemma{\skt{कफकराश्}}\Afootnote{N H; \skt{कृमिकफकरा}\ A1938}}}{\lemma{\skt{चेति}}\Afootnote{N H; \skt{श्चेति । ते चानेकविधाः । तद्यथा- ॥१४८॥}\ A1938}}\skt{ ॥}\pend

\pstart\versepar \ledsidenote{\passref{1.45.155}}\edtext{\skt{स्नेहप्रसादमाधुर्यगुणोत्कर्षप्रकारकः}}{\Afootnote{H; \skt{स्नेहप्रसादमाधुर्यगुणोत्कर्षप्रकारतः ।}\ N}}\skt{ । }\pend
\pstart\versepar \skt{कान्तारकाद्वरः }\edtext{\skt{पौण्ड्रः}}{\Afootnote{H; \skt{पौण्ठः}\ N}}\skt{ }\edtext{\skt{पौण्ड्रकाद्वंसकोवरः}}{\Afootnote{H; \skt{पौण्ड्रकाद्वशकोवरं ॥}\ N}}\skt{ ॥}\edtext{}{\lemma{\skt{… पौण्ड्रकाद्वंसको}}\Afootnote{\emph{om.}\ A1938}}\pend

\pstart \ledsidenote{\passref{1.45.157}}\edtext{\skt{शर्क्करासमवीर्यस्तु}}{\Afootnote{N H; \skt{अविदाही कफकरो वातपित्तनिवारणः । वक्त्रप्रह्लादनो वृष्यो}\ A1938}}\skt{ दन्तनिष्पीडितो रसः ।}\pend

\pstart \ledsidenote{\passref{1.45.158}}\skt{गुरुर्विदाही विष्टम्भी }\edtext{\skt{यान्त्रिकस्तु}}{\Afootnote{A1938 H; \skt{यान्तिकस्तु}\ N}}\skt{ }\edtext{\skt{प्रकीर्त्तितः}}{\Afootnote{N H; \skt{प्रकीर्तितः पक्वो गुरुः सरः स्निग्धः सतीक्ष्णः कफवातनुत् ॥१५८॥}\ A1938}}\skt{ ॥}\pend

\pstart \ledsidenote{\passref{1.45.159}}\edtext{\skt{फाणितम्मधुरमभिष्यन्दि}}{\lemma{\skt{मधुरम्}}\Afootnote{N H; \skt{गुरु मधुरम्}\ A1938}}\skt{ }\edtext{\skt{बृंहणं}}{\Afootnote{A1938 H; \skt{बृंह}\ N}}\skt{ । }\edtext{\skt{शुक्रकरं}}{\Afootnote{H; \skt{अवृष्यं त्रिदोषकृच्च ॥१५९॥}\ A1938; \skt{मणं शुक्रकफकरं}\ N\textsuperscript{ac}; \skt{णं शुक्रकफकरं}\ N\textsuperscript{pc}}}\skt{ }\edtext{\edtext{\skt{पित्तघ्नञ्च}}{\lemma{\skt{पित्तघ्नञ्}}\Afootnote{N H; \emph{om.}\ A1938}}}{\lemma{\skt{च}}\Afootnote{N H; \emph{om.}\ A1938}}\skt{ ।}\pend

\pstart \ledsidenote{\passref{1.45.160}}\skt{गुडः सक्षारमधुरो नातिशीतः स्निग्धो }\edtext{\skt{मूत्ररक्तविशोधनः}}{\Afootnote{N H; \skt{मूत्ररक्तशोधनो}\ A1938}}\skt{ । }\edtext{\skt{पित्तघ्नङ्कफकरो}}{\lemma{\skt{पित्तघ्नङ्}}\Afootnote{H; \skt{नातिपित्तजिद्वातघ्नो मेदःकृमि}\ A1938; \skt{पित्तघ्नः ।}\ N}}\skt{ }\edtext{\skt{वृष्यश्च}}{\lemma{\skt{वृष्यश्}}\Afootnote{N H; \skt{बल्यो वृष्यश्}\ A1938}}\skt{ ॥}\pend

\pstart \ledsidenote{\passref{1.45.162}}\skt{मत्स्यण्डिकाखण्डशर्क्करा }\edtext{\skt{विमला}}{\Afootnote{N H; \skt{विमलजाता}\ A1938}}\skt{ }\edtext{\skt{उत्तरोत्तरे च}}{\Afootnote{H; \skt{उत्तरोत्तरं}\ A1938; \skt{उत्तरोत्तरश्च}\ N}}\skt{ }\edtext{\edtext{\skt{शीतस्निग्धगुरुसरमधुरारक्तपित्तप्रशमनाश्चेति}}{\lemma{\skt{शीतस्निग्धगुरुसरमधुरारक्तपित्तप्रशमनाश्}}\Afootnote{N H\textsuperscript{pc}; \skt{शीताः स्निग्धाः गुरुतरा मधुरतरा वृष्या रक्तपित्तप्रशमनास्तृष्णाप्रशमनाश्}\ A1938; \skt{शीतस्निग्धगुरुसमधुरारक्तपित्तप्रशमनाश्}\ H\textsuperscript{ac}}}}{\lemma{\skt{चेति}}\Afootnote{N H; \skt{च ॥१६२॥}\ A1938}}\skt{ ॥}\pend

\pstart\versepar \ledsidenote{\passref{1.45.163}}\skt{यथा यथैषाम्वैमल्यम्मधुरत्वन्तथा तथा । }\pend
\pstart\versepar \skt{}\edtext{\skt{स्नेहलाघवशैत्यानि}}{\Afootnote{N H; \skt{स्नेहगौरवशैत्यानि}\ A1938}}\skt{ सरत्वञ्च तथा तथा ॥}\pend

\pstart \ledsidenote{\passref{1.45.166}}\skt{मधुशर्क्करा }\edtext{\skt{छर्द्द्यतीसारहरा}}{\Afootnote{N H; \skt{पुनश्छर्द्यतीसारहरी}\ A1938}}\skt{ रूक्षा }\edtext{\skt{च्छेदनी}}{\Afootnote{N H; \skt{छेदनी प्रसादनी}\ A1938}}\skt{ }\edtext{\skt{च}}{\Afootnote{N H; \skt{कषायमधुरा मधुरविपाका च ॥१६६॥}\ A1938}}\skt{ । }\edtext{\skt{वृष्या}}{\Afootnote{N H; \emph{om.}\ A1938}}\skt{ }\edtext{\skt{क्षीणक्षतसन्धानकृत्सस्नेहा}}{\Afootnote{H; \emph{om.}\ A1938; \skt{क्षीणक्षतसन्धानकृत् । सस्नेहेहा}\ N}}\skt{ }\edtext{\skt{गुडशर्क्करा}}{\Afootnote{N H; \emph{om.}\ A1938}}\skt{ ॥}\pend

\pstart \ledsidenote{\passref{1.45.167}}\edtext{\edtext{\skt{कषायशीतमधुरः}}{\Afootnote{H; \skt{यवा}\ A1938; \skt{कषायआशीतमधुर्ः}\ N\textsuperscript{ac}; \skt{कषायशीतमधुराः}\ N\textsuperscript{pc}}}\skt{ सतिक्तायावशर्क्कराः}}{\lemma{\skt{शीतमधुरः सतिक्तायाव}}\Afootnote{\emph{om.}\ A1938}}\skt{ ।}\pend

\pstart \ledsidenote{\passref{1.45.168}}\edtext{\skt{तृष्णा}}{\Afootnote{N H; \emph{om.}\ A1938}}\skt{ }\edtext{\skt{शोणितपित्तदाहशमनीसामान्यतः}}{\Afootnote{N H; \skt{यावत्यः शर्कराः}\ A1938}}\skt{ }\edtext{\skt{सर्व्वेण}}{\Afootnote{N H; \skt{प्रोक्ताः सर्वा दाहप्रणा}\ A1938}}\skt{ }\edtext{\skt{सर्व्वम्पित्तहरं}}{\Afootnote{N H; \skt{शनाः । रक्तपित्तप्रशमनाश्छर्दिमूर्च्छातृषापहाः ॥१६८॥}\ A1938}}\skt{ ॥}\pend

\pstart \ledsidenote{\passref{1.45.170ab}}\edtext{\edtext{\edtext{\skt{मद्यमम्लन्दीपनं}}{\lemma{\skt{मद्यम्}}\Afootnote{N H; \skt{अथ मद्यवर्गः । सर्वं पित्तकरं मद्यम्}\ A1938}}}{\lemma{\skt{अम्लन्}}\Afootnote{A1938 H; \emph{om.}\ N}}}{\lemma{\skt{दीपनं}}\Afootnote{H; \skt{रोचनदीपनम् ।}\ A1938; \skt{अम्ब्लं दीपणं}\ N}}\skt{ }\edtext{\skt{रोचनम्}}{\Afootnote{N H; \emph{om.}\ A1938}}\skt{ ।}\pend

\pstart \ledsidenote{\passref{1.45.171cd}}\edtext{\skt{विकासि}}{\Afootnote{A1938 H; \skt{विका}\ N}}\skt{ }\edtext{\skt{सृष्टविट्मूत्रं}}{\Afootnote{N H; \skt{सृष्टविण्मूत्रं}\ A1938}}\skt{ }\edtext{\skt{श्र्णु}}{\Afootnote{N H; \skt{शृणु}\ A1938}}\skt{ तस्य विशेषणं ॥}\pend

\pstart\versepar \ledsidenote{\passref{1.45.172}}\edtext{\edtext{\skt{माद्वीकमविदाहित्वात्मधुरानुरसस्तथा}}{\lemma{\skt{माद्वीकम्}}\Afootnote{N H; \skt{मार्द्वीकम्}\ A1938}}}{\lemma{\skt{अविदाहित्वात्मधुरानुरसस्}}\Afootnote{H; \skt{अविदाहित्वान्मधुरान्वयतस्}\ A1938; \skt{अविदाहित्वान्मधुरानुरसन्}\ N}}\skt{ । }\pend
\pstart\versepar \skt{रक्तपित्ते }\edtext{\skt{तु}}{\Afootnote{N H; \skt{ऽपि}\ A1938}}\skt{ सततम्बुधैर्न्न प्रतिषिध्यते ॥}\pend

\pstart \ledsidenote{\passref{1.45.174}}\edtext{\edtext{\skt{तस्मादल्पान्तरङ्किंचित्खार्ज्जूरम्वातकोपनं}}{\lemma{\skt{तस्मादल्पान्तरङ्}}\Afootnote{N H; \skt{मार्द्वीकाल्पान्तरं}\ A1938}}}{\lemma{\skt{वातकोपनं}}\Afootnote{N H; \skt{वातकोपनम् । तदेव विशदं रुच्यं कफघ्नं कर्शनं लघु ॥१७४॥}\ A1938}}\skt{ ।}\pend

\pstart \ledsidenote{\passref{1.45.182}}\edtext{\skt{कषायमधुरः}}{\Afootnote{N H; \skt{दीपनः सृष्टविण्मूत्रो विशदोऽल्पमदो गुरुः कषायो मधुरः}\ A1938}}\skt{ }\edtext{\edtext{\skt{सीधुर्गौडः}}{\lemma{\skt{सीधुर्}}\Afootnote{N H; \skt{शीधु}\ A1938}}}{\lemma{\skt{गौडः}}\Afootnote{N H; \skt{र्गौडः}\ A1938}}\skt{ पाचनदीपनः ॥}\pend

\pstart \ledsidenote{\passref{1.45.184}}\edtext{\skt{तद्वत्पक्वरसं}}{\Afootnote{N H; \skt{तद्वत्पक्वरसः}\ A1938}}\skt{ }\edtext{\edtext{\skt{सीधुर्ब्बलवर्ण्णकरः}}{\lemma{\skt{सीधुर्}}\Afootnote{N H; \skt{शीधु}\ A1938}}}{\lemma{\skt{ब्बलवर्ण्णकरः}}\Afootnote{N H; \skt{र्बलवर्णकरः}\ A1938}}\skt{ }\edtext{\skt{परः}}{\Afootnote{N H; \skt{सरः । शोफघ्नो दीपनो हृद्यो रुच्यः श्लेष्मार्शसां हितः ॥१८४॥}\ A1938}}\skt{ ॥}\pend

\pstart \ledsidenote{\passref{1.45.186}}\edtext{\skt{माक्षिकः}}{\Afootnote{N H; \skt{आक्षिकः}\ A1938}}\skt{ पाण्डुरोगघ्नो }\edtext{\skt{वृष्यः}}{\Afootnote{N H; \skt{व्रण्यः}\ A1938}}\skt{ }\edtext{\skt{सङ्ग्राहिको}}{\Afootnote{N H; \skt{सङ्ग्राहको}\ A1938}}\skt{ }\edtext{\skt{लघुः}}{\Afootnote{N H; \skt{लघुः । कषायमधुरः सीधुः पित्तघ्नोऽसृक्प्रसादनः ॥१८६॥}\ A1938}}\skt{ ॥}\pend

\pstart\versepar \ledsidenote{\passref{1.45.187}}\edtext{\skt{जाम्बवो}}{\Afootnote{A1938 H; \skt{जम्ब्लोष्ठ}\ N}}\skt{ }\edtext{\edtext{\skt{बद्धनिष्यन्दस्तौरवो}}{\lemma{\skt{बद्धनिष्यन्दस्}}\Afootnote{N H; \skt{बद्धनिस्यन्द}\ A1938}}}{\lemma{\skt{तौरवो}}\Afootnote{N H; \skt{स्तुवरो}\ A1938}}\skt{ वातकोपनः । }\pend
\pstart\versepar \skt{तीक्ष्णः सुरासवो }\edtext{\skt{द्यो}}{\Afootnote{H; \skt{हृद्यो}\ A1938 N}}\skt{ मूत्रलः }\edtext{\skt{कफवातहा}}{\Afootnote{N H; \skt{कफवातनुत् ॥१८७॥}\ A1938}}\skt{ ॥}\pend

\pstart \ledsidenote{\passref{1.45.188}}\edtext{\skt{मधुरो}}{\Afootnote{N H; \skt{मुखप्रियः स्थिरमदो विज्ञेयोऽनिलनाश}\ A1938}}\skt{ }\edtext{\edtext{\skt{गुडमैरेयश्च्छैदी}}{\lemma{\skt{गुडमैरेयश्}}\Afootnote{N H; \skt{नः लघु}\ A1938}}}{\lemma{\skt{च्छैदी}}\Afootnote{H; \emph{om.}\ A1938; \skt{च्छेदी}\ N}}\skt{ }\edtext{\skt{मध्वासवो}}{\Afootnote{N H; \skt{र्मध्वासवश्}\ A1938}}\skt{ }\edtext{\skt{लघुः}}{\Afootnote{N H; \skt{छेदी मेहकुष्ठविषापहः ॥१८८॥}\ A1938}}\skt{ ।}\pend

\pstart \ledsidenote{\passref{1.45.190}}\edtext{\skt{वल्यः}}{\Afootnote{N H; \skt{कृमिमेदोनिलहरो मैरेयो मधुरो गुरुः बल्यः}\ A1938}}\skt{ }\edtext{\skt{पित्तसहो}}{\Afootnote{N H; \skt{पित्तहरो}\ A1938}}\skt{ }\edtext{\skt{वृष्यो}}{\Afootnote{N H; \skt{वर्ण्यो}\ A1938}}\skt{ }\edtext{\skt{द्यश्चेक्षुरसासवः}}{\Afootnote{H; \skt{हृद्यश्चेक्षुरसासवः ॥१९०॥}\ A1938; \skt{हृद्यश्चेक्षरसासवः ।}\ N}}\skt{ ॥}\pend

\pstart \ledsidenote{\passref{1.45.190.1}}\edtext{\edtext{\skt{प्रजरणोऽरिष्टरसः}}{\lemma{\skt{प्रजरणो}}\Afootnote{N H\textsuperscript{pc}; \emph{om.}\ H\textsuperscript{ac}}}}{\lemma{\skt{ऽरिष्टरसः}}\Afootnote{H\textsuperscript{pc}; \skt{रिष्टरसः}\ N; \emph{om.}\ H\textsuperscript{ac}}}\skt{ }\edtext{\skt{कफहा}}{\Afootnote{N H\textsuperscript{pc}; \emph{om.}\ H\textsuperscript{ac}}}\skt{ }\edtext{\skt{भुक्तपाचनः}}{\Afootnote{N H\textsuperscript{pc}; \emph{om.}\ H\textsuperscript{ac}}}\skt{ ।}\pend

\pstart \ledsidenote{\passref{1.45.197}}\edtext{\skt{अरिष्टासवसीधूनां}}{\Afootnote{H; \skt{चिकित्सितेषु वक्ष्यन्तेऽरिष्टा रोगहराः पृथकरिष्टासवशीधूनां}\ A1938; \skt{अरिष्टासमसीधूनां}\ N}}\skt{ }\edtext{\skt{गुणान्कर्म्म्णि}}{\Afootnote{H; \skt{गुणान्कर्माणि}\ A1938; \skt{गुणां मि}\ N}}\skt{ चादिशेत् ॥}\pend

\pstart \ledsidenote{\passref{1.45.198}}\edtext{\skt{यथास्वौषधसंस्कारमवेक्ष्य}}{\lemma{\skt{यथास्वौषधसंस्कारम्}}\Afootnote{H; \skt{बुद्ध्या यथास्वं संस्कारम्}\ A1938; \skt{यथास्वौषधसस्कारम्}\ N}}\skt{ कुशलो }\edtext{\skt{भिषक्}}{\Afootnote{N H; \skt{भिषक्सान्द्रं विदाहि दुर्गन्धं विरसं कृमिलं गुरु ॥१९८॥}\ A1938}}\skt{ ।}\pend

\pstart \ledsidenote{\passref{1.45.176cd}}\edtext{\skt{कासार्शो}}{\Afootnote{H; \skt{कासार्शो}\ A1938; \skt{काशार्षो}\ N}}\skt{ }\edtext{\skt{ग्रहणीदोषप्रतिश्यायविनाशनी}}{\Afootnote{N H; \skt{ग्रहणीश्वासप्रतिश्यायविनाशनी ॥१७६॥}\ A1938}}\skt{ ॥}\pend

\pstart\versepar \ledsidenote{\passref{1.45.177}}\edtext{\skt{स्वेदमूत्रकफास्तन्यरक्तमान्सकरी}}{\Afootnote{H; \skt{श्वेता मूत्रकफस्तन्यरक्तमांसकरी}\ A1938; \skt{स्वेतामूत्रकफस्तन्यरक्तमान्सकरी}\ N}}\skt{ सुरा । }\pend
\pstart\versepar \skt{}\edtext{\skt{वम्य}}{\Afootnote{N H; \skt{छर्द्य}\ A1938}}\skt{ रोचकहृत्कुक्षि तोदशूलप्रमर्द्दनी ॥}\pend

\pstart \ledsidenote{\passref{1.45.178ab}}\skt{प्रसन्ना }\edtext{\skt{वातगुल्माऽर्शो}}{\Afootnote{N H; \skt{कफवातार्शो}\ A1938}}\skt{ }\edtext{\skt{विवन्धानाहनाशनी}}{\Afootnote{H; \skt{विबन्धानाहनाशनी}\ A1938 N}}\skt{ ।}\pend

\pstart \ledsidenote{\passref{1.45.180cd}}\edtext{\skt{ग्राह्युष्णो}}{\Afootnote{A1938 N H\textsuperscript{pc}; \emph{om.}\ H\textsuperscript{ac}}}\skt{ }\edtext{\skt{जगलो}}{\Afootnote{A1938 N H\textsuperscript{pc}; \emph{om.}\ H\textsuperscript{ac}}}\skt{ }\edtext{\skt{रूक्षः}}{\Afootnote{N H\textsuperscript{pc}; \skt{पक्ता}\ A1938; \emph{om.}\ H\textsuperscript{ac}}}\skt{ }\edtext{\skt{शोफहा}}{\Afootnote{N H\textsuperscript{pc}; \skt{रूक्षस्तृ}\ A1938; \emph{om.}\ H\textsuperscript{ac}}}\skt{ }\edtext{\skt{भुक्तपाचनः}}{\Afootnote{N H\textsuperscript{pc}; \skt{ट्कफशोफकृत् ॥१८०॥}\ A1938; \emph{om.}\ H\textsuperscript{ac}}}\skt{ ॥}\pend

\pstart \ledsidenote{\passref{1.45.181cd}}\edtext{\skt{बहुसो}}{\Afootnote{\skt{बक्कसो}\ A1938; \skt{वक्षशो}\ N; \skt{वहुसो}\ H}}\skt{ }\edtext{\skt{हृतसारत्वाद्विष्टम्भीदोषकोपनः}}{\Afootnote{H; \skt{हृतसारत्वाद्विष्टम्भी वातकोपनः ॥१८१॥}\ A1938; \skt{हृतसारत्वाद्विष्टम्भीदोषपानः ।}\ N}}\skt{ ।}\pend

\pstart \ledsidenote{\passref{1.45.192cd}}\edtext{\edtext{\skt{नवम्मद्यमभिष्यन्दि}}{\lemma{\skt{मद्यम्}}\Afootnote{N H; \skt{मद्य}\ A1938}}}{\lemma{\skt{अभिष्यन्दि}}\Afootnote{N H; \skt{मभिष्यन्दि}\ A1938}}\skt{ गुरुवातादिकोपनं ॥}\pend

\pstart \ledsidenote{\passref{1.45.194ab}}\edtext{\skt{स्फुटस्रोतकरञ्जीर्ण्णं}}{\lemma{\skt{स्फुटस्रोतकरञ्}}\Afootnote{N H; \skt{स्फुटस्रोतस्करं}\ A1938}}\skt{ लघुवातकफापहं ।}\pend

\pstart \ledsidenote{\passref{1.45.210ab}}\edtext{\skt{रक्तपित्तहरं}}{\Afootnote{H; \skt{रक्तपित्तकरं}\ A1938 N}}\skt{ }\edtext{\skt{शुक्तं}}{\Afootnote{A1938 H; \skt{शुक्लं}\ N}}\skt{ }\edtext{\skt{च्छेदि}}{\Afootnote{N H; \skt{छेदनं}\ A1938}}\skt{ }\edtext{\skt{शुक्तविपाचनं}}{\Afootnote{H; \skt{भुक्तपाचनम् ।}\ A1938; \skt{शुक्तविपानं ॥}\ N}}\skt{ ॥}\pend

\pstart \ledsidenote{\passref{1.45.211cd}}\edtext{\skt{तद्वत्तदाश्रुतं}}{\Afootnote{N H; \skt{तद्वत्तदासुतं}\ A1938}}\skt{ सर्व्वं }\edtext{\skt{रोचनन्तु}}{\Afootnote{N H; \skt{रोचनं च}\ A1938}}\skt{ }\edtext{\skt{विशेषतः}}{\Afootnote{A1938 H; \skt{विसेषतः ॥}\ N}}\skt{ ।}\pend

\pstart \ledsidenote{\passref{1.45.213}}\edtext{\skt{तुषाम्वु}}{\Afootnote{H; \skt{तुषाम्बु}\ A1938 N}}\skt{ दीपनं }\edtext{\skt{हृद्यमुक्तं}}{\lemma{\skt{उक्तं}}\Afootnote{N H; \skt{हृत्पाण्डुकृमिरोगनुत् । ग्रहण्यर्शोविकारघ्नं भेदि}\ A1938}}\skt{ सौवीरकन्तथा ॥}\pend

\pstart \ledsidenote{\passref{1.45.214.1}}\edtext{\edtext{\skt{धान्याम्लन्धान्ययोनित्वात्प्राणधारणमम्लत्वाद्वातघ्नम्विदाहित्वात्पित्तकरं}}{\lemma{\skt{धान्याम्लन्}}\Afootnote{H; \skt{धान्याम्ब्लं}\ N}}}{\lemma{\skt{अम्लत्वाद्वातघ्नम्}}\Afootnote{H; \skt{अम्ब्लत्वाद्वातघ्नम्}\ N}}\skt{ }\edtext{\skt{कफघ्नम्वहिः}}{\lemma{\skt{वहिः}}\Afootnote{H; \skt{बहिः ।}\ N}}\skt{ शीतं गुरुविपाकं हृद्यत्वात्तृष्णापनयनं }\edtext{\skt{रुचिजननं}}{\Afootnote{H; \skt{रुचिजनंनं}\ N}}\skt{ समुद्रान्तसंश्रितानां च जनानाम्परं सात्म्यं ॥}\pend

\pstart \ledsidenote{\passref{1.45.204cd}}\skt{तस्यानेकप्रकारस्य मद्यस्य रसवीर्यतः ।}\pend

\pstart\versepar \ledsidenote{\passref{1.45.205}}\edtext{\skt{सौक्ष्म्यादौष्णाच्च}}{\Afootnote{H; \skt{सौक्ष्म्यादौष्ण्याच्च}\ A1938; \skt{सौक्ष्मादौष्णाच्च}\ N}}\skt{ }\edtext{\skt{तैक्ष्ण्याच्च}}{\Afootnote{A1938 H; \skt{तैक्ष्णाच्च}\ N}}\skt{ विकासित्वाच्च वह्निना । }\pend
\pstart\versepar \skt{समेत्य }\edtext{\skt{हृदये}}{\Afootnote{N H; \skt{हृदयं}\ A1938}}\skt{ प्राप्य }\edtext{\edtext{\skt{धमनीमूर्ध्वमागतं}}{\lemma{\skt{धमनीम्}}\Afootnote{N H; \skt{धमनीर्}\ A1938}}}{\lemma{\skt{आगतं}}\Afootnote{N H; \skt{आगतम् । विचाल्येन्द्रियचेतांसि वीर्यं मदयतेऽचिरात् ॥२०५॥}\ A1938}}\skt{ ।}\pend

\pstart\versepar \ledsidenote{\passref{1.45.206}}\skt{विचाल्येन्द्रियचेतांसि वीर्यम्मदयते चिरात् ॥ }\pend
\pstart\versepar \skt{चिरेण श्लैष्मिके पुंसि पानतो जायते मदः । }\pend
\pstart\versepar \skt{}\edtext{\skt{वातिके}}{\Afootnote{N H; \skt{अचिराद्वातिके}\ A1938}}\skt{ }\edtext{\skt{जायते}}{\Afootnote{N H; \emph{om.}\ A1938}}\skt{ }\edtext{\skt{तीक्ष्णः}}{\Afootnote{N H; \skt{दृष्टः,}\ A1938}}\skt{ पैत्तिके शीघ्रमेव }\edtext{\skt{तु}}{\Afootnote{N H; \skt{च ॥२०६॥}\ A1938}}\skt{ ॥}\edtext{}{\lemma{\skt{… चि}}\Afootnote{\emph{om.}\ A1938}}\edtext{}{\lemma{\skt{चिरात् ॥ चिरेण}}\Afootnote{N; \skt{चिरेण}\ H}}\pend

\pstart \ledsidenote{\passref{1.45.207ab}}\skt{सात्त्विके }\edtext{\skt{सौचदाक्षिण्यहर्षमण्डनवत्स्थितः}}{\Afootnote{H; \skt{शौचदाक्षिण्यहर्षमण्डनलालसः ।}\ A1938; \skt{सौचदाक्षिण्यहर्षमण्डलत स्थितः ।}\ N}}\skt{ ।}\pend

\pstart\versepar \ledsidenote{\passref{1.45.208}}\skt{राजसे }\edtext{\skt{दुःखशीलत्वमात्मत्यागं}}{\lemma{\skt{दुःखशीलत्वम्}}\Afootnote{A1938 N H\textsuperscript{pc}; \skt{दुःशीलत्वम्}\ H\textsuperscript{ac}}}\skt{ }\edtext{\skt{सुसाहसं}}{\Afootnote{N H; \skt{ससाहसम् ।}\ A1938}}\skt{ । }\pend
\pstart\versepar \skt{}\edtext{\skt{कलहं}}{\Afootnote{A1938 H; \skt{कलहंश्}\ N}}\skt{ }\edtext{\skt{सानवस्थानं}}{\Afootnote{H; \skt{सानुबन्धं तु}\ A1938; \skt{चानवस्थानं}\ N}}\skt{ करोति पुरुषे मदः ।}\pend

\pstart\versepar \ledsidenote{\passref{1.45.209}}\edtext{\edtext{\skt{अशौचनिद्रामात्सर्यमगम्यगमनः}}{\lemma{\skt{अशौचनिद्रामात्सर्यम्}}\Afootnote{\skt{अशौचनिद्रामात्सर्या}\ A1938; \skt{असौचनिद्रामात्सर्यम्}\ N H}}}{\lemma{\skt{अगम्यगमनः}}\Afootnote{N H; \skt{गम्यागमन}\ A1938}}\skt{ }\edtext{\skt{सदा}}{\Afootnote{H; \skt{लोलताः ।}\ A1938; \skt{यदा ।}\ N}}\skt{ । }\pend
\pstart\versepar \skt{असत्यभाषणञ्चापि }\edtext{\skt{कुर्याद्वै}}{\Afootnote{N H; \skt{कुर्याद्धि}\ A1938}}\skt{ तामसे मदः ।}\pend

\pstart\versepar \ledsidenote{\passref{1.45.218}}\edtext{\skt{तीक्ष्णं}}{\Afootnote{N H; \skt{भवतश्चात्र तत्सर्वं}\ A1938}}\skt{ }\edtext{\skt{कटुकं}}{\Afootnote{N H; \skt{कटु ती}\ A1938}}\skt{ }\edtext{\skt{मूत्रं}}{\Afootnote{N H; \skt{क्ष्णोष्णं}\ A1938}}\skt{ लवणानुरसं }\edtext{\skt{लघुः}}{\Afootnote{N H; \skt{लघु ।}\ A1938}}\skt{ । }\pend
\pstart\versepar \skt{}\edtext{\skt{शोधनं}}{\Afootnote{A1938 H; \skt{सोधनं}\ N}}\skt{ कफवातघ्नं }\edtext{\skt{शृणु}}{\Afootnote{N H; \skt{कृमि}\ A1938}}\skt{ }\edtext{\skt{तस्य}}{\Afootnote{N H; \skt{मेदो}\ A1938}}\skt{ }\edtext{\skt{विशेषणं}}{\Afootnote{H; \skt{विषापहम् ॥२१८॥}\ A1938; \skt{विसेषणं ।}\ N}}\skt{ ।}\pend

\pstart\versepar \ledsidenote{\passref{1.45.222}}\edtext{\skt{आनाहशोफगुल्मेषु}}{\Afootnote{H; \skt{दुर्नामोदरशूलेषु कुष्ठमेहाविशुद्धिषु । आनाहशोफगुल्मेषु}\ A1938; \skt{आनाहसोफगुल्मेषु}\ N}}\skt{ }\edtext{\skt{पाण्डुरोगे तु}}{\Afootnote{N H; \skt{पाण्डुरोगे च}\ A1938}}\skt{ माहिषं । }\pend
\pstart\versepar \skt{}\edtext{\skt{शोफघ्नमाजमौरभ्रं}}{\Afootnote{H\textsuperscript{pc}; \emph{om.}\ A1938 H\textsuperscript{ac}; \skt{शोफघ्नं माजमौरभ्रं}\ N}}\skt{ }\edtext{\skt{कासश्वासविषापहं}}{\Afootnote{H\textsuperscript{pc}; \emph{om.}\ A1938 H\textsuperscript{ac}; \skt{कासस्वासविसापहं ।}\ N}}\skt{ ।}\pend

\pstart \ledsidenote{\passref{1.45.225}}\edtext{\skt{आश्वङ्कफहरं}}{\lemma{\skt{आश्वङ्}}\Afootnote{H; \skt{दीपनं कटु तीक्ष्णोष्णं वातचेतोविकारनुत् । आश्वं}\ A1938; \skt{आस्वं}\ N}}\skt{ }\edtext{\skt{मूत्रं}}{\Afootnote{A1938 H; \skt{मूत्र}\ N}}\skt{ }\edtext{\skt{क्रिमिदर्द्रुषु}}{\Afootnote{N H; \skt{कृमिदद्रुषु}\ A1938}}\skt{ }\edtext{\skt{शस्यते}}{\Afootnote{A1938 H; \skt{सस्यते ।}\ N}}\skt{ ॥}\pend

\pstart \ledsidenote{\passref{1.45.226}}\edtext{\skt{तीक्ष्णं}}{\Afootnote{N H; \skt{सतिक्तं लवणं भेदि वातघ्नं पित्तकोपनम् । तीक्ष्णं}\ A1938}}\skt{ क्षारे किलासे च }\edtext{\skt{नागमूत्रम्प्रयोजयेत्}}{\lemma{\skt{नागमूत्रम्}}\Afootnote{N H; \skt{नागं मूत्रं}\ A1938}}\skt{ ।}\pend

\pstart \ledsidenote{\passref{1.45.228}}\edtext{\skt{अर्शोघ्नङ्कारभं}}{\lemma{\skt{अर्शोघ्नङ्}}\Afootnote{H; \skt{शोफकुष्ठोदरोन्मादमारुतक्रिमिनाशनम् । अर्शोघ्नं}\ A1938; \skt{अशोघ्नं}\ N}}\skt{ }\edtext{\skt{मूत्रम्मानुषन्तुविषापहं}}{\lemma{\skt{तुविषापहं}}\Afootnote{N H; \skt{च विषापहम् ॥२२८॥}\ A1938}}\skt{ ॥}\pend

\pstart\versepar \ledsidenote{\passref{1.45.229}}\skt{द्रवद्रव्याणि सर्व्वाणि }\edtext{\skt{कीर्त्तितानि}}{\Afootnote{N H; \skt{समासात्कीर्तितानि}\ A1938}}\skt{ }\edtext{\skt{समासतः}}{\Afootnote{N H; \skt{तु । काल}\ A1938}}\skt{ । }\pend
\pstart\versepar \skt{}\edtext{\skt{देशकालविभागज्ञो}}{\Afootnote{N H; \skt{देशविभागज्ञो}\ A1938}}\skt{ नृपतेः }\edtext{\edtext{\skt{कर्त्तुमर्हति}}{\lemma{\skt{कर्त्तुम्}}\Afootnote{N H; \skt{दातुम्}\ A1938}}}{\lemma{\skt{अर्हति}}\Afootnote{A1938; \skt{अर्हतीति ॥}\ N H}}\pend

\pstart \ledsidenote{\passref{1.45.trailer}}\edtext{\skt{॥ ४५ ॥}}{\Afootnote{\skt{इति सुश्रुतसंहितायां सूत्रस्थाने द्रवद्रव्यविज्ञानीयो नाम पञ्चचत्वारिंशोऽध्यायः ॥४५॥}\ A1938; \skt{फ ह ४० ५ ॥ ० ॥}\ N; \skt{फ छ् ४० ४}\ H}}\pend

\endnumbering
\beginnumbering
\eledsection*{Adhyāya 46}
\pstart \ledsidenote{\passref{1.46.1}}\edtext{\skt{अथातोऽन्नपानविधिम्व्याख्यास्यामः}}{\lemma{\skt{व्याख्यास्यामः}}\Afootnote{H; \skt{अध्यायं व्याख्यास्यामः ॥}\ A1938; \emph{om.}\ N}}\skt{ ॥}\pend

\pstart\versepar \ledsidenote{\passref{1.46.3}}\skt{धन्वन्तरिमभिवाद्य सुश्रुत उवाच ॥ }\pend
\pstart\versepar \skt{}\edtext{\edtext{\skt{भगवन्प्रागभिहितम्प्राणिनां}}{\lemma{\skt{भगवन्प्रागभिहितम्}}\Afootnote{H; \skt{प्रागभिहितः}\ A1938; \skt{भगवां प्रागभिहित}\ N}}}{\lemma{\skt{प्राणिनां}}\Afootnote{A1938 H; \skt{प्राणि}\ N}}\skt{ }\edtext{\edtext{\skt{मूलमाहारो}}{\lemma{\skt{मूलम्}}\Afootnote{H; \skt{पुनर्मूल}\ A1938; \skt{ना मूलम्}\ N}}}{\lemma{\skt{आहारो}}\Afootnote{N H; \skt{माहारो}\ A1938}}\skt{ }\edtext{\skt{बलवर्ण्णौजसाञ्च सः}}{\Afootnote{\skt{बलवर्णौजसां च, स}\ A1938; \skt{वलवर्णोजसां च सः}\ N; \skt{वलवर्ण्णौजसाञ्च सः}\ H}}\skt{ षट्सु }\edtext{\edtext{\skt{रसेष्वायत्तो}}{\lemma{\skt{रसेष्व्}}\Afootnote{N H; \skt{रसे}\ A1938}}}{\lemma{\skt{आयत्तो}}\Afootnote{N H; \skt{ष्वायत्तः,}\ A1938}}\skt{ रसाः }\edtext{\edtext{\skt{पुनर्द्रव्याश्रयिणः}}{\lemma{\skt{पुनर्}}\Afootnote{N H; \skt{पुन}\ A1938}}}{\lemma{\skt{द्रव्याश्रयिणः}}\Afootnote{N H; \skt{र्द्रव्याश्रयिणः,’}\ A1938}}\skt{ । }\edtext{\skt{द्रव्यरसवीर्यविपाकनिमित्ते च}}{\Afootnote{N H; \skt{द्रव्यरसगुणवीर्यविपाकनिमित्ते च}\ A1938}}\skt{ क्षयवृद्धी दोषाणां }\edtext{\skt{सात्म्यञ्च}}{\Afootnote{H; \skt{साम्यं च}\ A1938 N}}\skt{ । }\edtext{\edtext{\skt{ब्रह्मादेरपि}}{\lemma{\skt{ब्रह्मादेर्}}\Afootnote{\skt{ब्रह्मादे}\ A1938; \skt{व्रह्मादेर्}\ N H}}\skt{ च}}{\lemma{\skt{अपि च}}\Afootnote{N H; \skt{रपि च}\ A1938}}\skt{ लोकस्याहारः }\edtext{\edtext{\edtext{\edtext{\skt{स्थित्युत्पत्तिहेतुराहारादेवाभिवृद्धिर्बलमारोग्यवर्णेन्द्रियप्रसादश्च}}{\lemma{\skt{स्थित्युत्पत्तिहेतुर्}}\Afootnote{H; \skt{स्थित्युत्पत्तिविनाशहेतुः,}\ A1938; \skt{स्थित्युत्पत्तिहेतुर्}\ N}}}{\lemma{\skt{आहारादेवाभिवृद्धिर्}}\Afootnote{\skt{आहारादेवाभिवृद्धि}\ A1938; \skt{आहारादेवाभिवृद्धिर्}\ N H}}}{\lemma{\skt{बलम्}}\Afootnote{\skt{र्बल}\ A1938; \skt{वलम्}\ N H}}}{\lemma{\skt{आरोग्यवर्णेन्द्रियप्रसादश्च}}\Afootnote{N H; \skt{मारोग्यं वर्णेन्द्रियप्रसादश्च,}\ A1938}}\skt{ । }\edtext{\skt{तथा }\edtext{\edtext{\skt{ह्याहारवैषम्पादस्वास्थ्यन्तस्यासितपीतलीढखादितस्य}}{\lemma{\skt{आहारवैषम्पादस्वास्थ्यन्}}\Afootnote{H; \skt{ह्याहारवैषम्यादस्वास्थ्यं;}\ A1938; \skt{सवैषम्यादस्वास्थ्यं}\ N}}}{\lemma{\skt{तस्यासितपीतलीढखादितस्य}}\Afootnote{H; \skt{तस्याशितपीतलीढखादितस्य}\ A1938 N}}}{\lemma{\skt{तथा ह्य्}}\Afootnote{H; \skt{तथा}\ A1938; \skt{तथा र}\ N}}\skt{ }\edtext{\skt{नानाद्रव्यात्मकस्यानेकविधविकल्पप्रभावस्य}}{\Afootnote{H; \skt{नानाद्रव्यात्मकस्यानेकविधविकल्पस्यानेकविधप्रभावस्य}\ A1938; \skt{नानाद्रव्यात्मकस्यानेकविधविकल्प्रभावस्य}\ N}}\skt{ }\edtext{\skt{पृथक्पृथक्द्रव्यरसवीर्यविपाककर्म्मेच्छां}}{\Afootnote{H; \skt{पृथक्पृथग्द्रव्यरसगुणवीर्यविपाककर्माणीच्छामि}\ A1938; \skt{पृथक्पृथक्द्रव्यरसवीर्यविपाकर्मेच्छां}\ N}}\skt{ ज्ञातुं न }\edtext{\skt{ह्यनवबुद्धद्रव्यस्वभावा}}{\lemma{\skt{अनवबुद्धद्रव्यस्वभावा}}\Afootnote{\skt{अनवबुद्धस्वभावा}\ A1938; \skt{अनववुद्ध द्रव्य स्वभावा}\ N H}}\skt{ }\edtext{\skt{भिषजो}}{\Afootnote{A1938 H; \skt{भिष}\ N}}\skt{ }\edtext{\skt{रोगनिग्रहङ्कर्त्तुं}}{\lemma{\skt{रोगनिग्रहङ्}}\Afootnote{H; \skt{स्वस्थानुवृत्तिं रोगनिग्रहणं च}\ A1938; \skt{जरोगनिग्रहणं}\ N}}\skt{ समर्था }\edtext{\skt{इत्याहारायत्ताश्च}}{\lemma{\skt{इत्य्}}\Afootnote{N H; \emph{om.}\ A1938}}\skt{ }\edtext{\skt{प्राणिनो}}{\Afootnote{N H; \skt{सर्वप्राणिनो}\ A1938}}\skt{ }\edtext{\edtext{\skt{यस्मात्तस्मादन्नपानविधिमुपदिशत्तु}}{\lemma{\skt{यस्मात्तस्मादन्नपानविधिम्}}\Afootnote{A1938 H; \skt{यस्मात्तस्मादन्नपान्नविधिम्}\ N}}\skt{ मे}}{\lemma{\skt{उपदिशत्तु मे}}\Afootnote{H; \skt{उपदिशतु मे}\ A1938 N}}\skt{ भगवानित्युक्तः प्रोवाच भगवान्धन्वन्तरिः }\pend
\pstart\versepar \skt{॥ अथ खलु वत्स }\edtext{\skt{सुश्रुत}}{\Afootnote{N H; \skt{सुश्रुत! यथाप्रश्नमुच्यमानमुपधारयस्व- ॥३॥}\ A1938}}\skt{ ॥}\pend

\pstart \ledsidenote{\passref{1.46.4}}\skt{तत्र }\edtext{\skt{लोहितकशालिकलमककर्द्दमकपाण्डुकसुगन्धकसकुनाहृतकुष्माण्डकपुण्डरीकमहाशालिशीतभीरुकलोध्रशूकदीर्घशूककाञ्चनकहायनकदूषीमहाख्यप्रभृतयः}}{\Afootnote{H; \skt{लोहितशालिकलमकर्दमकपाण्डुकसुगन्धकशकुनाहृतपुष्पाण्डकपुण्डरीकमहाशालिशीतभीरुकरोध्रपुष्पक- दीर्घशूककाञ्चनकमहिषमहाशूकहायनकदूषकमहादूषकप्रभृतयः}\ A1938; \skt{लोहितशालिकलमकर्दमकपाण्डुकसुगन्धकशकुनाहृतकुष्माण्डकपुण्डरीकमहाशालिशीतभीरूकलोध्रपुष्पकदीर्घशूककाञ्चनकहायनकदूषीमहीप्रभृतयः}\ N}}\skt{ शालयः ॥}\pend

\pstart\versepar \ledsidenote{\passref{1.46.5}}\skt{मधुरा वीर्यतः शीता }\edtext{\skt{विपाका}}{\Afootnote{H; \skt{लघुपा}\ A1938; \skt{विपाक}\ N}}\skt{ }\edtext{\skt{कटुकाः}}{\Afootnote{H; \skt{का बला}\ A1938; \skt{कटूकाः}\ N}}\skt{ }\edtext{\skt{स्मृताः}}{\Afootnote{N H; \skt{वहाः ।}\ A1938}}\skt{ । }\pend
\pstart\versepar \skt{}\edtext{\skt{पित्तघ्नाल्पानिलकफः}}{\Afootnote{H; \skt{पित्तघ्नाल्पानिलकफाः}\ A1938 N}}\skt{ स्निग्धा }\edtext{\skt{वद्धाल्पवर्च्चसः}}{\Afootnote{N H; \skt{बद्धाल्पवर्चसः ॥५॥}\ A1938}}\skt{ ॥}\pend

\pstart\versepar \ledsidenote{\passref{1.46.6}}\skt{तेषां }\edtext{\skt{लोहितकः}}{\Afootnote{A1938 H; \skt{लोहितक}\ N}}\skt{ श्रेष्ठो दोषघ्नः शुक्रमूत्रलः । }\pend
\pstart\versepar \skt{चक्षुष्यो }\edtext{\skt{वर्ण्णवलकृत्स्वर्यो}}{\Afootnote{N H; \skt{वर्णबलकृत्स्वर्यो}\ A1938}}\skt{ }\edtext{\edtext{\skt{हृद्यस्तृषापहः}}{\lemma{\skt{हृद्यस्}}\Afootnote{N H; \skt{हृद्य}\ A1938}}}{\lemma{\skt{तृषापहः}}\Afootnote{N H; \skt{स्तृषापहः ॥६॥}\ A1938}}\skt{ ॥}\pend

\pstart \ledsidenote{\passref{1.46.7cd}}\skt{तस्मादल्पान्तरगुणाः क्रमशः }\edtext{\skt{शालयो}}{\Afootnote{A1938 H; \skt{शालयो}\ N}}\skt{ }\edtext{\skt{परे}}{\Afootnote{H; \skt{ऽवराः ॥७॥}\ A1938; \skt{परे ॥}\ N}}\skt{ ॥}\pend

\pstart \ledsidenote{\passref{1.46.8}}\edtext{\skt{षष्टिककाङ्गकमुकुन्दकपीतकप्रमोदककालकासनपुष्पकमहास्वेतकमहाव्रीहिचूर्ण्णककुरवकेदारप्रभृतयो}}{\Afootnote{H; \skt{षष्टिककाङ्गुकमुकुन्दकपीतकप्रमोदककाकलकासनपुष्पकमहाषष्टिकचूर्णककुरवककेदारप्रभृतयः}\ A1938; \skt{षष्टिककाङ्गुकमुकुन्दकपीतकप्रमोदककालकासनपुष्पकमहास्वेतकमहाव्रीहिचूर्ण्णककुरवकेकेदारकप्रभृतयो}\ N}}\skt{ }\edtext{\skt{व्रीहयः}}{\Afootnote{N H; \skt{षष्टिकाः ॥८॥}\ A1938}}\skt{ ॥}\pend

\pstart\versepar \ledsidenote{\passref{1.46.9}}\skt{रसे पाके च मधुराः }\edtext{\skt{पित्तानिलहराः}}{\Afootnote{N H; \skt{शमना वातपित्त}\ A1938}}\skt{ }\edtext{\skt{स्मृताः}}{\Afootnote{N H; \skt{योः ।}\ A1938}}\skt{ । }\pend
\pstart\versepar \skt{}\edtext{\edtext{\skt{शालीनाङ्गुणतश्चापि}}{\lemma{\skt{शालीनाङ्गुणतश्}}\Afootnote{N H; \skt{शालीनां च गुणैस्तुल्या बृंहणाः कफशुक्रलाः ॥९॥}\ A1938}}}{\lemma{\skt{चापि}}\Afootnote{N H; \emph{om.}\ A1938}}\skt{ }\edtext{\skt{समाना}}{\Afootnote{N H; \emph{om.}\ A1938}}\skt{ }\edtext{\skt{वद्धवर्च्चसः}}{\Afootnote{N H; \emph{om.}\ A1938}}\skt{ ॥}\pend

\pstart\versepar \ledsidenote{\passref{1.46.10}}\edtext{\skt{षष्टिकः}}{\Afootnote{A1938 H; \skt{षष्ठिकः}\ N}}\skt{ प्रवरस्तेषाङ्कषायानुरसो लघुः । }\pend
\pstart\versepar \skt{मृदुः }\edtext{\edtext{\skt{स्निग्धस्त्रिदोषघ्नः}}{\lemma{\skt{स्निग्धस्}}\Afootnote{N H; \skt{स्निग्ध}\ A1938}}}{\lemma{\skt{त्रिदोषघ्नः}}\Afootnote{N H; \skt{स्त्रिदोषघ्नः}\ A1938}}\skt{ }\edtext{\skt{स्थैर्यकृद्वलवर्द्धनः}}{\Afootnote{H; \skt{स्थैर्यकृद्बलवर्धनः ॥१०॥}\ A1938; \skt{स्थैर्यकृद्वलवर्द्धनः ।}\ N}}\skt{ ॥}\pend

\pstart\versepar \ledsidenote{\passref{1.46.11}}\edtext{\skt{रसतो}}{\Afootnote{N H; \skt{विपाके}\ A1938}}\skt{ }\edtext{\skt{मधुरग्राही}}{\Afootnote{N H; \skt{मधुरो ग्राही}\ A1938}}\skt{ तुल्यो लोहितशालिभिः । }\pend
\pstart\versepar \skt{}\edtext{\edtext{\edtext{\skt{शेषास्त्वल्पान्तरास्त्वस्माद्व्रीहयः}}{\lemma{\skt{शेषास्}}\Afootnote{H; \skt{शेषा}\ A1938 N}}}{\lemma{\skt{त्वल्पान्तरास्}}\Afootnote{H; \skt{स्त्वल्पान्तरा}\ A1938; \skt{त्स्वल्पान्तरी}\ N}}}{\lemma{\skt{त्वस्माद्व्रीहयः}}\Afootnote{H; \skt{स्तस्मात्षष्टिकाः}\ A1938; \skt{स्मात्व्रीहयः}\ N}}\skt{ क्रमशो गुणैः ॥}\pend

\pstart \ledsidenote{\passref{1.46.12}}\edtext{\skt{कृष्णव्रीहिशालामुखलावाक्षजातुमुखनन्दीमुखत्वरितककुर्क्कुटाण्डकपारावतकपाटलाप्रभृतयो}}{\Afootnote{H; \skt{कृष्णव्रीहिशालामुखजतुमुखनन्दीमुखलावाक्षकत्वरितककुक्कुटाण्डकपारावतकपाटलप्रभृतयो}\ A1938; \skt{कृष्णव्रीहिशालामुखलावाक्षजतुमुखनन्दीमुखत्वरितककुटाण्डकपारावतकपाटलो प्रभृतयोः}\ N}}\skt{ व्रीहयः ॥}\pend

\pstart\versepar \ledsidenote{\passref{1.46.13}}\skt{कषायमधुराः पाके }\edtext{\skt{मधुरा}}{\Afootnote{N H; \skt{ऽमधुरा}\ A1938}}\skt{ }\edtext{\skt{वीर्यतोऽहिमाः}}{\lemma{\skt{ऽहिमाः}}\Afootnote{A1938 H; \skt{हिमाः ।}\ N}}\skt{ । }\pend
\pstart\versepar \skt{}\edtext{\skt{अल्पाभिष्यन्दिनस्तुल्याः}}{\lemma{\skt{अल्पाभिष्यन्दिनस्}}\Afootnote{A1938 H; \skt{अल्पाभिष्यन्तिनस्}\ N}}\skt{ }\edtext{\edtext{\skt{षष्टिकैर्व्वलवर्द्धनाः}}{\lemma{\skt{षष्टिकैर्}}\Afootnote{A1938 H; \skt{षष्ठिकैर्}\ N}}}{\lemma{\skt{व्वलवर्द्धनाः}}\Afootnote{N H; \skt{बद्धवर्चसः ॥१३॥}\ A1938}}\skt{ ॥}\pend

\pstart\versepar \ledsidenote{\passref{1.46.14}}\edtext{\skt{कृष्णव्रीहिर्व्वरस्तेषां}}{\lemma{\skt{कृष्णव्रीहिर्}}\Afootnote{A1938 H; \skt{कृष्णव्रीहि}\ N}}\skt{ }\edtext{\skt{कषायमधुरो}}{\Afootnote{H; \skt{कषायानुरसो}\ A1938 N}}\skt{ लघुः । }\pend
\pstart\versepar \skt{}\edtext{\skt{तस्मादल्पान्तरे}}{\Afootnote{H; \skt{तस्मादल्पान्तर}\ A1938 N}}\skt{ गुणाः क्रमशो }\edtext{\edtext{\skt{व्रीहयोऽपरे}}{\lemma{\skt{व्रीहयो}}\Afootnote{A1938 H; \skt{व्रीहयो}\ N}}}{\lemma{\skt{ऽपरे}}\Afootnote{A1938 H; \skt{परे ॥}\ N}}\skt{ ॥}\pend

\pstart\versepar \ledsidenote{\passref{1.46.15}}\edtext{\skt{दग्धायामवनौ}}{\Afootnote{H; \skt{दग्धायामवनौ}\ A1938; \skt{गग्धायामवनौ}\ N}}\skt{ जाताः शालयो लघुपाकिनः । }\pend
\pstart\versepar \skt{कषाया }\edtext{\skt{वद्धविण्मूत्रा}}{\Afootnote{H; \skt{बद्धविण्मूत्रा}\ A1938; \skt{वद्धविन्मूत्रा । ः}\ N}}\skt{ रूक्षाः श्लेष्मापकर्षणाः ॥}\pend

\pstart\versepar \ledsidenote{\passref{1.46.16}}\skt{स्थलजाः कफपित्तघ्नाः कषायाः कटुकान्वयाः । }\pend
\pstart\versepar \skt{किञ्चित्सतिक्तमधुराः पवनानलवर्द्धनाः ॥}\pend

\pstart\versepar \ledsidenote{\passref{1.46.17}}\skt{कैदारा मधुरा वृष्या }\edtext{\skt{बल्याः}}{\Afootnote{A1938 H; \skt{वल्याः}\ N}}\skt{ }\edtext{\skt{पित्तनिवर्हणाः}}{\Afootnote{N H; \skt{पित्तनिबर्हणाः ।}\ A1938}}\skt{ । }\pend
\pstart\versepar \skt{}\edtext{\skt{ईषत्कषायाल्पवला}}{\Afootnote{N H; \skt{ईषत्कषायाल्पमला}\ A1938}}\skt{ गुरवः कफशुक्रलाः ॥}\pend

\pstart\versepar \ledsidenote{\passref{1.46.18}}\edtext{\skt{रोप्यातिरोप्या}}{\Afootnote{A1938 H; \skt{रोप्यादिरोप्य}\ N}}\skt{ लघवः }\edtext{\skt{शीघ्रपाका}}{\Afootnote{A1938 H; \skt{शीघ्रपाक}\ N}}\skt{ गुणोत्तराः । }\pend
\pstart\versepar \skt{}\edtext{\skt{अविदाहिनो}}{\Afootnote{N H; \skt{अदाहिनो}\ A1938}}\skt{ }\edtext{\skt{वातहरा}}{\Afootnote{N H; \skt{दोषहरा}\ A1938}}\skt{ }\edtext{\skt{वल्या}}{\Afootnote{N H; \skt{बल्या}\ A1938}}\skt{ मूत्रविवर्धनाः ॥}\pend

\pstart\versepar \ledsidenote{\passref{1.46.19}}\skt{शालयश्छिन्नरूढा ये }\edtext{\skt{रूक्षा}}{\Afootnote{H; \skt{रूक्षास्ते}\ A1938 N}}\skt{ }\edtext{\skt{वद्धाल्पवर्च्चसः}}{\Afootnote{N H; \skt{बद्धवर्चसः ।}\ A1938}}\skt{ । }\pend
\pstart\versepar \skt{}\edtext{\skt{तिक्ताकषायाः}}{\Afootnote{H; \skt{तिक्ताः कषायाः}\ A1938 N}}\skt{ पित्तघ्ना }\edtext{\skt{लघुपाकाः}}{\Afootnote{A1938 H; \skt{लघुपाको}\ N}}\skt{ }\edtext{\skt{कफापहाः}}{\Afootnote{N H; \skt{कफापहाः ॥१९॥ इति शालिवर्गः ।}\ A1938}}\skt{ ॥}\pend

\pstart\versepar \ledsidenote{\passref{1.46.20}}\edtext{\edtext{\skt{विस्तरेणायमुद्दिष्टः}}{\lemma{\skt{विस्तरेणायम्}}\Afootnote{N H; \skt{विस्तरेणाय}\ A1938}}}{\lemma{\skt{उद्दिष्टः}}\Afootnote{N H; \skt{मुद्दिष्टः}\ A1938}}\skt{ शालिवर्ग्गो हिताहितः । }\pend
\pstart\versepar \skt{}\edtext{\skt{तद्वत्कुधान्यमुद्गानाम्माषादीनाञ्च}}{\lemma{\skt{तद्वत्कुधान्यमुद्गानाम्}}\Afootnote{H; \skt{अथ कुधान्यवर्गः । तद्वत्कुधान्यमुद्गादि}\ A1938; \skt{तद्वत्कुधान्यमुद्गानाम्}\ N}}\skt{ वक्ष्यते ॥}\pend

\pstart \ledsidenote{\passref{1.46.21}}\edtext{\skt{कोरदूषश्यामाकनीवारशान्तनुवरकोद्दालकप्रियङ्गुमधूलिकानन्दीमुखकुरविन्दकसक}}{\Afootnote{H; \skt{कोरदूषकश्यामाकनीवारशान्तनुवरकोद्दालकप्रियङ्गुमधूलिकानन्दीमुखी- कुरुविन्दगवेधुकसर}\ A1938; \skt{कोरदूषकश्यामाकनीवारशान्तनुवरकोद्दालकप्रियङ्गुमधूलिकानन्दीमुखीकुरूविन्दकसक}\ N}}\skt{ }\edtext{\skt{वरूकतोलपर्ण्णामुकुन्दकवेणुयवप्रभृतयः}}{\Afootnote{H; \skt{बरुकतोद(य)पर्णीमुकुन्दकवेणुयवप्रभृतयः}\ A1938; \skt{वस्ताकतोलपर्ण्णीमुकुन्दकवेणुयवप्रभृतयः ।}\ N}}\skt{ कुधान्यविशेषाः ॥}\pend

\pstart\versepar \ledsidenote{\passref{1.46.22}}\skt{उष्णाः कषायमधुरा रूक्षाः }\edtext{\skt{कटुविपाकिनः}}{\Afootnote{A1938 H; \skt{कटुपाकिनः}\ N}}\skt{ । }\pend
\pstart\versepar \skt{श्लेष्मघ्ना }\edtext{\skt{वातनिष्यन्दा}}{\Afootnote{N H; \skt{बद्धनिस्यन्दा}\ A1938}}\skt{ }\edtext{\skt{वातपित्तप्रकोपनाः}}{\Afootnote{N H; \skt{वातपित्तप्रकोपणाः ॥२२॥}\ A1938}}\skt{ ॥}\pend

\pstart\versepar \ledsidenote{\passref{1.46.23}}\edtext{\skt{कषायमधूराः}}{\Afootnote{H; \skt{कषायमधुरस्}\ A1938; \skt{कषायमधुराः}\ N}}\skt{ }\edtext{\skt{शीतास्तेषां}}{\Afootnote{H; \skt{तेषां}\ A1938; \skt{शीतास्तेषां}\ N}}\skt{ }\edtext{\skt{पित्तसपाः}}{\Afootnote{H; \skt{शीतः पित्तापहः}\ A1938; \skt{पित्तसहाः}\ N}}\skt{ }\edtext{\skt{स्मृताः}}{\Afootnote{N H; \skt{स्मृतः ।}\ A1938}}\skt{ । }\pend
\pstart\versepar \skt{}\edtext{\skt{सकोरदूषणामाको}}{\Afootnote{\skt{कोद्रवश्च सनीवारः}\ A1938; \skt{सकोरदूषश्यामाकको नीवारश्}\ N; \skt{सकोरदूषणामा}\ H}}\skt{ }\edtext{\edtext{\skt{नीरश्चैव}}{\lemma{\skt{नीरश्}}\Afootnote{\skt{श्यामाकश्}\ A1938; \emph{om.}\ N; \skt{कोनीरश्}\ H}}}{\lemma{\skt{चैव}}\Afootnote{H; \skt{च स}\ A1938; \skt{च}\ N}}\skt{ शान्तनुः ॥}\pend

\pstart\versepar \ledsidenote{\passref{1.46.24}}\skt{कृष्णा रक्ताश्च पीताश्च श्वेताश्चैव प्रियङ्गवः । }\pend
\pstart\versepar \skt{यथोत्तरं }\edtext{\skt{प्रधानाः स्यूः}}{\Afootnote{\skt{प्रधानाः स्यू}\ A1938; \skt{प्रधानाः स्यु}\ N; \skt{प्रधानाः स्युः}\ H}}\skt{ }\edtext{\skt{स्निग्धाः}}{\Afootnote{N H; \skt{रूक्षाः}\ A1938}}\skt{ }\edtext{\skt{कफकराः}}{\Afootnote{N H; \skt{कफहराः}\ A1938}}\skt{ }\edtext{\skt{सराः}}{\Afootnote{N\textsuperscript{pc} H; \skt{स्मृताः ॥२४॥}\ A1938; \skt{साराः ।}\ N\textsuperscript{ac}}}\skt{ ।}\pend

\pstart\versepar \ledsidenote{\passref{1.46.25}}\skt{मधूली }\edtext{\skt{मधुराः}}{\Afootnote{N H; \skt{मधुरा}\ A1938}}\skt{ }\edtext{\skt{शीताः}}{\Afootnote{N H; \skt{शीता}\ A1938}}\skt{ स्निग्धा नन्दीमुखी तथा । }\pend
\pstart\versepar \skt{विशोषी तत्र }\edtext{\skt{भूयिष्ठवरकः}}{\Afootnote{H; \skt{भूयिष्ठं वरुकः}\ A1938; \skt{भूयिष्ठवरुकः}\ N}}\skt{ समुकुन्दकः ॥}\pend

\pstart\versepar \ledsidenote{\passref{1.46.26}}\skt{रूक्षा }\edtext{\skt{वेणुयवा}}{\Afootnote{A1938 H; \skt{वेणु ये}\ N}}\skt{ ज्ञेया वीर्योष्णा कटुपाकिनः । }\pend
\pstart\versepar \skt{}\edtext{\skt{वद्धमूत्रः}}{\Afootnote{\skt{बद्धमूत्राः}\ A1938; \skt{वद्धमूत्राः}\ N H}}\skt{ कफहराः कषाया वातकोपनाः ॥}\pend

\pstart \ledsidenote{\passref{1.46.27}}\edtext{\skt{मुद्गवनमुद्गसूरमुकुष्ठकलायहरेण्वाढकीसतीना}}{\Afootnote{H; \skt{मुद्गवनमुद्गकलायमकुष्ठमसूरमङ्गल्यचणकसतीनत्रिपुटकहरेण्वाढकीप्रभृतयो}\ A1938; \skt{मुद्गवनमुद्गमसूरमुद्गकुष्ठककलायहरेण्ढकीसतीना}\ N}}\skt{ वैदलाः ॥}\pend

\pstart\versepar \ledsidenote{\passref{1.46.28}}\edtext{\skt{कषायमधुराः}}{\Afootnote{A1938 H; \skt{कषायो मधुराः}\ N}}\skt{ }\edtext{\skt{शीता}}{\Afootnote{H\textsuperscript{ac}; \skt{शीताः}\ A1938 N; \skt{शीताᳵ}\ H\textsuperscript{pc}}}\skt{ }\edtext{\skt{कटुपाकानिलापहाः}}{\Afootnote{N H; \skt{कटुपाका मरुत्कराः ।}\ A1938}}\skt{ । }\pend
\pstart\versepar \skt{}\edtext{\skt{वद्धमूत्रपुरीषाश्च}}{\Afootnote{N H; \skt{बद्धमूत्रपुरीषाश्च}\ A1938}}\skt{ पित्तश्लेष्महरास्तथा ॥}\pend

\pstart\versepar \ledsidenote{\passref{1.46.29}}\edtext{\edtext{\edtext{\skt{नात्यर्थवातलास्तेषाम्मुद्गादृष्ठि}}{\lemma{\skt{नात्यर्थवातलास्}}\Afootnote{H; \skt{नात्यर्थं वातला}\ A1938; \skt{नात्यर्थं वातलास्}\ N}}}{\lemma{\skt{तेषाम्}}\Afootnote{N H; \skt{स्तेषु}\ A1938}}}{\lemma{\skt{मुद्गादृष्ठि}}\Afootnote{H; \skt{मुद्गा दृष्टि}\ A1938; \skt{मुद्गदृष्टि}\ N}}\skt{ प्रसादनाः । }\pend
\pstart\versepar \skt{प्रधाना हरितास्तत्र वन्या मुद्गसमाः स्मृताः ॥}\pend

\pstart\versepar \ledsidenote{\passref{1.46.30}}\skt{विपाके }\edtext{\skt{मधुराः}}{\Afootnote{A1938; \skt{मधुरो}\ N; \skt{मधुराः प्रो}\ H}}\skt{ }\edtext{\skt{प्रोक्ता}}{\Afootnote{A1938 N; \skt{क्ता}\ H}}\skt{ मसूरा }\edtext{\skt{वद्ध वर्च्चसः}}{\Afootnote{H\textsuperscript{pc}; \skt{बद्धवर्चसः}\ A1938; \skt{वर्द्धवर्च्चसः ।}\ N; \skt{वर्च्चसः ।}\ H\textsuperscript{ac}}}\skt{ । }\pend
\pstart\versepar \skt{}\edtext{\skt{मौकुष्थकाः}}{\Afootnote{H; \skt{मकुष्ठकाः}\ A1938 N}}\skt{ }\edtext{\skt{क्रिमिहराः}}{\Afootnote{H; \skt{कृमिकराः}\ A1938; \skt{क्रिमिहः}\ N}}\skt{ कलायाः }\edtext{\skt{प्रचूरानिलाः}}{\Afootnote{H; \skt{प्रचुरानिलाः}\ A1938 N}}\skt{ ॥}\pend

\pstart\versepar \ledsidenote{\passref{1.46.33}}\skt{हरेणवः सतीनाश्च विज्ञेया }\edtext{\skt{भिन्नवर्च्चसः}}{\Afootnote{N H; \skt{बद्धवर्चसः}\ A1938}}\skt{ । }\pend
\pstart\versepar \skt{ऋते }\edtext{\edtext{\skt{मुद्गमसूराभ्यामन्ये}}{\lemma{\skt{मुद्गमसूराभ्याम्}}\Afootnote{N H; \skt{मुद्गमसूराभ्या}\ A1938}}}{\lemma{\skt{अन्ये}}\Afootnote{N H; \skt{मन्ये}\ A1938}}\skt{ }\edtext{\skt{चाध्मानकारकाः}}{\Afootnote{H; \skt{त्वाध्मानकारकाः ॥३३॥}\ A1938; \skt{त्वाध्य्मानकारकाः ।}\ N}}\skt{ ॥}\pend

\pstart\versepar \ledsidenote{\passref{1.46.34}}\edtext{\skt{माषो}}{\Afootnote{A1938 H; \skt{माषो}\ N}}\skt{ }\edtext{\skt{गुरुर्ब्भिन्नपुरीषमूत्रः}}{\Afootnote{A1938; \skt{गुरूर्भिन्नपुरीषमूत्रः ।}\ N; \skt{गुरुर्ब्भिन्नपुरीषमूत्रः}\ H}}\skt{ }\edtext{\edtext{\skt{स्निग्धोऽथ}}{\lemma{\skt{स्निग्धो}}\Afootnote{H; \skt{स्निग्धो}\ A1938 N}}}{\lemma{\skt{ऽथ}}\Afootnote{H; \skt{ष्ण}\ A1938; \skt{ष}\ N}}\skt{ वृष्यो }\edtext{\edtext{\skt{मधुरोऽनिलघ्नः}}{\lemma{\skt{मधुरो}}\Afootnote{A1938 H; \skt{मधुरो}\ N}}}{\lemma{\skt{ऽनिलघ्नः}}\Afootnote{A1938 H; \skt{निलघ्नः ।}\ N}}\skt{ । }\pend
\pstart\versepar \skt{}\edtext{\skt{सन्तर्प्पणः}}{\Afootnote{A1938 H; \skt{सन्तर्पण}\ N}}\skt{ स्तन्यकरो }\edtext{\skt{विशेषाद्वलप्रदः}}{\Afootnote{N H; \skt{विशेषाद्बलप्रदः}\ A1938}}\skt{ शुक्रकफावहश्च ॥}\pend

\pstart\versepar \ledsidenote{\passref{1.46.35}}\edtext{\skt{कषायभावान्न}}{\Afootnote{A1938 H; \skt{कषायो भावान्न}\ N}}\skt{ पुरीषभेदी न मूत्रलो नैव }\edtext{\skt{लवास}}{\Afootnote{H; \skt{कफस्य}\ A1938; \skt{वलास}\ N}}\skt{ कर्त्ता । }\pend
\pstart\versepar \skt{}\edtext{\skt{स्वादुर्व्रिपाके}}{\Afootnote{H; \skt{स्वादुर्विपाके}\ A1938; \skt{स्वादुर्विपाको}\ N}}\skt{ }\edtext{\skt{मधुरोनसान्द्रः}}{\Afootnote{H; \skt{मधुरोऽलसान्द्रः}\ A1938; \skt{मधुरोलसान्द्रः}\ N}}\skt{ }\edtext{\skt{सन्तर्प्पणस्तन्यरुचिप्रदश्च}}{\lemma{\skt{सन्तर्प्पणस्तन्यरुचिप्रदश्}}\Afootnote{N H; \skt{सन्तर्पणः स्तन्यरुचिप्रदश्}\ A1938}}\skt{ ॥}\pend

\pstart\versepar \ledsidenote{\passref{1.46.36}}\skt{माषैः समानं }\edtext{\edtext{\edtext{\skt{फलमात्मगुप्तमुक्तञ्च}}{\lemma{\skt{फलम्}}\Afootnote{N H; \skt{फल}\ A1938}}}{\lemma{\skt{आत्मगुप्तम्}}\Afootnote{N H; \skt{मात्मगुप्त}\ A1938}}}{\lemma{\skt{उक्तञ्च}}\Afootnote{N H; \skt{मुक्तं च}\ A1938}}\skt{ }\edtext{\skt{काकाण्डुफलं}}{\Afootnote{H; \skt{काकाण्डफलं}\ A1938; \skt{काकण्डफलन्}\ N}}\skt{ तथैव । }\pend
\pstart\versepar \skt{आरण्यमाषा गुणतः प्रदिष्टा }\edtext{\skt{रूक्षाः}}{\Afootnote{A1938 H; \skt{रूक्षा}\ N}}\skt{ कषाया अविदाहिनश्च ॥}\pend

\pstart\versepar \ledsidenote{\passref{1.46.37}}\edtext{\skt{उष्णः}}{\Afootnote{A1938 H; \skt{उष्ण}\ N}}\skt{ कुलत्थो रसतः }\edtext{\skt{कषायः}}{\Afootnote{A1938 H; \skt{कर्षयः}\ N}}\skt{ }\edtext{\skt{कटूर्व्विपाके}}{\Afootnote{H; \skt{कटुर्विपाके}\ A1938; \skt{कटुर्विपाके}\ N}}\skt{ कफमारुतघ्नः । }\pend
\pstart\versepar \skt{}\edtext{\skt{शिक्राश्मरीगुल्मनिसूदनश्च}}{\Afootnote{H; \skt{शुक्राश्मरीगुल्मनिषूदनश्च}\ A1938 N}}\skt{ }\edtext{\skt{साङ्ग्राहिकः}}{\Afootnote{A1938 H; \skt{संग्राहिकः}\ N}}\skt{ पीनसकासहन्ता ॥}\pend

\pstart\versepar \ledsidenote{\passref{1.46.38}}\edtext{\skt{आनाहमेदो}}{\Afootnote{H; \skt{आनाहमेदो}\ A1938 N}}\skt{ }\edtext{\skt{गुदकीलहिक्काश्वासापहः}}{\Afootnote{A1938 H; \skt{गुदकीलहिक्का । स्वासापहः}\ N}}\skt{ शोणितपित्तकृच्च । }\pend
\pstart\versepar \skt{}\edtext{\skt{वलास}}{\Afootnote{N H; \skt{कफस्य}\ A1938}}\skt{ हन्ता नयनामयघ्नो विशेषतो वन्यकुलत्थ उक्तः ॥}\pend

\pstart\versepar \ledsidenote{\passref{1.46.39}}\edtext{\skt{ईषत्कषायो}}{\Afootnote{A1938 H; \skt{ईषत्कषायोः}\ N}}\skt{ मधुरः सतिक्तः }\edtext{\skt{सांग्राहिकः}}{\Afootnote{A1938 H; \skt{संग्राहिकः}\ N}}\skt{ }\edtext{\edtext{\skt{पित्तकरास्तथोष्णः}}{\lemma{\skt{पित्तकरास्}}\Afootnote{H; \skt{पित्तकर}\ A1938 N}}}{\lemma{\skt{तथोष्णः}}\Afootnote{H; \skt{स्तथोष्णः}\ A1938 N}}\skt{ । }\pend
\pstart\versepar \skt{तिलो }\edtext{\skt{विपके}}{\Afootnote{H; \skt{विपाके}\ A1938 N}}\skt{ }\edtext{\skt{कटूको}}{\Afootnote{H; \skt{मधुरो}\ A1938; \skt{कटुको}\ N}}\skt{ }\edtext{\skt{वलिष्ठः}}{\Afootnote{N H; \skt{बलिष्ठः}\ A1938}}\skt{ स्निग्धो }\edtext{\skt{व्रणो}}{\Afootnote{H; \skt{व्रणे}\ A1938 N}}\skt{ लेपन }\edtext{\skt{पथ्यादुक्तः}}{\Afootnote{H; \skt{एव पथ्यः ॥३९॥}\ A1938; \skt{पथ्य उक्तः ।}\ N}}\skt{ ॥}\pend

\pstart\versepar \ledsidenote{\passref{1.46.40}}\edtext{\edtext{\edtext{\skt{दन्त्योऽग्निमेधाजननोऽल्पमूत्रः}}{\lemma{\skt{दन्त्यो}}\Afootnote{A1938 H; \skt{दन्त्यो}\ N}}}{\lemma{\skt{ऽग्निमेधाजननो}}\Afootnote{A1938 H; \skt{ग्निमेधाजननो}\ N}}}{\lemma{\skt{ऽल्पमूत्रः}}\Afootnote{H; \skt{ऽल्पमूत्र}\ A1938; \skt{ल्पमूत्रः ।}\ N}}\skt{ }\edtext{\edtext{\skt{स्वर्योऽथ}}{\lemma{\skt{स्वर्यो}}\Afootnote{H; \skt{स्त्वच्यो}\ A1938; \emph{om.}\ N}}}{\lemma{\skt{ऽथ}}\Afootnote{A1938 H; \skt{र्याष}\ N}}\skt{ }\edtext{\edtext{\skt{केश्योऽनिलहा}}{\lemma{\skt{केश्यो}}\Afootnote{A1938 H; \skt{केश्यो}\ N}}}{\lemma{\skt{ऽनिलहा}}\Afootnote{A1938 H; \skt{निलहा}\ N}}\skt{ गुरुश्च । }\pend
\pstart\versepar \skt{तिलेषु }\edtext{\edtext{\skt{सर्व्वषेव}}{\lemma{\skt{सर्व्वष्}}\Afootnote{H; \skt{सर्वे}\ A1938; \skt{सर्वेष्}\ N}}}{\lemma{\skt{एव}}\Afootnote{N H; \skt{ष्व}\ A1938}}\skt{ सितः प्रधानो मध्यः सितो }\edtext{\edtext{\skt{हीनतरास्तथान्ये}}{\lemma{\skt{हीनतरास्}}\Afootnote{A1938 H; \skt{हीनतरा}\ N}}}{\lemma{\skt{तथान्ये}}\Afootnote{A1938 H; \skt{स्तथान्ये ॥}\ N}}\skt{ ॥}\pend

\pstart\versepar \ledsidenote{\passref{1.46.41}}\edtext{\skt{यवxकषायो}}{\Afootnote{H; \skt{यवः कषायो}\ A1938 N}}\skt{ मधुरो हिमश्च }\edtext{\edtext{\skt{कटुर्व्विपाके}}{\lemma{\skt{कटुर्}}\Afootnote{N H; \skt{कटु}\ A1938}}}{\lemma{\skt{व्विपाके}}\Afootnote{N H; \skt{र्विपाके}\ A1938}}\skt{ }\edtext{\skt{कफपित्तहन्ता}}{\Afootnote{N H; \skt{कफपित्तहारी ।}\ A1938}}\skt{ । }\pend
\pstart\versepar \skt{व्रणेषु }\edtext{\skt{पथ्यास्तिलवच्च}}{\Afootnote{H; \skt{पथ्यस्तिलवच्च}\ A1938 N}}\skt{ }\edtext{\skt{नित्यं}}{\Afootnote{A1938 H; \skt{नित्य्}\ N}}\skt{ }\edtext{\skt{प्रवृद्धमूत्रो}}{\Afootnote{H; \skt{प्रबद्धमूत्रो}\ A1938; \skt{अप्रवद्धमूत्रो}\ N}}\skt{ }\edtext{\skt{बहुवातवर्च्चाः}}{\Afootnote{A1938; \skt{वहुवातवर्च्चाः}\ N H}}\skt{ ॥}\pend

\pstart\versepar \ledsidenote{\passref{1.46.42}}\skt{स्थैर्याग्निमेधास्वरवर्ण्णकृच्च }\edtext{\skt{सपिच्छिलः}}{\Afootnote{A1938 H; \skt{सपिच्छिलो}\ N}}\skt{ स्थूलविलेखनश्च । }\pend
\pstart\versepar \skt{}\edtext{\edtext{\edtext{\skt{मेदोऽनिलस्तृतृरणोऽथ}}{\lemma{\skt{मेदो}}\Afootnote{H; \skt{मेदो}\ A1938 N}}}{\lemma{\skt{ऽनिलस्तृतृरणो}}\Afootnote{H; \skt{मरुत्तृढ्हरणो}\ A1938; \skt{निल}\ N}}}{\lemma{\skt{ऽथ}}\Afootnote{H; \skt{ऽति}\ A1938; \skt{स्रणो}\ N}}\skt{ }\edtext{\skt{रूक्षः}}{\Afootnote{A1938 H; \skt{रूक्ष्यः}\ N}}\skt{ प्रसादनः शोणितपित्तयोश्च ॥}\pend

\pstart\versepar \ledsidenote{\passref{1.46.43}}\edtext{\skt{एभिर्ग्गुणैर्हीनतरैश्च}}{\lemma{\skt{हीनतरैश्च}}\Afootnote{N H; \skt{हीनतरैस्तु}\ A1938}}\skt{ }\edtext{\edtext{\skt{किञ्चिद्विद्याद्यवेभ्योऽपि}}{\lemma{\skt{किञ्चिद्विद्याद्यवेभ्यो}}\Afootnote{A1938 H; \skt{किञ्चिद्विद्याद्यवेभ्यो}\ N}}}{\lemma{\skt{ऽपि}}\Afootnote{H; \skt{ऽति}\ A1938; \skt{पि}\ N}}\skt{ यवानशेषान् । }\pend
\pstart\versepar \skt{गोधूम उक्तो मधुरो गुरुश्च }\edtext{\skt{वल्यः}}{\Afootnote{N H; \skt{बल्यः}\ A1938}}\skt{ स्थिरः }\edtext{\skt{शुक्रवलप्रदश्च}}{\lemma{\skt{शुक्रवलप्रदश्}}\Afootnote{N H; \skt{शुक्ररुचिप्रदश्}\ A1938}}\skt{ ॥}\edtext{}{\lemma{\skt{यवानशेषान् । गोधूम}}\Afootnote{H; \skt{यवानशेषैः गोधूम}\ A1938; \skt{यवानशेषान ॥ गोधूम}\ N}}\pend

\pstart\versepar \ledsidenote{\passref{1.46.44}}\edtext{\skt{स्निग्धोऽथ}}{\lemma{\skt{ऽथ}}\Afootnote{N H; \skt{ऽति}\ A1938}}\skt{ शीतोऽनिलपित्तहन्ता }\edtext{\skt{सन्धानकृच्छ्लेष्मकरः}}{\Afootnote{H; \skt{सन्धानकृच्छ्लेष्मकरः}\ A1938; \skt{सन्धानकृच्छ श्लेष्मकरः}\ N}}\skt{ सरश्च । }\pend
\pstart\versepar \skt{रूक्षः कषायो }\edtext{\skt{विषशोखशुक्रवलासदृष्टिक्षयकृद्विदाही}}{\Afootnote{H; \skt{विषशोषशुक्रबलासदृष्टिक्षयकृद्विदाही ॥४४॥}\ A1938; \skt{विषशोषशुक्रवलासदृष्टिक्षयकृद्विदाही ॥}\ N}}\skt{ ॥}\pend

\pstart\versepar \ledsidenote{\passref{1.46.45}}\edtext{\skt{कटुर्व्विपाके}}{\Afootnote{A1938 H; \skt{कटुर्विपाके}\ N}}\skt{ }\edtext{\skt{मधुरश्च}}{\Afootnote{H; \skt{मधुरस्तु}\ A1938 N}}\skt{ }\edtext{\skt{सिम्वः}}{\Afootnote{H; \skt{शिम्बः}\ A1938 N}}\skt{ }\edtext{\skt{प्रभिन्नविण्मारुतपित्तलश्च}}{\Afootnote{A1938 H; \skt{प्रभिन्नविट्मारुतपित्तलश्च ।}\ N}}\skt{ । }\pend
\pstart\versepar \skt{}\edtext{\skt{शिताशिता}}{\Afootnote{H; \skt{सितासिताः}\ A1938; \skt{सितासिता}\ N}}\skt{ }\edtext{\skt{पीतकिनः}}{\Afootnote{H; \skt{पीतकर}\ A1938; \skt{पीतः}\ N}}\skt{ }\edtext{\skt{कुवर्ण्णा}}{\Afootnote{N H; \skt{क्तवर्णा}\ A1938}}\skt{ भवन्ति }\edtext{\skt{येऽन्येकिसराश्च}}{\lemma{\skt{ऽन्येकिसराश्च}}\Afootnote{H; \skt{ऽनेकविधास्तु}\ A1938; \skt{न्येकिशराश्च}\ N}}\skt{ }\edtext{\skt{सिम्वाः}}{\Afootnote{H; \skt{शिम्बाः ॥४५॥}\ A1938; \skt{शिम्वाः ॥}\ N}}\skt{ ॥}\pend

\pstart\versepar \ledsidenote{\passref{1.46.46}}\edtext{\skt{यथोदितस्ते}}{\Afootnote{H; \skt{यथादितस्ते}\ A1938; \skt{यथादितास्ते}\ N}}\skt{ गुणतः प्रधाना }\edtext{\skt{ज्ञेयास्तथार्द्द्रा}}{\lemma{\skt{तथार्द्द्रा}}\Afootnote{H; \skt{कटूष्णा}\ A1938; \skt{तथाद्रा}\ N}}\skt{ रसपाकयोश्च । }\pend
\pstart\versepar \skt{सहाद्वयं }\edtext{\skt{मूलकपादिका च}}{\Afootnote{N H; \skt{मूलकजाश्च शिम्बाः}\ A1938}}\skt{ }\edtext{\skt{कुसिद्धवल्लीप्रभवाश्च}}{\Afootnote{N H; \skt{कुशिम्बिवल्लीप्रभवास्तु}\ A1938}}\skt{ }\edtext{\skt{शिम्बाः}}{\Afootnote{A1938 H; \skt{शिम्वाः ॥}\ N}}\skt{ ॥}\pend

\pstart\versepar \ledsidenote{\passref{1.46.47}}\skt{ज्ञेया विपाके मधुरा रसे च }\edtext{\skt{वलप्रदाः}}{\Afootnote{N H; \skt{बलप्रदाः}\ A1938}}\skt{ }\edtext{\skt{पित्तनिबर्हणाश्च}}{\Afootnote{A1938 H; \skt{पित्तनिवर्हणाश्च ।}\ N}}\skt{ । }\pend
\pstart\versepar \skt{विदाहवन्तश्च }\edtext{\skt{भृशञ्च}}{\Afootnote{H; \skt{भृशं वि}\ A1938; \skt{भृषञ्च}\ N}}\skt{ }\edtext{\skt{रूक्षा}}{\Afootnote{A1938 H; \skt{रुक्षा}\ N}}\skt{ }\edtext{\skt{विष्टभ्य}}{\Afootnote{A1938 H; \skt{विष्टंभ्य}\ N}}\skt{ }\edtext{\edtext{\skt{जीर्यन्त्यनिलप्रदाश्च}}{\lemma{\skt{जीर्यन्त्य्}}\Afootnote{N H; \skt{जीर्यन्}\ A1938}}}{\lemma{\skt{अनिलप्रदाश्}}\Afootnote{N H; \skt{त्यनिलप्रदाश्}\ A1938}}\skt{ ॥}\pend

\pstart\versepar \ledsidenote{\passref{1.46.48}}\skt{रुचिप्रदाश्चैव सुदुर्ज्जराश्च }\edtext{\skt{सर्व्वाः}}{\Afootnote{N H; \skt{सर्वे}\ A1938}}\skt{ स्मृता }\edtext{\skt{वैदलिकाश्च}}{\Afootnote{N H; \skt{वैदलिकास्तु}\ A1938}}\skt{ }\edtext{\skt{शिम्बाः}}{\Afootnote{A1938 H; \skt{शिम्वाः ।}\ N}}\skt{ । }\pend
\pstart\versepar \skt{}\edtext{\edtext{\skt{कटुर्व्विपाके}}{\lemma{\skt{कटुर्}}\Afootnote{N H; \skt{कटु}\ A1938}}}{\lemma{\skt{व्विपाके}}\Afootnote{N H; \skt{र्विपाके}\ A1938}}\skt{ }\edtext{\skt{मधुरानिलघ्नो}}{\Afootnote{H; \skt{कटुकः कफघ्नो}\ A1938; \skt{मधुरोऽनिलघ्न}\ N}}\skt{ }\edtext{\skt{विदाहिभावादहितः}}{\Afootnote{A1938 H; \skt{विदाहिभावादहितः}\ N}}\skt{ }\edtext{\skt{कुसुम्बः}}{\Afootnote{H; \skt{कुसुम्भः}\ A1938 N}}\skt{ ॥}\pend

\pstart\versepar \ledsidenote{\passref{1.46.49}}\skt{उष्णातसी स्वादुरसानिलघ्नी }\edtext{\skt{पित्तोल्वणा}}{\Afootnote{N H; \skt{पित्तोल्बणा}\ A1938}}\skt{ स्यात्कटुका विपाके । }\pend
\pstart\versepar \skt{पाके रसे चापि }\edtext{\skt{कटुप्रदिष्ट स्सिद्धार्थकः}}{\Afootnote{N H; \skt{कटुः प्रदिष्टः सिद्धार्थकः}\ A1938}}\skt{ शोणितपित्तकोपी ॥ }\pend
\pstart\versepar \skt{}\edtext{\skt{स्निग्धोष्णतिक्तः}}{\Afootnote{H; \skt{तीक्ष्णोष्णरूक्षः}\ A1938; \skt{स्निग्धोष्णतीक्ष्णः}\ N}}\skt{ }\edtext{\skt{कफवातहन्ता}}{\Afootnote{N H; \skt{कफमारुतघ्नस्}\ A1938}}\skt{ तथा }\edtext{\skt{गुणश्चाशितसर्षपो}}{\lemma{\skt{चाशितसर्षपो}}\Afootnote{H; \skt{चासितसर्षपो}\ A1938; \skt{चासितसर्षपो}\ N}}\skt{ }\edtext{\skt{पि}}{\Afootnote{A1938 H; \skt{हि ॥}\ N}}\skt{ ॥}\pend

\pstart\versepar \ledsidenote{\passref{1.46.50}}\skt{अनार्त्तवं }\edtext{\edtext{\skt{व्याधिहतमपर्यागतमेव}}{\lemma{\skt{व्याधिहतम्}}\Afootnote{N H; \skt{व्याधिहत}\ A1938}}}{\lemma{\skt{अपर्यागतम्}}\Afootnote{H; \skt{मपर्यागतम्}\ A1938; \skt{अपर्यातम्}\ N}}\skt{ च । }\pend
\pstart\versepar \skt{}\edtext{\skt{अभूमिजन्नवम्वापि न}}{\lemma{\skt{वापि न}}\Afootnote{H; \skt{चापि न}\ A1938 N}}\skt{ धान्यङ्गुणवत्स्मृतं ॥}\pend

\pstart\versepar \ledsidenote{\passref{1.46.51}}\edtext{\edtext{\skt{नवन्धान्यमभिष्यन्दि}}{\lemma{\skt{धान्यम्}}\Afootnote{N H; \skt{धान्य}\ A1938}}}{\lemma{\skt{अभिष्यन्दि}}\Afootnote{N H; \skt{मभिष्यन्दि}\ A1938}}\skt{ लघु }\edtext{\skt{सम्वत्सरोषितं}}{\Afootnote{A1938 H; \skt{संवत्सरोसितं ।}\ N}}\skt{ । }\pend
\pstart\versepar \skt{विदाहि गुरु विष्टम्भि }\edtext{\skt{विरूढम्वातकोपनं}}{\lemma{\skt{वातकोपनं}}\Afootnote{N H; \skt{दृष्टिदूषणम् ॥५१॥}\ A1938}}\skt{ ॥}\pend

\pstart\versepar \ledsidenote{\passref{1.46.52}}\edtext{\skt{शाल्यादेः}}{\Afootnote{A1938 H; \skt{शाल्यादः}\ N}}\skt{ सर्षपान्तस्य }\edtext{\skt{द्विविधस्यास्य}}{\Afootnote{N H; \skt{विविधस्यास्य}\ A1938}}\skt{ }\edtext{\skt{भागशः}}{\Afootnote{A1938 H; \skt{भागतः ।}\ N}}\skt{ । }\pend
\pstart\versepar \skt{}\edtext{\skt{कालप्रमाणसंस्कारो}}{\Afootnote{H; \skt{कालप्रमाणसंस्कार}\ A1938; \skt{कालप्रमाणसस्कारो}\ N}}\skt{ }\edtext{\skt{मात्रा}}{\Afootnote{N H; \skt{मात्राः}\ A1938}}\skt{ }\edtext{\skt{चास्मिन्परीक्ष्यते}}{\Afootnote{H; \skt{सम्परिकीर्तिताः ॥५२॥ इति कुधान्यवर्गः ।}\ A1938; \skt{चास्मिम्परीक्ष्यते ॥ {\symfont ❈} ॥}\ N}}\skt{ ॥}\pend

\pstart\versepar \ledsidenote{\passref{1.46.53}}\edtext{\skt{अत}}{\Afootnote{N H; \skt{अथ मांसवर्गः । अत}\ A1938}}\skt{ }\edtext{\edtext{\edtext{\skt{ऊर्द्ध्वम्मांसवर्ग्गमुपदेक्ष्यामः}}{\lemma{\skt{ऊर्द्ध्वम्}}\Afootnote{A1938 H; \skt{ऊर्द्ध्व}\ N}}}{\lemma{\skt{मांसवर्ग्गम्}}\Afootnote{N H; \skt{मांसवर्गा}\ A1938}}}{\lemma{\skt{उपदेक्ष्यामः}}\Afootnote{N H; \skt{नुपदेक्ष्यामः;}\ A1938}}\skt{ ॥ }\pend
\pstart\versepar \skt{}\edtext{\skt{तत्र}}{\Afootnote{H; \skt{तद्यथा-}\ A1938; \skt{दत्र}\ N}}\skt{ जलेशया }\edtext{\skt{अनूपा}}{\Afootnote{H; \skt{आनूपा}\ A1938 N}}\skt{ ग्राम्याः । क्रव्यभुज । }\edtext{\edtext{\skt{एकसफजाङ्गलाश्चेति}}{\lemma{\skt{एकसफजाङ्गलाश्}}\Afootnote{H; \skt{एकशफा, जाङ्गला}\ A1938; \skt{एकशफा जङ्गलाश्}\ N}}}{\lemma{\skt{चेति}}\Afootnote{N H; \skt{श्चेति}\ A1938}}\skt{ । }\edtext{\skt{षण्मांसवर्ग्गा}}{\Afootnote{A1938 H; \skt{षड्मांसवर्गा}\ N}}\skt{ भवन्ति । }\edtext{\edtext{\skt{तेषामुत्तरोत्तराः}}{\lemma{\skt{तेषाम्}}\Afootnote{N H; \skt{एतेषां वर्गाणा}\ A1938}}}{\lemma{\skt{उत्तरोत्तराः}}\Afootnote{N H; \skt{मुत्तरोत्तरा}\ A1938}}\skt{ प्रधानतमास्ते }\edtext{\edtext{\skt{पुनर्द्विविधा}}{\lemma{\skt{पुनर्}}\Afootnote{N H; \skt{पुन}\ A1938}}}{\lemma{\skt{द्विविधा}}\Afootnote{N H; \skt{र्द्विविधा}\ A1938}}\skt{ }\edtext{\skt{आनूपजाङ्गलाश्च}}{\Afootnote{H; \skt{जाङ्गला आनूपाश्चेति ।}\ A1938; \skt{आनूपा जाङ्गलश्च ।}\ N}}\skt{ । तत्र }\edtext{\edtext{\skt{जाङ्गलवर्ग्गोऽष्टविधः}}{\lemma{\skt{जाङ्गलवर्ग्गो}}\Afootnote{A1938 H; \skt{जा ङ्गलवर्गो}\ N}}}{\lemma{\skt{ऽष्टविधः}}\Afootnote{A1938 H; \skt{ष्टविधः ।}\ N}}\skt{ । }\edtext{\skt{तत्र}}{\Afootnote{N H; \skt{तद्यथा-}\ A1938}}\skt{ }\edtext{\skt{जंघाला}}{\Afootnote{A1938 H; \skt{जाङ्गला}\ N}}\skt{ विष्किराः }\edtext{\skt{प्रनुदा}}{\Afootnote{H; \skt{प्रतुदा,}\ A1938; \skt{प्रत्युदा}\ N}}\skt{ }\edtext{\skt{गुहेशयाः}}{\Afootnote{H; \skt{गुहाशयाः,}\ A1938; \skt{गुहशयाः}\ N}}\skt{ प्रसहाः पर्ण्णमृगाः । }\edtext{\skt{विलेशया}}{\Afootnote{N H; \skt{बिलेशया,}\ A1938}}\skt{ }\edtext{\edtext{\skt{ग्राम्याश्चेति}}{\lemma{\skt{ग्राम्याश्}}\Afootnote{N H; \skt{ग्राम्या}\ A1938}}}{\lemma{\skt{चेति}}\Afootnote{N H; \skt{श्चेति ।}\ A1938}}\skt{ ॥ }\edtext{\edtext{\skt{तेषञ्जङ्घालविष्किरौ}}{\lemma{\skt{तेषञ्}}\Afootnote{H; \skt{तेषां}\ A1938; \skt{ए षां}\ N}}}{\lemma{\skt{जङ्घालविष्किरौ}}\Afootnote{A1938 H; \skt{जागलविष्किरौ}\ N}}\skt{ प्रधानतमौ ।}\pend

\pstart \ledsidenote{\passref{1.46.54}}\edtext{\skt{तत्र}}{\Afootnote{N H; \emph{om.}\ A1938}}\skt{ }\edtext{\skt{जङ्घला}}{\Afootnote{H; \skt{ता}\ A1938; \skt{जङ्गला}\ N}}\skt{ }\edtext{\skt{एणहरिणर्कुरङ्ग}}{\Afootnote{H; \skt{वेणहरिणर्क्षकुरङ्ग}\ A1938; \skt{एणहरिणकुरङ्ग}\ N}}\skt{ }\edtext{\skt{ऋष्य}}{\Afootnote{H; \skt{क}\ A1938; \skt{दष्य}\ N}}\skt{ }\edtext{\skt{रालकृतमालशलभश्वदंष्ट्रीचारुष्कवृषतेमृगमातृकाप्रभृतयः}}{\Afootnote{H; \skt{रालकृतमालशरभश्वदंष्ट्रापृषतचारुष्करमृगमातृकाप्रभृतयो}\ A1938; \skt{शलकृतमालशरभश्वदंष्टीचारुष्कवृषत मृगामातृकाप्रभृतयः}\ N}}\skt{ }\edtext{\skt{कषाय}}{\Afootnote{H; \skt{जङ्घाला मृगाः कषाया}\ A1938; \skt{कषाया}\ N}}\skt{ मधुरा लघवो }\edtext{\edtext{\skt{वातपित्तहरास्तीक्ष्णा}}{\lemma{\skt{वातपित्तहरास्}}\Afootnote{N H; \skt{वातपित्तहरा}\ A1938}}}{\lemma{\skt{तीक्ष्णा}}\Afootnote{H; \skt{स्तीक्ष्णा}\ A1938; \skt{तीक्ष्ण्या}\ N}}\skt{ }\edtext{\skt{बस्तिविशोधनाः}}{\Afootnote{\skt{हृद्या बस्तिशोधनाश्च ॥५४॥}\ A1938; \skt{वस्तिशोधनः ॥}\ N; \skt{वस्तिविशोधनाः ॥}\ H}}\skt{ ॥}\pend

\pstart\versepar \ledsidenote{\passref{1.46.55}}\edtext{\skt{कषायमधुरा}}{\Afootnote{H; \skt{कषायमधुरो}\ A1938; \skt{कषायो मधुरो}\ N}}\skt{ हृद्यः }\edtext{\skt{पित्तासृक्कफरोगहा}}{\Afootnote{A1938 H; \skt{पित्तसृक्कफरोगहा ।}\ N}}\skt{ । }\pend
\pstart\versepar \skt{}\edtext{\skt{साङ्ग्राही}}{\Afootnote{N H; \skt{सङ्ग्राही}\ A1938}}\skt{ रोचनो }\edtext{\skt{वल्यस्तेषामेणो}}{\lemma{\skt{वल्यस्}}\Afootnote{N H; \skt{बल्यस्}\ A1938}}\skt{ ज्वरापहः ॥}\pend

\pstart\versepar \ledsidenote{\passref{1.46.56}}\skt{मधुरो मधुरः पाके दोषघ्नो }\edtext{\skt{लघुदीपनः}}{\Afootnote{H; \skt{ऽनलदीपनः ।}\ A1938; \skt{लघुदपिनः ।.}\ N}}\skt{ । }\pend
\pstart\versepar \skt{शीतलो }\edtext{\skt{वद्धविट्मूत्रः}}{\Afootnote{H; \skt{बद्धविण्मूत्र}\ A1938; \skt{वद्धविड्मूत्रः}\ N}}\skt{ }\edtext{\edtext{\skt{सुगन्धिर्हरिणः}}{\lemma{\skt{सुगन्धिर्}}\Afootnote{N H; \skt{सुगन्धि}\ A1938}}}{\lemma{\skt{हरिणः}}\Afootnote{H; \skt{र्हरिणो}\ A1938; \skt{हरिणाः}\ N}}\skt{ }\edtext{\skt{स्मृतः}}{\Afootnote{N H; \skt{लघुः ॥५६॥}\ A1938}}\skt{ ॥}\pend

\pstart \ledsidenote{\passref{1.46.57ab}}\skt{एणः }\edtext{\skt{कृष्णस्तयोर्ज्ञेयो}}{\lemma{\skt{कृष्णस्}}\Afootnote{A1938 H; \skt{कृष्ण}\ N}}\skt{ हरिणस्ताम्र उच्यते ॥}\pend

\pstart \ledsidenote{\passref{1.46.59}}\edtext{\skt{लावतित्तिरिकपिञ्जलवर्त्तीरकवर्त्तकवातीकाचकोरकलविङ्कमयूरक्रकरोपचक्रककुक्कुटा}}{\Afootnote{H; \skt{लावतित्तिरिकपिञ्जलवर्तीरवर्तिकावर्तकनप्तृकावार्तीकचकोरकलविङ्कमयूरक्रकरोपचक्र- कुक्कुट}\ A1938; \skt{लावत्तित्तिरिकपिञ्जलवर्त्तीरकवर्त्तकवातीकचकोरकलविड्कमयूरक्रकरोपचक्र तकुर्क्कुट}\ N}}\skt{ }\edtext{\skt{विष्किरा}}{\Afootnote{H; \skt{सारङ्ग}\ A1938; \emph{om.}\ N}}\skt{ }\edtext{\skt{लघवः}}{\Afootnote{H; \skt{शतपत्रकुतित्तिरिकुरुवाहकयवालकप्रभृतयस्त्र्याहला विष्किराः ॥५९॥}\ A1938; \skt{विष्किरा}\ N}}\skt{ ।}\pend

\pstart\versepar \ledsidenote{\passref{1.46.60}}\edtext{\skt{शीतलामधुराः}}{\Afootnote{H; \skt{लघवः शीतमधुराः}\ A1938; \skt{लघवः । शीतला मधुराः}\ N}}\skt{ कषाया }\edtext{\skt{दोषशमनाश्च}}{\Afootnote{N H; \skt{दोषनाशनाः}\ A1938}}\skt{ ॥ }\pend
\pstart\versepar \skt{}\edtext{\skt{सांग्राही}}{\Afootnote{N H; \skt{सङ्ग्राही}\ A1938}}\skt{ दीपनः }\edtext{\skt{शीतः}}{\Afootnote{N H; \skt{चैव}\ A1938}}\skt{ }\edtext{\skt{कषायमधुरस्तथा}}{\lemma{\skt{तथा}}\Afootnote{N H; \skt{लघुः ।}\ A1938}}\skt{ । }\pend
\pstart\versepar \skt{लावः कटुविपाकश्च }\edtext{\skt{सन्निपातेषु}}{\Afootnote{H; \skt{सन्निपाते तु}\ A1938; \skt{सन्निपाते च}\ N}}\skt{ पूजितः ॥}\pend

\pstart\versepar \ledsidenote{\passref{1.46.61}}\edtext{\skt{ईषद्गुरूष्णमधुरो}}{\Afootnote{A1938 H; \skt{इषद्गुरूष्णमधुरो}\ N}}\skt{ वृष्यो मेधाग्निवर्द्धनः । }\pend
\pstart\versepar \skt{तित्तिरिः सर्व्वदोषघ्नो ग्राही }\edtext{\skt{वर्ण्णप्रसादनः}}{\Afootnote{A1938 H; \skt{वर्ण्णप्रसाधनः ॥}\ N}}\skt{ ॥}\pend

\pstart\versepar \ledsidenote{\passref{1.46.62}}\skt{रक्तपित्तहरः शीतो लघुश्चापि कपिञ्जलः । }\pend
\pstart\versepar \skt{कफोत्थेषु च रोगेषु मन्दवाते च शस्यते ॥}\pend

\pstart \ledsidenote{\passref{1.46.63}}\edtext{\skt{वातपित्तहरा}}{\Afootnote{N H; \skt{हिक्काश्वासानिलहरो विशेषाद्गौरतित्तिरिः वातपित्तहरा}\ A1938}}\skt{ वृष्या }\edtext{\skt{मेधाग्निबलवर्द्धनाः}}{\Afootnote{A1938 H; \skt{मेधाग्निवलवर्धनाः ।}\ N}}\skt{ ।}\pend

\pstart\versepar \ledsidenote{\passref{1.46.64}}\skt{लघवः क्रकरा }\edtext{\skt{हृद्यास्तथाश्चैवोपचक्रकाः}}{\lemma{\skt{तथाश्}}\Afootnote{H; \skt{तथा}\ A1938 N}}\skt{ ॥ }\pend
\pstart\versepar \skt{}\edtext{\skt{कषायः}}{\Afootnote{A1938 H; \skt{कषाय}\ N}}\skt{ }\edtext{\skt{स्वादुलवणः}}{\Afootnote{N H; \skt{स्वादुलवण}\ A1938}}\skt{ }\edtext{\skt{स्वर्यः}}{\Afootnote{H; \skt{स्त्वच्यः}\ A1938; \skt{त्वर्यः}\ N}}\skt{ }\edtext{\skt{केश्यो}}{\Afootnote{A1938 H; \skt{केश्यो}\ N}}\skt{ }\edtext{\skt{रुचिप्रदः}}{\Afootnote{N H; \skt{ऽरुचौ हितः ॥६४॥}\ A1938}}\skt{ ।}\pend

\pstart\versepar \ledsidenote{\passref{1.46.65}}\edtext{\skt{मयूरः}}{\Afootnote{A1938 H; \skt{मयूर}\ N}}\skt{ }\edtext{\skt{स्वरमेधाग्निर्द्दृक्छ्रोत्रेन्द्रियदार्ढ्यकृत्}}{\lemma{\skt{स्वरमेधाग्निर्}}\Afootnote{H; \skt{स्वरमेधाग्नि}\ A1938 N}}\skt{ ॥ }\pend
\pstart\versepar \skt{}\edtext{\skt{स्निग्धोष्णोऽनिलहा}}{\lemma{\skt{ऽनिलहा}}\Afootnote{A1938 H; \skt{निलहा}\ N}}\skt{ वृष्यः }\edtext{\skt{स्वेदः}}{\Afootnote{N H; \skt{स्वेद}\ A1938}}\skt{ }\edtext{\skt{स्वरबलावहः}}{\Afootnote{A1938 H; \skt{स्वरवलावहः ।}\ N}}\skt{ ।}\edtext{}{\lemma{\skt{द्दृक्छ्रोत्रेन्द्रियदार्ढ्यकृत् ॥ स्निग्धोष्णो}}\Afootnote{A1938 H; \skt{दृक्च्छ्रोत्रन्द्रियदार्ढ्यकृत् ।. स्निग्धोष्णो}\ N}}\pend

\pstart \ledsidenote{\passref{1.46.66ab}}\skt{बृंहणः कुक्कुटो }\edtext{\edtext{\skt{वन्यस्तद्वद्ग्राम्यो}}{\lemma{\skt{वन्यस्}}\Afootnote{H; \skt{वन्य}\ A1938}}}{\lemma{\skt{तद्वद्ग्राम्यो}}\Afootnote{H; \skt{स्तद्वद्ग्राम्यो}\ A1938}}\skt{ गुरुस्तु सः ॥}\pend

\pstart \ledsidenote{\passref{1.46.67}}\edtext{\skt{कपोतपारावतभृङ्गराजपरभृतयष्टीमधुकुलिङ्गगोक्ष्वेलकडिण्डिमाणशतपत्रमातृलिङ्गभेदाशिशुकसारिकावग्गुलीलट्वालदूषकसूगृहाखञ्जरीटकदात्यूहप्रभृतयः}}{\Afootnote{\skt{कपोतपारावतभृङ्गराजपरभृतकोयष्टिककुलिङ्गगृहकुलिङ्गगोक्ष्वेडकडिण्डिमाणवकशतपत्रक- मातृनिन्दकभेदाशिशुकसारिकावल्गुलीगिरिशालट्वालट्टूषकसुगृहाखञ्जरीटहारीतदात्यूहप्रभृतयः ,}\ A1938; \skt{कपोतपारावतभृङ्गराजपरभृतयष्टीमधुकुलिंगोक्ष्वेलडिण्डिमाणश तपत्रमातृलिङ्गभेदाशीशुकसारिकावङ्गुलीलट्वालदृषकसूगृहाखंजरीटकदात्यूह प्रभृतयः}\ N; \skt{कपोतपारावतभृङ्गराजपरभृतयष्टीमककुलिङ्गगोक्ष्वेलकडिण्डिमाणशतपत्रमातृलिङ्गभेदाशिशुकसारिकावग्गुलीलट्वालदूषकसूगृहाखञ्जरीटकदात्यूहप्रभृतयः}\ H}}\skt{ }\edtext{\skt{प्रनुदाः}}{\Afootnote{H; \skt{प्रतुदाः}\ A1938 N}}\skt{ ॥}\pend

\pstart\versepar \ledsidenote{\passref{1.46.68}}\skt{कषायमधुरा }\edtext{\skt{रूक्षाः}}{\Afootnote{A1938 H; \skt{रूक्षा}\ N}}\skt{ }\edtext{\skt{फलाहारानिलावहाः}}{\Afootnote{N H; \skt{फलाहारा मरुत्कराः ।}\ A1938}}\skt{ । }\pend
\pstart\versepar \skt{}\edtext{\skt{श्लेष्मपित्तहराः}}{\Afootnote{N H; \skt{पित्तश्लेष्महराः}\ A1938}}\skt{ शीता }\edtext{\skt{बद्धमूत्राल्पवर्च्चसः}}{\Afootnote{A1938 H; \skt{वद्धमूत्राल्पवर्चसः}\ N}}\skt{ ॥}\pend

\pstart\versepar \ledsidenote{\passref{1.46.69}}\edtext{\skt{सर्व्वदोषकरस्तेषाम्भेदासी}}{\lemma{\skt{भेदासी}}\Afootnote{H; \skt{भेदाशी}\ A1938 N}}\skt{ }\edtext{\skt{मलदूषणः}}{\Afootnote{N H; \skt{मलदूषकः}\ A1938}}\skt{ । }\pend
\pstart\versepar \skt{कषायस्वादुलवणो }\edtext{\skt{गुरुः}}{\Afootnote{A1938 H; \skt{गुरु}\ N}}\skt{ काणकपोतकः ॥}\pend

\pstart\versepar \ledsidenote{\passref{1.46.70}}\skt{रक्तपित्तप्रशमनः }\edtext{\skt{कषायविषदो}}{\Afootnote{N H; \skt{कषायविशदो}\ A1938}}\skt{ पि च । }\pend
\pstart\versepar \skt{}\edtext{\skt{रसतो}}{\Afootnote{N H; \skt{विपाके}\ A1938}}\skt{ मधुरश्चापि गुरुः पारावतः }\edtext{\skt{स्मृताः}}{\Afootnote{H; \skt{स्मृतः}\ A1938 N}}\skt{ ॥}\pend

\pstart\versepar \ledsidenote{\passref{1.46.71}}\skt{कुलिङ्गो }\edtext{\skt{मधुर}}{\Afootnote{N H; \skt{मधुरः}\ A1938}}\skt{ स्निग्धः कफशुक्रविवर्द्धनः । }\pend
\pstart\versepar \skt{रक्तपित्तहरो }\edtext{\edtext{\skt{वेश्मकुलिङ्गस्त्वतिशुक्रलः}}{\lemma{\skt{वेश्मकुलिङ्गस्}}\Afootnote{H; \skt{वेश्मकुलिङ्ग}\ A1938; \skt{वेष्मकुलिङ्ग}\ N}}}{\lemma{\skt{त्वतिशुक्रलः}}\Afootnote{H; \skt{स्त्वतिशुक्रलः}\ A1938 N}}\skt{ ॥}\pend

\pstart \ledsidenote{\passref{1.46.72}}\edtext{\skt{सिंहव्याघ्रतरक्ष्वृक्षकद्वीपिमार्ज्जारसृगालमृगेर्व्वारुकप्रभृतयो}}{\Afootnote{H; \skt{सिंहव्याघ्रवृकतरक्ष्वृक्षद्वीपिमार्जारशृगालमृगेर्वारुकप्रभृतयो}\ A1938; \skt{सिंहव्याघ्रतरक्वृक्षकद्वीपिमार्जारशृगालमृगेर्वारुकप्रभृतयो}\ N}}\skt{ }\edtext{\skt{गुहेशयाः}}{\Afootnote{H; \skt{गुहाशयाः ॥७२॥}\ A1938; \skt{गुहशयाः ॥}\ N}}\skt{ ॥}\pend

\pstart\versepar \ledsidenote{\passref{1.46.73}}\skt{मधुरा }\edtext{\skt{गुरवः}}{\Afootnote{A1938 H; \skt{गुरुव}\ N}}\skt{ स्निग्धा }\edtext{\skt{बल्या}}{\Afootnote{A1938 H; \skt{वल्या}\ N}}\skt{ मारुतनाशनाः । }\pend
\pstart\versepar \skt{उष्णवीर्या }\edtext{\skt{हिते}}{\Afootnote{H; \skt{हिता}\ A1938; \skt{हित}\ N}}\skt{ नित्यं }\edtext{\skt{नेत्रगुह्ये च}}{\Afootnote{N H; \skt{नेत्रगुह्यवि}\ A1938}}\skt{ }\edtext{\skt{रोगिणां}}{\Afootnote{N H; \skt{कारिणाम् ॥७३॥}\ A1938}}\skt{ ॥}\pend

\pstart \ledsidenote{\passref{1.46.74}}\edtext{\skt{काककङ्ककुररभाससमघात्युल्लूकचीरिल्लश्येनप्रभृतयः}}{\Afootnote{H; \skt{काककङ्ककुररचाषभासशशघात्युलूकचिल्लिश्येनगृध्रप्रभृतयः}\ A1938; \skt{काककङ्ककुररभासशसघात्यूल्लूश्येनप्रभृतयः}\ N}}\skt{ प्रसहाः ॥}\pend

\pstart\versepar \ledsidenote{\passref{1.46.75}}\skt{एते सिंहादिभिः सर्व्वे समाना वायसादयः । }\pend
\pstart\versepar \skt{}\edtext{\skt{रसे}}{\Afootnote{H; \skt{रस}\ A1938 N}}\skt{ }\edtext{\skt{पाके च}}{\Afootnote{N H; \skt{वीर्यविपाके}\ A1938}}\skt{ }\edtext{\skt{वीर्ये च}}{\Afootnote{H; \skt{षु}\ A1938; \skt{वीर्यं च}\ N}}\skt{ }\edtext{\skt{विशेषाच्छोषिणो}}{\Afootnote{H; \skt{विशेषाच्छोषिणे}\ A1938 N}}\skt{ हिताः ॥}\pend

\pstart \ledsidenote{\passref{1.46.76}}\edtext{\skt{मद्गुमूषिकवृक्षशायिकवक्कसपूतीघसवानरप्रभृतयः}}{\Afootnote{H; \skt{मद्गुमूषिकवृक्षशायिकावकुशपूतिघासवानरप्रभृतयः}\ A1938; \skt{मद्गुमूषिकवृक्षशायिकवकुशपूतीघसवानरप्रभृतयः}\ N}}\skt{ पर्ण्णमृगाः ॥}\pend

\pstart\versepar \ledsidenote{\passref{1.46.77}}\skt{मधुरा गुरवो वृष्याश्चक्षुष्याः }\edtext{\skt{शोषिणां}}{\Afootnote{N H; \skt{शोषिणे}\ A1938}}\skt{ हिताः । }\pend
\pstart\versepar \skt{सृष्टमूत्रपुरीषाश्च }\edtext{\edtext{\skt{कासश्वासार्शसान्तथा}}{\lemma{\skt{कासश्वासार्शसान्}}\Afootnote{N H; \skt{कासार्शःश्वा}\ A1938}}}{\lemma{\skt{तथा}}\Afootnote{N H; \skt{सनाशनाः ॥७७॥}\ A1938}}\skt{ ॥}\pend

\pstart \ledsidenote{\passref{1.46.78}}\edtext{\skt{श्वाविच्छल्यकगोधाशशवृषदंशलोपाकलोमशकर्ण्णकदलीमृगप्रियकाजगरसर्प्पप्रभृतयो}}{\Afootnote{H; \skt{श्वाविच्छल्यकगोधाशशवृषदंशलोपाकलोमशकर्णकदलीमृगप्रियकाजगरसर्पमूषिकनकुलमहाबभ्रुप्रभृतयो}\ A1938; \skt{श्वाविट्छल्यकगोधाशशवृषदंशलोपाकलोमशकर्ण्णकदलीमृगप्रियकाजगरसर्पप्रभृतयो}\ N}}\skt{ }\edtext{\skt{विलेशयाः}}{\Afootnote{N H; \skt{बिलेशयाः ॥७८॥}\ A1938}}\skt{ ॥}\pend

\pstart\versepar \ledsidenote{\passref{1.46.79}}\edtext{\edtext{\skt{साङ्ग्राहिकावद्धविट्मूत्रास्तथैते}}{\lemma{\skt{साङ्ग्राहिकावद्धविट्मूत्रास्}}\Afootnote{N H; \skt{वर्चोमूत्रं सं}\ A1938}}}{\lemma{\skt{तथैते}}\Afootnote{N H; \skt{हतं कुर्युरेते}\ A1938}}\skt{ }\edtext{\skt{वीर्योष्णाः}}{\Afootnote{N H; \skt{वीर्ये चोष्णाः}\ A1938}}\skt{ }\edtext{\skt{पूर्व्ववत्स्वादुकाः}}{\Afootnote{H; \skt{पूर्ववत्स्वादुपा}\ A1938; \skt{पूर्व्ववत्स्वदुकाः}\ N}}\skt{ }\edtext{\skt{स्मृताः}}{\Afootnote{N H; \skt{काः ।}\ A1938}}\skt{ । }\pend
\pstart\versepar \skt{वातं हन्युः श्लेष्मपित्ते च कुर्युः }\edtext{\skt{स्न्गिधाः}}{\Afootnote{H; \skt{स्निग्धाः}\ A1938; \skt{स्निग्धा}\ N}}\skt{ कासश्वासकार्श्यापहाश्च ॥}\pend

\pstart\versepar \ledsidenote{\passref{1.46.80}}\skt{कषायमधुरस्तेषां }\edtext{\skt{शशः}}{\Afootnote{A1938 H; \skt{शश}\ N}}\skt{ पित्तकफापहः । }\pend
\pstart\versepar \skt{}\edtext{\skt{नातिशीतलवीर्यत्वाद्वातसाधारणो}}{\Afootnote{A1938 H; \skt{नातिशीतगुरुस्निग्धो मन्दपित्तकफः}\ N}}\skt{ }\edtext{\skt{मतः}}{\Afootnote{A1938 H; \skt{स्मृतः ।}\ N}}\skt{ ॥}\pend

\pstart\versepar \ledsidenote{\passref{1.46.83}}\skt{दुर्न्नामानिलदोषघ्नाः }\edtext{\skt{क्रिमिदूषीविषापहः}}{\Afootnote{H; \skt{कृमिदूषीविषापहाः ।}\ A1938}}\skt{ । }\pend
\pstart\versepar \skt{चक्षुष्या }\edtext{\skt{मधुरा}}{\Afootnote{H; \skt{मधुराः}\ A1938}}\skt{ }\edtext{\skt{र्पाके}}{\Afootnote{H; \skt{पाके}\ A1938}}\skt{ सर्प्पा }\edtext{\skt{मेधाकराः}}{\Afootnote{H; \skt{मेधाग्निवर्धनाः ॥८३॥}\ A1938}}\skt{ }\edtext{\skt{स्मृताः}}{\Afootnote{H; \emph{om.}\ A1938}}\skt{ ॥}\pend

\pstart\versepar \ledsidenote{\passref{1.46.84}}\skt{दर्व्वीकरा }\edtext{\skt{दीपकश्च}}{\Afootnote{H; \skt{दीपकाश्च}\ A1938}}\skt{ तेषूक्ताः कटुपाकिनः । }\pend
\pstart\versepar \skt{}\edtext{\skt{मधुरात्यर्थ}}{\Afootnote{H; \skt{मधुराश्चाति}\ A1938}}\skt{ चक्षुष्याः }\edtext{\skt{सर्प्पा}}{\Afootnote{H; \skt{सृष्ट}\ A1938}}\skt{ }\edtext{\skt{मेधाकराः}}{\Afootnote{H; \skt{विण्मूत्रमा}\ A1938}}\skt{ }\edtext{\skt{स्मृताः}}{\Afootnote{H; \skt{रुताः ॥८४॥}\ A1938}}\skt{ ॥}\pend

\pstart \ledsidenote{\passref{1.46.85}}\edtext{\skt{अश्वाश्वतरगोखरोष्ट्रवस्तोरभ्रमेदःपुच्छकप्रभृतयो}}{\Afootnote{H; \skt{अश्वाश्वतरगोखरोष्ट्रबस्तोरभ्रमेदःपुच्छकप्रभृतयो}\ A1938}}\skt{ ग्राम्याः ॥}\pend

\pstart\versepar \ledsidenote{\passref{1.46.86}}\skt{ग्राम्या वातहराः सर्व्वे }\edtext{\skt{वृंहणाः}}{\Afootnote{H; \skt{बृंहणाः}\ A1938}}\skt{ कफपित्तलाः । }\pend
\pstart\versepar \skt{मधुरा रसपाकाभ्यान्दीपना }\edtext{\skt{वलवर्द्धनाः}}{\Afootnote{H; \skt{बलवर्धनाः ॥८६॥}\ A1938}}\skt{ ॥}\pend

\pstart\versepar \ledsidenote{\passref{1.46.87}}\edtext{\skt{नातिशीत}}{\Afootnote{H; \skt{नातिशीतो}\ A1938}}\skt{ }\edtext{\skt{गुरू}}{\Afootnote{H; \skt{गुरुः}\ A1938}}\skt{ स्निग्धो मन्दपित्तकफः स्मृतः । }\pend
\pstart\versepar \skt{}\edtext{\skt{छगलस्त्वनभिष्यन्दि}}{\Afootnote{H; \skt{छगलस्त्वनभिष्यन्दी}\ A1938; \skt{छगलस्त्वनभिष्यन्दि}\ N}}\skt{ तेषां पीनसनाशनः ॥}\edtext{}{\lemma{\skt{… छ}}\Afootnote{\emph{om.}\ N}}\pend

\pstart\versepar \ledsidenote{\passref{1.46.88}}\edtext{\edtext{\edtext{\skt{वृंहणम्मांसमौरभ्रं}}{\lemma{\skt{वृंहणम्}}\Afootnote{N H; \skt{बृंहणं}\ A1938}}}{\lemma{\skt{मांसम्}}\Afootnote{N; \skt{मांस}\ A1938 H}}}{\lemma{\skt{औरभ्रं}}\Afootnote{N; \skt{मौरभ्रं}\ A1938 H}}\skt{ }\edtext{\edtext{\skt{पित्तश्लेष्माकरङ्गुरुः}}{\lemma{\skt{पित्तश्लेष्माकरङ्}}\Afootnote{N H; \skt{पित्तश्लेष्मावहं}\ A1938}}}{\lemma{\skt{गुरुः}}\Afootnote{N H; \skt{गुरु}\ A1938}}\skt{ । }\pend
\pstart\versepar \skt{मेदः पुच्छोद्भवं }\edtext{\edtext{\skt{वृष्यमौरभ्रसदृशङ्गुणैः}}{\lemma{\skt{वृष्यम्}}\Afootnote{N H; \skt{वृष्य}\ A1938}}}{\lemma{\skt{औरभ्रसदृशङ्}}\Afootnote{H; \skt{मौरभ्रसदृशं}\ A1938; \skt{औरभ्हसदृशं}\ N}}\skt{ ॥}\pend

\pstart\versepar \ledsidenote{\passref{1.46.89}}\edtext{\skt{शोषकासप्रतिश्यायविषमज्वरनाशनं}}{\Afootnote{H; \skt{श्वासकासप्रतिश्यायविषमज्वरनाशनम् ।}\ A1938; \skt{सोषकासप्रतिश्यायविशमज्वरनाशनं ।.}\ N}}\skt{ । }\pend
\pstart\versepar \skt{}\edtext{\skt{गव्यं}}{\Afootnote{H; \emph{om.}\ A1938; \skt{गंव्य}\ N}}\skt{ श्रमात्यग्निहितं }\edtext{\edtext{\skt{पवित्रमनिलापहं}}{\lemma{\skt{पवित्रम्}}\Afootnote{H; \skt{गव्यं पवित्र}\ A1938; \skt{पवित्र}\ N}}}{\lemma{\skt{अनिलापहं}}\Afootnote{H; \skt{मनिलापहम्}\ A1938 N}}\skt{ ॥}\pend

\pstart\versepar \ledsidenote{\passref{1.46.90}}\edtext{\edtext{\skt{औरभ्रवत्सलवणम्मांसमेकसफोद्भवं}}{\lemma{\skt{मांसम्}}\Afootnote{N H; \skt{मांस}\ A1938}}}{\lemma{\skt{एकसफोद्भवं}}\Afootnote{H; \skt{मेकशफोद्भवम्}\ A1938; \skt{एकशफोद्भवं ।.}\ N}}\skt{ । }\pend
\pstart\versepar \skt{}\edtext{\skt{अल्पाभिष्यन्दिवर्ग्गोयं}}{\Afootnote{H; \skt{अल्पाभिष्यन्द्ययं वर्गो}\ A1938; \skt{अल्पाभिष्यन्दि वर्गोय}\ N}}\skt{ जाङ्गलः समुदाहृतः ॥}\pend

\pstart\versepar \ledsidenote{\passref{1.46.91}}\skt{दूरे }\edtext{\skt{जनानुनिलया}}{\Afootnote{H; \skt{जनान्तनिलया}\ A1938 N}}\skt{ दूरे पानीय गोचराः । }\pend
\pstart\versepar \skt{ये मृगाश्च विहङ्गाश्च }\edtext{\skt{तेऽल्पाभिष्यन्दिनो}}{\lemma{\skt{ऽल्पाभिष्यन्दिनो}}\Afootnote{A1938 H; \skt{ल्पाभिष्यन्दिनो}\ N}}\skt{ मताः ॥}\pend

\pstart\versepar \ledsidenote{\passref{1.46.92}}\edtext{\skt{अतीवासन्ननिलयाः}}{\Afootnote{A1938 H; \skt{अतीवासन्ननिलया}\ N}}\skt{ समीपोदकगोचराः । }\pend
\pstart\versepar \skt{ये मृगाश्च विहङ्गाश्च महाभिष्यन्दिनस्तु ते ॥}\pend

\pstart \ledsidenote{\passref{1.46.93}}\skt{आनूपवर्ग्गस्तु }\edtext{\skt{पञ्चविधस्तद्यथा}}{\lemma{\skt{तद्यथा}}\Afootnote{A1938 H; \skt{तवोद्यथा ।}\ N}}\skt{ ॥ }\edtext{\skt{कूलचराः}}{\Afootnote{A1938 H; \skt{कुलचाराः}\ N}}\skt{ प्लवाः }\edtext{\skt{कोशस्थाः}}{\Afootnote{N H; \skt{कोषस्थाः,}\ A1938}}\skt{ पादिनो मत्स्याश्चेति ॥}\pend

\pstart \ledsidenote{\passref{1.46.94}}\skt{तत्र }\edtext{\skt{गजगवयमहिषरुरुचमररोहितवराहखङ्गगोकर्ण्णकालपुच्छकोद्रन्यङ्कुकुरङ्गप्रभृतयः}}{\Afootnote{H; \skt{गजगवयमहिषरुरुचमरसृमररोहितवराहखड्गिगोकर्णकालपुच्छकोद्रन्यङ्क्वरण्यगवयप्रभृतयः}\ A1938; \skt{गजगवयमहिषरुरुचमररोहितवराहखड्गगोकर्ण्णकालपुच्छकोद्रन्यङ्कु कुरङ्गप्रभृतयः}\ N}}\skt{ }\edtext{\skt{कुलचराः}}{\Afootnote{N H; \skt{कूलचराः}\ A1938}}\skt{ पशवः ॥}\pend

\pstart\versepar \ledsidenote{\passref{1.46.95}}\skt{वातपित्तहरावृष्या मधुरा रसपाकयोः ॥ }\pend
\pstart\versepar \skt{}\edtext{\skt{शीताः}}{\Afootnote{H; \skt{शीतला बलिनः}\ A1938; \skt{शीता}\ N}}\skt{ }\edtext{\skt{स्निग्धाश्च}}{\Afootnote{N H; \skt{स्निग्धा}\ A1938}}\skt{ }\edtext{\skt{वल्याश्च}}{\Afootnote{N H; \emph{om.}\ A1938}}\skt{ }\edtext{\skt{मूत्रलाःकफवर्द्धनाः}}{\Afootnote{H; \skt{मूत्रलाः कफवर्धनः ॥९५॥}\ A1938; \skt{मूत्रला कफवर्द्धनाः ॥}\ N}}\skt{ ॥}\pend

\pstart\versepar \ledsidenote{\passref{1.46.96}}\skt{विरूक्षणो लेखनश्च वीर्योष्णः पित्तदूषणः । }\pend
\pstart\versepar \skt{}\edtext{\skt{स्वाद्वम्ललवणस्तेषां}}{\lemma{\skt{स्वाद्वम्ललवणस्}}\Afootnote{A1938 H; \skt{स्वाद्वम्व्ललवनस्}\ N}}\skt{ }\edtext{\skt{गजः}}{\Afootnote{A1938 H; \skt{गज}\ N}}\skt{ श्लेष्मानिलापहः ॥}\pend

\pstart\versepar \ledsidenote{\passref{1.46.98}}\edtext{\edtext{\edtext{\edtext{\skt{स्निग्धोष्णलवणम्वृष्यम्माहिषन्तर्प्पणङ्गुरुः}}{\lemma{\skt{स्निग्धोष्णलवणम्}}\Afootnote{N H; \skt{स्निग्धोष्णमधुरो}\ A1938}}}{\lemma{\skt{वृष्यम्}}\Afootnote{N H; \skt{वृष्यो}\ A1938}}}{\lemma{\skt{माहिषन्}}\Afootnote{N H; \skt{महिष}\ A1938}}}{\lemma{\skt{तर्प्पणङ्}}\Afootnote{N H; \skt{स्तर्पणो}\ A1938}}\skt{ । }\pend
\pstart\versepar \skt{}\edtext{\skt{निद्रापुंस्त्ववलष्यन्दि}}{\Afootnote{H; \skt{निद्रापुंस्त्वबलस्तन्य}\ A1938; \skt{निद्रापुंस्त्ववलस्यन्दि}\ N}}\skt{ }\edtext{\skt{वर्द्धनम्मान्सदार्ढ्यकृत्}}{\lemma{\skt{वर्द्धनम्}}\Afootnote{N H; \skt{वर्धनो}\ A1938}}\skt{ ॥}\pend

\pstart\versepar \ledsidenote{\passref{1.46.102}}\skt{स्वेदनं }\edtext{\skt{वृंहणम्वृष्यं}}{\lemma{\skt{वृंहणम्}}\Afootnote{H; \skt{बृंहणं}\ A1938; \skt{वृहणं}\ N}}\skt{ }\edtext{\edtext{\edtext{\skt{रोचनन्तर्प्पणङ्गुरुः}}{\lemma{\skt{रोचनन्}}\Afootnote{H; \skt{शीतलं}\ A1938; \skt{रोधनं}\ N}}}{\lemma{\skt{तर्प्पणङ्}}\Afootnote{A1938 N H\textsuperscript{pc}; \skt{तर्प्पङ्}\ H\textsuperscript{ac}}}}{\lemma{\skt{गुरुः}}\Afootnote{N H; \skt{गुरु ।}\ A1938}}\skt{ । }\pend
\pstart\versepar \skt{}\edtext{\skt{स्निग्धं}}{\Afootnote{N H; \emph{om.}\ A1938}}\skt{ श्रमानिलहरं }\edtext{\skt{वाराहं}}{\Afootnote{N H; \skt{स्निग्धं वाराहं}\ A1938}}\skt{ }\edtext{\skt{बलवर्द्धनं}}{\Afootnote{A1938 H; \skt{वलवर्द्धनं ॥}\ N}}\skt{ ॥}\pend

\pstart\versepar \ledsidenote{\passref{1.46.103}}\edtext{\skt{कफघ्नम्खड्गपिशितं}}{\lemma{\skt{खड्गपिशितं}}\Afootnote{H; \skt{खड्गिपिशितं}\ A1938; \skt{खङ्गपिशितं}\ N}}\skt{ कषायमनिलापहं । }\pend
\pstart\versepar \skt{पित्र्यं }\edtext{\skt{पवित्रमायुष्यं}}{\Afootnote{A1938 H; \skt{पवित्रमायुष्य}\ N}}\skt{ }\edtext{\skt{बद्धमूत्रं}}{\Afootnote{A1938 H; \skt{वद्धमूत्रं}\ N}}\skt{ विरूक्षणं ॥}\pend

\pstart \ledsidenote{\passref{1.46.105}}\edtext{\skt{हंससारसक्रोञ्चचक्रवाककुररकारण्डबकादम्बजीवञ्जीवकबकबलाकापुण्डरीकाप्लवशरारीमद्गुक्रोशकाकाक्षपुष्करशायिकाकुनालकम्बुकुक्कुटिकमेघरावस्वेतवालप्रभृतयः}}{\Afootnote{H; \skt{हंससारसक्रौञ्चचक्रवाककुररकादम्बकारण्डवजीवञ्जीवकबकबलाकापुण्डरीकप्लवशरारीमुख- नन्दीमुखमद्गूत्क्रोशकाचाक्षमल्लिकाक्षशुक्लाक्षपुष्करशायिकाकोनालकाम्बुकुक्कुटिकामेघरावश्वेतवारलप्रभृतयः}\ A1938; \skt{हंससारसक्रोञ्चचक्रवाककुररकारण्डवकादम्वजीवंजीववकवलाकापुण्डरीका प्लवशरारीमद्गूक्रोशकाकाक्षपुष्करशायिकाकुनालकाम्वुकुर्कुटिकाघरावश्वेतवालप्रभृतयः}\ N}}\skt{ प्लवाः संघातचारिणः ॥}\pend

\pstart\versepar \ledsidenote{\passref{1.46.106}}\skt{रक्तपित्तहराः शीताः स्निग्धा }\edtext{\skt{वृष्यानिलापहाः}}{\Afootnote{N H; \skt{वृष्या मरुज्जितः ।}\ A1938}}\skt{ । }\pend
\pstart\versepar \skt{सृष्टमूत्रपुरीषाश्च मधुरा रसपाकयोः ॥}\pend

\pstart\versepar \ledsidenote{\passref{1.46.107}}\edtext{\skt{गुरूष्णस्निग्धमधुरः}}{\Afootnote{N H; \skt{गुरूष्णमधुरः}\ A1938}}\skt{ }\edtext{\skt{स्वरवर्ण्णबलप्रदः}}{\Afootnote{H; \skt{स्निग्धः स्वरवर्णबलप्रदः ।}\ A1938; \skt{स्वरवर्ण्णवलप्रदः ।}\ N}}\skt{ । }\pend
\pstart\versepar \skt{}\edtext{\skt{बृंहणः}}{\Afootnote{A1938 H; \skt{वृंहणः}\ N}}\skt{ शुक्रलस्तेषां हंसो }\edtext{\skt{मारुतनाशनः}}{\Afootnote{N H; \skt{वातविकारनुत् ॥१०७॥}\ A1938}}\skt{ ॥}\pend

\pstart \ledsidenote{\passref{1.46.108}}\edtext{\skt{शङ्खनखशुक्तिशम्बूकबल्लूरप्रभृतयः}}{\Afootnote{H; \skt{शङ्खशङ्खनकशुक्तिशम्बूकभल्लूकप्रभृतयः}\ A1938; \skt{शंखनखशुक्तिशम्वूकवल्लूरप्रभृतयः}\ N}}\skt{ कोशस्थाः ॥}\pend

\pstart \ledsidenote{\passref{1.46.109}}\edtext{\skt{कूर्म्मकुम्भीरकर्क्कटप्रभृतयः}}{\Afootnote{\skt{कूर्मकुम्भीरकर्कटककृष्णकर्कटकशिशुमारप्रभृतयः}\ A1938; \skt{कुर्मकुम्भीरकर्कटप्रभृतयः}\ N; \skt{कूर्म्मकुम्भीरकर्क्कट…प्रभृतयः}\ H}}\skt{ पादिनः ॥}\pend

\pstart\versepar \ledsidenote{\passref{1.46.110}}\skt{शङ्खकूर्म्मादयः }\edtext{\skt{स्वादुरसपाकानिलापहाः}}{\Afootnote{H; \skt{स्वादुरसपाका मरुन्नुदः ।}\ A1938; \skt{स्वादुरसपाकानिलावहाः ।}\ N}}\skt{ । }\pend
\pstart\versepar \skt{शीताः }\edtext{\skt{स्निग्धाहिता}}{\Afootnote{N H; \skt{स्निग्धा हिताः}\ A1938}}\skt{ पित्ते }\edtext{\skt{वर्च्चस्याः}}{\Afootnote{A1938 H; \skt{वर्च्चस्या}\ N}}\skt{ }\edtext{\skt{शुक्रवर्द्धनाः}}{\Afootnote{N H; \skt{श्लेष्मवर्धनाः ॥११०॥}\ A1938}}\skt{ ॥}\pend

\pstart\versepar \ledsidenote{\passref{1.46.111}}\edtext{\edtext{\skt{कृष्णकर्क्कटकस्तेषाम्बल्यः}}{\lemma{\skt{कृष्णकर्क्कटकस्तेषाम्}}\Afootnote{H; \skt{कृष्णकर्कटकस्तेषां}\ A1938 N}}}{\lemma{\skt{बल्यः}}\Afootnote{A1938 H; \skt{वल्यः}\ N}}\skt{ }\edtext{\edtext{\skt{कोष्णोऽनिलापहः}}{\lemma{\skt{कोष्णो}}\Afootnote{A1938 H; \skt{कोष्णो}\ N}}}{\lemma{\skt{ऽनिलापहः}}\Afootnote{A1938 H; \skt{निलापहः ।.}\ N}}\skt{ । }\pend
\pstart\versepar \skt{}\edtext{\skt{शुक्लः}}{\Afootnote{A1938 H; \skt{शुक्ल}\ N}}\skt{ }\edtext{\edtext{\skt{सन्धानकृत्सृष्टविण्मूत्रोऽनिलपित्तहाः}}{\lemma{\skt{सन्धानकृत्सृष्टविण्मूत्रो}}\Afootnote{A1938 H; \skt{संधानकृत्सृष्टविट्मूट्रो}\ N}}}{\lemma{\skt{ऽनिलपित्तहाः}}\Afootnote{H; \skt{ऽनिलपित्तहा ॥१११॥}\ A1938; \skt{निलपित्तहाः ॥}\ N}}\skt{ ॥}\pend

\pstart \ledsidenote{\passref{1.46.112}}\skt{मत्स्यास्तु द्विविधा नादेयाः सामुद्राश्च ॥}\pend

\pstart \ledsidenote{\passref{1.46.113}}\skt{तत्र }\edtext{\skt{नादेयाः}}{\Afootnote{N H; \skt{नादेयाः-}\ A1938}}\skt{ }\edtext{\skt{पाठीनरोहिपाटलावर्म्मिगोमत्स्यवाकुचमुरलसहस्रदन्तप्रभृतयः}}{\Afootnote{H; \skt{रोहितपाठीनपाटलाराजीववर्मिगोमत्स्यकृष्णमत्स्यवागुञ्जारमुरलसहस्रदंष्ट्रप्रभृतयः ॥११३॥}\ A1938; \skt{पाठीनरोहिपाटलावर्म्मिगोमत्स्यवाकुचमुरलसह सुदत्तप्रभृतयः ॥}\ N}}\skt{ ॥}\pend

\pstart\versepar \ledsidenote{\passref{1.46.114}}\skt{नादेया }\edtext{\skt{गुरवो}}{\Afootnote{N H; \emph{om.}\ A1938}}\skt{ }\edtext{\skt{मत्स्या}}{\Afootnote{N H; \emph{om.}\ A1938}}\skt{ }\edtext{\skt{मधुरावतनाशनाः}}{\Afootnote{H; \skt{मधुरा मत्स्या गुरवो मारुतापहाः ।}\ A1938; \skt{मधुरा वातनाशनाः ।}\ N}}\skt{ । }\pend
\pstart\versepar \skt{}\edtext{\edtext{\skt{रक्त्रपित्तहरास्तूष्णा}}{\lemma{\skt{रक्त्रपित्तहरास्}}\Afootnote{H; \skt{रक्तपित्तकराश्}\ A1938 N}}}{\lemma{\skt{तूष्णा}}\Afootnote{H; \skt{चोष्णा}\ A1938; \skt{तृष्णा}\ N}}\skt{ }\edtext{\skt{वृष्या}}{\Afootnote{N H; \skt{वृष्याः}\ A1938}}\skt{ स्निग्धाल्पवर्च्चसः ॥}\pend

\pstart\versepar \ledsidenote{\passref{1.46.115}}\skt{कषायानुरसस्तेषां }\edtext{\skt{शष्पशैवलभोजनः}}{\Afootnote{H; \skt{शष्पशैवालभोजनः ।}\ A1938; \skt{शस्पशैवालभोजनाः ।}\ N}}\skt{ । }\pend
\pstart\versepar \skt{}\edtext{\skt{रोहीत}}{\Afootnote{N H; \skt{रोहितो मारुत}\ A1938}}\skt{ }\edtext{\skt{मत्स्यौ}}{\Afootnote{H; \skt{हरो}\ A1938; \skt{मत्स्यो}\ N}}\skt{ नात्यर्थं }\edtext{\skt{रक्तपित्तप्रकोपनः}}{\Afootnote{H; \skt{पित्तकोपनः ॥११५॥}\ A1938; \skt{रक्तपित्तप्रकोपिनः ।}\ N}}\skt{ ॥}\pend

\pstart\versepar \ledsidenote{\passref{1.46.117}}\edtext{\skt{सरस्तडागसंभूताः}}{\Afootnote{A1938 H; \skt{सरस्तडागसंभूताः}\ N}}\skt{ स्निग्धाः स्वादुरसाः स्मृताः । }\pend
\pstart\versepar \skt{}\edtext{\skt{महाह्रैदेषु}}{\Afootnote{H\textsuperscript{ac}; \skt{महाह्रदेषु}\ A1938 N H\textsuperscript{pc}}}\skt{ }\edtext{\skt{बलिनः}}{\Afootnote{A1938 H; \skt{वलिनः}\ N}}\skt{ }\edtext{\edtext{\skt{स्वल्पेऽम्भस्यबलामताः}}{\lemma{\skt{स्वल्पे}}\Afootnote{A1938 H; \skt{स्वल्पेम्}\ N}}}{\lemma{\skt{ऽम्भस्यबलामताः}}\Afootnote{H; \skt{ऽम्भस्यबलाः स्मृताः ॥११७॥}\ A1938; \skt{भस्य वला मताः ॥}\ N}}\skt{ ॥}\pend

\pstart \ledsidenote{\passref{1.46.118}}\edtext{\edtext{\skt{तिमितिमिङ्गिलकुलिशकाकमत्स्यभिरलनन्दीवरमकरगर्ग्गरचन्द्रकमहामीनराजमत्स्यप्रभृतयः}}{\lemma{\skt{तिमितिमिङ्गिलकुलिशकाकमत्स्यभिर्}}\Afootnote{H; \skt{तिमितिमिङ्गिलकुलिशपाकमत्स्यनि}\ A1938; \skt{तिमितिमिंगिलकुलिशकाकमत्स्यरि}\ N}}}{\lemma{\skt{अलनन्दीवरमकरगर्ग्गरचन्द्रकमहामीनराजमत्स्यप्रभृतयः}}\Afootnote{H; \skt{रुलनन्दिवार(रु)लकमकरगर्गरचन्द्रकमहामीनराजीवप्रभृतयः}\ A1938; \skt{रलनंदीवरमकरगर्गरचन्द्रकमहामीनराजमत्स्य प्रभृतयः}\ N}}\skt{ सामुद्राः ॥}\pend

\pstart\versepar \ledsidenote{\passref{1.46.119}}\skt{सामुद्रा }\edtext{\skt{गुरवः}}{\Afootnote{A1938 H; \skt{गुरव}\ N}}\skt{ स्निग्धा मधुरा नातिपित्तलाः । }\pend
\pstart\versepar \skt{उष्णा वातहरा वृष्या वर्च्चस्याः श्लेष्मवर्द्धनाः ॥}\pend

\pstart\versepar \ledsidenote{\passref{1.46.120}}\edtext{\skt{बलावहा}}{\Afootnote{A1938 H; \skt{वलावहा}\ N}}\skt{ }\edtext{\skt{विशेषेण}}{\Afootnote{A1938 N\textsuperscript{pc} H; \skt{विशेषेमाण}\ N\textsuperscript{ac}}}\skt{ मांसाशित्वात्समुद्रजाः । }\pend
\pstart\versepar \skt{समुद्रजेभ्यो नादेया }\edtext{\skt{बृंहणत्वाद्गुणोत्तराः}}{\Afootnote{A1938 H; \skt{वृहणत्वाद्गुणोत्तराः ।}\ N}}\skt{ ॥}\pend

\pstart\versepar \ledsidenote{\passref{1.46.121}}\edtext{\edtext{\skt{तस्मादत्यनिलघ्नत्वादौत्सो}}{\lemma{\skt{तस्मादत्य्}}\Afootnote{H; \skt{तेषामप्य्}\ A1938; \skt{तस्माद}\ N}}}{\lemma{\skt{अनिलघ्नत्वादौत्सो}}\Afootnote{H; \skt{अनिलघ्नत्वाच्चौण्ट्य}\ A1938; \skt{त्यनिलघ्नत्वादौत्सो}\ N}}\skt{ }\edtext{\skt{ज्ञेया}}{\Afootnote{H; \skt{कौप्यौ}\ A1938; \skt{ज्ञेयो}\ N}}\skt{ }\edtext{\skt{गुणोत्तरः}}{\Afootnote{H; \skt{गुणोत्तरौ ।}\ A1938; \skt{गुणोत्तराः ।}\ N}}\skt{ । }\pend
\pstart\versepar \skt{}\edtext{\skt{स्निग्धत्वात्स्वादुपाकत्वात्तयोर्वाप्यो}}{\lemma{\skt{वाप्यो}}\Afootnote{H; \skt{वाप्या}\ A1938; \skt{वाप्या}\ N}}\skt{ }\edtext{\skt{गुणाधिकः}}{\Afootnote{H; \skt{गुणाधिकाः}\ A1938 N}}\skt{ ॥}\pend

\pstart\versepar \ledsidenote{\passref{1.46.122}}\skt{नादेया गुरवो }\edtext{\skt{मेध्यो}}{\Afootnote{H; \skt{मध्ये}\ A1938; \skt{मेध्या}\ N}}\skt{ }\edtext{\skt{यस्मात्पुच्छास्यचारिणः}}{\Afootnote{A1938 H; \skt{यस्मात्पुच्छांस्य चारिणः ।}\ N}}\skt{ । }\pend
\pstart\versepar \skt{सरस्तडागजानान्तु विशेषेण }\edtext{\skt{शिरो}}{\Afootnote{A1938 H; \skt{सिरो}\ N}}\skt{ }\edtext{\skt{लघुः}}{\Afootnote{N H; \skt{लघु ॥१२२॥}\ A1938}}\skt{ ॥}\pend

\pstart\versepar \ledsidenote{\passref{1.46.123}}\skt{अदूरगोचरा }\edtext{\skt{यस्मात्तस्मादौत्सोदपानजाः}}{\Afootnote{H; \skt{यस्मात्तस्मादुत्सोदपानजाः ।}\ A1938; \skt{यस्मात्तस्मादौत्सानपानजाः ।}\ N}}\skt{ । }\pend
\pstart\versepar \skt{किञ्चिन्मुक्त्वा }\edtext{\skt{शिरोदेशमत्यर्थङ्गुरवस्तु}}{\lemma{\skt{शिरोदेशम्}}\Afootnote{A1938 H; \skt{सिरोदेशम्}\ N}}\skt{ ते ॥}\pend

\pstart\versepar \ledsidenote{\passref{1.46.124}}\edtext{\skt{अधस्ताङ्गुरवो}}{\Afootnote{H; \skt{अधस्ताद्गुरवो}\ A1938 N}}\skt{ ज्ञेया }\edtext{\skt{मत्स्याः}}{\Afootnote{A1938 H; \skt{मत्स्या}\ N}}\skt{ सागरसम्भवाः । }\pend
\pstart\versepar \skt{उरोविचरणात्तेषां }\edtext{\edtext{\skt{पूर्व्वमङ्गं}}{\lemma{\skt{पूर्व्वम्}}\Afootnote{H; \skt{पूर्व}\ A1938 N}}}{\lemma{\skt{अङ्गं}}\Afootnote{H; \skt{मङ्गं}\ A1938 N}}\skt{ लघु स्मृतं ॥}\pend

\pstart \ledsidenote{\passref{1.46.125}}\skt{इत्यानूपो }\edtext{\skt{महास्यन्दी}}{\Afootnote{N H; \skt{महाष्यन्दी}\ A1938}}\skt{ }\edtext{\skt{मांसवर्ग}}{\Afootnote{A1938 H; \skt{मांसवर्गम्}\ N}}\skt{ उदीरितः ॥}\pend

\pstart \ledsidenote{\passref{1.46.126}}\skt{तत्र }\edtext{\edtext{\skt{शुष्कपूतिदिग्धविद्धसर्प्पापराद्धजीर्णकृशवालानामसात्य्मचारिणाञ्च}}{\lemma{\skt{शुष्कपूतिदिग्धविद्धसर्प्पापराद्धजीर्णकृशवालानाम्}}\Afootnote{H; \skt{शुष्कपूतिव्याधिविषसर्पहतदिग्धविद्धजीर्णकृशबालाना}\ A1938; \skt{शुष्कपूतिविद्धदिग्धसर्पापराद्धजीर्ण्णकृशवालानाम्}\ N}}}{\lemma{\skt{असात्य्मचारिणाञ्च}}\Afootnote{H; \skt{मसात्म्यचारिणां च}\ A1938; \skt{असात्म्यचारिणां च}\ N}}\skt{ }\edtext{\skt{मांसान्यभक्षाणि}}{\lemma{\skt{अभक्षाणि}}\Afootnote{H; \skt{अभक्ष्याणि}\ A1938 N}}\skt{ }\edtext{\skt{भवन्ति}}{\Afootnote{N H; \emph{om.}\ A1938}}\skt{ ॥ }\edtext{\skt{कस्मात्विगतव्यापन्नापरिणताल्पासम्पूर्ण्णा}}{\Afootnote{H; \skt{यस्माद्विगतव्यापन्नापहतपरिणताल्पासम्पूर्ण}\ A1938; \skt{कस्माद्विगतव्यापन्नापरिणताल्पासंपूर्ण्ण}\ N}}\skt{ }\edtext{\skt{यथार्थकरत्वाद्दोषकराणि}}{\Afootnote{H; \skt{वीर्यत्वाद्दोषकराणि}\ A1938; \skt{यथार्थकरत्वाद्दोषकरानि}\ N}}\skt{ भवन्ति । }\edtext{\edtext{\skt{एभ्योऽन्येषामुपनादेयम्मांसमिति}}{\lemma{\skt{ऽन्येषाम्}}\Afootnote{H; \skt{ऽन्येषा}\ A1938; \skt{न्येषाम्}\ N}}}{\lemma{\skt{उपनादेयम्}}\Afootnote{H; \skt{मुपादेयं}\ A1938; \skt{उपादेयम्}\ N}}\skt{ ॥}\pend

\pstart \ledsidenote{\passref{1.46.129}}\edtext{\skt{स्त्रियश्चतुष्पदेषु}}{\lemma{\skt{चतुष्पदेषु}}\Afootnote{H; \skt{चतुष्पात्सु,}\ A1938; \skt{चततुष्पदेषु}\ N}}\skt{ }\edtext{\skt{मांसो}}{\Afootnote{H; \skt{पुमांसो}\ A1938 N}}\skt{ विहङ्गेषु महाशरीरेष्वल्पशरीरा अल्पशरीरेषु महाशरीराः प्रधानतमा }\edtext{\skt{भवन्ति}}{\Afootnote{N H; \skt{एवम्}\ A1938}}\skt{ ॥ }\edtext{\skt{एकजातीयानामपि}}{\lemma{\skt{अपि}}\Afootnote{N H; \emph{om.}\ A1938}}\skt{ महाशरीरेभ्यः कृशशरीराः }\edtext{\edtext{\skt{प्रधानतम्भवन्ति}}{\lemma{\skt{प्रधानतम्}}\Afootnote{H; \skt{प्रधानतमाः ॥१२९॥}\ A1938; \skt{प्रधानतमा}\ N}}}{\lemma{\skt{भवन्ति}}\Afootnote{N H; \emph{om.}\ A1938}}\skt{ ॥}\pend

\pstart \ledsidenote{\passref{1.46.130}}\edtext{\skt{तत्र}}{\Afootnote{N H; \emph{om.}\ A1938}}\skt{ }\edtext{\skt{स्थानादिकृतं तु}}{\Afootnote{H; \skt{स्थानादिकृतं}\ A1938; \skt{स्थानादिकृतन्दु}\ N}}\skt{ मांसस्य गुरुलाघवमुपदेक्ष्यामः । तद्यथा }\edtext{\skt{रक्तादिषु}}{\Afootnote{N H; \skt{रक्तादिषु शुक्रान्तेषु}\ A1938}}\skt{ }\edtext{\edtext{\skt{धातुषूत्तरोत्तरोस्तथान्यकृत्कालेयसक्थिकटिपृष्ठचरणशिरांसि}}{\lemma{\skt{धातुषूत्तरोत्तरोस्}}\Afootnote{H; \skt{धातुषूत्तरोत्तरा गुरु}\ A1938; \skt{धातुषूत्तरोत्त}\ N}}}{\lemma{\skt{तथान्यकृत्कालेयसक्थिकटिपृष्ठचरणशिरांसि}}\Afootnote{H; \skt{तराः, तथा सक्थिस्कन्धक्रोडशिरःपादकरकटीपृष्ठचर्मकालेयकयकृदन्त्राणि ॥१३०॥}\ A1938; \skt{रास्तथान्तयकृत्कालेयक सक्थ्यिकटिपृष्ठचरणशिरांसि}\ N}}\pend

\pstart\versepar \ledsidenote{\passref{1.46.131}}\edtext{\skt{उरस्कन्धौ}}{\Afootnote{N H; \skt{शिरःस्कन्धं कटी पृष्ठं}\ A1938}}\skt{ सक्थिनी }\edtext{\skt{चात्मपक्षयोः}}{\Afootnote{A1938; \skt{चापक्वयोः ।}\ N; \skt{चामपक्वयोः ॥}\ H}}\skt{ ॥ }\pend
\pstart\versepar \skt{}\edtext{\skt{गुरवश्च}}{\Afootnote{N H; \skt{गुरु}\ A1938}}\skt{ }\edtext{\skt{यथा}}{\Afootnote{N H; \emph{om.}\ A1938}}\skt{ }\edtext{\skt{पूर्वन्धातवश्च}}{\lemma{\skt{धातवश्च}}\Afootnote{N H; \skt{विजानीयाद्धातवस्तु}\ A1938}}\skt{ यथोत्तरं ।}\pend

\pstart \ledsidenote{\passref{1.46.132cd}}\edtext{\skt{पूर्व्व}}{\Afootnote{A1938 H; \skt{पूव्व}\ N}}\skt{ }\edtext{\skt{भागो}}{\Afootnote{A1938 H; \skt{भागे ।}\ N}}\skt{ गुरुः }\edtext{\skt{पुंसाम्मध्यभागश्च}}{\lemma{\skt{मध्यभागश्च}}\Afootnote{N H; \skt{अधोभागस्तु}\ A1938}}\skt{ }\edtext{\skt{योषितां}}{\Afootnote{A1938 H; \skt{योषिता ।}\ N}}\skt{ ॥}\pend

\pstart\versepar \ledsidenote{\passref{1.46.133}}\edtext{\skt{उरूग्रीवा}}{\Afootnote{H; \skt{उरोग्रीवं}\ A1938; \skt{उरोग्रीवा}\ N}}\skt{ विहङ्गानाम्विशेषेण }\edtext{\skt{गुरू}}{\Afootnote{H; \skt{गुरु}\ A1938 N}}\skt{ स्मृतं । }\pend
\pstart\versepar \skt{पक्षोत्क्षेपात्समो दृष्टो }\edtext{\skt{मध्यभागश्च}}{\Afootnote{N H; \skt{मध्यभागस्तु}\ A1938}}\skt{ पक्षिणां ॥}\pend

\pstart\versepar \ledsidenote{\passref{1.46.134}}\skt{अतीव }\edtext{\skt{रूक्षम्मान्सान्तु}}{\lemma{\skt{मान्सान्तु}}\Afootnote{H; \skt{मांसं तु}\ A1938; \skt{मांम्सं तु}\ N}}\skt{ विहंगानां फलाशिनां । }\pend
\pstart\versepar \skt{}\edtext{\skt{वृंहणं}}{\Afootnote{H; \skt{बृंहणं}\ A1938; \skt{वृंहनम्}\ N}}\skt{ }\edtext{\edtext{\edtext{\skt{मांसमत्यर्थम्विहङ्गानाम्मांसभाजिनां}}{\lemma{\skt{मांसम्}}\Afootnote{A1938 H; \skt{मासम्}\ N}}}{\lemma{\skt{विहङ्गानाम्}}\Afootnote{N H; \skt{खगानां}\ A1938}}}{\lemma{\skt{मांसभाजिनां}}\Afootnote{H; \skt{पिशिताशिनाम् ॥१३४॥}\ A1938; \skt{मांसभोजिनां ।}\ N}}\skt{ ॥}\pend

\pstart\versepar \ledsidenote{\passref{1.46.135}}\edtext{\edtext{\edtext{\skt{मत्याशिनाम्पित्तहरम्वातघ्नन्धान्यचारिणां}}{\lemma{\skt{मत्याशिनाम्}}\Afootnote{H; \skt{मत्स्याशिनां}\ A1938 N}}}{\lemma{\skt{पित्तहरम्}}\Afootnote{N H; \skt{पित्तकरं}\ A1938}}}{\lemma{\skt{धान्यचारिणां}}\Afootnote{N H; \skt{धान्यचारिनाम्}\ A1938}}\skt{ । }\pend
\pstart\versepar \skt{जलजानूपजा }\edtext{\skt{ग्राम्या}}{\Afootnote{A1938 H; \skt{ग्राम्याः}\ N}}\skt{ }\edtext{\skt{क्रव्यादैकसफास्तथा}}{\lemma{\skt{क्रव्यादैकसफास्}}\Afootnote{\skt{क्रच्यादैकशफास्}\ A1938; \skt{क्रव्यादैकशफास्}\ N; \skt{क्रदैकसफास्}\ H}}\skt{ ॥}\pend

\pstart\versepar \ledsidenote{\passref{1.46.136}}\edtext{\skt{प्रहअ}}{\Afootnote{H; \skt{प्रसहा}\ A1938; \skt{प्रवहा}\ N}}\skt{ }\edtext{\skt{विलवासाश्च}}{\Afootnote{N H; \skt{बिलवासाश्च}\ A1938}}\skt{ }\edtext{\skt{तथा}}{\Afootnote{N H; \skt{ये च}\ A1938}}\skt{ }\edtext{\skt{जङ्घालसंज्ञिताः}}{\Afootnote{A1938 H; \skt{जंगलसंज्ञिताः ।}\ N}}\skt{ । }\pend
\pstart\versepar \skt{}\edtext{\skt{प्रनुदा}}{\Afootnote{H; \skt{प्रतुदा}\ A1938 N}}\skt{ विष्किराश्चैव }\edtext{\skt{लघवः स्युर्यथोत्तरं}}{\lemma{\skt{लघवः स्युर्}}\Afootnote{A1938 H; \skt{लघवः स्यु}\ N}}\skt{ ॥ }\pend
\pstart\versepar \skt{अल्पाभिष्यन्दिनश्चैव }\edtext{\edtext{\skt{यथापूर्व्वमतो}}{\lemma{\skt{यथापूर्व्वम्}}\Afootnote{N H; \skt{यथापूर्व}\ A1938}}}{\lemma{\skt{अतो}}\Afootnote{N H; \skt{मतो}\ A1938}}\skt{ }\edtext{\skt{न्यथा}}{\Afootnote{N H; \skt{ऽन्यथा ॥१३६॥}\ A1938}}\skt{ ।}\pend

\pstart \ledsidenote{\passref{1.46.137}}\edtext{\skt{सर्वशरीरेभ्यः}}{\Afootnote{N H; \skt{प्रामाणाधिकास्तु स्वजातावल्पसारा गुरवश्च । सर्वप्राणिनां सर्वशरीरेभ्यः}\ A1938}}\skt{ }\edtext{\edtext{\skt{प्रधानतमावर्त्तिभ्यामादद्यात्}}{\lemma{\skt{प्रधानतमावर्त्तिभ्याम्}}\Afootnote{H; \skt{प्रधानतमा भवन्ति यकृत्प्रदेशवर्तिन}\ A1938; \skt{प्रधानतमा वर्त्तिन्याम्}\ N}}\skt{ ॥ प्रधानालाभे}}{\lemma{\skt{आदद्यात् ॥ प्रधानालाभे}}\Afootnote{N H; \skt{स्तानाददीत; प्रधानालाभे}\ A1938}}\skt{ }\edtext{\skt{मध्यमवयसस्तु}}{\Afootnote{N H; \skt{मध्यमवय}\ A1938}}\skt{ }\edtext{\skt{मान्सं}}{\Afootnote{N H; \skt{स्कं}\ A1938}}\skt{ }\edtext{\edtext{\edtext{\edtext{\skt{सद्यस्कमक्लिष्टमुपादेयमिति}}{\lemma{\skt{सद्यस्कम्}}\Afootnote{N H; \skt{सद्यस्क}\ A1938}}}{\lemma{\skt{अक्लिष्टम्}}\Afootnote{N H; \skt{मक्लिष्ट}\ A1938}}}{\lemma{\skt{उपादेयम्}}\Afootnote{N H; \skt{मुपादेयं}\ A1938}}}{\lemma{\skt{इति}}\Afootnote{N H; \skt{मांसमिति ॥१३७॥}\ A1938}}\skt{ ॥}\pend

\pstart\versepar \ledsidenote{\passref{1.46.138}}\edtext{\skt{चरः}}{\Afootnote{N H; \skt{भवति चात्र चरः}\ A1938}}\skt{ शरीरावयवाः स्वभावो धातवः क्रिया । }\pend
\pstart\versepar \skt{}\edtext{\skt{लिङ्गम्प्रणां}}{\lemma{\skt{प्रणां}}\Afootnote{H; \skt{प्रमाणं}\ A1938 N}}\skt{ संस्कारो मात्रा }\edtext{\skt{चास्मिन्परीक्षते}}{\Afootnote{H; \skt{चास्मिन्परीक्ष्यते ॥१३८॥ इति मांसवर्गः ।}\ A1938; \skt{चास्मि परीक्ष्यते ॥}\ N}}\skt{ ॥ ॥}\pend

\pstart\versepar \ledsidenote{\passref{1.46.139}}\edtext{\skt{अत}}{\Afootnote{N H; \skt{अथ फलवर्गः । अत}\ A1938}}\skt{ }\edtext{\skt{ऊर्ध्वं}}{\Afootnote{A1938 H; \skt{ऊर्ध्वं ल}\ N}}\skt{ }\edtext{\edtext{\skt{फलान्युपदेक्ष्यामः}}{\lemma{\skt{फलान्य्}}\Afootnote{N H; \skt{फला}\ A1938}}}{\lemma{\skt{उपदेक्ष्यामः}}\Afootnote{N H; \skt{न्युपदेक्ष्यामः ।}\ A1938}}\skt{ ॥ }\pend
\pstart\versepar \skt{तद्यथा ॥ }\edtext{\edtext{\skt{दाडिमामलकवदरकोलकर्कन्धुकपित्थमातुलुङ्गाम्प्रमाणं}}{\lemma{\skt{दाडिमामलकवदरकोलकर्कन्धुकपित्थमातुलुङ्गाम्}}\Afootnote{H; \skt{दाडिमामलकबदरकोलकर्कन्धुसौवीरसिञ्चितिकाफलकपित्थमातुलुङ्}\ A1938; \skt{दाडिमामलकवदरकोलकर्क्कन्दुकपित्थमातुलुं}\ N}}}{\lemma{\skt{प्रमाणं}}\Afootnote{H; \skt{गा}\ A1938 N}}\skt{ }\edtext{\skt{म्राम्रातकलकुचकर}}{\Afootnote{H\textsuperscript{pc}; \skt{म्राम्रातककर}\ A1938; \skt{म्रातकलकुचकर}\ N H\textsuperscript{ac}}}\skt{ }\edtext{\skt{मर्द्दनभव्यपियालपारावतवेत्रफलप्राचीनो}}{\Afootnote{H; \skt{मर्दप्रियालनारङ्ग- जम्बीरलकुचभव्यपारावतवेत्रफलप्राचीना}\ A1938; \skt{मर्द्नभव्यपियालपारावतवेत्रफलप्राचीनाम}\ N}}\skt{ }\edtext{\skt{मलकतिन्तिडीकनीपकुशाम्राम्लीतकनागरङ्गप्रभृतीनि}}{\Afootnote{H; \skt{मलकतिन्तिडीकनीपकोशाम्राम्लीकाप्रभृतीनि ॥१३९॥}\ A1938; \skt{मलकतिन्तिडीकनीपकुशाम्राम्लीतकनागर प्रभृतीनि ।}\ N}}\skt{ ।}\pend

\pstart \ledsidenote{\passref{1.46.139.add}}\skt{एतान्यम्लान्यनिलं हन्युरोष्णाद्विपाके मधुराणि च ॥}\pend

\pstart\versepar \ledsidenote{\passref{1.46.141}}\edtext{\edtext{\edtext{\skt{कषायानुरसस्तेषांन्दाडिमन्नातिपित्तलं}}{\lemma{\skt{कषायानुरसस्}}\Afootnote{H; \skt{कषायानुरसं}\ A1938 N}}}{\lemma{\skt{तेषांन्}}\Afootnote{H; \skt{तेषां}\ A1938 N}}}{\lemma{\skt{दाडिमन्}}\Afootnote{A1938 H; \skt{डाडिमन्}\ N}}\skt{ । }\pend
\pstart\versepar \skt{दीपनीयं रुचिकरं हृद्यं }\edtext{\skt{वर्च्चो}}{\Afootnote{N H; \skt{वर्चो}\ A1938}}\skt{ }\edtext{\skt{निवन्धनं}}{\Afootnote{N H; \skt{विबन्धनम् ॥१४१॥}\ A1938}}\skt{ ॥}\pend

\pstart\versepar \ledsidenote{\passref{1.46.142}}\skt{द्विविधन्तत्तु विज्ञेयम्मधुरञ्चाम्लमेव च । }\pend
\pstart\versepar \skt{त्रिदोषघ्नं }\edtext{\edtext{\skt{मधुरमम्लम्वातकफापहं}}{\lemma{\skt{मधुरम्}}\Afootnote{H; \skt{तु मधुर}\ A1938; \skt{समधुरं}\ N}}}{\lemma{\skt{अम्लम्}}\Afootnote{H; \skt{मम्लं}\ A1938; \skt{अम्व्लम्}\ N}}\skt{ ॥}\pend

\pstart\versepar \ledsidenote{\passref{1.46.143}}\edtext{\skt{अम्लं}}{\Afootnote{A1938 H; \skt{अम्व्लं}\ N}}\skt{ समधुरन्तिक्तं }\edtext{\skt{कषायं}}{\Afootnote{A1938 H; \skt{कषाय}\ N}}\skt{ कटुकं सरं । }\pend
\pstart\versepar \skt{चक्षुष्यं सर्व्वदोषघ्नं }\edtext{\edtext{\skt{वृष्यमामलकं}}{\lemma{\skt{वृष्यम्}}\Afootnote{N H; \skt{वृष्य}\ A1938}}}{\lemma{\skt{आमलकं}}\Afootnote{H; \skt{मामलकी}\ A1938; \skt{आमलक}\ N}}\skt{ फलं ॥}\pend

\pstart\versepar \ledsidenote{\passref{1.46.144}}\skt{हन्ति }\edtext{\edtext{\skt{वातन्तदम्लत्वात्पित्तम्माधुर्यशत्यतः}}{\lemma{\skt{तदम्लत्वात्पित्तम्}}\Afootnote{H; \skt{तदम्लत्वात्पित्तं}\ A1938; \skt{तदम्व्लत्वात्पित्तम्}\ N}}}{\lemma{\skt{माधुर्यशत्यतः}}\Afootnote{H; \skt{माधुर्यशैत्यतः ।}\ A1938; \skt{माधूर्यशैत्यतः ।}\ N}}\skt{ । }\pend
\pstart\versepar \skt{कफं }\edtext{\edtext{\edtext{\skt{रूक्षकषायत्वात्फलेभ्योऽभ्यधिकन्ततः}}{\lemma{\skt{रूक्षकषायत्वात्फलेभ्यो}}\Afootnote{A1938 H; \skt{रूक्षकषायत्वात्फलेभ्यो}\ N}}}{\lemma{\skt{ऽभ्यधिकन्}}\Afootnote{A1938 H; \skt{भ्यधिकन्}\ N}}}{\lemma{\skt{ततः}}\Afootnote{N H; \skt{च तत् ॥१४४॥}\ A1938}}\skt{ ॥}\pend

\pstart\versepar \ledsidenote{\passref{1.46.145}}\edtext{\edtext{\edtext{\skt{कर्क्कन्धुकोलवरमम्लम्वातकफापहं}}{\lemma{\skt{कर्क्कन्धुकोलवरम्}}\Afootnote{H; \skt{कर्कन्धुकोलबदर}\ A1938; \skt{कर्कन्धुकोलवदरम्}\ N}}}{\lemma{\skt{अम्लम्}}\Afootnote{H; \skt{मामं}\ A1938; \skt{अम्व्लम्}\ N}}}{\lemma{\skt{वातकफापहं}}\Afootnote{N H; \skt{पित्तकफावहम् ।}\ A1938}}\skt{ । }\pend
\pstart\versepar \skt{}\edtext{\skt{पक्क्वम्पित्तानिलहरं}}{\lemma{\skt{पक्क्वम्}}\Afootnote{H; \skt{पक्वं}\ A1938 N}}\skt{ स्निग्धं समधुरं सरं ॥}\pend

\pstart\versepar \ledsidenote{\passref{1.46.146}}\edtext{\skt{पुराणं}}{\Afootnote{N H; \skt{पुरात}\ A1938}}\skt{ }\edtext{\skt{भग्नशमनं}}{\Afootnote{H; \skt{नं तृट्शमनं}\ A1938; \skt{तृप्रशमनं}\ N}}\skt{ }\edtext{\skt{श्रमघं}}{\Afootnote{H; \skt{श्रमघ्नं}\ A1938 N}}\skt{ }\edtext{\skt{लघुदीपनं}}{\Afootnote{N H; \skt{दीपनं}\ A1938}}\skt{ । }\pend
\pstart\versepar \skt{}\edtext{\skt{सौवीरं}}{\Afootnote{H; \skt{लघु । सौवीरं}\ A1938; \skt{सौवीर}\ N}}\skt{ }\edtext{\skt{वदरं}}{\Afootnote{N H; \skt{बदरं}\ A1938}}\skt{ }\edtext{\edtext{\skt{स्निग्धम्मधुरम्वातपित्तजित्}}{\lemma{\skt{स्निग्धम्}}\Afootnote{A1938 H; \skt{स्निघ्नं}\ N}}}{\lemma{\skt{वातपित्तजित्}}\Afootnote{A1938 H; \skt{वात्तपित्तजित् ।}\ N}}\skt{ ॥}\pend

\pstart\versepar \ledsidenote{\passref{1.46.147}}\skt{कषायं स्वा दु }\edtext{\skt{सांग्राही}}{\Afootnote{H; \skt{सङ्ग्राहि}\ A1938; \skt{संग्राही}\ N}}\skt{ शीतं सिञ्चितिकाफलम् ॥ }\pend
\pstart\versepar \skt{}\edtext{\skt{आमम्विषघ्नमस्वर्य}}{\lemma{\skt{विषघ्नमस्वर्य}}\Afootnote{H; \skt{कपित्थमस्वर्यं}\ A1938; \skt{विषघ्नमस्वर्यं}\ N}}\skt{ }\edtext{\skt{कपित्थं}}{\Afootnote{N H; \skt{कफघ्नं}\ A1938}}\skt{ ग्राहि वातलं ॥}\pend

\pstart\versepar \ledsidenote{\passref{1.46.148}}\edtext{\edtext{\edtext{\skt{कफानिलहरम्पक्वन्मधुराम्लरसङ्गुरुः}}{\lemma{\skt{कफानिलहरम्}}\Afootnote{A1938 H; \skt{कफनीलहरं}\ N}}}{\lemma{\skt{पक्वन्मधुराम्लरसङ्}}\Afootnote{H; \skt{पक्वं मधुराम्लरसं}\ A1938; \skt{पक्वं म्मधुराम्व्लरसं}\ N}}}{\lemma{\skt{गुरुः}}\Afootnote{N H; \skt{गुरु ।}\ A1938}}\skt{ ॥ }\pend
\pstart\versepar \skt{}\edtext{\skt{श्वासकासोरुचिहरन्घ्नं}}{\Afootnote{H; \skt{श्वासकासारुचिहरं तृष्णाघ्नं}\ A1938 N}}\skt{ }\edtext{\skt{कण्ठविशोधनं}}{\Afootnote{N H; \skt{कण्ठशोधनम् ॥१४८॥}\ A1938}}\skt{ ।}\pend

\pstart\versepar \ledsidenote{\passref{1.46.149}}\edtext{\skt{लघ्वम्लं}}{\Afootnote{A1938 H; \skt{लघ्वम्व्ल}\ N}}\skt{ दीपनं }\edtext{\edtext{\edtext{\skt{स्निग्धम्मातुलुङ्गम्मुदाहृतं}}{\lemma{\skt{स्निग्धम्}}\Afootnote{H; \skt{हृद्यं}\ A1938; \skt{स्निग्धां}\ N}}}{\lemma{\skt{मातुलुङ्गम्}}\Afootnote{A1938 H; \skt{मातुलङ्गम्}\ N}}}{\lemma{\skt{मुदाहृतं}}\Afootnote{H; \skt{उदाहृतम्}\ A1938 N}}\skt{ ॥ }\pend
\pstart\versepar \skt{त्वक्तिक्ता }\edtext{\skt{दुर्ज्जरास्तस्य}}{\lemma{\skt{दुर्ज्जरास्}}\Afootnote{N H; \skt{दुर्जरा}\ A1938}}\skt{ }\edtext{\skt{वातक्रिमिकफापहाः}}{\Afootnote{N H; \skt{वातक्रिमिकफापहा ॥१४९॥}\ A1938}}\skt{ ।}\pend

\pstart\versepar \ledsidenote{\passref{1.46.150}}\skt{स्वादु शीतङ्गुरुस्निग्धम्मांसम्मारुतपित्तजित् ॥ }\pend
\pstart\versepar \skt{मेध्यं }\edtext{\edtext{\skt{शूलानिलश्च्चर्द्दिकफारोचकनाशनं}}{\lemma{\skt{शूलानिलश्}}\Afootnote{H; \skt{शूलानिल}\ A1938 N}}}{\lemma{\skt{च्चर्द्दिकफारोचकनाशनं}}\Afootnote{H; \skt{च्छर्दिकफारोचकनाशनम् ॥१५०॥}\ A1938; \skt{च्छर्दिकफरोचकनाशनं ।}\ N}}\skt{ ।}\pend

\pstart \ledsidenote{\passref{1.46.151ab}}\skt{दीपनं लघु }\edtext{\skt{साङ्ग्रहि}}{\Afootnote{H; \skt{सङ्ग्राहि}\ A1938 N}}\skt{ }\edtext{\skt{गुल्मार्शोघ्नं}}{\Afootnote{\skt{गुल्मार्शोघ्नं तु}\ A1938; \skt{गुल्माषोघ्नं तु}\ N; \skt{गुल्मार्शोघ्न}\ H}}\skt{ }\edtext{\skt{केशरं}}{\Afootnote{\skt{केसरम्}\ A1938 N; \skt{न्केशरं ॥}\ H}}\skt{ ॥}\pend

\pstart \ledsidenote{\passref{1.46.152cd}}\edtext{\edtext{\edtext{\skt{पित्तमारुतकृद्वालम्पित्तलम्वद्धकेशरं}}{\lemma{\skt{पित्तम्}}\Afootnote{N H; \skt{पित्ता}\ A1938}}}{\lemma{\skt{आरुतकृद्वालम्}}\Afootnote{H; \skt{निलकरं बालं}\ A1938; \skt{आरुदकृद्वालं}\ N}}}{\lemma{\skt{वद्धकेशरं}}\Afootnote{N H; \skt{बद्धकेसरम् ॥१५२॥}\ A1938}}\skt{ ।}\pend

\pstart\versepar \ledsidenote{\passref{1.46.153}}\edtext{\skt{हृद्यम्वर्ण्णकरम्वृष्यं}}{\lemma{\skt{वृष्यं}}\Afootnote{N H; \skt{रुच्यं}\ A1938}}\skt{ }\edtext{\edtext{\skt{रुच्यम्मांसवलप्रदं}}{\lemma{\skt{रुच्यम्}}\Afootnote{N H; \skt{रक्त}\ A1938}}}{\lemma{\skt{मांसवलप्रदं}}\Afootnote{H; \skt{मांसबलप्रदम् ।}\ A1938; \skt{मासवलप्रदं ।}\ N}}\skt{ ॥ }\pend
\pstart\versepar \skt{कषायानुरसं स्वादु वातघ्नं }\edtext{\edtext{\skt{वृंहणङ्गुरुः}}{\lemma{\skt{वृंहणङ्}}\Afootnote{N H; \skt{बृंहणं}\ A1938}}}{\lemma{\skt{गुरुः}}\Afootnote{N H; \skt{गुरु ॥१५३॥}\ A1938}}\skt{ ॥}\pend

\pstart\versepar \ledsidenote{\passref{1.46.154}}\skt{पित्ताविरोधि }\edtext{\edtext{\skt{सम्पक्क्वमाम्रं}}{\lemma{\skt{सम्पक्क्वम्}}\Afootnote{H; \skt{सम्पक्वम्}\ A1938; \skt{संप}\ N}}}{\lemma{\skt{आम्रं}}\Afootnote{A1938 H; \skt{माम्ल्व}\ N}}\skt{ }\edtext{\skt{शुल्रविवर्द्धनं}}{\Afootnote{H; \skt{शुक्रविवर्धनम् ।}\ A1938; \skt{शुक्रविवर्द्ध ।}\ N}}\skt{ ॥ }\pend
\pstart\versepar \skt{}\edtext{\skt{वृंहणम्मधुरं}}{\lemma{\skt{वृंहणम्}}\Afootnote{H; \skt{बृंहणं}\ A1938; \skt{वृंहन}\ N}}\skt{ }\edtext{\skt{वल्यं}}{\Afootnote{N H; \skt{बल्यं}\ A1938}}\skt{ गुरु विष्टभ्य }\edtext{\skt{जीर्यति}}{\Afootnote{A1938 H; \skt{र्यति ।}\ N}}\skt{ ।}\pend

\pstart\versepar \ledsidenote{\passref{1.46.155}}\edtext{\skt{आम्रातकफलं}}{\Afootnote{A1938 H; \skt{अमात्रकफलं}\ N}}\skt{ वृष्यं सस्नेहं श्लेष्मवर्द्धनं ॥ }\pend
\pstart\versepar \skt{त्रिदोषविष्टम्भकरं लकुचं }\edtext{\skt{शुक्रदूषणं}}{\Afootnote{N; \skt{शुक्रनाशनम् ॥१५५॥}\ A1938; \skt{शुक्रदुषलणं ।}\ H}}\skt{ ।}\pend

\pstart\versepar \ledsidenote{\passref{1.46.156}}\edtext{\skt{अम्लम्तृषापहं}}{\lemma{\skt{अम्लम्}}\Afootnote{A1938 H; \skt{अम्व्लं}\ N}}\skt{ }\edtext{\skt{रुच्यम्पित्तकृत्करमर्द्दकं}}{\lemma{\skt{पित्तकृत्करमर्द्दकं}}\Afootnote{A1938 H; \skt{पित्तं कृत्करमर्दकं ॥}\ N}}\skt{ ॥ }\pend
\pstart\versepar \skt{वातपित्तहरं }\edtext{\skt{वृष्यम्पिपालं}}{\lemma{\skt{पिपालं}}\Afootnote{H; \skt{प्रियालं}\ A1938; \skt{प्रिपालं}\ N}}\skt{ गुरु शीतलं ।}\pend

\pstart\versepar \ledsidenote{\passref{1.46.157}}\skt{हृद्यं स्वादु }\edtext{\skt{कषायाम्लम्भव्यमास्यविशोधनं}}{\lemma{\skt{कषायाम्लम्}}\Afootnote{A1938 H; \skt{कषायाम्ल्वं}\ N}}\skt{ ॥ }\pend
\pstart\versepar \skt{पित्तश्लेष्महरम्ग्राहि गुरु विष्टम्भि शीतलं ।}\pend

\pstart\versepar \ledsidenote{\passref{1.46.158}}\skt{पारावतं समधुरं रुच्यमत्यग्निवातनुत् ॥ }\pend
\pstart\versepar \skt{}\edtext{\skt{गरदोषहरन्नीपम्प्राचीनामलकन्तथा}}{\lemma{\skt{प्राचीनामलकन्}}\Afootnote{A1938 H; \skt{प्राचीअमलकन्}\ N}}\skt{ ।}\pend

\pstart\versepar \ledsidenote{\passref{1.46.159}}\edtext{\skt{वातापहन्तिन्तिडीकमामं}}{\lemma{\skt{तिन्तिडीकमामं}}\Afootnote{A1938 H; \skt{तिन्तिडीकमामं}\ N}}\skt{ पित्तवलासकृत् ॥ }\pend
\pstart\versepar \skt{}\edtext{\skt{ग्राह्युष्णन्दीपनं}}{\lemma{\skt{दीपनं}}\Afootnote{A1938 H; \skt{दीपणं}\ N}}\skt{ रुच्यं }\edtext{\skt{सम्पक्क्वं}}{\Afootnote{H; \skt{सम्पक्वं}\ A1938 N}}\skt{ कफवातनुत् ।}\edtext{}{\lemma{\skt{पित्तवलासकृत् ॥ ग्राह्युष्णन्}}\Afootnote{H; \skt{पित्तबलासकृत् । ग्राह्युष्णं}\ A1938; \skt{पित्तवलासकृत् ॥ ग्राह्युष्णं}\ N}}\pend

\pstart \ledsidenote{\passref{1.46.160cd}}\edtext{\skt{अम्लीकाया}}{\Afootnote{H; \skt{अम्लीकायाः}\ A1938; \skt{अम्व्लीकायाः}\ N}}\skt{ }\edtext{\edtext{\skt{फलम्पक्क्वन्तद्वद्भेदि}}{\lemma{\skt{पक्क्वन्}}\Afootnote{H; \skt{पक्वं}\ A1938 N}}\skt{ तु}}{\lemma{\skt{तद्वद्भेदि तु}}\Afootnote{A1938 H; \skt{तद्वद्भेदि तु}\ N}}\skt{ केवलं ॥}\pend

\pstart\versepar \ledsidenote{\passref{1.46.161}}\edtext{\skt{अम्लं}}{\Afootnote{A1938 H; \skt{अम्व्लं}\ N}}\skt{ समधुरं }\edtext{\skt{हृद्यम्विषदम्भक्तरोचनं}}{\lemma{\skt{विषदम्}}\Afootnote{H; \skt{विशदं}\ A1938; \skt{विषडं}\ N}}\skt{ । }\pend
\pstart\versepar \skt{}\edtext{\edtext{\skt{वातघ्नन्दुर्जरञ्चोक्तन्नागरङ्गफलं}}{\lemma{\skt{चोक्तन्}}\Afootnote{N H; \skt{प्रोक्तं}\ A1938}}}{\lemma{\skt{नागरङ्गफलं}}\Afootnote{H; \skt{नारङ्गस्य फलं}\ A1938; \skt{नागरं फलं}\ N}}\skt{ }\edtext{\skt{गुरुः}}{\Afootnote{N H; \skt{गुरु ॥१६१॥}\ A1938}}\skt{ ॥}\pend

\pstart \ledsidenote{\passref{1.46.163}}\edtext{\skt{क्षीरवृक्षफलजाम्वूराजादनतोदनन्दुकवकुलधन्वश्मन्तकफल्गुफरूषकचाङ्गेरुकपुष्करवर्त्तिबिल्वप्रभृतीन्येतानि}}{\Afootnote{\skt{क्षीरवृक्षफलजाम्बवराजादनतोदनशीतफलतिन्दुकबकुलधन्वनाश्मन्तकाश्वकर्णफल्गु- परूषकगाङ्गेरुकीपुष्करवर्तिबिल्वबिम्बीप्रभृतीनि ॥१६३॥}\ A1938; \skt{क्षीरवृक्षफलजाम्वुगजादनतोदनतिन्दुकेवलकुलधव्न्वनाश्मन्तकफल्गुफरूषकचाङ्गेरुकपुष्करवर्त्तिविल्वप्रभृतीन्येतानि}\ N; \skt{क्षीरवृक्षफलजाम्वूराजादनतोदनन्दुकवकुलधन्वश्मन्तकफल्गुफरूषकचाङ्गेरुकपुष्करवर्त्तिविल्वप्रभृतीन्येतानि}\ H}}\skt{ }\edtext{\skt{शीतानि}}{\Afootnote{N H; \emph{om.}\ A1938}}\skt{ }\edtext{\skt{कफपित्तहराणि}}{\Afootnote{N H; \emph{om.}\ A1938}}\skt{ }\edtext{\skt{च}}{\Afootnote{N H; \emph{om.}\ A1938}}\skt{ ॥}\pend

\pstart \ledsidenote{\passref{1.46.164cd}}\edtext{\skt{सांग्राहिकानि}}{\Afootnote{N H; \skt{सङ्ग्राहकाणि}\ A1938}}\skt{ }\edtext{\skt{रूक्षाणि}}{\Afootnote{A1938 H; \skt{रुक्षाणि}\ N}}\skt{ }\edtext{\skt{कयमधुराणि}}{\Afootnote{H; \skt{कषायमधुराणि}\ A1938 N}}\skt{ च ।}\pend

\pstart\versepar \ledsidenote{\passref{1.46.165}}\skt{क्षीरवृक्षफलन्तेषां गुरु विष्टम्भि शीतलं ॥ }\pend
\pstart\versepar \skt{}\edtext{\skt{कषायमधुरं}}{\Afootnote{H; \skt{कषायं मधुरं}\ A1938; \skt{कषायमधुर}\ N}}\skt{ }\edtext{\skt{साम्लन्नातिमारुतकोपनं}}{\lemma{\skt{साम्लन्}}\Afootnote{A1938 H; \skt{साम्व्लं}\ N}}\skt{ ।}\pend

\pstart\versepar \ledsidenote{\passref{1.46.166}}\skt{अत्यर्थम्वातलं ग्राहि }\edtext{\skt{जाम्ववङ्कफपित्तजित्}}{\lemma{\skt{जाम्ववङ्}}\Afootnote{H; \skt{जाम्बवं}\ A1938; \skt{जाम्वं}\ N}}\skt{ ॥ }\pend
\pstart\versepar \skt{स्निग्धस्वादुकषायञ्च }\edtext{\skt{राकादनफलं}}{\Afootnote{\skt{राजादनफलं}\ A1938 N H}}\skt{ }\edtext{\skt{गुरुः}}{\Afootnote{N H; \skt{गुरु ॥१६६॥}\ A1938}}\skt{ ।}\edtext{}{\lemma{\skt{कफपित्तजित् ॥ स्निग्धस्वादुकषायञ्च}}\Afootnote{H; \skt{कफपित्तजित्स्निग्धं स्वादु कषायं च}\ A1938 N}}\pend

\pstart\versepar \ledsidenote{\passref{1.46.167}}\edtext{\skt{कषायमधुरं}}{\Afootnote{N H; \skt{कषायं मधुरं}\ A1938}}\skt{ }\edtext{\skt{प्रोक्तं}}{\Afootnote{N H; \skt{रूक्षं}\ A1938}}\skt{ तोदनं कफवातजित् ॥ }\pend
\pstart\versepar \skt{अम्लोष्णं लघु }\edtext{\skt{सांग्राहि}}{\Afootnote{H; \skt{सङ्ग्राहि}\ A1938 N}}\skt{ स्निग्धं पित्ताग्निवर्द्धनं ।}\edtext{}{\lemma{\skt{कफवातजित् ॥ अम्लोष्णं}}\Afootnote{A1938 H; \skt{कफवातजित् ॥ अम्व्लोष्णं}\ N}}\pend

\pstart\versepar \ledsidenote{\passref{1.46.168}}\skt{आमं }\edtext{\skt{कषायं}}{\Afootnote{A1938 H; \skt{कषायां}\ N}}\skt{ }\edtext{\skt{सांग्राहि}}{\Afootnote{H; \skt{सङ्ग्राहि}\ A1938 N}}\skt{ तिन्दुकम्वातकोपनं ॥ }\pend
\pstart\versepar \skt{विपाके गुरु }\edtext{\edtext{\skt{सम्पक्क्वम्मधुरङ्कफपित्तजित्}}{\lemma{\skt{सम्पक्क्वम्}}\Afootnote{H; \skt{सम्पक्वं}\ A1938 N}}}{\lemma{\skt{मधुरङ्कफपित्तजित्}}\Afootnote{A1938 H; \skt{मधुरं कफपित्तचित् ।}\ N}}\skt{ ।}\pend

\pstart\versepar \ledsidenote{\passref{1.46.169}}\skt{मधुरञ्च कषायञ्च }\edtext{\skt{स्निग्धङ्ग्राहि च}}{\lemma{\skt{ग्राहि च}}\Afootnote{N H; \skt{सङ्ग्राहि}\ A1938}}\skt{ }\edtext{\skt{वाकुलं}}{\Afootnote{N H; \skt{बाकुलम् ।}\ A1938}}\skt{ ॥ }\pend
\pstart\versepar \skt{स्थिरीकरञ्च }\edtext{\edtext{\edtext{\skt{दन्तानाम्विषदम्फलमुच्यते}}{\lemma{\skt{विषदम्}}\Afootnote{H; \skt{विशदं}\ A1938 N}}}{\lemma{\skt{फलम्}}\Afootnote{N H; \skt{फल}\ A1938}}}{\lemma{\skt{उच्यते}}\Afootnote{N H; \skt{मुच्यते ॥१६९॥}\ A1938}}\skt{ ।}\pend

\pstart\versepar \ledsidenote{\passref{1.46.170}}\skt{सकषायं हिमं स्वादु }\edtext{\skt{धान्वनङ्कफपित्तजित्}}{\lemma{\skt{धान्वनङ्}}\Afootnote{A1938 H; \skt{धान्वानं}\ N}}\skt{ ॥ }\pend
\pstart\versepar \skt{}\edtext{\skt{तद्वद्गाङ्गेरुकम्विद्यादश्मन्तकफलानि}}{\lemma{\skt{विद्यादश्मन्तकफलानि}}\Afootnote{A1938; \skt{विंद्यादश्मन्तकफलानि}\ N; \skt{विद्या अश्मन्तकफलानि}\ H}}\skt{ च ।}\edtext{}{\lemma{\skt{कफपित्तजित् ॥ तद्वद्गाङ्गेरुकम्}}\Afootnote{H; \skt{कफवातजित् । तद्वद्गाङ्गेरुकं}\ A1938; \skt{कफपित्तजित् ॥ तद्वद्भाङ्गेरुकं}\ N}}\pend

\pstart\versepar \ledsidenote{\passref{1.46.171}}\skt{विष्टम्भि मधुरं }\edtext{\skt{स्निग्धम्फल्गु}}{\lemma{\skt{फल्गु}}\Afootnote{A1938 H; \skt{फल्गुजुः}\ N}}\skt{ }\edtext{\edtext{\skt{तर्प्पणङ्गुरुः}}{\lemma{\skt{तर्प्पणङ्}}\Afootnote{H; \skt{जं तर्पणं}\ A1938; \skt{संतर्पणं}\ N}}}{\lemma{\skt{गुरुः}}\Afootnote{N H; \skt{गुरु ।}\ A1938}}\skt{ ॥ }\pend
\pstart\versepar \skt{}\edtext{\skt{अत्यम्लमीषत्मधुरङ्कषायानुरसं}}{\lemma{\skt{अत्यम्लमीषत्मधुरङ्}}\Afootnote{H; \skt{अत्यम्लमीषन्मधुरं}\ A1938; \skt{अत्यम्व्लमीषन्मधुरं}\ N}}\skt{ }\edtext{\skt{लघुः}}{\Afootnote{N H; \skt{लघु ॥१७१॥}\ A1938}}\skt{ ।}\pend

\pstart\versepar \ledsidenote{\passref{1.46.172}}\edtext{\edtext{\edtext{\skt{वातघ्नम्पित्तजनमामम्विद्यात्फरूषकं}}{\lemma{\skt{पित्तजनम्}}\Afootnote{N H; \skt{पित्तजन}\ A1938}}}{\lemma{\skt{आमम्}}\Afootnote{N H; \skt{मामं}\ A1938}}}{\lemma{\skt{विद्यात्फरूषकं}}\Afootnote{H; \skt{विद्यात्परुषकम् ।}\ A1938; \skt{विद्यात्पकं ॥}\ N}}\skt{ ॥ }\pend
\pstart\versepar \skt{तदेव }\edtext{\edtext{\skt{पक्क्वम्मधुरम्वातपित्तनिवर्हणं}}{\lemma{\skt{पक्क्वम्}}\Afootnote{H; \skt{पक्वं}\ A1938 N}}}{\lemma{\skt{वातपित्तनिवर्हणं}}\Afootnote{H; \skt{वातपित्तनिबर्हणम् ॥१७२॥}\ A1938; \skt{वातपित्तनिवर्हणां ।}\ N}}\skt{ ॥}\pend

\pstart\versepar \ledsidenote{\passref{1.46.173}}\skt{विपाके मधुरं शीतं }\edtext{\skt{रक्तप्रसादनं}}{\Afootnote{H; \skt{रक्तपित्तप्रसादनम्}\ A1938 N}}\skt{ ॥ }\pend
\pstart\versepar \skt{पौष्करं स्वादु विष्टम्भि }\edtext{\skt{वल्यं}}{\Afootnote{N H; \skt{बल्यं}\ A1938}}\skt{ कफकरं फलं ।}\pend

\pstart\versepar \ledsidenote{\passref{1.46.174}}\skt{कफानिलहरन्तीक्ष्णं }\edtext{\skt{स्निग्धं}}{\Afootnote{A1938 H; \skt{स्निग्धां}\ N}}\skt{ }\edtext{\skt{साङ्ग्राहि}}{\Afootnote{N H; \skt{सङ्ग्राहि}\ A1938}}\skt{ दीपनं ॥ }\pend
\pstart\versepar \skt{}\edtext{\edtext{\edtext{\skt{कटुतिक्तकषायोष्णम्वालम्बिल्वमुदाहृतं}}{\lemma{\skt{वालम्}}\Afootnote{N H; \skt{बालं}\ A1938}}}{\lemma{\skt{बिल्वम्}}\Afootnote{\skt{बिल्व}\ A1938; \skt{विल्वम्}\ N H}}}{\lemma{\skt{उदाहृतं}}\Afootnote{N H; \skt{मुदाहृतम् ॥१७४॥}\ A1938}}\skt{ ।}\pend

\pstart\versepar \ledsidenote{\passref{1.46.175}}\edtext{\skt{तदेव}}{\Afootnote{N H; \skt{विद्यात्तदेव}\ A1938}}\skt{ }\edtext{\edtext{\edtext{\skt{विद्यात्सम्पक्वम्मधुरानुरसङ्गुरुः}}{\lemma{\skt{विद्यात्सम्पक्वम्}}\Afootnote{N H; \skt{सम्पक्वं}\ A1938}}}{\lemma{\skt{मधुरानुरसङ्}}\Afootnote{A1938 H; \skt{मधुरानुरस}\ N}}}{\lemma{\skt{गुरुः}}\Afootnote{N H; \skt{गुरु ।}\ A1938}}\skt{ ॥ }\pend
\pstart\versepar \skt{विदाहि विष्टम्भकरन्दोषकृत्पूतिमारुतं ॥}\pend

\pstart \ledsidenote{\passref{1.46.177}}\skt{तालनालिकेरपनसमोचप्रभृतीनि ॥}\pend

\pstart\versepar \ledsidenote{\passref{1.46.178}}\skt{स्वादुपाकरसान्याहुर्व्वातपित्तहराणि च । }\pend
\pstart\versepar \skt{}\edtext{\skt{वलप्रदानि}}{\Afootnote{N H; \skt{बलप्रदानि}\ A1938}}\skt{ }\edtext{\skt{स्निग्धानि}}{\Afootnote{A1938 H; \skt{स्ग्निग्धानि}\ N}}\skt{ }\edtext{\skt{वृंहणानि}}{\Afootnote{N H; \skt{बृंहणानि}\ A1938}}\skt{ हिमानि च ॥}\pend

\pstart\versepar \ledsidenote{\passref{1.46.179}}\skt{फलं स्वादुरसम्तेषान्तालजं गुरुपित्तजित् । }\pend
\pstart\versepar \skt{तद्वीजं स्वादुपाकन्तु }\edtext{\skt{मूत्रलङ्कफपित्तजित्}}{\lemma{\skt{कफपित्तजित्}}\Afootnote{N H; \skt{वातपित्तजित् ॥१७९॥}\ A1938}}\skt{ ॥}\edtext{}{\lemma{\skt{गुरुपित्तजित् । तद्वीजं}}\Afootnote{N H; \skt{गुरु पित्तजित् । तद्बीजं}\ A1938}}\pend

\pstart\versepar \ledsidenote{\passref{1.46.180}}\skt{नालिकेरं गुरु }\edtext{\skt{स्निग्धं}}{\Afootnote{A1938 H; \skt{स्ग्निधं}\ N}}\skt{ पित्तघ्नं स्वादु शीतलं । }\pend
\pstart\versepar \skt{}\edtext{\skt{वलमांसप्रदं}}{\Afootnote{N H; \skt{बलमांसप्रदं}\ A1938}}\skt{ }\edtext{\edtext{\skt{हृद्यम्वृंहणम्बस्तिशोधनं}}{\lemma{\skt{वृंहणम्}}\Afootnote{N H; \skt{बृंहणं}\ A1938}}}{\lemma{\skt{बस्तिशोधनं}}\Afootnote{A1938; \skt{वस्तिशोधन ।}\ N; \skt{वस्तिशोधनं ॥}\ H}}\skt{ ॥}\pend

\pstart\versepar \ledsidenote{\passref{1.46.181}}\skt{पनसं सकषायम्तु स्निग्धं स्वादु }\edtext{\skt{हिमं}}{\Afootnote{N H; \skt{रसं}\ A1938}}\skt{ }\edtext{\skt{गुरु}}{\Afootnote{A1938 H; \skt{गुरुः ।}\ N}}\skt{ । }\pend
\pstart\versepar \skt{मौचंस्वादुरसम्प्रोक्तं }\edtext{\edtext{\skt{कषायान्नातितलं}}{\lemma{\skt{कषायान्}}\Afootnote{H; \skt{कषायं}\ A1938 N}}}{\lemma{\skt{नातितलं}}\Afootnote{H; \skt{नातिशीतलम्}\ A1938 N}}\skt{ ॥ }\pend
\pstart\versepar \skt{रक्तपित्तहरम्वृष्यं }\edtext{\skt{रुच्यं}}{\Afootnote{A1938 H; \emph{om.}\ N}}\skt{ }\edtext{\skt{श्लेष्मकरङ्गुरुः}}{\lemma{\skt{गुरुः}}\Afootnote{N H; \skt{गुरु ॥१८१॥}\ A1938}}\skt{ ॥}\pend

\pstart \ledsidenote{\passref{1.46.182}}\edtext{\skt{द्राक्षाकाश्मर्यमधूकपुष्पखर्जूरप्रभृतीनि}}{\Afootnote{A1938 H; \skt{दाक्षाकाश्मर्यमधूकपुष्पखर्जुरप्रभृतीनि}\ N}}\skt{ ॥}\pend

\pstart\versepar \ledsidenote{\passref{1.46.183}}\edtext{\edtext{\edtext{\skt{रक्तपित्तहराण्याहुर्गुरूणि}}{\lemma{\skt{रक्तपित्तहराण्य्}}\Afootnote{H; \skt{रक्तपित्तहरा}\ A1938 N}}}{\lemma{\skt{आहुर्}}\Afootnote{H; \skt{ण्याहुर्}\ A1938; \skt{ण्या}\ N}}}{\lemma{\skt{गुरूणि}}\Afootnote{A1938 H; \emph{om.}\ N}}\skt{ }\edtext{\skt{च}}{\Afootnote{H; \skt{मधुराणि च}\ A1938; \skt{हुगुरुणि च ।}\ N}}\skt{ ॥ }\pend
\pstart\versepar \skt{तेषान्द्राक्षा सरा स्वर्या }\edtext{\skt{मधुरा}}{\Afootnote{A1938 H; \skt{मधुराः}\ N}}\skt{ }\edtext{\skt{स्निग्धशीतला}}{\Afootnote{A1938 H; \skt{स्निग्धशीतलाः ।}\ N}}\skt{ ।}\pend

\pstart\versepar \ledsidenote{\passref{1.46.184}}\skt{रक्तपित्तज्वरश्वासतृष्णादाहक्षयापहा ॥ }\pend
\pstart\versepar \skt{हृद्यं }\edtext{\edtext{\skt{मूत्रविवन्धघ्नम्पित्तासृग्दाहनाशनं}}{\lemma{\skt{मूत्रविवन्धघ्नम्}}\Afootnote{N H; \skt{मूत्रविबन्धघ्नं}\ A1938}}}{\lemma{\skt{पित्तासृग्दाहनाशनं}}\Afootnote{N H; \skt{पित्तासृग्वातनाशनम् ॥१८४॥}\ A1938}}\skt{ ।}\pend

\pstart\versepar \ledsidenote{\passref{1.46.185}}\skt{केश्यं }\edtext{\edtext{\skt{रसायनम्मेध्यङ्काश्मर्यम्फलमुच्यते}}{\lemma{\skt{फलम्}}\Afootnote{N H; \skt{फल}\ A1938}}}{\lemma{\skt{उच्यते}}\Afootnote{N H; \skt{मुच्यते ।}\ A1938}}\skt{ ॥ }\pend
\pstart\versepar \skt{}\edtext{\skt{क्षतक्षयापहं}}{\Afootnote{A1938 H; \skt{क्षतक्षयायापहं}\ N}}\skt{ }\edtext{\edtext{\skt{ह्र्̥द्यम्वृंहणं}}{\lemma{\skt{ह्र्̥द्यम्}}\Afootnote{H; \skt{हृद्यं}\ A1938 N}}}{\lemma{\skt{वृंहणं}}\Afootnote{N H; \emph{om.}\ A1938}}\skt{ शीतलं }\edtext{\skt{गुरु}}{\Afootnote{H; \skt{तर्पणं गुरु ॥१८५॥}\ A1938; \skt{गुरुः ।}\ N}}\skt{ ।}\pend

\pstart\versepar \ledsidenote{\passref{1.46.186}}\edtext{\skt{वृष्यं}}{\Afootnote{N H; \skt{र}\ A1938}}\skt{ }\edtext{\skt{स्निग्धं}}{\Afootnote{H; \skt{से पाके}\ A1938; \skt{स्निग्ध}\ N}}\skt{ }\edtext{\skt{समधुरं खा}}{\Afootnote{H; \skt{च मधुरं खा}\ A1938; \skt{समधुरं ख}\ N}}\skt{ }\edtext{\skt{र्ज्जुरं}}{\Afootnote{H; \skt{र्जूरं}\ A1938 N}}\skt{ रक्तपित्तहृत् ॥ }\pend
\pstart\versepar \skt{वृंहणीयहृद्यञ्च }\edtext{\skt{मधूककुसुमङ्गुरु}}{\lemma{\skt{गुरु}}\Afootnote{A1938 H; \skt{गुरुः}\ N}}\skt{ । }\pend
\pstart\versepar \skt{}\edtext{\skt{वातपित्तोपशमनम्फलन्तस्योपदिश्यते}}{\lemma{\skt{फलन्}}\Afootnote{A1938 H; \skt{फल}\ N}}\skt{ ॥}\edtext{}{\lemma{\skt{रक्तपित्तहृत् ॥ वृंहणीयहृद्यञ्च}}\Afootnote{H; \skt{रक्तपित्तनुत् । बृंहणीयमहृद्यं च}\ A1938; \skt{रक्तपित्तत् । वृंहणीयं अहृद्यञ्च}\ N}}\pend

\pstart \ledsidenote{\passref{1.46.187}}\edtext{\skt{वातमाक्षोडातिकमुकुलनिचुलप्रभृतीनि}}{\Afootnote{H; \skt{वातामाक्षोडाभिषुकनिचुलपिचुनिकोचकोरुमाणप्रभृतीनि ॥१८७॥}\ A1938; \skt{वातामाक्षोडभिषुकनिचूलप्रभृतीनि ।}\ N}}\skt{ ॥}\pend

\pstart\versepar \ledsidenote{\passref{1.46.188}}\skt{पित्तश्लेष्महराण्याहुः }\edtext{\skt{स्निग्धोष्णानि}}{\Afootnote{A1938 N H\textsuperscript{pc}; \skt{स्निग्धोनि}\ H\textsuperscript{ac}}}\skt{ }\edtext{\skt{गुरूणि च}}{\Afootnote{A1938 H; \skt{गुरुणि च ।}\ N}}\skt{ । }\pend
\pstart\versepar \skt{}\edtext{\skt{वृंहणान्यनिलघ्नानि}}{\Afootnote{N H; \skt{बृंहणान्यनिलघ्नानि}\ A1938}}\skt{ }\edtext{\skt{वल्यानि}}{\Afootnote{N H; \skt{बल्यानि}\ A1938}}\skt{ मधुराणि च ॥}\pend

\pstart\versepar \ledsidenote{\passref{1.46.189}}\skt{कषायङ्कफपित्तघ्नं }\edtext{\skt{किञ्चित्तिक्तं}}{\Afootnote{H; \skt{किञ्चित्तिक्तं}\ A1938 N}}\skt{ रुचिप्रदं । }\pend
\pstart\versepar \skt{हृद्यं सुगंधि }\edtext{\skt{मधुरं}}{\Afootnote{N H; \skt{विशदं}\ A1938}}\skt{ }\edtext{\edtext{\skt{लवलीफलमुच्यते}}{\lemma{\skt{लवलीफलम्}}\Afootnote{N H; \skt{लवलीफल}\ A1938}}}{\lemma{\skt{उच्यते}}\Afootnote{N H; \skt{मुच्यते ॥१८९॥}\ A1938}}\skt{ ॥}\pend

\pstart\versepar \ledsidenote{\passref{1.46.190}}\skt{वसिरं }\edtext{\skt{शीतपाकं तु}}{\Afootnote{N H; \skt{शीतपाक्यं च}\ A1938}}\skt{ }\edtext{\skt{सारुः}}{\Afootnote{N H; \skt{सारु}\ A1938}}\skt{ }\edtext{\skt{करनिवन्धनं}}{\Afootnote{N H; \skt{ष्करनिबन्धनम् ।}\ A1938}}\skt{ । }\pend
\pstart\versepar \skt{विष्टम्भि }\edtext{\skt{शीतं}}{\Afootnote{N H; \skt{दुर्जरं}\ A1938}}\skt{ }\edtext{\skt{रूक्षङ्च}}{\Afootnote{N H; \skt{रूक्षं शी}\ A1938}}\skt{ }\edtext{\skt{वातपित्तप्रकोपनं}}{\Afootnote{N H; \skt{तलं वातकोपनम् ॥१९०॥}\ A1938}}\skt{ ॥}\pend

\pstart\versepar \ledsidenote{\passref{1.46.191}}\skt{विपाके मधुरञ्चापि }\edtext{\skt{रक्तपित्तप्र क्त}}{\Afootnote{H; \skt{रक्तपित्त}\ A1938 N}}\skt{ प्रसादनं । }\pend
\pstart\versepar \skt{}\edtext{\edtext{\edtext{\skt{ऐरावतन्दन्तशथमम्लं}}{\lemma{\skt{ऐरावतन्}}\Afootnote{H; \skt{ऐरावतं}\ A1938; \skt{ऐवदतन्}\ N}}}{\lemma{\skt{दन्तशथम्}}\Afootnote{H; \skt{दन्तशठ}\ A1938; \skt{दन्तशठम्}\ N}}}{\lemma{\skt{अम्लं}}\Afootnote{H; \skt{मम्लं}\ A1938; \skt{अम्व्लं}\ N}}\skt{ शोणितपित्तकृत् ॥}\pend

\pstart\versepar \ledsidenote{\passref{1.46.192}}\edtext{\edtext{\skt{शीतङ्कषायमधुरण्टङ्कम्मारुतकृद्गुरु}}{\lemma{\skt{कषायमधुरण्टङ्कम्}}\Afootnote{H; \skt{कषायं मधुरं टङ्कं}\ A1938; \skt{कषायं मधुरं टङ्क}\ N}}}{\lemma{\skt{मारुतकृद्गुरु}}\Afootnote{H; \skt{मारुतकृद्गुरु ।}\ A1938; \skt{मारुतकृद्गुतुः ।}\ N}}\skt{ ॥ }\pend
\pstart\versepar \skt{}\edtext{\skt{स्निग्धोष्णन्तिक्तमधुरम्वातश्लेष्मघ्नमिङ्मुदं}}{\lemma{\skt{वातश्लेष्मघ्नमिङ्मुदं}}\Afootnote{H; \skt{वातश्लेष्मघ्नमैङ्गुदम् ॥१९२॥}\ A1938; \skt{वातश्लेष्मघ्नमिङ्गुदं ।}\ N}}\skt{ ॥}\pend

\pstart\versepar \ledsidenote{\passref{1.46.193}}\skt{शमीफलं गुरु स्वादु }\edtext{\skt{रूक्षोष्णं}}{\Afootnote{A1938 H; \skt{रुक्षोष्णं}\ N}}\skt{ केशनाशनं । }\pend
\pstart\versepar \skt{}\edtext{\edtext{\skt{गुरुश्लेष्मान्तकफलङ्कफकृन्मधुरं}}{\lemma{\skt{गुरुश्लेष्मान्तकफलङ्}}\Afootnote{H; \skt{गुरु श्लेष्मातकफलं}\ A1938; \skt{गुरु श्लेष्मातकओहलं}\ N}}}{\lemma{\skt{कफकृन्मधुरं}}\Afootnote{A1938 H; \skt{कफकृन्मधुरं}\ N}}\skt{ हिमं ॥}\pend

\pstart\versepar \ledsidenote{\passref{1.46.194}}\edtext{\skt{करीराक्षकपीलूनि}}{\Afootnote{H; \skt{करीराक्षिकपीलूनि}\ A1938 N}}\skt{ }\edtext{\skt{तृणशूण्यफलानि च}}{\Afootnote{H; \skt{तृणशून्यफलानि च ।}\ A1938 N}}\skt{ । }\pend
\pstart\versepar \skt{स्वादुतिक्तकटूष्णानि कफवातहराणि च ॥}\pend

\pstart\versepar \ledsidenote{\passref{1.46.195}}\edtext{\skt{रक्तपित्तहरम्तेषां}}{\lemma{\skt{रक्तपित्तहरम्}}\Afootnote{N H; \skt{तिक्तं पिन्त्तकरं}\ A1938}}\skt{ }\edtext{\skt{रसं}}{\Afootnote{N H; \skt{सरं}\ A1938}}\skt{ }\edtext{\skt{कटुविपाकिनः}}{\Afootnote{H; \skt{कटुविपाकि च}\ A1938 N}}\skt{ । }\pend
\pstart\versepar \skt{}\edtext{\edtext{\skt{तीक्षोष्णङ्कटूकम्पीलु}}{\lemma{\skt{तीक्षोष्णङ्}}\Afootnote{H; \skt{तीक्ष्णोष्णं}\ A1938 N}}}{\lemma{\skt{कटूकम्}}\Afootnote{H; \skt{कटुकं}\ A1938 N}}\skt{ }\edtext{\skt{सस्नेहङ्कफपित्तजित्}}{\lemma{\skt{कफपित्तजित्}}\Afootnote{N H; \skt{कफवातजित् ॥१९५॥}\ A1938}}\skt{ ॥}\pend

\pstart\versepar \ledsidenote{\passref{1.46.196}}\edtext{\edtext{\skt{अरुष्करन्तौरवकङ्कषायं}}{\lemma{\skt{अरुष्करन्}}\Afootnote{N H; \skt{आरुष्करं}\ A1938}}}{\lemma{\skt{तौरवकङ्}}\Afootnote{N H; \skt{तौवरकं}\ A1938}}\skt{ }\edtext{\skt{लघुपाकिनः}}{\Afootnote{H; \skt{कटुपाकि च ।}\ A1938; \skt{लघुपाकि च}\ N}}\skt{ । }\pend
\pstart\versepar \skt{}\edtext{\skt{उष्णक्रिमिहरो}}{\Afootnote{N H; \skt{उष्णं कृमिज्वरा}\ A1938}}\skt{ }\edtext{\skt{मेहशाफदुर्नमनाशनं}}{\Afootnote{H; \skt{नाहमेहोदावर्तनाशनम् ।}\ A1938; \skt{मेहशाफदुर्नामनाशनं ।}\ N}}\skt{ ॥ }\pend
\pstart\versepar \skt{}\edtext{\skt{कुष्ठगुल्मोदरार्शोघ्नं}}{\Afootnote{A1938 H; \skt{कुष्ठगुल्मोदरार्शोघ्नं}\ N}}\skt{ }\edtext{\skt{लघुपाकि}}{\Afootnote{H; \skt{कटुपाकि}\ A1938; \skt{लघुपाके}\ N}}\skt{ तथैव च ।}\pend

\pstart \ledsidenote{\passref{1.46.197ab}}\edtext{\skt{करञ्जकिङ्शुकारिष्टफलञ्जन्तुप्रमेहनुत्}}{\lemma{\skt{करञ्जकिङ्शुकारिष्टफलञ्}}\Afootnote{A1938 H; \skt{करञ्जकिंशुकारिष्टफल}\ N}}\skt{ ॥}\pend

\pstart\versepar \ledsidenote{\passref{1.46.198}}\edtext{\skt{रुक्षोष्णङ्कटुकम्पाके}}{\lemma{\skt{रुक्षोष्णङ्}}\Afootnote{H; \skt{रूक्षोष्णं}\ A1938; \skt{रुक्ष्णोष्णां}\ N}}\skt{ लघु वातकफापहं । }\pend
\pstart\versepar \skt{}\edtext{\edtext{\skt{तिक्तमीषद्विषहितम्विडङ्गङ्गं}}{\lemma{\skt{तिक्तमीषद्विषहितम्}}\Afootnote{A1938 H; \skt{तिक्तमीषद्विषहितं}\ N}}}{\lemma{\skt{विडङ्गङ्गं}}\Afootnote{H; \skt{विडङ्गं}\ A1938 N}}\skt{ }\edtext{\skt{क्रिमिनाशनं}}{\Afootnote{N H; \skt{कृमिनाशनम् ॥१९८॥}\ A1938}}\skt{ ॥}\pend

\pstart\versepar \ledsidenote{\passref{1.46.199}}\edtext{\edtext{\skt{व्रण्यमुष्णं}}{\lemma{\skt{व्रण्यम्}}\Afootnote{N H; \skt{व्रण्य}\ A1938}}}{\lemma{\skt{उष्णं}}\Afootnote{N H; \skt{मुष्णं}\ A1938}}\skt{ }\edtext{\skt{सनं}}{\Afootnote{H; \skt{सरं}\ A1938 N}}\skt{ }\edtext{\skt{मेध्यं}}{\Afootnote{A1938 H; \skt{मेदो}\ N}}\skt{ दोषघ्नम्मेहकुष्ठनुत् । }\pend
\pstart\versepar \skt{कषायन्दीपनं }\edtext{\edtext{\skt{साम्लञ्चक्षुष्यञ्हरीतकं}}{\lemma{\skt{साम्लञ्}}\Afootnote{H; \skt{चाम्लं}\ A1938; \skt{साम्व्लं}\ N}}}{\lemma{\skt{चक्षुष्यञ्हरीतकं}}\Afootnote{H; \skt{चक्षुष्यं चाभयाफलम् ॥१९९॥}\ A1938; \skt{चक्षुष्यञ्च हरीतकं ॥}\ N}}\skt{ ॥}\edtext{}{\lemma{\skt{मेहकुष्ठनुत् । कषायन्दीपनं}}\Afootnote{N H; \skt{शोफकुष्ठनुत् । कषायं दीपनं}\ A1938}}\pend

\pstart\versepar \ledsidenote{\passref{1.46.200}}\edtext{\edtext{\skt{भेदनङ्कटुरूक्षोष्णम्वैवस्वर्यक्रिमिशोधनं}}{\lemma{\skt{कटुरूक्षोष्णम्}}\Afootnote{H; \skt{लघु रूक्षोष्णं}\ A1938; \skt{कटु रुक्षोष्णं}\ N}}}{\lemma{\skt{वैवस्वर्यक्रिमिशोधनं}}\Afootnote{H; \skt{वैस्वर्यक्रिमिनाशनम् ।}\ A1938; \skt{वैस्वर्यं क्रिमिशोधनं ।}\ N}}\skt{ । }\pend
\pstart\versepar \skt{}\edtext{\skt{चक्षुषां}}{\Afootnote{H; \skt{चक्षुष्यं}\ A1938 N}}\skt{ }\edtext{\skt{लघुपाकाक्षं}}{\Afootnote{H; \skt{स्वादुपाक्याक्षं}\ A1938; \skt{लघुपाकक्ष}\ N}}\skt{ कषायं कफपित्तजित् ॥}\pend

\pstart\versepar \ledsidenote{\passref{1.46.201}}\skt{कफपित्तहरं रूक्षं }\edtext{\skt{वक्त्रक्लेदमलापहं}}{\Afootnote{A1938 H; \skt{वक्रक्लेदमालापहं ।}\ N}}\skt{ । }\pend
\pstart\versepar \skt{}\edtext{\edtext{\skt{कषायमीषत्मधुरं}}{\lemma{\skt{कषायम्}}\Afootnote{N H; \skt{कषाय}\ A1938}}}{\lemma{\skt{ईषत्मधुरं}}\Afootnote{N H; \skt{मीषन्मधुरं}\ A1938}}\skt{ किञ्चित्पूगफलं सरं ॥}\pend

\pstart\versepar \ledsidenote{\passref{1.46.202}}\edtext{\skt{जातीकोशोथ}}{\Afootnote{H; \skt{जातीकोशोऽथ}\ A1938; \skt{जातिकोषोथ}\ N}}\skt{ कर्पूरं जातीकटुकयोः फलं । }\pend
\pstart\versepar \skt{}\edtext{\skt{कक्कोलकं}}{\Afootnote{N H; \skt{कक्को(ङ्को)लकं}\ A1938}}\skt{ लवङ्गञ्च तिक्तङ्कटु कफापहं ॥}\pend

\pstart \ledsidenote{\passref{1.46.203ab}}\skt{लघु }\edtext{\skt{तृष्णापहं}}{\Afootnote{A1938 H; \skt{तृषापहं}\ N}}\skt{ }\edtext{\skt{वक्त्रक्लेददौर्गन्ध्यनाशनं}}{\Afootnote{A1938 H; \skt{वक्रक्लेदौर्गन्ध्यनाशनां ।}\ N}}\skt{ ।}\pend

\pstart\versepar \ledsidenote{\passref{1.46.205}}\edtext{\skt{पियालमज्जा}}{\Afootnote{H; \skt{प्रियालमज्जा}\ A1938; \skt{पिपालमज्जा}\ N}}\skt{ }\edtext{\skt{मधुरा}}{\Afootnote{H; \skt{मधुरो}\ A1938 N}}\skt{ }\edtext{\skt{वृष्या}}{\Afootnote{N H; \skt{वृष्यः}\ A1938}}\skt{ }\edtext{\skt{पित्तानिलापहा}}{\Afootnote{N H; \skt{पित्तानिलापहः ।}\ A1938}}\skt{ ॥ }\pend
\pstart\versepar \skt{}\edtext{\skt{बैभीतकी}}{\Afootnote{H; \skt{बैभीतको}\ A1938; \skt{वैभीतकी}\ N}}\skt{ }\edtext{\skt{मदकरी}}{\Afootnote{N H; \skt{मदकरः}\ A1938}}\skt{ }\edtext{\skt{कफमारुतनाशनी}}{\Afootnote{N H; \skt{कफमारुतनाशनः ॥२०५॥}\ A1938}}\skt{ ।}\pend

\pstart\versepar \ledsidenote{\passref{1.46.206}}\edtext{\skt{कषायमधुरा}}{\Afootnote{H; \skt{कषायमधुरो}\ A1938 N}}\skt{ मज्जा }\edtext{\skt{कोलानाम्पित्तनाशनी}}{\lemma{\skt{पित्तनाशनी}}\Afootnote{N H; \skt{पित्तनाशनः ।}\ A1938}}\skt{ ॥ }\pend
\pstart\versepar \skt{}\edtext{\skt{तृष्णाच्छर्द्यनिलघ्ना च}}{\Afootnote{H; \skt{तृष्णाच्छर्द्यनिलघ्नश्च}\ A1938; \skt{तृष्णाच्छर्द्यनिलघ्ना च}\ N}}\skt{ }\edtext{\skt{तद्वदामलकस्य}}{\Afootnote{A1938 H; \skt{तद्वदामलकस्य}\ N}}\skt{ च ।}\pend

\pstart\versepar \ledsidenote{\passref{1.46.207}}\edtext{\skt{बीजपूरकसम्याके}}{\Afootnote{H; \skt{बीजपूरकशम्याक}\ A1938; \skt{वीजपूरकशम्याक}\ N}}\skt{ मज्जा }\edtext{\skt{कोशांम्रसम्भवाः}}{\Afootnote{H; \skt{कोशाम्रसम्भवः ।}\ A1938; \skt{कोशाम्रसम्भवाः ॥}\ N}}\skt{ ॥ }\pend
\pstart\versepar \skt{}\edtext{\skt{स्वादुपाकाग्निबलदा}}{\Afootnote{H; \skt{स्वादुपाकोऽग्निबलदः}\ A1938; \skt{स्वादुपाकोग्निवलदा}\ N}}\skt{ }\edtext{\skt{स्निग्धा}}{\Afootnote{N H; \skt{स्निग्धः}\ A1938}}\skt{ }\edtext{\skt{पित्तानिलापहा}}{\Afootnote{N H; \skt{पित्तानिलापहः ॥२०७॥}\ A1938}}\skt{ ।}\pend

\pstart\versepar \ledsidenote{\passref{1.46.208}}\skt{यस्य यस्य }\edtext{\skt{फलस्येह}}{\Afootnote{A1938 H; \skt{फलस्येहं}\ N}}\skt{ वीर्यम्भवति यादृशं ॥ }\pend
\pstart\versepar \skt{तस्य तस्यैव वीर्येण }\edtext{\edtext{\skt{मज्जानामपि}}{\lemma{\skt{मज्जानाम्}}\Afootnote{H; \skt{मज्जानम्}\ A1938; \skt{मज्जाना}\ N}}}{\lemma{\skt{अपि}}\Afootnote{A1938 H; \skt{मेपि}\ N}}\skt{ निर्द्दिशेत् ।}\pend

\pstart\versepar \ledsidenote{\passref{1.46.209abcd}}\skt{फलेषु }\edtext{\skt{परिपक्क्वं}}{\Afootnote{H; \skt{परिपक्वं}\ A1938; \skt{पक्वं}\ N}}\skt{ यद्गुणवत्तदुदाहृतं ॥ }\pend
\pstart\versepar \skt{}\edtext{\skt{बिल्वादन्यत्र}}{\Afootnote{A1938 H; \skt{विल्वादन्यत्र}\ N}}\skt{ }\edtext{\skt{विज्ञेयमाममेतद्गुणोत्तरं}}{\lemma{\skt{एतद्गुणोत्तरं}}\Afootnote{N H; \skt{तद्धि गुणोत्तरम् ।}\ A1938}}\skt{ ।}\pend

\pstart\versepar \ledsidenote{\passref{1.46.210}}\edtext{\skt{व्याधितङ्क्रिमिदुष्टञ्च}}{\lemma{\skt{क्रिमिदुष्टञ्च}}\Afootnote{N H; \skt{कृमिजुष्टं च}\ A1938}}\skt{ }\edtext{\edtext{\skt{पाकातीतमदेशजं}}{\lemma{\skt{पाकातीतम्}}\Afootnote{N H; \skt{पाकातीत}\ A1938}}}{\lemma{\skt{अदेशजं}}\Afootnote{H; \skt{मकालजम् ।}\ A1938; \skt{अदेशज ॥}\ N}}\skt{ ॥ }\pend
\pstart\versepar \skt{वर्ज्जनीयं }\edtext{\edtext{\edtext{\skt{हितत्सर्व्वमपर्यागतमेव}}{\lemma{\skt{हितत्सर्व्वम्}}\Afootnote{H; \skt{फलं सर्व}\ A1938; \skt{हितं सर्व्वम्}\ N}}}{\lemma{\skt{अपर्यागतम्}}\Afootnote{N H; \skt{मपर्यागतम्}\ A1938}}\skt{ च}}{\lemma{\skt{एव च}}\Afootnote{N H; \skt{एव च ॥२१०॥ इति फलवर्गः ।}\ A1938}}\skt{ ॥}\pend

\pstart \ledsidenote{\passref{1.46.210.add1}}\edtext{\skt{विदारीकन्दशतावरीविशमृणालशृङ्गाटककशेरुपिण्डालुमध्वालुहम्त्या}}{\Afootnote{H; \skt{विदारीकन्दशभावरीविशमृणालशृङ्गाटलकशेरुपिण्डालुमध्वालुह}\ N}}\skt{ लुकप्रभृतीनि ॥}\pend

\pstart\versepar \ledsidenote{\passref{1.46.210.add2}}\edtext{\skt{रक्तपित्तहराण्याहुः}}{\Afootnote{H; \skt{रक्तं पित्तहराण्याहुः}\ N}}\skt{ शीतानि मधुराणि च । }\pend
\pstart\versepar \skt{गुरूणि }\edtext{\skt{बहुशुक्राणि}}{\Afootnote{\skt{वहुशुक्राणि}\ N H}}\skt{ स्तन्यवृद्धिकराणि च ॥}\pend

\pstart\versepar \ledsidenote{\passref{1.46.210.add3}}\skt{मधुरो }\edtext{\skt{बृंहणो}}{\Afootnote{H; \skt{वृंहणो}\ N}}\skt{ वृष्यः }\edtext{\edtext{\skt{शीतस्वर्योऽतिमूत्रलः}}{\lemma{\skt{शीतस्वर्यो}}\Afootnote{H; \skt{शीतस्वर्यो}\ N}}}{\lemma{\skt{ऽतिमूत्रलः}}\Afootnote{H; \skt{ति मूत्रलः ।}\ N}}\skt{ । }\pend
\pstart\versepar \skt{}\edtext{\skt{विदारिकन्दाओव्ल्पश्च}}{\Afootnote{H; \skt{विदारिकन्दो वल्यश्च}\ N}}\skt{ वातपित्तहरश्च सः ॥}\pend

\pstart\versepar \ledsidenote{\passref{1.46.210.add4}}\skt{वातपित्तहरावृष्या स्वादुशीताशतावरी । }\pend
\pstart\versepar \skt{महती }\edtext{\skt{सैव}}{\Afootnote{H; \skt{सौव}\ N}}\skt{ हृद्या तु }\edtext{\skt{मेषाग्निबलवर्द्धनी}}{\Afootnote{H; \skt{मेधाग्निबलवर्द्धनी ॥}\ N}}\skt{ ॥}\pend

\pstart\versepar \ledsidenote{\passref{1.46.210.add5}}\edtext{\skt{ग्रहण्य}}{\Afootnote{H; \skt{ग्रहण्य}\ N}}\skt{ }\edtext{\skt{विकारघ्नी}}{\Afootnote{H; \skt{ष्टौविकारघ्नी}\ N}}\skt{ वृष्याशीतारसायनी । }\pend
\pstart\versepar \skt{कफपित्तहरास्तिक्तास्तस्या एवाङ्कुराः स्मृताः ॥}\pend

\pstart\versepar \ledsidenote{\passref{1.46.210.add6}}\skt{अविदाहिरसम्प्रोक्तं रक्तपित्तप्रसादनं । }\pend
\pstart\versepar \skt{विष्टम्भिमधुरं रूक्षं दुर्ज्जरम्वातकोपनं ॥}\pend

\pstart\versepar \ledsidenote{\passref{1.46.210.add7}}\edtext{\skt{गुरुविष्टम्भिशीतौ}}{\Afootnote{H; \skt{गुरु विष्टम्भि शीतौ तु}\ N}}\skt{ }\edtext{\skt{शृङ्गाटककशेरुकौ}}{\Afootnote{N H\textsuperscript{pc}; \skt{शृङ्गाटककरुकौ ।}\ H\textsuperscript{ac}}}\skt{ । }\pend
\pstart\versepar \skt{पिण्डोलुकं }\edtext{\skt{कफकरं}}{\Afootnote{N H\textsuperscript{pc}; \skt{कफरं}\ H\textsuperscript{ac}}}\skt{ गुरुवातप्रकोपनं ॥}\pend

\pstart\versepar \ledsidenote{\passref{1.46.210.add8}}\skt{सुरेन्द्रकन्दश्लेष्मघ्नो विपाके कटुपित्तकृत् ॥ }\pend
\pstart\versepar \skt{शलरणमाणकप्रभृतयः कन्दाः ॥}\edtext{}{\lemma{\skt{कटुपित्तकृत् ॥ शलरणमाणकप्रभृतयः}}\Afootnote{H; \skt{कटुपित्तकृत् ॥ शलशूरण माणकप्रभृतयः}\ N}}\pend

\pstart\versepar \ledsidenote{\passref{1.46.210.add9}}\skt{ईषत्कषायाः कटुकाविष्टम्भिनो गुरवः । }\pend
\pstart\versepar \skt{कफकृद्वातलाः पित्तहराश्च ॥}\pend

\pstart\versepar \ledsidenote{\passref{1.46.210.add10}}\skt{माणकं स्वादुपाकं तु गुरु }\edtext{\skt{वापि}}{\Afootnote{H; \skt{चापि}\ N}}\skt{ }\edtext{\skt{प्रकीर्त्तितं}}{\Afootnote{H; \skt{प्रकीर्तित ।}\ N}}\skt{ । }\pend
\pstart\versepar \skt{शलकन्दस्तु नात्युष्णः }\edtext{\skt{शूरणोगुदकीलहा}}{\Afootnote{H; \skt{शूरणो गुदकीलहा ॥}\ N}}\skt{ ॥}\pend

\pstart\versepar \ledsidenote{\passref{1.46.210.add11}}\edtext{\skt{कुमुदोत्पलपद्मानाङ्कन्दमारुतकोपनाः}}{\lemma{\skt{कुमुदोत्पलपद्मानाङ्}}\Afootnote{H; \skt{कुमोदोत्पलपद्मानां}\ N}}\skt{ । }\pend
\pstart\versepar \skt{कषायाः पित्तशमना विपाके मधुराहिमाः ।}\pend

\pstart\versepar \ledsidenote{\passref{1.46.210.add12}}\edtext{\skt{वराहकन्दश्लेष्मघ्नः}}{\Afootnote{H; \skt{वराहकन्दः श्लेष्मघ्नः}\ N}}\skt{ }\edtext{\skt{कटुकोरसपाकतः}}{\Afootnote{H; \skt{कटुरसपाकतः ॥}\ N}}\skt{ ॥ }\pend
\pstart\versepar \skt{कुष्ठमेहक्रिमिहरो वृष्योष्णः पित्तवर्द्धनः ।}\pend

\pstart\versepar \ledsidenote{\passref{1.46.210.add13}}\skt{तालतात्रकनालिकेरखर्ज्जूरप्रभृतीनाम्मस्तकजातानि ॥ }\pend
\pstart\versepar \skt{स्वादुपाकरसान्याहुः पित्तरक्तहराणि च ।}\pend

\pstart\versepar \ledsidenote{\passref{1.46.210.add14}}\edtext{\skt{शुक्रलाभ्यनिलघ्नानि}}{\Afootnote{H; \skt{शुक्रलायनिलघ्नानि}\ N}}\skt{ कफवृद्धिकराणि च ॥ }\pend
\pstart\versepar \skt{वालं ह्यनार्त्तवञ्जीर्ण्णं व्याधितं क्रिमिभक्षितं ।}\pend

\pstart \ledsidenote{\passref{1.46.210.add15}}\skt{कन्दम्विवर्ज्जयेत्सर्व्वं यो वा }\edtext{\skt{सम्यक्रोहति}}{\Afootnote{H; \skt{सम्यक्रोहति ॥ {\symfont ❈} ॥}\ N}}\skt{ ॥}\pend

\pstart \ledsidenote{\passref{1.46.211}}\skt{अथ }\edtext{\edtext{\skt{शाकवर्ग्गमुपदेक्ष्यामः}}{\lemma{\skt{शाकवर्ग्गम्}}\Afootnote{N H; \skt{शाकवर्गः ।}\ A1938}}}{\lemma{\skt{उपदेक्ष्यामः}}\Afootnote{N H; \skt{शाकान्यत ऊर्ध्वं वक्ष्यामः ।}\ A1938}}\skt{ ॥ }\edtext{\skt{पुष्पफलालाबुकलिङ्गप्रभृतीनि}}{\Afootnote{H; \skt{तत्र पुष्पफलालाबुकालिन्दकप्रभृतीनि ॥२११॥}\ A1938; \skt{तत्र पुष्पफलाबुकालिङ्गप्रभृतीनि ।}\ N}}\skt{ ॥}\pend

\pstart\versepar \ledsidenote{\passref{1.46.212}}\edtext{\skt{पित्तघ्नान्यनिलं}}{\Afootnote{N H; \skt{पित्तघ्नान्यनिलं}\ A1938}}\skt{ कुर्युस्तथा मन्दकफानि च । }\pend
\pstart\versepar \skt{}\edtext{\skt{सृष्टमूत्रपुरीषाणि}}{\Afootnote{A1938 H; \skt{शृष्टमूत्रपुरीषाणि}\ N}}\skt{ स्वादुपाकरसानि च ॥}\pend

\pstart\versepar \ledsidenote{\passref{1.46.213}}\skt{पित्तघ्नन्तेषु }\edtext{\edtext{\skt{कुष्माण्डम्बालम्मध्यङ्कफापहं}}{\lemma{\skt{कुष्माण्डम्}}\Afootnote{H; \skt{कूष्माण्डं}\ A1938 N}}}{\lemma{\skt{कफापहं}}\Afootnote{N H; \skt{कफावहम् ।}\ A1938}}\skt{ । }\pend
\pstart\versepar \skt{}\edtext{\skt{पक्वं}}{\Afootnote{N H; \skt{शुक्लं}\ A1938}}\skt{ लघूष्णं सक्षारन्दीपनं बस्तिशोधनं ॥}\pend

\pstart \ledsidenote{\passref{1.46.214ab}}\skt{सर्व्वदोषहरं हृद्यं पथ्यञ्चेतोविकारिणां ।}\pend

\pstart \ledsidenote{\passref{1.46.216}}\edtext{\edtext{\skt{त्रपुषेर्व्वारुक्केर्क्कारुशीर्ण्णवृत्तप्रभृतीनि}}{\lemma{\skt{त्रपुषेर्}}\Afootnote{H; \skt{त्रपुसैर्}\ A1938; \skt{त्रपुसेर्}\ N}}}{\lemma{\skt{व्वारुक्केर्क्कारुशीर्ण्णवृत्तप्रभृतीनि}}\Afootnote{H; \skt{वारुकर्कारुकशीर्णवृन्तप्रभृतीनि ॥२१६॥}\ A1938; \skt{वारुकेङ्क्कारुसीर्ण्णवृत्तप्रभृतीनि ॥}\ N}}\skt{ ॥}\pend

\pstart\versepar \ledsidenote{\passref{1.46.217}}\edtext{\skt{स्वादुतिक्तरसान्याहुस्तथावातहराणि च}}{\lemma{\skt{तथावातहराणि च}}\Afootnote{N H; \skt{कफवातकराणि च ।}\ A1938}}\skt{ । }\pend
\pstart\versepar \skt{सृष्टमूत्रपुरीषाणि }\edtext{\skt{कफपित्तहराणि}}{\Afootnote{N H; \skt{रक्तपित्तहराणि}\ A1938}}\skt{ च ॥}\pend

\pstart\versepar \ledsidenote{\passref{1.46.218}}\skt{बालं }\edtext{\edtext{\skt{सनीलन्त्रपुषन्तेषाम्पित्तहरं}}{\lemma{\skt{त्रपुषन्}}\Afootnote{H; \skt{त्रपुसं}\ A1938 N}}}{\lemma{\skt{पित्तहरं}}\Afootnote{A1938 H; \skt{पित्तकरं}\ N}}\skt{ स्मृतं । }\pend
\pstart\versepar \skt{}\edtext{\edtext{\skt{तत्पाण्डुङ्कफकृज्जीर्णमम्लम्वातकफापहं}}{\lemma{\skt{तत्पाण्डुङ्}}\Afootnote{H; \skt{तत्पाण्डु}\ A1938; \skt{तत्पाण्डुः}\ N}}}{\lemma{\skt{अम्लम्}}\Afootnote{A1938 H; \skt{अम्व्लम्}\ N}}\skt{ ॥}\pend

\pstart\versepar \ledsidenote{\passref{1.46.219}}\edtext{\edtext{\skt{एर्व्वारुकंसकेर्क्कारुः}}{\lemma{\skt{एर्व्वारुकंसकेर्}}\Afootnote{H; \skt{एर्वारुकं सक}\ A1938; \skt{एर्वारुकं सके}\ N}}}{\lemma{\skt{क्कारुः}}\Afootnote{H; \skt{र्कारु}\ A1938; \skt{र्क्कारुः}\ N}}\skt{ सम्पक्वङ्कफवातकृत् । }\pend
\pstart\versepar \skt{सक्षारं मधुरं रुच्यं दीपनन्नातिपित्तलं ॥}\pend

\pstart \ledsidenote{\passref{1.46.221}}\edtext{\skt{पिप्पलीमरिचशृङ्गवेरहिंगुजीरककुस्तुम्बुरुजम्बीरसुमुखसुरसार्ज्जकभूस्तृणसुगन्धककालमालकुठेरकक्षवकखरपुषकासमर्द्दमधुशिग्रुफणिज्जकसर्षपवेत्रकुलथगण्डीरतिलपर्ण्णिवर्षाभूचित्रकमूलकरसोनपलाण्डुप्रभृतीनि}}{\Afootnote{H; \skt{पिप्पलीमरिचशृङ्गवेरार्द्रकहिङ्गुजीरककुस्तुम्बुरुजम्बीरसुमुखसुरसार्जकभूस्तृणसुगन्धक- कासमर्दककालमालकुठेरकक्षवकखरपुष्पशिग्रुमधुशिग्रुफणिज्झकसर्षपराजिकाकुलाहलावगुत्थगण्डीरतिलपर्णिका- वर्षाभूचित्रकमूलकलशुनकलायपलाण्डुप्रभृतीनि ॥२२१॥}\ A1938; \skt{पिप्पलीमरिचसृङ्गवेरहिंगुजीरककुस्तुम्बुरुजम्बीरसुमुखसुरससार्ज्जकभूस्तृणसुगन्धकर्कलमाकुठेरकक्षवकखरपुसकासमर्दमधुसिग्रुफणिज्ज्ञकषर्सपचेत्रकुलत्थगण्डीरतिलपर्ण्णिवर्षाभूचित्रकमूलकरसोलप्लापण्डुप्रभृती}\ N\textsuperscript{ac}; \skt{पिप्पलीमरिचसृङ्गवेरहिंगुजीरककुस्तुम्बुरुजम्बीरसुमुखसुरससार्ज्जकभूस्तृणसुगन्धकर्कलमाकुठेरकक्षवकखरपुसकासमर्दमधुसिग्रुफणिज्ज्ञकषर्सपचेत्रकुलत्थगण्डीरतिलपर्ण्णिवर्षाभूचित्रकमूलकरसोलप्लाण्डुप्रभृतीः}\ N\textsuperscript{pc}}}\skt{ ॥}\pend

\pstart\versepar \ledsidenote{\passref{1.46.222}}\skt{कटून्युष्णानि रुच्यानि वातश्लेष्महराणि च । }\pend
\pstart\versepar \skt{कृतान्नेषूपयुज्यन्ते संस्कारार्थमनेकधा ॥}\pend

\pstart\versepar \ledsidenote{\passref{1.46.223}}\edtext{\skt{तेषु}}{\Afootnote{N H; \skt{तेषां}\ A1938}}\skt{ गुर्व्वी }\edtext{\skt{स्वादुशीता}}{\Afootnote{A1938 H; \skt{स्वादुसीता}\ N}}\skt{ }\edtext{\skt{पिप्पल्यार्द्रा}}{\Afootnote{A1938 H; \skt{पिप्पलयार्द्रा}\ N}}\skt{ कफावहा । }\pend
\pstart\versepar \skt{}\edtext{\skt{शुष्का}}{\Afootnote{A1938 H; \skt{शुश्का}\ N}}\skt{ }\edtext{\skt{कफानिलहनी}}{\Afootnote{H; \skt{कफानिलघ्नी सा}\ A1938; \skt{कफानिलहरी}\ N}}\skt{ वृष्या }\edtext{\skt{पित्तविरोधनी}}{\Afootnote{H; \skt{पित्ताविरोधिनी ॥२२३॥}\ A1938; \skt{पित्तविरोधिनी ।}\ N}}\skt{ ॥}\pend

\pstart\versepar \ledsidenote{\passref{1.46.224}}\edtext{\skt{स्वादुपाकार्द्रमरिचं}}{\Afootnote{N H; \skt{स्वादुपाक्यार्द्रमरिचं}\ A1938}}\skt{ गुरु }\edtext{\skt{श्लेष्मप्रसेकि च}}{\Afootnote{A1938 H; \skt{श्लेष्मप्रशेकि च ।}\ N}}\skt{ । }\pend
\pstart\versepar \skt{कटूष्णं लघु }\edtext{\skt{तच्छुष्कमवृष्यङ्कफवातजित्}}{\lemma{\skt{तच्छुष्कमवृष्यङ्}}\Afootnote{A1938 H; \skt{तच्छुश्कमवृष्यं}\ N}}\skt{ ॥}\pend

\pstart\versepar \ledsidenote{\passref{1.46.225}}\edtext{\skt{नात्युष्णम्नातिरूक्षञ्च}}{\lemma{\skt{नातिरूक्षञ्च}}\Afootnote{N H; \skt{नातिशीतं च}\ A1938}}\skt{ वीर्यतो मरिचं सितं । }\pend
\pstart\versepar \skt{}\edtext{\skt{गुणवत्मरिचेभ्यश्च}}{\Afootnote{H; \skt{गुणवन्मरिचेभ्यश्च}\ A1938 N}}\skt{ चक्षुष्यञ्च }\edtext{\skt{विशेषतः}}{\Afootnote{A1938 H; \skt{विसेषतः ।}\ N}}\skt{ ॥}\pend

\pstart\versepar \ledsidenote{\passref{1.46.226}}\skt{नागरङ्कफवातघ्नम्विपाके }\edtext{\skt{मधुरङ्कटुः}}{\lemma{\skt{कटुः}}\Afootnote{N H; \skt{कटु ।}\ A1938}}\skt{ । }\pend
\pstart\versepar \skt{}\edtext{\skt{वृष्योष्णरोचनं}}{\Afootnote{N H; \skt{वृष्योष्णं रोचनं}\ A1938}}\skt{ हृद्यं सस्नेहं }\edtext{\skt{स्वादु}}{\Afootnote{N H; \skt{लघु}\ A1938}}\skt{ दीपनं ॥}\pend

\pstart\versepar \ledsidenote{\passref{1.46.227}}\skt{कफानिलहरं स्वर्यम्विबन्धानाहशूलनुत् । }\pend
\pstart\versepar \skt{कटूष्णं }\edtext{\skt{रोचनम्वृष्यं}}{\lemma{\skt{वृष्यं}}\Afootnote{N H; \skt{हृद्यं}\ A1938}}\skt{ }\edtext{\skt{हृद्यञ्चैवार्द्रकं}}{\lemma{\skt{हृद्यञ्}}\Afootnote{N H; \skt{वृष्यं}\ A1938}}\skt{ स्मृतं ॥}\pend

\pstart\versepar \ledsidenote{\passref{1.46.228}}\edtext{\skt{लघूष्णम्पाचनं}}{\lemma{\skt{लघूष्णम्}}\Afootnote{A1938 N\textsuperscript{pc} H; \skt{लघूघूष्णं}\ N\textsuperscript{ac}}}\skt{ हिंगु दीपनङ्कफवातजित् । }\pend
\pstart\versepar \skt{}\edtext{\skt{स्निग्धन्तीक्ष्णङ्कटुरसं}}{\lemma{\skt{कटुरसं}}\Afootnote{N H; \emph{om.}\ A1938}}\skt{ }\edtext{\skt{शूलाजीर्ण्णविबन्धजित्}}{\Afootnote{H; \skt{शूलाजीर्णविबन्धनुत् ॥२२८॥}\ A1938; \skt{शूलाजीर्ण्णविबन्धहित् ।}\ N}}\skt{ ॥}\edtext{}{\lemma{\skt{कफवातजित् । स्निग्धन्}}\Afootnote{N H; \skt{कफवातजित् । कटु स्निग्धं सरं}\ A1938}}\pend

\pstart\versepar \ledsidenote{\passref{1.46.229}}\skt{तीक्ष्णोष्णङ्कटुकम्पाके रुच्यम्पित्ताग्निवर्द्धनं । }\pend
\pstart\versepar \skt{कटु श्लेष्मानिलहरं गन्धाढ्यञ्जीरकद्वयं ॥}\pend

\pstart\versepar \ledsidenote{\passref{1.46.230}}\edtext{\skt{काकारवी}}{\Afootnote{H; \skt{कारवी}\ A1938 N}}\skt{ करवी तद्वद्विज्ञेया सोपकुञ्चिका । }\pend
\pstart\versepar \skt{भक्ष्यव्यञ्जनभोज्येषु }\edtext{\skt{विविधेष्ववचारिता}}{\lemma{\skt{अवचारिता}}\Afootnote{A1938 H; \skt{अचारिता ।}\ N}}\skt{ ॥}\pend

\pstart\versepar \ledsidenote{\passref{1.46.231}}\skt{आर्द्रा }\edtext{\skt{कुस्तुम्बरी}}{\Afootnote{A1938 H; \skt{कुस्तुम्बुरी}\ N}}\skt{ कुर्यात्सौगन्ध्यस्वादुहृद्यतां । }\pend
\pstart\versepar \skt{सा }\edtext{\skt{शुष्का}}{\Afootnote{A1938 H; \skt{शुस्का}\ N}}\skt{ }\edtext{\skt{मधुराः}}{\Afootnote{N H; \skt{मधुरा}\ A1938}}\skt{ पाके स्निग्धा }\edtext{\skt{दाहस्तृषापहा}}{\Afootnote{H; \skt{तृड्दाहनाशनी ॥२३१॥}\ A1938; \skt{दाहतृषापहा ।}\ N}}\skt{ ॥}\edtext{}{\lemma{\skt{कुर्यात्सौगन्ध्यस्वादुहृद्यतां । सा}}\Afootnote{H; \skt{कुर्यात्स्वादुसौगन्ध्यहृद्यताम् । सा}\ A1938; \skt{कुर्यात्स्वौगन्ध्यस्वादुहृद्यतां । सा}\ N}}\pend

\pstart\versepar \ledsidenote{\passref{1.46.232}}\edtext{\skt{दोषघ्ना}}{\Afootnote{H; \skt{दोषघ्नी}\ A1938; \skt{दोषघ्नाः}\ N}}\skt{ कटुका }\edtext{\skt{किञ्चित्तिक्ताः}}{\Afootnote{N H; \skt{किञ्चित्तिक्ता}\ A1938}}\skt{ }\edtext{\skt{स्रोतोविशोधनी}}{\Afootnote{A1938 H; \skt{स्रोतोविसोधनी ॥}\ N}}\skt{ । }\pend
\pstart\versepar \skt{}\edtext{\skt{जम्बीरः}}{\Afootnote{A1938 H; \skt{जम्बीर}\ N}}\skt{ }\edtext{\edtext{\skt{पाचनस्तीक्ष्णः}}{\lemma{\skt{पाचनस्}}\Afootnote{N; \skt{पाचन}\ A1938 H}}}{\lemma{\skt{तीक्ष्णः}}\Afootnote{N; \skt{स्तीक्ष्णः}\ A1938 H}}\skt{ }\edtext{\skt{कृमिवातकफापहः}}{\Afootnote{A1938 H; \skt{क्रिमिवातकफापहः ।}\ N}}\skt{ ॥}\pend

\pstart\versepar \ledsidenote{\passref{1.46.233}}\edtext{\skt{सुरभिर्द्दीपनो}}{\lemma{\skt{सुरभिर्}}\Afootnote{H; \skt{सुराभिर्}\ A1938; \skt{सुर}\ N}}\skt{ }\edtext{\skt{हृद्यो}}{\Afootnote{N H; \skt{रुच्यो}\ A1938}}\skt{ }\edtext{\skt{मुखवैशद्यकारकः}}{\Afootnote{A1938 H; \skt{मुखवैषद्यकारकः ।}\ N}}\skt{ । }\pend
\pstart\versepar \skt{}\edtext{\skt{कफानिलविषश्वासकासदौर्गन्ध्यनाशनः}}{\Afootnote{A1938 H; \skt{कफानिलविषश्वासकासदौर्गन्ध्यनासनः ।}\ N}}\skt{ ॥}\pend

\pstart \ledsidenote{\passref{1.46.234ab}}\skt{पित्तकृत्पार्श्वशूलघ्नः }\edtext{\skt{सुरसः}}{\Afootnote{A1938 H; \skt{सुरस्}\ N}}\skt{ समुदाहृतः ।}\pend

\pstart\versepar \ledsidenote{\passref{1.46.235}}\edtext{\skt{रूक्षाः}}{\Afootnote{H; \emph{om.}\ A1938; \skt{रूक्षा}\ N}}\skt{ }\edtext{\skt{कफघ्ना}}{\Afootnote{A1938 H; \skt{काफघ्ना}\ N}}\skt{ लघवः । }\pend
\pstart\versepar \skt{}\edtext{\skt{सुरसार्ज्जकभूस्तृणाः}}{\Afootnote{N H; \skt{रूक्षास्तीक्ष्णोष्णाः पित्तवर्धनाः । कटुपाकरसाश्चैव सुरसार्जकभूस्तृणाः ॥२३५॥}\ A1938}}\skt{ ॥}\pend

\pstart\versepar \ledsidenote{\passref{1.46.236}}\skt{मधुरः }\edtext{\skt{कफवातघ्नः}}{\Afootnote{A1938 N; \skt{कफवातघ्नᳶ}\ H}}\skt{ पाचनः कण्ठशोधनः । }\pend
\pstart\versepar \skt{विशेषतो }\edtext{\skt{रुचिकरः}}{\Afootnote{N H; \skt{पित्तहरः}\ A1938}}\skt{ सतिक्तः कासमर्द्दकः ॥}\pend

\pstart\versepar \ledsidenote{\passref{1.46.237}}\skt{कटुः सक्षारमधुरः }\edtext{\edtext{\skt{सिग्रुस्तिक्तोऽथ}}{\lemma{\skt{सिग्रुस्तिक्तो}}\Afootnote{\skt{शिग्रुस्तिक्तो}\ A1938; \skt{सिग्रुस्तिक्तो}\ N; \skt{सिग्रुस्तिक्तो}\ H}}}{\lemma{\skt{ऽथ}}\Afootnote{A1938 H; \skt{थ}\ N}}\skt{ पिच्छिलः । }\pend
\pstart\versepar \skt{}\edtext{\skt{मधुसिग्रुः}}{\Afootnote{H; \skt{मधुशिग्रुः}\ A1938; \skt{मधुशिग्रु}\ N}}\skt{ }\edtext{\skt{सरस्तिक्तः}}{\Afootnote{A1938 H; \skt{सरस्तिक्तः}\ N}}\skt{ शोफघ्नो दीपनः कटुः ॥}\pend

\pstart\versepar \ledsidenote{\passref{1.46.238}}\skt{विदाहि }\edtext{\skt{बद्धविण्मूत्रं}}{\Afootnote{A1938 H; \skt{बद्धविड्मूत्रं}\ N}}\skt{ रूक्षन्तीक्ष्णोष्णमेव च । }\pend
\pstart\versepar \skt{}\edtext{\skt{कफघ्नं}}{\Afootnote{N H; \skt{त्रिदोषं}\ A1938}}\skt{ सार्षपं शाकं }\edtext{\skt{मासूरं}}{\Afootnote{N H; \skt{गाण्डीरं}\ A1938}}\skt{ }\edtext{\edtext{\skt{शाकमेव}}{\lemma{\skt{शाकम्}}\Afootnote{H; \skt{वेगना}\ A1938; \skt{च शाकम्}\ N}}}{\lemma{\skt{एव}}\Afootnote{N H; \skt{म}\ A1938}}\skt{ च ॥}\pend

\pstart\versepar \ledsidenote{\passref{1.46.239}}\edtext{\edtext{\skt{चित्रकस्तिलपर्ण्णी}}{\lemma{\skt{चित्रकस्}}\Afootnote{N H; \skt{चित्रक}\ A1938}}\skt{ च}}{\lemma{\skt{तिलपर्ण्णी च}}\Afootnote{N H; \skt{स्तिलपर्णी च}\ A1938}}\skt{ कफशोफहरे }\edtext{\skt{लघू}}{\Afootnote{A1938 H; \skt{लघूः ।}\ N}}\skt{ । }\pend
\pstart\versepar \skt{}\edtext{\skt{वर्षा}}{\Afootnote{H; \skt{वर्षाभूः}\ A1938; \skt{वर्षाभौ}\ N}}\skt{ }\edtext{\skt{कफवातघ्नौ}}{\Afootnote{N H; \skt{कफवातघ्नी}\ A1938}}\skt{ }\edtext{\skt{हितौ}}{\Afootnote{H; \skt{हिता}\ A1938; \skt{हिनौ}\ N}}\skt{ शोफोदरार्शसां ॥}\pend

\pstart\versepar \ledsidenote{\passref{1.46.240}}\edtext{\skt{कटुस्तिक्तरसा}}{\Afootnote{H; \skt{कटुतिक्तरसा}\ A1938 N}}\skt{ हृद्या रोचनी वह्निदीपनी । }\pend
\pstart\versepar \skt{सर्व्वदोषहरा }\edtext{\skt{लघ्वी}}{\Afootnote{A1938 H; \skt{लघ्नी}\ N}}\skt{ कण्ठ्या मूलकपोतिका ॥}\pend

\pstart\versepar \ledsidenote{\passref{1.46.241}}\edtext{\skt{महन्तं}}{\Afootnote{H; \skt{महत्त}\ A1938 N}}\skt{ }\edtext{\skt{गुरु}}{\Afootnote{H; \skt{द्गुरु}\ A1938 N}}\skt{ }\edtext{\skt{विष्टम्भि}}{\Afootnote{A1938 H; \skt{विष्टम्भी}\ N}}\skt{ }\edtext{\edtext{\skt{तीक्ष्णमामं}}{\lemma{\skt{तीक्ष्णम्}}\Afootnote{N H; \skt{तीक्ष्ण}\ A1938}}}{\lemma{\skt{आमं}}\Afootnote{N H; \skt{मामं}\ A1938}}\skt{ त्रिदोषकृत् । }\pend
\pstart\versepar \skt{तदेव }\edtext{\skt{स्निग्धसिद्धन्तु}}{\Afootnote{N H; \skt{स्नेहसिद्धं तु}\ A1938}}\skt{ }\edtext{\skt{श्लेष्मकृद्वातपित्तजित्}}{\Afootnote{H; \skt{पित्तनुत्कफवातजित् ॥२४१॥}\ A1938; \skt{श्लेष्म कृद्धात पित्तजित् ।}\ N}}\skt{ ॥}\pend

\pstart\versepar \ledsidenote{\passref{1.46.242}}\edtext{\skt{शुष्कन्तु}}{\Afootnote{H; \skt{त्रि}\ A1938; \skt{शुस्कं तु}\ N}}\skt{ }\edtext{\skt{शोफशमनं}}{\Afootnote{H; \skt{दोषशमनं}\ A1938; \skt{सोफशमनं}\ N}}\skt{ }\edtext{\skt{विषदोषहरं}}{\Afootnote{N H; \skt{शुष्कं विषदोषहरं}\ A1938}}\skt{ }\edtext{\skt{लघुः}}{\Afootnote{N H\textsuperscript{ac}; \skt{लघु ।}\ A1938 H\textsuperscript{pc}}}\skt{ । }\pend
\pstart\versepar \skt{विष्टम्भि वातलं शाकं }\edtext{\skt{शुष्कमन्यत्र}}{\lemma{\skt{शुष्कम्}}\Afootnote{A1938 H; \skt{शुस्कम्}\ N}}\skt{ मूलकात् ॥}\pend

\pstart\versepar \ledsidenote{\passref{1.46.243}}\skt{पुष्पञ्च पत्रञ्च फलन्तथैव यथोत्तरन्ते }\edtext{\skt{लघवः}}{\Afootnote{N H; \skt{गुरवः}\ A1938}}\skt{ प्रदिष्टाः । }\pend
\pstart\versepar \skt{तेषां तु }\edtext{\skt{पुष्पङ्कफवातहन्तृ}}{\lemma{\skt{कफवातहन्तृ}}\Afootnote{N H; \skt{कफपित्तहन्तृ}\ A1938}}\skt{ फलन्निहन्यात्कफमारुतौ तु ॥}\pend

\pstart\versepar \ledsidenote{\passref{1.46.244}}\edtext{\skt{स्निग्धोष्णतीक्ष्णः}}{\Afootnote{A1938 H; \skt{स्निग्धौष्णतीक्ष्णः}\ N}}\skt{ कटुपिच्छिलश्च गुरुः सरः }\edtext{\edtext{\skt{स्वादुरसोऽथ}}{\lemma{\skt{स्वादुरसो}}\Afootnote{A1938 H; \skt{स्वादुरसो}\ N}}}{\lemma{\skt{ऽथ}}\Afootnote{H; \skt{च}\ A1938; \skt{थ}\ N}}\skt{ बल्यः । }\pend
\pstart\versepar \skt{वृष्यश्च }\edtext{\skt{मेधास्वरवर्ण्णचक्षुर्ब्भग्नास्थिसन्धानकरो}}{\Afootnote{A1938 H; \skt{मेधाश्वरवर्ण्णश्चक्षुर्भग्नास्थिसन्धानकरो}\ N}}\skt{ रसोनः ॥}\pend

\pstart\versepar \ledsidenote{\passref{1.46.245}}\edtext{\skt{हृद्रोगजीर्ण्णज्वरकुक्षिशूलविबन्धगुल्मारुचिकासशोफान्}}{\Afootnote{A1938 H; \skt{हृद्रोगजीर्ण्णज्वरकुक्षिशूलविबन्धगुल्मारुचिकासशोपान् ।}\ N}}\skt{ । }\pend
\pstart\versepar \skt{}\edtext{\skt{दुर्न्नामकुष्ठानलसादजन्तून्समीरणश्वासकफांश्च}}{\Afootnote{H; \skt{दुर्नामकुष्ठानलसादजन्तुसमीरणश्वासकफांश्च}\ A1938; \skt{दुर्ण्णामकुष्ठानलसादजन्तूं समीरणश्वासकफाश्च}\ N}}\skt{ हन्ति ॥}\pend

\pstart\versepar \ledsidenote{\passref{1.46.246}}\edtext{\skt{नात्युष्णवीर्योऽनिलहा}}{\lemma{\skt{ऽनिलहा}}\Afootnote{A1938 H; \skt{निलहा}\ N}}\skt{ कटुश्च तीक्ष्णो }\edtext{\edtext{\skt{गुरुर्न्नाति}}{\lemma{\skt{गुरुर्}}\Afootnote{N H; \skt{गुरु}\ A1938}}}{\lemma{\skt{न्नाति}}\Afootnote{N H\textsuperscript{pc}; \skt{र्नाति}\ A1938; \skt{न्नाम}\ H\textsuperscript{ac}}}\skt{ कफावहश्च । }\pend
\pstart\versepar \skt{बलावहः }\edtext{\edtext{\skt{पित्तकरोऽथ}}{\lemma{\skt{पित्तकरो}}\Afootnote{A1938 H; \skt{पित्तकरो}\ N}}\skt{ }\edtext{\edtext{\skt{किञ्चित्पलाण्डुरग्निञ्च}}{\lemma{\skt{चित्पलाण्डुर्}}\Afootnote{N H; \skt{चित्पलाण्डु}\ A1938}}}{\lemma{\skt{अग्निञ्च}}\Afootnote{N H; \skt{रग्निं च}\ A1938}}}{\lemma{\skt{ऽथ किञ्}}\Afootnote{A1938 H; \skt{थ किञ्}\ N}}\skt{ विवर्द्धयेत्तु ॥}\pend

\pstart\versepar \ledsidenote{\passref{1.46.247}}\skt{स्निग्धो रुचिष्यः }\edtext{\skt{स्थिरधातुकर्त्ता}}{\Afootnote{H; \skt{स्थिरधातुकारी}\ A1938; \skt{स्थिरधातुकर्त्रा}\ N}}\skt{ }\edtext{\edtext{\skt{बल्योऽथ}}{\lemma{\skt{बल्यो}}\Afootnote{A1938 H; \skt{बल्यो}\ N}}}{\lemma{\skt{ऽथ}}\Afootnote{A1938 H; \skt{थ}\ N}}\skt{ मेधाकफपुष्टिदश्च । }\pend
\pstart\versepar \skt{}\edtext{\edtext{\skt{स्वादुर्ग्गुरुः}}{\lemma{\skt{स्वादुर्}}\Afootnote{H; \skt{स्वादु}\ A1938 N}}}{\lemma{\skt{ग्गुरुः}}\Afootnote{N H; \skt{र्गुरुः}\ A1938}}\skt{ }\edtext{\skt{शोणितपित्तशस्तः स}}{\Afootnote{A1938 H; \skt{सोणितपित्तशस्तः स}\ N}}\skt{ पिच्छिलः }\edtext{\edtext{\skt{क्षीरपलाण्डुरुक्तः}}{\lemma{\skt{क्षीरपलाण्डुर्}}\Afootnote{H; \skt{क्षीरपलाण्डु}\ A1938 N}}}{\lemma{\skt{उक्तः}}\Afootnote{H; \skt{रुक्तः}\ A1938 N}}\skt{ ॥}\pend

\pstart \ledsidenote{\passref{1.46.249}}\edtext{\skt{चुच्चूयूथिकावरुणजीवन्तीनदीभल्लातकच्छगलान्त्रीवृक्षादनीफञ्जीतण्डुलीयकपोतकाश्च}}{\Afootnote{H; \skt{चुच्चूयूथिकातरुणीजीवन्तीबिम्बीतिकानदी(न्दी)भल्लातकच्छगलान्त्री- वृक्षादनीफञ्जीशाल्मलीशेलुवनस्पतिप्रस}\ A1938; \skt{चुच्चूयूथिकावरुणजीवन्तीनंदीभल्लातकच्छगलान्त्रीवृक्षादनीफञ्जीतण्डुलीय}\ N}}\skt{ }\edtext{\skt{वलचिल्लीपालब्दावास्तूकप्रभृतीनि}}{\Afootnote{H; \skt{वशणकर्बुदारकोविदारप्रभृतीनि ॥२४९॥}\ A1938; \skt{कोपातकासुबलविल्लीपालद्कृवाश्तूकप्रभृतीनि ।}\ N}}\skt{ ॥}\pend

\pstart \ledsidenote{\passref{1.46.257ab}}\skt{सृष्टमूत्रपुरीषाणि सक्षारमधुराणि च ।}\pend

\pstart\versepar \ledsidenote{\passref{1.46.258}}\skt{मधुरो रसपाकाभ्यां }\edtext{\skt{रक्तपित्तकफापहः}}{\Afootnote{N H; \skt{रक्तपित्तमदापहः ।}\ A1938}}\skt{ ॥ }\pend
\pstart\versepar \skt{तेषां }\edtext{\skt{शीततरो}}{\Afootnote{N H; \skt{शीततमो}\ A1938}}\skt{ रूक्षस्तण्डुलीयो विषापहः ।}\pend

\pstart\versepar \ledsidenote{\passref{1.46.259}}\skt{स्वादुपाकरसा वृष्या वातपित्तमदापहा ॥ }\pend
\pstart\versepar \skt{}\edtext{\skt{उपोतका}}{\Afootnote{N H; \skt{उपोदका}\ A1938}}\skt{ सरा स्निग्धा बल्या श्लेष्मकरी हिमा ।}\pend

\pstart\versepar \ledsidenote{\passref{1.46.260}}\edtext{\edtext{\skt{लघुर्व्विपाकक्रिमिहा}}{\lemma{\skt{लघुर्}}\Afootnote{H; \skt{कटु}\ A1938; \skt{लघु}\ N}}}{\lemma{\skt{व्विपाकक्रिमिहा}}\Afootnote{N H; \skt{र्विपाके कृमिहा}\ A1938}}\skt{ मेधाग्निबलवर्द्धनः ॥ }\pend
\pstart\versepar \skt{सक्षारः सर्व्वदोषघ्नो वास्तूको रोचनः सरः ।}\pend

\pstart\versepar \ledsidenote{\passref{1.46.261}}\skt{चिल्ली वास्तूकवज्ज्ञेया पालङ्क्या तण्डुलीयवत् ॥ }\pend
\pstart\versepar \skt{वातकृद्बद्धविण्मूत्रा }\edtext{\skt{रूक्षाः}}{\Afootnote{N H; \skt{रूक्षा}\ A1938}}\skt{ पित्तकफे }\edtext{\skt{हिताः}}{\Afootnote{N H; \skt{हिता । शाकमाश्वबलं रूक्षं बद्धविण्मूत्रमारुतम् ॥२६१॥}\ A1938}}\skt{ ॥}\pend

\pstart \ledsidenote{\passref{1.46.262}}\edtext{\skt{माण्डूकपर्ण्णीसप्तलासुनिषण्णकसुवर्च्चलापिप्पलीगुडूचीगोजिह्वाप्रपुनाटावल्गुजसतीनबृहतीफलपटोलवार्त्ताकुकारवेल्लकटुकिकाकेवुकोरुबूकपर्प्पटकिराततिक्तकर्क्कोटकारिष्टकोषातकीवेत्रकरीराटरूषकार्क्कपुष्पीप्रभृतीनि}}{\Afootnote{H; \skt{मण्डूकपर्णीसप्तलासुनिषण्णकसुवर्चलापिप्पलीगुडूचीगोजिह्वा- काकमाचीप्रपुन्नाडावल्गुजसतीनबृहतीकण्टकारिकाफल- पटोलवार्ताककारवेल्लककटुकिकाकेबु(म्बु)कोरुबूकपर्पटककिराततिक्त- कर्कोटकारिष्टकोशातकीवेत्रकरीराटरूषकार्कपुष्पीप्रभृतीनि ॥२६२॥}\ A1938; \skt{मण्डूकपर्ण्णीसप्तलासुनिषर्ण्णकसुवर्च्चला ॥ पिप्पलीगुडूचीगोजिह्वाप्रपुनाटावल्गुजसतीनबृहतीफलपटोलवार्ताकुकारवेल्लककटुकिकाकेवुकौओरुबूकपर्पटकिरार्ततिक्तकर्कोटकारिष्ठकोसातकीवेत्रकरीराटरूषकार्कपुष्पीप्रभृतीनि ।}\ N}}\skt{ ॥}\pend

\pstart\versepar \ledsidenote{\passref{1.46.263}}\edtext{\edtext{\skt{कफपित्तहराण्याहुर्हृद्यानि}}{\lemma{\skt{कफपित्तहराण्य्}}\Afootnote{N\textsuperscript{pc} H; \skt{रक्तपित्तहराण्य्}\ A1938; \skt{कफपित्तहरा}\ N\textsuperscript{ac}}}}{\lemma{\skt{आहुर्}}\Afootnote{A1938 N\textsuperscript{pc} H; \skt{ण्य्हुः}\ N\textsuperscript{ac}}}\skt{ सुलघूनि च । }\pend
\pstart\versepar \skt{कुष्ठमेहज्वरश्वासकासारुचिहराणि च ॥}\pend

\pstart\versepar \ledsidenote{\passref{1.46.264}}\edtext{\skt{कषायानुहिता}}{\Afootnote{N H; \skt{कषाया तु हिता}\ A1938}}\skt{ पित्ते }\edtext{\skt{तिक्ता}}{\Afootnote{H; \emph{om.}\ A1938; \skt{तिक्ताः}\ N}}\skt{ }\edtext{\skt{स्वादुरसा}}{\Afootnote{H; \skt{स्वादुपाकरसा}\ A1938; \skt{श्वादुरसा}\ N}}\skt{ हिमाः । }\pend
\pstart\versepar \skt{}\edtext{\skt{लघ्वी}}{\Afootnote{A1938 H; \skt{लघ्नी}\ N}}\skt{ मण्डूकपर्ण्णी तु }\edtext{\skt{तेषां}}{\Afootnote{N H; \skt{तद्व}\ A1938}}\skt{ }\edtext{\skt{गोजिह्विका}}{\Afootnote{H; \skt{द्गोजिह्विका}\ A1938; \skt{गोजिह्वका}\ N}}\skt{ }\edtext{\skt{तथा}}{\Afootnote{N H; \skt{मता ॥२६४॥}\ A1938}}\skt{ ॥}\pend

\pstart\versepar \ledsidenote{\passref{1.46.265}}\skt{अविदाही त्रिदोषघ्नः }\edtext{\skt{सांग्राही}}{\Afootnote{H; \skt{सङ्ग्राही}\ A1938 N}}\skt{ सुनिषण्णकः । }\pend
\pstart\versepar \skt{}\edtext{\skt{अवल्गुजः}}{\Afootnote{A1938 H; \skt{अवल्भुजः}\ N}}\skt{ }\edtext{\skt{कटुपाकी तु}}{\Afootnote{N H; \skt{कटुः पाके}\ A1938}}\skt{ तिक्तः पित्तकफापहः ॥}\pend

\pstart\versepar \ledsidenote{\passref{1.46.266}}\edtext{\skt{ईषत्तिक्तं}}{\Afootnote{A1938 H; \skt{ईषतिक्तं}\ N}}\skt{ त्रिदोषघ्नं शाकं कटु सतीनजं । }\pend
\pstart\versepar \skt{}\edtext{\skt{नात्युष्णं}}{\Afootnote{H; \skt{नात्युष्ण}\ A1938 N}}\skt{ शीतं कुष्ठघ्नं }\edtext{\skt{काकमाच्याश्च}}{\Afootnote{N H; \skt{काकमाच्यास्तु}\ A1938}}\skt{ तद्विधं ॥}\pend

\pstart\versepar \ledsidenote{\passref{1.46.267}}\edtext{\skt{कण्डुकुष्ठक्रिमिघ्नानि}}{\Afootnote{H; \skt{कण्डूकुष्ठकृमिघ्नानि}\ A1938; \skt{कण्डूकुष्ठक्रिमिघ्नानि}\ N}}\skt{ कफवातहराणि च । }\pend
\pstart\versepar \skt{फलानि बृहतीनां तु कटुतिक्तलघूनि च ॥}\pend

\pstart\versepar \ledsidenote{\passref{1.46.268}}\edtext{\edtext{\edtext{\skt{कफपित्तहरम्व्रण्यमुष्णन्तिक्तमवातलं}}{\lemma{\skt{उष्णन्}}\Afootnote{A1938 H; \skt{उष्ण}\ N}}}{\lemma{\skt{तिक्तम्}}\Afootnote{N H; \skt{तिक्त}\ A1938}}}{\lemma{\skt{अवातलं}}\Afootnote{N H; \skt{मवातलम् ।}\ A1938}}\skt{ । }\pend
\pstart\versepar \skt{पटोलङ्कटुकम्पाके वृष्यं रोचनदीपनं ॥}\pend

\pstart\versepar \ledsidenote{\passref{1.46.269}}\skt{कफवातहरन्तिक्तं रोचनङ्कटुकं }\edtext{\skt{लघुः}}{\Afootnote{N H; \skt{लघु ।}\ A1938}}\skt{ । }\pend
\pstart\versepar \skt{}\edtext{\skt{वार्त्ताकुर्द्दीपनं}}{\lemma{\skt{वार्त्ताकुर्}}\Afootnote{H; \skt{वार्ताकं}\ A1938; \skt{वार्त्ताक्तं}\ N}}\skt{ प्रोक्तं जीर्ण्णं }\edtext{\skt{सक्षारपित्तलं}}{\Afootnote{N H; \skt{सक्षारपित्तलम् । तद्वत्कर्कोटकं विद्यात्कारवेल्लकमेव च ॥२६९॥}\ A1938}}\skt{ ॥}\pend

\pstart\versepar \ledsidenote{\passref{1.46.270}}\edtext{\skt{अटरूषकवेत्राग्रगुडूचीनिम्बपर्प्पटाः}}{\Afootnote{H; \skt{आटरूषकवेत्राग्रगुडूचीनिम्बपर्पटाः ।}\ A1938; \skt{अटरूषकवेत्रोग्रगुडूचीनिम्बपर्पटाः ।}\ N}}\skt{ । }\pend
\pstart\versepar \skt{}\edtext{\edtext{\skt{किराततिक्तसहितास्तिक्ताः}}{\lemma{\skt{किराततिक्तसहितास्}}\Afootnote{N H; \skt{किराततिक्तसहिता}\ A1938}}}{\lemma{\skt{तिक्ताः}}\Afootnote{N H; \skt{स्तिक्ताः}\ A1938}}\skt{ पित्तकफापहाः ॥}\pend

\pstart\versepar \ledsidenote{\passref{1.46.271}}\skt{कफापहं }\edtext{\skt{शाकमुक्तं}}{\lemma{\skt{उक्तं}}\Afootnote{A1938 H; \skt{उष्णं}\ N}}\skt{ }\edtext{\skt{वरुणप्रपुनाटयोः}}{\Afootnote{N H; \skt{वरुणप्रपुनाड(ट)योः ।}\ A1938}}\skt{ । }\pend
\pstart\versepar \skt{रूक्षं लघु च शीतञ्च }\edtext{\skt{वातपित्तप्रकोपणं}}{\Afootnote{A1938 H; \skt{वातपित्तप्रकोपनं ।}\ N}}\skt{ ॥}\pend

\pstart\versepar \ledsidenote{\passref{1.46.272}}\skt{दीपनङ्कालशाकं तु गरदोषहरं }\edtext{\skt{लघु}}{\Afootnote{H; \skt{कटु ।}\ A1938; \skt{कटुः ।}\ N}}\skt{ । }\pend
\pstart\versepar \skt{कौसुम्भम्मधुरं }\edtext{\edtext{\skt{रूक्षमुष्णं}}{\lemma{\skt{रूक्षम्}}\Afootnote{N H; \skt{रूक्ष}\ A1938}}}{\lemma{\skt{उष्णं}}\Afootnote{H; \skt{मुष्णं}\ A1938 N}}\skt{ }\edtext{\skt{श्लेष्मकरं}}{\Afootnote{N H; \skt{श्लेष्महरं}\ A1938}}\skt{ }\edtext{\skt{लघु}}{\Afootnote{A1938 H; \skt{लघुः ।}\ N}}\skt{ ॥}\pend

\pstart\versepar \ledsidenote{\passref{1.46.273}}\edtext{\skt{वातघ्नन्नालिकाशाकम्पित्तघ्नम्मधुरञ्च}}{\lemma{\skt{वातघ्नन्}}\Afootnote{N H; \skt{वातलं}\ A1938}}\skt{ तत् । }\pend
\pstart\versepar \skt{ग्रहण्यर्शोविकारघ्नी }\edtext{\skt{साम्ला}}{\Afootnote{A1938 H; \skt{सम्व्लोः}\ N}}\skt{ }\edtext{\skt{वातकफे}}{\Afootnote{A1938 H; \skt{पित्तकफे}\ N}}\skt{ }\edtext{\skt{हिताः}}{\Afootnote{A1938 H; \skt{लाः ।}\ N}}\skt{ ॥ }\pend
\pstart\versepar \skt{}\edtext{\skt{उष्णाः}}{\Afootnote{H; \skt{उष्णा}\ A1938 N}}\skt{ कषायमधुरा चाङ्गेरी }\edtext{\skt{चापि}}{\Afootnote{N H; \skt{चाग्नि}\ A1938}}\skt{ }\edtext{\skt{दीपनी}}{\Afootnote{A1938 H; \skt{दीपिनी ।}\ N}}\skt{ ।}\edtext{}{\lemma{\skt{तत् । ग्रहण्यर्शोविकारघ्नी}}\Afootnote{A1938 H; \skt{तत् ॥ ग्रहण्योशं विकारघ्नी}\ N}}\pend

\pstart\versepar \ledsidenote{\passref{1.46.274}}\edtext{\skt{त्रिपर्ण्णीलोणिकापीलुपर्ण्णीपत्तूरजीवकाः}}{\Afootnote{H; \skt{लोणिकाजातुकत्रिपर्णिकापत्तूरजीवक}\ A1938; \skt{त्रिपर्ण्णिलोणिकापीलुपर्ण्णीपत्तूरजीवकः ।}\ N}}\skt{ ॥ }\pend
\pstart\versepar \skt{}\edtext{\skt{सुवर्च्चलासञ्जणककुटुम्बककुर्टिञ्जराः}}{\Afootnote{H; \skt{सुवर्चलाडुडुरककुतुम्बककुठिञ्जरकुन्तलिकाकुरण्टिकाप्रभृतयः ॥२७४॥}\ A1938; \skt{सुवर्च्चलाझुणकुटुम्बककुटिजराः ।}\ N\textsuperscript{ac}; \skt{सुवर्च्चलाझुणककुटुम्बककुटिजराः ।}\ N\textsuperscript{pc}}}\skt{ ।}\pend

\pstart\versepar \ledsidenote{\passref{1.46.275}}\skt{स्वादुपाकरसाः शीताः कफघ्ना नातिपित्तलाः ॥ }\pend
\pstart\versepar \skt{लवणानुरसा रूक्षाः सक्षारा वातलाः सराः ॥}\pend

\pstart\versepar \ledsidenote{\passref{1.46.276}}\edtext{\skt{स्वादुतिक्ताः}}{\Afootnote{N H; \skt{स्वादुतिक्ता}\ A1938}}\skt{ }\edtext{\skt{कुनलिका}}{\Afootnote{N H; \skt{कुन्तलिका}\ A1938}}\skt{ }\edtext{\skt{कषायाः}}{\Afootnote{N H; \skt{कषाया}\ A1938}}\skt{ }\edtext{\skt{सकुरूटिका}}{\Afootnote{H\textsuperscript{pc}; \skt{सकुरण्टिका}\ A1938; \skt{सकुराटिका ।}\ N; \skt{सकुटिका ॥}\ H\textsuperscript{ac}}}\skt{ ॥ }\pend
\pstart\versepar \skt{}\edtext{\skt{संग्राहि}}{\Afootnote{A1938 N\textsuperscript{ac} H; \skt{संग्राही}\ N\textsuperscript{pc}}}\skt{ शीतलं चापि लघु }\edtext{\skt{दोषाविरोधिकः}}{\Afootnote{N H; \skt{दोषापहं तथा ।}\ A1938}}\skt{ । }\pend
\pstart\versepar \skt{राजक्षवकशाकं तु }\edtext{\skt{शटीशाकं तु}}{\Afootnote{H; \skt{शटीशाकं च}\ A1938; \skt{सटीशाकन्तु}\ N}}\skt{ तद्विधं ॥}\pend

\pstart\versepar \ledsidenote{\passref{1.46.277}}\skt{स्वादुपाकरसं }\edtext{\skt{शीतं}}{\Afootnote{N H; \skt{शाकं}\ A1938}}\skt{ दुर्ज्जरं हरिमन्थजं । }\pend
\pstart\versepar \skt{भेदनं }\edtext{\skt{रूक्षमधुरं}}{\Afootnote{N H; \skt{मधुरं}\ A1938}}\skt{ }\edtext{\skt{कलायमतिवातलं}}{\Afootnote{H; \skt{रूक्षं कालायमतिवातलम् ॥२७७॥}\ A1938; \skt{कालायमतिवातलं ।}\ N}}\skt{ ॥}\pend

\pstart\versepar \ledsidenote{\passref{1.46.278}}\edtext{\skt{भेदनङ्कटुकम्पाके}}{\lemma{\skt{भेदनङ्}}\Afootnote{N H; \skt{स्रंसनं}\ A1938}}\skt{ }\edtext{\edtext{\skt{कफघ्नमनिलापहं}}{\lemma{\skt{कफघ्नम्}}\Afootnote{H; \skt{लघु वा}\ A1938; \skt{कफग्नम्}\ N}}}{\lemma{\skt{अनिलापहं}}\Afootnote{N H; \skt{तकफापहम् ।}\ A1938}}\skt{ । }\pend
\pstart\versepar \skt{}\edtext{\edtext{\skt{शोफघ्नमुष्णवीर्यञ्च}}{\lemma{\skt{शोफघ्नम्}}\Afootnote{N H; \skt{शोफघ्न}\ A1938}}}{\lemma{\skt{उष्णवीर्यञ्च}}\Afootnote{N H; \skt{मुष्णवीर्यं च}\ A1938}}\skt{ }\edtext{\skt{पत्रम्पूतिकरञ्जकं}}{\lemma{\skt{पूतिकरञ्जकं}}\Afootnote{H; \skt{पूतिकरञ्जजम्}\ A1938 N}}\skt{ ॥}\pend

\pstart\versepar \ledsidenote{\passref{1.46.279}}\skt{ताम्बूलपत्रं }\edtext{\edtext{\skt{कटुकन्तीक्ष्णोष्णम्पित्तकोपणं}}{\lemma{\skt{कटुकन्}}\Afootnote{N H; \emph{om.}\ A1938}}}{\lemma{\skt{पित्तकोपणं}}\Afootnote{H; \skt{कटु पित्तप्रकोपणम् ।}\ A1938; \skt{पित्तकोपनं ।}\ N}}\skt{ । }\pend
\pstart\versepar \skt{}\edtext{\skt{तिक्तं}}{\Afootnote{N H; \emph{om.}\ A1938}}\skt{ सुगन्धि }\edtext{\skt{विषदं}}{\Afootnote{N H; \skt{विशदं}\ A1938}}\skt{ }\edtext{\skt{स्वर्यम्वातकफापहम्}}{\lemma{\skt{स्वर्यम्}}\Afootnote{N H; \skt{तिक्तं स्वर्यं}\ A1938}}\skt{ ॥}\pend

\pstart\versepar \ledsidenote{\passref{1.46.280}}\skt{स्रन्सनङ्कटुकम्पाके कषायम्वह्निदीपनं । }\pend
\pstart\versepar \skt{}\edtext{\skt{वक्त्रकण्डूगलक्लेददौर्ग्गन्ध्यादिविनाशनं}}{\Afootnote{H; \skt{वक्त्रकण्डूमलक्लेददौर्गन्ध्यादिविशोधनम् ॥२८०॥}\ A1938; \skt{वक्त्रकण्डूमलक्लेददौर्गन्ध्यादिविनाशआनं ॥}\ N\textsuperscript{ac}; \skt{वक्त्रकण्डूमलक्लेददौर्गन्ध्यादिविनाशनं ॥}\ N\textsuperscript{pc}}}\skt{ ॥}\pend

\pstart \ledsidenote{\passref{1.46.281}}\edtext{\skt{कोविदारशणशाल्मलीपुष्पाणि}}{\Afootnote{H; \skt{अथ पुष्पवर्गः । कोविदारशणशाल्मलीपुष्पाणि मधुराणि}\ A1938; \skt{कोविदारशरणशाल्मलीपुष्पाणि मधुराणि}\ N\textsuperscript{ac}; \skt{कोविदारशणशाल्मलीपुष्पाणि मधुराणि}\ N\textsuperscript{pc}}}\skt{ मधुरविपाकानि रक्तपित्तहराणि च । }\edtext{\skt{वृषागस्तिकयोः}}{\Afootnote{N H; \skt{वृषागस्त्ययोः}\ A1938}}\skt{ पुष्पाणि तिक्तानि }\edtext{\skt{कटुपाकीनि}}{\Afootnote{H; \skt{कटुविपाकानि}\ A1938; \skt{कटुपाकानि ।}\ N}}\skt{ क्षयकासापहानि च ।}\pend

\pstart \ledsidenote{\passref{1.46.289cd}}\edtext{\skt{मधुशिग्रुकरीराणि}}{\Afootnote{A1938 N; \skt{मधुसिग्रुकरीराणि}\ H}}\skt{ कटूनि श्लेष्महराणि }\edtext{\skt{चेति}}{\Afootnote{N H; \skt{च ॥२८९॥}\ A1938}}\skt{ ।}\pend

\pstart \ledsidenote{\passref{1.46.290}}\edtext{\skt{क्षवककुलेवरवंशकरीरप्रभृतीनि}}{\Afootnote{A1938 H; \skt{क्षवककुलेवरवंसकरीरप्रभृतीनि}\ N}}\skt{ }\edtext{\skt{कफपित्तहराणि}}{\Afootnote{H; \skt{कफकराणि}\ A1938; \skt{कफपित्तहराणि ं}\ N}}\skt{ }\edtext{\skt{सृष्टमूत्रपुरीषाणि}}{\Afootnote{H; \skt{सृष्टमूत्रपुरीषाणि च ॥२९०॥}\ A1938; \skt{ष्टमूत्रपुरीषाणि ।}\ N}}\skt{ ॥}\pend

\pstart\versepar \ledsidenote{\passref{1.46.291}}\edtext{\skt{क्षवकङ्क्रिमिलन्तेषु}}{\lemma{\skt{क्रिमिलन्तेषु}}\Afootnote{N H; \skt{कृमिलं तेषु}\ A1938}}\skt{ स्वादुपाकं }\edtext{\skt{सपिच्छिलं}}{\Afootnote{N H; \skt{सपिच्छलम् ।}\ A1938}}\skt{ । }\pend
\pstart\versepar \skt{}\edtext{\edtext{\skt{विष्यन्दिवालतन्नातिपित्तश्लेष्मकरञ्च}}{\lemma{\skt{विष्यन्दिवालतन्}}\Afootnote{H; \skt{विस्यन्दि वातलं}\ A1938; \skt{विष्यन्दिवातलं}\ N}}}{\lemma{\skt{नातिश्लेष्मकरञ्च}}\Afootnote{N H; \skt{नातिपित्तश्लेष्मकरं च}\ A1938}}\skt{ तत् ॥}\pend

\pstart\versepar \ledsidenote{\passref{1.46.292}}\skt{वेणोः करीराः }\edtext{\skt{श्लेष्मघ्ना}}{\Afootnote{N H; \skt{कफला}\ A1938}}\skt{ मधुरा रसपाकतः । }\pend
\pstart\versepar \skt{विदाहिनो वातकराः }\edtext{\skt{सकषाया}}{\Afootnote{A1938 H; \skt{सकाषाया}\ N}}\skt{ विरूक्षणाः ॥}\pend

\pstart \ledsidenote{\passref{1.46.293}}\edtext{\skt{उद्भिदानि तु}}{\Afootnote{N H; \skt{उद्भिदानि}\ A1938}}\skt{ }\edtext{\skt{पलालेक्षुकरीषवेणुजातानि}}{\Afootnote{N H; \skt{पलालेक्षुकरीषवेणुक्षितिजानि ।}\ A1938}}\skt{ । तत्र पलालजातम्मधुरं मधुरविपाकं }\edtext{\skt{रूक्षन्दोषकरञ्च}}{\lemma{\skt{दोषकरञ्च}}\Afootnote{N H; \skt{दोषप्रशमनं च ।}\ A1938}}\skt{ । इक्षुजम्मधुरं कषायानुरसं कटुपाकं शीतलञ्च । }\edtext{\edtext{\edtext{\skt{तद्वदेवोष्णङ्करीषम्वेणुजातं}}{\lemma{\skt{तद्वदेवोष्णङ्}}\Afootnote{A1938 H; \skt{तद्वदेवोष्णं}\ N}}}{\lemma{\skt{करीषम्}}\Afootnote{H; \skt{कारीषं}\ A1938 N}}}{\lemma{\skt{वेणुजातं}}\Afootnote{N H; \skt{कषायं वातकोपनं च, (वेणुजातं}\ A1938}}\skt{ कषायम्वातकोपनञ्च । भूमिजं }\edtext{\edtext{\skt{गुरुर्न्नातिवातलमभूमिजश्चास्यानुरसः}}{\lemma{\skt{गुरुर्}}\Afootnote{H; \skt{गुरु}\ A1938 N}}}{\lemma{\skt{न्नातिवातलमभूमिजश्}}\Afootnote{H; \skt{नातिवातलं भूमितश्}\ A1938; \skt{नातिवातलं । अभूमिजम्}\ N}}\skt{ ॥}\pend

\pstart \ledsidenote{\passref{1.46.294}}\edtext{\skt{पिण्याकतिलकल्कस्थूणिकाशुष्कशाकानि}}{\Afootnote{A1938 H; \skt{पिण्याकस्तिलकल्कस्थूनि काशुष्कशाकानि ।}\ N}}\skt{ }\edtext{\skt{सर्व्वदोषप्रकोपनानि}}{\Afootnote{N H; \skt{सर्वदोषप्रकोपणानि ॥२९४॥}\ A1938}}\skt{ ॥}\pend

\pstart\versepar \ledsidenote{\passref{1.46.295}}\skt{विष्टम्भिनः स्मृताः }\edtext{\skt{सर्व्ववटका}}{\Afootnote{H; \skt{सर्वे वटका}\ A1938 N}}\skt{ }\edtext{\skt{वातकोपनाः}}{\Afootnote{A1938 H; \skt{वातप्रकोपनाः ।}\ N}}\skt{ । }\pend
\pstart\versepar \skt{}\edtext{\skt{पिण्याकी}}{\Afootnote{N H; \skt{सिण्डाकी}\ A1938}}\skt{ वातला सार्द्द्रा रुचिष्यानलदीपनी ॥}\pend

\pstart\versepar \ledsidenote{\passref{1.46.296abcd}}\edtext{\skt{विभेदि}}{\Afootnote{N H; \skt{विड्भेदि}\ A1938}}\skt{ गुरुरूक्षञ्च प्रायो विष्टम्भि }\edtext{\skt{शीतलं}}{\Afootnote{N H; \skt{दुर्जरम् ।}\ A1938}}\skt{ । }\pend
\pstart\versepar \skt{सकषायञ्च सर्व्वं हि स्वादु }\edtext{\edtext{\skt{शाकमुदाहृतं}}{\lemma{\skt{शाकम्}}\Afootnote{H; \skt{शाक}\ A1938 N}}}{\lemma{\skt{उदाहृतं}}\Afootnote{H; \skt{मुदाहृतम्}\ A1938 N}}\skt{ ॥}\pend

\pstart \ledsidenote{\passref{1.46.313}}\edtext{\skt{सैन्धवसामुद्रविडसौवर्च्चलरोमकोद्भिदप्रभृतीनि}}{\Afootnote{N H; \skt{अथ लवणवर्गः । अथ लवणानि- सैन्धवसामुद्रविडसौवर्चलरोमकौद्भिदप्रभृतीनि}\ A1938}}\skt{ }\edtext{\edtext{\skt{यथोत्तरमुष्णानि}}{\lemma{\skt{यथोत्तरम्}}\Afootnote{N; \skt{यथोत्तर}\ A1938 H}}}{\lemma{\skt{उष्णानि}}\Afootnote{N; \skt{मुष्णानि}\ A1938 H}}\skt{ वातहराणि कफपित्तकराणि ॥ }\edtext{\skt{कटुपाकीनि}}{\Afootnote{A1938 H; \skt{कटुपाकानि}\ N}}\skt{ यथापूर्व्वं स्निग्धानि }\edtext{\skt{स्वादूनि}}{\Afootnote{A1938 H; \skt{स्वादुनि}\ N}}\skt{ सृष्टमूत्रपुरीषाणि चेति ॥}\pend

\pstart\versepar \ledsidenote{\passref{1.46.314}}\edtext{\skt{चक्षुष्यं}}{\Afootnote{A1938 H; \skt{चक्षुष्य}\ N}}\skt{ सैन्धवं हृद्यं रुच्यंलघ्वग्निदीपनं ॥ }\pend
\pstart\versepar \skt{}\edtext{\skt{स्निग्धम्वृष्यं}}{\lemma{\skt{वृष्यं}}\Afootnote{N H; \emph{om.}\ A1938}}\skt{ समधुरं }\edtext{\skt{शीतन्दोषघ्नमुत्तमं}}{\lemma{\skt{शीतन्}}\Afootnote{H; \skt{वृष्यं शीतं}\ A1938; \skt{शीत}\ N}}\skt{ ॥}\pend

\pstart\versepar \ledsidenote{\passref{1.46.315}}\skt{सामुद्रं मधुरम्पाके }\edtext{\edtext{\skt{नात्युष्णमविदाहि}}{\lemma{\skt{नात्युष्णम्}}\Afootnote{N H; \skt{नात्युष्ण}\ A1938}}\skt{ च}}{\lemma{\skt{अविदाहि च}}\Afootnote{N H; \skt{मविदाहि च ।}\ A1938}}\skt{ । }\pend
\pstart\versepar \skt{भेदनं }\edtext{\skt{स्निग्धमीषच्च}}{\Afootnote{A1938 H; \skt{स्निग्धमीषच्च}\ N}}\skt{ }\edtext{\edtext{\skt{शीतघ्नन्नातिपित्तलं}}{\lemma{\skt{शीतघ्नन्}}\Afootnote{N H; \skt{शूलघ्नं}\ A1938}}}{\lemma{\skt{नातिपित्तलं}}\Afootnote{A1938 H; \skt{नातिपित्तलं ॥ ॥ ।…॥}\ N}}\skt{ ॥}\pend

\pstart\versepar \ledsidenote{\passref{1.46.316}}\skt{सक्षारन्दीपनं }\edtext{\skt{शूक्ष्मं}}{\Afootnote{H; \skt{सूक्ष्मं}\ A1938 N}}\skt{ }\edtext{\skt{हृद्रोगकफनाशनं}}{\Afootnote{N H; \skt{शूलहृद्रोगनाशनम् ।}\ A1938}}\skt{ । }\pend
\pstart\versepar \skt{रोचनं }\edtext{\edtext{\skt{तीक्ष्णमुष्णञ्च}}{\lemma{\skt{तीक्ष्णम्}}\Afootnote{N H; \skt{तीक्ष्ण}\ A1938}}}{\lemma{\skt{उष्णञ्च}}\Afootnote{N H; \skt{मुष्णं च}\ A1938}}\skt{ विडम्वातानुलोमनं ॥}\pend

\pstart\versepar \ledsidenote{\passref{1.46.317}}\skt{लघु सौवर्च्चलं पाके }\edtext{\skt{वीर्योष्णम्विषदं}}{\lemma{\skt{विषदं}}\Afootnote{N H; \skt{विशदं}\ A1938}}\skt{ }\edtext{\skt{कटुः}}{\Afootnote{N H; \skt{कटु ।}\ A1938}}\skt{ । }\pend
\pstart\versepar \skt{}\edtext{\skt{गुल्मशूलविवन्धघ्नं}}{\Afootnote{N H; \skt{गुल्मशूलविबन्धघ्नं}\ A1938}}\skt{ हृद्यं }\edtext{\skt{सुरभि}}{\Afootnote{A1938 H; \skt{सुरभिर्}\ N}}\skt{ }\edtext{\skt{रोचनं}}{\Afootnote{A1938 H; \skt{ओचनं ॥}\ N}}\skt{ ॥}\pend

\pstart\versepar \ledsidenote{\passref{1.46.318}}\edtext{\edtext{\skt{रोमकंतीक्ष्णमत्युष्णम्व्यवायि}}{\lemma{\skt{रोमकंतीक्ष्णम्}}\Afootnote{N H; \skt{रोमकं तीक्ष्ण}\ A1938}}}{\lemma{\skt{अत्युष्णम्}}\Afootnote{H; \skt{मत्युष्णं}\ A1938; \skt{अत्युष्णं}\ N}}\skt{ }\edtext{\skt{कटुदीपनं}}{\Afootnote{N H; \skt{कटुपाकि च ।}\ A1938}}\skt{ । }\pend
\pstart\versepar \skt{वातघ्नं }\edtext{\skt{लघुविष्यन्दि}}{\Afootnote{N H; \skt{लघु विस्यन्दि}\ A1938}}\skt{ सूक्ष्मं विड्भेदि मूत्रलं ॥}\pend

\pstart\versepar \ledsidenote{\passref{1.46.319}}\skt{लघुतीक्ष्णोष्णमुत्क्लेदि सूक्ष्मम्वातानुलोमनं । }\pend
\pstart\versepar \skt{}\edtext{\skt{सतिक्तकटुकं}}{\Afootnote{N H; \skt{सतिक्तं कटु स}\ A1938}}\skt{ }\edtext{\edtext{\skt{क्षारम्विद्याल्लवणमुद्भिदं}}{\lemma{\skt{विद्याल्लवणम्}}\Afootnote{N H; \skt{विद्याल्लवण}\ A1938}}}{\lemma{\skt{उद्भिदं}}\Afootnote{N H; \skt{मौद्भिदम् ॥३१९॥}\ A1938}}\skt{ ॥}\pend

\pstart\versepar \ledsidenote{\passref{1.46.320}}\edtext{\skt{कफवातक्रिमिहरं}}{\Afootnote{N H; \skt{कफवातक्रिमिघ्नं च}\ A1938}}\skt{ लेखनम्पित्तकोपनं । }\pend
\pstart\versepar \skt{दीपनं पाचनम्भेदि }\edtext{\skt{लवणङ्गुडिकाह्वयं}}{\lemma{\skt{गुडिकाह्वयं}}\Afootnote{N H; \skt{गुटिकाह्वयम् ॥३२०॥}\ A1938}}\skt{ ॥}\pend

\pstart\versepar \ledsidenote{\passref{1.46.321}}\edtext{\edtext{\skt{ऊषप्रसूतम्वालार्क्कमलमूत्रकरोद्भवं}}{\lemma{\skt{ऊषप्रसूतम्}}\Afootnote{N H; \skt{ऊषसूतं}\ A1938}}}{\lemma{\skt{वालार्क्कमलमूत्रकरोद्भवं}}\Afootnote{H; \skt{वालुकैलं शैलमूलाकरोद्भवम् ।}\ A1938; \skt{वालाकमलमूकाकरोद्भवं ।}\ N}}\skt{ । }\pend
\pstart\versepar \skt{}\edtext{\skt{लवणङ्कटुकं}}{\lemma{\skt{कटुकं}}\Afootnote{A1938 H; \skt{कटुक}\ N}}\skt{ }\edtext{\skt{च्छेदि}}{\Afootnote{A1938 H; \skt{च्छेदि}\ N}}\skt{ }\edtext{\skt{विहिमं}}{\Afootnote{N H; \skt{विहितं}\ A1938}}\skt{ }\edtext{\skt{लघु}}{\Afootnote{N H; \skt{कटु}\ A1938}}\skt{ चोच्यते ॥}\pend

\pstart\versepar \ledsidenote{\passref{1.46.322}}\edtext{\skt{यवेक्षारस्वर्ज्जिकापाचिमटङ्गणक्षाराः}}{\Afootnote{\skt{यवक्षारस्वर्जिकाक्षारोषक्षारपाकिमटङ्कणक्षार}\ A1938; \skt{यवक्षारस्वर्ज्जिकापाचिमटंगणक्षाराः}\ N H}}\skt{ ॥ }\pend
\pstart\versepar \skt{}\edtext{\skt{गुल्मार्शोग्रहणीरोगशर्क्कराश्मरि}}{\Afootnote{H; \skt{प्रभृतयः । गुल्मार्शोग्रहणीदोषशर्कराश्मरि}\ A1938; \skt{गुल्मार्शोग्रहणं रोगशर्क्कराम्परि}\ N}}\skt{ नाशनाः । }\pend
\pstart\versepar \skt{क्षारास्तु }\edtext{\skt{पाचनाः}}{\Afootnote{A1938 H; \skt{पाचना}\ N}}\skt{ सर्व्वे }\edtext{\skt{अग्निदीप्तिकरा}}{\Afootnote{N H; \skt{रक्तपित्तकराः}\ A1938}}\skt{ }\edtext{\skt{स्मराः}}{\Afootnote{H; \skt{सराः ।}\ A1938 N}}\skt{ ॥}\pend

\pstart\versepar \ledsidenote{\passref{1.46.323}}\skt{ज्ञेयौ वह्निसमौ क्षारौ स्वर्ज्जिकायावशूकजौ । }\pend
\pstart\versepar \skt{}\edtext{\skt{शुक्रश्लेष्मविवन्धार्शोगुल्मप्लीहविनाशनौ}}{\Afootnote{H; \skt{शुक्रश्लेष्मविबन्धार्शोगुल्मप्लीहविनाशनौ ॥३२३॥}\ A1938; \skt{शुक्रश्लेष्मविवन्धार्शो गुल्मप्लीहविनाशनौ ।}\ N}}\skt{ ॥}\pend

\pstart\versepar \ledsidenote{\passref{1.46.324}}\edtext{\edtext{\skt{उष्णोऽनिलघ्नः}}{\lemma{\skt{उष्णो}}\Afootnote{A1938 H; \skt{उष्णो}\ N}}}{\lemma{\skt{ऽनिलघ्नः}}\Afootnote{A1938 H; \skt{निलघ्नः}\ N}}\skt{ }\edtext{\skt{प्रक्लेद्यूषक्षारो}}{\Afootnote{H; \skt{प्रक्लेदी चोषक्षारो}\ A1938; \skt{प्रक्लेदद्युषाक्षारो}\ N}}\skt{ }\edtext{\skt{वलनाशनः}}{\Afootnote{N H; \skt{बलापहः ।}\ A1938}}\skt{ । }\pend
\pstart\versepar \skt{मेदोघ्नः }\edtext{\edtext{\skt{पाचिमक्षारस्तेषाम्बस्तिविशोधनः}}{\lemma{\skt{पाचिमक्षारस्तेषाम्}}\Afootnote{H; \skt{पाकिमः क्षारस्तेषां}\ A1938; \skt{पिचम क्षारस्तेषां}\ N}}}{\lemma{\skt{बस्तिविशोधनः}}\Afootnote{A1938; \skt{वस्तिविशोधनः}\ N H}}\skt{ ॥}\pend

\pstart\versepar \ledsidenote{\passref{1.46.325}}\skt{विरूक्षणोऽनिलकरः श्लेष्मघ्नः पित्तदूषणः । }\pend
\pstart\versepar \skt{}\edtext{\skt{अग्निदीप्तिकरस्तीक्ष्णष्टङ्गणक्षार}}{\Afootnote{\skt{अग्निदीप्तिकरस्तीक्ष्णष्टङ्कणः क्षार}\ A1938; \skt{अग्निदीप्तितरस्तीक्ष्णष्टङ्कणक्षार}\ N; \skt{अग्निदीप्तिकरस्तीक्ष्णष्टङ्गणक्षार}\ H}}\skt{ उच्यते ॥}\pend

\pstart\versepar \ledsidenote{\passref{1.46.331}}\skt{धान्येषु मांसेषु फलेषु चैव शाकेषु }\edtext{\skt{चानुक्तमिहाप्रमोहात्}}{\lemma{\skt{चानुक्तम्}}\Afootnote{H; \skt{चानुक्त}\ A1938 N}}\skt{ । }\pend
\pstart\versepar \skt{आस्वादतो }\edtext{\edtext{\skt{भूतगणैर्ग्गृहित्वा}}{\lemma{\skt{भूतगणैर्}}\Afootnote{H; \skt{भूतगुणैश्}\ A1938 N}}\skt{ }\edtext{\edtext{\skt{तमादिशेद्द्रव्यमनल्पबुद्धिः}}{\lemma{\skt{आदिशेद्द्रव्यम्}}\Afootnote{H; \skt{दादिशेद्द्रव्य}\ A1938; \skt{दादिशेद्द्रव्य}\ N}}}{\lemma{\skt{अनल्पबुद्धिः}}\Afootnote{H; \skt{मनल्पबुद्धिः ॥३३१॥}\ A1938; \skt{मनल्पवुद्धिः ॥}\ N}}}{\lemma{\skt{ग्गृहित्वा तम्}}\Afootnote{H; \skt{च मत्वा त}\ A1938; \skt{गृहीत्वा त}\ N}}\skt{ ॥}\edtext{}{\lemma{\skt{इहाप्रमोहात् । आस्वादतो}}\Afootnote{H; \skt{मिहाप्रमेयात् । आस्वादतो}\ A1938; \skt{मिहं प्रमोहात् । आस्वादतो}\ N}}\pend

\pstart\versepar \ledsidenote{\passref{1.46.332}}\edtext{\skt{षष्टिका}}{\Afootnote{A1938 H; \skt{यष्टिका}\ N}}\skt{ }\edtext{\skt{यवगोधूमलोहिता ये च}}{\Afootnote{N H; \skt{यवगोधूमा लोहिता ये च}\ A1938}}\skt{ शालयः । }\pend
\pstart\versepar \skt{मुद्गाढकीमसूराश्च }\edtext{\skt{धन्येषु}}{\Afootnote{H; \skt{धान्येषु}\ A1938 N}}\skt{ प्रवराः स्मृताः ॥}\pend

\pstart\versepar \ledsidenote{\passref{1.46.333}}\edtext{\skt{एणः}}{\Afootnote{N H; \emph{om.}\ A1938}}\skt{ }\edtext{\skt{कुरङ्गो}}{\Afootnote{H; \emph{om.}\ A1938; \skt{कुरंगो}\ N}}\skt{ }\edtext{\edtext{\edtext{\skt{हरिणस्तित्तिरिर्लाव}}{\lemma{\skt{हरिणस्}}\Afootnote{N H; \skt{लाव}\ A1938}}}{\lemma{\skt{तित्तिरिर्}}\Afootnote{H; \skt{तित्तिरि}\ A1938 N}}}{\lemma{\skt{लाव}}\Afootnote{N H; \skt{सारङ्गकुरङ्गैणक}\ A1938}}\skt{ }\edtext{\skt{एव च}}{\Afootnote{N H; \skt{पिञ्जलाः ।}\ A1938}}\skt{ । }\pend
\pstart\versepar \skt{मयूरवर्म्मिकूर्म्माश्च श्रेष्ठा }\edtext{\skt{मांसगणेषु वै}}{\Afootnote{N H; \skt{मांसगणेष्विह ॥३३३॥}\ A1938}}\skt{ ॥}\pend

\pstart\versepar \ledsidenote{\passref{1.46.334}}\skt{दाडिमामलकन्द्राक्षा खर्ज्जूरं सपरूषकं । }\pend
\pstart\versepar \skt{राजादनम्मातुलुङ्गम्फलवर्ग्गे प्रशस्यते ॥}\pend

\pstart\versepar \ledsidenote{\passref{1.46.335}}\edtext{\skt{चुच्चुसतीना}}{\Afootnote{H; \skt{सतीनो}\ A1938; \skt{चुच्चुसतीनो}\ N}}\skt{ }\edtext{\skt{वास्तूके}}{\Afootnote{H; \skt{वास्तुक}\ A1938; \skt{वस्तूक}\ N}}\skt{ }\edtext{\skt{चिल्लीमूलकपोतिका}}{\Afootnote{N H; \skt{श्चुच्चूचिल्लीमूलकपोतिकाः ।}\ A1938}}\skt{ । }\pend
\pstart\versepar \skt{}\edtext{\skt{माण्डूकपर्ण्णी}}{\Afootnote{H; \skt{मण्डूकपर्णी}\ A1938 N}}\skt{ जीवन्ती शाकवर्ग्गे प्रशस्यते ॥}\pend

\pstart\versepar \ledsidenote{\passref{1.46.336}}\skt{गव्यं }\edtext{\skt{क्षीरघृतं}}{\Afootnote{N H; \skt{क्षीरं घृतं}\ A1938}}\skt{ }\edtext{\skt{शस्तं}}{\Afootnote{N H; \skt{श्रेष्ठं,}\ A1938}}\skt{ सैन्धवं लवणेषु च । }\pend
\pstart\versepar \skt{}\edtext{\skt{धात्रीदाडिममम्लेषु}}{\Afootnote{A1938 H; \skt{धात्रीदाडिममम्व्लेषु}\ N}}\skt{ पिप्पली नागरङ्कटौ ॥}\pend

\pstart\versepar \ledsidenote{\passref{1.46.337}}\skt{तिक्ते }\edtext{\skt{पटोलवार्त्ताके}}{\Afootnote{N H; \skt{पटोलवार्ताकं,}\ A1938}}\skt{ }\edtext{\skt{मधुरघृत}}{\Afootnote{H; \skt{मधुरे घृत}\ A1938 N}}\skt{ }\edtext{\skt{उच्यते}}{\Afootnote{N H; \skt{मुच्यते ।}\ A1938}}\skt{ । }\pend
\pstart\versepar \skt{क्षौद्रम्पूगफलं श्रेष्ठङ्कषाये सपरूषकं ॥}\pend

\pstart\versepar \ledsidenote{\passref{1.46.338}}\edtext{\skt{शर्क्करेक्षुविकारेषु}}{\Afootnote{A1938 H; \skt{शर्करेक्षु विपाकेषु}\ N}}\skt{ }\edtext{\skt{पानजातौ}}{\Afootnote{N H; \skt{पाने मध्वा}\ A1938}}\skt{ }\edtext{\skt{सुरासवौ}}{\Afootnote{N H; \skt{सवौ तथा ।}\ A1938}}\skt{ । }\pend
\pstart\versepar \skt{परिसम्वत्सरन्धान्यम्मांसम्वयसि मध्यमे ॥}\pend

\pstart \ledsidenote{\passref{1.46.339cd}}\skt{फलम्पर्यागतं }\edtext{\edtext{\skt{शाकमशुष्कन्तरुणन्नवम्}}{\lemma{\skt{शाकमशुष्कन्}}\Afootnote{A1938 H; \skt{शाकमशुष्कं}\ N}}}{\lemma{\skt{नवम्}}\Afootnote{A1938 H; \skt{नवं ॥ {\symfont ❈} ॥}\ N}}\skt{ ॥}\pend

\pstart\versepar \ledsidenote{\passref{1.46.340}}\edtext{\skt{अतः}}{\Afootnote{H; \skt{अथ कृतान्नवर्गः । अतः}\ A1938; \skt{अथः}\ N}}\skt{ }\edtext{\skt{पदम्प्रवक्ष्यामि}}{\lemma{\skt{पदम्}}\Afootnote{H; \skt{परं}\ A1938 N}}\skt{ कृतान्नगुणविस्तरं । }\pend
\pstart\versepar \skt{}\edtext{\skt{लाजामण्डो}}{\Afootnote{H; \skt{लाजमण्डो}\ A1938; \skt{लाजमण्डौ}\ N}}\skt{ }\edtext{\skt{विशुद्धानाम्पथ्यः}}{\lemma{\skt{पथ्यः}}\Afootnote{A1938 H; \skt{पथ्य}\ N}}\skt{ पाचनदीपनः ॥}\pend

\pstart\versepar \ledsidenote{\passref{1.46.341}}\edtext{\skt{वातानुलोमनो}}{\Afootnote{A1938 H; \skt{वातान्लोमनो}\ N}}\skt{ हृद्यः पिप्पलीनागरायुतः । }\pend
\pstart\versepar \skt{स्वेदाग्निजननी लघ्वी दीपनी }\edtext{\skt{बस्तिशोधनी}}{\Afootnote{A1938; \skt{वस्तिशोधनी ॥}\ N H}}\skt{ ॥}\pend

\pstart\versepar \ledsidenote{\passref{1.46.342}}\edtext{\skt{क्षुत्तृट्छ्रमग्लानिहरा}}{\Afootnote{H; \skt{क्षुत्तृट्श्रमग्लानिहरी}\ A1938; \skt{क्षुत्तृमग्लानिहरा}\ N}}\skt{ }\edtext{\skt{पेयो}}{\Afootnote{H; \skt{पेया}\ A1938 N}}\skt{ वातानुलोमनी । }\pend
\pstart\versepar \skt{विलेपी तर्प्पणी हृद्या ग्राहिणी }\edtext{\skt{वलवर्द्धनी}}{\Afootnote{N H; \skt{बलवर्धनी ॥३४२॥}\ A1938}}\skt{ ॥}\pend

\pstart\versepar \ledsidenote{\passref{1.46.343}}\skt{पथ्या }\edtext{\skt{स्वादुरसा}}{\Afootnote{N H; \skt{स्यादुरसा}\ A1938}}\skt{ लघ्वी दीपनी क्षुत्तृषापहा । }\pend
\pstart\versepar \skt{हृद्या सन्तर्प्पणी वृष्या }\edtext{\skt{वृंहणी}}{\Afootnote{N H; \skt{बृंहणी}\ A1938}}\skt{ }\edtext{\skt{बलवर्द्धनी}}{\Afootnote{A1938 N; \skt{वलवर्द्धनी ॥}\ H}}\skt{ ॥}\pend

\pstart \ledsidenote{\passref{1.46.344}}\edtext{\skt{शाकमांसफलैर्युक्ता}}{\lemma{\skt{शाकमांसफलैर्}}\Afootnote{A1938 H; \skt{शाकमांसफल}\ N}}\skt{ }\edtext{\edtext{\skt{विलेप्योऽन्याश्च}}{\lemma{\skt{विलेप्यो}}\Afootnote{N H; \skt{विलेप्य}\ A1938}}}{\lemma{\skt{ऽन्याश्च}}\Afootnote{H; \skt{म्ला च}\ A1938; \skt{न्यान्यास्च}\ N}}\skt{ }\edtext{\skt{दुर्ज्जराः}}{\Afootnote{H; \skt{दुर्जरा । सिक्थैर्विरहितो मण्डः पेया सिक्थसमन्विता ॥३४४॥}\ A1938; \skt{दुर्जरा ॥}\ N}}\skt{ ।}\pend

\pstart \ledsidenote{\passref{1.46.345}}\edtext{\skt{विष्टम्भी}}{\Afootnote{N H; \skt{विलेपी बहुसिक्था स्याद्यवागूर्विरलद्रवा । विष्टम्भी}\ A1938}}\skt{ पायसो }\edtext{\skt{वल्य}}{\Afootnote{H; \skt{बल्यो}\ A1938; \skt{वल्यः}\ N}}\skt{ मेदः कफकरो गुरुः ॥}\pend

\pstart\versepar \ledsidenote{\passref{1.46.346}}\skt{कफपित्तकरी }\edtext{\skt{वल्या}}{\Afootnote{N H; \skt{बल्या}\ A1938}}\skt{ कृशरानिलनाशनी । }\pend
\pstart\versepar \skt{धौतस्तु विमलः }\edtext{\skt{शुद्धौ}}{\Afootnote{H; \skt{शुद्धो}\ A1938 N}}\skt{ मनोज्ञः सुरभिः समः ॥}\pend

\pstart\versepar \ledsidenote{\passref{1.46.347}}\edtext{\skt{स्निन्नः}}{\Afootnote{H; \skt{स्विन्नः}\ A1938 N}}\skt{ }\edtext{\skt{सुप्रस्रुतस्तूष्णो}}{\Afootnote{A1938 H; \skt{सुप्रस्रुतस्तूष्णो}\ N}}\skt{ }\edtext{\edtext{\skt{विशदश्चोदनो}}{\lemma{\skt{विशदश्}}\Afootnote{H; \skt{विशद}\ A1938; \skt{विशडस्}\ N}}}{\lemma{\skt{चोदनो}}\Afootnote{H; \skt{स्त्वोदनो}\ A1938; \skt{त्वोदनो}\ N}}\skt{ लघुः । }\pend
\pstart\versepar \skt{}\edtext{\skt{अधौतो}}{\Afootnote{A1938 H; \skt{धौतो}\ N}}\skt{ }\edtext{\skt{प्रश्रुतो}}{\Afootnote{H; \skt{ऽप्रस्रुतो}\ A1938 N}}\skt{ स्विन्नः }\edtext{\edtext{\skt{शीतश्चाप्योदनो}}{\lemma{\skt{चाप्य्}}\Afootnote{A1938 H; \skt{चापो}\ N}}}{\lemma{\skt{ओदनो}}\Afootnote{A1938 H; \skt{दनो}\ N}}\skt{ गुरुः ।}\pend

\pstart\versepar \ledsidenote{\passref{1.46.348}}\skt{लघुः सुगन्धिः कफहा विज्ञेयो भृष्टतण्डुलः । }\pend
\pstart\versepar \skt{}\edtext{\edtext{\skt{स्नैहैर्म्मांसैः}}{\lemma{\skt{स्नैहैर्}}\Afootnote{H; \skt{स्नेहैर्}\ A1938 N}}}{\lemma{\skt{म्मांसैः}}\Afootnote{A1938 H; \skt{मासैः}\ N}}\skt{ फलैः }\edtext{\edtext{\edtext{\skt{स्कन्दैर्व्वैदलैश्चापि}}{\lemma{\skt{स्कन्दैर्}}\Afootnote{H; \skt{कन्दैर्}\ A1938 N}}}{\lemma{\skt{व्वैदलैश्}}\Afootnote{H; \skt{विदला}\ A1938; \skt{वेदलैश्}\ N}}}{\lemma{\skt{चापि}}\Afootnote{N H; \skt{म्लैश्च}\ A1938}}\skt{ }\edtext{\skt{संस्कृताः}}{\Afootnote{H; \skt{संयुताः ॥३४८॥}\ A1938; \skt{संताः ।}\ N}}\skt{ ॥}\pend

\pstart\versepar \ledsidenote{\passref{1.46.349}}\edtext{\edtext{\skt{गुरवोर्वृहणा}}{\lemma{\skt{गुरवोर्}}\Afootnote{H; \skt{गुरवो}\ A1938 N}}}{\lemma{\skt{वृहणा}}\Afootnote{H; \skt{बृंहणा}\ A1938; \skt{वृंहणा}\ N}}\skt{ }\edtext{\skt{वल्या ये च}}{\Afootnote{N H; \skt{बल्या ये च}\ A1938}}\skt{ क्षीरोपसाधिताः । }\pend
\pstart\versepar \skt{}\edtext{\skt{सुस्निन्ना}}{\Afootnote{H; \skt{सुस्विन्नो}\ A1938 N}}\skt{ निस्तुषो }\edtext{\skt{भृष्टो}}{\Afootnote{N H; \skt{भृष्ट}\ A1938}}\skt{ }\edtext{\skt{माषसूपो}}{\Afootnote{N H; \skt{ईषत्सूपो}\ A1938}}\skt{ }\edtext{\skt{लघुर्हितः}}{\Afootnote{A1938 H; \skt{लघुर्हितः ।}\ N}}\skt{ ॥}\pend

\pstart\versepar \ledsidenote{\passref{1.46.350}}\edtext{\skt{स्निन्ननिःपीडितं}}{\Afootnote{H; \skt{स्विन्नं निष्पीडितं}\ A1938; \skt{स्विन्ननिःपीडितं}\ N}}\skt{ }\edtext{\skt{शाकं}}{\Afootnote{A1938 H; \skt{शाकᳵ}\ N}}\skt{ }\edtext{\skt{हितन्तत्स्नैहसंस्कृतं}}{\lemma{\skt{तत्स्नैहसंस्कृतं}}\Afootnote{H; \skt{स्यात्स्नेहसंस्कृतम् ।}\ A1938; \skt{हसंस्कृतं ॥}\ N}}\skt{ । }\pend
\pstart\versepar \skt{}\edtext{\skt{अस्निन्नं}}{\Afootnote{H; \skt{अस्विन्नं}\ A1938 N}}\skt{ }\edtext{\edtext{\edtext{\edtext{\skt{स्नैहरहितमपीडितमतोऽन्यथा}}{\lemma{\skt{स्नैहरहितम्}}\Afootnote{H; \skt{स्नेहरहित}\ A1938; \skt{स्नेहरहितम्}\ N}}}{\lemma{\skt{अपीडितम्}}\Afootnote{H; \skt{मपीडित}\ A1938; \skt{अपीडित}\ N}}}{\lemma{\skt{अतो}}\Afootnote{H; \skt{मतो}\ A1938; \skt{मनो}\ N}}}{\lemma{\skt{ऽन्यथा}}\Afootnote{A1938 H; \skt{न्येषा ॥}\ N}}\skt{ ॥}\pend

\pstart \ledsidenote{\passref{1.46.351}}\edtext{\skt{स्नेहगोरसधान्याम्लफलाम्लकटुकैः}}{\Afootnote{H; \skt{मांसं स्वभावतो वृष्यं स्नेहनं बलवर्धनम् । स्नेहगोरसधान्याम्लफलाम्लकटुकैः}\ A1938; \skt{स्नेहगोरसधान्याम्व्लफला सुकटुकैः}\ N}}\skt{ सह ।}\pend

\pstart\versepar \ledsidenote{\passref{1.46.352}}\skt{सिद्धं }\edtext{\skt{मांसं}}{\Afootnote{A1938 H; \skt{मांस}\ N}}\skt{ }\edtext{\skt{हितम्वल्यं}}{\lemma{\skt{वल्यं}}\Afootnote{N H; \skt{बल्यं}\ A1938}}\skt{ }\edtext{\skt{वृंहणं}}{\Afootnote{H; \skt{रोचनं}\ A1938; \skt{वृहण}\ N}}\skt{ }\edtext{\skt{रोचनं}}{\Afootnote{N H; \skt{बृंहणं}\ A1938}}\skt{ }\edtext{\skt{लघु}}{\Afootnote{H; \skt{गुरु ।}\ A1938; \skt{लघुः ।}\ N}}\skt{ ॥ }\pend
\pstart\versepar \skt{तदेव }\edtext{\edtext{\skt{गोरसादानसुरभिर्द्द्रव्यसंस्कृतं}}{\lemma{\skt{गोरसादानसुरभिर्}}\Afootnote{H; \skt{गोरसादानं सुरभि}\ A1938; \skt{गोरसादानसुरभि}\ N}}}{\lemma{\skt{द्द्रव्यसंस्कृतं}}\Afootnote{A1938 H; \skt{द्रव्यं संस्कृतां ।}\ N}}\skt{ ।}\pend

\pstart\versepar \ledsidenote{\passref{1.46.353}}\edtext{\skt{विद्यात्पित्तकफोद्रेकि}}{\Afootnote{N H; \skt{विद्यात्पित्तकफोद्रेकी}\ A1938}}\skt{ }\edtext{\skt{वलमांसाग्निवर्द्धनं}}{\Afootnote{N H; \skt{बलमांसाग्निवर्धनम् ।}\ A1938}}\skt{ ॥ }\pend
\pstart\versepar \skt{परिशुष्कं स्थिरं स्निग्धं }\edtext{\skt{हर्षणं}}{\Afootnote{A1938 H; \skt{हर्षं}\ N}}\skt{ प्रीणनं }\edtext{\skt{गुरुः}}{\Afootnote{N H; \skt{गुरु ॥३५३॥}\ A1938}}\skt{ ।}\pend

\pstart\versepar \ledsidenote{\passref{1.46.354}}\edtext{\edtext{\skt{रोचनम्वलमेधाग्निर्म्मांसोजःशुक्रवर्द्धनः}}{\lemma{\skt{वलमेधाग्निर्}}\Afootnote{H; \skt{बलमेधाग्नि}\ A1938; \skt{वलमेधाग्नि}\ N}}}{\lemma{\skt{म्मांसोजःशुक्रवर्द्धनः}}\Afootnote{H; \skt{मांसौजःशुक्रवर्धनम् ।}\ A1938; \skt{मांसोजः शुक्रवर्द्धनं ।}\ N}}\skt{ ॥ }\pend
\pstart\versepar \skt{}\edtext{\skt{तदेवोलुप्तपिष्टत्वादुल्लुप्तमिति}}{\lemma{\skt{तदेवोलुप्तपिष्टत्वादुल्लुप्तम्}}\Afootnote{H; \skt{तदेवोल्लुप्तपिष्टत्वादुल्लुप्तम्}\ A1938; \skt{तदेवोप्लुप्तपिष्टत्वादुल्लुप्तम्}\ N}}\skt{ पाचकाः ।}\pend

\pstart\versepar \ledsidenote{\passref{1.46.355}}\skt{परिशुष्कगुणैर्युक्तं }\edtext{\skt{वह्नेः}}{\Afootnote{H; \skt{वह्नौ}\ A1938; \skt{वन्हेः}\ N}}\skt{ }\edtext{\skt{पथ्यतमं}}{\Afootnote{H; \skt{पक्वमतो}\ A1938; \skt{पथ्यातमं}\ N}}\skt{ }\edtext{\skt{गुरुः}}{\Afootnote{N H; \skt{लघु ।}\ A1938}}\skt{ ॥ }\pend
\pstart\versepar \skt{तदेव }\edtext{\skt{शूलिकाप्रोतमङ्गारपरिपाचितं}}{\Afootnote{A1938 H; \skt{शूलिकाप्रोक्तमङ्गारपरिपाचितं ।}\ N}}\skt{ ।}\pend

\pstart\versepar \ledsidenote{\passref{1.46.356}}\skt{ज्ञेयं गुरुतरं }\edtext{\skt{किञ्चित्प्रतप्तं}}{\Afootnote{N H; \skt{किञ्चित्प्रदिग्धं}\ A1938}}\skt{ }\edtext{\skt{कन्दुपाकतः}}{\Afootnote{N H; \skt{गुरुपाकतः ।}\ A1938}}\skt{ ॥ }\pend
\pstart\versepar \skt{उल्लुप्तम्भर्ज्जितम्पिष्टं प्रतप्तं कन्दुपाचितं ।}\pend

\pstart\versepar \ledsidenote{\passref{1.46.357}}\skt{परिशुष्कम्प्रदिग्धञ्च }\edtext{\skt{शूलं}}{\Afootnote{N H; \skt{शूल्यं}\ A1938}}\skt{ }\edtext{\skt{यच्चान्यदीदृशम्}}{\Afootnote{H; \skt{यच्चान्यदीदृशम् ।}\ A1938; \skt{यच्चान्यदीदृश ।}\ N}}\skt{ ॥ }\pend
\pstart\versepar \skt{}\edtext{\skt{मांसं}}{\Afootnote{A1938 H; \skt{मांस}\ N}}\skt{ }\edtext{\skt{यत्तैलसिद्धन्तु}}{\Afootnote{N H; \skt{यत्तैलसिद्धं त}\ A1938}}\skt{ }\edtext{\edtext{\skt{वीर्योष्णम्पित्तकृद्गुरु}}{\lemma{\skt{वीर्योष्णम्}}\Afootnote{N H; \skt{द्वीर्योष्णं}\ A1938}}}{\lemma{\skt{पित्तकृद्गुरु}}\Afootnote{H; \skt{पित्तकृद्गुरु ॥३५७॥}\ A1938; \skt{पित्तकृद्गुरुः ।}\ N}}\skt{ ।}\pend

\pstart\versepar \ledsidenote{\passref{1.46.358}}\skt{लघ्वग्निदीपनं }\edtext{\skt{हृद्यं}}{\Afootnote{A1938 H; \skt{हृद्य}\ N}}\skt{ }\edtext{\skt{संस्कृतं}}{\Afootnote{H; \skt{रुच्यं दृ}\ A1938; \skt{न्संस्कृतं}\ N}}\skt{ }\edtext{\skt{त्वक्प्रसादनम्}}{\Afootnote{H; \skt{ष्टिप्रसादनम् ।}\ A1938; \skt{त्वत्प्रसादनं ।}\ N}}\skt{ ॥ }\pend
\pstart\versepar \skt{अनुष्णवीर्यम्पित्तघ्नम्मनोज्ञं }\edtext{\skt{घृतसाधितम्}}{\Afootnote{A1938 H; \skt{घृतसाधितᳵ ।}\ N}}\skt{ ।}\pend

\pstart\versepar \ledsidenote{\passref{1.46.359}}\skt{प्रीणनः }\edtext{\skt{प्राणजननः}}{\Afootnote{A1938 H; \skt{प्राणजनः}\ N}}\skt{ श्वासकासक्षयापहः ॥ }\pend
\pstart\versepar \skt{}\edtext{\skt{रक्तपित्तश्रमहरो}}{\Afootnote{H; \skt{वातपित्तश्रमहरो}\ A1938; \skt{रक्तपित्तप्रशमनो}\ N}}\skt{ हृद्यो मांसरसः स्मृतः ।}\pend

\pstart\versepar \ledsidenote{\passref{1.46.360}}\edtext{\skt{स्मृत्योजःस्वरहीनानां}}{\Afootnote{A1938 H; \skt{स्मृत्योजस्वरहीनानां}\ N}}\skt{ }\edtext{\skt{ज्वरक्षीणक्षतौजसां}}{\Afootnote{H; \skt{ज्वरक्षीणक्षतोरसाम् ।}\ A1938; \skt{ज्वरक्षीणक्षतौरुशां ।}\ N}}\skt{ ॥ }\pend
\pstart\versepar \skt{भग्नविश्लिष्टसन्धीनां }\edtext{\edtext{\skt{कृशानामल्परेतसां}}{\lemma{\skt{कृशानाम्}}\Afootnote{N H; \skt{कृशाना}\ A1938}}}{\lemma{\skt{अल्परेतसां}}\Afootnote{N H; \skt{मल्परेतसाम् ॥३६०।}\ A1938}}\skt{ ।}\pend

\pstart\versepar \ledsidenote{\passref{1.46.361}}\skt{आप्यायनः संहननः शुक्रदो बलवर्द्धनः ॥ }\pend
\pstart\versepar \skt{स दाडिमयुतो वृष्यः }\edtext{\skt{संस्कृतो}}{\Afootnote{A1938 H; \skt{सस्कृतो}\ N}}\skt{ दोषनाशनः ।}\edtext{}{\lemma{\skt{बलवर्द्धनः ॥ स}}\Afootnote{A1938 H; \skt{वलवर्द्धनः । स}\ N}}\pend

\pstart\versepar \ledsidenote{\passref{1.46.363}}\edtext{\edtext{\skt{यन्मांसन्निःसृतरसन्न}}{\lemma{\skt{यन्मांसन्}}\Afootnote{N H; \skt{यन्मांस}\ A1938}}}{\lemma{\skt{निःसृतरसन्न}}\Afootnote{N H; \skt{मुद्धृतरसं न}\ A1938}}\skt{ }\edtext{\skt{तत्पुष्टिवलावहं}}{\Afootnote{N H; \skt{तत्पुष्टिबलावहम् ।}\ A1938}}\skt{ ॥ }\pend
\pstart\versepar \skt{विष्टम्भि दुर्ज्जरं रूक्षं }\edtext{\skt{विरसम्मारुतावहं}}{\lemma{\skt{मारुतावहं}}\Afootnote{A1938 H; \skt{मारुतावहᳵ ॥}\ N}}\skt{ ।}\pend

\pstart \ledsidenote{\passref{1.46.366}}\edtext{\skt{कफघ्नो}}{\Afootnote{N H; \skt{वेसवारो गुरुः स्निग्धो बल्यो वातरुजापहः । कफघ्नो}\ A1938}}\skt{ दीपनो हृद्यः शुद्धानाम्व्रणिनामपि ॥}\pend

\pstart\versepar \ledsidenote{\passref{1.46.367}}\skt{ज्ञेयः }\edtext{\skt{पथ्यतमश्चापि}}{\lemma{\skt{चापि}}\Afootnote{N H; \skt{चैव}\ A1938}}\skt{ }\edtext{\skt{मुद्गयूषः}}{\Afootnote{A1938 H; \skt{मुद्गयूष}\ N}}\skt{ कृताकृतः । }\pend
\pstart\versepar \skt{स तु }\edtext{\skt{दाडिममृद्वीकायुक्तः}}{\Afootnote{A1938 H; \skt{दाषिममृद्वीकायुक्तः}\ N}}\skt{ }\edtext{\skt{स्याद्रागवाडवः}}{\Afootnote{H; \skt{स्याद्रागषाडवः ॥३६७॥}\ A1938; \skt{स्याद्रागषाडवः ।}\ N}}\skt{ ॥}\edtext{}{\lemma{\skt{कृताकृतः । स तु}}\Afootnote{A1938 H; \skt{न्कृताकृतः । स तु}\ N}}\pend

\pstart\versepar \ledsidenote{\passref{1.46.368}}\skt{रुचिष्यो लघुपाकश्च दोषाणाम्चाविरोधकृत् ॥ }\pend
\pstart\versepar \skt{मसूरमुद्गगोधूमकुलत्थलवणैः }\edtext{\skt{कृतं}}{\Afootnote{N H; \skt{कृतः ॥३६८॥}\ A1938}}\skt{ ।}\edtext{}{\lemma{\skt{चाविरोधकृत् ॥ मसूरमुद्गगोधूमकुलत्थलवणैः}}\Afootnote{A1938 H; \skt{चाविरोधकृत् । मसूरमुद्गगोधूम कुलत्थलवणै}\ N}}\pend

\pstart\versepar \ledsidenote{\passref{1.46.369}}\skt{कफपित्ताविरोधी }\edtext{\skt{स्याद्वातव्याधौ}}{\Afootnote{N H; \skt{स्याद्वातव्याधौ}\ A1938}}\skt{ }\edtext{\skt{प्रशस्यते}}{\Afootnote{N H; \skt{च शस्यते ।}\ A1938}}\skt{ ॥ }\pend
\pstart\versepar \skt{मृद्वीका }\edtext{\skt{दाडिमयुतः स}}{\Afootnote{N H; \skt{दाडिमैर्युक्तः स}\ A1938}}\skt{ }\edtext{\edtext{\skt{चाप्युक्तोऽनिलार्द्दिते}}{\lemma{\skt{उक्तो}}\Afootnote{A1938 N; \skt{उक्तो}\ H}}}{\lemma{\skt{ऽनिलार्द्दिते}}\Afootnote{A1938; \skt{निलार्दिते}\ N H}}\skt{ ।}\pend

\pstart\versepar \ledsidenote{\passref{1.46.370}}\skt{रोचनो दीपनो हृद्यो }\edtext{\edtext{\skt{लघुपाक्युपदिश्यते}}{\lemma{\skt{लघुपाक्य्}}\Afootnote{N H; \skt{लघुपा}\ A1938}}}{\lemma{\skt{उपदिश्यते}}\Afootnote{N H; \skt{क्युपदिश्यते ।}\ A1938}}\skt{ ॥ }\pend
\pstart\versepar \skt{}\edtext{\skt{पटोलनिम्बयूषौ तु}}{\Afootnote{A1938 H; \skt{पटोलनिम्वयूषौ तु}\ N}}\skt{ }\edtext{\skt{कफमेदो}}{\Afootnote{H; \skt{कफमेदो}\ A1938 N}}\skt{ }\edtext{\skt{विशोषणौ}}{\Afootnote{H; \skt{विशोषिणौ}\ A1938 N}}\skt{ ।}\pend

\pstart\versepar \ledsidenote{\passref{1.46.371}}\skt{पित्तघ्नौ दीपनौ हृद्यौ }\edtext{\skt{क्रिमिकुष्ठज्वरापहौ}}{\Afootnote{N H; \skt{कृमिकुष्ठज्वरापहौ ।}\ A1938}}\skt{ ॥ }\pend
\pstart\versepar \skt{}\edtext{\skt{श्वासकासप्रतिस्याय}}{\Afootnote{H; \skt{श्वासकासप्रतिश्याय}\ A1938 N}}\skt{ प्रसेकारोचकज्वरान् ।}\pend

\pstart\versepar \ledsidenote{\passref{1.46.372}}\skt{हन्तिमूलकयूषस्तु कफमेदोगलग्रहान् ॥ }\pend
\pstart\versepar \skt{}\edtext{\skt{कुलत्थयूषोऽनिलहा}}{\lemma{\skt{ऽनिलहा}}\Afootnote{A1938 H; \skt{ऽनिला}\ N}}\skt{ }\edtext{\skt{शर्क्कराश्मरिनाशनः}}{\Afootnote{N H; \skt{श्वासपीनसनाशनः ॥३७२॥}\ A1938}}\skt{ ।}\pend

\pstart\versepar \ledsidenote{\passref{1.46.373}}\edtext{\skt{तूनीप्रतूनिकासादगुल्ममेदः}}{\Afootnote{H; \skt{तूणीप्रतूणीकासाश्मगुल्ममेदः}\ A1938; \skt{तूनीप्रतुनिकासादगुल्ममेदः}\ N}}\skt{ }\edtext{\skt{कफापहः}}{\Afootnote{N H; \skt{कफावहः ।}\ A1938}}\skt{ ॥ }\pend
\pstart\versepar \skt{दाडिमामलकैर्यूषो हृद्यः संशमनो लघुः ।}\pend

\pstart\versepar \ledsidenote{\passref{1.46.374}}\skt{प्राणाग्निजननो मूर्च्छा }\edtext{\skt{मदघ्नः}}{\Afootnote{H; \skt{मेदोघ्नो}\ A1938; \skt{मेदघ्नः}\ N}}\skt{ पित्तवातजित् ॥ }\pend
\pstart\versepar \skt{मुद्गामलकयूषस्तु ग्राही }\edtext{\skt{पित्तकफापहः}}{\Afootnote{H; \skt{पित्तकफे हितः ॥३७४॥}\ A1938; \skt{पित्तकफोपहाः ॥}\ N}}\skt{ ।}\edtext{}{\lemma{\skt{पित्तवातजित् ॥ मुद्गामलकयूषस्तु}}\Afootnote{N H; \skt{वातपित्तजित् । मुद्गामलकयूषस्तु}\ A1938}}\pend

\pstart\versepar \ledsidenote{\passref{1.46.375}}\edtext{\skt{यवकोलकुलत्थानां}}{\Afootnote{A1938 H; \skt{…कोलकुलत्थानां}\ N}}\skt{ यूषः }\edtext{\edtext{\skt{कण्ठ्योऽनिलापहः}}{\lemma{\skt{कण्ठ्यो}}\Afootnote{A1938 H; \skt{कण्ठो}\ N}}}{\lemma{\skt{ऽनिलापहः}}\Afootnote{A1938 H; \skt{निलापहः ।}\ N}}\skt{ ॥ }\pend
\pstart\versepar \skt{}\edtext{\edtext{\skt{सर्व्वधान्यकृतस्तद्वद्बृंहणः}}{\lemma{\skt{सर्व्वधान्यकृतस्}}\Afootnote{N H; \skt{सर्वधान्यकृत}\ A1938}}}{\lemma{\skt{तद्वद्बृंहणः}}\Afootnote{H; \skt{स्तद्वद्बृंहणः}\ A1938; \skt{तद्वद्वृंहणः}\ N}}\skt{ प्राणवर्द्धनः ।}\pend

\pstart\versepar \ledsidenote{\passref{1.46.376}}\edtext{\skt{खलकावलिकौ}}{\Afootnote{H; \skt{खडकाम्बलिकौ}\ A1938; \skt{खलुकवलिकौ}\ N}}\skt{ }\edtext{\skt{द्यौ}}{\Afootnote{H; \skt{हृद्यौ}\ A1938; \skt{हृद्यौर्}\ N}}\skt{ }\edtext{\skt{च्छेदीवातकफे}}{\Afootnote{H; \skt{तथा वातकफे}\ A1938; \skt{च छर्दीवातकफे}\ N}}\skt{ हितौ ॥ }\pend
\pstart\versepar \skt{}\edtext{\skt{वल्यः}}{\Afootnote{N H; \skt{बल्यः}\ A1938}}\skt{ कफानिलौ हन्ति }\edtext{\edtext{\skt{दाडिमाम्लोऽग्निदीपनः}}{\lemma{\skt{दाडिमाम्लो}}\Afootnote{A1938 H; \skt{दाडिमम}\ N}}}{\lemma{\skt{ऽग्निदीपनः}}\Afootnote{A1938 H; \skt{ष्टौग्निदीपनः ।}\ N}}\skt{ ।}\pend

\pstart\versepar \ledsidenote{\passref{1.46.377}}\edtext{\skt{धान्याम्लो}}{\Afootnote{H\textsuperscript{pc}; \skt{धान्याम्व्लो}\ N}}\skt{ दीपनो हृद्यः }\edtext{\skt{पित्तकृद्वातनाशनः}}{\Afootnote{H\textsuperscript{pc}; \skt{पित्तकृद्वाभिनाशनः ।}\ N}}\skt{ ॥ }\pend
\pstart\versepar \skt{}\edtext{\skt{दध्यम्लः}}{\Afootnote{A1938 H; \skt{दध्यम्व्लः}\ N}}\skt{ }\edtext{\skt{कफकृद्वल्यः}}{\Afootnote{N H; \skt{कफकृद्बल्यः}\ A1938}}\skt{ }\edtext{\skt{पित्तलो}}{\Afootnote{H; \skt{स्निग्धो}\ A1938; \skt{पित्तला}\ N}}\skt{ }\edtext{\skt{वातहा}}{\Afootnote{H; \skt{वातहरो}\ A1938; \skt{वाततहा}\ N}}\skt{ गुरुः । }\pend
\pstart\versepar \skt{}\edtext{\skt{तक्राम्लः}}{\Afootnote{A1938 H; \skt{तक्राम्व्लः}\ N}}\skt{ पित्तकृत्प्रोक्तो विषरक्तप्रदूषणः ॥}\edtext{}{\lemma{\skt{… द}}\Afootnote{\emph{om.}\ H\textsuperscript{ac}}}\edtext{}{\lemma{\skt{… द}}\Afootnote{\emph{om.}\ A1938}}\pend

\pstart \ledsidenote{\passref{1.46.381}}\edtext{\skt{तिलपिण्याकविकृतं}}{\Afootnote{N H; \skt{दधिमस्त्वम्लसिद्धस्तु यूषः काम्बलिकः स्मृतः । तिलपिण्याकविकृतिः}\ A1938}}\skt{ }\edtext{\edtext{\skt{शुष्कशोकम्विरूक्षणं}}{\lemma{\skt{शुष्कशोकम्}}\Afootnote{H; \skt{शुष्कशाकं}\ A1938 N}}}{\lemma{\skt{विरूक्षणं}}\Afootnote{H; \skt{विरूढकम् ॥३८१॥}\ A1938; \skt{विरूणां ।}\ N}}\skt{ ।}\pend

\pstart \ledsidenote{\passref{1.46.382}}\edtext{\skt{शाण्डकी तु}}{\Afootnote{H; \skt{सिण्डाकी च}\ A1938; \skt{शाण्डाकी च}\ N}}\skt{ }\edtext{\skt{गुरूणि स्युः}}{\Afootnote{A1938 H; \skt{गुरुणि स्युः}\ N}}\skt{ कफपित्तकराणि }\edtext{\skt{च}}{\Afootnote{N H; \skt{च । तद्वच्च वटकान्याहुर्विदाहीनि गुरूणि च ॥३८२॥}\ A1938}}\skt{ ॥}\pend

\pstart\versepar \ledsidenote{\passref{1.46.379}}\edtext{\skt{अस्नेहलवणं}}{\Afootnote{A1938 H; \skt{अहलवणं}\ N}}\skt{ }\edtext{\edtext{\skt{सर्व्वमकृतङ्कटुकैर्व्विना}}{\lemma{\skt{सर्व्वम्}}\Afootnote{N H; \skt{सर्व}\ A1938}}}{\lemma{\skt{अकृतङ्}}\Afootnote{N H; \skt{मकृतं}\ A1938}}\skt{ । }\pend
\pstart\versepar \skt{}\edtext{\skt{विज्ञेयङ्कटुकस्नेहलवणैः}}{\lemma{\skt{कटुकस्नेहलवणैः}}\Afootnote{N H; \skt{लवणस्नेहकटुकैः}\ A1938}}\skt{ }\edtext{\skt{संयुतङ्कृतं}}{\lemma{\skt{संयुतङ्}}\Afootnote{A1938 N\textsuperscript{pc} H; \skt{संयु क्ततं}\ N\textsuperscript{ac}}}\skt{ ॥}\pend

\pstart\versepar \ledsidenote{\passref{1.46.380}}\edtext{\edtext{\skt{यूषम्विद्यात्फलाम्लैस्तु}}{\lemma{\skt{यूषम्}}\Afootnote{H; \emph{om.}\ A1938; \skt{यूषां}\ N}}}{\lemma{\skt{विद्यात्फलाम्लैस्तु}}\Afootnote{H; \skt{अथ गोरस}\ A1938; \skt{विद्यात्फलाम्व्लैस्तु}\ N}}\skt{ }\edtext{\skt{धान्याम्लेनाम्लितञ्च}}{\Afootnote{H; \skt{धान्याम्लफलाम्लैरम्लितं च}\ A1938; \skt{धान्याम्व्लनाम्व्लितश्च}\ N}}\skt{ यत् । }\pend
\pstart\versepar \skt{यथोत्तरं }\edtext{\skt{गुरु}}{\Afootnote{N H; \skt{लघु}\ A1938}}\skt{ }\edtext{\skt{तथा}}{\Afootnote{N H; \skt{हितं}\ A1938}}\skt{ संस्कृतासंस्कृतं रसं ॥}\pend

\pstart\versepar \ledsidenote{\passref{1.46.383}}\skt{लघवो }\edtext{\skt{बृंहणो}}{\Afootnote{H; \skt{बृंहणा}\ A1938; \skt{वृहणा}\ N}}\skt{ }\edtext{\skt{वृष्यो}}{\Afootnote{H; \skt{वृष्या}\ A1938 N}}\skt{ हृद्या }\edtext{\skt{रोचनदीपनाः}}{\Afootnote{N H; \skt{रोचनदीपनाः ।}\ A1938}}\skt{ ॥ }\pend
\pstart\versepar \skt{}\edtext{\skt{भ्रममूर्च्छातृषाच्छर्द्दि}}{\Afootnote{N H; \skt{तृष्णामूर्च्छाभ्रमच्छर्दि}\ A1938}}\skt{ श्रमघ्ना }\edtext{\skt{रागषाडवाः}}{\Afootnote{A1938 H; \skt{रागशाडवाः ।}\ N}}\skt{ ॥}\pend

\pstart\versepar \ledsidenote{\passref{1.46.384}}\skt{रसाला }\edtext{\skt{रोचना}}{\Afootnote{N H; \skt{बृंहणी}\ A1938}}\skt{ }\edtext{\skt{वल्या}}{\Afootnote{N H; \skt{बल्या}\ A1938}}\skt{ स्निग्धा }\edtext{\skt{वृष्याथ}}{\Afootnote{H; \skt{वृष्या च}\ A1938; \skt{वृष्याथो}\ N}}\skt{ }\edtext{\skt{वृंहणी}}{\Afootnote{H; \skt{रोचनी ।}\ A1938; \skt{वृंहणीः ॥}\ N}}\skt{ । }\pend
\pstart\versepar \skt{स्नेहनं गुडसंयुक्तं हृद्यन्दध्यनिलापहं ॥}\pend

\pstart\versepar \ledsidenote{\passref{1.46.385}}\edtext{\skt{शक्तवः}}{\Afootnote{H; \skt{सक्तवः}\ A1938 N}}\skt{ सर्प्पिषाभ्यक्ताः शीतवारिपरिप्लुताः । }\pend
\pstart\versepar \skt{}\edtext{\skt{नात्यच्छो}}{\Afootnote{N H; \skt{नातिद्रवा}\ A1938}}\skt{ }\edtext{\skt{नातिसान्द्रश्च}}{\Afootnote{N H; \skt{नातिसान्द्रा}\ A1938}}\skt{ मन्थ }\edtext{\skt{इत्यभिधीयते}}{\lemma{\skt{अभिधीयते}}\Afootnote{N H; \skt{उपदिश्यते ॥३८५॥}\ A1938}}\skt{ ॥}\pend

\pstart\versepar \ledsidenote{\passref{1.46.386}}\skt{मन्थः }\edtext{\skt{सद्योवलच्छर्द्दि}}{\Afootnote{N H; \skt{सद्योबलकरः}\ A1938}}\skt{ }\edtext{\skt{पिपासास्रमनाशनः}}{\Afootnote{H; \skt{पिपासाश्रमनाशनः ।}\ A1938; \skt{पिपासाभ्रमनाशनः ॥}\ N}}\skt{ । }\pend
\pstart\versepar \skt{}\edtext{\skt{साम्लं}}{\Afootnote{H; \skt{साम्ल}\ A1938; \skt{साम्व्लं}\ N}}\skt{ स्नेहगुडो मूत्रकृच्छ्रोदावर्त्तनाशनः ॥}\pend

\pstart\versepar \ledsidenote{\passref{1.46.387}}\edtext{\skt{शर्क्करेक्षुरसद्राक्षायुक्तः}}{\Afootnote{A1938 H; \skt{शर्क्करेक्षेरसद्राक्षायुक्तः}\ N}}\skt{ पित्तविकारनुत् । }\pend
\pstart\versepar \skt{द्राक्षामधुसमायुक्तः }\edtext{\skt{कफरोगहरः}}{\Afootnote{H; \skt{कफरोगनिबर्हणः ॥३८७॥}\ A1938; \skt{कफरोगहरहरस्}\ N}}\skt{ }\edtext{\skt{स्मृतः}}{\Afootnote{N H; \emph{om.}\ A1938}}\skt{ ॥}\edtext{}{\lemma{\skt{पित्तविकारनुत् । द्राक्षामधुसमायुक्तः}}\Afootnote{N H; \skt{पित्तविकारनुत् । द्राक्षामधूदकयुतः}\ A1938}}\pend

\pstart\versepar \ledsidenote{\passref{1.46.388}}\skt{वर्ग्गत्रयेणोपहितो मलदोषानुलोमनः । }\pend
\pstart\versepar \skt{}\edtext{\skt{गौडमम्लमनम्लम्वा}}{\Afootnote{A1938 H; \skt{गौडमम्व्लमनम्व्लं वा}\ N}}\skt{ पानकं गुरुमूत्रलं ॥}\pend

\pstart\versepar \ledsidenote{\passref{1.46.389}}\skt{तदेव खण्डमृद्वीका शर्क्करासहितम्पुनः । }\pend
\pstart\versepar \skt{}\edtext{\skt{साम्लं}}{\Afootnote{A1938 H; \skt{साम्व्लं}\ N}}\skt{ सतीक्ष्णं सहिमं }\edtext{\skt{पानकं }\edtext{\skt{स्यान्निरत्ययं}}{\lemma{\skt{निरत्ययं}}\Afootnote{A1938 H; \skt{निरत्ययं ।}\ N}}}{\lemma{\skt{पानकं स्यान्}}\Afootnote{A1938 H; \skt{पानकस्यान्}\ N}}\skt{ ॥}\pend

\pstart\versepar \ledsidenote{\passref{1.46.390}}\edtext{\skt{माद्धीकन्तु}}{\Afootnote{H; \skt{मार्द्वीकं तु}\ A1938 N}}\skt{ श्रमहरं मूर्च्छा }\edtext{\skt{दाहतृषापहं}}{\Afootnote{A1938 H; \skt{दाहतृषापह ॥}\ N}}\skt{ । }\pend
\pstart\versepar \skt{}\edtext{\skt{फरूषकानां}}{\Afootnote{H; \skt{परूषकाणां}\ A1938 N}}\skt{ कोलानां हृद्यम्विष्टम्भिपानकं ॥}\pend

\pstart\versepar \ledsidenote{\passref{1.46.391}}\edtext{\skt{द्रव्यसंयोगसंस्कारं}}{\Afootnote{A1938 H; \skt{द्रव्यं संयोगसंस्कारं}\ N}}\skt{ ज्ञात्वा मात्राञ्च सर्व्वतः । }\pend
\pstart\versepar \skt{पानकानां यथायोगं }\edtext{\skt{गुरुलाघवमादिशेत्}}{\lemma{\skt{आदिशेत्}}\Afootnote{H; \skt{आदिशेत् ॥३९१॥ इति कृतान्नवर्गः ।}\ A1938; \skt{आदिशेत् ॥ {\symfont ❈} ॥}\ N}}\skt{ ॥}\pend

\pstart \ledsidenote{\passref{1.46.392}}\edtext{\edtext{\skt{वक्ष्याम्यतः}}{\lemma{\skt{वक्ष्याम्य्}}\Afootnote{N; \skt{अथ भक्ष्यवर्गः । वक्ष्या}\ A1938; \skt{वक्ष्या}\ H}}}{\lemma{\skt{अतः}}\Afootnote{N; \skt{म्यतः}\ A1938 H}}\skt{ परं }\edtext{\skt{कृत्स्नान्रसवीर्यविपाकतः}}{\Afootnote{\skt{भक्ष्यान्रसवीर्यविपाकतः । भक्ष्याः क्षीरकृता बल्या वृष्या हृद्याः सुगन्धिनः ॥३९२॥}\ A1938; \skt{कृत्स्नान्रसवीर्यविपाकतः भक्ष्याः क्षीरकृतावल्या वृष्या हृद्याः सुगन्धिनः}\ N H}}\skt{ ।}\pend

\pstart\versepar \ledsidenote{\passref{1.46.393}}\skt{भक्ष्याः क्षीरकृतावल्या वृष्या हृद्याः सुगन्धिनः । }\pend
\pstart\versepar \skt{अविदाहिनः पुष्टिकरा दीपनाः }\edtext{\skt{पित्तनाशनाः}}{\Afootnote{\skt{पित्तनाशनाः । तेषां प्राणकरा हृद्या घृतपूराः कफावहाः ॥३९३॥}\ A1938; \skt{पित्तनाशनाः तेषाम्प्राणकरा हृद्याः घृतपूराः कफापहाः}\ N H}}\skt{ ॥}\edtext{}{\lemma{\skt{… अ}}\Afootnote{\emph{om.}\ H}}\edtext{}{\lemma{\skt{… अ}}\Afootnote{\emph{om.}\ N}}\edtext{}{\lemma{\skt{… अ}}\Afootnote{\emph{om.}\ A1938}}\pend

\pstart\versepar \ledsidenote{\passref{1.46.394}}\skt{तेषां प्राणकरा हृद्या घृतपूराः कफापहाः । }\pend
\pstart\versepar \skt{वातपित्तहरा वृष्या गुरवो }\edtext{\skt{रक्तमांसलाः}}{\Afootnote{\skt{रक्तमांसलाः । बृंहणा गौडिका भक्ष्या गुरवोऽनिलनाशनाः ॥३९४॥}\ A1938; \skt{रक्तमांसलाः । बृंहणा गैडिका भक्ष्या गुररवोऽनिलनाशनाः ॥}\ N\textsuperscript{ac}; \skt{रक्तमांसलाः । बृंहणा गैडिका भक्ष्या गुरवोऽनिलनाशनाः ॥}\ N\textsuperscript{pc}; \skt{रक्तमांसलाः ॥ वृंहणा गौडिका भक्ष्या गुरवोऽनिलनाशनाः ।}\ H}}\skt{ ॥}\edtext{}{\lemma{\skt{… वा}}\Afootnote{\emph{om.}\ A1938}}\edtext{}{\lemma{\skt{… वा}}\Afootnote{\emph{om.}\ H}}\edtext{}{\lemma{\skt{… वा}}\Afootnote{\emph{om.}\ N\textsuperscript{ac}}}\edtext{}{\lemma{\skt{… वा}}\Afootnote{\emph{om.}\ N\textsuperscript{pc}}}\pend

\pstart\versepar \ledsidenote{\passref{1.46.395}}\skt{वृंहणा गौडिका भक्ष्या गुरवोऽनिलनाशनाः । }\pend
\pstart\versepar \skt{अविदाहिनः }\edtext{\skt{पित्तसहाः}}{\Afootnote{N H; \skt{पित्तहराः}\ A1938}}\skt{ शुक्रलाः }\edtext{\skt{कफवर्द्धनाः}}{\Afootnote{\skt{कफवर्धनाः मधुशीर्षकसंयावाः पूपा ये ते विशेषतः}\ A1938 H; \skt{कफवर्द्धनाः । मधुशीर्षकसंयावाः पू यौ ते विशेषतः ॥}\ N\textsuperscript{ac}; \skt{कफवर्द्धनाः । मधुशीर्षकसंयावाः पूवी ये ते विशेषतः ॥}\ N\textsuperscript{pc}}}\skt{ ॥}\edtext{}{\lemma{\skt{… अ}}\Afootnote{\emph{om.}\ N\textsuperscript{ac}}}\edtext{}{\lemma{\skt{… अ}}\Afootnote{\emph{om.}\ A1938}}\edtext{}{\lemma{\skt{… अ}}\Afootnote{\emph{om.}\ H}}\edtext{}{\lemma{\skt{… अ}}\Afootnote{\emph{om.}\ N\textsuperscript{pc}}}\pend

\pstart\versepar \ledsidenote{\passref{1.46.396}}\skt{मधुशीर्षकसंयावाः पूपा ये ते विशेषतः । }\pend
\pstart\versepar \skt{गुरवो }\edtext{\skt{वृंहणाश्चैव}}{\lemma{\skt{वृंहणाश्}}\Afootnote{H; \skt{बृंहणाश्}\ A1938 N}}\skt{ मोदकास्तु }\edtext{\skt{सुदुर्ज्जराः}}{\Afootnote{\skt{सुदुर्जराः रोचनो दीपनः स्वर्यः पित्तघ्नः पवनापहः}\ A1938 N H}}\skt{ ॥}\edtext{}{\lemma{\skt{… गु}}\Afootnote{\emph{om.}\ H}}\edtext{}{\lemma{\skt{… गु}}\Afootnote{\emph{om.}\ A1938}}\edtext{}{\lemma{\skt{… गु}}\Afootnote{\emph{om.}\ N}}\pend

\pstart\versepar \ledsidenote{\passref{1.46.397}}\skt{रोचनो दीपनः स्वर्यः पित्तघ्नः पवनापहः । }\pend
\pstart\versepar \skt{}\edtext{\skt{गुरुर्मृष्टतमश्चैव}}{\lemma{\skt{मृष्टतमश्}}\Afootnote{N H; \skt{र्मृष्टतरमश्}\ A1938}}\skt{ }\edtext{\skt{सद्यकः}}{\Afootnote{N H; \skt{सट्टकः}\ A1938}}\skt{ }\edtext{\skt{प्राणवर्द्धनः}}{\Afootnote{\skt{प्राणवर्धनः । हृद्यः सुगन्धिर्मधुरः स्निग्धः कफकरो गुरुः ॥३९७॥}\ A1938; \skt{प्राणवर्द्धनः । हृद्यः सुगन्धिमधुरः स्निग्धः श्लेष्मकरो गुरुः ॥}\ N; \skt{प्राणवर्द्धनः ॥ हृद्यः सुगन्धमधुरः स्निग्धः श्लेष्मकरो गुरुः ।}\ H}}\skt{ ॥}\edtext{}{\lemma{\skt{… गु}}\Afootnote{\emph{om.}\ H}}\edtext{}{\lemma{\skt{… गु}}\Afootnote{\emph{om.}\ N}}\edtext{}{\lemma{\skt{… गु}}\Afootnote{\emph{om.}\ A1938}}\pend

\pstart\versepar \ledsidenote{\passref{1.46.398}}\skt{हृद्यः सुगन्धमधुरः स्निग्धः श्लेष्मकरो गुरुः । }\pend
\pstart\versepar \skt{}\edtext{\skt{पित्तपहस्तृप्तिकरो}}{\lemma{\skt{तृप्तिकरो}}\Afootnote{N H; \skt{स्तृप्तिकरो}\ A1938}}\skt{ }\edtext{\skt{वल्यो}}{\Afootnote{H; \skt{बल्यो}\ A1938 N}}\skt{ }\edtext{\skt{विष्यन्दि}}{\Afootnote{N H; \skt{विष्यन्द}\ A1938}}\skt{ }\edtext{\skt{उच्यते}}{\Afootnote{\skt{उच्यते बृंहणा वातपित्तघ्ना भक्ष्या बल्यास्तु सामिताः ॥३९८॥}\ A1938; \skt{उच्यते । बृंहणा वातपित्तघ्ना भक्ष्या बल्यास्तु सम्मिताः ॥}\ N; \skt{उच्यते ॥ वृङ्ङणा वातपित्तघ्ना भक्ष्या वल्यास्तु सम्मिताः ।}\ H}}\skt{ ॥}\edtext{}{\lemma{\skt{… पि}}\Afootnote{\emph{om.}\ H}}\edtext{}{\lemma{\skt{… पि}}\Afootnote{\emph{om.}\ A1938}}\edtext{}{\lemma{\skt{… पि}}\Afootnote{\emph{om.}\ N}}\pend

\pstart\versepar \ledsidenote{\passref{1.46.399}}\skt{वृङ्ङणा वातपित्तघ्ना भक्ष्या वल्यास्तु सम्मिताः । }\pend
\pstart\versepar \skt{हृद्याः पथ्यतमास्तेषां लघवः }\edtext{\skt{फेनकादयः}}{\Afootnote{\skt{फेनकादयः । मुद्गादिवेसवाराणां पूर्णा विष्टम्भिनो मताः ॥३९९॥}\ A1938; \skt{फेनकादयः मुद्गादिवेशवारैश्च पूर्ण्णा विष्टम्भिनो मताः}\ N H}}\skt{ ॥}\edtext{}{\lemma{\skt{… हृ}}\Afootnote{\emph{om.}\ A1938}}\edtext{}{\lemma{\skt{… हृ}}\Afootnote{\emph{om.}\ H}}\edtext{}{\lemma{\skt{… हृ}}\Afootnote{\emph{om.}\ N}}\pend

\pstart\versepar \ledsidenote{\passref{1.46.400}}\skt{मुद्गादिवेशवारैश्च पूर्ण्णा विष्टम्भिनो मताः । }\pend
\pstart\versepar \skt{वेशवारैः }\edtext{\skt{सपिसितैः}}{\Afootnote{H; \skt{सपिशितैः}\ A1938 N}}\skt{ सम्पूर्ण्णा }\edtext{\skt{गुरुवृङ्हणाः}}{\Afootnote{\skt{गुरुबृंहणाः पाललाः श्लेष्मजननाः, शष्कुल्यः कफपित्तलाः ॥४००॥}\ A1938; \skt{गुरुबृंहणाः । पालालाः श्लेष्मजननाः षष्कुल्यः कफपित्तलाः ॥}\ N; \skt{गुरुवृङ्हणाः ॥ पालालाः श्लेष्मजननाः षष्कुल्यः कफपित्तलाः ।}\ H}}\skt{ ॥}\edtext{}{\lemma{\skt{… वे}}\Afootnote{\emph{om.}\ H}}\edtext{}{\lemma{\skt{… वे}}\Afootnote{\emph{om.}\ A1938}}\edtext{}{\lemma{\skt{… वे}}\Afootnote{\emph{om.}\ N}}\pend

\pstart\versepar \ledsidenote{\passref{1.46.401}}\skt{पालालाः श्लेष्मजननाः षष्कुल्यः कफपित्तलाः । }\pend
\pstart\versepar \skt{वीर्योष्णाः }\edtext{\skt{पैत्तिकाः}}{\Afootnote{H; \skt{पैष्टिका}\ A1938; \skt{पैष्टकाः}\ N}}\skt{ भक्ष्याः }\edtext{\skt{कफपित्तप्रकोपणाः}}{\Afootnote{\skt{कफपित्तप्रकोपणाः विदाहिनो नातिबला गुरवश्च विशेषतः}\ A1938 N; \skt{कफपित्तप्रकोपणाः ॥ विदाहिनो नातिवला गुरवश्च विशेषतः ।}\ H}}\skt{ ॥}\edtext{}{\lemma{\skt{… वी}}\Afootnote{\emph{om.}\ H}}\edtext{}{\lemma{\skt{… वी}}\Afootnote{\emph{om.}\ A1938}}\edtext{}{\lemma{\skt{… वी}}\Afootnote{\emph{om.}\ N}}\pend

\pstart\versepar \ledsidenote{\passref{1.46.402}}\skt{विदाहिनो नातिवला गुरवश्च विशेषतः । }\pend
\pstart\versepar \skt{वैदला }\edtext{\skt{गुरवो}}{\Afootnote{N H; \skt{लघवो}\ A1938}}\skt{ भक्ष्याः कषायाः }\edtext{\skt{सृष्टमारुताः}}{\Afootnote{\skt{सृष्टमारुताः । विष्टम्भिनः पित्तसमाः श्लेष्मघ्ना भिन्नवर्चसः ॥४०२॥}\ A1938; \skt{सृष्टमारुताः विष्टम्भिनः श्लेष्महराः पित्तघ्नाभिन्नवर्च्चसः}\ N H}}\skt{ ॥}\edtext{}{\lemma{\skt{… वै}}\Afootnote{\emph{om.}\ H}}\edtext{}{\lemma{\skt{… वै}}\Afootnote{\emph{om.}\ N}}\edtext{}{\lemma{\skt{… वै}}\Afootnote{\emph{om.}\ A1938}}\pend

\pstart \ledsidenote{\passref{1.46.403}}\edtext{\edtext{\edtext{\skt{विष्टम्भिनः}}{\Afootnote{\skt{बल्या वृष्यास्}\ A1938}}\skt{ }\edtext{\skt{श्लेष्महराः}}{\Afootnote{\skt{तु गुरवो}\ A1938}}\skt{ }\edtext{\skt{पित्तघ्नाभिन्नवर्च्चसः}}{\Afootnote{\skt{विज्ञेया माषसाधिताः । कूर्चिकाविकृता भक्ष्या गुरवो नातिपित्तलाः ॥४०३॥}\ A1938}}\skt{ ।}}{\Afootnote{\emph{om.}\ H}}}{\Afootnote{\emph{om.}\ N}}\pend

\pstart \ledsidenote{\passref{1.46.404}}\edtext{\edtext{\edtext{\skt{कुञ्चिताविकृता}}{\Afootnote{\skt{विरूढककृता}\ A1938}}\skt{ भक्ष्या }\edtext{\skt{गुरवोऽनिलपित्तलाः}}{\lemma{\skt{ऽनिलपित्तलाः}}\Afootnote{\skt{ऽनिलपित्तलाः । विदाहोत्क्लेशजनना रूक्षा दृष्टिप्रदूषणाः ॥४०४॥}\ A1938}}\skt{ ॥}}{\Afootnote{\emph{om.}\ H}}}{\Afootnote{\emph{om.}\ N}}\pend

\pstart\versepar \ledsidenote{\passref{1.46.405}}\skt{विदाहः क्लेदजननाः रूक्षा दृष्टिप्रदूषणाः । }\pend
\pstart\versepar \skt{हृद्याः सुगन्धिनो भक्ष्या लघवो }\edtext{\skt{घृतपाचिमाः}}{\Afootnote{\skt{घृतपाचिताः । वातपित्तहरा बल्या वर्णदृष्टिप्रसादनाः ॥४०५॥}\ A1938; \skt{घृतपाचिमाः । वातपित्तहरा बल्या वर्ण्णदृष्टिप्रसादनाः ॥}\ N; \skt{घृतपाचिमाः ॥ वातपित्तहरा वल्या वर्ण्णदृष्टिप्रसादनाः ।}\ H}}\skt{ ॥}\edtext{}{\lemma{\skt{… हृ}}\Afootnote{\emph{om.}\ A1938}}\edtext{}{\lemma{\skt{… हृ}}\Afootnote{\emph{om.}\ N}}\edtext{}{\lemma{\skt{… हृ}}\Afootnote{\emph{om.}\ H}}\pend

\pstart\versepar \ledsidenote{\passref{1.46.406}}\skt{वातपित्तहरा वल्या वर्ण्णदृष्टिप्रसादनाः । }\pend
\pstart\versepar \skt{विदाहिनस्तैलकृता }\edtext{\skt{भक्ष्यास्तु}}{\Afootnote{N H; \skt{गुरवः}\ A1938}}\skt{ }\edtext{\skt{गुरुपाकिनः}}{\Afootnote{\skt{कटुपाकिनः । उष्णा मारुतदृष्टिघ्नाः पित्तलास्त्वक्प्रदूषणाः ॥४०६॥}\ A1938; \skt{गुरुपाकिनः उष्णा मारुतदृष्टिघ्नाः पित्तलासृक्प्रदूषणाः}\ N H}}\skt{ ॥}\edtext{}{\lemma{\skt{… वि}}\Afootnote{\emph{om.}\ H}}\edtext{}{\lemma{\skt{… वि}}\Afootnote{\emph{om.}\ N}}\edtext{}{\lemma{\skt{… वि}}\Afootnote{\emph{om.}\ A1938}}\pend

\pstart\versepar \ledsidenote{\passref{1.46.407}}\skt{उष्णा मारुतदृष्टिघ्नाः पित्तलासृक्प्रदूषणाः । }\pend
\pstart\versepar \skt{फलमांसेषुविकृतीस्तिलमाषोपसंस्कृताः ॥}\edtext{}{\lemma{\skt{… फ}}\Afootnote{\emph{om.}\ H}}\edtext{}{\lemma{\skt{… फ}}\Afootnote{\emph{om.}\ A1938}}\edtext{}{\lemma{\skt{… फ}}\Afootnote{\emph{om.}\ N}}\pend

\pstart\versepar \ledsidenote{\passref{1.46.408}}\edtext{\skt{भक्ष्या}}{\Afootnote{\skt{कपालाङ्गारपक्वास्तु लघवो वातकोपनाः ।}\ A1938}}\skt{ }\edtext{\skt{वल्याथ गुरवो}}{\Afootnote{\skt{सुपक्वास्तनवश्}\ A1938}}\skt{ }\edtext{\skt{वृंहणा हृदयप्रियाः}}{\Afootnote{\skt{चैव भूयिष्ठं लघवो}\ A1938}}\skt{ । }\pend
\pstart\versepar \skt{}\edtext{\skt{कपालाङ्गारपक्वाः स्युः}}{\Afootnote{N H; \skt{मताः ॥४०८॥}\ A1938}}\skt{ }\edtext{\skt{किञ्चिल्लघुतरास्तु}}{\Afootnote{N H; \emph{om.}\ A1938}}\skt{ }\edtext{\skt{ते}}{\Afootnote{N H; \emph{om.}\ A1938}}\skt{ ॥}\edtext{}{\lemma{\skt{… क}}\Afootnote{\emph{om.}\ H}}\edtext{}{\lemma{\skt{… क}}\Afootnote{\emph{om.}\ N}}\pend

\pstart\versepar \ledsidenote{\passref{1.46.409}}\skt{सकिलाटादयो भक्ष्या गुरवः कफवर्धनाः । }\pend
\pstart\versepar \skt{कुल्माषा वातला रूक्षा गुरवो भिन्नवर्च्चसः ॥}\pend

\pstart\versepar \ledsidenote{\passref{1.46.410}}\skt{उदावर्त्तहरो }\edtext{\skt{वाद्यः}}{\Afootnote{N H; \skt{वाट्यः}\ A1938}}\skt{ }\edtext{\skt{प्रतिस्यामेहनाशनः}}{\Afootnote{N H; \skt{कासपीनसमेहनु}\ A1938}}\skt{ । }\pend
\pstart\versepar \skt{}\edtext{\skt{धानालुम्बास्तु}}{\Afootnote{N H; \skt{त्धानोलुम्बास्तु}\ A1938}}\skt{ लघवः }\edtext{\skt{कफमेदोविशेषणाः}}{\Afootnote{H; \skt{कफमेदोविशोषणाः}\ A1938 N}}\skt{ ॥}\pend

\pstart\versepar \ledsidenote{\passref{1.46.411}}\edtext{\skt{शक्तवस्तर्प्पणा}}{\Afootnote{H; \skt{शक्तवो बृंहणा}\ A1938; \skt{सक्तवस्तर्पणा}\ N}}\skt{ }\edtext{\edtext{\skt{हृद्यास्तृष्णापित्तकफापहाः}}{\lemma{\skt{हृद्यास्}}\Afootnote{N H; \skt{वृष्या}\ A1938}}}{\lemma{\skt{तृष्णापित्तकफापहाः}}\Afootnote{N H; \skt{स्तृष्णापित्तकफापहाः ।}\ A1938}}\skt{ । }\pend
\pstart\versepar \skt{पीताः }\edtext{\skt{सद्योवलकरा}}{\Afootnote{H; \skt{सद्योबलकरा}\ A1938 N}}\skt{ भेदिनः पवनापहाः ॥}\pend

\pstart\versepar \ledsidenote{\passref{1.46.412}}\skt{गुर्व्वी पिण्डी खराऽत्यर्थं }\edtext{\skt{लघ्वी सा तु}}{\Afootnote{N H; \skt{लघ्वी सैव}\ A1938}}\skt{ विपर्यया । }\pend
\pstart\versepar \skt{}\edtext{\skt{शक्तूनामाशु}}{\Afootnote{H; \skt{त् । शक्तूनामाशु}\ A1938; \skt{सक्तूनामासु}\ N}}\skt{ जीर्येत मृदुत्वादवलेहिका ॥}\pend

\pstart\versepar \ledsidenote{\passref{1.46.413}}\edtext{\skt{लाजाच्छर्द्यतिसारघ्ना}}{\Afootnote{H; \skt{लाजाश्छर्द्यतिसारघ्ना}\ A1938; \skt{लाजाच्छर्द्यातीसारघ्ना}\ N}}\skt{ दीपनाः कफनाशनाः । }\pend
\pstart\versepar \skt{बल्याः कषायमधुरा }\edtext{\skt{लघवस्तृण्मलापहाः}}{\Afootnote{A1938 H; \skt{लघवस्तृन्मलापहाः ।}\ N}}\skt{ ॥}\pend

\pstart\versepar \ledsidenote{\passref{1.46.415}}\skt{पृथुका गुरवः स्निग्धा }\edtext{\skt{वृंहणाः}}{\Afootnote{H; \skt{बृंहणाः}\ A1938 N}}\skt{ कफवर्द्धनाः । }\pend
\pstart\versepar \skt{}\edtext{\skt{बल्याः}}{\Afootnote{A1938; \skt{बल्या}\ N; \skt{वल्याः}\ H}}\skt{ }\edtext{\skt{सक्षीरभावत्वाद्वातघ्ना}}{\Afootnote{N H; \skt{सक्षीरभावात्तु वातघ्ना}\ A1938}}\skt{ भिन्नवर्च्चसः ॥}\pend

\pstart\versepar \ledsidenote{\passref{1.46.416}}\edtext{\edtext{\skt{सन्धानकृत्पिष्टमामन्ताण्डुलं}}{\lemma{\skt{सन्धानकृत्पिष्टमामन्}}\Afootnote{A1938 H; \skt{सन्धानकृत्पिष्टमामं}\ N}}}{\lemma{\skt{ताण्डुलं}}\Afootnote{N H; \skt{तण्डुलं}\ A1938}}\skt{ कफमेदकृत् । }\pend
\pstart\versepar \skt{सुदुर्ज्जरः स्वादुरसो }\edtext{\edtext{\skt{वृंहणस्तण्डुलो}}{\lemma{\skt{वृंहणस्}}\Afootnote{H; \skt{बृंहण}\ A1938; \skt{बृंहणस्}\ N}}}{\lemma{\skt{तण्डुलो}}\Afootnote{N H; \skt{स्तण्डुलो}\ A1938}}\skt{ नवः ॥}\edtext{}{\lemma{\skt{कफमेदकृत् । सुदुर्ज्जरः}}\Afootnote{N H; \skt{कृमिमेहनुत् । सुदुर्जरः}\ A1938}}\pend

\pstart \ledsidenote{\passref{1.46.417}}\edtext{\skt{द्रव्यसंयोगसंस्कारविकारान्समवेक्ष्य}}{\Afootnote{H; \skt{सन्धानकृन्मेहहरः पुराणस्तण्डुलः स्मृतः । द्रव्यसंयोगसंस्कारविकारान्समवेक्ष्य}\ A1938; \skt{द्रव्यं संयोगसंस्कारविकारां समवेक्ष्य}\ N}}\skt{ }\edtext{\skt{तु}}{\Afootnote{N H; \skt{च ॥४१७॥}\ A1938}}\skt{ ।}\pend

\pstart\versepar \ledsidenote{\passref{1.46.417.add1}}\edtext{\skt{भिषग्यथास्वम्भक्ष्याणामादिशेद्गुरुलाघवं}}{\lemma{\skt{आदिशेद्गुरुलाघवं}}\Afootnote{\skt{आदिशेद्गुरुलाघवं}\ N H}}\skt{ ॥ }\pend
\pstart\versepar \skt{खलाखलयवाग्वश्च }\edtext{\skt{रागषाडवषटूकाः}}{\Afootnote{H; \skt{रागशाडवषट्टकाः ॥}\ N}}\skt{ । }\pend
\pstart\versepar \skt{}\edtext{\skt{पानकानि च}}{\Afootnote{H; \skt{पाणकानि च}\ N}}\skt{ चित्राणि यूषाश्चानेकयोनयः ॥}\pend

\pstart\versepar \ledsidenote{\passref{1.46.417.add2}}\edtext{\skt{कट्वम्लस्वादुलवणा}}{\Afootnote{H; \skt{कट्वम्ब्लस्वादुलवणा}\ N}}\skt{ लघवो ये फलोद्भवाः । }\pend
\pstart\versepar \skt{}\edtext{\skt{एवमादीनि}}{\lemma{\skt{एवम्}}\Afootnote{N\textsuperscript{pc} H; \skt{एवीम्}\ N\textsuperscript{ac}}}\skt{ चान्यानि }\edtext{\skt{क्रियस्ते}}{\Afootnote{H; \skt{क्रियन्ते}\ N}}\skt{ वैद्यवाक्यतः॥}\pend

\pstart\versepar \ledsidenote{\passref{1.46.418}}\edtext{\skt{यदा}}{\Afootnote{N H; \skt{यथा}\ A1938}}\skt{ }\edtext{\edtext{\skt{कारणमासाद्य}}{\lemma{\skt{कारणम्}}\Afootnote{N H; \skt{कारण}\ A1938}}}{\lemma{\skt{आसाद्य}}\Afootnote{N H; \skt{मासाद्य}\ A1938}}\skt{ }\edtext{\skt{भोक्तॄणाञ्च्छन्दतोऽपि}}{\lemma{\skt{भोक्तॄणाञ्}}\Afootnote{H; \skt{भोक्तॄणां}\ A1938; \skt{भोर्क्तॄणां}\ N}}\skt{ वा । }\pend
\pstart\versepar \skt{}\edtext{\edtext{\skt{अनेकद्रव्ययोनित्वाच्छास्त्रतस्तान्विनिर्द्दिशेत्}}{\lemma{\skt{अनेकद्रव्ययोनित्वाच्छास्त्रतस्}}\Afootnote{N H; \skt{अनेकद्रव्ययोनित्वाच्छास्त्रत}\ A1938}}}{\lemma{\skt{तान्विनिर्द्दिशेत्}}\Afootnote{H; \skt{स्तान्विनिर्दिशेत् ॥४१८॥ इति भक्ष्यवर्गः ।}\ A1938; \skt{तान्विनिर्द्दिशेत् ॥ {\symfont ❈} ॥}\ N}}\skt{ ॥ ॥}\pend

\pstart\versepar \ledsidenote{\passref{1.46.419}}\skt{अतः }\edtext{\skt{सर्व्वानुपानानुपदेक्ष्यामः}}{\Afootnote{N H; \skt{सर्वानुपानान्युपदेक्ष्यामः ।}\ A1938}}\skt{ ॥ }\pend
\pstart\versepar \skt{}\edtext{\skt{अम्लेन}}{\Afootnote{A1938 H; \skt{अम्ब्लेन}\ N}}\skt{ }\edtext{\skt{केचिद्विहता}}{\Afootnote{H; \skt{केचिद्विहता}\ A1938; \skt{केचिद्विहिता}\ N}}\skt{ मनुष्या माधुर्ययोगे प्रणयीभवन्ति । }\pend
\pstart\versepar \skt{}\edtext{\skt{तथाम्ल}}{\Afootnote{A1938 H; \skt{तथाम्ब्ल}\ N}}\skt{ }\edtext{\skt{एके}}{\Afootnote{N H; \skt{योगे}\ A1938}}\skt{ मधुरेण तृप्तास्तेषां यथेष्टं प्रवदन्ति पथ्यम् ॥}\pend

\pstart\versepar \ledsidenote{\passref{1.46.420}}\edtext{\skt{शीतोष्णतोयासवमद्ययूषफलाम्लधान्याम्लपयोरसानां}}{\Afootnote{H; \skt{शीतोष्णतोयासवमद्ययूषफलाम्लधान्याम्लपयोरसानाम् ।}\ A1938; \skt{शीतोष्णतोयासवमद्ययूषफलाम्बधान्याम्ब्लपयोरसानां ।}\ N}}\skt{ । }\pend
\pstart\versepar \skt{यस्यानुपानं तु }\edtext{\skt{हितम्भवेत}}{\lemma{\skt{भवेत}}\Afootnote{N H; \skt{भवेद्य}\ A1938}}\skt{ }\edtext{\skt{तस्मै}}{\Afootnote{N H; \skt{त्तस्मै}\ A1938}}\skt{ प्रदेयं त्विह मात्रया तत् ॥}\pend

\pstart \ledsidenote{\passref{1.46.421}}\skt{व्याधिञ्च कालञ्च विभाव्य }\edtext{\skt{धीरो}}{\Afootnote{N H; \skt{धीरैर्}\ A1938}}\skt{ }\edtext{\skt{द्रव्यानि}}{\Afootnote{H; \skt{द्रव्याणि}\ A1938 N}}\skt{ }\edtext{\skt{योज्यानि च}}{\Afootnote{N H; \skt{भोज्यानि च}\ A1938}}\skt{ तानि }\edtext{\skt{तानि}}{\Afootnote{N H; \skt{तानि । सर्वानुपानेषु वरं वदन्ति मैध्यं यदम्भः शुचिभाजनस्थम् ॥४२१॥}\ A1938}}\skt{ ।}\pend

\pstart \ledsidenote{\passref{1.46.422}}\edtext{\skt{संक्षेप}}{\Afootnote{N H; \skt{लोकस्य जन्मप्रभृति प्रशस्तं तोयात्मकाः सर्वरसाश्च दृष्टाः । सङ्क्षेप}\ A1938}}\skt{ }\edtext{\skt{एषोऽभिहितो}}{\lemma{\skt{ऽभिहितो}}\Afootnote{A1938 H; \skt{भिहितो}\ N}}\skt{ }\edtext{\skt{नुपानेष्वतः}}{\lemma{\skt{नुपानेष्व्}}\Afootnote{N H; \skt{ऽनुपानेष्व्}\ A1938}}\skt{ परम्विस्तरतोऽभिधास्ये ॥}\pend

\pstart\versepar \ledsidenote{\passref{1.46.423}}\skt{उष्णोदकानुपानन्तु }\edtext{\edtext{\skt{स्नेहानामथ}}{\lemma{\skt{स्नेहानाम्}}\Afootnote{N H; \skt{स्नेहाना}\ A1938}}}{\lemma{\skt{अथ}}\Afootnote{N H; \skt{मथ}\ A1938}}\skt{ शस्यते । }\pend
\pstart\versepar \skt{ऋते }\edtext{\skt{भल्लातकस्नेहात्तत्र}}{\Afootnote{N H; \skt{भल्लातकस्नेहात्स्नेहात्तौवरकात्तथा ॥४२३॥}\ A1938}}\skt{ }\edtext{\skt{तोयं}}{\Afootnote{N H; \emph{om.}\ A1938}}\skt{ }\edtext{\skt{स्मशीतलं}}{\Afootnote{H; \emph{om.}\ A1938; \skt{सुशीतलं ॥}\ N}}\skt{ ॥}\pend

\pstart\versepar \ledsidenote{\passref{1.46.424}}\edtext{\edtext{\skt{अनुपानम्वदन्त्येके}}{\lemma{\skt{वदन्त्य्}}\Afootnote{N H; \skt{वदन्}\ A1938}}}{\lemma{\skt{एके}}\Afootnote{N H; \skt{त्येके}\ A1938}}\skt{ }\edtext{\skt{तैले}}{\Afootnote{A1938 H; \skt{तैल}\ N}}\skt{ }\edtext{\skt{यूषाम्लकाञ्जिकं}}{\Afootnote{A1938 H; \skt{यूषाम्ब्लकाञ्जिकं ।}\ N}}\skt{ । }\pend
\pstart\versepar \skt{शीतोदकम्माक्षिकस्य पिष्टान्नस्य च सर्व्वशः ॥}\pend

\pstart\versepar \ledsidenote{\passref{1.46.425}}\edtext{\skt{दधिपायसमद्यार्त्तिविषजुष्टे}}{\Afootnote{A1938 H; \skt{दधिपायसमद्यार्त्तिविषर्जुष्टे}\ N}}\skt{ तथैव च । }\pend
\pstart\versepar \skt{}\edtext{\edtext{\edtext{\skt{केचित्पिष्टपयस्याहुरनुपानं}}{\lemma{\skt{केचित्पिष्टपयस्य्}}\Afootnote{N; \skt{केचित्पिष्टमय}\ A1938; \skt{केचित्पिष्टपय}\ H}}}{\lemma{\skt{आहुर्}}\Afootnote{N; \skt{स्याहु}\ A1938 H}}}{\lemma{\skt{अनुपानं}}\Afootnote{N; \skt{रनुपानं}\ A1938 H}}\skt{ सुखोदकं ॥}\pend

\pstart \ledsidenote{\passref{1.46.426}}\edtext{\skt{पयो}}{\Afootnote{A1938 H; \skt{प्र या}\ N}}\skt{ मांसरसो वापि }\edtext{\skt{शालिमुद्गादिभोजिनां}}{\Afootnote{N H; \skt{शालिमुद्गादिभोजिनाम् । युद्धाध्वातपसन्तापविषमद्यरुजासु च ॥४२६॥}\ A1938}}\skt{ ।}\pend

\pstart\versepar \ledsidenote{\passref{1.46.427}}\edtext{\edtext{\edtext{\skt{माषादीनामनुपानन्धान्याम्लन्दधिमस्तु}}{\lemma{\skt{माषादीनाम्}}\Afootnote{N H; \skt{माषादे}\ A1938}}}{\lemma{\skt{अनुपानन्}}\Afootnote{N H; \skt{रनुपानं तु}\ A1938}}}{\lemma{\skt{धान्याम्लन्}}\Afootnote{A1938 H; \skt{धान्याम्ब्लं}\ N}}\skt{ वा ॥ }\pend
\pstart\versepar \skt{मद्यं }\edtext{\skt{मद्योचितानान्तु}}{\Afootnote{A1938 H; \skt{मद्यो चितानान्तु}\ N}}\skt{ सर्व्वमांसेषु पूजितं ।}\pend

\pstart\versepar \ledsidenote{\passref{1.46.428}}\edtext{\edtext{\edtext{\skt{अमद्यपानामुदकम्फलाम्लम्वा}}{\lemma{\skt{अमद्यपानाम्}}\Afootnote{N H; \skt{अमद्यपाना}\ A1938}}}{\lemma{\skt{उदकम्}}\Afootnote{N H; \skt{मुदकं}\ A1938}}}{\lemma{\skt{फलाम्लम्वा}}\Afootnote{A1938 H; \skt{फलाम्ब्लम्वा}\ N}}\skt{ प्रशस्यते ॥ }\pend
\pstart\versepar \skt{क्षीरं }\edtext{\edtext{\skt{घर्म्माध्वभाष्यस्त्रीक्लान्तानाममृतोपमं}}{\lemma{\skt{घर्म्माध्वभाष्यस्त्रीक्लान्तानाम्}}\Afootnote{H; \skt{घर्माध्वभाष्यस्त्रीक्लान्ताना}\ A1938; \skt{घर्माधूभाष्यस्त्रीक्लान्ताना}\ N}}}{\lemma{\skt{अमृतोपमं}}\Afootnote{H; \skt{ममृतोपमम् ॥४२८॥}\ A1938; \skt{मु न्मृतोपमं ।}\ N}}\skt{ ।}\pend

\pstart\versepar \ledsidenote{\passref{1.46.429}}\skt{सुरा }\edtext{\skt{कृशानां}}{\Afootnote{A1938 H; \skt{कृषाणां}\ N}}\skt{ }\edtext{\skt{स्थूलानामनुशस्तम्मधूदकं}}{\lemma{\skt{स्थूलानामनुशस्तम्}}\Afootnote{\skt{स्थूलानामनुपानं}\ A1938; \skt{स्थ्नूलानामनुषस्त}\ N; \skt{स्थूलानामनुशस्तम्}\ H}}\skt{ ॥ }\pend
\pstart\versepar \skt{निरामयानां }\edtext{\skt{चित्रकन्तु}}{\Afootnote{H; \skt{चित्रं तु}\ A1938 N}}\skt{ }\edtext{\skt{फलमध्ये}}{\Afootnote{H; \skt{भुक्तमध्ये}\ A1938; \skt{भक्तमध्ये}\ N}}\skt{ प्रकीर्त्तितं ।}\pend

\pstart\versepar \ledsidenote{\passref{1.46.430}}\skt{स्निग्धोष्णम्मारुते }\edtext{\skt{शस्तं}}{\Afootnote{N H; \skt{पथ्यं,}\ A1938}}\skt{ कफे रूक्षोष्णमिष्यते ॥ }\pend
\pstart\versepar \skt{}\edtext{\skt{अनुपानं}}{\Afootnote{A1938 H; \skt{अनुपान}\ N}}\skt{ }\edtext{\skt{हितञ्चापि}}{\lemma{\skt{हितञ्}}\Afootnote{A1938 H; \skt{ङ्हितञ्}\ N}}\skt{ पित्ते मधुरशीतलं ।}\pend

\pstart \ledsidenote{\passref{1.46.431}}\skt{हितं }\edtext{\skt{शोणितपित्ते तु}}{\Afootnote{N H; \skt{शोणितपित्तिभ्यः}\ A1938}}\skt{ }\edtext{\edtext{\edtext{\skt{क्षीरमिक्षुरसन्तथा}}{\lemma{\skt{क्षीरम्}}\Afootnote{N; \skt{क्षीर}\ A1938 H}}}{\lemma{\skt{इक्षुरसन्}}\Afootnote{N; \skt{मिक्षुरसस्}\ A1938; \skt{मिक्षुरसन्}\ H}}\skt{ ॥ {\symfont ❈}}}{\lemma{\skt{तथा ॥ {\symfont ❈}}}\Afootnote{N; \skt{तथा । अर्कशेलुशिरीषाणामासवास्तु विषार्तिषु ॥४३१॥}\ A1938; \skt{तथा ॥ ॥}\ H}}\skt{ ॥}\pend

\pstart\versepar \ledsidenote{\passref{1.46.432}}\skt{अतः परं तु }\edtext{\edtext{\skt{वर्गाणामनुपानम्पृथक्पृथक्}}{\lemma{\skt{वर्गाणाम्}}\Afootnote{H; \skt{वर्गाणा}\ A1938; \skt{वर्गानाम्}\ N}}}{\lemma{\skt{अनुपानम्}}\Afootnote{N H; \skt{मनुपानं}\ A1938}}\skt{ । }\pend
\pstart\versepar \skt{}\edtext{\skt{प्रवक्ष्याम्यनुपूर्व्वेण}}{\lemma{\skt{अनुपूर्व्वेण}}\Afootnote{N H; \skt{म्यनुपूर्वेण}\ A1938}}\skt{ सर्व्वेषामेव }\edtext{\skt{मे}}{\Afootnote{A1938 H; \skt{तत}\ N}}\skt{ शृणु ॥}\edtext{}{\lemma{\skt{पृथक्पृथक् । प्रवक्ष्याम्य्}}\Afootnote{H; \skt{पृथक्पृथक् । प्रवक्ष्या}\ A1938; \skt{पृठक्पृथक् । प्रवक्ष्याम्य्}\ N}}\pend

\pstart \ledsidenote{\passref{1.46.433}}\skt{तत्र }\edtext{\skt{पूर्व्वशस्य जा}}{\Afootnote{A1938 H; \skt{पूर्व्वसस्यजा}\ N}}\skt{ }\edtext{\skt{तानाम्वेदराम्लं}}{\lemma{\skt{वेदराम्लं}}\Afootnote{H; \skt{बदराम्लं,}\ A1938; \skt{बदराम्ब्लं ।}\ N}}\skt{ । }\edtext{\edtext{\skt{वैदलानान्धाम्लं}}{\lemma{\skt{वैदलानान्}}\Afootnote{A1938 H; \skt{वैदलाना}\ N}}}{\lemma{\skt{धाम्लं}}\Afootnote{H; \skt{धान्याम्लं,}\ A1938; \skt{धान्याम्ब्लं ।}\ N}}\skt{ । }\edtext{\edtext{\skt{जाङ्गलानान्मृगानान्धन्वजानाम्पक्षिणाञ्च}}{\lemma{\skt{जाङ्गलानान्मृगानान्}}\Afootnote{H; \skt{जङ्घालानां}\ A1938; \skt{जंगालानां मृगानां}\ N}}}{\lemma{\skt{पक्षिणाञ्च}}\Afootnote{N H; \skt{च}\ A1938}}\skt{ पिप्पल्यासवः । }\edtext{\skt{विष्किराणां}}{\Afootnote{A1938 H; \skt{विष्ठिराणां}\ N}}\skt{ }\edtext{\skt{कोलवदरासवः}}{\Afootnote{H; \skt{कोलबदरासवः}\ A1938 N}}\skt{ । प्रतुदानां }\edtext{\skt{क्षीरवृक्सासवः}}{\Afootnote{H; \skt{क्षीरवृक्षासवः}\ A1938 N}}\skt{ । }\edtext{\skt{गुहे}}{\Afootnote{N H; \skt{गुहा}\ A1938}}\skt{ शयानां }\edtext{\skt{खर्ज्जूरनालिकेरासवः}}{\Afootnote{H; \skt{खर्जूरनालिकेरासवः, प्रसहानामश्वगन्धासवः, पर्णमृगाणां कृष्णगन्धासवः, बिलेशयानां फलासवः , एकशफानां त्रिफलासवः, अनेकशफानां खदिरासवः, कूलचराणां शृङ्गाटककशेरुकासवः, कोशवासिनां पादिनां च स एव, प्लवानामिक्षुरसासवः, नादेयानां मत्स्यानां मृणालासवः,}\ A1938; \skt{खर्ज्जूरनाडिकेसरासवः ।}\ N}}\skt{ । }\edtext{\skt{सामुद्राणाम्मातुलुङ्गासवः}}{\lemma{\skt{मातुलुङ्गासवः}}\Afootnote{N H; \skt{मातुलुङ्गा}\ A1938}}\skt{ । }\edtext{\edtext{\edtext{\skt{कुष्माण्डानामम्लानाम्सार्द्वीकासवः}}{\lemma{\skt{अम्लानाम्}}\Afootnote{H; \skt{फलानां}\ N}}}{\lemma{\skt{सार्द्वीकासवः}}\Afootnote{H; \skt{गार्द्वीकासवः ।}\ N}}\skt{ । }\edtext{\skt{अम्ला अम्लानां}}{\Afootnote{\skt{अम्लानां}\ A1938 H; \skt{अम्ब्लानां}\ N}}\skt{ फलानां पद्मोत्पलकन्दासवः}}{\lemma{\skt{सवः… सार्द्वीकास}}\Afootnote{\emph{om.}\ A1938}}\skt{ । कषायाणां दाडिमवेत्रासवः । मधुराणां त्रिकटुकयुक्तः }\edtext{\skt{खदिरासवः}}{\Afootnote{N H; \skt{खण्डासवः,}\ A1938}}\skt{ । तालफलादीनां }\edtext{\skt{धान्याम्लं}}{\Afootnote{A1938 H; \skt{धान्याम्ब्लं ।}\ N}}\skt{ । }\edtext{\skt{कन्दानां}}{\Afootnote{N H; \skt{कटुकानां}\ A1938}}\skt{ दूर्व्वानलवेत्रासवः । पिप्पल्यादीनां }\edtext{\skt{स्वदंष्ट्रावसुकासवः}}{\Afootnote{H; \skt{श्वदंष्ट्रावसुकासवः, कूष्माण्डादीनां दार्वीकरीरासवः,}\ A1938; \skt{स्वद्रष्ट्रावसुकासवः ।}\ N}}\skt{ । }\edtext{\skt{चुच्चूप्रभृतीनां}}{\Afootnote{N H; \skt{चुच्चुप्रभृतीनां}\ A1938}}\skt{ }\edtext{\skt{लोध्रासवः}}{\Afootnote{N H; \skt{लोध्रासवः, जीवन्त्यादीनां त्रिफलासवः,}\ A1938}}\skt{ । कुसुम्भशाकस्य }\edtext{\skt{तदेव}}{\Afootnote{N H; \skt{स एव,}\ A1938}}\skt{ । }\edtext{\skt{मण्डूकपर्ण्यादीनाम्महापञ्चमूल्यासवः}}{\lemma{\skt{महापञ्चमूल्यासवः}}\Afootnote{H; \skt{महापञ्चमूलासवः,}\ A1938; \skt{महापञ्चशूल्यासवः ।}\ N}}\skt{ । }\edtext{\edtext{\skt{तालमस्तकादीनामम्लफलासवः}}{\lemma{\skt{तालमस्तकादीनाम्}}\Afootnote{N H; \skt{तालमस्तकादीना}\ A1938}}}{\lemma{\skt{अम्लफलासवः}}\Afootnote{H; \skt{मम्लफलासवः,}\ A1938; \skt{अम्ब्लफलासवः ।}\ N}}\skt{ । सैन्धवादीनां }\edtext{\edtext{\edtext{\skt{सुरासवमारनालन्तोयम्वा}}{\lemma{\skt{सुरासवम्}}\Afootnote{N H; \skt{सुरासव}\ A1938}}}{\lemma{\skt{आरनालन्}}\Afootnote{N H; \skt{आरनालं च,}\ A1938}}}{\lemma{\skt{तोयम्वा}}\Afootnote{N H; \skt{तोयं वा सर्वत्रेति ॥४३३॥}\ A1938}}\skt{ ॥}\pend

\pstart\versepar \ledsidenote{\passref{1.46.434}}\edtext{\skt{भवन्ति}}{\Afootnote{A1938 H; \skt{भवति}\ N}}\skt{ चात्र }\edtext{\skt{श्लोकाः}}{\Afootnote{H; \emph{om.}\ A1938; \skt{श्लोकः ॥}\ N}}\skt{ ॥ }\pend
\pstart\versepar \skt{सर्व्वेषामनुपानानां }\edtext{\skt{माहेन्द्रन्तोयमुत्तमं}}{\lemma{\skt{माहेन्द्रन्}}\Afootnote{A1938 H; \skt{माहेन्द्र}\ N}}\skt{ । }\pend
\pstart\versepar \skt{सात्म्यम्वा यस्य यत्तोयं }\edtext{\skt{तत्स्मै}}{\Afootnote{H; \skt{तत्तस्मै}\ A1938 N}}\skt{ }\edtext{\edtext{\skt{हितमुच्यते}}{\lemma{\skt{हितम्}}\Afootnote{N H; \skt{हित}\ A1938}}}{\lemma{\skt{उच्यते}}\Afootnote{N H; \skt{मुच्यते ॥४३४॥}\ A1938}}\skt{ ॥}\pend

\pstart \ledsidenote{\passref{1.46.435}}\edtext{\skt{दोषवद्गुरु वा}}{\Afootnote{N H; \skt{उष्णं वाते कफे तोयं पित्ते रक्ते च शीतलम्दोषवद्गुरु वा}\ A1938}}\skt{ }\edtext{\edtext{\edtext{\skt{भुक्तमतिमात्रमथापि}}{\lemma{\skt{भुक्तम्}}\Afootnote{N H; \skt{भुक्त}\ A1938}}}{\lemma{\skt{अतिमात्रम्}}\Afootnote{N H; \skt{मतिमात्र}\ A1938}}}{\lemma{\skt{अथापि}}\Afootnote{N H; \skt{मथापि}\ A1938}}\skt{ वा ॥}\pend

\pstart\versepar \ledsidenote{\passref{1.46.436}}\skt{यथोक्तेनानुपानेन }\edtext{\edtext{\skt{सुखमन्नं}}{\lemma{\skt{सुखम्}}\Afootnote{N H; \skt{सुख}\ A1938}}}{\lemma{\skt{अन्नं}}\Afootnote{N H; \skt{मन्नं}\ A1938}}\skt{ }\edtext{\skt{प्रजीर्यते}}{\Afootnote{N H; \skt{प्रजीर्यति ।}\ A1938}}\skt{ ॥ }\pend
\pstart\versepar \skt{रोचनं }\edtext{\edtext{\skt{वृंहणम्वृष्यन्दोषसंघातभेदनं}}{\lemma{\skt{वृंहणम्}}\Afootnote{H; \skt{बृंहणं}\ A1938; \skt{बृंहनं}\ N}}}{\lemma{\skt{वृष्यन्}}\Afootnote{A1938 H; \skt{वृष्य}\ N}}\skt{ ।}\pend

\pstart\versepar \ledsidenote{\passref{1.46.437}}\skt{तर्प्पणम्मार्द्दवकरं श्रमक्लमहरं सुखं ॥ }\pend
\pstart\versepar \skt{दीपनन्दोषशमनं पिपासाच्छेदनम्परं ।}\pend

\pstart\versepar \ledsidenote{\passref{1.46.438}}\edtext{\skt{वल्यं}}{\Afootnote{H; \skt{बल्यं}\ A1938 N}}\skt{ }\edtext{\skt{वर्ण्णकरञ्चापि}}{\lemma{\skt{चापि}}\Afootnote{N H; \skt{सम्य}\ A1938}}\skt{ }\edtext{\skt{अनुपानं}}{\Afootnote{N H; \skt{गनुपानं}\ A1938}}\skt{ सदोच्यते ॥ }\pend
\pstart\versepar \skt{}\edtext{\skt{तदादौ}}{\Afootnote{A1938 H; \skt{तदादौ}\ N}}\skt{ }\edtext{\skt{कर्षयेत्पीतं}}{\Afootnote{N H; \skt{कर्शयेत्पीतं}\ A1938}}\skt{ }\edtext{\skt{स्थापयेत्मध्यसेवितं}}{\Afootnote{H; \skt{स्थापयेन्मध्यसेवितम् ॥४३८॥}\ A1938; \skt{स्थापये मध्यसेवितं ।}\ N}}\skt{ ।}\pend

\pstart\versepar \ledsidenote{\passref{1.46.439}}\skt{पश्चात्पीतं }\edtext{\skt{वृंहयति}}{\Afootnote{H; \skt{बृंहयति}\ A1938 N}}\skt{ }\edtext{\skt{तत्समीक्ष्य}}{\Afootnote{N H; \skt{तस्माद्वीक्ष्य}\ A1938}}\skt{ प्रयोजयेत् ॥ }\pend
\pstart\versepar \skt{}\edtext{\edtext{\edtext{\edtext{\skt{स्थिरताङ्गतमक्लिन्नमन्नमद्रवपायिनः}}{\lemma{\skt{गतम्}}\Afootnote{N H; \skt{गत}\ A1938}}}{\lemma{\skt{अक्लिन्नम्}}\Afootnote{N H; \skt{मल्किन्न}\ A1938}}}{\lemma{\skt{अन्नम्}}\Afootnote{N H; \skt{मन्न}\ A1938}}}{\lemma{\skt{अद्रवपायिनः}}\Afootnote{N H; \skt{मद्रवपायिनाम् ॥४३९॥}\ A1938}}\skt{ ।}\pend

\pstart\versepar \ledsidenote{\passref{1.46.440}}\edtext{\edtext{\edtext{\edtext{\skt{भवत्यावाधजनमनुपानमतः}}{\lemma{\skt{भवत्य्}}\Afootnote{N H; \skt{भव}\ A1938}}}{\lemma{\skt{आवाधजनम्}}\Afootnote{H; \skt{त्याबाधजनन}\ A1938; \skt{आवाधजननम्}\ N}}}{\lemma{\skt{अनुपानम्}}\Afootnote{N H; \skt{मनुपान}\ A1938}}}{\lemma{\skt{अतः}}\Afootnote{N H; \skt{मतः}\ A1938}}\skt{ पिवेत् ॥ }\pend
\pstart\versepar \skt{न }\edtext{\skt{पिवेच्छ्वासकासार्तरोगे}}{\Afootnote{H; \skt{पिबेच्छ्वासकासार्तो रोगे}\ A1938 N}}\skt{ }\edtext{\edtext{\skt{वाप्युर्ध्वजत्रुगे}}{\lemma{\skt{वाप्य्}}\Afootnote{N H; \skt{चाप्य्}\ A1938}}}{\lemma{\skt{उर्ध्वजत्रुगे}}\Afootnote{H; \skt{ऊर्ध्वजत्रुगे ॥४४०॥}\ A1938; \skt{ऊर्द्वजन्त्रुकेजे ।}\ N\textsuperscript{ac}; \skt{ऊर्द्वजन्त्रुजे ।}\ N\textsuperscript{pc}}}\skt{ ।}\edtext{}{\lemma{\skt{पिवेत् ॥ न}}\Afootnote{H; \skt{पिबेत्न}\ A1938 N}}\pend

\pstart\versepar \ledsidenote{\passref{1.46.441}}\edtext{\skt{क्षतोरस्कप्रसेकी च}}{\Afootnote{N H; \skt{क्षतोरस्कः प्रसेकी च}\ A1938}}\skt{ यस्य }\edtext{\skt{चोपहतस्वरः}}{\Afootnote{N H; \skt{चोपहतः स्वरः ।}\ A1938}}\skt{ ॥ }\pend
\pstart\versepar \skt{}\edtext{\skt{पीत्वा च}}{\Afootnote{H; \skt{पीत्वाऽध्व}\ A1938; \skt{पीत्वाव}\ N}}\skt{ }\edtext{\skt{भाष्याध्ययन}}{\Afootnote{A1938 H; \skt{भाष्यध्ययन}\ N}}\skt{ }\edtext{\skt{स्वप्रगेयान्न}}{\Afootnote{H; \skt{गेयस्वप्नान्न}\ A1938; \skt{स्वप्नगेयान्न}\ N}}\skt{ शीलयेत् ।}\pend

\pstart\versepar \ledsidenote{\passref{1.46.442}}\edtext{\edtext{\skt{प्रदूष्यामाशयन्तद्धि}}{\lemma{\skt{प्रदूष्यामाशयन्}}\Afootnote{A1938 H; \skt{प्रदूष्यामाषयं}\ N}}}{\lemma{\skt{तद्धि}}\Afootnote{A1938 H; \skt{तद्धि}\ N}}\skt{ तस्य कण्ठोरसि स्थितं ॥ }\pend
\pstart\versepar \skt{}\edtext{\skt{स्यन्दाग्निसादच्छर्द्यादीनजन}}{\Afootnote{H; \skt{स्यन्दाग्निसादच्छर्द्यादीनामयाञ्जन}\ A1938; \skt{स्यन्दाग्निस्यादच्छर्द्यादीञ्जान}\ N}}\skt{ }\edtext{\skt{यदा}}{\Afootnote{H; \skt{ये}\ A1938; \skt{येदा}\ N}}\skt{ }\edtext{\skt{मयान्वहून्}}{\Afootnote{H; \skt{द्बहून् ॥४४२॥}\ A1938; \skt{मयास्बहून् ।}\ N}}\skt{ ।}\pend

\pstart\versepar \ledsidenote{\passref{1.46.443}}\edtext{\skt{गुरूलाघवचिन्तेयं}}{\Afootnote{H; \skt{गुरुलाघवचिन्तेयं}\ A1938 N}}\skt{ स्वभावन्नातिवर्त्तते ॥ }\pend
\pstart\versepar \skt{तथा }\edtext{\skt{संस्करमात्रान्न}}{\Afootnote{H; \skt{संस्कारमात्रान्न}\ A1938 N}}\skt{ }\edtext{\edtext{\skt{कालाश्चाप्युत्तरोत्तरां}}{\lemma{\skt{कालाश्}}\Afootnote{H; \skt{कालांश्}\ A1938 N}}}{\lemma{\skt{उत्तरोत्तरां}}\Afootnote{N H; \skt{उत्तरोत्तरम् ॥४४३॥}\ A1938}}\skt{ ।}\pend

\pstart\versepar \ledsidenote{\passref{1.46.444}}\skt{मन्दकर्म्मानलारोग्याः सुकुमाराः }\edtext{\skt{सुखोचिताः}}{\Afootnote{A1938 H; \skt{सुखे चिताः ॥}\ N}}\skt{ ॥ }\pend
\pstart\versepar \skt{जन्तवो ये तु तेषां हि }\edtext{\skt{चिन्तेयम्परिकीर्त्तिते}}{\lemma{\skt{परिकीर्त्तिते}}\Afootnote{H; \skt{परिकीर्त्यते}\ A1938 N}}\skt{ ।}\pend

\pstart\versepar \ledsidenote{\passref{1.46.445}}\edtext{\skt{वलिनः}}{\Afootnote{H; \skt{बलिनः}\ A1938 N}}\skt{ }\edtext{\skt{खरभक्ष्याश्च ये च}}{\Afootnote{N H; \skt{खरभक्ष्या ये ये च}\ A1938}}\skt{ दीप्ताग्नयो नराः ॥ }\pend
\pstart\versepar \skt{कर्म्मनित्याश्च ये }\edtext{\skt{तेषान्नावश्यम्परिचिन्त्यते}}{\lemma{\skt{परिचिन्त्यते}}\Afootnote{N H; \skt{परिकीर्त्यते ॥४४५॥ इत्यनुपानवर्गः ।}\ A1938}}\skt{ ॥ ॥}\pend

\pstart\versepar \ledsidenote{\passref{1.46.446}}\edtext{\skt{अथाहारविधिविस्तरेणाखिलं}}{\Afootnote{N\textsuperscript{pc} H; \skt{अथाहारविधिं वत्स! विस्तरेणाखिलं}\ A1938; \skt{अथाहारिविधिविस्तरेणाखिलं}\ N\textsuperscript{ac}}}\skt{ शृणु ॥ }\pend
\pstart\versepar \skt{}\edtext{\edtext{\skt{आप्तास्थितमसंकीर्ण्णे}}{\lemma{\skt{आप्तास्थितम्}}\Afootnote{N; \skt{आप्तास्थित}\ A1938 H}}}{\lemma{\skt{असंकीर्ण्णे}}\Afootnote{\skt{मसङ्कीर्णं}\ A1938; \skt{असंकीर्ण्णं}\ N; \skt{मसंकीर्ण्णे}\ H}}\skt{ }\edtext{\skt{शुचि}}{\Afootnote{A1938 H; \skt{सुचि}\ N}}\skt{ }\edtext{\skt{कार्यं}}{\Afootnote{A1938 H; \skt{द्कायाअं}\ N\textsuperscript{ac}; \skt{द्कायअं}\ N\textsuperscript{pc}}}\skt{ महानसं ।}\pend

\pstart\versepar \ledsidenote{\passref{1.46.447}}\edtext{\edtext{\skt{तत्राप्तैर्गुणसम्पन्नम्भक्ष्यादिषु}}{\lemma{\skt{तत्राप्तैर्}}\Afootnote{A1938 H; \skt{यत्राप्तै}\ N}}}{\lemma{\skt{भक्ष्यादिषु}}\Afootnote{N H; \skt{भक्ष्यमन्नं}\ A1938}}\skt{ सुसंस्कृतं ॥ }\pend
\pstart\versepar \skt{}\edtext{\skt{शुचौ}}{\Afootnote{A1938 H; \skt{सुचौ}\ N}}\skt{ }\edtext{\skt{देशेषु}}{\Afootnote{H; \skt{देशे सु}\ A1938 N}}\skt{ संगुप्तं }\edtext{\skt{समुपस्थापयेद्भिषक्}}{\Afootnote{A1938 H; \skt{समुपास्थापयेद्भिषक् ।}\ N}}\skt{ ।}\pend

\pstart\versepar \ledsidenote{\passref{1.46.448}}\skt{विषघ्नैरगदैः स्पृष्टम्प्रोक्षितं व्यजनोदकैः ॥ }\pend
\pstart\versepar \skt{सिद्धैर्मन्त्रैर्हतविषं }\edtext{\edtext{\edtext{\skt{सिद्धमन्नन्निदेदयेत्}}{\lemma{\skt{सिद्धम्}}\Afootnote{N H; \skt{सिद्ध}\ A1938}}}{\lemma{\skt{अन्नन्}}\Afootnote{N H; \skt{मन्नं}\ A1938}}}{\lemma{\skt{निदेदयेत्}}\Afootnote{H; \skt{निवेदयेत्}\ A1938 N}}\skt{ ॥ ० ॥}\pend

\pstart\versepar \ledsidenote{\passref{1.46.449}}\edtext{\skt{वक्ष्याम्यतः}}{\Afootnote{A1938; \skt{वक्ष्याम्यतः}\ N H}}\skt{ }\edtext{\skt{परं}}{\Afootnote{A1938 H; \skt{पर}\ N}}\skt{ }\edtext{\edtext{\skt{कृत्स्नामाहारस्योपकल्पनां}}{\lemma{\skt{कृत्स्नाम्}}\Afootnote{H; \skt{कृत्स्ना}\ A1938; \skt{कृष्णाम्}\ N}}}{\lemma{\skt{आहारस्योपकल्पनां}}\Afootnote{N H; \skt{माहारस्योपकल्पनाम् ।}\ A1938}}\skt{ । }\pend
\pstart\versepar \skt{घृतं कार्ष्णायसे देयं पेया }\edtext{\skt{देया तु}}{\Afootnote{A1938 H; \skt{देये तु}\ N}}\skt{ राजते ॥}\pend

\pstart\versepar \ledsidenote{\passref{1.46.450}}\edtext{\skt{फलानि}}{\Afootnote{A1938 H; \skt{फलास्च}\ N}}\skt{ }\edtext{\skt{सर्व्वभक्ष्यांश्च}}{\Afootnote{A1938 H; \skt{सर्व्वभक्ष्यांस्च}\ N}}\skt{ }\edtext{\skt{प्रदया}}{\Afootnote{H; \skt{प्रदद्या}\ A1938; \skt{प्रदेया}\ N}}\skt{ }\edtext{\skt{वैदलेषु तु}}{\Afootnote{H; \skt{द्वै दलेषु तु ।}\ A1938; \skt{वैदलीषु तु ।}\ N}}\skt{ । }\pend
\pstart\versepar \skt{}\edtext{\skt{परिशुष्कं}}{\Afootnote{N H; \skt{परिशुष्क}\ A1938}}\skt{ प्रदिग्धानि }\edtext{\skt{सौवर्ण्णेषूपहारयेत्}}{\Afootnote{H; \skt{सौवर्णेषु प्रकल्पयेत् ॥४५०॥}\ A1938; \skt{सौवर्ण्णेषूपकल्पयेत् ॥}\ N}}\skt{ ॥}\pend

\pstart\versepar \ledsidenote{\passref{1.46.451}}\edtext{\skt{द्रवाणि}}{\Afootnote{H; \skt{प्रद्रवाणि}\ A1938; \skt{द्रवानि तु}\ N}}\skt{ }\edtext{\edtext{\skt{रसाश्चैव}}{\lemma{\skt{रसाश्}}\Afootnote{N H; \skt{रसांश्}\ A1938}}}{\lemma{\skt{चैव}}\Afootnote{N H; \skt{चैवं}\ A1938}}\skt{ राजतेषूपहारयेत् । }\pend
\pstart\versepar \skt{कट्वराणि }\edtext{\skt{खलाश्चैव}}{\lemma{\skt{खलाश्}}\Afootnote{N H; \skt{खडांश्}\ A1938}}\skt{ }\edtext{\skt{सर्व्वाञ्च्छैलेषु}}{\lemma{\skt{सर्व्वाञ्}}\Afootnote{A1938 H; \skt{सर्वा}\ N}}\skt{ दापयेत् ॥}\edtext{}{\lemma{\skt{राजतेषूपहारयेत् । कट्वराणि}}\Afootnote{A1938 H; \skt{राजतेषूपकल्पयेत् । कट्वराणि}\ N}}\pend

\pstart\versepar \ledsidenote{\passref{1.46.452}}\skt{दद्यात्ताम्रमये पात्रेसु }\edtext{\skt{शीतं}}{\Afootnote{A1938 H; \skt{सीतं}\ N}}\skt{ }\edtext{\skt{सुश्रितं}}{\Afootnote{H; \skt{सुशृतं}\ A1938; \skt{शुश्रितम्}\ N}}\skt{ पयः । }\pend
\pstart\versepar \skt{}\edtext{\skt{पानीयम्पानकम्मद्यम्मृन्मयेषु}}{\lemma{\skt{मद्यम्}}\Afootnote{A1938 H; \skt{मद्य}\ N}}\skt{ प्रदापयेत् ॥}\pend

\pstart\versepar \ledsidenote{\passref{1.46.453}}\skt{काचस्फटिकपात्रेषु }\edtext{\skt{शीतलेषु}}{\Afootnote{A1938 H; \skt{सीतलेषु}\ N}}\skt{ }\edtext{\skt{शुभेषु}}{\Afootnote{H; \skt{शुभेषु च ।}\ A1938 N}}\skt{ । }\pend
\pstart\versepar \skt{}\edtext{\skt{वज्रवैदूर्यचित्रेषु}}{\Afootnote{H; \skt{दद्याद्वैदूर्यचित्रेषु}\ A1938; \skt{वज्रवैडूर्यचित्रेषु}\ N}}\skt{ }\edtext{\skt{रागषाडवषदूकान्}}{\Afootnote{H; \skt{रागषाडवसट्टकान् ॥४५३॥}\ A1938; \skt{रागसाडवषट्टकान् ॥}\ N}}\skt{ ॥}\pend

\pstart\versepar \ledsidenote{\passref{1.46.454}}\edtext{\skt{पुरस्ताद्विमले}}{\Afootnote{N H; \skt{पुरस्ताद्विमले}\ A1938}}\skt{ }\edtext{\skt{पारे}}{\Afootnote{H; \skt{पात्रे}\ A1938 N}}\skt{ सुविस्तीर्ण्णे मनोरमे । }\pend
\pstart\versepar \skt{सूदः }\edtext{\edtext{\skt{सूपोदनन्दद्यात्प्रदेहाश्च}}{\lemma{\skt{सूपोदनन्}}\Afootnote{N H; \skt{सूपौदनं}\ A1938}}\skt{ सु}}{\lemma{\skt{दद्यात्प्रदेहाश्च सु}}\Afootnote{H; \skt{दद्यात्प्रदेहांश्च सु}\ A1938 N}}\skt{ }\edtext{\skt{संस्कृतं}}{\Afootnote{H; \skt{संस्कृतान् ॥४५४॥}\ A1938; \skt{संस्कृतां ॥}\ N}}\skt{ ॥}\pend

\pstart\versepar \ledsidenote{\passref{1.46.455}}\skt{फलानि सर्व्व }\edtext{\skt{भक्साश्च}}{\Afootnote{H; \skt{भक्ष्यांश्च}\ A1938; \skt{भक्स्याश्च}\ N}}\skt{ }\edtext{\skt{परिशुष्काणि}}{\Afootnote{A1938 H; \skt{परिसुष्कानि}\ N}}\skt{ यानि च । }\pend
\pstart\versepar \skt{तानि दक्षिणपार्श्वे तु भुञ्जानस्योपकल्पयेत् ॥}\pend

\pstart\versepar \ledsidenote{\passref{1.46.456}}\edtext{\skt{प्रद्रवाणि}}{\Afootnote{A1938 H; \skt{प्रदवानि}\ N}}\skt{ }\edtext{\skt{रसाश्चैव}}{\lemma{\skt{रसाश्}}\Afootnote{H; \skt{रसांश्}\ A1938 N}}\skt{ }\edtext{\skt{पानीयम्पानकम्पयः}}{\lemma{\skt{पानकम्}}\Afootnote{A1938 H; \skt{पानम्}\ N}}\skt{ । }\pend
\pstart\versepar \skt{}\edtext{\skt{खलान्यूषाश्च}}{\Afootnote{H; \skt{खडान्यूषांश्च}\ A1938; \skt{खलान्यूषांश्च}\ N}}\skt{ }\edtext{\skt{पेयाश्च}}{\Afootnote{H; \skt{पेयांश्च}\ A1938 N}}\skt{ }\edtext{\skt{सर्व्वे}}{\Afootnote{H; \skt{सव्ये}\ A1938; \skt{सर्वो}\ N}}\skt{ पार्श्वे प्रदापयेत् ॥}\pend

\pstart\versepar \ledsidenote{\passref{1.46.457}}\edtext{\skt{सर्व्वान्गुडविकाराश्च}}{\Afootnote{H; \skt{सर्वान्गुडविकारांश्च}\ A1938; \skt{सर्वां गुडविकारांश्च}\ N}}\skt{ रागषाडवषटूकान् । }\pend
\pstart\versepar \skt{पुरस्थात्स्थापयेत्प्राज्ञो द्वयोरपि च }\edtext{\skt{मधयतः}}{\Afootnote{N H; \skt{मध्यतः ॥४५७॥}\ A1938}}\skt{ ॥}\edtext{}{\lemma{\skt{रागषाडवषटूकान् । पुरस्थात्स्थापयेत्प्राज्ञो}}\Afootnote{H; \skt{रागषाडवसट्टकान् । पुरस्तात्स्थापयेत्प्राज्ञो}\ A1938; \skt{रागषाडवषट्टकं ॥ पुरस्थात्स्थापयेत्प्राज्ञो}\ N}}\pend

\pstart\versepar \ledsidenote{\passref{1.46.458}}\skt{एवं विज्ञाय }\edtext{\skt{मतिमान्भोजनस्योपकल्पनां}}{\Afootnote{A1938 H; \skt{मतिमां भोजनस्योपकल्पनां ॥}\ N}}\skt{ । }\pend
\pstart\versepar \skt{}\edtext{\edtext{\skt{भोक्तारम्विजने}}{\lemma{\skt{भोक्तारम्}}\Afootnote{A1938 H; \skt{भोजये}\ N}}}{\lemma{\skt{विजने}}\Afootnote{A1938 H; \skt{द्विजने}\ N}}\skt{ रम्ये निःसम्पाते शुभे शुचौ ॥}\pend

\pstart\versepar \ledsidenote{\passref{1.46.459}}\skt{सुगन्धपुष्परचिते समे }\edtext{\skt{देशेऽथ}}{\lemma{\skt{ऽथ}}\Afootnote{\skt{च}\ A1938; \skt{थ}\ N H}}\skt{ भोजयेत् । }\pend
\pstart\versepar \skt{}\edtext{\skt{विशिष्टमिष्टसंस्कारैः}}{\lemma{\skt{इष्टसंस्कारैः}}\Afootnote{H; \skt{मिष्टसंस्कारैः}\ A1938; \skt{इष्टसंस्कारै}\ N}}\skt{ }\edtext{\edtext{\edtext{\skt{पथैर्हृद्यैरसादिभिः}}{\lemma{\skt{पथैर्}}\Afootnote{H; \skt{पथ्यैर्}\ A1938; \skt{पेथ्यैर्}\ N}}}{\lemma{\skt{हृद्यैर्}}\Afootnote{H; \skt{इष्टै}\ A1938; \skt{हृद्यै}\ N}}}{\lemma{\skt{असादिभिः}}\Afootnote{H; \skt{रसादिभिः}\ A1938 N}}\skt{ ॥}\edtext{}{\lemma{\skt{भोजयेत् । विशिष्टम्}}\Afootnote{H; \skt{भोजयेत्विशिष्ट}\ A1938; \skt{भोजयेत् ॥ विसिष्टम्}\ N}}\pend

\pstart\versepar \ledsidenote{\passref{1.46.460}}\skt{मनोज्ञं शुचि }\edtext{\skt{नात्युष्णं}}{\Afootnote{A1938; \skt{नान्युष्णम्}\ N; \skt{नात्युष्णं}\ H}}\skt{ }\edtext{\edtext{\skt{प्रत्यग्रमशनं}}{\lemma{\skt{प्रत्यग्रम्}}\Afootnote{N H; \skt{प्रत्यग्र}\ A1938}}}{\lemma{\skt{अशनं}}\Afootnote{N H; \skt{मशनं}\ A1938}}\skt{ हितं । }\pend
\pstart\versepar \skt{}\edtext{\edtext{\skt{पूर्व्वम्मधुरमश्नीयाल्लवणोम्लौतः}}{\lemma{\skt{मधुरम्}}\Afootnote{N H; \skt{मधुर}\ A1938}}}{\lemma{\skt{अश्नीयाल्लवणोम्लौतः}}\Afootnote{H; \skt{मश्नीयान्मध्येऽम्ललवणौ रसौ ॥४६०॥}\ A1938; \skt{अश्नीयाल्लवण्यम्ब्लौ ततः}\ N}}\skt{ }\edtext{\skt{परं}}{\Afootnote{N H; \emph{om.}\ A1938}}\skt{ ॥}\pend

\pstart\versepar \ledsidenote{\passref{1.46.461}}\edtext{\skt{पश्चादेषान्रसान्वै द्यो}}{\Afootnote{H; \skt{पश्चाच्छेषान्रसान्वैद्यो}\ A1938; \skt{पश्चाच्छेषां रसा वैद्यो}\ N}}\skt{ }\edtext{\edtext{\skt{भोजनेषेव}}{\lemma{\skt{भोजनेष्}}\Afootnote{H; \skt{भोजनेष्व्}\ A1938 N}}}{\lemma{\skt{एव}}\Afootnote{H; \skt{अव}\ A1938 N}}\skt{ चारयेत् । }\pend
\pstart\versepar \skt{पादौ फलानि }\edtext{\skt{युञ्जीत}}{\Afootnote{N H; \skt{भुञ्जीत}\ A1938}}\skt{ दाडिमादीनि }\edtext{\skt{वुद्धिमान्}}{\Afootnote{H; \skt{बुद्धिमान्}\ A1938 N}}\skt{ ॥}\edtext{}{\lemma{\skt{चारयेत् । पादौ}}\Afootnote{H; \skt{चारयेतादौ}\ A1938 N}}\pend

\pstart\versepar \ledsidenote{\passref{1.46.462}}\skt{ततः }\edtext{\edtext{\skt{पेयान्ततो}}{\lemma{\skt{पेयान्}}\Afootnote{A1938 H; \skt{पेया}\ N}}}{\lemma{\skt{ततो}}\Afootnote{H; \skt{स्ततो}\ A1938; \skt{त्ततो}\ N}}\skt{ }\edtext{\edtext{\edtext{\skt{भोज्याञ्श्चित्रांस्ततः}}{\lemma{\skt{भोज्याञ्श्}}\Afootnote{H; \skt{भोज्यान्भक्ष्यां}\ A1938; \skt{भोज्यञ्}\ N}}}{\lemma{\skt{चित्रांस्}}\Afootnote{H; \skt{श्चित्रां}\ A1938; \skt{चित्राम्भक्ष्यास्}\ N}}}{\lemma{\skt{ततः}}\Afootnote{H; \skt{स्ततः}\ A1938; \skt{त्वतः}\ N}}\skt{ परं । }\pend
\pstart\versepar \skt{घनं पूर्व्वं }\edtext{\skt{समश्नीयादिति}}{\Afootnote{N H; \skt{समश्नीया}\ A1938}}\skt{ }\edtext{\skt{केचिदवस्थिताः}}{\Afootnote{H; \skt{त्, केचिदाहुर्विपर्ययम् ॥४६२॥}\ A1938; \skt{केचिद्व्यवस्थिताः ।}\ N}}\skt{ ॥}\pend

\pstart\versepar \ledsidenote{\passref{1.46.463}}\edtext{\skt{आदावन्ते च}}{\Afootnote{\skt{आदावन्ते च}\ A1938 N H}}\skt{ मध्ये च }\edtext{\skt{भोजने च}}{\Afootnote{H; \skt{भोजनस्य}\ A1938; \skt{भोजने व}\ N}}\skt{ }\edtext{\skt{प्रशस्यते}}{\Afootnote{H; \skt{तु शस्यते ।}\ A1938; \skt{प्रस्यते ॥}\ N}}\skt{ । }\pend
\pstart\versepar \skt{}\edtext{\edtext{\skt{निरत्ययन्दोषहरम्फलेष्वामलकं}}{\lemma{\skt{फलेष्व्}}\Afootnote{N\textsuperscript{pc} H; \skt{फलेष्वा}\ A1938; \skt{फलेओष्व्}\ N\textsuperscript{ac}}}}{\lemma{\skt{आमलकं}}\Afootnote{N H; \skt{मलकं}\ A1938}}\skt{ नृणां ॥}\pend

\pstart\versepar \ledsidenote{\passref{1.46.464}}\edtext{\skt{मृणालविसशालूककन्देक्षुप्रभृतीनि तु}}{\Afootnote{H; \skt{मृणालबिसशालूककन्देक्षुप्रभृतीनि च ।}\ A1938; \skt{मृणालविससालूककन्देक्षुप्रभृतीनि तु ॥}\ N}}\skt{ । }\pend
\pstart\versepar \skt{पूर्व्वं योज्यानि भिषजा न तु }\edtext{\skt{भुक्ते}}{\Afootnote{A1938 H; \skt{भक्ते}\ N}}\skt{ }\edtext{\skt{कथञ्चन}}{\Afootnote{H; \skt{कदाचन ॥४६४॥}\ A1938; \skt{कथञ्चनः ।}\ N}}\skt{ ॥}\pend

\pstart\versepar \ledsidenote{\passref{1.46.465}}\edtext{\edtext{\skt{सुखमुच्चैः}}{\lemma{\skt{सुखम्}}\Afootnote{N H; \skt{सुख}\ A1938}}}{\lemma{\skt{उच्चैः}}\Afootnote{H; \skt{मुच्चैः}\ A1938; \skt{उप्वैस्}\ N}}\skt{ }\edtext{\skt{समासीनं}}{\Afootnote{H; \skt{समासीनः}\ A1938; \skt{सामानीनं}\ N}}\skt{ }\edtext{\skt{समदेहोन्नतत्परम्}}{\Afootnote{H; \skt{समदेहोऽन्नतत्परः ।}\ A1938; \skt{समदेहोन्नतत्परः ॥}\ N}}\skt{ । }\pend
\pstart\versepar \skt{काले }\edtext{\skt{सात्म्यं}}{\Afootnote{A1938 H; \skt{सात्म्य}\ N}}\skt{ लघु }\edtext{\edtext{\skt{स्निग्धमुष्णं}}{\lemma{\skt{स्निग्धम्}}\Afootnote{N H; \skt{स्निग्ध}\ A1938}}}{\lemma{\skt{उष्णं}}\Afootnote{N H; \skt{मुष्णं}\ A1938}}\skt{ क्षिप्रं द्रवोत्तरं ॥}\pend

\pstart\versepar \ledsidenote{\passref{1.46.466}}\edtext{\edtext{\skt{वुभुक्षितोन्नमश्नीयात्मात्रावद्विदितासनः}}{\lemma{\skt{वुभुक्षितोन्नम्}}\Afootnote{H; \skt{बुभुक्षितोऽन्न}\ A1938; \skt{बुभुक्षितोन्नम्}\ N}}}{\lemma{\skt{अश्नीयात्मात्रावद्विदितासनः}}\Afootnote{N H; \skt{मश्नीयान्मात्रावद्विदितागमः}\ A1938}}\skt{ । }\pend
\pstart\versepar \skt{काले }\edtext{\edtext{\skt{भुक्तम्प्रीणयति}}{\lemma{\skt{भुक्तम्}}\Afootnote{N H; \emph{om.}\ A1938}}}{\lemma{\skt{प्रीणयति}}\Afootnote{N H; \skt{प्रीणयते}\ A1938}}\skt{ }\edtext{\edtext{\skt{सात्म्यमन्नन्न}}{\lemma{\skt{सात्म्यम्}}\Afootnote{N H; \skt{भुक्तं सात्म्य}\ A1938}}}{\lemma{\skt{अन्नन्न}}\Afootnote{N H; \skt{मन्नं न}\ A1938}}\skt{ }\edtext{\skt{वाधते}}{\Afootnote{H; \skt{बाधते}\ A1938 N}}\skt{ ॥}\pend

\pstart\versepar \ledsidenote{\passref{1.46.467}}\edtext{\skt{लघुशीघ्रम्व्रजेत्पाकं}}{\lemma{\skt{लघुशीघ्रम्}}\Afootnote{A1938 H; \skt{लघुसीघ्रं}\ N}}\skt{ स्निग्धोष्णं }\edtext{\skt{वलवधृदं}}{\Afootnote{H; \skt{बलवह्निदम् ।}\ A1938; \skt{बलवह्निदः ॥}\ N}}\skt{ । }\pend
\pstart\versepar \skt{}\edtext{\skt{क्षिप्रमुक्तं}}{\lemma{\skt{उक्तं}}\Afootnote{H; \skt{भुक्तं}\ A1938; \skt{भुंक्तं}\ N}}\skt{ समम्पाकं }\edtext{\edtext{\skt{यात्यदुष्टन्द्रवोत्तरं}}{\lemma{\skt{यात्य्}}\Afootnote{N H; \skt{या}\ A1938}}}{\lemma{\skt{अदुष्टन्}}\Afootnote{N H; \skt{त्यदोषं}\ A1938}}\skt{ ॥}\pend

\pstart\versepar \ledsidenote{\passref{1.46.468}}\skt{सुखञ्जीर्यति }\edtext{\skt{मात्रावद्धातुसात्म्यं}}{\Afootnote{N H; \skt{मात्रावद्धातुसाम्यं}\ A1938}}\skt{ करोति च । }\pend
\pstart\versepar \skt{}\edtext{\skt{अतीवायतमात्राम्यः}}{\Afootnote{H; \skt{अतीवायतयामास्तु}\ A1938; \skt{अतीवायतमात्रा स्युः}\ N}}\skt{ क्षपा येष्वृतुषु स्मृताः ॥}\pend

\pstart\versepar \ledsidenote{\passref{1.46.469}}\skt{तेषु }\edtext{\skt{तत्प्रत्यनीकार्थं}}{\Afootnote{N H; \skt{तत्प्रत्यनीकाढ्यं}\ A1938}}\skt{ भुञ्जीत }\edtext{\skt{प्रातरेव तु}}{\Afootnote{A1938 H; \skt{प्रातरेव तु ॥}\ N}}\skt{ । }\pend
\pstart\versepar \skt{येषु चापि }\edtext{\skt{भवेयुस्तदिव सा}}{\lemma{\skt{तदिव सा}}\Afootnote{H; \skt{च दिवसा}\ A1938; \skt{ते दिवसा}\ N}}\skt{ }\edtext{\skt{भृशमायताः}}{\Afootnote{A1938 H; \skt{भृषमायताः ।}\ N}}\skt{ ॥}\pend

\pstart\versepar \ledsidenote{\passref{1.46.470}}\skt{तेषु }\edtext{\skt{तत्कालवि}}{\Afootnote{A1938 H; \skt{तत्कालभि}\ N}}\skt{ }\edtext{\edtext{\skt{हितमपराह्णे}}{\lemma{\skt{हितम्}}\Afootnote{N H; \skt{हित}\ A1938}}}{\lemma{\skt{अपराह्णे}}\Afootnote{N H; \skt{मपराह्णे}\ A1938}}\skt{ प्रशस्यते । }\pend
\pstart\versepar \skt{रजन्यो दिवसाश्चैव येषु }\edtext{\skt{वापि}}{\Afootnote{H; \skt{चापि}\ A1938 N}}\skt{ समाः स्मृताः ॥}\pend

\pstart\versepar \ledsidenote{\passref{1.46.471}}\skt{कृत्वा }\edtext{\edtext{\skt{सममहोरात्रं}}{\lemma{\skt{समम्}}\Afootnote{N H; \skt{सम}\ A1938}}}{\lemma{\skt{अहोरात्रं}}\Afootnote{N H; \skt{महोरात्रं}\ A1938}}\skt{ तेषु भुञ्जीत भोजनं । }\pend
\pstart\versepar \skt{}\edtext{\skt{अप्राप्तातीतकालम्वा}}{\Afootnote{N H; \skt{नाप्राप्तातीतकालं वा}\ A1938}}\skt{ }\edtext{\skt{नाश्नीयात्तु}}{\Afootnote{N H; \skt{हीनाधिकमथापि वा ॥४७१॥}\ A1938}}\skt{ }\edtext{\skt{यथा}}{\Afootnote{N H; \emph{om.}\ A1938}}\skt{ }\edtext{\skt{तथं}}{\Afootnote{N H; \emph{om.}\ A1938}}\skt{ ॥}\pend

\pstart\versepar \ledsidenote{\passref{1.46.472}}\skt{अप्राप्तकालम्भुञ्जानः शरीरे }\edtext{\skt{ह्यघौ}}{\lemma{\skt{अघौ}}\Afootnote{H; \skt{अलघौ}\ A1938; \skt{अलघु}\ N}}\skt{ नरः । }\pend
\pstart\versepar \skt{}\edtext{\skt{तास्तान्व्याधीनवाप्नोति}}{\Afootnote{N H; \skt{तांस्तान्व्याधीनवाप्नोति}\ A1938}}\skt{ मरणम्वा नियच्छति ॥}\pend

\pstart\versepar \ledsidenote{\passref{1.46.473}}\edtext{\skt{अतीतकालम्भुक्तन्तु}}{\lemma{\skt{भुक्तन्तु}}\Afootnote{N H; \skt{भुञ्जानो}\ A1938}}\skt{ }\edtext{\skt{वायुनोपहतेऽनले}}{\lemma{\skt{ऽनले}}\Afootnote{A1938 H; \skt{नले ॥}\ N}}\skt{ । }\pend
\pstart\versepar \skt{}\edtext{\skt{कृच्छाद्विपच्यते}}{\Afootnote{H; \skt{कृच्छ्राद्विपच्यते}\ A1938; \skt{कृच्छ्राद्विपच्यते}\ N}}\skt{ }\edtext{\edtext{\skt{भुक्तन्द्वितीयन्न}}{\lemma{\skt{भुक्तन्}}\Afootnote{A1938 H; \skt{भक्तं}\ N}}\skt{ च}}{\lemma{\skt{द्वितीयन्न च}}\Afootnote{N H; \skt{द्वितीयं च न}\ A1938}}\skt{ कांक्षति ॥}\pend

\pstart \ledsidenote{\passref{1.46.474cd}}\edtext{\skt{आलस्यगौरवाटोपमचिङ्कुरुतेऽधिकं}}{\lemma{\skt{आलस्यगौरवाटोपमचिङ्}}\Afootnote{H; \skt{आलस्यगौरवाटोपसादांश्च}\ A1938; \skt{आलस्यं गौरवाटोपमरुचिं}\ N}}\skt{ ।}\pend

\pstart \ledsidenote{\passref{1.46.474ab}}\edtext{\edtext{\skt{हीनमात्रमसन्तोषं}}{\lemma{\skt{हीनमात्रम्}}\Afootnote{N H; \skt{हीनमात्र}\ A1938}}}{\lemma{\skt{असन्तोषं}}\Afootnote{N; \skt{मसन्तोषं}\ A1938; \skt{असन्तोष}\ H}}\skt{ }\edtext{\skt{करोति च}}{\Afootnote{A1938 N; \skt{न्करोति च}\ H}}\skt{ }\edtext{\skt{वलक्षयं}}{\Afootnote{H; \skt{बलक्षयम्}\ A1938 N}}\skt{ ॥}\pend

\pstart\versepar \ledsidenote{\passref{1.46.475}}\edtext{\skt{तस्मात्सुसंस्कृतं}}{\Afootnote{A1938 H; \skt{तस्मात्सुसस्कृतं}\ N}}\skt{ युक्त्या }\edtext{\edtext{\skt{दोषैरेभिर्व्विवर्जितं}}{\lemma{\skt{दोषैर्}}\Afootnote{A1938 H; \skt{दौषैर्}\ N\textsuperscript{ac}; \skt{दोषैर्}\ N\textsuperscript{pc}}}}{\lemma{\skt{एभिर्}}\Afootnote{N H; \skt{एतैर्}\ A1938}}\skt{ । }\pend
\pstart\versepar \skt{यथोक्तगुणसम्पन्नमुपसेवेत भोजनं ॥}\pend

\pstart\versepar \ledsidenote{\passref{1.46.476}}\skt{विभज्य }\edtext{\edtext{\skt{कालदोषादीन्कालयोरुभयोरपि}}{\lemma{\skt{कालदोषादीन्कालयोर्}}\Afootnote{H; \skt{दोषकालादीन्कालयोर्}\ A1938; \skt{कालदोषादीङ्कालयोर्}\ N}}}{\lemma{\skt{अपि}}\Afootnote{A1938 H; \skt{अपिं ॥}\ N}}\skt{ । }\pend
\pstart\versepar \skt{}\edtext{\edtext{\edtext{\skt{अचोक्षन्दुष्टमुत्सृष्टम्पाषाणतृणलोष्टवत्}}{\lemma{\skt{अचोक्षन्}}\Afootnote{A1938 H; \skt{अचौक्ष}\ N}}}{\lemma{\skt{दुष्टम्}}\Afootnote{N H; \skt{दुष्ट}\ A1938}}}{\lemma{\skt{उत्सृष्टम्}}\Afootnote{H; \skt{मुत्सृष्टं}\ A1938; \skt{उच्छिष्टम्}\ N}}\skt{ ॥}\pend

\pstart\versepar \ledsidenote{\passref{1.46.477}}\skt{द्विष्टं }\edtext{\edtext{\skt{व्युषितमस्वादु}}{\lemma{\skt{व्युषितम्}}\Afootnote{N H; \skt{व्युषित}\ A1938}}}{\lemma{\skt{अस्वादु}}\Afootnote{N H; \skt{मस्वादु}\ A1938}}\skt{ }\edtext{\edtext{\skt{पूतिमन्नम्विवर्ज्जयेत्}}{\lemma{\skt{पूतिम्}}\Afootnote{N H; \skt{पूति}\ A1938}}}{\lemma{\skt{अन्नम्}}\Afootnote{N H; \skt{चान्नं}\ A1938}}\skt{ । }\pend
\pstart\versepar \skt{चिरसिद्धं स्थिरं }\edtext{\edtext{\edtext{\skt{शीतमन्नमुष्णीकृतम्पुनः}}{\lemma{\skt{शीतम्}}\Afootnote{N H; \skt{शीत}\ A1938}}}{\lemma{\skt{अन्नम्}}\Afootnote{N H; \skt{मन्न}\ A1938}}}{\lemma{\skt{उष्णीकृतम्}}\Afootnote{N H; \skt{मुष्णीकृतं}\ A1938}}\skt{ ॥}\pend

\pstart\versepar \ledsidenote{\passref{1.46.478}}\edtext{\skt{अशान्तमुपदग्धञ्च न}}{\lemma{\skt{उपदग्धञ्च न}}\Afootnote{N H; \skt{उपदग्धं च}\ A1938}}\skt{ तथा स्वादु न लक्ष्यते । }\pend
\pstart\versepar \skt{}\edtext{\edtext{\skt{यद्यत्स्वादुतरन्तत्तद्विदध्यादुत्तरोत्तरं}}{\lemma{\skt{यद्यत्स्वादुतरन्}}\Afootnote{N H; \skt{यद्यत्स्वादुतरं}\ A1938}}}{\lemma{\skt{तत्तद्विदध्यादुत्तरोत्तरं}}\Afootnote{N; \skt{तत्तद्विदध्यादुत्तरोत्तरम् ॥४७८॥}\ A1938; \skt{तत्तद्विदध्यादुत्तरोत्तरं ॥}\ H}}\skt{ ॥}\pend

\pstart\versepar \ledsidenote{\passref{1.46.479}}\edtext{\skt{प्रक्षालयेद्भिरास्यं}}{\Afootnote{H; \skt{प्रक्षालयेदद्भिरास्यं}\ A1938 N}}\skt{ }\edtext{\skt{भुञ्जानस्य}}{\Afootnote{A1938 H; \skt{भुंञ्जानस्य}\ N}}\skt{ }\edtext{\edtext{\skt{मुहुर्म्मुहुः}}{\lemma{\skt{मुहुर्}}\Afootnote{N H; \skt{मुहु}\ A1938}}}{\lemma{\skt{म्मुहुः}}\Afootnote{N H; \skt{र्मुहुः ।}\ A1938}}\skt{ । }\pend
\pstart\versepar \skt{विशुद्धे }\edtext{\skt{रसने }\edtext{\skt{ह्यस्मै}}{\lemma{\skt{अस्मै}}\Afootnote{N H; \skt{तस्य}\ A1938}}}{\lemma{\skt{रसने ह्य्}}\Afootnote{N H; \skt{रसने}\ A1938}}\skt{ }\edtext{\edtext{\edtext{\skt{रोचतेऽन्नमपूर्ववत्}}{\lemma{\skt{रोचते}}\Afootnote{A1938 H; \skt{रोचते न्}\ N}}}{\lemma{\skt{ऽन्नम्}}\Afootnote{H; \skt{ऽन्न}\ A1938; \skt{न}\ N}}}{\lemma{\skt{अपूर्ववत्}}\Afootnote{H; \skt{मपूर्ववत्}\ A1938 N}}\skt{ ॥}\pend

\pstart\versepar \ledsidenote{\passref{1.46.480}}\skt{स्वादुना तस्य रसनं }\edtext{\skt{प्रथमेनाथ}}{\Afootnote{H; \skt{प्रथमेनाति}\ A1938; \skt{पथमेनाथ}\ N}}\skt{ तर्प्पितं । }\pend
\pstart\versepar \skt{}\edtext{\skt{तथा न}}{\Afootnote{H; \skt{न तथा}\ A1938; \skt{तथान्न}\ N}}\skt{ }\edtext{\edtext{\edtext{\skt{स्वादयेदन्नन्तस्मात्प्रक्षाल्यमन्तरा}}{\lemma{\skt{स्वादयेदन्नन्}}\Afootnote{H; \skt{स्वादयेदन्य}\ A1938; \skt{स्वादुयेदन्नन्}\ N}}}{\lemma{\skt{तस्मात्प्रक्षाल्यम्}}\Afootnote{N H; \skt{त्तस्मात्प्रक्षाल्य}\ A1938}}}{\lemma{\skt{अन्तरा}}\Afootnote{N H; \skt{मन्तरा ॥४८०॥}\ A1938}}\skt{ ॥}\pend

\pstart\versepar \ledsidenote{\passref{1.46.481}}\edtext{\skt{सौमनस्यम्वलं}}{\lemma{\skt{वलं}}\Afootnote{H; \skt{बलं}\ A1938 N}}\skt{ }\edtext{\edtext{\skt{तुष्टिमुत्साहं}}{\lemma{\skt{तुष्टिम्}}\Afootnote{N H; \skt{पुष्टि}\ A1938}}}{\lemma{\skt{उत्साहं}}\Afootnote{N H; \skt{मुत्साहं}\ A1938}}\skt{ हर्षणं सुखं । }\pend
\pstart\versepar \skt{स्वादु }\edtext{\edtext{\edtext{\skt{सञ्जनयत्यन्नमस्वादु}}{\lemma{\skt{सञ्जनयत्य्}}\Afootnote{N; \skt{सञ्जनय}\ A1938 H}}}{\lemma{\skt{अन्नम्}}\Afootnote{N; \skt{त्यन्न}\ A1938; \skt{त्यन्नम्}\ H}}\skt{ तु}}{\lemma{\skt{अस्वादु तु}}\Afootnote{N H; \skt{मस्वादु च}\ A1938}}\skt{ }\edtext{\skt{विवर्जयेत्}}{\Afootnote{N H; \skt{विपर्ययम् ॥४८१॥}\ A1938}}\skt{ ॥}\pend

\pstart \ledsidenote{\passref{1.46.482}}\skt{भुक्त्वा }\edtext{\skt{पुनः}}{\Afootnote{N H; \skt{पि य}\ A1938}}\skt{ }\edtext{\skt{प्रार्थयते}}{\Afootnote{N H; \skt{त्प्रार्थयते}\ A1938}}\skt{ }\edtext{\skt{भूयस्तत्स्वादु}}{\Afootnote{A1938 H; \skt{भूयस्तत्स्वादु}\ N}}\skt{ }\edtext{\skt{भोजनं}}{\Afootnote{N H; \skt{भोजनमशितश्चोदकं युक्त्या भुञ्जानश्चान्तरा पिबेत् ॥४८२॥}\ A1938}}\skt{ ।}\pend

\pstart\versepar \ledsidenote{\passref{1.46.483}}\skt{दन्तान्तरगतं चान्नं }\edtext{\skt{शोधनेनाहरेच्छनैः}}{\Afootnote{A1938; \skt{सोधनेनाहरेच्छन्नैः ।}\ N; \skt{शोधनेनाहरेच्छन्नैः ॥}\ H}}\skt{ ॥ }\pend
\pstart\versepar \skt{}\edtext{\skt{कुर्यादनिर्हृतन्तद्धि}}{\lemma{\skt{तद्धि}}\Afootnote{A1938 H; \skt{तद्धि}\ N}}\skt{ }\edtext{\skt{मुखस्यानिष्टगन्धितां}}{\Afootnote{N H; \skt{मुखस्यानिष्टगन्धताम् ॥४८३॥}\ A1938}}\skt{ ।}\pend

\pstart\versepar \ledsidenote{\passref{1.46.484}}\edtext{\edtext{\skt{जीर्ण्णेऽन्ने}}{\lemma{\skt{जीर्ण्णे}}\Afootnote{A1938; \skt{जीर्णे न्}\ N H}}}{\lemma{\skt{ऽन्ने}}\Afootnote{A1938; \skt{ने}\ N H}}\skt{ }\edtext{\skt{वर्ते}}{\Afootnote{H; \skt{वर्धते}\ A1938 N}}\skt{ वायुर्व्विदग्धे पित्तमेव तु ॥ }\pend
\pstart\versepar \skt{भुक्तमात्रे कफश्चापि }\edtext{\edtext{\skt{तस्माद्भुक्तेरितङ्कफं}}{\lemma{\skt{तस्माद्भुक्तेरितङ्}}\Afootnote{A1938; \skt{तस्माद्भुंक्तेरितङ्}\ N; \skt{तस्माद्भुक्तेरितङ्}\ H}}}{\lemma{\skt{कफं}}\Afootnote{A1938 H; \skt{फलं ॥}\ N}}\skt{ ॥}\pend

\pstart\versepar \ledsidenote{\passref{1.46.485}}\edtext{\skt{हरेद्धूमेन}}{\Afootnote{N H; \skt{धूमेना}\ A1938}}\skt{ }\edtext{\skt{हृद्यैर्व्वा}}{\Afootnote{N H; \skt{पोह्य हृद्यैर्वा}\ A1938}}\skt{ कषायकटुतिक्तकैः । }\pend
\pstart\versepar \skt{}\edtext{\skt{पूगकङ्कोलकर्पूरलवङ्गसुमनः}}{\Afootnote{A1938 H; \skt{पूगकक्कोलकर्प्पूरलवंगसुमनः}\ N}}\skt{ फलैः ।}\pend

\pstart\versepar \ledsidenote{\passref{1.46.486}}\edtext{\skt{फलैः}}{\Afootnote{A1938 N\textsuperscript{pc} H; \emph{om.}\ N\textsuperscript{ac}}}\skt{ }\edtext{\skt{कटुकसार्व्वा}}{\Afootnote{H; \skt{कटुकषायैर्वा}\ A1938; \skt{कटुकसाह्वैर्वा}\ N}}\skt{ }\edtext{\skt{मुखवैषद्यकारकैः}}{\Afootnote{H; \skt{मुखवैशद्यकारकैः ।}\ A1938; \skt{मुखवैसद्यकारकैः ।}\ N}}\skt{ ॥ }\pend
\pstart\versepar \skt{}\edtext{\skt{ताम्वुलपत्रसहितैः}}{\Afootnote{H; \skt{ताम्बूलपत्रसहितैः}\ A1938 N}}\skt{ सुगन्धैर्व्वा विचक्षणः ।}\pend

\pstart \ledsidenote{\passref{1.46.487}}\edtext{\skt{ततः}}{\Afootnote{N H; \skt{भुक्त्वा राजवदासीत यावदन्नक्लमो गतः । ततः}\ A1938}}\skt{ }\edtext{\skt{पदशतं}}{\Afootnote{H; \skt{पादशतं}\ A1938; \skt{पदं शतं}\ N}}\skt{ गत्वा वामपार्श्वेन }\edtext{\skt{सम्विशेत्}}{\Afootnote{A1938 H; \skt{सम्विसेत् ।}\ N}}\skt{ ॥}\pend

\pstart\versepar \ledsidenote{\passref{1.46.488}}\edtext{\skt{शव्दं}}{\Afootnote{H; \skt{शब्दा}\ A1938; \skt{सब्दं}\ N}}\skt{ }\edtext{\skt{रूपं}}{\Afootnote{N H; \skt{न्रूपा}\ A1938}}\skt{ }\edtext{\edtext{\skt{रसङ्गन्धं}}{\lemma{\skt{रसङ्}}\Afootnote{N H; \skt{न्रसा}\ A1938}}}{\lemma{\skt{गन्धं}}\Afootnote{N H; \skt{न्गन्धा}\ A1938}}\skt{ }\edtext{\skt{सेवेत}}{\Afootnote{N H; \skt{न्स्पर्शांश्च}\ A1938}}\skt{ मनसः }\edtext{\skt{प्रियं}}{\Afootnote{H; \skt{प्रिया}\ A1938; \skt{यियं ॥}\ N}}\skt{ । }\pend
\pstart\versepar \skt{}\edtext{\skt{भुक्तमात्रः}}{\Afootnote{N H; \skt{न् । भुक्तवानु}\ A1938}}\skt{ }\edtext{\skt{शुचौ}}{\Afootnote{N H; \skt{पसे}\ A1938}}\skt{ }\edtext{\skt{देशे}}{\Afootnote{N H; \skt{वेत ते}\ A1938}}\skt{ नान्नं }\edtext{\skt{चेत्साधु}}{\Afootnote{N H; \skt{साधु}\ A1938}}\skt{ तिष्ठति ॥}\pend

\pstart\versepar \ledsidenote{\passref{1.46.489}}\edtext{\skt{शव्दं}}{\Afootnote{H; \skt{शब्द}\ A1938; \skt{सब्दं}\ N}}\skt{ }\edtext{\skt{रूपं}}{\Afootnote{N H; \skt{रूप}\ A1938}}\skt{ }\edtext{\skt{रसं}}{\Afootnote{N H; \skt{रसा}\ A1938}}\skt{ }\edtext{\skt{गन्धं}}{\Afootnote{N H; \skt{गन्धाः}\ A1938}}\skt{ }\edtext{\skt{स्पर्शश्चापि}}{\lemma{\skt{स्पर्शश्}}\Afootnote{H; \skt{स्पर्शाश्}\ A1938; \skt{स्पर्शञ्}\ N}}\skt{ }\edtext{\skt{जुगुप्सितं}}{\Afootnote{H; \skt{जुगुप्सिताः ।}\ A1938; \skt{जुगुप्सितां ॥}\ N}}\skt{ । }\pend
\pstart\versepar \skt{}\edtext{\skt{अशुच्यन्नन्तथा}}{\lemma{\skt{अशुच्यन्नन्}}\Afootnote{A1938 H; \skt{अशुचियन्नन्}\ N\textsuperscript{ac}; \skt{अशुच्यन्नन्}\ N\textsuperscript{pc}}}\skt{ }\edtext{\edtext{\skt{भुक्तमतिमात्रञ्च}}{\lemma{\skt{भुक्तम्}}\Afootnote{H; \skt{भुक्त}\ A1938; \skt{भुक्ताम्}\ N}}}{\lemma{\skt{अतिमात्रञ्च}}\Afootnote{N H; \skt{मतिहास्यं च}\ A1938}}\skt{ वामयेत् ॥}\pend

\pstart \ledsidenote{\passref{1.46.490}}\edtext{\edtext{\skt{शयनम्वासनम्वापि}}{\lemma{\skt{शयनम्}}\Afootnote{A1938 H; \skt{शयनाम्}\ N}}}{\lemma{\skt{वासनम्}}\Afootnote{H; \skt{चासनं}\ A1938; \skt{वासनाम्}\ N}}\skt{ }\edtext{\skt{नेच्छेद्वापि}}{\Afootnote{N H; \skt{नेच्छेद्वाऽपि}\ A1938}}\skt{ }\edtext{\skt{च तल्लक्षणं}}{\Afootnote{H; \skt{द्रवोत्तरम् । नाग्न्यातपौ न प्लवनं न यानं नापि वाहनम् ॥४९०॥}\ A1938; \skt{च तत्क्षणं ॥}\ N}}\skt{ ।}\pend

\pstart\versepar \ledsidenote{\passref{1.46.491}}\skt{न चैकरससेवायां }\edtext{\skt{प्रसज्ये च}}{\Afootnote{H; \skt{प्रसज्येत}\ A1938; \skt{प्रसजेत}\ N}}\skt{ }\edtext{\skt{कदाचन}}{\Afootnote{A1938 H; \skt{कदाचनः ।}\ N}}\skt{ ॥ }\pend
\pstart\versepar \skt{}\edtext{\skt{शाकावरान्नभूयिष्ठं}}{\Afootnote{N\textsuperscript{pc} H; \skt{शाकावरान्नभूयिष्ठ}\ A1938; \skt{शाकावारान्नभूयिष्ठं}\ N\textsuperscript{ac}}}\skt{ }\edtext{\skt{व्यम्लञ्च न}}{\Afootnote{H; \skt{मम्लं च न}\ A1938; \skt{व्यम्ब्लञ्च न}\ N}}\skt{ समाचरेत् ।}\pend

\pstart\versepar \ledsidenote{\passref{1.46.492}}\skt{एकैकशः }\edtext{\skt{समस्ताम्वा}}{\Afootnote{N H; \skt{समस्तान्वा}\ A1938}}\skt{ }\edtext{\skt{नाश्नीयाच्च}}{\Afootnote{N H; \skt{नाध्यश्नीयाद्र}\ A1938}}\skt{ }\edtext{\skt{कदाचन}}{\Afootnote{H; \skt{सान्सदा}\ A1938; \skt{कदाचनः ।}\ N}}\skt{ ॥ }\pend
\pstart\versepar \skt{}\edtext{\edtext{\edtext{\skt{प्राग्भुक्तेऽप्यविविक्तेऽग्नौ}}{\lemma{\skt{ऽप्य्}}\Afootnote{H; \skt{त्व}\ A1938; \skt{प्य्}\ N}}}{\lemma{\skt{अविविक्ते}}\Afootnote{N H; \skt{विविक्ते}\ A1938}}}{\lemma{\skt{ऽग्नौ}}\Afootnote{A1938 H; \skt{ग्नौ}\ N}}\skt{ द्विरन्नन्न समाचरेत् ।}\pend

\pstart \ledsidenote{\passref{1.46.493}}\edtext{\edtext{\skt{मात्रागुरूम्परिहरेदाहारं}}{\lemma{\skt{मात्रागुरूम्}}\Afootnote{H; \skt{पूर्वभुक्ते विदग्धेऽन्ने भुञ्जानो हन्ति पावकम् । मात्रागुरुं}\ A1938; \skt{मात्र्रात्रागुरुम्}\ N}}}{\lemma{\skt{परिहरेदाहारं}}\Afootnote{N H; \skt{परिहरेदाहारं}\ A1938}}\skt{ }\edtext{\skt{द्रव्यतोऽपि च}}{\lemma{\skt{ऽपि च}}\Afootnote{N H; \skt{च यः ॥४९३॥}\ A1938}}\skt{ ॥}\pend

\pstart\versepar \ledsidenote{\passref{1.46.494}}\edtext{\skt{पिष्टान्न}}{\Afootnote{H; \skt{पिष्टान्नं}\ A1938 N}}\skt{ नैव }\edtext{\skt{सेवत}}{\Afootnote{N H; \skt{भुञ्जीत}\ A1938}}\skt{ मात्रया वा }\edtext{\skt{क्षुधातुरः}}{\Afootnote{N H; \skt{बुभुक्षितः । द्विगुणं च पिबेत्तोयं सुखं}\ A1938}}\skt{ । }\pend
\pstart\versepar \skt{}\edtext{\skt{पेयलेह्याद्यभक्ष्याणां}}{\Afootnote{H; \skt{सम्यक्प्रजीर्यति पेयलेह्याद्यभक्ष्याणां}\ A1938; \skt{पेयलेह्यादिभक्षाणां}\ N}}\skt{ }\edtext{\skt{गुरूविद्याद्यथोत्तरं}}{\Afootnote{N H; \skt{गुरु विद्याद्यथोत्तरम् ॥४९४॥}\ A1938}}\skt{ ॥}\pend

\pstart\versepar \ledsidenote{\passref{1.46.495}}\edtext{\edtext{\skt{गुरूणामधसौहित्यं}}{\lemma{\skt{गुरूणाम्}}\Afootnote{N H; \skt{गुरूणा}\ A1938}}}{\lemma{\skt{अधसौहित्यं}}\Afootnote{H; \skt{मर्धसौहित्यं}\ A1938; \skt{अर्धशौनित्यं}\ N}}\skt{ }\edtext{\skt{लfहूनां}}{\Afootnote{H; \skt{लघूनां}\ A1938 N}}\skt{ }\edtext{\edtext{\skt{वृत्तिरिष्यते}}{\lemma{\skt{वृत्तिर्}}\Afootnote{N H; \skt{तृप्ति}\ A1938}}}{\lemma{\skt{इष्यते}}\Afootnote{N H; \skt{रिष्यते ।}\ A1938}}\skt{ । }\pend
\pstart\versepar \skt{द्रवोत्तरो द्रवश्चापि न मात्रा }\edtext{\edtext{\skt{गुरुरिष्यते}}{\lemma{\skt{गुरुर्}}\Afootnote{N H; \skt{गुरु}\ A1938}}}{\lemma{\skt{इष्यते}}\Afootnote{N H; \skt{रिष्यते ॥४९५॥}\ A1938}}\skt{ ॥}\pend

\pstart\versepar \ledsidenote{\passref{1.46.496}}\edtext{\edtext{\skt{द्रवाद्यमविशुष्कन्तु}}{\lemma{\skt{द्रवाद्यम्}}\Afootnote{N H; \skt{द्रवाढ्यम्}\ A1938}}}{\lemma{\skt{अविशुष्कन्तु}}\Afootnote{N H; \skt{अपि शुष्कं तु}\ A1938}}\skt{ }\edtext{\edtext{\skt{सम्यक्कान्युपपद्यते}}{\lemma{\skt{सम्यक्कान्य्}}\Afootnote{H; \skt{सम्यगेवो}\ A1938; \skt{सम्यक्खान्य्}\ N}}}{\lemma{\skt{उपपद्यते}}\Afootnote{N H; \skt{पपद्यते}\ A1938}}\skt{ । }\pend
\pstart\versepar \skt{}\edtext{\edtext{\edtext{\skt{विशुष्कमन्नमभ्यस्तन्न}}{\lemma{\skt{विशुष्कम्}}\Afootnote{N H; \skt{विशुष्क}\ A1938}}}{\lemma{\skt{अन्नम्}}\Afootnote{N H; \skt{मन्न}\ A1938}}}{\lemma{\skt{अभ्यस्तन्न}}\Afootnote{N H; \skt{मभ्यस्तं न}\ A1938}}\skt{ पाकं साधु गच्छति ॥}\pend

\pstart\versepar \ledsidenote{\passref{1.46.497}}\edtext{\skt{पिण्डीकृतसमं क्लिन्नम्विदाहमुपगच्छति}}{\lemma{\skt{पिण्डीकृतसमं क्लिन्नम्}}\Afootnote{H; \skt{पिण्डीकृतमसङ्क्लिन्नं}\ A1938; \skt{विशुष्कपिण्डीकृतमसंक्लिन्नम्}\ N\textsuperscript{ac}; \skt{पिण्डीकृतमसंक्लिन्नम्}\ N\textsuperscript{pc}}}\skt{ । }\pend
\pstart\versepar \skt{}\edtext{\edtext{\skt{श्रोतस्यन्नवहे}}{\lemma{\skt{श्रोतस्य्}}\Afootnote{N H; \skt{स्रोत}\ A1938}}}{\lemma{\skt{अन्नवहे}}\Afootnote{N H; \skt{स्यन्नवहे}\ A1938}}\skt{ पित्तम्पक्तौ वा यस्य तिष्ठति ॥}\pend

\pstart\versepar \ledsidenote{\passref{1.46.498}}\skt{विदाहि }\edtext{\edtext{\skt{भुक्तमन्यद्वा}}{\lemma{\skt{भुक्तम्}}\Afootnote{N H; \skt{भुक्त}\ A1938}}}{\lemma{\skt{अन्यद्वा}}\Afootnote{N H; \skt{मन्यद्वा}\ A1938}}\skt{ }\edtext{\edtext{\skt{तस्याप्यन्नम्विदह्यते}}{\lemma{\skt{तस्याप्य्}}\Afootnote{H; \skt{तस्या}\ A1938; \skt{तथाप्य्}\ N}}}{\lemma{\skt{अन्नम्}}\Afootnote{N H; \skt{प्यन्नं}\ A1938}}\skt{ । }\pend
\pstart\versepar \skt{}\edtext{\edtext{\skt{शुष्कम्भुक्तम्विदग्धं}}{\lemma{\skt{भुक्तम्}}\Afootnote{N H; \skt{विरुद्धं}\ A1938}}}{\lemma{\skt{विदग्धं}}\Afootnote{\skt{विष्टम्}\ A1938; \skt{विदग्ध}\ N; \skt{विदग्धं ः}\ H}}\skt{ }\edtext{\edtext{\skt{स्यादग्धेर्व्यापादकारकः}}{\lemma{\skt{स्यादग्धेर्}}\Afootnote{\skt{भि वह्नि}\ A1938; \skt{स्यादग्नेर्}\ N; \skt{स्यादग्धेर्}\ H}}}{\lemma{\skt{व्यापादकारकः}}\Afootnote{H; \skt{व्यापदमावहेत् ॥४९८॥}\ A1938; \skt{व्यापादकारकं ।}\ N}}\skt{ ॥}\pend

\pstart\versepar \ledsidenote{\passref{1.46.499}}\edtext{\edtext{\skt{आमम्विदधम्विष्टच्च}}{\lemma{\skt{विदधम्}}\Afootnote{H; \skt{विदग्धं}\ A1938; \skt{विदध}\ N}}}{\lemma{\skt{विष्टच्च}}\Afootnote{H; \skt{विष्टब्धं}\ A1938 N}}\skt{ }\edtext{\edtext{\skt{कफपित्तानिलैस्त्रिभिः}}{\lemma{\skt{कफपित्तानिलैस्}}\Afootnote{H; \skt{कफपित्तानिलै}\ A1938; \skt{कफपित्तनिलेस्}\ N}}}{\lemma{\skt{त्रिभिः}}\Afootnote{N H; \skt{स्त्रिभिः ।}\ A1938}}\skt{ । }\pend
\pstart\versepar \skt{}\edtext{\skt{अजीर्ण्णे}}{\Afootnote{H; \skt{अजीर्णं}\ A1938 N}}\skt{ }\edtext{\skt{केचिदिच्छन्ति}}{\Afootnote{N H; \skt{केचिदिच्छन्ति}\ A1938}}\skt{ चतुर्थं रसशेषतः ॥}\pend

\pstart\versepar \ledsidenote{\passref{1.46.500}}\edtext{\edtext{\skt{अत्यम्बुपानाद्विषमाशनाच्च}}{\lemma{\skt{अत्यम्बुपानाद्विषम्}}\Afootnote{N; \skt{अत्यम्बुपानाद्विष}\ A1938; \skt{अत्यम्वुपानाद्विषम्}\ H}}}{\lemma{\skt{आशनाच्च}}\Afootnote{H; \skt{माशनाद्वा}\ A1938; \skt{आसनाच्च}\ N}}\skt{ }\pend
\pstart\versepar \skt{}\edtext{\skt{सन्धारणात्स्वप्नविपर्ययाद्वा}}{\Afootnote{N H; \skt{सन्धारणात्स्वप्नविपर्ययाच्च ।}\ A1938}}\skt{ । }\pend
\pstart\versepar \skt{}\edtext{\skt{कालेऽपि}}{\lemma{\skt{ऽपि}}\Afootnote{A1938 H; \skt{पि}\ N}}\skt{ सात्म्यं लघु चापि }\edtext{\skt{भुक्तम्}}{\Afootnote{N H; \skt{भुक्त}\ A1938}}\skt{ }\pend
\pstart\versepar \skt{}\edtext{\skt{अन्नन्न}}{\Afootnote{N H; \skt{मन्नं न}\ A1938}}\skt{ पाकं भजते नरस्य ॥}\pend

\pstart\versepar \ledsidenote{\passref{1.46.502}}\edtext{\edtext{\skt{माधुर्यमन्नङ्गतमांससंज्ञाम्}}{\lemma{\skt{माधुर्यम्}}\Afootnote{H; \skt{माधुर्य}\ A1938; \skt{म्धूयम्}\ N\textsuperscript{ac}; \skt{माधूयम्}\ N\textsuperscript{pc}}}}{\lemma{\skt{अन्नङ्गतमांससंज्ञाम्}}\Afootnote{H; \skt{मन्नं गतमामसञ्ज्ञं}\ A1938; \skt{अन्नं गतमांससंज्ञाम्}\ N}}\skt{ }\pend
\pstart\versepar \skt{}\edtext{\skt{विदग्धसंज्ञां}}{\Afootnote{H; \skt{विदग्धसञ्ज्ञं}\ A1938; \skt{विदग्धसंज्ञा}\ N}}\skt{ }\edtext{\edtext{\skt{गतमम्लभावं}}{\lemma{\skt{गतम्}}\Afootnote{N H; \skt{गत}\ A1938}}}{\lemma{\skt{अम्लभावं}}\Afootnote{H; \skt{मम्लभावम् ।}\ A1938; \skt{अम्ब्लभावः ।}\ N}}\skt{ । }\pend
\pstart\versepar \skt{}\edtext{\skt{किञ्चिद्विदग्धं}}{\Afootnote{N H; \skt{किञ्चिद्विपक्वं,}\ A1938}}\skt{ भृशतोदशूलम् }\pend
\pstart\versepar \skt{}\edtext{\skt{विष्टब्धमानद्धनिरुद्धवातं}}{\Afootnote{H; \skt{विष्टब्धमानद्धविरुद्धवातम् ॥५०२॥}\ A1938; \skt{विष्टब्धमानन्तु निरुद्धवातं ॥}\ N}}\skt{ ॥}\pend

\pstart\versepar \ledsidenote{\passref{1.46.503}}\edtext{\edtext{\skt{उद्गारशुद्धावपि}}{\lemma{\skt{उद्गारशुद्धाव्}}\Afootnote{N; \skt{उद्गारशुद्धा}\ A1938 H}}}{\lemma{\skt{अपि}}\Afootnote{N; \skt{वपि}\ A1938 H}}\skt{ भक्षकांक्षा }\pend
\pstart\versepar \skt{न जायते हृद्गुरुता च यस्य । }\pend
\pstart\versepar \skt{रसावशेषेण तु सप्रसेकं }\pend
\pstart\versepar \skt{}\edtext{\edtext{\skt{चतुर्थमेतन्प्रवदन्त्यजीर्ण्णं}}{\lemma{\skt{चतुर्थम्}}\Afootnote{N H; \skt{चतुर्थ}\ A1938}}}{\lemma{\skt{एतन्प्रवदन्त्य्}}\Afootnote{H; \skt{मेतत्प्रवदन्त्य्}\ A1938; \skt{एतत्प्रवदन्त्य्}\ N}}\skt{ ॥}\edtext{}{\lemma{\skt{भक्षकांक्षा न}}\Afootnote{H; \skt{भक्तकाङ्क्षा न}\ A1938 N}}\pend

\pstart\versepar \ledsidenote{\passref{1.46.504}}\skt{मूर्च्छा प्रलापो वमथुः प्रसेको }\pend
\pstart\versepar \skt{}\edtext{\edtext{\skt{ज्वरोऽतिसारः}}{\lemma{\skt{ज्वरो}}\Afootnote{N H; \emph{om.}\ A1938}}}{\lemma{\skt{ऽतिसारः}}\Afootnote{N H; \emph{om.}\ A1938}}\skt{ }\edtext{\skt{सदनम्भ्रमश्च}}{\lemma{\skt{भ्रमश्च}}\Afootnote{N H; \skt{भ्र}\ A1938}}\skt{ ॥ }\pend
\pstart\versepar \skt{शिरोरुजा पृष्ठकटिग्रहश्च }\pend
\pstart\versepar \skt{तृष्णा }\edtext{\edtext{\skt{विपाकोऽथ}}{\lemma{\skt{विपाको}}\Afootnote{H; \skt{विपाको}\ N}}}{\lemma{\skt{ऽथ}}\Afootnote{H; \skt{थ}\ N}}\skt{ विजृम्भिका च ॥ }\pend
\pstart\versepar \skt{उपद्रवा भवन्त्येते }\edtext{\skt{मरणम्वाप्यजीर्ण्णतः}}{\lemma{\skt{वाप्य्}}\Afootnote{N H; \skt{चाप्य्}\ A1938}}\skt{ ।}\edtext{}{\lemma{\skt{मश्… च}}\Afootnote{\emph{om.}\ A1938}}\pend

\pstart\versepar \ledsidenote{\passref{1.46.505}}\skt{तत्रामे }\edtext{\skt{लंघनङ्कार्यम्विदग्धे}}{\lemma{\skt{लंघनङ्}}\Afootnote{A1938 H; \skt{लघनङ्}\ N}}\skt{ वमनं हितं ॥ }\pend
\pstart\versepar \skt{विष्टब्धे स्वेदनं }\edtext{\skt{शस्तं}}{\Afootnote{H; \skt{पथ्यं,}\ A1938; \skt{सस्तं}\ N}}\skt{ }\edtext{\skt{रसशेषे}}{\Afootnote{A1938 H; \skt{रससेषे}\ N}}\skt{ शयीत च ।}\pend

\pstart\versepar \ledsidenote{\passref{1.46.506}}\edtext{\skt{वामयेदाशु }\edtext{\skt{तन्तस्मादुष्णेन}}{\lemma{\skt{तस्मादुष्णेन}}\Afootnote{A1938 H; \skt{त्तस्मादुष्णेन}\ N}}}{\lemma{\skt{वामयेदाशु तन्}}\Afootnote{A1938 H; \skt{वामयेदासु त}\ N}}\skt{ }\edtext{\skt{लवणाम्बुना}}{\Afootnote{A1938; \skt{लवणाबुना ॥}\ N; \skt{लवणाम्वुना ॥}\ H}}\skt{ ॥ }\pend
\pstart\versepar \skt{}\edtext{\edtext{\skt{कार्यन्चानशनन्तावद्यावन्न}}{\lemma{\skt{कार्यन्चानशनन्}}\Afootnote{H; \skt{कार्यं वाऽनशनं}\ A1938 N}}}{\lemma{\skt{तावद्यावन्न}}\Afootnote{N H; \skt{तावद्यावन्न}\ A1938}}\skt{ }\edtext{\skt{प्रकृतिम्व्रजेत्}}{\lemma{\skt{व्रजेत्}}\Afootnote{N H; \skt{भजेत् ॥५०६॥}\ A1938}}\skt{ ।}\pend

\pstart\versepar \ledsidenote{\passref{1.46.507}}\edtext{\skt{लघुकायमनश्चैनं}}{\lemma{\skt{लघुकायमनश्}}\Afootnote{H; \skt{लघुकायमतश्}\ A1938; \skt{लघुकायमतश्}\ N}}\skt{ }\edtext{\skt{लघ्वन्नैः}}{\Afootnote{N H; \skt{लङ्घनैः}\ A1938}}\skt{ समुपाचरेत् ॥ }\pend
\pstart\versepar \skt{यावन्न प्रकृतिस्थः }\edtext{\skt{स्याद्दोषतः}}{\Afootnote{N H; \skt{स्याद्दोषतः}\ A1938}}\skt{ }\edtext{\skt{घ्राणतस्तथा}}{\lemma{\skt{घ्राणतस्}}\Afootnote{H; \skt{प्राणतस्}\ A1938 N}}\skt{ ।}\pend

\pstart\versepar \ledsidenote{\passref{1.46.508}}\edtext{\edtext{\skt{हिताहितोपसंयुक्तमन्नसंशमनं}}{\lemma{\skt{हिताहितोपसंयुक्तम्}}\Afootnote{N H; \skt{हिताहितोपसंयुक्त}\ A1938}}}{\lemma{\skt{अन्नसंशमनं}}\Afootnote{N H; \skt{मन्नं समशनं}\ A1938}}\skt{ स्मृतं ॥ }\pend
\pstart\versepar \skt{}\edtext{\skt{बहु}}{\Afootnote{A1938 N; \skt{वहु}\ H}}\skt{ }\edtext{\edtext{\skt{स्तोकमकाले}}{\lemma{\skt{स्तोकम्}}\Afootnote{N H; \skt{स्तोक}\ A1938}}\skt{ वा तं}}{\lemma{\skt{अकाले वा तं}}\Afootnote{H; \skt{मकाले वा त}\ A1938; \skt{अकाले वा त}\ N}}\skt{ }\edtext{\edtext{\skt{ज्ञेयम्विषमाशनं}}{\lemma{\skt{ज्ञेयम्}}\Afootnote{N H; \skt{ज्ज्ञेयं}\ A1938}}}{\lemma{\skt{विषमाशनं}}\Afootnote{A1938 H; \skt{विषमासनं ।}\ N}}\skt{ ।}\pend

\pstart\versepar \ledsidenote{\passref{1.46.509}}\edtext{\skt{साजीर्ण्णे}}{\Afootnote{N H; \skt{अजीर्णे}\ A1938}}\skt{ }\edtext{\skt{भुज्यते यस्तु}}{\Afootnote{N H; \skt{भुज्यते यत्तु}\ A1938}}\skt{ }\edtext{\edtext{\skt{तदध्यशनमुच्यते}}{\lemma{\skt{तदध्यशनम्}}\Afootnote{H; \skt{तदध्यशन}\ A1938; \skt{तदध्यसनम्}\ N}}}{\lemma{\skt{उच्यते}}\Afootnote{N H; \skt{मुच्यते ।}\ A1938}}\skt{ ॥ }\pend
\pstart\versepar \skt{}\edtext{\skt{तद्वदेनन्निहन्त्याशु}}{\lemma{\skt{तद्वदेनन्}}\Afootnote{N H; \skt{त्रयमेतन्}\ A1938}}\skt{ }\edtext{\skt{वहून्व्याधीन्करोति}}{\Afootnote{N H; \skt{बहून्व्याधीन्करोति}\ A1938}}\skt{ }\edtext{\skt{च}}{\Afootnote{N H; \skt{वा ॥५०९॥}\ A1938}}\skt{ ॥}\pend

\pstart\versepar \ledsidenote{\passref{1.46.510}}\skt{अन्नम्विदग्धं हि नरस्यशीघ्रं }\edtext{\skt{शीताम्वुना वै}}{\Afootnote{H; \skt{शीताम्बुना वै}\ A1938 N}}\skt{ }\edtext{\edtext{\skt{परिपाकमेति}}{\lemma{\skt{परिपाकम्}}\Afootnote{N H; \skt{परिपाक}\ A1938}}}{\lemma{\skt{एति}}\Afootnote{N H; \skt{मेति ।}\ A1938}}\skt{ । }\pend
\pstart\versepar \skt{}\edtext{\edtext{\skt{तद्ध्यस्य}}{\lemma{\skt{तद्ध्य्}}\Afootnote{\skt{त}\ A1938 H; \skt{तद्ध्य्}\ N}}}{\lemma{\skt{अस्य}}\Afootnote{N; \skt{द्ध्यस्य}\ A1938 H}}\skt{ }\edtext{\edtext{\skt{शैत्येन}}{\lemma{\skt{शैत्य्}}\Afootnote{A1938 H; \skt{शै}\ N}}}{\lemma{\skt{एन}}\Afootnote{A1938 H; \skt{त्येन}\ N}}\skt{ निहन्ति }\edtext{\edtext{\skt{पित्तमाक्लेदिभावाच्च}}{\lemma{\skt{पित्तम्}}\Afootnote{N H; \skt{पित्त}\ A1938}}\skt{ न}}{\lemma{\skt{आक्लेदिभावाच्च न}}\Afootnote{N H; \skt{माक्लेदिभावाच्च न}\ A1938}}\skt{ }\edtext{\edtext{\skt{यान्त्यधस्तात्}}{\lemma{\skt{यान्त्य्}}\Afootnote{H; \skt{य}\ A1938; \skt{यत्य्}\ N}}}{\lemma{\skt{अधस्तात्}}\Afootnote{N H; \skt{त्यधस्तात् ॥५१०॥}\ A1938}}\skt{ ॥}\pend

\pstart\versepar \ledsidenote{\passref{1.46.511}}\skt{विदह्यते यस्य तु }\edtext{\skt{भुक्तमात्रं}}{\Afootnote{N H; \skt{भुक्तमात्रे}\ A1938}}\skt{ }\edtext{\skt{दह्यन्ति}}{\Afootnote{N H; \skt{दह्येत}\ A1938}}\skt{ }\edtext{\skt{हृत्कोष्ठगलाश्च}}{\Afootnote{H; \skt{हृत्कोष्ठगलं च}\ A1938; \skt{हृत्कोष्ठगलाश्च}\ N}}\skt{ यस्य । }\pend
\pstart\versepar \skt{}\edtext{\skt{द्राक्षासितां}}{\Afootnote{H; \skt{द्राक्षाभयां}\ A1938; \skt{द्राक्षासिता}\ N}}\skt{ }\edtext{\skt{माक्षिकसंप्रयुक्ता}}{\Afootnote{N H; \skt{माक्षिकसम्प्रयुक्तां}\ A1938}}\skt{ }\edtext{\skt{लीढ्वाभयाम्वै}}{\Afootnote{N H; \skt{लीढ्वाऽभयां वा}\ A1938}}\skt{ ससुखं लभेत ॥}\pend

\pstart\versepar \ledsidenote{\passref{1.46.512}}\skt{भवेदजीर्ण्णं प्रति यस्य }\edtext{\skt{शङ्का}}{\Afootnote{A1938 H; \skt{सङ्का}\ N}}\skt{ }\edtext{\skt{स्निग्धस्य}}{\Afootnote{A1938 H; \skt{श्निग्धस्य}\ N}}\skt{ }\edtext{\edtext{\edtext{\skt{जन्तोर्व्वलिनोऽन्नकाले}}{\lemma{\skt{जन्तोर्}}\Afootnote{N H; \skt{जन्तो}\ A1938}}}{\lemma{\skt{व्वलिनो}}\Afootnote{H; \skt{र्बलिनो}\ A1938; \skt{बलि न्}\ N}}}{\lemma{\skt{ऽन्नकाले}}\Afootnote{A1938 H; \skt{नकाले ।}\ N}}\skt{ । }\pend
\pstart\versepar \skt{प्रातः }\edtext{\edtext{\edtext{\skt{सशुण्ठीमभयामशङ्को}}{\lemma{\skt{सशुण्ठीम्}}\Afootnote{H; \skt{सशुण्ठी}\ A1938; \skt{सगुण्ठीम्}\ N}}}{\lemma{\skt{अभयाम्}}\Afootnote{N H; \skt{मभया}\ A1938}}}{\lemma{\skt{अशङ्को}}\Afootnote{H; \skt{मशङ्को}\ A1938; \skt{असङ्कौ}\ N}}\skt{ भुञ्जीत }\edtext{\skt{सम्प्राश्य}}{\Afootnote{A1938 H; \skt{सम्प्रास्य}\ N}}\skt{ हितं हितार्थी ॥}\pend

\pstart\versepar \ledsidenote{\passref{1.46.513}}\skt{स्वल्पं यदा }\edtext{\skt{दोषविवन्धमामं}}{\Afootnote{H; \skt{दोषविबद्धमामं}\ A1938; \skt{दोषविवन्धमामं}\ N}}\skt{ }\pend
\pstart\versepar \skt{लीनन्न }\edtext{\skt{तेजः}}{\Afootnote{A1938 H; \skt{तेजं}\ N}}\skt{ }\edtext{\edtext{\skt{पथमावृणोति}}{\lemma{\skt{पथम्}}\Afootnote{N H; \skt{पथ}\ A1938}}}{\lemma{\skt{आवृणोति}}\Afootnote{N H; \skt{मावृणोति ।}\ A1938}}\skt{ । }\pend
\pstart\versepar \skt{}\edtext{\edtext{\edtext{\skt{भवन्त्यजीर्ण्णेऽपि}}{\lemma{\skt{भवन्त्य्}}\Afootnote{H; \skt{भव}\ A1938; \skt{भवत्य्}\ N}}}{\lemma{\skt{अजीर्ण्णे}}\Afootnote{N H; \skt{त्यजीर्णे}\ A1938}}}{\lemma{\skt{ऽपि}}\Afootnote{A1938 H; \skt{पि}\ N}}\skt{ तदा बुभुक्षा }\pend
\pstart\versepar \skt{सा }\edtext{\skt{मन्दवृद्धिम्विषवन्निहन्ति}}{\lemma{\skt{मन्दवृद्धिम्}}\Afootnote{H; \skt{मन्दबुद्धिं}\ A1938 N}}\skt{ ॥}\edtext{}{\lemma{\skt{बुभुक्षा सा}}\Afootnote{N; \skt{बुभुक्षा या}\ A1938; \skt{वुभुक्षा सा}\ H}}\pend

\pstart\versepar \ledsidenote{\passref{1.46.514}}\skt{अत }\edtext{\skt{ऊर्द्ध्वम्प्रवक्ष्यामि}}{\lemma{\skt{ऊर्द्ध्वम्}}\Afootnote{N H; \skt{उर्ध्वं}\ A1938}}\skt{ गुणानाङ्कर्म्मविस्तरं । }\pend
\pstart\versepar \skt{कर्म्मभिस्त्वनुमीयन्ते नाना द्रव्याश्रया गुणाः ॥}\pend

\pstart \ledsidenote{\passref{1.46.519cd}}\skt{दशाद्याः कर्म्मतः }\edtext{\skt{प्रोक्तास्तेषाङ्कर्म्मविशेषणैः}}{\lemma{\skt{कर्म्मविशेषणैः}}\Afootnote{A1938 H; \skt{कर्मविसेषणैः ।}\ N}}\skt{ ।}\pend

\pstart \ledsidenote{\passref{1.46.520ab}}\edtext{\edtext{\skt{दशैवाद्याम्प्रवक्ष्यामि}}{\lemma{\skt{दशैवाद्याम्}}\Afootnote{N H; \skt{दशैवान्या}\ A1938}}}{\lemma{\skt{प्रवक्ष्यामि}}\Afootnote{N H; \skt{न्प्रवक्ष्यामि}\ A1938}}\skt{ }\edtext{\edtext{\skt{द्रवादीस्तांश्छृणुश्व}}{\lemma{\skt{द्रवादीस्तांश्}}\Afootnote{H; \skt{द्रवादींस्तान्}\ A1938; \skt{द्रवादीन्स्तां}\ N}}}{\lemma{\skt{छृणुश्व}}\Afootnote{H; \skt{निबोध}\ A1938; \skt{च्छृणुष्व}\ N}}\skt{ मे ॥}\pend

\pstart\versepar \ledsidenote{\passref{1.46.515}}\edtext{\skt{शीतः}}{\Afootnote{H; \emph{om.}\ A1938; \skt{सीतः}\ N}}\skt{ }\edtext{\skt{प्रह्लादनस्तम्भी}}{\Afootnote{N H; \skt{ह्लादनः स्तम्भनः शीतो}\ A1938}}\skt{ मूर्च्छातृट्स्वेददाहजित् ॥१॥ }\pend
\pstart\versepar \skt{}\edtext{\skt{उष्णस्तद्विपरीतः}}{\lemma{\skt{तद्विपरीतः}}\Afootnote{N; \skt{स्तद्विपरीतः}\ A1938; \skt{स्तद्विपरीतः}\ H}}\skt{ स्यात्पाचनश्च }\edtext{\skt{विशेषतः}}{\Afootnote{A1938 H; \skt{विसेषतः ॥२ २}\ N}}\skt{ ॥२॥}\edtext{}{\lemma{\skt{मूर्च्छातृट्स्वेददाहजित् ॥१॥ उष्णस्}}\Afootnote{N; \skt{मूर्च्छातृट्स्वेददाहजितुष्ण}\ A1938 H}}\pend

\pstart\versepar \ledsidenote{\passref{1.46.516}}\edtext{\skt{स्नेहो}}{\Afootnote{N H; \skt{स्नेह}\ A1938}}\skt{ मार्द्दवकृत्स्निग्धो }\edtext{\skt{बलवर्ण्णकरस्तथा}}{\lemma{\skt{बलवर्ण्णकरस्}}\Afootnote{A1938 N; \skt{वलवर्ण्णकरस्}\ H}}\skt{ ॥३॥ }\pend
\pstart\versepar \skt{}\edtext{\edtext{\skt{रूक्षस्तद्विपरीतः}}{\lemma{\skt{रूक्षस्}}\Afootnote{N H; \skt{रूक्ष}\ A1938}}}{\lemma{\skt{तद्विपरीतः}}\Afootnote{N H; \skt{स्तद्विपरीतः}\ A1938}}\skt{ }\edtext{\skt{स्याद्विशेषात्स्तम्भनः}}{\Afootnote{N H; \skt{स्याद्विशेषात्स्तम्भनः}\ A1938}}\skt{ }\edtext{\skt{स्मृतः}}{\Afootnote{N H; \skt{खरः ॥५१६॥}\ A1938}}\skt{ ॥४॥}\pend

\pstart\versepar \ledsidenote{\passref{1.46.517}}\skt{पिच्छिलो }\edtext{\skt{जीवनश्लेषी}}{\Afootnote{N H; \skt{जीवनो बल्यः}\ A1938}}\skt{ सन्धानो }\edtext{\edtext{\skt{वृंहणस्तथा}}{\lemma{\skt{वृंहणस्}}\Afootnote{H; \skt{श्लेष्मलो}\ A1938; \skt{बृंहणस्}\ N}}}{\lemma{\skt{तथा}}\Afootnote{N H; \skt{गुरुः ।}\ A1938}}\skt{ ॥५॥ }\pend
\pstart\versepar \skt{}\edtext{\skt{विषदो}}{\Afootnote{N H; \skt{विशदो}\ A1938}}\skt{ }\edtext{\skt{विपरीतोऽस्य}}{\lemma{\skt{ऽस्य}}\Afootnote{N H; \skt{ऽस्मा}\ A1938}}\skt{ }\edtext{\skt{भेदी}}{\Afootnote{N H; \skt{त्क्लेदा}\ A1938}}\skt{ }\edtext{\skt{शोधनरोपणः}}{\Afootnote{H; \skt{चूषणरोपणः ॥५१७॥}\ A1938; \skt{सोधनरोपनः ॥६ ६॥}\ N}}\skt{ ॥६॥}\pend

\pstart\versepar \ledsidenote{\passref{1.46.518}}\edtext{\edtext{\skt{दाहपाककरस्तीक्ष्णः}}{\lemma{\skt{दाहपाककरस्}}\Afootnote{N H; \skt{दाहपाककर}\ A1938}}}{\lemma{\skt{तीक्ष्णः}}\Afootnote{N H; \skt{स्तीक्ष्णः}\ A1938}}\skt{ }\edtext{\skt{श्रावणो}}{\Afootnote{H; \skt{स्रावणो}\ A1938 N}}\skt{ }\edtext{\edtext{\skt{मृदुरन्यथा}}{\lemma{\skt{मृदुर्}}\Afootnote{N H; \skt{मृदु}\ A1938}}}{\lemma{\skt{अन्यथा}}\Afootnote{N H; \skt{रन्यथा ।}\ A1938}}\skt{ ॥८॥ }\pend
\pstart\versepar \skt{}\edtext{\skt{सान्द्रोपलेपः}}{\Afootnote{N H; \skt{सादोपलेप}\ A1938}}\skt{ }\edtext{\skt{कफकृद्गुरुः}}{\Afootnote{N H; \skt{बलकृद्गुरु}\ A1938}}\skt{ }\edtext{\skt{प्रीणनबृंहणः}}{\Afootnote{N; \skt{स्तर्पणबृंहणः ॥५१८॥}\ A1938; \skt{प्रीणनवृंहणः ॥१०॥}\ H}}\skt{ ॥१०॥}\pend

\pstart \ledsidenote{\passref{1.46.519ab}}\edtext{\edtext{\skt{लघुस्तद्विपरीतः}}{\lemma{\skt{लघुस्}}\Afootnote{N H; \skt{लघु}\ A1938}}}{\lemma{\skt{तद्विपरीतः}}\Afootnote{N H; \skt{स्तद्विपरीतः}\ A1938}}\skt{ स्याल्लेखनो रोपणस्तथा ॥११॥}\pend

\pstart \ledsidenote{\passref{1.46.520cd}}\skt{द्रवः प्रक्लेदनः सान्द्रः }\edtext{\skt{शुष्कः}}{\Afootnote{H; \skt{स्थूलः}\ A1938; \skt{शुश्कः}\ N}}\skt{ }\edtext{\skt{स्याद्द्वन्द्वकारकः}}{\Afootnote{N H; \skt{स्याद्बन्धकारकः ॥५२०॥}\ A1938}}\skt{ ॥१३॥}\pend

\pstart\versepar \ledsidenote{\passref{1.46.521}}\skt{श्लक्ष्णः }\edtext{\skt{पिच्छिलवज्ज्ञेयः}}{\Afootnote{N H; \skt{पिच्छिलवज्ज्ञेयः}\ A1938}}\skt{ }\edtext{\skt{कर्क्कशो}}{\Afootnote{A1938 H; \skt{कर्कसो}\ N}}\skt{ }\edtext{\skt{विषदो}}{\Afootnote{N H; \skt{विशदो}\ A1938}}\skt{ यथा ॥१५॥ }\pend
\pstart\versepar \skt{}\edtext{\skt{सुखानुवर्त्ती}}{\Afootnote{N H; \skt{सुखानुबन्धी}\ A1938}}\skt{ }\edtext{\skt{सूक्ष्मञ्च}}{\Afootnote{H; \skt{सूक्ष्मश्च}\ A1938 N}}\skt{ }\edtext{\skt{सुगन्धीरोचनो}}{\Afootnote{N H; \skt{सुगन्धो रोचनो}\ A1938}}\skt{ }\edtext{\skt{मतः}}{\Afootnote{N H; \skt{मृदुः ॥५२१॥}\ A1938}}\skt{ ॥१६॥}\pend

\pstart\versepar \ledsidenote{\passref{1.46.522}}\skt{दुर्ग्गन्धो विपरीतो }\edtext{\skt{स्य}}{\Afootnote{H; \skt{ऽस्मा}\ A1938; \skt{स्या}\ N}}\skt{ }\edtext{\skt{प्रकाशौ}}{\Afootnote{N H; \skt{द्धृल्लासा}\ A1938}}\skt{ }\edtext{\edtext{\edtext{\skt{चाप्युभावपि}}{\lemma{\skt{चाप्य्}}\Afootnote{N H; \skt{रुचि}\ A1938}}}{\lemma{\skt{उभाव्}}\Afootnote{N H; \skt{कार}\ A1938}}}{\lemma{\skt{अपि}}\Afootnote{N H; \skt{कः ।}\ A1938}}\skt{ ॥१७॥ }\pend
\pstart\versepar \skt{}\edtext{\edtext{\skt{सरोऽनुलोमनः}}{\lemma{\skt{सरो}}\Afootnote{A1938 H; \skt{सषे}\ N}}}{\lemma{\skt{ऽनुलोमनः}}\Afootnote{A1938 H; \skt{नुलोमनोः}\ N}}\skt{ प्रोक्तो मन्दो }\edtext{\skt{यात्राकरः}}{\Afootnote{A1938 H; \skt{यात्राकर}\ N}}\skt{ स्मृतः ॥१९॥}\pend

\pstart \ledsidenote{\passref{1.46.523}}\skt{व्यवायी }\edtext{\edtext{\skt{देहमखिलं}}{\lemma{\skt{देहम्}}\Afootnote{N H; \skt{चाखिलं देहं}\ A1938}}}{\lemma{\skt{अखिलं}}\Afootnote{N H; \emph{om.}\ A1938}}\skt{ व्याप्य पाकाय }\edtext{\skt{कल्पते}}{\Afootnote{N H; \skt{कल्पते । विकासी विकसन्नेवं धातुबन्धान्विमोक्षयेत् ॥५२३॥}\ A1938}}\skt{ ॥२०॥}\pend

\pstart \ledsidenote{\passref{1.46.525}}\skt{गुणा }\edtext{\edtext{\edtext{\skt{विंशतिरित्येते}}{\lemma{\skt{विंशतिर्}}\Afootnote{N H; \skt{विंशति}\ A1938}}}{\lemma{\skt{इत्य्}}\Afootnote{N H; \skt{रित्य्}\ A1938}}}{\lemma{\skt{एते}}\Afootnote{N H; \skt{एवं}\ A1938}}\skt{ }\edtext{\skt{कर्म्मतः}}{\Afootnote{N H; \skt{यथाव}\ A1938}}\skt{ }\edtext{\skt{परिकीर्त्तिताः}}{\Afootnote{N H; \skt{त्परिकीर्तिताः । सम्प्रवक्ष्याम्यतश्चोर्ध्वमाहारगतिनिश्चयम् ॥५२५॥}\ A1938}}\skt{ ॥}\pend

\pstart\versepar \ledsidenote{\passref{1.46.526}}\skt{अत }\edtext{\skt{ऊर्द्धम्प्रवक्ष्यामि}}{\lemma{\skt{ऊर्द्धम्}}\Afootnote{H; \skt{ऊद्धंम्}\ N}}\skt{ आहारगतिनिश्चयं । }\pend
\pstart\versepar \skt{पञ्चभूतात्मके देहे }\edtext{\skt{आहारः}}{\Afootnote{H; \skt{ह्याहारः}\ A1938; \skt{आहार}\ N}}\skt{ }\edtext{\skt{पाकभौतिकः}}{\Afootnote{H; \skt{पाञ्चभौतिकः}\ A1938 N}}\skt{ ॥ }\pend
\pstart\versepar \skt{विपक्वः पञ्चधा }\edtext{\edtext{\skt{सम्यग्गुणांस्तानभिवर्द्धयेत्}}{\lemma{\skt{सम्यग्गुणांस्}}\Afootnote{H; \skt{सम्यग्गुणा}\ A1938; \skt{सम्यग्गुणास्}\ N}}}{\lemma{\skt{तानभिवर्द्धयेत्}}\Afootnote{N H; \skt{न्स्वानभिवर्धयेत् ॥५२६॥}\ A1938}}\skt{ ॥}\edtext{}{\lemma{\skt{… पञ्}}\Afootnote{\emph{om.}\ A1938}}\pend

\pstart\versepar \ledsidenote{\passref{1.46.527}}\edtext{\skt{अविदग्धं}}{\Afootnote{N H; \skt{अविदग्धः}\ A1938}}\skt{ कफं पित्तम्विदग्धः पवनम्पुनः ॥ }\pend
\pstart\versepar \skt{}\edtext{\skt{सम्यग्विपक्षो}}{\Afootnote{N H; \skt{सम्यग्विपक्वो}\ A1938}}\skt{ निःसार आहारः }\edtext{\skt{परिबृंहयेत्}}{\Afootnote{A1938; \skt{परिबृङ्हयेत् ।}\ N; \skt{परिवृंहयेत् ।}\ H}}\skt{ ।}\pend

\pstart\versepar \ledsidenote{\passref{1.46.528}}\edtext{\edtext{\skt{विण्मूत्रमाहारमलः}}{\lemma{\skt{विण्मूत्रम्}}\Afootnote{H; \skt{विण्मूत्र}\ A1938; \skt{विट्मूत्रम्}\ N}}}{\lemma{\skt{आहारमलः}}\Afootnote{N H; \skt{माहारमलः}\ A1938}}\skt{ }\edtext{\skt{सारः}}{\Afootnote{A1938 N\textsuperscript{pc} H; \skt{सार}\ N\textsuperscript{ac}}}\skt{ प्रागीरितो रसः ॥ }\pend
\pstart\versepar \skt{स तु व्यानेन विक्षिप्तः }\edtext{\skt{सर्व्वान्धातून्विसर्प्पति}}{\lemma{\skt{धातून्विसर्प्पति}}\Afootnote{N H; \skt{धातून्प्रतर्पयेत् ॥५२८॥}\ A1938}}\skt{ ।}\pend

\pstart\versepar \ledsidenote{\passref{1.46.529}}\edtext{\skt{कफपित्तमलः}}{\Afootnote{N H; \skt{कफः पित्तं मलः}\ A1938}}\skt{ खेषु स्वेदः स्यान्नखरोम च ॥ }\pend
\pstart\versepar \skt{नेत्रविट्त्वक्षु च }\edtext{\skt{स्नेहो}}{\Afootnote{A1938 H; \skt{स्तेहो}\ N}}\skt{ धातूनां क्रमशो मलाः ।}\pend

\pstart\versepar \ledsidenote{\passref{1.46.530}}\edtext{\skt{दिवाविवुद्धहृदये}}{\Afootnote{H; \skt{दिवा विबुद्धे हृदये}\ A1938; \skt{दिवा विबुद्धहृदये}\ N}}\skt{ जाग्रतः पुण्डरीकवत् ॥ }\pend
\pstart\versepar \skt{अनुपक्लिन्न }\edtext{\edtext{\skt{धात्वन्नमजीर्ण्णेऽपि}}{\lemma{\skt{धात्वन्नम्}}\Afootnote{N\textsuperscript{pc} H; \skt{धातुत्वा}\ A1938; \skt{धात्वईन्नम्}\ N\textsuperscript{ac}}}}{\lemma{\skt{अजीर्ण्णेऽपि}}\Afootnote{\skt{दजीर्णेऽपि}\ A1938; \skt{अजीर्ण्णेपि}\ N H}}\skt{ }\edtext{\skt{हितं}}{\Afootnote{A1938 N\textsuperscript{pc} H; \skt{हित}\ N\textsuperscript{ac}}}\skt{ }\edtext{\skt{निशि}}{\Afootnote{A1938 H; \skt{न् ।}\ N\textsuperscript{ac}; \skt{दिवा ।}\ N\textsuperscript{pc}}}\skt{ ।}\edtext{}{\lemma{\skt{पुण्डरीकवत् ॥ अनुपक्लिन्न}}\Afootnote{N H; \skt{पुण्डरीकवत् । अन्नमक्लिन्न}\ A1938}}\pend

\pstart\versepar \ledsidenote{\passref{1.46.531}}\edtext{\skt{हृदयं}}{\Afootnote{N H; \skt{हृदि सम्}\ A1938}}\skt{ }\edtext{\skt{लीयते}}{\Afootnote{N H; \skt{मीलिते}\ A1938}}\skt{ रात्रौ प्रसुप्तस्य }\edtext{\skt{विशेषतः}}{\Afootnote{A1938 H; \skt{विसेषतः ॥}\ N}}\skt{ ॥ }\pend
\pstart\versepar \skt{}\edtext{\edtext{\skt{समुपक्लिन्नधात्वन्नमजीर्ण्णे}}{\lemma{\skt{समुपक्लिन्नधात्वन्नम्}}\Afootnote{N H; \skt{क्लिन्नविस्रस्तधातुत्वा}\ A1938}}}{\lemma{\skt{अजीर्ण्णे}}\Afootnote{N H; \skt{दजीर्णे}\ A1938}}\skt{ }\edtext{\skt{त्वहितं}}{\Afootnote{N\textsuperscript{pc} H; \skt{न हितं}\ A1938; \skt{त्वहित}\ N\textsuperscript{ac}}}\skt{ }\edtext{\skt{दिवा}}{\Afootnote{A1938 H; \skt{न् ॥}\ N\textsuperscript{ac}; \skt{निशि ॥}\ N\textsuperscript{pc}}}\skt{ ॥}\pend

\pstart\versepar \ledsidenote{\passref{1.46.532}}\edtext{\edtext{\skt{इमम्विधिं}}{\lemma{\skt{इमम्}}\Afootnote{A1938 H; \skt{इयम्}\ N}}\skt{ }\edtext{\edtext{\skt{योऽनुमतम्महामुनेर्}}{\lemma{\skt{ऽनुमतम्}}\Afootnote{A1938 N\textsuperscript{pc} H; \skt{ऽनुमहतम्}\ N\textsuperscript{ac}}}}{\lemma{\skt{महामुनेर्}}\Afootnote{H; \skt{महामुने}\ A1938; \skt{महानुमुनेर्}\ N}}}{\lemma{\skt{विधिं यो}}\Afootnote{A1938 N\textsuperscript{pc} H; \skt{विधिईं यो}\ N\textsuperscript{ac}}}\skt{ }\pend
\pstart\versepar \skt{}\edtext{\skt{म्महर्षिमुख्यस्य}}{\Afootnote{N H; \skt{र्नृपर्षिमुख्यस्य}\ A1938}}\skt{ }\edtext{\skt{पठेत्तु}}{\Afootnote{N H; \skt{पठेद्धि}\ A1938}}\skt{ यत्नतः । }\pend
\pstart\versepar \skt{स }\edtext{\skt{भूमिपालो य}}{\Afootnote{\skt{भूमिपालाय}\ A1938 N; \skt{भूमिपालोय}\ H}}\skt{ }\edtext{\edtext{\skt{विधातुमौषधम्}}{\lemma{\skt{विधातुम्}}\Afootnote{N H; \skt{विधातु}\ A1938}}}{\lemma{\skt{औषधम्}}\Afootnote{N H; \skt{मौषधं}\ A1938}}\skt{ }\pend
\pstart\versepar \skt{महात्मनाञ्चार्हति }\edtext{\skt{वैद्यसत्तम}}{\Afootnote{H; \skt{सूरिसत्तमः ॥५३२॥}\ A1938; \skt{वैद्यसत्तमम्}\ N}}\skt{ }\edtext{\skt{इति}}{\Afootnote{N H; \emph{om.}\ A1938}}\skt{ ॥ }\edtext{\skt{ओ}}{\Afootnote{N H; \emph{om.}\ A1938}}\skt{ ॥}\pend

\pstart\versepar \ledsidenote{\passref{1.46.532.add1}}\edtext{\skt{द्रव्यज्ञानं}}{\Afootnote{N H\textsuperscript{pc}; \emph{om.}\ H\textsuperscript{ac}}}\skt{ }\edtext{\edtext{\skt{रसज्ञानम्वसनञ्च}}{\lemma{\skt{रसज्ञानम्}}\Afootnote{N H\textsuperscript{pc}; \emph{om.}\ H\textsuperscript{ac}}}}{\lemma{\skt{वसनञ्च}}\Afootnote{H\textsuperscript{pc}; \skt{वमनञ्च}\ N; \emph{om.}\ H\textsuperscript{ac}}}\skt{ }\edtext{\skt{विरेचनं}}{\Afootnote{N H\textsuperscript{pc}; \emph{om.}\ H\textsuperscript{ac}}}\skt{ । }\pend
\pstart\versepar \skt{}\edtext{\edtext{\skt{द्रवद्रव्यपरिज्ञानमन्नपानेन}}{\lemma{\skt{द्रवद्रव्यपरिज्ञानम्}}\Afootnote{N H\textsuperscript{pc}; \emph{om.}\ H\textsuperscript{ac}}}}{\lemma{\skt{अन्नपानेन}}\Afootnote{N H\textsuperscript{pc}; \emph{om.}\ H\textsuperscript{ac}}}\skt{ }\edtext{\skt{षट्स्मृतः}}{\Afootnote{H\textsuperscript{pc}; \skt{षट्स्मृतः ॥}\ N; \emph{om.}\ H\textsuperscript{ac}}}\skt{ ॥ }\edtext{\skt{ओ}}{\Afootnote{H\textsuperscript{pc}; \emph{om.}\ N H\textsuperscript{ac}}}\skt{ ॥}\pend

\pstart\versepar \ledsidenote{\passref{1.46.532.add2}}\edtext{\skt{सूत्रस्थाने}}{\Afootnote{N H\textsuperscript{pc}; \emph{om.}\ H\textsuperscript{ac}}}\skt{ }\edtext{\skt{पुरा}}{\Afootnote{N H\textsuperscript{pc}; \emph{om.}\ H\textsuperscript{ac}}}\skt{ }\edtext{\edtext{\skt{प्रोक्तञ्चत्वारिंशत्षडुत्तरं}}{\lemma{\skt{प्रोक्तञ्}}\Afootnote{H\textsuperscript{pc}; \skt{प्रोक्तंञ्}\ N; \emph{om.}\ H\textsuperscript{ac}}}}{\lemma{\skt{चत्वारिंशत्षडुत्तरं}}\Afootnote{N H\textsuperscript{pc}; \emph{om.}\ H\textsuperscript{ac}}}\skt{ । }\pend
\pstart\versepar \skt{}\edtext{\skt{अध्यायाः}}{\Afootnote{N H\textsuperscript{pc}; \emph{om.}\ H\textsuperscript{ac}}}\skt{ }\edtext{\skt{काशिराजेन}}{\Afootnote{H\textsuperscript{pc}; \skt{कासिराजेन}\ N; \emph{om.}\ H\textsuperscript{ac}}}\skt{ }\edtext{\skt{पूर्ण्णं}}{\Afootnote{H\textsuperscript{pc}; \skt{पूर्ण}\ N; \emph{om.}\ H\textsuperscript{ac}}}\skt{ }\edtext{\skt{सर्व्वं}}{\Afootnote{H\textsuperscript{pc}; \skt{न्सर्वआं}\ N\textsuperscript{ac}; \skt{न्सर्वं}\ N\textsuperscript{pc}; \emph{om.}\ H\textsuperscript{ac}}}\skt{ }\edtext{\skt{सविस्तरं}}{\Afootnote{H\textsuperscript{pc}; \skt{सविस्तरं ॥ {\symfont ❈} ॥}\ N; \emph{om.}\ H\textsuperscript{ac}}}\skt{ ॥}\pend

\pstart \ledsidenote{\passref{1.46.trailer}}\edtext{\edtext{\skt{॥ इति}}{\Afootnote{H; \skt{इति भगवता श्रीधन्वन्तरिणोपदिष्टायां तच्छिष्येण महर्षिणा सुश्रुतेन विरचितायां}\ A1938}}\skt{ }\edtext{\skt{सुश्रुते}}{\Afootnote{H; \skt{सुश्रुत}\ A1938}}\skt{ }\edtext{\skt{शल्यतन्त्रे}}{\Afootnote{H; \skt{संहितायां}\ A1938}}\skt{ }\edtext{\skt{सूत्रस्थानं}}{\Afootnote{H; \skt{सूत्रस्थाने}\ A1938; \skt{सूत्रस्थान}\ N}}}{\lemma{\skt{।… सू}}\Afootnote{\emph{om.}\ N}}\skt{ }\edtext{\skt{समाप्तं}}{\Afootnote{H; \skt{षट्चत्वारिंशत्तमोऽध्यायः ॥४६॥}\ A1938; \skt{समाप्तं ॥ {\symfont ❈} ॥}\ N}}\skt{ ॥ ॥}\pend

\endnumbering
